% La Longue Marche — Volume 140-3
% AI Transcription (Gemini 3.1 Pro, medium thinking)
% 696 of 696 pages transcribed successfully
% 114 diagram pages re-transcribed with tikz-cd enhanced prompt
% Rebuilt: 2026-04-17

%% ===== Page 1 =====

UNIVERSITÉ MONTPELLIER

IMAG
INSTITUT MONTPELLIERAIN
ALEXANDER GROTHENDIECK

La "Longue Marche" à travers la théorie de Galois

Deuxième partie : variations sur l'équation des lacets [suite à La "Longue Marche" à travers la théorie de Galois. . ., pages 293 à 650, dont table des matières]*: notes et copies de notes manuscrites (s.d.).

Cote n° 140-3, 695 pages

Alexander GROTHENDIECK

s.d.

*Lacunes entre les pages 371 et 379, la page 414 est manquante (ou erreur de pagination), la page 558 bis semble manquer. La page 518 est une photocopie.

\newpage

%% ===== Page 2 =====

Pages 291-650 + table des matières

NB Les pages 291, 292 figurent deux fois
(erreur de numérotation)

\newpage

%% ===== Page 3 =====

$2^{\text{ème}}$ partie | Variations sur l'équation des lacets |

38 Détermination de $\hat{\Pi}_{0,3} \xrightarrow{\sim} \hat{\mathcal{C}}_{0,3}^+$ 291

39 Propriétés de congruences de l'opération
de $\hat{\Gamma}_{\mathbb{Q}}$ sur $\hat{\Pi}_{0,3}$, liées à l'hom. $N_{\mathbb{Q}} \to \operatorname{Gl}(2,\hat{\mathbb{Z}})$ 299

40 Paraphrases topologico-isotopiques de la
théorie de [unclear: l'équation] 307'

41 Conditions de comm - Action des
$\hat{\Pi}_{0,3} \xrightarrow{\sim} \hat{\mathcal{C}}_{0,3}^+$ aux op. ext. de $\hat{\Gamma}_{\mathbb{Q}}$.
Calcul des autom. ext. de $\hat{\Pi}_{0,3}$ et
de $\hat{\mathcal{C}}_{0,3}^+$. 315

42 Récapitulation sur le groupe $\hat{\mathcal{C}}_{0,3}^+$
$\simeq ...$ et sur les autom. "locaux" de
$\hat{\mathcal{C}}_{0,3}^+$ et de $\hat{\Pi}_{0,3}$ : 331

I $\hat{\mathcal{C}}_{0,3}^+$ 331
II $\hat{\Pi}_{0,3}$ et $\hat{\mathcal{D}}_{0,3}$ 332
III relations avec $\hat{\Pi}_{0,3}'$ et $\widehat{Sl}(2,\mathbb{Z})'$ 334
IV [unclear: scratched text] $\widehat{Sl}(2,\mathbb{Z})'$ 335 bis
V Les autom. de $\hat{\mathcal{C}}_{0,3}$ 336'

[MARGIN: NB les calculs des Sections VII, VIII en partie refaits dans § 43 pour des raisons de notations plus simples...]

VI Autom. "locaux" de $\hat{\Pi}_{0,3}$ et $\hat{\mathcal{C}}_{0,3}^+$ 337'
VII Calculs dans $\widehat{Sl}(2,\mathbb{Z})'$ : cocycles sur $\hat{\Gamma}_{\mathbb{Q}}$ 342
VIII Calculs dans $\widehat{Sl}(2,\mathbb{Z})'$ : cocycles sur $\mathfrak{S}_{\infty}$ 346

43 Digression sur les 1-cocycles des groupes cycl.
(les [unclear: $\beta,\alpha,f,\alpha,\beta,...$] ?) 350

44 Retour sur notations. Mises en équations
dans $\widehat{Sl}(2,\mathbb{Z})$. 356

45 L'hom. $\mathcal{M}_{0,3} \to \widehat{Gl}(2,\hat{\mathbb{Z}})'$, et construction
de l'extension $\tilde{\mathcal{M}}_{0,3}$ de $\hat{\Gamma}$ par $\widehat{Sl}(2,\mathbb{Z})^\wedge$. 384

46 Notations définitives (?) - récapitulation
des calculs sur les autom. "locaux"
de $\hat{\Pi}_{0,3}$ et de $\hat{\mathcal{C}}_{0,3}^+$ 390

\newpage

%% ===== Page 4 =====

I Les groupes $M_{0,3}^{\sim}$ et $\mathcal{N}_{0,3}^{\sim}$   390
II Calcul des éléments de $\mathcal{N}_{0,3}^{\sim}$   392
1) Via $f, g, h$   392
2) Via $\alpha, \beta, f$   393
3) Via $\alpha, \beta, g$   393
4) Via $\xi, \eta, \zeta$   393
5) Congruence entre $\alpha, \beta$, relations dans cocycles   395
III Calculs dans $\operatorname{Gl}(2, \hat{\mathbb{Z}})'$ et cocycles dans $\operatorname{Gl}(2, \hat{\mathbb{Z}})$   395
IV Les trois cas fondamentaux [unclear]
$(\alpha=1, \beta=1, f=1)$, et les relèvements $K_{1/2}, K_{2}, K_{-1}$. Les sous-groupes $M_{0,3}^{\sim}(q)$, $q \in \{ \frac{1}{2}, 2, -1, \frac{1}{3} \}$   402
V Récapitulation en termes de $M_{0,3}^{\sim}[2]$   404
[MARGIN: calculs sur [unclear]] VI Les sous-groupes $M_{0,3}^{\sim}(q)$, $q \in \{2, -1, \frac{1}{2}\}$   407
VII Retour sur l'introduction de $\alpha$ et $\beta$ autour de $\{\xi,\eta\}$.   414

47. Retour sur l'homomorphisme $\Pi_{0,3}^{\sim} \hat{\to} \hat{\mathfrak{S}}_{0,3}^{+}$, les groupes $M_{0,3}^{\sim}, \mathcal{N}_{0,3}^{\sim}, M_{0,3}^{\sim}(2)$   415

48 Construction d'un schéma en groupes   420
I) Étude du morphisme $\alpha \mapsto \alpha(p(\alpha)^{-1} f_1 p)$   420
II) L'équation $\xi = \eta\varepsilon_1(\eta \varepsilon - 1)$, et le groupe $Bo_{\varepsilon_1, p, \varepsilon}$   434
III) L'équation $\alpha p(\alpha) = \beta p(\beta) (\beta \varepsilon p - 1)$,   443
et le groupe $B_{p, \varepsilon}$ de $\hat{\mathfrak{S}}_{0,3}^{+}$
(IV) Un épinglage canonique associé par   445
(V) L'équation $\alpha p(\alpha) = \varepsilon(\beta \varepsilon_1 p(\beta)^a) \varepsilon_2 p(\gamma)$
[crossed out: de l'équation $\alpha p(\alpha) = \beta p(\beta) \varepsilon_1(\beta \varepsilon_2 p(\gamma)-1)$]   454
Le groupe $B_{\underline{\varepsilon}, \sigma}$, la condition $Tr \sigma = 0$   457
(VI) Caractérisation du triplet $(\varepsilon_1, \varepsilon_2, p)$ modulaire   462

VII Récapitulation sur les épinglages.   475
VIII Récapitulation sur l'équation.
Le groupe $SG_{p,\varepsilon}$ et ses variantes   485
IX Les groupes $M_{p,\sigma}$ et $SM_{p,\sigma}$, et les homomorphismes
$SM_{0,3}^{\sim} \to SM_{p,\sigma} \hookleftarrow M_{0,3}^{\sim} \to M_{p,\sigma} (G/2)$   490'
X Structure de $M_{p,\sigma}$, homomorphisme canonique $M_{p,\sigma} \to GM_{p,\sigma} \simeq \operatorname{Gl}(2)$   499
XI Récapitulation sur $M_{p,\sigma}$, les deux systèmes de paramètres $(\varepsilon_1, \varepsilon_2, \varepsilon_3)$ et $(f_1, f_2, f_3)$
[crossed out text]
XII Le cas modulaire.
Les groupes $M_{p}(\sigma)$, $M_{p,\sigma}(a)$, $\Lambda_{p,\sigma}$, $\Lambda_{p,\sigma}^{e}$   504
$AM_{p,\sigma}$, $\mathbb{C}_{p,\sigma}$, $I_{p,\sigma}^{+}$ etc.   507
[MARGIN: depuis des [unclear]] $M_{p,\sigma}$, $M_{p,\sigma}$ calculs de p.ts critiques)   511
6°) Le hom $\Theta : M_{p,\sigma} \to G$, $\bar{\Theta} : M_{p,\sigma} \to P G$   513
5) $R_{p,\sigma}$, $SR_{p,\sigma}$ J   513
[crossed out text] Centres de $M_{p,\sigma}$, $M_{p,\sigma}'$. Le sous-groupe $Z$
$\Sigma \simeq Z \cap E \times Z \cap \Sigma' \times Z \cap \Sigma_{1}$.   517
7°) Diviseurs de $R_{p,\sigma}$ et de $M_{p,\sigma}$   516
8°) Les paramètres $\xi_i = \varepsilon_i K_i p_i \varepsilon_i p_i \varepsilon_i$, et $\underline{\alpha}, \underline{f}, \alpha, \beta, \gamma$, $\underline{\beta_0}$ ($f,g,h$?)   519
9°) Les sous-groupes $M_{p,\sigma}(\alpha), M_{p,\sigma}(\beta)$, et le groupe $SM_{p,\sigma}$ lissé $S\underline{M}_{p,\sigma} \rightleftharpoons \bar{S}\bar{\Sigma}_{p,\sigma}^{\Delta}$   522
10°) Les sous-groupes $\bar{M}_{p,\sigma}(\delta)$, $M_{p,\sigma}(p)$   530'

XII Rétrospective.   542

\newpage

%% ===== Page 5 =====

49. Les groupes $M_{p,\sigma}$ généraux 550

[ I Préliminaires heuristiques ]

II Les groupes $M_{p,\sigma}$ "sans signes" - nouvelles variantes sur la relation des lacets 553

III Passage de $M_{0,3}^\sim$ à $N_{1,1}^\sim$ - et de $\operatorname{Gl}(2,\mathbb{Z})^\wedge/\pm 1$ à $\operatorname{Gl}(2,\mathbb{Z})^\wedge$ 565

[ IV Où on se bat avec "les signes" $\epsilon_2, \epsilon_3$ 568
dans la relation des lacets - calculs -
brouillons ... ] [unclear]

V Les groupes $M_{p,\sigma}$ généraux 577
(contexte pro-fini) 577
1°) Données 577
2°) $Z_{p,\sigma}$, $\Gamma_{p,\sigma}$ 577
3°) $Sol_{p,\sigma}$, $C_{p,\sigma}$ 579
4°) $Z_{\rho}$, $Z_{\sigma}$, $N_{\rho}$, $N_{\sigma}$, $N_{\rho}^\Lambda$, $N_{\sigma}^\Lambda$, $SN_{\rho}^\Lambda$, $SN_{\sigma}^\Lambda$ 580
5°) $SolM_{p,\sigma}$, $M_{p,\sigma}$, $\widetilde{Sol}_{p,\sigma}$, $\widetilde{\Gamma}_{p,\sigma}$ 585
6°) $\Sigma_{p,\sigma}$, $Sg_{p,\sigma}$, $\widetilde{\Sigma}_{p,\sigma}$, $\widetilde{Sg}_{p,\sigma}$ 586
7°) Fonctorialités. 588

[MARGIN: 7°) $M_{p,\sigma}(\rho)$, $M_{p,\sigma}[\sigma]$ (612)
8°) Retour sur le cas universel : $M^{\wedge}_{\rho}$ etc.
9°) Digression sur la flèche $M_{\rho}^\wedge \to M_{\rho}^\wedge[\sigma]$ (613)
10°) Lien avec §47 des groupes 648
rectification de la "relation des lacets" page 622]

8°) Le cas universel $G^\wedge = \widehat{Gl}(2,\mathbb{Z})^\wedge$. Le [unclear: rayé]
sous-groupe $K_{\bullet}$. Étude de l'extension
$Sol_{p,\sigma}$ de $Z_{p,\sigma} = M_{p,\sigma}^\wedge$ par $L_{\rho} \times L_{\sigma}$, et [unclear: rayé]
$S\Gamma_{p,\sigma}^\wedge$ de $Z_{p,\sigma}$ par $\mu = \pm 1$. 590
9°) Récapitulation sur le cas universel. [unclear: rayé]
Structure de $N_0^\wedge$, $N_0^{\Lambda \wedge}$. L'hom. can.
$M \xrightarrow{\sim} M_{p,\sigma} / \mu_{\rho} \hookrightarrow M_{p,\sigma} / \mu_{\rho}$. 614

VI Les groupes $M_{p,\sigma}$ généraux : récapitulation
et raffinement de la
construction "sans lacets". 628
1°) Données (628) 2°) $Z_{p,\sigma}$, $\Gamma_{p,\sigma}$ (629)
3°) $C_{p,\sigma}$, $C_{p,\sigma}^\sim$... (632) 4°) $N_{\rho}$, $N_{\sigma}$, $N_{\rho}^\Lambda$, $N_{\sigma}^\Lambda$... (635)
5°) $M_{p,\sigma}$, $\mu_{p,\sigma}$... (639) 6°) Fonctorialités (640)

\newpage

%% ===== Page 6 =====

38 Détermination explicite 29. (bis) de l'hom $\Pi'_{1,3} \overset{\sim}{\to} \underline{C}_{0,3}^+$
(et de $\Pi'_{0,3}$ ($\simeq$ gr. catégor. plan [unclear: plein] $\overset{\sim}{\to} \underline{C}_{0,3}$ ).
(bi-gradué non orienté)

J'ai envie de revenir sur la prop. du § précédent (p. 262) - une surface $X$, groupe discret $G$ opérant proprement, les stabilisateurs $G_x$ des $x \in X$ conservent l'orientation et opèrent fidèlem. - voisinage de $x$ ... en prenant $Y = X/G$, muni de la "signature" $\underline{d} = d(X/Y)$ définie par la ramification - d'où un iso

(1) $$ \Pi = \Pi_1((X,G), x) \overset{\sim}{\to} \Pi_1((Y,\underline{d}), y) = \Pi' $$
($x \in X$, $y$ image de $x$ dans $Y$)
que je voudrais étudier de plus près.

On a dans $\Pi_1((Y,\underline{d}), y)$, pour tout $b \in \text{supp } \underline{d} = Y'$ des "sous-groupes locaux - tronqués", plus précisément une classe de conjugaison de sous-groupes quotients de $\mathbb{Z}/d_b \mathbb{Z}$, et plus précisém. encore, un hom canonique

(2) $$ T_b \otimes_{\mathbb{Z}} \mathbb{Z}/d_b \mathbb{Z} \to \Pi' $$
déduit par passage au quotient de l'hom nat
(3) $$ T_b \overset{\tau_y}{\to} \Pi_1(Y \setminus Y') \quad (\to \Pi_1((Y,\underline{d}))) $$

(où $T_b$ est le $\mathbb{Z}$-module libre de rang $1$ des "entiers tordus" en $b$ sur $Y$). D'autre part, si $a \in X$ est au dessus de $b$, le stabilisateur $G_a$ est [unclear] un groupe cyclique d'ordre $d_a = d_b$ [unclear] - on a un isom. can. évident

(4) $$ G_a \simeq T_a \otimes_{\mathbb{Z}} \mathbb{Z}/d_b \mathbb{Z} \simeq T_b \otimes \mathbb{Z}/d_b \mathbb{Z} $$
(NB $T_a \simeq T_b$ canoniquement)

\newpage

%% ===== Page 7 =====

D'autre part, à $a$ est associé une classe de
conjugaison de scindages de l'extension

$$(5) \qquad 1 \to \Pi_1(X, x) \to \Pi_1((X, G), x) \to G \to 1 \ ,$$

(cf § 15)

$$(6) \qquad k_a : G_a \to \Pi_1((X, G), x)$$
(à $\Pi_1(X,G)$-conj. près).

Ces classes de scindages (si $X \neq S^2$) corre-
spondent [unclear: biunivoque] aux pts fixes de
$G_a$ opérant sur $X$. Si on prend l'hom.
extérieur défini par (6), (celui-ci [unclear: mod conjugaison]
non seulement par un $\in \Pi_1(X)$, mais
par un $\in \Pi_1(X,G)$ tout entier) on
trouve une corres. 1-1 entre l'ens des hom
ext. (6) qui relèvent l'inclusion
$$G_a \subset G$$
(considérés [unclear: modulo conjugaison] intérieurs), et
l'ens des orbites de $G$ dans $X!$,
donc entre les pts de $Y!$. En résumé -

Proposition. On a une bijection entre $Y!$
et l'ens des hom ext. [MARGIN: $\varphi = \bar{\varphi}_0 : \Gamma \to \Pi$] "maximaux" d'un
ss-groupe $\Gamma \subset G$, (entendu $\Gamma \neq 1$, [MARGIN: $\Gamma \subset G$]) dans
$\Pi_1(X, G)$ qui relèvent l'inclusion $\Gamma \hookrightarrow G$ (les dits
$\Gamma$ sont nécessairement cycliques, les $\Gamma$

\newpage

%% ===== Page 8 =====

292 (bis)

en question corresp. à : un $b \in Y!$ et à
les stabilisateurs $G_a$ dans $G$ des
$a \in X_b = X!_b$. De plus, le composé des
homom.
(7) $$T_b \otimes \mathbb{Z}/n\mathbb{Z} \overset{(4)}{\simeq} G_a \xrightarrow{\varphi_a = \varphi_b} \Pi_1(X, G) \overset{(1)}{\simeq} \Pi_1((Y, \underline{d}))$$
n'est autre (on s'en serait douté !) que "l'hom
lacets tronqué" (2).

Supposons maintenant que $X$ soit de
la forme $\hat{X} \setminus S$, $\hat{X}$ surface compacte,
$S$ partie finie de $\hat{X}$,
(8)
$$ \begin{cases} X = \hat{X} \setminus S \text{ et } Y = \hat{Y} \setminus T \text{ où} \\ \hat{Y} = \hat{X}/G \end{cases} $$

Considérons alors, pour tout $a \in S$, l'hom
extérieur
(9) $$k_a : T_a \longrightarrow \Pi_1(X) \longrightarrow \Pi_1(X, G)$$
que je vais comparer au composé
(10) $$T_b \xrightarrow{k_b} \Pi_1(Y) \longrightarrow \Pi_1(Y, \underline{d})$$
(pour $b = \hat{f}(a)$), moyennant l'isom. can.
(11) $$T_a \simeq T_b$$
Si $d_a$ désigne encore l'indice de ramifica-
tion en $a$, on aura un diagramme
commutatif de compatibilité :

\newpage

%% ===== Page 9 =====

(12)
\begin{tikzcd}
T_a \ar[rr, "K_a"] \ar[d, "\text{[unclear]} \text{ can}"'] & & \pi_1(X) \ar[r] & \pi_1(X, G) \ar[dd, "\wr"] \\
\tilde{T}_b \ar[d, "d_a \cdot id"'] & & & \\
T_b \ar[rrr, "K_b"'] & & & \pi_1(Y, \underline{d})
\end{tikzcd}

Nous allons appliquer ceci à la situation
où $X = U_{0,3}$ , $G = \mathfrak{S}_3$ , $\hat{X} = \mathbb{P}^{1 \text{ an}}$, $\hat{Y} \simeq \mathbb{P}^{1 \text{ an}}$,
où l'application de passage au quotient est
donnée par (p. 288' (58))

(13) $$ f(z) = \frac{3^3}{2^2} \frac{z^2(z-1)^2}{(z^2-z+1)^3} $$
$$ f : \begin{matrix} 0 \\ 1 \end{matrix} \mapsto 0 \qquad \begin{matrix} 1/2 \\ 2 \\ -1 \end{matrix} \mapsto 1 \qquad \begin{matrix} j \\ \bar{j} \end{matrix} \mapsto \infty $$

On voit donc, avec les notations de ,
que les générateurs $\ell'_0, \ell'_1, \ell'_\infty$ de $\Pi'_{0,3}$
(satisfaisant
(14) $$ \ell'_\infty \ell'_1 \ell'_0 = 1 \qquad {\ell'_1}^2 = {\ell'_\infty}^3 = 1 \quad ) $$
correspondent respectivement, par l'isomor-
phisme
(15) $$ \overline{\mathfrak{S}}^+_{0,3} = \Pi_1(U_{0,3}, \mathfrak{S}_3) \stackrel{\sim}{\longrightarrow} \Pi'_{0,3} = \overline{C}^+_{1,1} $$
[unclear: groupe cartographique bi-parti pondéré or.] \quad \quad \quad [unclear: groupe cartographique triangulé orienté $\simeq \operatorname{Sl}(2,\mathbb{Z})$]

aux éléments notés

(16)
$$ \begin{array}{rcccl}
\varepsilon_0, \varepsilon_1, \varepsilon_\infty & & & & \\
\text{(tels que } \varepsilon_i^2 = 1 \text{)} & & \longmapsto & \ell'_0 & \\
& & & & \\
\sigma_0, \sigma_1, \sigma_\infty & & \longmapsto & \ell'_1 & \quad \text{mod conjugaison} \\
& & & & \\
\text{[MARGIN: ou } \rho^{-1} \text{ ? ]} \quad \quad \rho & & \longmapsto & \ell'_\infty &
\end{matrix} $$

\newpage

%% ===== Page 10 =====

293

[MARGIN: de $\pi = \overline{C}_{0,3}^+$]
Ceci montre que l'on peut trouver dans 
un conjugué $\varepsilon_0'$ de $\varepsilon_0$, $\sigma_0'$ de $\sigma_0$, tels qu'on ait

(1) $\rho \sigma_0' \varepsilon_0' = 1$, en plus de $\rho^3 = {\sigma_0'}^2 = 1$

et de telle façon que ces relations constituent
une présentation de $\pi$ (donc en particulier,
que $\varepsilon_0', \sigma_0', \rho$ engendrent $\pi$). S'inspirant d'
ailleurs de l'isomorphisme plus gros qui
prolonge (15) et s'obtient :

(16) $$\overline{C}_{0,3} \simeq \pi_1 ( U_{0,3}, \underset{\overline{C}_{0,3}}{\underbrace{\mathfrak{S}_3 \times \mathbb{Z}/2}} ) \longrightarrow \overline{C}_{1,1}'$$
[MARGIN: groupe cartographique triangulaire étendu $\simeq \operatorname{Gl}(2, \mathbb{Z})$]

on devrait trouver dans le premier
trois générateurs $\tau_0', \tau_1', \tau_\infty'$ (correspondant aux
gén. analogues $\tau_0, \tau_1, \tau_\infty$ du groupe cartogra-
phique correspondant) satisfaisant

(17) $${\tau_0'}^2 = {\tau_1'}^2 = {\tau_\infty'}^2 = 1, \quad (\underbrace{\tau_0' \tau_1'}_{\text{corr. à } l_1'})^2 = (\underbrace{\tau_1' \tau_\infty'}_{\text{corr. à } l_\infty'})^3 = 1$$

de telle façon que ce soient là des
relations de présentation de $\overline{C}_{0,3}$. Il
faudrait que $\tau_\infty'$ soit conjugué de $\sigma_0$.

(18) $$
\begin{matrix}
\tau_\infty' \tau_1' = \varepsilon_0' & ; & \tau_0' \tau_\infty' = \sigma_0' & , & \tau_1' \tau_0' = \rho \\
\downarrow & & \downarrow & & \downarrow \\
l_0' & & l_1' & & l_\infty'
\end{matrix}
$$

\newpage

%% ===== Page 11 =====

$X$ surface topologique

$\mathcal{G}$ groupe discret opérant sur $X$

$G$ sous-groupe invariant dans $\mathcal{G}$, opérant
proprement sur $X$, de sorte que $\forall x \in X$
les stabilisateurs $G_x$ ($x \in X$) opèrent fidèle-
ment ([unclear: action] au v. de $x$, qui conserve
l'orientation locale en $x$)

$Y = X/G$ (surface topologique), $f: X \to Y$

$X'$ partie discrète de $X$, stable sous $\mathcal{G}$,
$d_X$ "signature" sur $X$ de support $X'$, inv. par $\mathcal{G}$
$d_X : X' \to \mathbb{N}$

$Y' = X'/G \subset Y$ image de $X'$ dans $Y$

On suppose $x \in X, G_x \neq \{1\} \implies x \in X'$

$d_Y : Y' \to \mathbb{N}$, $d_Y(y) = \nu_x d_X(x)$
si $y = f(x)$, $\nu_x = [G_x]$

$Rev(X, d_X, \mathcal{G})$ catégorie des $\mathcal{G}$-revêtements
ramifiés de $X$, subordonnés à $d_X$

$Rev(Y, d_Y, \mathcal{H})$ catégorie des $\mathcal{H}$-revêtements
ramifiés de $Y$, subordonnés à $d_Y$

Pour $X' \in Ob Rev(X, d_X, \mathcal{G})$, considérons
$X'/G$ comme surface sur laquelle opère
$\mathcal{H} = \mathcal{G}/G$, au dessus de $Y = X/G$. Alors

\newpage

%% ===== Page 12 =====

a) $X'/G \in \operatorname{Ob}\ Rev(Y, d_Y, \mathcal{H})$

b) on trouve ainsi un foncteur

(19) $Rev(X, d_X, \mathcal{G}) \to Rev(Y, d_Y, \mathcal{H})$

c) Ce foncteur est une équivalence de
catégories.

d) Description d'un foncteur quasi-inverse
par image inverse "normalisée".

Remarques. OPS $d_X$ à valeurs $>0$ i.e. $d_Y$ à valeurs $>0$
(sinon, si $S = \{x \in X' \mid d_X(x) = 0\}$, remplacer
$X$ par $X \setminus S$, $Y$ par $Y \setminus T$ ($T=S/G$)).
OPS $d_X(x) = 1 \implies [G_x] \neq 1$ [unclear: (sinon, enlever x de X! ...)]
Le cas le plus intéressant peut-être est
celui où $d_X(x) = 1 \quad \forall x \in X'$, où $G_x \neq \{1\}$
(donc $X' = \{x \in X \mid G_x \neq \{1\}\}$), dans
ce cas on a
(20) $Rev(X, d_X, \mathcal{G}) = Rev(X, \mathcal{G})$
et le $Rev(Y, d_Y, \mathcal{H})$ permet de
s'interpréter comme un $Rev(X, \mathcal{G})$
(sans donnée de ramification)

\newpage

%% ===== Page 13 =====

Corollaire Soit $x \in X \setminus X'$, $y$ son image dans $Y \setminus Y'$.
on a un isomorphisme canonique

(21) $$ \Pi_1((X, d_X, G), x) \xrightarrow{\sim} \Pi_1((Y, d_Y, \mathcal{H}), y) $$

Nous allons expliciter cet isomorphisme
en termes des classes de chemins donnés.
Un élément $\alpha$ de $\Pi_1((X, d_X, G), x)$ provient d'un
couple $(\gamma, l)$, où
$\gamma \in G$, $l$ classe de chemins dans
$X \setminus X'$ de $x$ vers $\gamma^{-1} x$.

(avec $(\gamma, l)$ et $(\gamma', l')$ définissent le même
élément
ssi $\gamma' = \gamma$, et $l'^{-1} l$ dans $\Pi_1(X \setminus X', x)$ est
dans le s-groupe invariant noyau de
$\Pi_1(X \setminus X', x) \to \Pi_1((X, d_X), x)$ - ie le s-groupe
normal engendré par les $k \lambda_a^{d_X(a)} k^{-1}$,
pour $a \in X'$, $\lambda_a$ désigne la cl dans
$\Pi_1(X \setminus X', x)$ d'un "lacet en $x$ autour de $a$")

Soit alors $j$ l'image de $\gamma$ dans $\mathcal{H}$ -
& soit $f_U : X \setminus X' \overset{dfn}{=} U \to Y \setminus Y' \overset{dfn}{=} V \simeq U/G$ l'application.

\newpage

%% ===== Page 14 =====

canonique, d'où l'im. par $f_*$ de $\alpha \in \Pi_1(X, d_X, \mathcal{G}; x)$ défini par $(\gamma, l)$ dans $\Pi_1(Y, d_Y, \mathcal{H}; y)$ est $(\dot{\gamma}, f_u(l))$

(NB $f_u(l)$ est une classe de chemins dans $V = Y \setminus Y!$ de $y = f_u(x)$ vers $\dot{\gamma}^{-1}(y) = f_u(\gamma^{-1}x)$).

[MARGIN: NB.] On voit tout de suite que si $(\gamma, l)$ et $(\gamma', l')$ définissent le même $\alpha$, alors ils définissent le même $\beta$, ce qui provient du fait que 

$$\Pi_1(f_u) : \Pi_1(X \setminus X!, x) \longrightarrow \Pi_1(Y \setminus Y!, y)$$

applique un $\lambda_a$ dans $\lambda_b^{\nu(a)}$

($b=f(a)$, $\nu(a) = [G_a] = \text{degré de ramification de } f \text{ en } a$)

donc $\lambda_a^{d_X(a)}$ dans $\lambda_b^{d_X(a) \nu(a) = d_Y(b)}$

Revenons à la situation de (13) :

$X = U_{0,3}$, $G = \widetilde{\mathcal{T}}_{0,3} = \widetilde{\mathfrak{S}}_3 \times \mathbb{Z}/2$, $\mathcal{H} = \mathbb{Z}/2$, $X! = \{ \infty, 2, -1, 1/2, \dots \}$

$Y \simeq \mathbb{P}^1 \setminus \{0\}$, $Y! = \{1, \infty\}$, (l'application

provenant au quotient de $\hat{X} = \mathbb{P}^1$ sur $\hat{Y} = \mathbb{P}^1$

donné par (13) :

$$f(z) = \frac{3^3}{2^2} \frac{z^2(z-1)^2}{(z^2 - z + 1)^3}$$

\newpage

%% ===== Page 15 =====

Pour calculer $\Pi_1(Y, d_Y, \mathbb{Z}/2\mathbb{Z}) = \tilde{\Pi}_{0,3}$ (groupe carto-
graphique tri-gentil non orienté), on a pris
le point base $y = -\vec{j}$ , et on a choisi
un $x \in U_{0,3}$ au dessus de $y$, tel que
$f(x) = y = -\vec{j}$
- il y a [unclear] choix d'en [unclear] un $f^{-1}(\{y\})$ à six
éléments, dont trois dans l'hémisphère nord
(les trois autres dans l'hémisphère sud),
ceux-ci étant permutés transitivement par $\rho$.
Il y a de même un relèvement dans les
tri-gles sphériques $(-\vec{j}, 0, 1/2)$, $(-\vec{j}, 1, 2)$,
$(-\vec{j}, \infty, -1)$ (appliqués par $f$ homéomorphiquement
sur le tri-gle sphérique $(0,1,\infty)$ supérieur,
[unclear: (qui est l'hémisphère sup)]
de centre $\vec{j}$ ). Le choix de $x$ revient à un
choix de l'un de ces trois tri-gles sphériques, ou encore un des 3
choix de l'un des trois segments sphériques
$(0, 1/2)$, $(1, 2)$, $(\infty, -1)$ - j'ai envie de prendre
$(0, 1/2)$, correspondant au tri-gle $(-\vec{j}, 0, 1/2)$.

Les trois générateurs standard $\tau'_0, \tau'_1, \tau'_\infty$ de $\Pi_1(Y, d_Y, \mathbb{Z}/2\mathbb{Z})$
au dessus de $\tau \in \mathbb{Z}/2\mathbb{Z}$
s'obtiennent en prenant des chemins $l'_0, l'_1, l'_\infty$
[unclear]
de $y = -\vec{j}$ vers $\tau y = \vec{j}$ , qui sont les arcs du
grand cercle coupant l'équateur resp. en $2, -1, 1/2$.

\newpage

%% ===== Page 16 =====

Ces chemins sont les im. géo par $f$ de chemins
sur $X$, soient $\lambda_0 : x \overset{\lambda_0}{\to} x_0$, $\lambda_1 : x \overset{\lambda_1}{\to} x_1$, $\lambda_\infty : x \overset{\lambda_\infty}{\to} x_\infty$, qui coupent
la subdivision barycentrique des doubles
triangles sphériques $(0, 1, \infty)$ en exactement un point,
situés respectivement sur le côté $(1/2, \bar{j})$, $(\bar{j}, 0)$
$(0, 1/2)$ des triangles sphériques $(0, 1/2, \bar{j})$.

[MARGIN: NB Les points $x_0, x_1, x_\infty$
(outre $x$)
sont certains points
au dessus de $y = -j$
du bord du [unclear]
sphérique [unclear]
situés dans les
triangles sphériques
$(-j, 1/2, 1)$,
$(\infty, -1, -j)$,
$(-1, 0, -j)$]

points $x_0, x_1, x_\infty$ [unclear] au dessus de $y$
i.e. $g_0 x = x_0, g_1 x = x_1, g_\infty x = x_\infty$
avec $g_0 = \sigma_0 \tau, g_1 = \sigma_1 \tau, g_\infty = \sigma_\infty \tau$.
Ainsi les générateurs involutifs
$\tau'_0, \tau'_1, \tau'_\infty$ de $\tilde{\Pi}_{0, \bar{S}} = \Pi_1(Y, \Delta_Y, \mathbb{Z}/2\mathbb{Z}; y)$
sont les images des éléments [unclear: (notés de même)]
(notés par le même signe)
$$\tau'_0 = (\sigma_0 \tau, \lambda_0) ; \tau'_1 = (\sigma_1 \tau, \lambda_1) , \tau'_\infty = (\sigma_\infty \tau, \lambda_\infty)$$
de $\Pi_1(X, \mathcal{G}; x)$. Nous allons utiliser [unclear]
de grand cercle pour [unclear: point base] $x = -j = P$ (origine
habituelle pour le calcul de $\Pi_1(X, \mathcal{G})$)
et nous identifierons $\Pi_1(X, \mathcal{G}; x)$ à $\Pi_1(X, \mathcal{G}; P)$.

On peut alors prendre

(21) $\tau'_0 = (\sigma_0 \tau, \lambda'_0)$ , $\tau'_1 = (\sigma_1 \tau, \lambda'_1)$ , $\tau'_\infty = (\sigma_\infty \tau, \lambda'_\infty)$

où $\lambda'_0, \lambda'_1, \lambda'_\infty$ sont les chemins de
de $P = -j$ vers $x_0, x_1, x_\infty$ [unclear: situés dans les bords du] [unclear: quadrangle].

\newpage

%% ===== Page 17 =====

[unclear: tien des demi-cercles en arêtes] pour le graphe barycentrique des doubles triangles sphériques à [unclear: sommet] $P=-j$, [unclear] les points, [unclear: situés] sur les segments sphériques $(0, 1/2)$ ! Les points $z_0, z_1, z_\infty$ sont construits respectivement comme les points $z$ tels que $f(z)=-j$ et qui se trouvent sur les triangles sphériques

$$([unclear: 1/2], j, 1), (-j, 0, -1), (0, 1/2, -j).$$

donnés par

$$\lambda'_i = [unclear](\lambda_i) l$$

On a donc

$$\lambda'_0 = (\sigma_0 \tau)(l') \cdot \lambda_0 \cdot l \quad \text{[MARGIN: homotope par un lacet trivial dans } U_{0s} \text{]}$$
[unclear: symétrie par rapport équateur] $(\infty, -j, 1/2, -j)$

$$\lambda'_1 = \sigma_1 \tau(l') \cdot \lambda_1 \cdot l \quad \text{[MARGIN: homotope en lacet trivial dans } U_{0s} \text{]}$$
[unclear: symétrie par l'équateur] $(0, -j, 1, -j)$

$$\lambda'_\infty = (\tau(l')) \cdot \lambda_\infty \cdot l \quad \text{[MARGIN: homotope au chemin du grand cercle } (-j, 1/2, j) \text{ de } P=\bar{j} \text{ vers } \tau P = -j \text{]}$$

\newpage

%% ===== Page 18 =====

297

[C] (42)

En résumant, avec les notations du
§ 7 f), on aura donc

$$ (22) \quad \left\{ \begin{array}{l} \tau'_0 = \tilde{\sigma}_{\infty} \quad (\text{sur } \dot{\sigma}_0 i) \\ \tau'_1 = \tilde{\sigma}_0 \quad (\text{sur } \dot{\sigma}_1 i) \\ \tau'_{\infty} = \tau_{\infty} \quad (\text{sur } \dot{\tau}) \end{array} \right. $$

[MARGIN: NB $\tau_\infty = \tilde{\sigma}_{\infty} \tilde{\sigma}_0 = \tilde{\sigma}_0 \tilde{\sigma}_{\infty}$]

Il faut vérifier les relations essentielles

$$ (23) \quad {\tau'_0}^2 = {\tau'_1}^2 = {\tau'_{\infty}}^2 = (\underbrace{\tau'_{\infty} \tau'_0}_{\tilde{\sigma}_{\infty}})^2 = (\underbrace{\tau'_1 \tau'_0}_{\rho^{-1}})^3 = 1 $$

les quatre premières sont claires, la dernière
s'écrit

$$ (\tilde{\sigma}_0 \tilde{\sigma}_{\infty} \rho)^3 = 1 $$
$$ \tilde{\sigma}_0 \tilde{\sigma}_{\infty} = \rho^{-1} \quad ([C], \text{ formule } 13_1) $$

d'où ce qu'on veut.

Calculons (conformément à (18)) $\tau'_{\infty} \tau'_1$, $\tau'_1 \tau'_0$, $\tau'_0 \tau'_{\infty}$
qui devraient correspondre à $l'_0, l'_1, l'_{\infty}$ :
$$ \underbrace{\tau'_1 \tau'_0}_{\rho^{-1}} \quad $$
$$ l'_0 = \tau'_{\infty} \tau'_1 = \tau_{\infty} \tilde{\sigma}_0 = \tilde{\sigma}_0 \underbrace{\tilde{\sigma}_{\infty} \tilde{\sigma}_0}_{\rho} = \tilde{\sigma}_0 \rho $$
$$ l'_1 = \tau'_0 \tau'_{\infty} = \tilde{\sigma}_{\infty} $$
$$ l'_{\infty} = \tau'_1 \tau'_0 = \rho^{-1} $$

i.e.

$$ (24) \quad l'_0 = \tilde{\sigma}_0 \rho, \quad l'_1 = \tilde{\sigma}_{\infty}, \quad l'_{\infty} = \rho^{-1} $$

avec les relations caractéristiques

$$ (25) \quad l'_{\infty} l'_1 l'_0 = 1, \quad {l'_1}^2 = 1, \quad {l'_{\infty}}^3 = 1 $$

Tout est OK !

[MARGIN:
On exprime les anciens générateurs $\rho, \sigma_0, \sigma_1, \sigma_{\infty}, l_0, l_1, l_{\infty}, \tau_0, \tau_1, \tau_{\infty}$ de $M_{0,3} ( \dots )$ en fct de ces $\tau'_0, \tau'_1, \tau'_{\infty}$ ainsi:
$\rho = \tau'_0 \tau'_1$
$l_{\infty} = \tau'_1 \tau'_0, \quad l_0 = \rho \tau'_{\infty} \rho^{-1} = \tau'_0 \tau'_1 \tau'_{\infty} \tau'_1 \tau'_0$
$l_1 = \tau'_0 \tau'_{\infty} \tau'_0$
$\tau_{\infty} = \tau'_{\infty} \dots$
$\dots$ [unclear]
]

\newpage

%% ===== Page 19 =====

Je voudrais encore expliciter la relation
de $\Pi_{0,3} \simeq \widehat{\mathfrak{S}}_{0,3}$ avec $GP(1, \mathbb{Z}) \simeq \operatorname{Gl}(2, \mathbb{Z})'$
[unclear: qu. prov. auto. graphes n. orientés]

$\simeq \operatorname{Gl}(2, \mathbb{Z})/\{\pm 1\}$. Rappelons d'abord l'opération
de
(26) $$GP(1, \mathbb{R}) \simeq \operatorname{Gl}(2, \mathbb{R})/\mathbb{R}^* \simeq \operatorname{Gl}(2, \mathbb{R})^{\pm} / \{\pm 1\}$$
[MARGIN: NB $GP(1, \mathbb{R})$ a pour sous-groupe connexe d'ind. 2 l'image de $\operatorname{Sl}(2, \mathbb{R})$ (cad est isom. à $\operatorname{Gl}(2, \mathbb{R})^+/\mathbb{R}^{*+}$ ou $\operatorname{Sl}(2, \mathbb{R})/\{\pm 1\}$]
sur le demi-plan de Poincaré $D$
(27) $$D = \{ z \in \mathbb{C} \mid \\operatorname{Im}(z) > 0 \} = \{ x+iy \mid x \in \mathbb{R}, y \in \mathbb{R}^{*+} \}$$
fait opérer la matrice $u = \begin{pmatrix} a & b \\ c & d \end{pmatrix}$ par
(28) $$f_u(z) = \begin{cases} \frac{az+b}{cz+d} & \text{si } \det u > 0 \\ \frac{a\bar{z}+b}{c\bar{z}+d} & \text{si } \det u < 0 \end{cases}$$

De cette façon, $GP(1, \mathbb{R})$ est isomorphe au groupe
des automorphismes conformes de $D$.

[MARGIN: équivalences de catégories]
D'ailleurs, le donné d'une courbe elliptique
$E$ sur $\mathbb{C}$ équivaut à celle d'un espace
vectoriel $V$ de dim $1$ sur $\mathbb{C}$, et d'un
"réseau" $\pi \subset V$ (sous-groupe discret de rang $2$).
L'orientation canonique de $V$ correspond à
un générateur de $\bigwedge_{\mathbb{Z}}^2 \pi \simeq \mathbb{Z}$, c-à-d. à
une "orientation". Le choix d'une rigidification

\newpage

%% ===== Page 20 =====

de [unclear: forme] uni-modulaire de $\pi$, équivaut à 
(comme $\mathbb{Z}$-module)
celui d'une base $(e_1, e_2)$ de $\pi$ telle que
(comme base de $V$ sur $\mathbb{C}$)
$e_1 \wedge e_2 = \eta$. L'élément $e_1$ permet d'ailleurs
d'identifier $V$ à $\mathbb{C}$, $e_1$ devenant $1$. Ainsi
la donnée de $E$ muni de rigidifica-
tion [unclear: de forme] uni-modulaire, équivaut à celle
d'un élément $e_2 = \tau$ de $\mathbb{C}$, tel que

1) $(1, \tau)$ engendre un sous-groupe discret
de rang $2$, i.e. $\tau \notin \mathbb{R}$ i.e. $\Im \tau \neq 0$

2) l'orientation induite sur ce $\mathbb{R}$-e.v. défini par
la base $(1, \tau)$ est "directe", i.e. $\Im \tau > 0$,

i.e. $\tau \in D$.

Quand on regarde l'opération de $\operatorname{Sl}(2, \mathbb{Z})$
sur les classes d'isomorphie de courbes ell.
rigidifiées, s'identifiant à $D$, on trouve
que $\begin{pmatrix} a & b \\ c & d \end{pmatrix}$ opère par $\tau \mapsto \frac{a \tau + b}{c \tau + d}$, ce
qui est l'opération habituelle de
$\operatorname{Sl}(2, \mathbb{Z})$ sur $D$, à travers $\operatorname{Sl}(2, \mathbb{Z}) \simeq GP(1, \mathbb{Z})^+$.

Ce groupe opère proprement, il a exacte-
ment deux trajectoires critiques, qui corres-

\newpage

%% ===== Page 21 =====

pondant respectivement aux $\tau$ qui, avec
$1$ engendrent le $\mathbb{Z}$-réseau des entiers de
$\mathbb{Q}(i)$ (soit $\mathbb{Z}+i\mathbb{Z}$), et le
réseau des entiers de $\mathbb{Q}(j)$, formé des
nombres de la forme $a + bj$ ($j = \frac{-1 + i\sqrt{3}}{2}$
$= \exp \frac{2 i \pi}{3}$) ; la ramification correspondante étant d'ordre
$2$ et $3$.

On a un isom. canonique (analytique p. ex.)
(29) $$D / \operatorname{Gl}(2, \mathbb{Z})^+ \simeq \mathbb{P}^1_{\mathbb{C}} \setminus \{ \infty \} = Y,$$
caractérisé par la condition que l'orbite
critique de ramification $3$ s'envoie sur le
point $0$, et celle de ramification $2$
point $1$. On trouve ainsi un isom.
(posant $g = \operatorname{Gl}(2, \mathbb{Z})$, $G = \operatorname{Gl}(2, \mathbb{Z})^+ = \operatorname{Sl}(2, \mathbb{Z})$)

(30) $$\pi_1(D, G, x) \xrightarrow{\sim} \pi_1(Y, y)$$

[MARGIN: $( x \in D \setminus D' )$]

$= G$

que je voudrais expliciter, en identifiant
les $\tau_0, \tau_1, \tau_\infty$ comme éléments de $\operatorname{Gl}(2, \mathbb{Z})^+$.

\newpage

%% ===== Page 22 =====

39 Les conditions de congruence pour l'opération extérieure de $\Gamma_{\mathbb{Q}}$ sur $\hat{\Pi}_{0,3}$.

Considérons le diagramme

(1)
\begin{tikzcd}
1 \ar[r] & \underset{\simeq \operatorname{Sl}(2, \mathbb{Z})^\wedge}{\hat{\mathcal{T}}_{1,1}^+} \ar[r] \ar[d] & \mathcal{N}_{1,1} \ar[r] \ar[d] & \Gamma_{\mathbb{Q}} \ar[r] \ar[d, "\chi"] & 1 \\
1 \ar[r] & \operatorname{Sl}(2, \hat{\mathbb{Z}}) \ar[r] & \operatorname{Gl}(2, \hat{\mathbb{Z}}) \ar[r, "det"] & \hat{\mathbb{Z}}^* \ar[r] & 1
\end{tikzcd}

Il montre que l'opération extérieure de $\Gamma_{\mathbb{Q}}$ sur $\hat{\mathcal{T}}_{1,1}^+ \simeq \operatorname{Sl}(2, \mathbb{Z})^\wedge$ passe au quotient en une opération extérieure sur $\operatorname{Sl}(2, \hat{\mathbb{Z}})$, et cette opération extérieure, de plus, peut s'expliciter entièrement à l'aide du caractère cyclotomique

(2)
$$ \chi : \Gamma_{\mathbb{Q}} \longrightarrow \hat{\mathbb{Z}}^* $$

Pour préciser ce dernier point, notons que la deuxième ligne de (1) donne un homomorphisme

(3)
$$ \hat{\mathbb{Z}}^* \longrightarrow Aut ext (\operatorname{Sl}(2, \hat{\mathbb{Z}})) , $$

et l'hom canonique

(4)
$$ \Gamma_{\mathbb{Q}} \longrightarrow Aut ext (\operatorname{Sl}(2, \hat{\mathbb{Z}})) $$

\newpage

%% ===== Page 23 =====

n'est autre que le composé de (2) et (3).

Par ailleurs, l'hom. (3) se factorise

en

(5) $$ \hat{\mathbb{Z}}^* / \hat{\mathbb{Z}}^{*2} \longrightarrow Autext(\operatorname{Sl}(2,\hat{\mathbb{Z}})) $$

(provient du fait que les éléments de
$\operatorname{Gl}(2,\hat{\mathbb{Z}})$ qui sont dans le centre
$\hat{\mathbb{Z}}^*$ opèrent trivialement sur $\operatorname{Sl}(2,\hat{\mathbb{Z}})$,
et que le composé
$$ \hat{\mathbb{Z}}^* (\text{centre}) \subset \operatorname{Gl}(2,\hat{\mathbb{Z}}) \xrightarrow{det} \hat{\mathbb{Z}}^* $$
n'est autre que $\lambda \mapsto \lambda^2$ ). Plus
précisément, on voit que pour que l'
autom. de $\operatorname{Sl}(2,\hat{\mathbb{Z}})$ induit par un
$u \in \operatorname{Gl}(2,\hat{\mathbb{Z}})$ soit intérieur, i.e. de la forme
$int(v)$ pour un $v \in \operatorname{Sl}(2,\hat{\mathbb{Z}})$, il f. et. s.
qu'il existe $v_0 \in \operatorname{Sl}(2,\hat{\mathbb{Z}})$ tel que
$u v_0^{-1}$ commute à $\operatorname{Sl}(2,\hat{\mathbb{Z}})$, et cela signifie
que c'est un scalaire - i.e.
$u \in \operatorname{Sl}(2,\hat{\mathbb{Z}}) . \hat{\mathbb{Z}}^*$. En d'autres termes
l'hom (5) est injectif, et on comprend

\newpage

%% ===== Page 24 =====

300

le noyau de (4) n'est autre que le
noyau du composé
$$ (6) \qquad \Gamma_Q \xrightarrow{\chi} \hat{\mathbb{Z}}^* \rightarrow \hat{\mathbb{Z}}^* / \hat{\mathbb{Z}}^{* 2} $$
i.e. (par la théorie des corps de classes pour
$\mathbb{Q}$) l'opération ext. de $\Gamma_Q$ sur $\operatorname{Sl}(2, \hat{\mathbb{Z}})$ se
factorise par une opération de $(\Gamma_{Q,ab})_2 =$ le
[MARGIN: NB l'op. ext. de $\Gamma_Q$ sur $\operatorname{Sl}(2, \hat{\mathbb{Z}})$ se fait par l'intermédiaire de $\hat{\mathbb{Z}}^* / \hat{\mathbb{Z}}^{* 2}$]
quotient "multiquadratique" de $\Gamma_Q$ - et cette
dernière est fidèle.

Ainsi, si on identifie $\Gamma_Q$ (par abus)
à un sous-groupe de $Autext\ \operatorname{Sl}(2, \hat{\mathbb{Z}})$,
ce sous-groupe $\Gamma$ est astreint à des
conditions assez fortes : conserver invariant
le sous-groupe distingué noyau de
$$ (7) \qquad \operatorname{Sl}(2, \mathbb{Z})^\wedge \rightarrow \operatorname{Sl}(2, \hat{\mathbb{Z}}) $$
et induire sur $\operatorname{Sl}(2, \hat{\mathbb{Z}})$ l'opération extérieure
provenant d'un hom. $\Gamma \rightarrow \hat{\mathbb{Z}}^* / \hat{\mathbb{Z}}^{* 2}$.

On peut d'ailleurs préciser cet hom.
ainsi. Considérons dans $\operatorname{Sl}(2, \hat{\mathbb{Z}})$ le
"sous-groupes-lacets" provenant des points :

\newpage

%% ===== Page 25 =====

l'infini de $M_{1,1\mathbb{Q}}$, il correspond
: une classe d'hom. ext.

(8) $$T = \Gamma_{\infty}(\bar{\mathbb{Q}}) \longrightarrow S\ell(2,\mathbb{Z})^\wedge$$
(ou $S\ell(2,\hat{\mathbb{Z}})$)

dont l'image [unclear] exact, est (à conj. près)
un sous-groupe "unipotent"

(9) $$\begin{pmatrix} 1 & 0 \\ n & 1 \end{pmatrix} \quad , \quad n \in \hat{\mathbb{Z}}$$

Cet hom. extérieur est respecté par
l'opération de $\Gamma$, via une opération
bien déterminée de $\Gamma$ sur $T$

(10) $$\Gamma \xrightarrow{\chi} \operatorname{Aut}(T) \simeq \hat{\mathbb{Z}}^\times$$

qui sera justement le caractère cyclo-
tomique. De façon plus précise, pour
un automorphisme extérieur de
$S\ell(2,\mathbb{Z})^\wedge$ (ou seulement de $S\ell(2,\hat{\mathbb{Z}})$)
qui invarie la classe de "conjugaison"
de l'image de (8) [ou seulement de
[MARGIN: [unclear]] l'image de $T \to S\ell(2,\hat{\mathbb{Z}})$], il existe un
$\chi(u) \in \operatorname{Aut}(T) \simeq \hat{\mathbb{Z}}^\times$ unique, tel que l'on ait
commutativité dans

\newpage

%% ===== Page 26 =====

(11)
\begin{tikzcd}
T \ar[r, "k"] \ar[d, "\chi(u)"'] & \operatorname{Sl}(2, \mathbb{Z})^\wedge \ar[d, "u"] \\
T \ar[r] & \operatorname{Sl}(2, \mathbb{Z})^\wedge
\end{tikzcd}

- ou mieux, seulement dans

(11)
\begin{tikzcd}
T \ar[r, "k"] \ar[d, "\chi(u)"'] & \operatorname{Sl}(2, \hat{\mathbb{Z}}) \ar[d, "u"] \\
T \ar[r] & \operatorname{Sl}(2, \hat{\mathbb{Z}})
\end{tikzcd}

- et l'on a une action $\chi$ (10) sur le
groupe de tous les autom. ext. $u$ de
$\operatorname{Sl}(2, \hat{\mathbb{Z}})$ qui ont la propriété envisagée.

Ceci dit, $\Gamma$ est contenu dans ce
sous-groupe $\Gamma'$, et de plus dans le sous-
-groupe des autom. ext. qui de plus
préservent le noyau de (7), et qui
opèrent (dans le quotient $\operatorname{Sl}(2, \hat{\mathbb{Z}})$)
suivant le composé
$$ \Gamma \xrightarrow{\chi} \hat{\mathbb{Z}}^* \to \hat{\mathbb{Z}}^*/\hat{\mathbb{Z}}^{*2} \to Aut ext_{loc!} (\operatorname{Sl}(2, \hat{\mathbb{Z}})) $$

Car d'ailleurs notons que le centre $\{\pm 1\}$
de $\operatorname{Sl}(2, \mathbb{Z})^\wedge$ n'étant pas trivial, l'op. ext.

\newpage

%% ===== Page 27 =====

de $\Gamma_Q$ sur $\hat{\mathcal{C}}_{1,1}^+$ ne permet pas à
priori de reconstituer (à iso unique
près) l'ext. $\mathcal{N}_{1,1}$ de $\Gamma_Q$ par $\hat{\mathcal{C}}_{1,1}^+$
de (1). Elle donne seulement l'
extension $\mathcal{N}_{1,1}'$ de $\Gamma_Q$ par
$\hat{\mathcal{C}}_{1,1}^{' +} \simeq \operatorname{Sl}(2,\hat{\mathbb{Z}})^{\wedge} / \{\pm 1\}$ $= \widehat{Sl}(2,\mathbb{Z})'$ qui est
image inverse de l'extension

(13) $$ 1 \to \widehat{Sl}(2,\hat{\mathbb{Z}})' \to \operatorname{Aut}_{ext, c}(\operatorname{Sl}(2,\hat{\mathbb{Z}})^{\wedge}) \to \operatorname{Aut}_{ext,c}([unclear]) $$

Cette extension $\mathcal{N}_{1,1}'$ a été ainsi interprétée
comme l'extension $\mathcal{M}_{0,3} \simeq \mathcal{N}_{0,4}'$ de
$\Gamma_Q$ par $\hat{\mathcal{C}}_{0,3} \simeq \hat{\mathcal{C}}_{0,4}'$ (groupe qui est
lui-même ext de $\mathfrak{S}_3 \simeq \operatorname{Sl}(2,\mathbb{Z}/2\mathbb{Z})$
par $\hat{\Pi}_{0,3}$). Nous savons déjà que
l'opération ext. de $\Gamma_Q$ sur le noyau
$\hat{\Pi}_{0,3}$ de $\operatorname{Sl}(2,\hat{\mathbb{Z}})^{\wedge} \to \operatorname{Sl}(2,\mathbb{Z}/2\mathbb{Z})$ est
fidèle, et [unclear: commutante] à l'op. ext.
de $\mathfrak{S}_3 \simeq \operatorname{Sl}(2,\mathbb{Z}/2\mathbb{Z}) \simeq \operatorname{Sl}(2,\mathbb{Z})^{\wedge} / \hat{\Pi}_{0,3}$.
Mais cette [unclear: commutativité] peut [unclear: masquer]

\newpage

%% ===== Page 28 =====

nant s'interpréter [unclear: par] comme conséquence du
fait que l'opération ext. de $\Gamma_Q$ sur
$\mathfrak{S}_3 \simeq \operatorname{Sl}(2, \mathbb{Z}/2\mathbb{Z})$ est triviale, elle ---
(dans l'optique présente)
conséquence du fait qu'elle se fait par
l'intermédiaire d'une op. (fidèle) de
$(\mathbb{Z}/n\mathbb{Z})^\ast_{/2}$ , ou pour $n=2$ $(\mathbb{Z}/n\mathbb{Z})^\ast_{/2} = 1$
(assertion qui est peut être fausse pour
les (valeurs, $n=1,2$ i.e. $n \ge 3$). Cette
relation de commutativité, de nature
assez triviale, apparait donc à présent
comme un cas très particulier de
relations congruentielles nettement
plus subtiles. On peut [unclear: en exprimer de]
façon directe en termes de l'
opération ext. de $\Gamma_Q$ sur $\hat{\Pi}_{0,3}$, en
introduisant les noyaux des hom.
$\hat{\mathcal{C}}_{1,1}^{+ \prime} \to \operatorname{Sl}(2, m)$, qui sont des
sous-groupes ouverts

(14) $$ \Pi(2m) \subset \hat{\Pi}_{0,3} $$

distingués et invariants par $\mathfrak{S}_3$, noyaux

\newpage

%% ===== Page 29 =====

d'homomorphismes

(15) $$ \varphi_{2m} : \hat{\Pi}_{0,3} \rightarrow \operatorname{Sl}(2, \mathbb{Z}/2m\mathbb{Z})' $$

qu'il ne doit pas être bien difficile d'expliciter, provenant par "réduction mod $2m$" d'un hom.

(16) $$ \varphi_{\infty} : \hat{\Pi}_{0,3} \rightarrow \operatorname{Sl}(2, \hat{\mathbb{Z}})' $$

dont il a été question déjà précédem-ment - hom. dont l'image est le sous-groupe noyau de $\operatorname{Sl}(2, \hat{\mathbb{Z}})' \rightarrow \operatorname{Sl}(2, \mathbb{Z}/2\mathbb{Z})$, i.e. on a des suites exactes

(17) $$ 1 \rightarrow \Pi(\infty) \rightarrow \hat{\Pi}_{0,3} \rightarrow \operatorname{Sl}(2, \hat{\mathbb{Z}})' \rightarrow \operatorname{Sl}(2, \mathbb{Z}/2\mathbb{Z}) \rightarrow 1 $$
(18) $$ 1 \rightarrow \Pi(2m) \rightarrow \hat{\Pi}_{0,3} \rightarrow \operatorname{Sl}(2, \mathbb{Z}/2m\mathbb{Z})' \rightarrow \operatorname{Sl}(2, \mathbb{Z}/2\mathbb{Z}) \rightarrow 1 $$

Ceci posé, ces suites exactes permettent de façon naturelle de définir une opération ext. de $\hat{\mathbb{Z}}^* / \hat{\mathbb{Z}}^{*2}$ sur $\hat{\Pi}_{0,3} / \Pi(\infty)$, ou encore de $(\mathbb{Z}/2m\mathbb{Z})^*_2$ sur le $\hat{\Pi}_{0,3} / \Pi(2m)$, et la "condition de congruences" s'exprime en disant que l'action cyclotomique $\chi$ de $\Gamma_Q$ (déduite de l'opération ext. de

\newpage

%% ===== Page 30 =====

$\Gamma_{\mathbb{Q}}$ sur les [unclear: gr. vr.] : (base $\hat{\Pi}_{0,3}$), prov. d'un
[unclear]

$$ (19) \quad \Gamma_{\mathbb{Q}} \xrightarrow{\chi} \hat{\mathbb{Z}}^* \xrightarrow{can.} \hat{\mathbb{Z}}^* / \hat{\mathbb{Z}}^{* 2} $$

donne une op. ext. de $\Gamma_{\mathbb{Q}}$ sur
$\hat{\Pi}_{0,3} / \pi(\infty)$, est compatible avec
l'op. ext. de $\Gamma_{\mathbb{Q}}$ sur $\hat{\Pi}_{0,3}$ lui m.
Pour chaque entier $m$, cette condition
s'explicite par le fait que $\pi(2m)$ est
invariant par l'op. ext. de $\Gamma_{\mathbb{Q}}$, et que
l'opération ext. induite sur
$\hat{\Pi}_{0,3} / \pi(2m)$ est celle déduite du
caractère composé de (19) avec
$\hat{\mathbb{Z}}^* / \hat{\mathbb{Z}}^{* 2} \to (\mathbb{Z}/2m\mathbb{Z})^*_2$ :

$$ (20) \quad \Gamma_{\mathbb{Q}} \to (\mathbb{Z}/2m\mathbb{Z})^*_2 $$

On devrait regarder si cette formulation
est vraiment équivalente à celle des
conditions de congruences de $\hat{C}_{1,1}^{+ '}$ :
les entraine, (équival. à y rajouter une
conséquence ; exercice !) Par contre, la

\newpage

%% ===== Page 31 =====

question de savoir si ces conditions
constituent bel et bien une restriction
effective, sur un autom. ext. du
gr. vo: lente $\hat{\pi}_{0,3}$, commutant à
l'action de $\mathfrak{S}_3$, risque fort d'être
hors de portée, faute de savoir
construire effectivement de tels autom.
(ou ne fût-ce des éléments extérieurs à $\widehat{\Gamma}_{0,3}$!)
extérieurs) [unclear]. Au fond je ne
justement que par voie arithmétique,
à l'aide de l'opérat. de $\Gamma_{\mathbb{Q}}$ sur
$\hat{\pi}_{0,3}$ ! (Il faudrait quand même
refaire à l'occasion un effort en ce
sens ...)

Je me pose aussi la question, dans
quelle mesure on peut reconstituer, en
termes de la théorie des groupes profinis
le centre et le [unclear] d'autom. ext. de
$\hat{\pi}_{0,3}$, ou de [crossed out: $\widehat{\Gamma}_{1,1}$] $\simeq \widehat{C}_{1,1}^+$, l'homo-

\newpage

%% ===== Page 32 =====

a) l'extension plus fine $N_{1,1}$ de $\Gamma_{\mathbb{Q}}$ par $\hat{C}^+_{1,1}$, qui est donc une ext. de $N'_{1,1}$ (censé être connu) par le gr. [unclear: vu] $\{\pm 1\}$

b) l'hom. $N_{1,1} \xrightarrow{(21)} \operatorname{Gl}(2, \hat{\mathbb{Z}})$ [de (11)]
et à défaut, l'hom. [crossed out text]

(22) $$N'_{1,1} \simeq \hat{M}_{0,3}^{pur} \longrightarrow \operatorname{Gl}(2, \hat{\mathbb{Z}})/\{\pm 1\}$$

Je commencerai par la question b).
Si on dispose de l'extension $N_{1,1}$ de $\Gamma_{\mathbb{Q}}$ par $\hat{C}^+_{1,1} \simeq \widehat{Sl(2, \mathbb{Z})}$, on a un hom.

(23) $$N_{1,1} \longrightarrow \operatorname{Aut}_{ext}(\hat{C}^+_{1,1})$$

qui par un quotient en donne un hom.

(24) $$N_{1,1} \longrightarrow \operatorname{Aut}_{ext}(\operatorname{Sl}(2, \hat{\mathbb{Z}}))$$

on a un hom. évident

$$\operatorname{Gl}(2, \hat{\mathbb{Z}}) \longrightarrow \operatorname{Aut}_{ext}(\operatorname{Sl}(2, \hat{\mathbb{Z}}))$$

qui se factorise en

(25) $$\operatorname{Gl}(2, \hat{\mathbb{Z}}) \longrightarrow \operatorname{Gl}(2, \hat{\mathbb{Z}})/\hat{\mathbb{Z}}^* \simeq PGl(2, \hat{\mathbb{Z}}) \hookrightarrow \operatorname{Aut}_{ext}(\operatorname{Sl}(2, \hat{\mathbb{Z}}))$$

et il se trouve que (24) se factorise en

(26) $$N_{1,1} \longrightarrow \underset{\simeq \operatorname{Gl}(2, \hat{\mathbb{Z}})/\hat{\mathbb{Z}}^*}{PGl(2, \hat{\mathbb{Z}})} \hookrightarrow \operatorname{Aut}_{ext}(\operatorname{Sl}(2, \hat{\mathbb{Z}}))$$

\newpage

%% ===== Page 33 =====

ce qui permet donc de reconstituer,
à partir de l'extension $\mathcal{N}_{1,1}$, un
homom

(27) $$ \mathcal{N}_{1,1} \xrightarrow{\varphi} GP(1,\hat{\mathbb{Z}}) \simeq \operatorname{Gl}(2,\hat{\mathbb{Z}})/\hat{\mathbb{Z}}^\times \ , $$

un peu plus faible que l'homom (21)
qu'on cherchait à construire. Mais
notons que l'homom
$u \mapsto (\bar{u}, \det(u))$

(28) $$ \operatorname{Gl}(2,\hat{\mathbb{Z}}) \longrightarrow GP(1,\hat{\mathbb{Z}}) \times \hat{\mathbb{Z}}^\times $$

a comme noyau l'ens des
$\lambda \in \hat{\mathbb{Z}}^\times$ tels que $\lambda^2 = 1$ , i.e. le groupe

(29) $$ _2\hat{\mathbb{Z}}^\times \simeq (\mathbb{Z}/2\mathbb{Z})^P \qquad P = \text{ens des nbs premiers} $$

Ainsi, un homom
$u \mapsto (\varphi(u), \chi(u))$

(30) $$ \mathcal{N}_{1,1} \longrightarrow GP(1,\hat{\mathbb{Z}}) \times \hat{\mathbb{Z}}^\times $$

(déduit de (27) et de $\chi$) [MARGIN: c'est-à-dire compte tenu de l'hom $\mathcal{N}_{1,1} \to \Gamma_{\mathbb{Q}}$ définissant la cat. ext. $\mathcal{N}_{1,1}$ de $\Gamma_{\mathbb{Q}}$ par $\hat{\mathcal{C}}_{1,1}^+$, et de l'hom $\Gamma_{\mathbb{Q}} \to \hat{\mathbb{Z}}^\times$ i.e. $\chi$] se factorise
[unclear: par l'image] de (28), et définit donc
un hom

(31) $$ \mathcal{N}_{1,1} \xrightarrow{\psi} \operatorname{Gl}(2,\hat{\mathbb{Z}})/_2\hat{\mathbb{Z}}^\times $$

qui est encore une version affaiblie de (21).

\newpage

%% ===== Page 34 =====

[MARGIN: 325]

A vrai dire, l'homomorphisme (31)
se factorise en

(32) $$[N_{1,1} \to] \underset{N_{1,1}/\pm 1}{N'_{1,1}} \xrightarrow{\Psi'} \operatorname{Gl}(2,\hat{\mathbb{Z}})/\hat{\mathbb{Z}}^*$$

et l'hom $\Psi'$ peut sûrement se construire
directement à partir de l'extension $N'_{1,1}$
de $\Gamma_{\mathbb{Q}}$ par $\pi'_{1,1} \simeq \hat{\mathcal{S}}_{0,3}^+ \simeq \operatorname{Sl}(2,\hat{\mathbb{Z}})'$, tout comme
tantôt $\Psi$ : il suffit de vérifier que l'hom
canonique

(33) $$\underset{\underset{\operatorname{Gl}(2,\hat{\mathbb{Z}})/\hat{\mathbb{Z}}^*}{\simeq}}{GP(1,\hat{\mathbb{Z}})} \longrightarrow \operatorname{Aut}_{ext}(\operatorname{Sl}(2,\hat{\mathbb{Z}})')$$

est injectif. La question (que je n'ai pas approfondie...)
explicite [crossed out: de relier] comment relier [crossed out: le $\Psi'$] de (31) avec
l'hom. canonique

(34) $$N'_{1,1} \longrightarrow \operatorname{Gl}(2,\hat{\mathbb{Z}})' \hookrightarrow \operatorname{Gl}(2,\hat{\mathbb{Z}})/\pm 1$$

- une fois cet hom. obtenu, on en
déduit l'ext. $N_{1,1}$ par image inverse de
la situation d'ext.

(35) $$1 \to \pm 1 \longrightarrow \operatorname{Gl}(2,\hat{\mathbb{Z}}) \longrightarrow \operatorname{Gl}(2,\hat{\mathbb{Z}})' \to 1,$$

ainsi que l'hom (21) $N_{1,1} \to \operatorname{Gl}(2,\hat{\mathbb{Z}})$.

\newpage

%% ===== Page 35 =====

Donc en résumé, la question essentielle
semble ici la détermination du
relèvement (34) de (32).

Reprenons nos arguments. On se
donne un sous-groupe $\Gamma$ de $\operatorname{Aut}_{ext}(\hat{\Pi}_{0,3})$,
contenant $\mathfrak{S}_3$, d'où une extension
[unclear: $\tilde{\Gamma}$] (notée $\mathcal{M}_{0,3}$ ou $\mathcal{N}_{1,1}'$) de $\Gamma$ par
$\hat{\mathfrak{S}}_{0,3}^+$, ou encore de $\mathfrak{S}_3 \times \Gamma$ par $\hat{\Pi}_{0,3}$,
[crossed out: d'où un homomorphisme]

(36) $\mathcal{M}_{0,3} (\mathcal{N}_{1,1}') \to \operatorname{Aut}_{ext}(\hat{\mathfrak{S}}_{0,3}^+)$

induisant sur $\hat{\mathfrak{S}}_{0,3}^+$ l'opération [crossed out: l'opération] de
l'int. de ce groupe sur lui-m.
On a par ailleurs un hom. injectif

(37) $\operatorname{Gl}(2,\hat{\mathbb{Z}}) \hookrightarrow \operatorname{Aut}(\operatorname{Sl}(2,\hat{\mathbb{Z}})')$

et $\Gamma \subset \operatorname{Aut}_{ext}(\hat{\Pi}_{0,3})$ est astreint (si ce
n'est déjà une conséquence de l'hyp. de commu-
tation à $\mathfrak{S}_3$) à la cond. que l'action
ext. [crossed out: de $\Gamma$ sur] 
(38) $\hat{\mathfrak{S}}_{0,3}^+ \simeq \operatorname{Sl}(2,\hat{\mathbb{Z}})'$
$\simeq \operatorname{Sl}(2,\mathbb{Z})'^\wedge$

\newpage

%% ===== Page 36 =====

soit compatible aux actions ext. de
$\Gamma$, sur $\widehat{\mathfrak{S}}_{0,3}^+$ par tautologie (du fait
que $\Gamma$ est formé d'él. de $\operatorname{Aut}_{ext}(\hat{\Pi}_{0,3})$
qui commutent à $\mathfrak{S}_3$), sur $\operatorname{Sl}(2, \hat{\mathbb{Z}})'$
via $\hat{\mathbb{Z}}^\ast/\hat{\mathbb{Z}}^{\ast 2}$ et le caractère "cyclotomique"

(39) $\chi: \Gamma \to \hat{\mathbb{Z}}^\ast$

déduit de l'op. ext. tautologique de $\Gamma$
sur $\hat{\Pi}_{0,3}$ (respectant la structure : [unclear]).
Cette "condition de congruences"
assure que (36) se factorise par (37),
en un hom. bien déterminé

(40) $\varphi': M_{0,3} \to Ap(1, \hat{\mathbb{Z}})$

d'où, en utilisant de plus ledit caractère
cyclotomique, un hom. $u \mapsto (\varphi'(u), \chi(u))$

(41) $M_{0,3} \to Ap(1, \hat{\mathbb{Z}}) \times \hat{\mathbb{Z}}^\ast$

et la condition co-essentielle nous assure
un fait, plus précis, que ce dernier
se factorise via un hom.

(42) $M_{0,3} \xrightarrow{\Psi'} \operatorname{Gl}(2, \hat{\mathbb{Z}})/\hat{\mathbb{Z}}^\ast$

compte tenu de l'injection

\newpage

%% ===== Page 37 =====

(43) $$\operatorname{Gl}(2, \hat{\mathbb{Z}})/_2\hat{\mathbb{Z}}^* \hookrightarrow GPl(\hat{\mathbb{Z}}) \times \hat{\mathbb{Z}}^*$$
$$u \mod {}_2\hat{\mathbb{Z}}^* \mapsto (u \mod \hat{\mathbb{Z}}^*, \det u)$$ .

Considérons l'hom. can. (surjectif)
(44) $$\operatorname{Gl}(2, \hat{\mathbb{Z}})' \to \operatorname{Gl}(2, \hat{\mathbb{Z}})/_2\hat{\mathbb{Z}}^*$$
déf ||
$$\operatorname{Gl}(2, \hat{\mathbb{Z}})/\pm 1$$

dont le noyau est isom. à
(45) $$\nu = (\mathbb{Z}/2\mathbb{Z})^P / diag. \qquad (P = \text{ens. des nb. premiers}).$$

L'hom. canonique $\psi'$ (42), déduit
des conditions de congruence, définit
par image inverse une extension centrale
(46) $$1 \to \nu \to \widetilde{\mathcal{M}}_{0,3} \to \mathcal{M}_{0,3} \to 1 \quad ,$$

et la donnée d'un hom. [unclear: barré: (39)]
(47) $$\mathcal{M}_{0,3} \to \operatorname{Gl}(2, \hat{\mathbb{Z}})' \quad$$ [MARGIN: celle d' un relèvt à $\nu$]
qui relèverait (42), équivaut à un
scindage de cette extension (46). La
question à élucider à présent : est-il
possible d'expliciter un tel relèvement canonique
(42), ou scindage de (46), en termes de la

\newpage

%% ===== Page 38 =====

seule donnée du $\hat{\Gamma}$ Teichmuller $(\Pi_{0,3})$, satisfai-
sant aux conditions de congruences usuelles
énoncées plus haut ?

Notons que sur le sous-groupe
$\hat{\Gamma}_{0,3}^{+} \simeq \operatorname{Sl}(2, \mathbb{Z})^\wedge$ de $M_{0,3}$, le relèvement
cherché (47) est déjà connu, comme
étant $\operatorname{Sl}(2, \mathbb{Z})^\wedge \to \operatorname{Sl}(2, \hat{\mathbb{Z}})' \hookrightarrow \operatorname{Gl}(2, \hat{\mathbb{Z}})'$.
En d'autres termes, le scindage cherché
de (46) est déjà connu au-dessus
du sous-groupe $\hat{\Gamma}_{0,3}^{+}$ de $M_{0,3}$.

L'indétermination du rel. ou scindage
cherché a priori est un élément
$$\operatorname{Hom}(M_{0,3}, \mathcal{V})$$
Tenu compte de la condition précédente,
elle sera dans
(48) $$\operatorname{Ker}(\operatorname{Hom}(M_{0,3}, \mathcal{V}) \to \operatorname{Hom}(\hat{\Gamma}_{0,3}^{+}, \mathcal{V}))$$
$$\simeq \operatorname{Hom}(\Gamma, \mathcal{V})$$
si on suppose que $\Gamma / \Gamma_{ab} \overset{\chi}{\to} \hat{\mathbb{Z}}^*$
est un iso, l'indétermination serait encore
dans $\operatorname{Hom}((\hat{\mathbb{Z}}^*)_2, \mathcal{V})$ - c'est gros ! Il

[MARGIN: 37]

\newpage

%% ===== Page 39 =====

faut un élément un peu astucieux
pour justifier sa présence - on devine
qu'il est lié justement à l'interprétation
d'un $M_{0,3}$ comme un $M_{1,1}$ - i.e. à
l'introduction, dans une "théorie profinie
de courbes rationnelles", de la "théorie
profinie des courbes elliptiques".

40 Paraphernalia topologico-isotopiques (autour
de la théorie de Teichmüller des groupes fondamentaux)

Pour tâter le terrain, j'ai envie d'abord
de partir de la situation des groupes
discrets, qu'on a mieux en mains, c'est
le cas de le dire !

Soit $C_{0,3}$ la catégorie des groupes
(avec les isomorphismes comme flèches)
(discrets) : tous du type $0,3$, dont le
représentant-type est $\Pi_{0,3}$ - cette catégorie
est un groupoïde connexe : elle
est torseur sous le groupe

$$ (49) \quad \operatorname{Aut}_{ext}(\Pi_{0,3}) \simeq \mathfrak{S}_3 \simeq \operatorname{Gl}(2, \mathbb{Z})^* / \pm 1 \simeq P\operatorname{Gl}(2, \mathbb{Z}) $$
[MARGIN: sous Aut, l'auteur a écrit "extérieurs"]

Soit $C_{1,1}$ la catégorie des groupes : tous
de type $(1,1)$, dont le représentant-type est
$\Pi_{1,1} = \{ x,y,c \ | \ [x,y]c = 1 \}$ ; cette catégorie est

\newpage

%% ===== Page 40 =====

équivalente à celle des torseurs sous le groupe
$C_{1,1}^{ext} \simeq$
$\operatorname{Aut}\, ext(\Pi_{1,1}) \simeq \operatorname{Gl}(2,\mathbb{Z})$ - l'isomorphisme

(50) $$\operatorname{Aut}\, ext(\Pi_{1,1}) \longrightarrow \operatorname{Gl}(2,\mathbb{Z})$$

il s'obtient en associant à un automorphisme
extérieur de $\Pi_{1,1}$ son action sur

(51) $$(\Pi_{1,1})_{ab} \simeq \mathbb{Z}^2$$

On a donc un foncteur canonique
[MARGIN: $\mathcal{L}$, initiale de "Legendre"]
(52) $$\mathcal{L} : C_{1,1}^{ext} \longrightarrow C_{0,3}$$

défini (à isomorphisme près) par la condition
de rendre commutatif le diagramme

(53)
\begin{tikzcd}
C_{1,1}^{ext} \ar[r] \ar[d, "\wr"'] & C_{0,3} \ar[d] \\
Tors(\operatorname{Gl}(2,\mathbb{Z})) \ar[r] & Tors(GPl(1,\mathbb{Z}))
\end{tikzcd}

où
$$Tors(\operatorname{Gl}(2,\mathbb{Z})) \longrightarrow Tors(GPl(1,\mathbb{Z}))$$
est le foncteur de chgt de groupe d'opérateurs
relatif à :
(54) $$\operatorname{Gl}(2,\mathbb{Z}) \xrightarrow{can} GPl(1,\mathbb{Z}) \simeq \operatorname{Gl}(2,\mathbb{Z}) / \pm 1$$

\newpage

%% ===== Page 41 =====

On peut [unclear: aussi] introduire le groupoïde
$C_{1,1}^{ext \ \prime}$, déduit de $C_{1,1}^{ext}$ en prenant le
quotient par l'automorphisme universel $-id$
- correspondant au centre $(\pm 1)$ (des groupes
fondamentaux de ce groupoïde - on trouve
un groupoïde connexe équivalent canon.
[MARGIN: groupoïde] à celui des torseurs sous $\operatorname{Gl}(2,\mathbb{Z})'= \operatorname{Gl}(2,\mathbb{Z})/(\pm 1)$ ,

dont on tire maintenant une
équivalence de catégories

(55) $$C_{1,1}^{ext \ \prime} \overset{\simeq}{\longrightarrow} C_{0,3}$$

rendant commutatif

(56)
\begin{tikzcd}
C_{1,1}^{ext \ \prime} \ar[r, "\simeq"] \ar[d] & C_{0,3} \ar[d] \\
Tors(\operatorname{Gl}(2,\mathbb{Z})') \ar[r, "\simeq"] & Tors(GPl(1,\mathbb{Z}))
\end{tikzcd}
,

où $Tors(\operatorname{Gl}(2,\mathbb{Z})') \overset{\simeq}{\longrightarrow} Tors(GPl(1,\mathbb{Z}))$ est
le foncteur tautologique associé à l'iso

(57) $$\operatorname{Gl}(2,\mathbb{Z})' \overset{\simeq}{\longrightarrow} GPl(1,\mathbb{Z})$$ .

\newpage

%% ===== Page 42 =====

En somme, je suis en train de faire une
paraphrase topologique tautologique, en termes
des groupes fondamentaux : [unclear], de
la théorie de Legendre des relations entre
courbes elliptiques et courbes rationnelles.

J'aimerais expliciter maintenant l'équi-
valence de Legendre (55), qui avait
été obtenue grâce à la connaissance des
groupes de Teichmüller discrets relatifs
aux deux membres et d'un isomorphisme
ad-hoc entre ces deux groupes,
directement en termes des groupes
fondamentaux des types $\Pi_{0,3} - \Pi_{1,1}$.

Si $\pi'$ est un groupe de lacets de type
$1,1$, il convient d'abord de noter qu'il
y a dans $Autext(\pi')$ un élément canonique
[MARGIN: NB [unclear] conserve l'orientation i.e. [unclear] $+1$ [unclear] multiplicité 2 [unclear] centre]
d'ordre 2 - en fait l'unique élément non trivial
du centre, noté comme de juste $(-1)_{\pi'}$.

\newpage

%% ===== Page 43 =====

[unclear: const !]

Ainsi, on peut considérer $\pi'$ comme muni
d'une [unclear: certaine] opération extérieure de
$G = \mathbb{Z}/2\mathbb{Z}$ sur lui, d'où une extension
canonique
$$(58) \quad 1 \to \pi' \to \mathcal{E} \to \mathbb{Z}/2\mathbb{Z} \to 1 .$$

Le modèle topologique nous apprend
qu'il y a exact. trois classes de
conjugaison de scindages de cette
extension, correspondant aux trois
points d'ordre 2 d'un groupe topologique
$E \simeq U \times U$ - les trois points fixes de
$x \mapsto -x$ dans $E \setminus \{0\}$. La théorie des

"franges de [unclear: tresses] simultanées", relative à une
situation de groupes / d'opérateurs exté-
rieurs dans de tels groupes, à centre (théorie que
je n'ai [MARGIN: mais j'en ai eu gd besoin ds d'autres cas, pr dégager les pptés des groupes de tresses ds le plan [unclear] ] dans le cas discuté ici vraiment
écrite - il y a pour le moment des
[unclear] ) permet alors

\newpage

%% ===== Page 44 =====

310

de construire canoniquement un groupe
(noté $\tilde{\pi}'$)
extension de type $\Pi_{1,4}$, avec une
opération extérieure de $G = \mathbb{Z}/2\mathbb{Z}$ sur ce
groupe, et une partie $J$ de cardinal 3

[MARGIN: en fait $G$ opère trivialement sur $I = I(\tilde{\pi}')$]

de $I = I(\tilde{\pi}')$ (qui est de cardinal 4), $J$
stable par $G$ i.e. dc fixée par $G$,
et un iso. de type $\Pi_{1,1}$ du groupe $\tilde{\pi}'(J)$ (déduit
par "bouchage du trou de $I \setminus J$") avec $\pi'$

(59) $$ \tilde{\pi}'(J) \simeq \pi' $$ ,

qui soit compatible avec les opérations
de $G$. Dans le modèle topologique,
$\tilde{\pi}'$ est le groupe fondamental cat. de l'espace
$E \setminus {}_2 E$, sur lequel $G$ opère librement.
L'extension $\tilde{\mathcal{E}}$ de $G$ par $\tilde{\pi}'$ correspondante
à l'opération de $G$ sur $\tilde{\pi}'$

(60) $$ 1 \to \tilde{\pi}' \to \tilde{\mathcal{E}} \to \mathbb{Z}/2\mathbb{Z} \to 1 $$ ,

s'identifie dc au groupe fondamental
mixte de $(E \setminus {}_2 E, G)$, s'identifie de [unclear].

\newpage

%% ===== Page 45 =====

au groupe fondamental de 
$$(61) \quad V = (E \setminus_2 E)/\{\pm 1\}$$
qui est une surface top. de type $(0,4)$.
En d'autres termes, $\tilde{\mathcal{E}}$ peut être
considéré comme un groupe à lacets
de type $(0,4)$, [unclear: muni] d'une structure
à lacets forts : déduite de 
celle de $\tilde{\Pi}'$, et pour laquelle on a
un iso. canonique
$$(62) \quad I(\tilde{\mathcal{E}}) = I(\tilde{\Pi}')/G \simeq I(\Pi')$$
Ainsi on dispose d'un $i_0 \in I(\tilde{\mathcal{E}})$,
ce qui permet de "boucher le trou en
$i_0$" pour obtenir un groupe effectif
à lacets de type $(0,3)$, soit
$$(63) \quad \Pi = \mathcal{L}(\Pi')$$
Reprenons un peu encore la relation entre
$\Pi$ et $\Pi'$ on peut procéder ainsi : d'un
groupe [inserted: à lacets] effectif de type $(0,3)$, on [unclear] pour

\newpage

%% ===== Page 46 =====

le foncteur "bouchage de trou" [unclear: $\tilde{\Pi} \mapsto \Pi$] un groupe
: lacet extérieur de type $(0,4)$, muni d'
un élément $i_0 \in I(\tilde{\Pi})$ tel que [crossed out text]

(64) $$I(\tilde{\Pi}) \simeq \{i_0\} \amalg I(\Pi)$$
[MARGIN: en dessous de $I(\Pi)$: cardinal 3]

On a alors ($T$ est le module d'orientation
de $\Pi$)

(65) $$\tilde{\Pi}_{ab} \simeq T^{I(\tilde{\Pi})} \qquad (\tilde{\Pi}_{ab})_2 \simeq (\mathbb{Z}/2\mathbb{Z})^{I(\tilde{\Pi})}$$

et on a un hom. (ext.)

(66) $$\tilde{\Pi} \longrightarrow \mathbb{Z}/2\mathbb{Z}$$

canonique, défini par la condition d'être
"ramifié" en les $i \in I(\tilde{\Pi})$ - i.e. comme
composée de l'hom.

(67) $$(\tilde{\Pi}_{ab})_2 \simeq (\mathbb{Z}/2\mathbb{Z})^{I(\tilde{\Pi})} \longrightarrow \mathbb{Z}/2\mathbb{Z}$$

somme des coefficients. Le noyau $\tilde{\Pi}_1$ de
(66) est un groupe : lacet de type $(1,4)$
- mais il n'est pas tout à fait un groupe
extérieur - il n'est pas défini "à un
automorphisme intérieur" près, mais seulement modulo (en plus)
l'opération extérieure de symétrie dans $\mathbb{Z}/2\mathbb{Z}$

\newpage

%% ===== Page 47 =====

dessus. Si on décrit le groupe $\tilde{\pi}$
(dans sa "classe extérieure") [unclear] on
trouve une suite exacte

(68) $$1 \to \tilde{\pi}' \to \tilde{\pi} \to \mathbb{Z}/2\mathbb{Z} \to 1$$

où $\tilde{\pi}'$ est un groupe "limite", avec

(69) $$I(\tilde{\pi}') \simeq I(\tilde{\pi})$$

En bouchant les trous relatifs :
$I(\pi) \subset I(\tilde{\pi}) \simeq I(\tilde{\pi}')$ , [unclear: qui est de codim 1] ,
--- on trouve un groupe extérieur
de type (1,1) , une opération de $\mathbb{Z}/2\mathbb{Z}$
qui n'est autre que l'op. canonique
dessus. Mais la [unclear: conjonction] des deux faits, c'est à dire [unclear]
[unclear] un groupe extérieur de type
(1,1) i.e. un objet de $C_{1,1}^{ext}$ - c'est une
bel illusion, "mirage symétrique" - i.e. un
objet de $C_{1,1}^{ext}$ . Pour avoir un objet
de $C_{1,1}^{ext}$ , il faut avoir "choisi" un [unclear: noyau] $\tilde{\pi}'$
pour l'hom. extérieur $\tilde{\pi} \to \mathbb{Z}/2\mathbb{Z}$ , c'est à dire un
groupe extérieur c'est la classe d'isomorphie

\newpage

%% ===== Page 48 =====

donne une $(\mathbb{Z}/2\mathbb{Z})$-gerbe, i.e. un groupoïde
connexe [unclear: équivalent] à $\mathbb{Z}/2\mathbb{Z}$ comme groupe
fondamental.

Modulo une rédaction plus soignée, il me
semble que ces réflexions un-peu allègres
prouvent bel et bien qu'on a un bon isomorphisme quasi-
isomorphisme canonique (mod autom. intérieurs)

(70) $\tilde{\mathcal{C}}_{1,1} / \{\pm 1\} \xrightarrow{\sim} \widehat{\mathbb{S}}_{0,3}$ [MARGIN: en fait, par la précision qu'il n'admet aucun [unclear] provenant de $\approx \Pi_{0,3}$]

- l'indétermination [unclear: venant je crois du choix d'] un
isomorphisme entre $\mathcal{L}(\Pi_{1,1})$ et $\Pi_{0,3}$.

Ce fait est tout à fait indépendant, il me
semble, de celui que l'hom.
naturel $\tilde{\mathcal{C}}_{1,1} \xrightarrow{\sim} \operatorname{Gl}(2, \mathbb{Z})$
soit un isomorphisme. Je ne vois
aucune [unclear: inconvénient] à identifier purement
et simplement $\tilde{\mathcal{C}}_{1,1}$ à $\operatorname{Gl}(2, \mathbb{Z})$; l'occasion
par contre il est plus sage [unclear: de se contenter]
d'identifier $\widehat{\mathbb{S}}_{0,3}$ à $\tilde{\mathcal{C}}_{1,1}/\{\pm 1\}$ ou $\operatorname{Gl}(2, \mathbb{Z})'$

\newpage

%% ===== Page 49 =====

si ce n'est seulement comme groupe
extérieur, à cause de l'indétermination
provenant de (70).

[MARGIN: avec un peu plus de soin, il apparaît [unclear] inutile de [unclear] des structures complexes, sur des surfaces top. de base, ce qui rendrait [unclear]]

D'un point de vue des multiplicités
modulaires analytiques, on a un
isomorphisme canonique

(71) $$M_{1,1}^{an} \longrightarrow \mathcal{U}_{0,3}^{an} \simeq \mathcal{M}_{0,4}^{1 an} = (\widetilde{U}_{0,3}, \mathfrak{S}_3)$$
[sous $\mathcal{U}_{0,3}^{an}$ : courbe analytique complexe de type (0,3) "universelle"]

dont l'effet sur [unclear] induit isomorphisme

(72) $$M_{1,1}^{\sim} \xrightarrow{\sim} M_{0,3}^{\sim}$$

et l'effet de (71) sur les groupes fondamen-
taux extérieurs n'est autre que l'hom.
can.

(73) $$\pi_{1,1}^+ \longrightarrow \mathfrak{S}_{0,3}^+$$

induit par (70) - est-il possible de
le préciser de façon utile, en un sens
effectif, non seulement extérieur ? Il y
a un revêtement universel canonique
$\widetilde{M}_{1,1}$ de $M_{1,1}$, fourni par l'espace des

\newpage

%% ===== Page 50 =====

Teichmüller (le demi-plan de Poin-
caré $D$ du § 38) - relatif à ce revête-
ment universel on a bel et bien un iso. can.

$$ (74) \quad \pi_1(M_{1,1}^\sim) \simeq \operatorname{Sl}(2, Z) $$

(pas seulement extérieur), qui n'est autre
d'ailleurs, essentiellement, que (50) -
(moyennant un sous-groupe d'indice 2
$\tau_{1,1}^+$ de $\tau_{1,1}$ ). On a bel et bien un
morphisme canonique

[MARGIN: Pour la définition en termes fuchsiens), cf plus bas § ]

$$ (75) \quad D \longrightarrow U_{0,3}^{an} \quad (\simeq M_{1,1}[2]') $$

compatible avec un hom. sur les groupes
d'opérateurs

$$ (76) \quad \operatorname{Sl}(2, Z)' \longrightarrow \operatorname{Sl}(2, Z/2Z) \simeq \mathfrak{S}_3 $$
[unclear: détail, qui rend explicite ds § ]

de sorte que $D$ apparait comme un
rev. universel de $U_{0,3}^{an}$, de
$(U_{0,3}^{an}, \mathfrak{S}_3)$, et $\operatorname{Sl}(2, Z)'$ comme
le $\pi_1$ mixte relatif à ce rev. universel. Si
pour l'identifier à $\pi_1((U_{0,3}^{an}, \mathfrak{S}_3), P=\vec{s})$, il

\newpage

%% ===== Page 51 =====

faut donc choisir un point [MARGIN: de $D$] au dessus
de $-j$ (qui permettra ainsi d'identifier
$D$ à un rev. universel pointé en $P$ de
$\mathcal{U}_{0,3}$ ) - ce qui introduit justement
l'indétermination par un élément du
noyau de $\operatorname{Sl}(2,\mathbb{Z})' \to \operatorname{Sl}(2,\mathbb{Z}/2) \simeq \mathfrak{S}_3$,
lui-même isomorphe à $\Pi_{0,3}$. Mais
on peut prendre p. ex.

(77) $$j = \exp \frac{2 i \pi}{3} = \frac{-1 + i \sqrt{3}}{2}$$

comme représentant pt. au dessus de
$P = -j$ - en touchant du bois pour
qu'il soit bien au dessus de $P$, et non
de $\bar{P} = -\bar{\jmath}$, où $-\bar{\jmath} = \frac{1 + i \sqrt{3}}{2} = \exp \frac{2 i \pi}{6}$
qu'on choisirait. [NB On trouve que
les trois éléments de $\operatorname{Sl}(2,\mathbb{Z})/\pm 1$
qui transforment $j$ en $-\bar{\jmath}$ sont
$\begin{pmatrix} 1 & 1 \\ 0 & 1 \end{pmatrix}$, $\begin{pmatrix} 1 & 0 \\ 1 & 1 \end{pmatrix}$, $\begin{pmatrix} 0 & -1 \\ 1 & 0 \end{pmatrix}$, dont aucun
n'est $\equiv \begin{pmatrix} 1 & 0 \\ 0 & 1 \end{pmatrix} \bmod 2$, ce qui montre que
$j$ et $-\bar{\jmath}$ ont des images ds $\mathcal{U}_{0,3}$ non confondues. ]

\newpage

%% ===== Page 52 =====

314
[MARGIN: $\Gamma(2)$]
mod le noyau de $\operatorname{Sl}(2, \mathbb{Z}) \to \operatorname{Sl}(2, \mathbb{Z}/2\mathbb{Z})$,
[unclear: ce qui fait] qu'ils n'ont pas = images dans $U_{0,3}$
$\simeq D/\Gamma(2)$ , donc leurs images sont
$P = - \bar{j}$ et $\bar{P} = -j$ (ou inversement)$]$. Le
choix en question est alors caractérisé
par le fait que $\tau \in D$ au dessus de
$P$ est alors une racine de l'unité, ce qui
est quand même bien naturel. On peut
donc dire qu'on a "normalisé" l'hom.
(73) en un hom effectif - et de
même pour (70).

On aurait envie de reprendre les
constructions précédentes dans le contexte
des groupes profinis - malheureusement
je manque des fondements nécessaires,
et surtout d'une théorie des "fromages"
- "fromages de tresses" au niveau groupes
à lacets profinis - même des [unclear: vraies] [unclear: conf.].

\newpage

%% ===== Page 53 =====

locale. Pourtant on sent que c'est
bien dans cette direction que se trouvera
une compréhension "gruppentheoretisch"
de l'homomorphisme ~~[unclear]~~ (40)

$$\hat{M}_{0,3} \to \operatorname{Gl}(2, \hat{Z})/\pm 1 \simeq PGl(2, \hat{Z}),$$

qui est la motivation principale des
réflexions du présent §, et du précédent.
Ce devrait être une incitation pour une
réflexion renouvelée sur la question du
tressage dans le contexte profini... Mais
je me rends compte dès maintenant
(la questi- s'il y a lieu)
que ~~[unclear]~~ extension des constructions du
présent §, au contexte profini, indépen-
damment de l'introduction d' éléments
de structure "arithmétiques" (tels que
la donnée d' un sous-groupe $\Gamma \subset Out(\hat{\Pi}_{0,3})$)
satisfaisant à un certain nombre de
conditions $\pm$ subtiles, genre relations
bi-pentagonales) est subordonnée à celle,

\newpage

%% ===== Page 54 =====

[MARGIN: 315]

si ces relations congruentielles sont bel et
bien des relations de nature arithmétique
spécifiques aux autom. ext. provenant
de $\Gamma_Q$, ou au contraire des
conséquences $\pm$ automatiques d'un
purement "algébrique"
formalisme des groupes profinis, calqué
sur le cas discret. Sans préjuger aucunement
de la réponse et du degré de
plausibilité d'un tel formalisme, je
pourrais à titre expérimental essayer
d'en dégager un, pour tester où il
mènera.

## § 11 Conditions de commutation de $\hat{\Pi}_{0,3} \xrightarrow{\sim} \widehat{\mathfrak{S}}_3^+$ aux opérations de $\Gamma_Q$. Calcul d'autom. extérieurs de $\hat{\Pi}_{0,3}$ et $\widehat{\mathfrak{S}}_3^+$.

Avant de m'élancer dans cette
tentative, je voudrais encore élucider
une question, concernant d'autres
éventuelles conditions nécessaires (en plus
des conditions congruentielles du § 8)
pour qu'un autom. ext. de $\hat{\Pi}_{0,3}$ provienne

\newpage

%% ===== Page 55 =====

de $\Gamma_{\mathbb{Q}}$. Pour ceci, reprenons l'hom.
extérieur conjugué $\varphi$

$$ (1) \quad \Pi_{0,3} \xrightarrow{epi} \Pi_{0,3}' \xrightarrow{\sim} G_{0,3}^+ \simeq \Pi_1((\mathcal{M}_{0,3})_{\overline{\mathbb{Q}}}, P_0) $$
$$ = \Pi_{0,3} / \{ l_1^2, l_3^3 \} $$
$$ \simeq [unclear: gr. de rotations de pavages triangul. orientés] $$

défini par
$$ (2) \quad \begin{cases} \varphi(l_0) = l_0' = \sigma_\infty \rho \\ \varphi(l_1) = l_1' = \sigma_\infty \\ \varphi(l_\infty) = l_\infty' = \rho^{-1} \end{cases} $$

(voir les notes de $[C, ...]$ - -),
qui induit un homom. ext. surjectif

$$ (3) \quad \hat{\varphi} : \hat{\Pi}_{0,3} \xrightarrow{epi} \hat{G}_{0,3}^+ \simeq \pi_1((\mathcal{M}_{0,3})_{\overline{\mathbb{Q}}}, P). $$
$$ \cap $$
$$ \hat{\Gamma}_{0,3} = Aut ext_{loc} (\hat{\Pi}_{0,3}) $$

Cet homomorphisme est défini via les
"formules brutales" (2), sans
référence à un sous-groupe particulier
$\Gamma$ de $Aut ext_{loc}^+ (\hat{\Pi}_{0,3})$. Si on se donne
un tel sous-groupe, comme le
groupe $\Gamma_{\mathbb{Q}}$, celui-ci opère automatiquement
sur le groupe extérieur $\hat{G}_{0,3}^+$, car il

\newpage

%% ===== Page 56 =====

318

s'insère dans une suite exacte

(4) $1 \longrightarrow \widehat{C}_{0,3}^+ \longrightarrow M_{0,3} \longrightarrow \Gamma \longrightarrow 1$ .

comparons à la suite ex. de

(5)
\begin{tikzcd}
1 \ar[r] & \widehat{\Pi}_{0,3} \ar[r] \ar[d] & \mathcal{E}_{0,3} \ar[r] \ar[d, "p"] & \Gamma \ar[r] \ar[d, equal] & 1 \\
1 \ar[r] & \widehat{C}_{0,3}^+ \ar[r] & M_{0,3} \ar[r] & \Gamma \ar[r] & 1
\end{tikzcd}

Le fait que l'homomorphisme
(3) commute à l'action de $\Gamma_{\mathbb{Q}} = \Gamma$
me semble bien a priori une condi-
tion non triviale sur $\Gamma$ - plus
précisément, sur les éléments de
$\Gamma$ séparément, qu'il faudrait bien
arriver à expliciter ([unclear: c a d par les]
conditions congruentielles), en
termes des calculs $\pm$ explicites
de [C]. Cette condition équivaut
à celle que l'hom (3) s'étende
en un homomorphisme d'extensions

\newpage

%% ===== Page 57 =====

[MARGIN: NB. $\mathcal{E}_{0,3}$ est le groupe noté $\mathcal{M}_{0,3}$ dans le § 46]

(6)
\begin{tikzcd}
1 \ar[r] & \hat{\Pi}_{0,3} \ar[r] \ar[d, "\hat{p}"] & \mathcal{E}_{0,3} \ar[r] \ar[d, "\underline{\Phi}"] & \Gamma \ar[r] \ar[d, equal] & 1 \\
1 \ar[r] & \hat{C}_{0,3}^+ \ar[r] & \mathcal{M}_{0,3} \ar[r] & \Gamma \ar[r] & 1 \quad ,
\end{tikzcd}

qui s'identifie à
moyennant le fait que $\mathcal{M}_{0,3}$ (interprété comme
groupe fondamental "à signatures" de
$(\mathbb{P}_{\overline{\mathbb{Q}}}^1 - \{0\}, 2 \{1\} + 3 \{\infty\} ) )$ s'identifie au
quotient de $\hat{\Pi}_{0,3}$ par le sous-groupe
fermé
invariant engendré par $l_1^2, l_\infty^3$.

Je commence à soupçonner
qu'il y aura toute une kyrielle
de conditions de même nature sur
$\Gamma$, qu'il faudrait bien arriver
à dégager de façon plus générale.

\newpage

%% ===== Page 58 =====

317

Je vais reprendre les calculs faits dans
[C], concernant les autom. extérieurs de
$\hat{\Pi}_{0,3}$ qui commutent à $\mathfrak{S}_3$, [crossed out: Ils peuvent] qui
peuvent se réaliser de façon unique
comme automorphismes effectifs (qui
commutent à $\sigma_\infty$ (et qui de plus
commutent extérieurement à $\rho$).
Ils correspondent biunivoquement aux couples

(7) $$ (g_0, p) \quad \begin{cases} g_0 \in \hat{\Pi}_{0,3} \pmod{L_0} \\ p \in \hat{\mathbb{Z}}^* \text{ (multiplicateur)} \end{cases} $$

satisfaisant les conditions suivantes
(où [crossed out: $p \in \hat{\mathbb{Z}}^*$ est déterminé] $p$ est déterminé de façon
unique par $u$ (ou encore, par l'autom.
extérieure correspondant $x$), où on a posé

[MARGIN: NB Autom. ext. $x$ de $\hat{\Pi}_{0,3}$ comm. à $\mathfrak{S}_3 \iff x$ induit sur gr $p(p)$ ... [unclear] ... comm à $h_0$ ...]

$$ h_0 = \sigma_\infty(g_0)^{-1} g_0 = [\sigma_\infty, g_0^{-1}] : $$

[MARGIN: NB fait que $\rho$ op. à g. de $\hat{\Pi}$]
(9) $$ int(h_0)(l_0^p) int(l_1^p \rho(h_0)^{-1})(l_\infty^p) = 1 $$

(10) $$ h_0 l_0^{-p} \rho^{-1}(h_0) = l_1^p \rho(h_0)^{-1} h_0^{-\sigma} $$

[crossed out: L'hom. $u$ est donné par]

[MARGIN: NB (10) s'écrit aussi]
(10 bis) $$ (l_0^p h_0^{-1}) \cdot \rho^{-1}(l_0^p h_0^{-1}) \cdot \rho^{-2}(l_0^p h_0^{-1}) = 1 $$

\newpage

%% ===== Page 59 =====

318

$$ (11) \quad \begin{cases} u(l_0) = int(g_0)(l_0^\rho) \\ u(l_1) = int(\sigma_\infty(g_0))(l_1^\rho) \\ u(l_\infty) = int(g_0 l_0^{-\rho} \rho^{-1}(h_0))(l_\infty^\rho) \end{cases} $$

et il faut en toute rigueur, en plus des conditions (9), (10), exiger que les deux éléments $u(l_0), u(l_1)$ donnés par (11) engendrent [unclear: (forment)] le groupe profini $\hat{\Pi}_{0,3}$.
[MARGIN: (11 bis)]

Je vais calculer l'action de $u$ sur $\hat{\mathcal{C}}^\sim_{0,3}$, en considérant ce groupe comme engendré par les éléments

$$ (12) \quad l'_1 = \sigma_\infty, \quad l'_\infty = \rho^{-1} $$

donc il faut calculer simplement $u(\sigma_\infty)$ et $u(\rho)$. On a par construction ($u$ commutant à $\sigma_\infty$)

$$ (13) \quad u \sigma_\infty u^{-1} = \sigma_\infty $$

et d'autre part on vérifie que 

$$ (14) \quad \rho^{-1} u \rho = int(h) u \quad i.e. \ u \rho u^{-1} = int(\rho(h)) \rho $$
$$ h = (\rho^{-1}\sigma_\infty)(g_0) l_0^\rho g_0^{-1} \quad i.e. \ \rho(h) = \sigma_\infty(g_0) l_1^\rho \rho(g_0)^{-1} $$

donc on trouve, en termes d'une opération sur le groupe $\hat{\mathcal{C}}^\sim_{0,3}$ engendré par $l'_1 = \sigma_\infty$ et $l'_\infty = \rho^{-1}$ (avec $\rho$)

$$ (15) \quad \begin{cases} u(\sigma_\infty) = \sigma_\infty \\ u(\rho) = g \cdot \rho \qquad \text{avec } g = \sigma_\infty(g_0) l_1^\rho \rho(g_0)^{-1} \end{cases} $$

\newpage

%% ===== Page 60 =====

NB La relation $u(\rho)^3 = u(\sigma_\infty)^2 = 1$ est triviale,
la relation $u(\rho)^3 = 1$ se réduit à
$$g \rho(g) \rho^2(g) = 1$$
(16)
qui est conséquence de l'expression explicite
(15) de $g$, et de la formule (10).

Notre [unclear: $g_0$] devrait avoir, [unclear: à y repenser]
$$l'_0 = \sigma_\infty \rho = \varepsilon_0 \quad (\text{où } l'_0 = l_0)$$
[MARGIN: ?]
(17) $$u(l'_0) = int(g_0)(l'_0)$$
- cette relation [unclear: - on conviendra d'attendre que]
la première relation (11)
$$u(l_0) = int(g_0)(l_0 \rho)$$
-
(La relation (17) en est une conséquence,
si on admet que le centralisateur
de $L_0 = l_0^{\hat{\mathbb{Z}}}$ dans $\hat{\Pi}_{0,3}$ est réduit à $L_0$...)

On devrait avoir [unclear: aussi] les relations [unclear: congruences]
(18) $$u(\rho) \equiv \rho \mod \hat{\Pi}_{0,2}$$
sur lesquelles nous reviendrons -
il n'est pas clair pour moi qu'elle
soit conséquence des conditions envisagées
jusqu'à présent.

\newpage

%% ===== Page 61 =====

319

Considérons maintenant l'hom. (18)
$$
(18) \begin{cases}
\hat{\varphi} : \hat{\Pi}_{0,3} \to \hat{G}^+_{0,3} \\
l_0 \mapsto l'_0 = \sigma_0 \rho = \varepsilon_0 \\
l_1 \mapsto l'_1 = \sigma_\infty \\
l_\infty \mapsto l'_\infty = \rho^{-1}
\end{cases}
$$

et écrivons qu'il commute, ou t. au
moins qu'il commute, à un autom. int. près. Cela signifie
qu'il existe un
$$ (19) \quad \lambda \in \hat{G}^+_{0,3} $$
tel qu'on ait
$$ (20) \quad \hat{\varphi}(u(x)) = \lambda u(\hat{\varphi}(x)) \lambda^{-1} \quad \forall x \in \hat{\Pi}_{0,3} $$

et il suffit d'exprimer cette condition pour
un syst. des deux gén. $l_0, l_1$ de $\hat{\Pi}_{0,3}$.
Pour $x = l_0$ la relation (20) s'écrit
$$ (21) \quad u(l'_0) = int(\lambda^{-1} \hat{\varphi}(g_0)) (l'_0{}^\rho) $$

ce qui, comparé avec (17), donne
$$ int(g_0^{-1} \lambda^{-1} \hat{\varphi}(g_0)) (l'_0{}^\rho) = l'_0{}^\rho \quad \text{donc} $$
$$ g_0^{-1} \lambda^{-1} \hat{\varphi}(g_0) = l_0^{\prime p} \quad (p \in \hat{\mathbb{Z}}) $$

ou encore
$$ (22) \quad \lambda = \hat{\varphi}(g_0) l_0^{\prime -p} g_0^{-1} = \left[ \hat{\varphi}(g_0 l_0^{-p}) g_0^{-1} \right] $$
[MARGIN: cf plus bas (24 bis)]

Pour $x = l_1$ la relation (20) s'écrit

\newpage

%% ===== Page 62 =====

(23)
$$u(l'_1) = int(\lambda^{-1} \hat{\varphi}(g_1)) (l'_1)^p \qquad (g_1 = \sigma_\infty(g_0) = \sigma_0 \rho \sigma_0)$$

[MARGIN: Inutile ?]
[MARGIN: car $\sigma_\infty^p = \dots$ on s'en doutait un peu ! $(Z/nZ)$ [unclear]]

ou encore
$$\sigma_\infty = int(\lambda^{-1} \hat{\varphi}(g_1)) (\sigma_\infty)$$

ou encore $\sigma_\infty$ commute à :
$$\lambda^{-1} \hat{\varphi}(g_1) \overset{cf(22)}{=} g_0 {l'_0}^\beta \hat{\varphi}(\underbrace{g_0^{-1} \sigma_\infty(g_0)}_{h_0^{-1}}) = g_0 {l'_0}^\beta \hat{\varphi}(h_0^{-1})$$

ou encore à l'inverse $\hat{\varphi}(h_0) {l'_0}^{-\beta} g_0^{-1}$

ce qui s'écrit aussi
$$\sigma_\infty \left( \hat{\varphi}(h_0) {l'_0}^{-\beta} g_0^{-1} \right) = \hat{\varphi}(h_0) {l'_0}^{-\beta} g_0^{-1}$$

$$\sigma_\infty(\hat{\varphi}(h_0)) . \underbrace{\varepsilon_1^{-\beta}}_{(\rho \sigma_\infty)^{-\beta}} . \sigma_\infty(g_0)^{-1} = \hat{\varphi}(h_0) {l'_0}^{-\beta} g_0^{-1}$$

$$\sigma_\infty(\hat{\varphi}(h_0)) (\rho \sigma_\infty)^{-\beta} \cdot h_0 = \hat{\varphi}(h_0) \varepsilon_0^{-\beta}$$

[MARGIN: NB. Si $g_0$ est remplacé par $h_0^{-1}g_0$, le nouveau $h_0$ soit $h'_0 = h_0^{-1}$, la rel. (24) de $h_0$ donne la même rel. (25) (en sens inverse), de sorte que ... on a donc le choix du relèvement $g_0$ de ... pour expliciter (25) [unclear]]

[MARGIN: NB. Cette équation ne dépend pas du choix de $g_0$ (24 bis)]

(24)
$$\varepsilon_1^{-\beta} h_0 \varepsilon_0^\beta = \sigma_\infty(\hat{\varphi}(h_0))^{-1} \hat{\varphi}(h_0)$$

où $h_0$ donné par (+) en fonction de $g_0$
$$ \begin{cases} \varepsilon_0 = l'_0 = \sigma_0 \rho & (\varepsilon_0^2 = l_0) \\ \varepsilon_1 = \rho \sigma_\infty = \sigma_0 \rho & (\varepsilon_1^2 = l_1) \end{cases} $$

Posons, pour $\mu \in \widehat{G}_{0,3}^+$ ,

(25)
$$\tilde{\mu} = \sigma_\infty(\mu) \mu = \sigma_\infty \mu^{-1} \sigma_\infty \mu = [\sigma_\infty, \mu]$$

(26) (de sorte que $\tilde{h}_0 = \tilde{g}_0$) la formule (24) prend la forme

(27)
$$\varepsilon_1^{-\beta} \tilde{g}_0 \varepsilon_0^\beta = (\hat{\varphi}(\tilde{g}_0))^\sim$$

\newpage

%% ===== Page 63 =====

320

En résumé :

Proposition. Considérons une extension ext. $u$
d'un groupe à centre $\Pi_{0,3}$, réalisé par
[MARGIN: $\wedge$ mod $l_0$, ] $g_0 \in \widehat{\Pi}_{0,3}$, (et $p \in \hat{\Z}^*$ et $\gamma \in \hat{\Z}^*$ univoquement
déterminés par l'action de $g_0$ et $p$) par les
formules (11), au moyen d'un autom. effectif
(noté encore $u$)
qui commute à $\sigma_{\infty}$ — les données
$g_0$, $p$, (et $\gamma$ et $u$) assujetties à vérifier les
conditions (9) et (10) ci-dessus, plus la
(On prolonge $u$ en un autom. de $\widehat{S}_{0,3}^+$ par $int(u)$, cf (13))
condition d'engendrement (11 bis). Pour que
l'hom. extérieur $\hat{\varphi} : \widehat{\Pi}_{0,3} \to \widehat{S}_{0,3}^+$
défini par (18) commute à l'action
ext. de $u$, il faut et il suffit que l'on
puisse trouver un $\beta \in \hat{\Z}$, tel qu'on
ait (24), ou ce qui revient au même,
[MARGIN: mieux : 24 bis]
(27), (pour $\beta$ fixé)
Cette condition équivaut aussi à la formule
(20) pour $x=l_{\infty}$ ($\lambda$ étant tjrs donné par (22)),
c.à.d. à la formule
(28) $$u(l'_{\infty}) = int(\lambda^{-1}\hat{\varphi}(g_{\infty}))(l_{\infty}^{' \beta})$$ , où
(29) $$g_{\infty} = g_0 l_0^{-\beta} \hat{\varphi}^{-1}(h_0)$$ (et $\lambda$ défini par (22))

\newpage

%% ===== Page 64 =====

formule qui précise donc la formule
présomptive (28). Elle s'écrit aussi, plus
explicitement
$$ \rho^\gamma \rho^{-1} = int(g_0 \varepsilon_0^\beta \hat{\varphi}(l_0^\gamma p^{-1}(h_0))) (p^{-\gamma} p) $$
i.e. (en posant $\alpha = \beta + \gamma$ et explicitant $g$)

[MARGIN: $h_0^\gamma l_0^\gamma \rho$
NB On prouve l'existence de $\beta \in \hat{\mathbb{Z}}$ vérifiant (24), i.e. d'un $\alpha \in \hat{\mathbb{Z}}$ vérifiant (30).]

(30) $$ h_0^\gamma l_0^\gamma \rho^\gamma = int( \varepsilon_0^\alpha \hat{\varphi}(p^{-1}(h_0))) (p^{-\gamma}) $$

L'équivalence de (24) et (30) peut
paraître un peu étrange, du fait qu'
dans (24) l'exposant $p$ n'intervient pas
explicitement, mais seulement $h_0$ ,
tandis que la formule (30) fait intervenir
la valeur de $p \in \hat{\mathbb{Z}}^*$ mod 3 (qui donne
$\rho^p = \rho$ ou $\rho^p = \rho^{-1}$ ). Mais en tout état
de cause, qu'on soit sur la forme
(24) ou sur la forme (30), il semble
bien que la condition en question soit
hautement non tautologique !

Pour bien faire, il faudrait expliciter
explicitement les conditions supplémentaires
du § 39 en fonction de $g_0, p$.

\newpage

%% ===== Page 65 =====

321

Je voudrais comparer les autom. ext. de $\hat{\Pi}_{0,3}$ , et ceux de $\hat{\mathfrak{S}}_{0,3}^+$ , tenant compte de la suite exacte

(31) $$1 \to \hat{\Pi}_{0,3} \to \hat{\mathfrak{S}}_{0,3}^+ \to \mathfrak{S}_3 \to 1$$

Notons [unclear: crossed out] que tt autom. ext. $u$ de $\hat{\mathfrak{S}}_{0,3}^+$ qui invarie le s.-groupe distingué $\hat{\Pi}_{0,3}$ , et induit dans le quotient $\mathfrak{S}_3$ un autom. extérieur trivial, peut s'écrire comme déduit d'un autom. $u_0$ de $\hat{\mathfrak{S}}_{0,3}^+$ stabilisant $\hat{\Pi}_{0,3}$ , et induisant l'identité sur le quotient $\mathfrak{S}_3$ - 

et $u_0$ est donc déterminé modulo autom. int $(g)$, i.e. $g \in \hat{\mathfrak{S}}_{0,3}^+$ - tel que son image dans $\mathfrak{S}_3$ soit central, i.e. $=1$, 

i.e. $g \in \hat{\Pi}_{0,3}$ . Donc

(32) Groupe des automorphismes ext. de $\hat{\mathfrak{S}}_{0,3}^+$ qui induisent dans $\mathfrak{S}_3$ un autom. ext. trivial (donc invarie $\hat{\Pi}_{0,3}$ ) $\simeq$ Gr. d'autom. de $\hat{\mathfrak{S}}_{0,3}^+$ qui induisent l'identité dans le quotient $\mathfrak{S}_3$ mod le s.-groupe des autom. int. par des $g \in \hat{\Pi}_{0,3}$

\newpage

%% ===== Page 66 =====

Un tel autom $u_0$ est-il connu par sa restrict.
à $\hat{\Pi}_{0,3}$ ? En d'autres termes, un autom.
de $\widehat{\mathfrak{G}}_{0,3}^+$ (ou $= \operatorname{Sl}(2, \mathbb{Z})^\wedge$) qui
induit l'identité sur $\hat{\Pi}_{0,3}$ et sur $\widehat{\mathfrak{G}}_3$,
est-il trivial ? Il suffit de prouver que sa
restriction à l'image inverse de $\sigma_0 \in \widehat{\mathfrak{G}}_3$
est trivial (par symétrie en remplaçant $\sigma_0$ par $\sigma_0, \sigma_1, \sigma_2$,
car ces [unclear] engendrent $\widehat{\mathfrak{G}}_{0,3}^+$. On
aura $f(\sigma_0) = \sigma_0 g$ , avec $g \in \hat{\Pi}_{0,3}$ [ et
on doit avoir $\sigma_0 g \sigma_0 g = 1$ i.e.
$g . \sigma_0(g) = 1$   ]
de plus on devra avoir $int(\sigma_0 g) | \hat{\Pi}_{0,3}$
$=$ [unclear] $int(\sigma_0) | \hat{\Pi}_{0,3}$ , d'où
$int(g) | \hat{\Pi}_{0,3} = id$ i.e. $g \in Centre(\hat{\Pi}_{0,3})$ i.e. $g=1$ ,
OK. Mais : quelle condition un autom.
de $\hat{\Pi}_{0,3}$ peut se prolonger-t-il en un
autom. de $\widehat{\mathfrak{G}}_{0,3}^+$ qui induise l'identité sur
$\widehat{\mathfrak{G}}_3$ ? On voit aussitôt qu'il faut pour
cela que l'hom. ext. de $\hat{\Pi}_{0,3}$ qu'il
induit commute à l'action de $\widehat{\mathfrak{G}}_3$.

\newpage

%% ===== Page 67 =====

332

C'est suffisant, car soit $\Gamma$ ( $\subset \operatorname{Aut}\, ext(\hat{\Pi}_{0,3})$ ) le groupe des autom.
ext. de $\hat{\Pi}_{0,3}$ qui ont cette propriété (on ne
se préoccupe pas pour le moment de la
structure locale de $\hat{\Pi}_{0,3}$), on a
$\Gamma \cap \hat{\mathfrak{S}}_3 = (1)$ (on identifie $\hat{\mathfrak{S}}_3$ à un
[MARGIN: ceci est évident] sous-groupe de $\operatorname{Aut}\, ext(\hat{\Pi}_{0,3})$) car le centre
de $\hat{\mathfrak{S}}_3$ est trivial, donc
$$ \hat{\mathfrak{S}}_3 \times \Gamma \hookrightarrow \operatorname{Aut}\, ext(\hat{\Pi}_{0,3}) \ , $$
et on a donc une structure d'extension
(33) $$ 1 \to \hat{\Pi}_{0,3} \to \mathcal{E} \to \hat{\mathfrak{S}}_3 \times \Gamma \to 1 $$
(car le centre de $\hat{\Pi}_{0,3}$ est trivial), ou encore
une structure d'extension
(34) $$ 1 \to \hat{\mathfrak{S}}_{0,3}^+ \to \mathcal{E} \to \Gamma \to 1 $$
qui définit une action extérieure de $\Gamma$
sur $\hat{\mathfrak{S}}_{0,3}^+$ qui est l'action nat.
cherchée. Ainsi

Prop. Le groupe des autom. ext. de $\hat{\Pi}_{0,3}$ qui
[MARGIN: Mais tout autom. de $\hat{\mathfrak{S}}_3$ est intérieur] normalisent $\hat{\mathfrak{S}}_3$ est isomorph. au groupe
des autom. ext. de $\hat{\mathfrak{S}}_{0,3}^+$ qui induisent
sur $\hat{\mathfrak{S}}_3$ l'autom. ext. trivial.

66

\newpage

%% ===== Page 68 =====

Corollaire : Le groupe des autom. de $\hat{\Pi}_{0,3}$ qui
commutent aux extérieurs dans $\mathfrak{S}_3$, est isom.
au gr. des autom. de $\hat{\mathfrak{S}}_{0,3}^+$ qui
induisent sur $\mathfrak{S}_3$ l'autom. trivial (par
restriction de ces derniers).

(35)
\begin{tikzcd}
1 \ar[r] & \hat{\Pi}_{0,3} \ar[r] \ar[d, equal] & \operatorname{Aut}(\hat{\Pi}_{0,3} | \substack{ \text{comm.} \\ \text{à ext.} \\ \mathfrak{S}_3 }) \ar[r] \ar[d, "\simeq \text{ resp. } \hat{\Pi}_{0,3}"] & Aut ext(\hat{\Pi}_{0,3} | = \text{ dans } \mathfrak{S}_3) \ar[r] \ar[d, "\simeq"] & 1 \\
1 \ar[r] & \hat{\Pi}_{0,3} \ar[r] & \underbrace{\operatorname{Aut}(\hat{\mathfrak{S}}_{0,3}^+ | \substack{ \text{id. sur le} \\ \text{quot. } \mathfrak{S}_3 })}_{\mathcal{E}'} \ar[r] & Aut ext(\hat{\mathfrak{S}}_{0,3}^+ | \substack{ \text{triv.} \\ \text{sur quot.} \\ \mathfrak{S}_3 }) \ar[r] & 1
\end{tikzcd}

Cette extension $\mathcal{E}'$ de $\Gamma$ par $\hat{\Pi}_{0,3}$ est
en fait scindée, les [unclear] opèrent [unclear]
de $\hat{\Pi}_{0,3}$ qui commutent aux
[unclear] de $\hat{\mathfrak{S}}_{0,3}^+$ qui [unclear]
(par ex. : $\sigma_{\infty L}$ et [unclear] $\hat{\Pi}_{0,3}$)

(36)
$$ \begin{cases} u(\sigma_{\infty L}) = \sigma_{\infty L} \\ u(p) = g \cdot p \end{cases} \qquad (g \in \hat{\Pi}_{0,3}) $$

[MARGIN: provient au fait qu'il y a exact. un elt d'ordre 2 au dessus de $\sigma_\infty$...]

\newpage

%% ===== Page 69 =====

323

[unclear] C'est à dire que
pour $\hat{\mathbb{B}}_{0,3}^+$ est défini par les gén. $\sigma_\infty, \rho$
avec les relations

(37) $$ \sigma_\infty^2 = \rho^3 = 1 $$

[MARGIN: De plus il faut pour que $\sigma_\infty$ et $g\rho$ engendrent [unclear] : $\sigma_0 = (g\rho)^3 = 1 \dots$]

trouve la condition

(38) $$ (g\rho)^3 = 1 \quad i.e. \quad g \rho(g) \rho^2(g) = 1 $$

Corollaire Soit $u \in \mathcal{E}^1$ un autom. de $\hat{\mathbb{B}}_{0,3}^+$ qui induit l'identité sur $\mathbb{B}_3$.
Conditions équivalentes :
a) $\exists p \in \hat{\mathbb{Z}}^*$ tel que $u(l_0')$ conjugué dans $\hat{\mathbb{B}}_{0,3}^+$ à $l_0'^p$ (à $p$ près [unclear] unique, il est même conjugué par un élém de $\hat{\Pi}_{0,3}$)
b) $\exists p \in \hat{\mathbb{Z}}^*$ tel que $u(l_0)$ soit conjugué dans $\hat{\Pi}_{0,3}$ à $l_0^p$ (à $p$ près c'est unique)
c) $u_0 = u | \hat{\Pi}_{0,3}$ respecte la structure [unclear] de $\hat{\Pi}_{0,3}$.
De plus, le $p$ dans a) et b) est le même.

Dém Si $u(l_0') = int(g_0)(l_0'^p)$, avec $g_0 \in \hat{\mathbb{B}}_{0,3}^+$, 
élève [unclear] $u(l_0) = int(g_0)(l_0^p)$,
cel n'entraîne que l'image $\bar{g}_0$ de $g_0$ dans $\mathbb{B}_3$ 
(se réduit à $1$ ou $\bar{\sigma}_0$, do [unclear] 
[unclear] b), si [unclear] on sait que 
$g_0 = h \sigma_0$ (où $h \in \hat{\Pi}_{0,3}$) on a $u(l_0) = int(h) \sigma_0 (l_0^p)$
[unclear] $\sigma_0(l_0) = \dots$ (on a dans $\hat{\Pi}_{0,3}$ $\sigma_0 l_0 \sigma_0^{-1} = l_\infty$)
$u(l_0) = int(h a)(l_0^p)$, d'où b) [unclear] $= p$.

\newpage

%% ===== Page 70 =====

Comme on sait l'unicité de $p$ dans b),
[unclear]

(39) $$u(l_0) = int(g_0)(l_0^p) \quad (g_0 \in \hat{\Pi}_{0,3})$$
on déduit, écrivant
[crossed out: $\rho u_0 \rho^{-1}(l_0)$] $\rho u_0 \rho^{-1} = int(a) u_0 \quad (a \in \hat{\Pi}_{0,3})$

$$u_0(l_1) = int(a) \underbrace{\rho \underbrace{u_0 \underbrace{\rho^{-1} (l_1)}_{l_0}}_{int(g_0)(l_0^p)}}_{int(\rho(g_0))(l_1^p)} = int(a \rho(g_0)) (l_1^p)$$
i.e.
(40) $$u_0(l_1) = int(g_1)(l_1^p) \quad \dots \quad g_1 = a \rho(g_0)$$
et de même
$$u_0(l_\infty) = int(g_\infty)(l_\infty^p) \quad \dots \quad g_\infty = a \rho(g_1)$$
(41) $$= a \rho(a \rho(g_0))$$

donc b) $\implies$ c). Enfin prouvons c) $\implies$ a).

On a par hypothèse
$$u(l_0) = int(g) l_0^p$$
$$u(\varepsilon_0)^2 = int(g)((\varepsilon_0^p)^2) = (int(g) \varepsilon_0^p)^2$$

Je dis qu'on a alors
$$u(\varepsilon_0) = int(g) (\varepsilon_0^p)$$

Il suffit de [MARGIN: vérifier] sur le s-g de [crossed out: $\hat{\Gamma}_{0,3}^+$] [unclear] $l_0$ où $\varepsilon_0$ \dots OK

\newpage

%% ===== Page 71 =====

324

On doit donc [unclear: avoir]

Désormais nous nous bornerons aux - -

(eff. ou extérieurs) de $\hat{\Pi}_{0,3}$ qui respectent la str. de

la pointe, et coïncident sur $\mathcal{C}_3$ (entièrement)

(eff. ou ext.),

- qui correspondent aux autom. (des

$\widehat{\mathcal{C}}_{0,3}^+$ qui induisent l'identité sur $\mathcal{C}_3$,

et satisfont a) ci-dessus. Ils sont

les autom. extérieurs de cette str. donnés

correspondant à des autom. effectifs

de $\hat{\Pi}_{0,3}$ - - - invariant

par $\sigma_\infty$, et sont déterminés par un

$g \in \hat{\Pi}_{0,3}$ satisfaisant à (38), et tel que

de plus il existe pour $p \in \hat{\mathbb{Z}}^\times$, $g_0 \in \hat{\Pi}_{0,3}$ satisfaisant

(39) $u(l_1') = int(g_0) (l_1^{' p})$ i.e. $u(l_0) = int(g_0) (l_0^p)$

(où $g_0$ est déterminé mod $L_\infty = l_0^{\hat{\mathbb{Z}}}$ - ). Dans

ce calcul fait par ailleurs on trouve qu'on

doit avoir

(40) $g = \sigma_\infty(g_0) l_1^r \rho(g_0)^{-1}$ où $r \in \hat{\mathbb{Z}}$ bien

[MARGIN: déterminé (indépen-
dant du choix de $g_0$)]

i.e.

(41) $\sigma_\infty(g_0)^{-1} g \rho(g_0) \in L_1 = Cent_{\hat{\Pi}_{0,3}}(L_1)$

\newpage

%% ===== Page 72 =====

En effet, on a par hyp. sur $g_0$ (cf (39))

(42) $$u(l_0) = int(g_0)(l_0)$$

d'où

$$u \sigma_\infty(l_0) = u(l_1) = int(g_1)(l_1)$$
$||$ (car $u \sigma_\infty = \sigma_\infty u$)
$$\sigma_\infty u (l_0) \underset{(42)}{=} \sigma_\infty (int(g_0)(l_0)) = int(\sigma_\infty(g_0)) (\sigma_\infty(l_0)) = int(\sigma_\infty(g_0))(l_1)$$

d'où $$int (\sigma_\infty(g_0)^{-1} g_1) (l_1) = l_1$$

donc $$\sigma_\infty(g_0)^{-1} g_1 \in N_{\hat{\Pi}_{0,3}}(L_1) = L_1$$
[MARGIN: $g \rho(g_0)$]

d'où $$\sigma_\infty(g_0)^{-1} g \rho(g_0) = l_1^\gamma \quad (\gamma \in \hat{\mathbb{Z}})$$

D'où (40). On notera, pour se rassurer, que
l'élément $g$ défini au lemme [unclear: n'en de $g_0$]
ne dépend en effet pas du $g_0$ mod $L_0$. [unclear: Soit]

Ainsi ; la donnée d'un $u$ [unclear: autom. ext.]
de $\hat{\Pi}_{0,3}$ [unclear: extérieur à] $\hat{\mathcal{C}}_3$, ou encore
un autom. ext. de $\hat{\mathcal{C}}_{0,3}^+$ induisant l'identité
dans $\hat{\mathcal{C}}_3$ et transformant la classe de conjugaison
générée par $s_\infty$ [unclear: l'aut] $L_1' = l_0^{\hat{\mathbb{Z}}} = \ell_0^{\hat{\mathbb{Z}}} \hat{\rho}(L_0)$,
en encore [unclear: elle] d'un hom. [unclear: ext.]
$u$ de $\hat{\Pi}_{0,3}$ normalisant $\sigma_\infty$, et [unclear: ext.]
[unclear: on aura] d'un hom. $u$ de $\hat{\mathcal{C}}_{0,3}^+$ qui
fixe $\sigma_\infty$, et qui fixe [unclear] mod $\hat{\Pi}_{0,3}$

Spécialement : la donnée de $g_0 \in \hat{\Pi}_{0,3}$ mod $L_0$,
et de $\gamma \in \hat{\mathbb{Z}}$, satisfaisant la relation

\newpage

%% ===== Page 73 =====

(38) $g \rho g \rho g^2 \rho = 1$, où $g$ est donné par (40), et où
$u$ s'exprime en termes de $g$ de $g_0, \gamma$ par
(36)
$$
\begin{cases}
u(\sigma_{\infty}) = \sigma_{\infty} \\
u(g_1) = g\rho \qquad (= \sigma_{\infty}(g_0) l_1^\gamma \rho (g_0)^{-1} \rho)
\end{cases}
$$

la relation (42), s'écrit alors, [unclear] $L_0$ [unclear]
[unclear: $l_0 = l_0^{2\rho} = \dots$]
$(u(\sigma_{\infty}\rho))^3 \in int(g_0)(L_0')$
$\sigma_{\infty} g\rho = g_0 \sigma_{\infty} l_1^\gamma \rho g_0^{-1} = int(g_0) (\sigma_{\infty} l_1^\gamma \rho)$

donc (42) s'écrit
$(\sigma_{\infty} l_1^\gamma \rho)^2 \in L_0$

NB [unclear]
$(int(g_0)^{-1})(L_0) = (\sigma_{\infty} l_1^\gamma \rho)^2 = l_0^{2\gamma+1}$

soit $l_0^{2\gamma+1} \in L_0$

de [unclear] $u(L_0) =$
la relation (42), - la relation [unclear] plus forte
(43) $u(L_0') \subset int(g_0)(L_0')$

s'écrit
$int(g_0^{-1}) u(\sigma_{\infty}\rho) \in L_0'$
$\sigma_{\infty} g\rho = g_0 \sigma_{\infty} l_1^\gamma \rho g_0^{-1}$
$\sigma_{\infty} l_1^\gamma \rho = \sigma_{\infty} l_1^\gamma \sigma_{\infty} \sigma_{\infty} \rho$
$= \sigma_{\infty}(l_1)^\gamma \sigma_{\infty} \rho = l_0^{2\gamma} \varepsilon_0$

NB on a
$int(g_0^{-1}) u(l_0') = \varepsilon_0^{2\gamma+1}$

(44) $u(l_0') = u(\varepsilon_0) = int(g_0) (\varepsilon_0^{2\gamma+1})$

\newpage

%% ===== Page 74 =====

(44) $$u(l'_0) = int(g_0)(l_0^{\prime p})$$ avec $$p = 2\gamma+1$$

[unclear: d.i-(39), d'où on observe ...]
$$u(l_0) = int(q_0)(l_0^p)$$

[MARGIN: NB] On notera que si $p \in \hat{\mathbb{Z}}^*$, alors $p-1 \in 2\hat{\mathbb{Z}}$
i.e. $\frac{p-1}{2} \in \hat{\mathbb{Z}}$, car il suffit de vérifier la
relation analogue pour un $p \in \mathbb{Z}_l^*$, pour
$p-1 \in 2\mathbb{Z}_l$, or si $l \neq 2 \to 2\mathbb{Z}_l = \mathbb{Z}_l$ OK,
reste le cas $l=2$, $p \in \mathbb{Z}_2^*$ d'où l'image de $p$
dans $\mathbb{Z}/4\mathbb{Z}$ est $1$ ou $3$ donc
$p-1 \equiv 0$ ou $2 \mod 4$ donc pas de pb $\frac{p-1}{2} \in \mathbb{Z}_2$, OK.

Il reste donc à exprimer la condition
(38), et on l'écrit
$$\sigma_\infty(q_0) l_1^\gamma \rho(q_0)^\gamma \rho^{-1}(q_0) l_0^\gamma \rho(q_0) \rho^2(\sigma_\infty(q_0)) l_0^\gamma q_0^{-1} = 1$$
soit, en appliquant $int(q_0^{-1})$ et posant
(cf (8) )
(45) $$q_0^{-1} \sigma_\infty(q_0) = h_0^*$$  i.e.  $$h_0 = \sigma_\infty(q_0)^{-1} q_0$$
$$h_0^{*-1} l_1^\gamma \rho(h_0^*) l_0^\gamma \rho^2(h_0^*) l_0^\gamma = 1$$
ou encore, en appliquant $Int(l_0^\gamma)$
(46) $$(h_0^{*-1} l_1^\gamma) . \rho(l_1^\gamma h_0^{*-1}) . \rho^2(l_0^\gamma h_0^{*-1}) = 1 \quad (\text{où } \gamma = \frac{p-1}{2})$$
ce qui n'est autre que (10). Quant à la
condition 17), on voit qu'elle est [unclear: constamment]
vérifiée.

\newpage

%% ===== Page 75 =====

Résumons :
32 bis

Théorème Un autom. ext. de $\hat{\Pi}_{0,3}$, commutant à $\hat{\mathfrak{S}}_3$, se relève de façon unique en un autom. du groupe lui-même $\hat{\Pi}_{0,3}$, commutant à $\sigma_\infty$ (et commutant à $\rho$ mod autom. intérieurs), ou encore comme restriction d'un autom. $u$ de $\hat{\mathfrak{G}}_{0,3}^+$ qui fixe $\sigma_\infty$, conserve $\rho$ mod $\hat{\Pi}_{0,3}$ (i.e. $\bar{\rho}$ dans $\mathfrak{S}_3$), et qui conserve la classe de conjugaison du ss-groupe
[MARGIN: fermé]
$L_0^+ = \hat{\mathbb{Z}}(l_0)$ engendré par $l_0' = \tau_0 = \sigma_\infty \rho$. (ou encore du sous-groupe $L_0 = (L_0^+)'$ engendré par $l_0$). La donnée d'un tel $u$ équivaut à celle de son multiplicateur $p \in \hat{\mathbb{Z}}^*$, et d'un $g_0 \in \hat{\Pi}_{0,3} / L_0$ satisfaisant les conditions suivantes

a) Posant $h_0 = \sigma_\infty(g_0)^{-1} g_0 = [\sigma_\infty, g_0^{-1}] \in \hat{\Pi}_{0,3}$
(47) $\lbrace \gamma = \frac{p-1}{2} \in \hat{\mathbb{Z}}$
on a la relation (46) ci-dessus (invariante par remplacement de $g_0$ par $g_0 l_0^\lambda$, $\lambda \in \hat{\mathbb{Z}}$)
(48) $(l_0^\gamma \tilde{h}_0^{-1}) \rho (l_0^\gamma \tilde{h}_0^{-1}) \rho^2 (l_0^\gamma \tilde{h}_0^{-1}) = 1$

[MARGIN:
i.e.
[unclear: $\rho_0(g \rho^{-1} g^{-1}) = 1$]
i.e.
$g = [unclear]$]
b) $\sigma_\infty$ et $g \rho$ engendrent top. $\hat{\mathfrak{G}}_{0,3}^+$, et $g$ est donné en fonction de $g_0$ par
$[(40)]$ $g = \sigma_\infty(g_0) l_1^\gamma \rho(g_0)^{-1} = \sigma_\infty(g_0 l_0^\gamma) \rho(g_0)^{-1}$
$(g \rho)^3 = 1$ i.e. :

[MARGIN: $(\rho g \rho^{-1}) (\rho^2 g \rho^{-2}) = \rho g \rho^2 g \rho$]

NB La relation (46) équivaut à : (38) $g \rho(g) \rho^2(g) = 1$, exprimée directement en termes de $g_0, \gamma = \frac{p-1}{2}$, vu les relations génératrices $\sigma_\infty^2 = (g \rho)^3 = 1$.
(36) L'autom. $u$ est donné en termes de $g_0$ et de $p$ par

\newpage

%% ===== Page 76 =====

$$[(36)] \qquad \begin{cases} u(\sigma_\infty) = \sigma_\infty \\ u(\rho) = g \rho \end{cases} \qquad \text{où } g \text{ donné par } [(40)]$$

On a alors

$$(49) \qquad \begin{cases} u(l_0) = int(g_0) l_0^p \\ u(l_1) = int(g_1) l_1^p = int(\sigma_\infty(g_0))(l_1^p) \\ u(l_\infty) = int(g_\infty) (l_\infty^p) \end{cases}$$

$$g_1 = g \rho(g_0) = \sigma_\infty(g_0) l_1^\gamma = \sigma_\infty(g_0 l_0^\gamma)$$
$$g_\infty = g \rho(g_1) = g \rho(g) \rho^2(g_0) = \rho^2(g^{-1} g_0) = \rho(g_0 l_0^\gamma) h_0$$

[Le texte suivant est barré dans le manuscrit]

Corollaire. Pour que l'hom. surj. (18)
$$\hat{\varphi} : \hat{\Pi}_{0,3} \to \widehat{\mathcal{G}}_{0,3}^+$$
[MARGIN: cf page plus haut]

défini [unclear: par les relations] $l_0 \mapsto l'_0 = \sigma_\rho, l_1 \mapsto l'_1 = \sigma_\infty, l_\infty \mapsto l'_\infty = \rho^{-1}$,
se prolonge en un hom. $f$, il faut et il suffit qu'on puisse trouver $\beta \in \hat{\mathbb{Z}}$ tel qu'on ait
$$[(27)] \qquad \varepsilon_1^{-\beta} h_0 \varepsilon_0^\beta = \sigma_\infty(\hat{\varphi}(h_0)) \hat{\varphi}(h_0)$$
[MARGIN: $\varphi$ remplace $u$ de $Aut...$ déterminé plus haut]

(ce prolongement est alors unique), et il existe $\alpha \in \hat{\mathbb{Z}}$ tel qu'on ait
$$[(30)] \qquad h_0^{-1} l_1^\gamma \rho = int(\varepsilon_0^\alpha \hat{\varphi}(\hat{\varphi}^{-1}(h_0))) (\rho^p)$$
(où nécessairement $\alpha = \beta \gamma$).

\newpage

%% ===== Page 77 =====

Scolie. Dans le th. 2, $g \in \widehat{\Pi}_{0,3}$ (écrit $u$) s'exprime par l'intermédiaire de $g_0$ et $\gamma = \frac{p-1}{2}$, se calcule ($g_0$ défini mod $L_0$)
- la condition (46) était en fait une condition qui s'exprime directement en termes de $g_0$, ou mieux encore, en termes de

(50) $\xi = l_0^\gamma h_0^{-1} \in \widehat{\Pi}_{0,3}$ (d'où $h_0 = \xi^{-1} l_0^\gamma$)

comme
(51) $\xi \rho(\xi) \rho^2(\xi) = 1$,

qui équivaut à
(52) $(\xi \rho)^3 = 1$.

Les solutions de l'équation (51), ou ce qui revient au même de (52), correspondent aux scindages de l'extension de $\mathbb{Z}/3\mathbb{Z}$ par $\widehat{\Pi}_{0,3}$ engendrée par $\rho \in \mathfrak{S}_3$ , bien i.e. aux éléments d'ordre 3 de $\widehat{\mathfrak{S}}_{0,3}^+$ au dessus de $\rho$.
Sur l'ensemble des scindages resp. de ces éléments $\xi \rho$, $\widehat{\Pi}_{0,3}$ opère par autom. intérieurs :

(53) $int(\lambda)(\xi \rho) = \lambda \xi \rho \lambda^{-1} = \lambda \xi \rho \lambda^{-1} \rho^{-1} \rho = (\lambda \xi \rho(\lambda)^{-1}) \rho$

i.e. il opère sur l'ensemble des solutions de (51) par

(54) $\lambda \cdot \xi = \lambda \xi \rho(\lambda)^{-1}$

[MARGIN: NB $\lambda \mapsto \lambda \cdot \xi$ injectif]

\newpage

%% ===== Page 78 =====

Admettons (quitte à le vérifier par la suite)
qu'il y a exactement deux classes de conjugaison :

celle de $\{1\}$, correspondant aux scindages tautologiques,
donnés [unclear: donc par des] $\zeta$ de la forme
(55) $$ \zeta = \lambda \rho(\lambda)^{-1} \quad (\lambda \in \hat{\Pi}_{0,3}) $$ ,

et d'autre part la classe de conjugaison
des rel. d'ordre 3 $\sigma_\infty \hat{\rho}^{-1} \sigma_\infty^{-1} = \sigma_\infty \rho^{-1} \sigma_\infty$ de
$\hat{\rho}$ , correspondant à
$$ \sigma_\infty \rho^{-1} \sigma_\infty = \zeta_0 \rho $$ , avec
$$ \zeta_0 = \sigma_\infty \rho^{-1} \sigma_\infty \rho^{-1} = \sigma_\infty (\rho^{-1} \sigma_\infty) (\rho^{-1} \sigma_\infty) \sigma_\infty $$
$$ = \sigma_\infty \varepsilon_0^{-2} \sigma_\infty = \overbrace{\sigma_\infty^{-1}(l_0^{-1})}^{\varepsilon_0^{-1}} = l_1^{-1} $$

(et on a bien
$l_1^{-1} \rho(l_1^{-1}) \rho^2(l_1^{-1}) = l_1^{-1} l_\infty^{-1} l_0^{-1} = (l_0 l_\infty l_1)^{-1} = 1$ ) ,

cette orbite correspondant aux solutions de
la forme
(56) $$ \zeta = \lambda l_0^{-1} \rho(\lambda)^{-1} \quad (\lambda \in \hat{\Pi}_{0,3}) $$

[on notera que $l_0^{-1}$ et $l_1^{-1}$ sont sur la même
orbite $l_1^{-1} = l_0 (l_0^{-1}) \rho(l_0)^{-1}$ ]

i.e. [unclear]

A) Cas des $\zeta$ de la forme (55), d'où
(57) $$ h_0 = \rho(\lambda) \lambda^{-1} l_0^\gamma $$

\newpage

%% ===== Page 79 =====

328

Il reste donc à résoudre l'équation en $g_0$ (47)

(58) $$\sigma_\infty(g_0)^{-1} g_0 = \underbrace{\rho(\lambda) \lambda^{-1} l_0^{-\gamma}}_{h_0}$$

Pour un $h_0$ fixé, il y a au plus un seul $g_0$ qui satisfait à cette condition, et un tel $g_0$ 
[MARGIN: car commutant de $\sigma_\infty = \{1\}$ dans $\hat{\Pi}_{0,3}$]
existe ssi $\sigma_\infty(h_0) h_0 = 1$, ce qui s'écrit,

[rayé: $\sigma_\infty(h_0) h_0 = 1$ ou encore $h_0 \sigma_\infty(h_0) = 1$]

[rayé: (59) $\sigma_\infty(\rho(\lambda)) \sigma_\infty(\lambda^{-1}) l_1^{-\gamma} \rho(\lambda) \lambda^{-1} l_0^{-\gamma} = 1$]

qui s'écrit aussi, en reprenant l'expression $h_0 = \zeta^{-1} l_0^{-\gamma}$ :

[rayé: $\sigma_\infty(\zeta^{-1}) l_1^{-\gamma} \zeta^{-1} l_0^{-\gamma} = 1$ i.e.]

(60) $$\sigma_\infty(\zeta) = l_1^{-\gamma} \zeta^{-1} l_0^{-\gamma}$$

En résumé

Corollaire Les [unclear: u qui induisent pour la] classe de $\hat{\Pi}_{0,3}$-conjugaison de $\rho$ sont en correspondance 1-1 aux systèmes $(\lambda, \gamma) \in \hat{\Pi}_{0,3} \times \hat{\mathbb{Z}}^*$, tels que en posant $\zeta = \lambda \rho(\lambda)^{-1}$, on ait la relation (60).

On notera que pour $\lambda$ (donc $\zeta$) fixé, il existe au plus un $\gamma \in \hat{\mathbb{Z}}$ tel que la relation (60) soit valable, donc $\gamma$ est déterminé par $\lambda$ ; mieux : la condition implique que l'élément $2\gamma+1$ de $\hat{\mathbb{Z}}$ soit dans $\hat{\mathbb{Z}}^*$ (i.e. $\exists \gamma' \in \hat{\mathbb{Z}}^*$ / $2\gamma\gamma' + \gamma + \gamma' = 0$)

\newpage

%% ===== Page 80 =====

NB. Si $\lambda$ est remplacé par $\lambda' = l_0^\alpha \lambda$,
$\xi = \lambda \rho(\lambda)^{-1}$ est remplacé par $\xi' = l_0^\alpha \xi \rho(l_0)^{-\alpha}$,
et $h_0 = \xi^{-1} l_0^\gamma$ est remplacé par
$h'_0 = \rho(l_0)^\alpha h_0 l_0^{-\alpha}$, enfin $g_0$ est remplacé
par $g'_0 = g_0 l_0^{-\alpha}$, OK. (L'équation (60) pour $\xi$
et pour $\xi'$ est équivalente)

Finalement, pour [unclear: se dégager] on ne change
pas la classe de $\hat{\Pi}_{0,3}$-conj. de $\rho$, la
démarche la plus claire pour déterminer
les u en question me semble celle-ci :

1) On choisit un $\xi \in \hat{\Pi}_{0,3}$ satisfaisant à l'équation
(A) $\xi \rho(\xi) \rho^2(\xi) = 1$
(on le prendra soit sous la forme (55), soit
sous la forme (56) :
(B) $\xi = \lambda \rho(\lambda)^{-1}$ ou $\xi = \lambda l_0^\alpha \rho(\lambda)^{-1}$ , avec $\lambda \in \hat{\Pi}_{0,3}$)

2) On choisit $\gamma \in \hat{\mathbb{Z}}$ tel que
(C) $\sigma_\infty(\xi) = l_1^\gamma \xi^{-1} l_0^{-\gamma}$
[MARGIN: ici, $\xi$ est d'ailleurs déterminé en fonction de $\lambda$ par (B)]
ce qui détermine $\gamma$ de façon unique. (On
devra d'ailleurs se borner aux $\xi$ pour
lesquels un tel $\gamma$ existe, et tel de plus
que $p = 2\gamma + 1$ soit $\in \hat{\mathbb{Z}}^*$).

\newpage

%% ===== Page 81 =====

329

3°) On construit alors 
$$h_0 = \xi^{-1} l_0^\gamma \in \hat{\Pi}_{0,3}$$
La relation de 2°) assure que $h_0 \sigma_\infty(h_0) = 1$ i.e.
que $h_0$ peut se mettre sous la forme
$$h_0 = \sigma_\infty(g_0)^{-1} g_0$$
avec $g_0 \in \hat{\Pi}_{0,3}$.

[MARGIN: NB Il faudrait s'assurer par contre que la [unclear] pour la seule [unclear] l'équation $h_0 \sigma_\infty(h_0) = 1$ soit satisfait par [unclear], i.e. que (A) soit vérifiée !]

4°) On peut maintenant poser
$$g = \sigma_\infty(g_0) l_1^\gamma \rho(g_0)^{-1} = (g_0 l_0^{-\gamma}) \xi \rho(g_0 l_0^{-\gamma})^{-1}$$
définissant ce qu'il faut pour définir
$$u(\sigma_\infty) = \sigma_\infty$$
$$u(\rho) = g \rho$$

Ici, la seule équation délicate, pour laquelle l'existence [unclear: (non l'unicité, on l'a vu)] d'une solution, est problématique, est l'équation C (en $\gamma \in \hat{\mathbb{Z}}$) - plus la condition annexe que $2\gamma + 1 \in \hat{\mathbb{Z}}^*$. C'est donc une condition sur $\lambda \in \hat{\Pi}_{0,3}$ - ou plutôt deux types de conditions, suivant que $\xi$ est défini en termes de $\lambda$ par l'une ou l'autre formule (B).

Je prévois qu'il est dans l'une ...

\newpage

%% ===== Page 82 =====

l'autre des formules (B), suivant
que $p \pmod 3$ est congru à 1, ou à -1
(NB. il ne doit pas être congru à 0,
puisque $p \pmod 3 \in (\mathbb{Z}/3\mathbb{Z})^*$ ) i.e. suivant
que $2\gamma \equiv 0 \text{ ou } -2 \quad (3)$ i.e.
$\gamma \equiv 0 \text{ ou } -1 \quad (3)$ (NB. il ne faut pas
que $\gamma$ soit $\equiv 1 \quad (3)$, sinon $2\gamma+1 = p$
serait congru à $0 \pmod 3$). C'est
en tous cas ce qui se passe dans le cas où
$u \in \Gamma_{\mathbb{Q}_0}$.

Pour [unclear: trancher], réduisons le $\hat{\Gamma}_{0,3}^+$
mod les équations $l_0^3=1$ i.e. $(\sigma_\infty \rho)^6 = 1$
$$l_0^3 = l_1^3 = l_\infty^3 = 1$$
de sorte que $\hat{\Gamma}_{0,3}^+$ se [unclear] par une
extension de $\mathfrak{S}_3$ par le groupe
défini par les générateurs $l_0, l_1, l_\infty$
satisfaisant
$$l_0^3 = l_1^3 = l_\infty^3 = (l_\infty l_1 l_0) = 1$$
- ou mieux, utilisons
$$\hat{\Gamma}_{0,3}^+ \simeq \operatorname{Sl}(2, \mathbb{Z})^\wedge$$
et réduisons - [unclear: au vu] de -
$$\operatorname{Sl}(2, \hat{\mathbb{Z}})' \simeq \operatorname{Sl}(2, \mathbb{Z}/2\mathbb{Z}) \times \operatorname{Sl}(2, \mathbb{Z}/3\mathbb{Z})'.$$

\newpage

%% ===== Page 83 =====

Mais pour faire ces calculs, il faut pouvoir
identifier $\sigma_{\infty}, \rho$ comme éléments de
$\operatorname{Sl}(2, \hat{\mathbb{Z}})$ - explicitation que nous ferons
tantôt.

Revenant à (24 bis), et y posant
$h_0 = \xi^{-1} l_0^\gamma$, il vient

(61) $$q_0 = \hat{\rho}(\xi) \varepsilon_0^{\gamma - \beta}$$

On trouve donc que se donner
$(\xi, q_0, \gamma, \beta)$, avec $\xi, q_0 \in \widehat{\Pi}_{0,3}$ et $\gamma, \beta \in \hat{\mathbb{Z}}$
satisfaisant

(62) $$ \begin{cases} \xi \rho(\xi) \rho^2(\xi) = 1 \\ \sigma_{\infty}(q_0)^{-1} q_0 = l_0^\gamma \\ q_0 = \hat{\rho}(\xi) \varepsilon_0^{\gamma - \beta} \\ 2\gamma + 1 \in \hat{\mathbb{Z}}^* \end{cases} $$

revient, par élimination de $q_0$, à se donner des
systèmes $(\xi, r, s)$ ($r = \gamma - \beta$, $s = \gamma + \beta$)

(63) $$ \begin{cases} \xi \rho(\xi) \rho^2(\xi) = 1 \quad (\text{où } \xi = \lambda \rho(\lambda)^{-1} \text{ avec } \lambda = l_0^\alpha \rho(l_0)^{-\alpha} ) \\ \sigma_{\infty}(\hat{\rho}(\xi))^{-1} \hat{\rho}(\xi) = \varepsilon_1^{-r} \xi^{-s} \varepsilon_0^s \\ r + s + 1 \in \hat{\mathbb{Z}}^* \end{cases} $$

Ici, $\xi$ est déterminé mod. transformation
$\lambda \mapsto \lambda l_0^\alpha$ qui sur $\xi$ donne $\xi' = l_0^\alpha \xi l_1^{-\alpha}$.

\newpage

%% ===== Page 84 =====

## 42 Récapitulation sur le groupe

$$(\underline{\mathbb{G}}_{0,3})^\sim \simeq \Pi_1(U_{0,3}^{an}, \vec{v}_1 ; -\vec{j}) (= (\tilde{\Pi}'_{0,3} \simeq \Pi_1(\mathbb{P}^{1an} \setminus \{0\} + \{1\} + \{\infty\}), \vec{v}_1 ; -\vec{j}))$$
$$\simeq (\operatorname{Gl}(2, \mathbb{Z})' = \underline{G}\ell(2, \mathbb{Z}) / \pm 1 \simeq \underline{G}P(1, \mathbb{Z})) \simeq (\underline{\Gamma}_{1,1}' \simeq \underline{\Gamma}_{1,1}^\sim)$$
et sur les automorphismes "de lacets" de $\underline{\mathbb{G}}_{0,3}^\sim$ et de $\tilde{\Pi}_{0,3}$
(Formulaire fondamental).

[MARGIN: On a défini $\underline{\mathbb{G}}_{0,3}^\sim$ (ou $\tilde{\Pi}_{0,3}$) comme un groupe quotient différent d'un [unclear] (En fait, ils diffèrent par un automorphisme [unclear] au-dessus de [unclear]...)]

I) Le groupe $\underline{\mathbb{G}}_{0,3}^\sim$ est présenté comme extension :

$$(1) \quad 1 \to \tilde{\Pi}_{0,3} \to \underline{\mathbb{G}}_{0,3}^\sim \to \underline{\mathbb{C}}_{0,3} \to 1$$
$$ \Pi_1(U_{0,3}^{an} ; -\vec{j}) \quad \quad \simeq \mathfrak{S}_3 \times \mathbb{Z}/2\mathbb{Z}$$

où $\mathfrak{S}_3 \times \mathbb{Z}/2\mathbb{Z}$ est considéré comme groupe
opérant sur $U_{0,3}^{an} = \mathbb{P}^{1an} \setminus \{0, 1, \infty\} = \mathbb{C} \setminus \{0,1\}$ ,
$\mathfrak{S}_3$ opérant par permutations -
- c'est le groupe des automorphismes conformes du plan qui stabilisent $\{0, 1, \infty\}$. Le groupe $\mathfrak{S}_3$ est le sous-groupe des automorphismes conformes, il s'identifie au groupe des permutations de $\{0, 1, \infty\}$, ces six éléments sont
a) $\sigma_0, \sigma_1, \sigma_\infty$ (transpositions), $\rho, \rho^{-1}$ (permutations circulaires)
donnés par les formules explicites

$$(2) \quad \begin{cases} \sigma_0(z) = \frac{z}{z-1}, \quad \sigma_1(z) = \frac{1}{z}, \quad \sigma_\infty(z) = 1-z \\ \rho(z) = \frac{1}{1-z}, \quad \rho^{-1}(z) = 1 - \frac{1}{z} = \frac{z-1}{z} \end{cases}$$

$$(3) \quad \begin{cases} \sigma_0 \sigma_1 = \sigma_1 \sigma_\infty = \sigma_\infty \sigma_0 = \rho \\ \rho \sigma_0 \rho^{-1} = \sigma_1, \quad \rho \sigma_1 \rho^{-1} = \sigma_\infty, \quad \rho \sigma_\infty \rho^{-1} = \sigma_0 \end{cases}$$

$$(4) \quad \sigma_0^2 = \sigma_1^2 = \sigma_\infty^2 = \rho^3 = 1$$

Le groupe $C_{0,3}' \simeq \mathbb{Z}/2\mathbb{Z} \subset \underline{\mathbb{C}}_{0,3}$ apparaît comme le centre,

\newpage

%% ===== Page 85 =====

il est engendré par la conjugaison complexe $\tau$
(5) $$ \tau(z) = \bar{z} $$
les cinq autres éléments de $\mathcal{T}_{0,3}$ sont
(6) $$ \tau \sigma_0 = \tau'_0 \quad , \quad \tau \sigma_1 = \tau'_1 \quad , \quad \tau \sigma_\infty = \tau'_\infty \quad , \quad \tau \rho \quad , \quad \tau \rho^{-1} = (\tau \rho)^{-1} , $$
les anti-involutions $\tau'_i$ ($i = 0, 1, \infty$) en particulier --
sont des involutions anti-holomorphes, données par
(7) $$ \tau'_0(z) = \frac{\bar{z}}{\bar{z}-1} \quad , \quad \tau'_1(z) = \frac{1}{\bar{z}} \quad , \quad \tau'_\infty(z) = 1-\bar{z} . $$

Les points fixes des autom. continus de $(\mathbb{P}^{1 \text{an}} \setminus \{0,1,\infty\})$
sont les suivants

(8) $$
\left\{
\begin{aligned}
& (\mathbb{P}^1_\mathbb{C})^{\sigma_0} = \{0, 2\} \quad , \quad (\mathbb{P}^1_\mathbb{C})^{\sigma_1} = \{1, -1\} \quad , \quad (\mathbb{P}^1_\mathbb{C})^{\sigma_\infty} = \{\infty, 1/2\} \\
& (\mathbb{P}^1_\mathbb{C})^\rho = (\mathbb{P}^1_\mathbb{C})^{\rho^{-1}} = \{-j, -j^2\} \qquad \text{où } j = \exp 2i\pi/3 = \frac{-1+i\sqrt{3}}{2} \text{ racine primitive 3ème de 1} \\
& \qquad \qquad \qquad \qquad \qquad \qquad \quad -j = \exp 2i\pi/6 = \frac{1+i\sqrt{3}}{2} \text{ racine primitive 6ème de 1}
\end{aligned}
\right.
$$

[MARGIN: NB $(\mathbb{P}^1_\mathbb{C})^\sigma \cap (\mathbb{P}^1_\mathbb{C})^{\sigma'} = \emptyset$]

(9) $$
\left\{
\begin{aligned}
& (\mathbb{P}^1_\mathbb{C})^{\tau} = \mathbb{P}^1_\mathbb{R} \\
& (\mathbb{P}^1_\mathbb{C})^{\tau'_0} = \{ z \mid z + \bar{z} = z \bar{z} \} = \left\{ x+iy \mid (x-1)^2 + y^2 = 1 \right\} \\
& (\mathbb{P}^1_\mathbb{C})^{\tau'_1} = \{ z \mid z \bar{z} = 1 \} = \{ x+iy \mid x^2+y^2=1 \} = U \\
& (\mathbb{P}^1_\mathbb{C})^{\tau'_\infty} = \{ z \mid z + \bar{z} = 1 \} = \{ x+iy \mid x = 1/2 \} = 1/2 + i\mathbb{R}
\end{aligned}
\right.
$$

L'image la plus suggestive de tout cela est la
sphère de Riemann, où $\tau, \tau'_0, \tau'_1, \tau'_\infty$ apparaissent comme
des réflexions par r. aux "équateurs" correspondants, qui
sont 4 "lignes" (cercles ou droites)
passant toutes par les pôles $-j, -j^2$, qui
divisent le dit [unclear] de l'équateur en six
segments circulaires, par les pts de
subdivision $0, 1/2, 1, 2, \infty, -1$ .

Notons aussi les formules
(10) $$
\left\{
\begin{aligned}
& \rho \tau'_0 \rho^{-1} = \tau'_1 \quad , \quad \rho \tau'_1 \rho^{-1} = \tau'_\infty \quad , \quad \rho \tau'_\infty \rho^{-1} = \tau'_0 \\
& \tau'_0 \tau'_1 = \tau'_1 \tau'_\infty = \tau'_\infty \tau'_0 = \rho
\end{aligned}
\right.
$$

\newpage

%% ===== Page 86 =====

392

Les quatre géodésiques (9) découpent la sphère en
12 triangles sphériques, qui forment des domaines
fondamentaux pour le groupe $\widetilde{\mathcal{C}}_{0,3} \simeq \mathfrak{S}_3 \times \mathbb{Z}/2$
d'ordre 12, et sur lesquels il opère de façon simple-
ment transitive. (Provenant [unclear] de la subdivision
barycentrique des doubles triangles sphériques
déterminés par l'équateur $\mathbb{P}^1_{\mathbb{R}}$ et ses trois pts de subdi-
vision $0, 1, \infty$). Prenons un de ces triangles sphériques,
disons $(-\bar{j}, 0, 1/2)$, les réflexions p.r. à ses trois côtés, i.e.
$\tau, \tau_0', \tau_\infty'$, forment un syst. de générateurs de $\widetilde{\mathcal{C}}_{0,3}$,
avec les relations

(11) $$ \tau^2 = \tau_0'^2 = \tau_\infty'^2 = (\underbrace{\tau \tau_0'}_{\sigma_0'})^2 = (\underbrace{\tau_0' \tau_\infty'}_{\rho^{-1}})^3 = (\underbrace{\tau_\infty' \tau}_{\sigma_\infty'})^2 = 1 $$

quant à $\mathcal{C}_{0,3} = \widetilde{\mathcal{C}}_{0,3}^+$, il est engendré par
$\tau \tau_0' = \sigma_0'$, $\tau_0' \tau_\infty' = \rho^{-1}$ et $\tau_\infty' \tau = \sigma_\infty'$, avec les relations

(12) $$ \sigma_0' \rho^{-1} \sigma_\infty' = \sigma_0'^2 = \rho^3 = \sigma_\infty'^2 = 1 $$

II) Le groupe $\Pi_{0,3}$ est engendré par $l_0, l_1, l_\infty$,
liés par la relation génératrice

(13) $$ l_\infty l_1 l_0 = 1 $$

Le sous-groupe d'indice 2 $\mathcal{C}_{0,3}' = (\widetilde{\mathcal{C}}_{0,3})_{-\bar{j}} \simeq \mathcal{C}_{0,3}$,
stabilisateur de $-\bar{j}$, formé des $(g, \tau^\varepsilon) \in \mathfrak{S}_3 \times \mathbb{Z}/2\mathbb{Z}$
de dét $\text{sgn}(g) = (-1)^\varepsilon$, opère sur $\Pi_{0,3}$. Ce
groupe est formé de $1, \tau_0', \tau_1', \tau_\infty', \rho, \rho^{-1}$, et leurs
opérations sur $\Pi_{0,3}$ sont notées respectivement $1, \tau_0', \tau_1', \tau_\infty', \rho, \rho^{-1}$ !

[MARGIN: NB $\tau_0', \tau_1', \tau_\infty'$ inversent (?) les $l_0, l_1, l_\infty$ respectivement d'où les relations $\tau_0' l_0 \tau_0'^{-1} = l_0^{-1}$ etc.]

\newpage

%% ===== Page 87 =====

On identifie ces éléments à des éléments du groupe $\operatorname{Aut}_{loc}(\Pi_{0,3}) \simeq \hat{\mathcal{G}}_{0,3}$
par $g(x) = gxg^{-1}$ ; et on trouve : pour $g \in \mathcal{C}_{0,3}'$, $x \in \Pi_{0,3}$, ...
$$ \rho(l_0) = \rho l_0 \rho^{-1} = l_1 \quad , \quad \rho(l_1) = \rho l_1 \rho^{-1} = l_\infty \quad , \quad \rho(l_\infty) = \rho l_\infty \rho^{-1} = l_0 $$

(14)
$$ \tau_0'(l_0) = l_0^{-1} \quad , \quad \tau_0'(l_1) = l_\infty^{-1} \quad , \quad \tau_0'(l_\infty) = l_1^{-1} $$
$$ \tau_1'(l_1) = l_1^{-1} \quad , \quad \tau_1'(l_0) = l_\infty^{-1} \quad , \quad \tau_1'(l_\infty) = l_0^{-1} $$
$$ \tau_\infty'(l_\infty) = l_\infty^{-1} \quad , \quad \tau_\infty'(l_0) = l_1^{-1} \quad , \quad \tau_\infty'(l_1) = l_0^{-1} $$

(15) [crossed out: Dans] $\tau_0'\tau_1' = \tau_1'\tau_\infty' = \tau_\infty'\tau_0' = \rho$ \quad (on a bien sûr $\tau_0'^2 = \tau_1'^2 = \tau_\infty'^2 = \rho^3 = 1$)

Le groupe $\mathcal{G}_{0,3}$ en déduit des éléments $\sigma_0, \sigma_1, \sigma_\infty$ au dessus de $\tau_0, \tau_1, \tau_\infty$, par

(15')
$$ \sigma_0(l_1) = l_\infty \quad , \quad \sigma_0(l_\infty) = l_1 \quad , \quad \sigma_0(l_0) = (l_1 l_\infty)^{-1} $$
$$ \sigma_1(l_\infty) = l_0 \quad , \quad \sigma_1(l_0) = l_\infty \quad , \quad \sigma_1(l_1) = (l_0 l_\infty)^{-1} $$
$$ \sigma_\infty(l_0) = l_1 \quad , \quad \sigma_\infty(l_1) = l_0 \quad , \quad \sigma_\infty(l_\infty) = (l_0 l_1)^{-1} $$

On a

(17)
$$ \rho \sigma_0 \rho^{-1} = \sigma_1 \quad , \quad \rho \sigma_1 \rho^{-1} = \sigma_\infty \quad , \quad \rho \sigma_\infty \rho^{-1} = \sigma_0 $$
$$ \rho \tau_0' \rho^{-1} = \tau_1' \quad , \quad \rho \tau_1' \rho^{-1} = \tau_\infty' \quad , \quad \rho \tau_\infty' \rho^{-1} = \tau_0' $$

[MARGIN: NB $\sigma_0, \sigma_1, \sigma_\infty \in \mathcal{G}_{0,3}^+$
$\tau_0', \tau_1', \tau_\infty' \in \mathcal{G}_{0,3}^-$
$\tau_0, \tau_1, \tau_\infty \in \mathcal{G}_{0,3}^-$]

(18') $\sigma_0 \sigma_1 = \rho l_0 = l_1 \rho \quad , \quad \sigma_1 \sigma_\infty = \rho l_1 = l_\infty \rho \quad , \quad \sigma_\infty \sigma_0 = \rho l_\infty = l_0 \rho^{-1}$

On a ainsi des éléments de $\mathcal{G}_{0,3}$ au dessus de $1, \sigma_0, \sigma_1, \sigma_\infty, \rho, \rho^{-1}, \tau_0', \tau_1', \tau_\infty'$, mais pas encore au dessus de $\tau \in C$ qui permet d'en avoir un sur $\tau \rho, \tau \rho^{-1}$), on posera $\tau_i = \sigma_i \tau_i' = \tau_i' \sigma_i$ (qui est sur $\tau$, de carré 1,

[MARGIN: l'appellation proviendra des conditions:]

(18) $\tau_i$ est caractérisé par la condition d'être sur $\tau$ et de commuter à $\sigma_i$ (ou encore, à $\tau_i'$)

(19) $[\tau_0, \sigma_0] = [\tau_1, \sigma_1] = [\tau_\infty, \sigma_\infty] = 1$ \quad $\tau_0 = \sigma_0 \tau_0' \quad , \quad \tau_1 = \sigma_1 \tau_1' \quad , \quad \tau_\infty = \sigma_\infty \tau_\infty'$

(20) $\tau_0^2 = \tau_1^2 = \tau_\infty^2 = 1$ \quad et \quad $\tau_0, \tau_1, \tau_\infty$ sur $\tau$

(21) $\tau_0 \tau_1 = l_0 \quad , \quad \tau_0 \tau_\infty = l_1 \quad , \quad \tau_1 \tau_0 = l_\infty$

Pour bien faire, il faudrait expliciter la description de Malgrange (faite dans [unclear: SGA]) de $\sigma_0, \sigma_1, \sigma_\infty, \rho, \tau_0', \tau_1', \tau_\infty', \tau_0, \tau_1, \tau_\infty$,

\newpage

%% ===== Page 88 =====

333

(22)
$$ \begin{cases}
(\sigma_0 \rho)^2 = (\rho \sigma_1)^2 = l_0 = \epsilon_0^2 , & \epsilon_0 = \sigma_0 \rho = \rho \sigma_1 \text{ [unclear]} \quad \text{(NB } \sigma_\infty = \rho \sigma_1 \rho^{-1} \text{)} \\
(\sigma_1 \rho)^2 = (\rho \sigma_\infty)^2 = l_1 = \epsilon_1^2 , & \epsilon_1 = \sigma_1 \rho = \rho \sigma_\infty \\
(\sigma_\infty \rho)^2 = (\rho \sigma_0)^2 = l_\infty = \epsilon_\infty^2 , & \epsilon_\infty = \sigma_\infty \rho = \rho \sigma_0
\end{cases} $$

L'élément $\epsilon_i \in \widehat{G}_{0,3}$ (en fait $\epsilon_i \in \widehat{G}_{0,3}^+$)
au dessus de $\sigma_i$, il est caractérisé par la condition
$\epsilon_i^2 = l_i$.

Dans le groupe $G_{0,3}$ d'ordre 12, il y a sept éléments
involutifs $\sigma_0, \sigma_1, \sigma_\infty, \tau_0', \tau_1', \tau_\infty'$, (et bien sûr l'identité),
d'ordre 3 $\rho, \rho^{-1}$ [unclear: et $\rho \tau, \rho^{-1} \tau$] [unclear: (ces plans en plus de l'unité)].
[unclear: Sauf les \dots \dots \dots]

[MARGIN: en effet $\rho \tau \rho = \dots$ [unclear]]

Nous allons [unclear: relever] la structure des
éléments d'ordre fini de $\widehat{G}_{0,3}$. On sait que
les possibilités, mod $\widehat{\Pi}_{0,3}$-conjugaison, de
tels [unclear: relèvements], correspondent aux
composantes connexes des pts fixes dans $M_{0,3}$ des
automorphismes correspondants. - pour $\sigma_0, \sigma_1, \sigma_\infty$,
$\tau_0', \tau_1', \tau_\infty'$ il n'y a qu'une possibilité mod conjugai-
son (on a choisi $\sigma_0, \sigma_1, \sigma_\infty, \tau_0, \tau_1, \tau_\infty$ de la façon
qui a semblé la plus commode pour les calculs),
pour $\tau$ il y a trois [unclear] ce sont les trois
choix $\tau_0, \tau_1, \tau_\infty$ pour $\rho, \rho^{-1}$ il y en a deux,
correspondants aux deux pôles $J, \bar{J}$ - nous avons
choisi celle qui correspond au pôle $J$. [unclear: Les rela-
tions correspondantes] se [unclear] d'ail-
leurs [unclear] pour $\tau_i$ :

\newpage

%% ===== Page 89 =====

[MARGIN: 
* NB
$\rho \rho_0 \rho^{-1} = \rho_1$
$\rho \rho_1 \rho^{-1} = \rho_\infty$
$\rho \rho_\infty \rho^{-1} = \rho_0$]

$$= l_0^{-1} \rho l_0 = l_1^{-1} \rho l_1 = l_\infty^{-1} \rho l_\infty$$

(23) $\rho_0 = \tau_0 \rho \tau_0^{-1}$, $\rho_1 = \tau_1 \rho \tau_1^{-1}$, $\rho_\infty = \tau_\infty \rho \tau_\infty^{-1}$ (sont d'ordre 3,
au dessus de $\bar{\rho}$ ds $\bar{\Pi}_{0,3}$ , $\Pi_{0,3}$-conjugués entre eux *,
mais non $\Pi_{0,3}$-conjugués à $\rho$.

Notons aussi que par les formules usuelles de la
théorie générale et du fait que les composantes con-
nexes des pts fixes qui apparaissent ici sont
$1$-connexes :

(24) $\Pi_{0,3}^{\sigma_i} = \Pi_{0,3}^\rho = \Pi_{0,3}^{\tau_i} = \Pi_{0,3}^{\tau'_i} = \{1\}$

Si on regarde les pts critiques dans $\mathcal{U}_{0,3}$ ,
relatifs à l'opération de $\overline{G}_{0,3}$ , on trouve
deux orbites
les pts $(2, -1, 1/2)$ et $(-j, -\bar{j})$ , le
stabilisateur d'un point dans ces orbites est d'ordre
4 resp. 6 . Pour le pt 2 (opposé à 0) ,
ce stabilisateur est le groupe $D_2 \simeq \mathbb{Z}/2\mathbb{Z} \times \mathbb{Z}/2\mathbb{Z}$ ,
engendré par $\bar{\sigma}_0, \bar{\tau}_0$ (ou $\sigma_0, \tau_0$) — qui a "été
bel et bien relevé multiplicativement" dans $\widehat{G}_{0,3}$ .
$Et =$ , pour le point $-\bar{j}$ , c'est $\overline{G}'_3 \simeq D_3$ , qui a
été relevé également. Par la théorie générale
on sait que les classes de $\Pi_{0,3}$-conjugaison d'un
tel relèvement correspondent aux pts fixes
du groupe d'opérateurs (sur $\mathcal{U}_{0,3}$) ; pour $D_2$ il n'y a
qu'un seul point fixe, savoir 2 , pour
$G'_3 \simeq D_3$ il y en a deux, $-j$ et $-\bar{j}$ , - on avait choisi
$-\bar{j}$ . En prenant l'autre choix on conjugue

\newpage

%% ===== Page 90 =====

334

en conjuguant par l'une quelconque des $\tau_i$ 
$(\tau_0, \tau_1, \tau_\infty)$.

## III Relation avec le Gr. pro cartographique tri-graphe
[unclear: orienté], $\overline{\Pi}_{0,3} = \pi_1( (\mathbb{P}^1_{\mathbb{C}} \setminus \{0\}) + 2\{1\} + 3\{\infty\} ; -j )$ et
$= ( \Pi_{0,3} * \{1, \tau\} ) / ( l_1^2, l_\infty^3 )$ [unclear: avec $\tau^2=1$]

On a des générateurs involutifs $\sigma_0, \sigma_1, \sigma_2$ de $\overline{\Pi}_{0,3}$,
($\sigma_1$ [unclear: dans] un repère de centre [unclear], fixe [unclear] la "moitié" de
$\sigma_0$ et $\sigma_2$ [unclear: ont resp.] la "symétrie" d'un repère, de classe d'orientation [unclear: bilatère]
resp. transversale, $\sigma_1$ le passage au repère binaire "adjacent"
...) [unclear: avec] les relations génératrices

[MARGIN: on pourrait d'autres choix de $\sigma'_i$. On aurait un système de générateurs comme on l'a fait des $\tau_0, \tau_1, \tau_\infty$ (au lieu de $\sigma_0, \sigma_1, \sigma_2$)]

(25)
$$ \sigma_0^2 = \sigma_1^2 = \sigma_2^2 = (\sigma_0 \sigma_2)^2 = 1 $$
$$ (\sigma_0 \sigma_1)^3 = 1 $$

Le sous-groupe $\Pi'_{0,3}$ d'indice 2 (groupe cart. orienté) est engendré par
$$ \underbrace{\sigma_0 \sigma_1}_{l'_\infty} = f, \quad \underbrace{\sigma_1 \sigma_2}_{l'_0} = \rho_s, \quad \underbrace{\sigma_2 \sigma_0}_{l'_1} = \sigma, \quad \text{avec les} $$

relations génératrices :
(26)
$$ f \rho_s \sigma = 1 \qquad \sigma^2 = f^3 = 1 $$

C'est donc le quotient de $\Pi_{0,3}$ par $l_1^2, l_\infty^3$

$$ \Pi_{0,3} \longrightarrow \Pi'_{0,3} $$
(27)
$$ l_0 \longmapsto l'_0 = \rho_s $$
$$ l_1 \longmapsto l'_1 = \sigma $$
$$ l_\infty \longmapsto l'_\infty = f $$

Cet homomorphisme se prolonge en [unclear]

\newpage

%% ===== Page 91 =====

$$ (28) \quad \tilde{\Pi}_{0,3} = \Pi_{0,3} \cdot \{1, s\} \longrightarrow \tilde{\Pi}'_{0,3} $$
$$ \tau_0, \tau_1, \tau_\infty \longmapsto \bar{\sigma}_0, \bar{\sigma}_1, \bar{\sigma}_\infty $$

[MARGIN: Cet iso. s'en déduit en passant au quotient sur
$$ \tilde{\Pi}'_{0,3} / \{1, s\} \simeq \tilde{\mathcal{C}}_{0,3} / \pm 1 $$
i.e.
$$ \Pi'_{0,3} \simeq \mathcal{C}^+_{0,3} $$
donné d'ailleurs formellement par
$$ l_1^2 = l_\infty^3 = 1 \dots $$]

Ceci posé, on définit un iso. $\sim$

$$ (29) \quad \tilde{\Pi}'_{0,3} \stackrel{\sim}{\longrightarrow} \tilde{\mathcal{C}}_{0,3} $$
$$ (\bar{\sigma}_0, \bar{\sigma}_1, \bar{\sigma}_\infty) \longmapsto (\tau'_0, \tau'_1, \tau'_\infty) \quad \text{notés aussi } r_0, r_1, r_\infty $$
$$ \text{(réflexions fonda-} $$
$$ \text{mentales)} $$

qui induit un iso.

$$ (30) \quad \Pi'_{0,3} \stackrel{\sim}{\longrightarrow} \mathcal{C}^+_{0,3} $$
$$ \rho_s = l'_0 \longmapsto \varepsilon_0 = \tau_\infty' \tau_0' $$
$$ \sigma = l'_1 \longmapsto \sigma_\infty $$
$$ \rho_f = l'_\infty \longmapsto \rho^{-1} $$

[MARGIN: NB
$\tau_\infty' \tau_0' = \varepsilon_0 = l'_0$
$\tau_0' \tau_1' = \sigma_\infty = l'_1$
$\tau_1' \tau_\infty' = \rho^{-1} = l'_\infty$
i.e. $r_\infty r_0 = l'_0 (= \varepsilon_0)$
$r_0 r_1 = l'_1 (= \sigma_\infty)$
$r_1 r_\infty = l'_\infty (= \rho^{-1})$]

On trouve donc que $\tilde{\mathcal{C}}_{0,3}$ a des générateurs et relations
génératrices
$$ (31) \quad \tau_0'^2 = \tau_1'^2 = \tau_\infty'^2 = (\tau_\infty' \tau_0')^2 = (\underbrace{\tau_0' \tau_1'}_{\rho^{-1}})^3 = 1 $$

et $\mathcal{C}_{0,3}^+$ a des relations génératrices

$$ (32) \quad l'_0 l'_1 l'_\infty = l'_1{}^2 = l'_\infty{}^3 = 1 $$

ou aussi simplement, en éliminant $l'_0$ et
posant $l'_1 = \sigma_\infty$ et $l'_\infty = \rho$

$$ (33) \quad \sigma_\infty^2 = \rho^3 = 1 $$

Pour mémoriser, on exprime les éléments
remarquables de $\mathcal{C}_{0,3}^+$ en termes de $\sigma_\infty, \rho$ :

\newpage

%% ===== Page 92 =====

(34)
$$ \begin{cases}
\varepsilon_0 = \sigma_\infty \rho & l_0 = \sigma_\infty \rho \sigma_\infty \rho = (\sigma_\infty \rho)^2 \\
\varepsilon_1 = \rho \sigma_\infty & l_1 = \rho \sigma_\infty \rho \sigma_\infty = (\rho \sigma_\infty)^2 \\
\varepsilon_\infty = \rho^{-1} \sigma_\infty \rho^{-1} & l_\infty = \rho^{-1} \sigma_\infty \rho^{-1} \rho^{-1} \sigma_\infty \rho^{-1} = (\rho^{-1} \sigma_\infty \rho^{-1})^2
\end{cases} $$

(35)
$$ \begin{cases}
\sigma_0 = \rho \sigma_\infty \rho^{-1} & (= \varepsilon_1 \rho^{-1} = \rho \varepsilon_1^{-1}) \\
\sigma_1 = \rho^{-1} \sigma_\infty \rho & (= \rho^{-1} \varepsilon_0 = \varepsilon_\infty \rho^{-1})
\end{cases} $$

et dans $\widehat{G}_{0,3}$, en fonction de $\tau_\infty', \tau_0', \tau_\infty$
les expressions de $\tau_0, \tau_1, \tau_\infty'$ :
(36)
$$ \begin{cases}
\tau_0 = \rho \tau_\infty' \rho^{-1} = \tau_\infty' \tau_0' \tau_\infty' \tau_0' \tau_\infty' & (= r_0 r_1 r_\infty r_1 r_0) \\
\tau_1 = \rho^{-1} \tau_\infty' \rho = \tau_0' \tau_\infty' \tau_0' \tau_\infty' \tau_0' & (= r_1 r_0 r_\infty r_0 r_1) \\
\tau_\infty' = \rho \tau_0' \rho^{-1} = \tau_\infty' \tau_0' \tau_\infty' & (= r_0 r_1 r_0)
\end{cases} \qquad \begin{cases} \tau_\infty' = r_0 \\ \tau_0' = r_1 \\ \tau_\infty = r_\infty \end{cases} $$

## IV) Relations avec $\mathcal{G}l(2, \mathbb{Z}) = \operatorname{Gl}(2, \mathbb{Z})/\pm 1 \simeq \operatorname{Aut}(\mathbb{P}^1_{\mathbb{Z}})$.

Ce groupe opère sur le demi-plan de Poincaré
$$D = \{z \in \mathbb{C} \mid \Im z > 0\}$$

par
(37)
$$ \begin{cases}
u_g(z) = \frac{az+b}{cz+d} & \text{si } ad-bc = +1 \\
u_g(z) = \frac{a\bar{z}+b}{c\bar{z}+d} & \text{si } ad-bc = -1
\end{cases} \qquad g = \begin{pmatrix} a & b \\ c & d \end{pmatrix} \quad a,b,c,d \in \mathbb{Z}, \quad ad-bc \in \pm 1 $$

On a dans $D$ un domaine fondamental bien connu, qui correspond au triangle sphérique $(-J, 1, 0)$ pour l'homéomorphisme du demi-plan naturel à $\Delta$ envoyant $\infty$ à $1$, $\rho$ à $-J$, et le point $i$ à $0$. Ainsi [unclear: on a les symétries] par rapport aux trois côtés de ce triangle sont donnés par

(38)
$$ r_0(z) = -\bar{z}, \quad r_1(z) = 1-\bar{z}, \quad r_\infty(z) = 1/\bar{z} $$

Ils décrivent donc [unclear: le groupe de symétries]
et les produits deux à deux donnent les [unclear: rotations affines]

\newpage

%% ===== Page 93 =====

$$
(39) \quad \begin{cases}
\tau_\infty \tau_1(z) = l'_0(z) = z - 1 \\
\tau_0 \tau_\infty(z) = l'_1(z) = - \frac{1}{z} \\
\tau_1 \tau_0(z) = l'_\infty(z) = 1 - \frac{1}{z} = \frac{z-1}{z}
\end{cases}
$$

On trouve ainsi un isomorphisme

$$
(40) \quad \widehat{G}_{0,3} \xrightarrow{\sim} \operatorname{Gl}(2, \mathbb{Z})' \simeq \operatorname{Gl}(2, \mathbb{Z})/{\pm 1} \simeq GP(1, \mathbb{Z})
$$
$$
\begin{cases}
\tau_0 = \tau'_0 \longmapsto r_0 = \begin{pmatrix} 0 & 1 \\ 1 & 0 \end{pmatrix} \\
\tau_1 = \tau'_1 \longmapsto r_1 = \begin{pmatrix} -1 & 1 \\ 0 & 1 \end{pmatrix} & \text{[MARGIN: (il vaut mieux changer de signes cf plus loin)]} \\
\tau_\infty = \tau_\infty \longmapsto r_\infty = \begin{pmatrix} 1 & 0 \\ 0 & -1 \end{pmatrix}
\end{cases}
$$

et une présentation de $GP(1, \mathbb{Z})$ par

$$
(41) \quad r_0^2 = r_1^2 = r_\infty^2 = (r_0 r_\infty)^2 = (r_1 r_0)^3 = 1
$$

Cela donne par suite un isomorphisme

$$
(42) \quad \widehat{G}_{0,3}^+ \xrightarrow{\sim} \operatorname{Sl}(2, \mathbb{Z})' = \operatorname{Sl}(2, \mathbb{Z})/{\pm 1} \simeq GP(1, \mathbb{Z})^+
$$

$$
\begin{cases}
l'_1 = \sigma_\infty \longmapsto \begin{pmatrix} 0 & -1 \\ 1 & 0 \end{pmatrix} \\
l'_\infty = \rho^{-1} \longmapsto \begin{pmatrix} 1 & -1 \\ 1 & 0 \end{pmatrix}
\end{cases}
$$

$$
( \text{d'où} \quad \rho \longmapsto \begin{pmatrix} 0 & 1 \\ -1 & 1 \end{pmatrix} ) .
$$

On identifiera par la suite les éléments de $\widehat{G}_{0,3}$ avec leurs images dans $\operatorname{Gl}(2, \mathbb{Z})' \simeq GP(1, \mathbb{Z})$.
Notons la formule [unclear] qui est utile pour [unclear] une matrice $u = \begin{pmatrix} a & b \\ c & d \end{pmatrix} \in \operatorname{Gl}(2, \mathbb{Z})$, pour $\det u \in \mathbb{Z}^\times = \{ \pm 1 \}$ [unclear]...

\newpage

%% ===== Page 94 =====

[MARGIN: N.B. 2 : Si $u = \begin{pmatrix} a & b \\ c & d \end{pmatrix} \in GL_2(A)$, $A$ anneau [unclear: qcq], on a $\det(u) \cdot u^{-1} = \begin{pmatrix} d & -b \\ -c & a \end{pmatrix}$]

(43) $\quad u^{-1} = \begin{pmatrix} d & -b \\ -c & a \end{pmatrix} \quad :$

$\varepsilon_0 = \sigma_\infty \rho = \begin{pmatrix} 1 & -1 \\ 0 & 1 \end{pmatrix} \quad , \quad \varepsilon_0^{-1} = \begin{pmatrix} 1 & 1 \\ 0 & 1 \end{pmatrix}$

$\varepsilon_1 = \rho \sigma_\infty = \begin{pmatrix} 1 & 0 \\ 1 & 1 \end{pmatrix} \quad , \quad \varepsilon_1^{-1} = \begin{pmatrix} 1 & 0 \\ -1 & 1 \end{pmatrix}$

(44) $\quad \varepsilon_\infty = \rho^{-1} \sigma_\infty \rho^{-1} = \begin{pmatrix} 2 & -1 \\ 1 & 0 \end{pmatrix} \quad , \quad \varepsilon_\infty^{-1} = \begin{pmatrix} 0 & 1 \\ -1 & 2 \end{pmatrix}$

[MARGIN: cf "formulaire" pour liste des relations d'identités plus complète]

$l_0 = \varepsilon_0^2 = \begin{pmatrix} 1 & -2 \\ 0 & 1 \end{pmatrix} \qquad \qquad \qquad \qquad \sigma_0 = \begin{pmatrix} 1 & -1 \\ 2 & -1 \end{pmatrix}$

$l_1 = \varepsilon_1^2 = \begin{pmatrix} 1 & 0 \\ 2 & 1 \end{pmatrix} \qquad \qquad \qquad \qquad \sigma_1 = \begin{pmatrix} 1 & -2 \\ 1 & -1 \end{pmatrix}$

$l_\infty = \varepsilon_\infty^2 = \begin{pmatrix} 3 & -2 \\ 2 & -1 \end{pmatrix} \qquad \qquad \qquad \sigma_\infty = \begin{pmatrix} 0 & -1 \\ 1 & 0 \end{pmatrix}$

(45) $\quad \rho = \begin{pmatrix} 0 & 1 \\ -1 & 1 \end{pmatrix} \quad , \quad \rho^{-1} = \begin{pmatrix} 1 & -1 \\ 1 & 0 \end{pmatrix} \quad ,$

(46) $\quad l_\infty l_1 l_0 = \begin{pmatrix} -1 & 0 \\ 0 & -1 \end{pmatrix} \quad$ l'élément non trivial
du centre de $\operatorname{Sl}(2, \mathbb{Z})$
qui (heureusement !) donne
$1$ dans $\operatorname{Sl}(2, \mathbb{Z})'$

Le groupe $\Pi_{0,3}$ se récupère comme le noyau
de l'homomorphisme
$$\operatorname{Sl}(2, \mathbb{Z})' \to \underset{\simeq \mathfrak{S}_3}{\operatorname{Sl}(2, \mathbb{Z}/2\mathbb{Z})}$$

la suite exacte
(47) $\quad 1 \to \Pi_{0,3} \to \mathcal{C}_{0,3}^+ \to \underset{\simeq \mathfrak{S}_3}{\overline{\mathcal{C}}_{0,3}^+} \to 1$

se récupère comme

(48) $\quad 1 \to \underset{\simeq \Pi_{0,3}}{\pi(2)'} \to \operatorname{Sl}(2, \mathbb{Z})' \to \underset{\simeq \mathfrak{S}_3}{\operatorname{Sl}(2, \mathbb{Z}/2\mathbb{Z})} \to 1$

\newpage

%% ===== Page 95 =====

[unclear: Ca nous fourmille dans les mains de travailler]
dans $\operatorname{Sl}(2, \mathbb{Z})$ lui-même, extension centrale de
$\hat{C}_{0,3}$ - en fait, cette extension s'interprète
comme $\hat{C}_{1,1}$, mais nous y reviendrons
ultérieurement...

V) Résumons les avatars des divers groupes que nous
avons rencontrés jusqu'à présent

1) $\hat{C}_{0,3} \stackrel{dfn}{=} \operatorname{Aut}_{loc}(\Pi_{0,3})$ avec ss-groupe $\hat{C}_{0,3}^+$ (respectant orientation)

2) $\hat{C}_{0,4}' \simeq Autext_{loc}(\Pi_{0,4})'$ aut ext. de $\Pi_{0,4}$ respectant la structure de lacets et fixant la classe de lacets $L_{4}$ (et permutant $L_0, L_1, L_\infty$)
[unclear: conjecturé orienté] des comp. connexes de 

3) $\Lambda_{0,4}' / \Lambda_{0,4}^0$ groupe des comp. connexes de 
$\Lambda_{0,4}' \subset \operatorname{Aut}_{top}(U_{0,4}^{top})$ sous-groupe des autom. qui fixent le point "à l'infini" $a_{0,4} = 1/2$
avec $\Lambda_{0,4}^{\prime +} / \Lambda_{0,4}^0$ (respectant orientation)

4) $\pi_1((U_{0,3}^{top}, \hat{\mathcal{B}}_3), -J)$ (décrit à la J. Malgoire)
avec le sous-groupe $\pi_1((U_{0,3}^{top}, \hat{\mathcal{B}}_3^+), -J)$ ($B_3 = \hat{C}_{0,3}$)

[MARGIN: C'est l'utile groupe $\bar{\Pi}_{0,3}$ (à [unclear: l'équiv.] près), où $\bar{\Pi}_{0,3} = \Pi_{0,3} \cdot \{1, c\}$]
5) $\pi_1((\mathbb{P}_{\mathbb{C}}^1, \infty \cdot \{0\} + 2 \cdot \{1\} + 3 \cdot \{\infty\}), -J)$ (à la J. Malgoire)
avec le sous-groupe $\pi_1((\mathbb{P}_{\mathbb{C}}^1, \infty \cdot \{0\} + 2 \cdot \{1\} + 3 \cdot \{\infty\})^+, -J)$

6) Groupe auto-graphique triangulaire $\Gamma_{0,3}$
$\Gamma_{0,3} = \{ \tau_0, \tau_1, \tau_\infty \mid \tau_0^2 = \tau_1^2 = \tau_\infty^2 = (\tau_1 \tau_\infty)^2 = (\tau_\infty \tau_0)^3 = 1 \}$
(avec sous-groupe cart. triangulaire orienté
$\Gamma_{0,3}^+ = \{ \dots \mid \dots \}$ )
$\tau_0 \tau_1, \tau_0 \tau_\infty, \tau_1 \tau_\infty$

[MARGIN: NB $U_{0,3}^{conf} = (U_{0,3}, \Gamma_{0,3}^+)$ ce qui permet d'identifier 7) & 4)]
7) $\pi_1^{orb}(U_{0,3}^{conf}, pt(U_{0,3}, -J))$ avec sous-groupe $\pi_1(U_{0,3}^{an}, pt(U_{0,3}, -J))$
[unclear: sous-jacente la forme] universelle de type $(0,3)$
des revêtements universels de type $(0,3)$

8) $\hat{C}_{1,1}' = \hat{C}_{1,1} / \{\pm 1\}$ avec $\hat{C}_{1,1} \simeq Autext_{loc}(\Pi_{1,1})$
$\simeq \Lambda_{1,1} / \Lambda_{1,1}^0$
$(\Lambda_{1,1} \simeq$ autom. top. de $U_{1,1}$
avec sous groupe $\hat{C}_{1,1}^+ / \{\pm 1\}$)

\newpage

%% ===== Page 96 =====

337

9) $\pi_1(M_{1,1}^{conf}, pt)' \simeq \pi_1( ) / centre \{\pm 1\} \simeq \pi_1(M_{1,1}^{\prime conf}, pt)$

où $M_{1,1}$ est la multiplicité modulaire analytique [conforme] des courbes elliptiques (i.e. courbes [unclear] de type $(1,1)$)

$M_{1,1}'$ celle des courbes elliptiques analytiques "mod. symétrie" ; avec sous-groupe

$\pi_1(M_{1,1}^{an}, pt)' \simeq \pi_1( ) / centre \{\pm 1\} \simeq \pi_1(M_{1,1}^{\prime an}, pt)$

où le pt base est la courbe analytique [ell. conforme an.] , quotient de $\mathbb{C}$ par $\mathbb{Z} + \mathbb{Z}(-j)$ .

[MARGIN: NB les multiplicités modulaires sont $(\mathbb{D}, \operatorname{Gl}(2, \mathbb{Z}))$, $(\mathbb{D}, \operatorname{Sl}(2, \mathbb{Z}))$]

dans le cadre conforme resp analytique pour $M_{1,1}$ ,

et $(\mathbb{D}, \operatorname{Gl}(2, \mathbb{Z}))' = GPL(2, \mathbb{Z})$ resp. $(\mathbb{D}, \operatorname{Sl}(2, \mathbb{Z}))' = SPL(2, \mathbb{Z})^+$

dans le cadre conforme resp analytique pour $M_{1,1}'$ .

10) $\operatorname{Gl}(2, \mathbb{Z})' \simeq GPL(2, \mathbb{Z})$ , avec le sous-groupe $\operatorname{Sl}(2, \mathbb{Z})' \simeq SPL(2, \mathbb{Z})$

\newpage

%% ===== Page 97 =====

(VI) Les automorphismes à l'ext. de $\widehat{\Pi}_{0,3}$ et de $\widehat{\mathbb{C}}_{0,3}^+$

On a un diagramme commutatif canonique :

(49)
\begin{tikzcd}
\operatorname{Aut}\,ext(\widehat{\mathbb{C}}_{0,3}^+)' \ar[r, "\text{restriction}", "\sim"'] & \operatorname{Aut}\,ext^*(\widehat{\Pi}_{0,3})' \\
\operatorname{Aut}(\widehat{\mathbb{C}}_{0,3}^+, \sigma_\infty)' \ar[u, "\wr"] \ar[r, "\text{restriction}", "\sim"'] & \operatorname{Aut}(\widehat{\Pi}_{0,3}, \sigma_\infty) \ar[u, "\wr"']
\end{tikzcd}

où on a :
$\operatorname{Aut}\,ext(\widehat{\mathbb{C}}_{0,3}^+)' = \text{groupe des autom. ext. induisant l'identité sur le quotient } \mathfrak{S}_3 \simeq \widehat{\mathbb{C}}_{0,3}^+/\widehat{\Pi}_{0,3}$

$\operatorname{Aut}\,ext^*(\widehat{\Pi}_{0,3})' = \text{groupe des autom. extérieurs qui commutent à l'action nat. de } \mathfrak{S}_3$

$\operatorname{Aut}(\widehat{\mathbb{C}}_{0,3}^+, \sigma_\infty)' = \text{groupe des autom. de } \widehat{\mathbb{C}}_{0,3}^+ \text{ qui fixent } \sigma_\infty, \text{ et qui fixent } \rho \bmod \widehat{\Pi}_{0,3}$

$\operatorname{Aut}(\widehat{\Pi}_{0,3}, \sigma_\infty) = \text{groupe des aut. qui commutent à } \sigma_\infty \text{ et qui } [unclear] \text{ ext. } id$

Les $u \in \operatorname{Aut}(\widehat{\mathbb{C}}_{0,3}^+, \sigma_\infty)'$ sont de la forme

[MARGIN: On écrit les formules dans $\operatorname{Aut}(\widehat{\mathbb{C}}_{0,3}^+)' = \mathcal{G}$ groupe des [unclear: classes d'] autom. extérieurs qui ind. l'identité dans $\mathcal{G}/\widehat{\Pi}_{0,3} \simeq \mathfrak{S}_3$, et $\operatorname{Aut}(\widehat{\Pi}_{0,3})' = gr. des [unclear]$ ext. de $\dots$ NB. On admet ici que le centre de $\widehat{\mathbb{C}}_{0,3}^+$ est trivial.]

(50) 
$$
\begin{cases}
u(\sigma_\infty) = \sigma_\infty \\
u(\rho) = g \rho \qquad (i.e.\ u \rho u^{-1} = g)
\end{cases}
$$

où $g \in \widehat{\Pi}_{0,3}$ est assujetti à la condition de
$$ \boxed{g \rho g \rho g \rho = 1} $$
(51) exprimant la relation
(52) $$(g \rho)^3 = 1$$

On trouve ainsi une corresp. 1-1 entre les éléments de $\operatorname{Aut}(\widehat{\mathbb{C}}_{0,3}^+, \sigma_\infty)'$ et les $g \in \widehat{\Pi}_{0,3}$ satisfaisant (51), et tels de plus que

(53) $\sigma_\infty, g \rho \text{ avec les relations } \sigma_\infty^2 = (g \rho)^3 = 1 \text{ présentent } \widehat{\mathbb{C}}_{0,3}^+$

\newpage

%% ===== Page 98 =====

(que je vais appeler la "condition liminaire").

Pour [unclear: text barré]
Pour un $u \in \mathcal{G} = \operatorname{Aut}(\hat{C}_{0,3}^+)''$ (ne fixant pas
nécessairement $\varpi_\infty$), on a des conditions équivalentes

(54)
(i) $u(\varepsilon_0) \in$ conjugué de $\varepsilon_0^{\hat{\mathbb{Z}}} \simeq L'_0$ (dans $\hat{C}_{0,3}^+$)
(ii) $u(l_0) \in c$-conjugué de $l_0^{\hat{\mathbb{Z}}} = L_0$ (dans $\hat{\Pi}_{0,3}$)
(iii) $u(l_i)$ conjugué de $l_i$ pour $i=0, 1, \infty$
(iv) $\exists \rho \in \hat{\mathbb{Z}}^*$, $g_0 \in \hat{\Pi}_{0,3}$ \quad $u(\varepsilon_0) = int(g_0) \varepsilon_0^\rho$
(v) $\exists \rho \in \hat{\mathbb{Z}}^*$, $g_0 \in \hat{\Pi}_{0,3}$ \quad $u(l_0) = int(g_0)(l_0^\rho)$

On aura alors aussi

(55) $u(l_1) = int(g_1)(l_1^\rho)$ \quad $g_1 = g \rho(g_0)$
$u(l_\infty) = int(g_\infty)(l_\infty^\rho)$ \quad $g_\infty = g \rho(g_1) = g \rho(g) \rho^2(g_0)$

L'élément [unclear: text barré]
(55') $\rho \in \hat{\mathbb{Z}}^*$ \quad $\rho = \chi(u)$ \quad (où $u$ donné par (50)),
est bien déterminé par $u$, donc par $g$, on l'appelle
le multiplicateur de $u$. L'élément

(56) $g_0 \in \hat{\Pi}_{0,3}$
est déterminé modulo $L_0 = l_0^{\hat{\mathbb{Z}}}$ (changement
(57) $g_0 \mapsto g_0 l_0^\alpha$ \quad ).

[MARGIN: Notation:
$\operatorname{Aut}_{ext, loc}(\hat{C}_{0,3}^+) \subset \operatorname{Aut}_{ext}(\hat{C}_{0,3}^+)'$
$\operatorname{Aut}_{loc}(\hat{C}_{0,3}^+) \subset \operatorname{Aut}(\hat{C}_{0,3}^+)'$
$\operatorname{Aut}_{loc}(\hat{\Pi}_{0,3}, \varpi_\infty) \dots$]

Si les conditions [unclear: text barré] (54) sont satisfaites, on dit que
$u$ respecte la structure à lacets de $\hat{C}_{0,3}^+$ (ou de
$\hat{\Pi}_{0,3}$) - en sous-entendant par $\hat{C}_{0,3}^+$ avec
son "groupe d'inertie" $\varepsilon_0^{\hat{\mathbb{Z}}} \simeq L'_0$. On
a affaire par la suite à de tels autom.

[MARGIN: NB ces conditions ne dépendent que de l'autom. ext.
défini par $u$, de $\hat{\Pi}_{0,3}$ dans $\hat{C}_{0,3}^+$]

\newpage

%% ===== Page 99 =====

s'intéresse aux relèvements canoniques des
tels autom. extérieurs donnés par (49), (50),
en des [unclear: autom. effectifs].

La formule de définition (50) de $u$ en fait
intervenir (via $g_0 \in \widehat{\Pi}_{0,3}$) en principe $g_0 \pmod{L_0}$ dans $\widehat{\Pi}_{0,3}/L_0$
qui se déduit de $g$. Mais en pratique $g_0$ joue un
rôle aussi important que $g$, - on verra que dans
une large mesure, il détermine $g$ (formule (62)).

S'il n'y avait pas lieu d'introduire $g_0$, pour
exprimer les conditions (54), on pourrait aisé-
ment paramétrer les $g$ satisfaisant à (55) -
donc correspondants à des relèvements de
$\xi \in \widehat{G}_3$ en des éléments d'ordre 3 de $\widehat{\mathcal{B}}_{0,3}^+$ -
en notant qu'il y a exactement deux classes
de conjugaison de tels relèvements,
(correspondants aux pts fixes $-j, j$ de $\xi$
opérant sur $U_{0,3}^\sim$), qui se paramétrisent par

[MARGIN: ces représentations paramétriques sont univoques.]

(58') $$g = \lambda_0 \rho (\lambda_0^\xi)^{-1} \quad [= \lambda_0 \rho \lambda_0^{-1}]$$ (classe principale, correspondant au "pôle nord" $-j$)
$$g = \lambda_0' \rho^{-1} ({\lambda_0'}^\xi)^{-1}$$ (classe adjointe, corresp. au "pôle sud" $j$)

où $\lambda_0, \lambda_0' \in \widehat{\Pi}_{0,3}$. On verra plus bas que ces
cas correspondent respectivement aux cas

(59) \quad $p \equiv 1 \quad (3)$ \quad et \quad $p \equiv -1 \quad (3)$ .

\newpage

%% ===== Page 100 =====

339

Mais ici la situation se corse considérablement par la nécessité de tenir compte des conditions (54) (p. ex. sous la forme (iv)) - elles [unclear] d'introduire conjointement $g_0, p$. Il y a lieu alors d'introduire les quantités

(60) $$ \gamma = \frac{p-1}{2} \in \hat{\mathbb{Z}} $$
(NB comme $p \in \hat{\mathbb{Z}}^*$, sa classe dans $\hat{\mathbb{Z}}/2\hat{\mathbb{Z}}$ $(= \mathbb{Z}/2\mathbb{Z})$ est $1$, $\implies p-1 \in 2\hat{\mathbb{Z}}$, d'où la possibilité de trouver $\gamma$)

[MARGIN: $h_0$ relatif à l'autom. $\sigma_\infty$ d'ordre 2, joue un rôle symétrique pour $g_0$, au $h_\infty$ relatif à l'autom. $\rho$ d'ordre 3 - ceci justifie un peu nos notations.]

$$ h_0 = g_0 \sigma_\infty(g_0) $$

(que nous noterons, précédemment on avait pris $h_0 = \sigma_\infty(g_0) g_0$, dans - pour l'inversion p. r. aux notations précédentes). On vérifie alors que 

$g$ s'exprime à l'aide de $g_0$ et de $\gamma$ (donc de $p$) à l'aide de la formule (équivalente à (54) (iv) ou (v))

[MARGIN: unclear]

(62) $$ g = \sigma_\infty(g_0) h_0^\gamma \rho(g_0)^{-1} = (g_0 h_0^{-1}) (h_0^\gamma h_0) \rho(g_0 h_0^{-1})^{-1} $$

Dans le cas qui nous intéresse on le décrit par un couple $(g_0, \gamma) \quad (g_0 \in \widehat{\Pi}_{0,3}/L_0, \gamma \in \hat{\mathbb{Z}})$ satisfaisant les conditions

(63) $$ 2\gamma + 1 \in \hat{\mathbb{Z}}^* \quad (NB\ 2\gamma+1 = p = \chi(\dots)) $$

et la condition (51), où $g$ est exprimé en fonction de $g_0$ et de $\gamma$ par (62). Cette formule (51) équivaut à son tour à la relation

(64) $$ \xi \rho(\xi) \rho^2(\xi) = 1 $$

où 
(65) $$ \xi = h_0^\gamma h_0 \quad \text{avec } (66) \quad h_0 = g_0 \sigma_\infty(g_0) $$

[MARGIN: NB On a appelé ici $\xi$ ce qui était noté [unclear]]

\newpage

%% ===== Page 101 =====

Ainsi, les $u$ sont paramétrés par les couples
$(q_0 \in \hat{\Pi}_{0,3}/L_0, \gamma \in \hat{\mathbb{Z}})$, satisfaisant les conditions
(63) et (64), qui moyennant (65) s'écrit sous forme
explicite:

(66) $$(l_0^\gamma \underbrace{q_0^{-1} \sigma_\infty(q_0)}_{h_0}) \rho(l_0^\gamma \underbrace{q_0^{-1} \sigma_\infty(q_0)}_{h_0}) \rho^2(l_0^\gamma \underbrace{q_0^{-1} \sigma_\infty(q_0)}_{h_0}) = 1$$

[MARGIN: s'il y avait à choisir, [unclear: je préférerais] cette "équation en $h_0$" symétrique plutôt que l'"équation en $\xi$" (68)]

("équation en $q_0$" - par opposition avec l'"équation en $\xi$" plus bas). (De plus, il faut rajouter la condition linéaire, pour mémoire...)

On peut aussi considérer que $u$ (via $g$)
est défini par $(\xi \in \hat{\Pi}_{0,3}, \gamma \in \hat{\mathbb{Z}})$, qui
permettent en effet de retrouver
$$h_0 = l_0^{-\gamma} \xi$$
par (65), puis $q_0$ par (66):
$$q_0^{-1} \sigma_\infty(q_0) = h_0$$
qui détermine en effet $q_0$ de façon unique,
grâce au fait (qu'on admet provisoirement)
que $\hat{\Pi}_{0,3}^{\sigma_\infty} = \{1\}$. La possibilité de
trouver $q_0$ satisfaisant à cette condition
équivaut (en l'admettant encore ici) : la
condition

$$\sigma_\infty(h_0) h_0 = 1$$

(67)

qui, en termes de $\xi, \gamma$, s'écrit aussi sous la forme

\newpage

%% ===== Page 102 =====

[MARGIN: plutôt "l'équation en $h_0$"]

(68) $$ \sigma_{\infty}(\xi) = l_1^{\gamma} \xi^{-1} l_0^{\gamma} $$

Les équations (64), (68) sont appelées les "équations en $\xi$" - par opposition à "l'équation en $q_0$" (66), qui à première vue a l'avantage d'être une seule, et le désavantage d'être plus compliquée. En résumé, les $u$ qui nous intéressent (i.e. dans $\operatorname{Aut}((\hat{\Pi}_{0,3}^+, \sigma_{\infty})')$) sont donnés par les couples $(\xi \in \hat{\Pi}_{0,3}^+, \gamma \in \hat{\mathbb{Z}})$ satisfaisant les conditions :

a) $2\gamma + 1 \in \hat{\mathbb{Z}}^*$ (63)

b) $\xi \rho(\xi) \rho^2(\xi) = 1$ (64)

c) $\sigma_{\infty}(\xi) = l_1^{\gamma} \xi^{-1} l_0^{\gamma}$ (68)

d) Condition bilinéaire pour mémoire (53) (qui ne s'exprime apparemment pas directement de façon commode en termes de $\xi, \gamma$).

Pour que deux couples $(\xi, \gamma)$, $(\xi', \gamma')$ définissent le même $u$, il faut et il suffit qu'on ait :

(69) $\gamma = \gamma'$, $\xi' = l_0^{-\alpha} \xi l_1^{\alpha} \quad (\alpha \in \hat{\mathbb{Z}})$.

La condition (64) peut encore se remplacer par une représentation paramétrique

(70) $$ \begin{cases} \xi = \lambda \rho(\lambda)^{-1} \\ \xi = \lambda l_0^{-1} \rho(\lambda)^{-1} \end{cases} \quad \text{ou} $$

$\lambda \in \hat{\Pi}_{0,3}$, auquel cas l'équation (68)

\newpage

%% ===== Page 103 =====

peut s'interpréter comme une "équation en $\lambda$"
qui suivant les deux cas prend la forme
$$ (71) \quad \begin{cases} \sigma_\infty (\lambda \rho(\lambda)^{-1}) = l_1^\gamma ( \rho(\lambda)\lambda^{-1}) l_0^{-\gamma} \\ \text{ou} \\ \sigma_\infty (\lambda l_0^{-1} \rho(\lambda)^{-1}) = l_1^\gamma ( \rho(\lambda) l_0 \lambda^{-1} ) l_0^{-\gamma} \end{cases} $$

Pour que $(\lambda,\gamma)$ et $(\lambda',\gamma')$ définissent
le $= u$ , il faut et s. qu'on ait
$$ (72) \quad \lambda' = l_0^{-\alpha} \lambda \quad , \quad \gamma = \gamma' \ . $$

[MARGIN: NB L'expression (62) qui met en évidence l'étroite parenté de forme de ces 2 équations disparait dans la forme (73)]

Remarque 1) Au facteur "anodin" $l_0^\gamma$ près, qui
différencie $h_0$ et $\xi$ , (65) , on peut dire
que la recherche des $u \in \operatorname{Aut}_{loc}(\widehat{(\mathcal{C}_{0,3}^+, \sigma_\infty)})$
équivaut à celle des solutions simultanées
des deux équations de cocycles
$$ (73) \quad \begin{cases} \xi \rho(\xi) \rho^2(\xi) = 1 \\ h_0 \sigma_\infty(h_0) = 1 \end{cases} $$
relatifs respectivement aux groupes $\mathbb{Z}/3\mathbb{Z}$ , opérant
sur $\widehat{\Pi}_{0,3}$ via $\rho$ , et $\mathbb{Z}/2\mathbb{Z}$ , opérant via
$\sigma_\infty$ , dont les solutions [unclear] respectivement
$$ \begin{cases} \xi = \lambda \rho(\lambda)^{-1} \text{ ou } \quad (= \lambda l_0^{-1} \rho(\lambda)^{-1} \\ \quad \text{(suivant que } p \equiv 0 \ (3) \text{ ou } p \equiv 1 \ (3) \\ \quad \text{i.e. } p \equiv 1 \ (3) \text{ ou } p \equiv -1 \ (3)) \\ h_0 = g_0 \sigma_\infty(g_0)^{-1} \quad \text{([unclear: NB on pourrait]} \\ \quad \text{[unclear: aussi bien changer] } g_0 \text{ [unclear: en] } g_0^{-1} \dots ) \end{cases} $$

\newpage

%% ===== Page 104 =====

341

3) La relation entre le $\lambda_0$ dans la représentation paramétrique de $g$ (58), et le $\lambda$ dans celle de $\xi$ (74 a) est implicite dans la deuxième expression (62) de $g$, savoir

(75) $$ \lambda_0 = g_0 l_0^{-\gamma} \lambda $$

4) Notons ici [unclear: que] les quantités [unclear: précédemment considérées] (avec les indéterminations
$$g, \gamma, g_0, h_0, \xi, \lambda_0, \lambda$$
$$\in \quad \in \quad \in \quad \in \qquad \in$$
$$\hat{\mathbb{Z}}^\times \quad \hat{\mathbb{Z}} \quad \hat{\Pi}_{0,3} \quad \hat{\Pi}_{0,3}/L_0 \quad \hat{\Pi}_{0,3}$$

spécifiques à $g_0, h_0, \xi, \lambda$) peuvent être considérées, [unclear: pour un repérage] fixé, comme des fonctions sur le groupe 
$$Autext_{Loc}(\widehat{\mathfrak{G}}_{0,3}^+)' \simeq Autext_{Loc} (\hat{\Pi}_{0,3})'$$ (sous-groupe des autom. ext. continus du groupe $\hat{\Pi}_{0,3}$ qui commutent jusqu'à $\mathfrak{S}_3$), à valeurs dans les ensembles $\hat{\mathbb{Z}}^\times, \hat{\mathbb{Z}}, \hat{\Pi}_{0,3}, \hat{\Pi}_{0,3}/L_0, \, \cdot \, , \hat{\Pi}_{0,3}$.

En [unclear: particularisant], on trouve ainsi $8$ fonctions sur $\widehat{\Gamma}_{0,3}$ : à valeurs dans les espaces appropriés - [unclear: ... pour la connexion cyclotomique.]

[MARGIN: On verra [unclear: temps!] non ah ! que pour un [unclear: autom] $g$ et [unclear: par l'él...] $h_0, \xi, \lambda_0, \lambda$ seront définis de façon canon. et rigoureuses dans 2 paragraphes...]

\newpage

%% ===== Page 105 =====

Notons à ce propos les conditions équivalentes

(i) $p=1$
(ii) $\gamma=0$
(iii) $g = \sigma_{\infty}(g_0) p(g_0)^{-1}$
[MARGIN: (75 bis)] (iv) $h_0 = 1$
(v) $\sigma_{\infty}(\xi) = \xi^{-1}$ ("équation en $\xi$ symplectique !)
(vi) $h_0 p(h_0) p^2(h_0) = 1$

signifient toutes que $\chi(u)=1$. D'où la recherche
des $u \in \operatorname{Aut}_{ext}(\hat{\mathcal{G}}_{0,3}^+, \sigma_{\infty})'$ de multiplicateur $1$
($u \in \operatorname{Aut}_{loc}^+(\hat{\mathcal{G}}_{0,3}, \sigma_{\infty})'^{\circ}$) équivaut donc bel et
bien à celle des solutions simultanées
des équations (73), ou des équations paramétriques
("de cobord") (74).

2) Revenant au cas général (multiplicateur quelcon-
que), le seul [unclear] que je sache écrire explicite-
ment (c.à.d $u=t$) qui correspond à $\tau_\infty \in \hat{\mathcal{G}}_{0,3}$
qui a comme caractéristiques

[MARGIN: (75 ter)]
$$u = \tau_\infty \text{ , } \chi(u) = p = -1 \text{ , } \gamma = -1$$
$$g = l_1^{-1}$$
$$g_0 = 1 \text{ , } h_0 = 1 \text{ , } \xi = l_0^{-1}$$
$$\lambda' = 1 \text{ , } \lambda_0 = l_0^{-1}$$
$$[\lambda_0 = \tau_\infty \text{ , } \lambda = l_0^{-1} \tau_\infty] \leftarrow \text{ cf. ci-dessous}$$

Ce serait décidément pas un luxe d'essayer
d'écrire un paquet d'autres $u$, pour se
rendre compte notamment si les foules de
conditions nécessaires qu'on est amené à dégager
sur les éléments provenant de $\hat{\Pi}_{\mathbb{Q}}$, ne sont
pas automatiquement satisfaites !

\newpage

%% ===== Page 106 =====

[MARGIN: Calculs "par comité de censure" [unclear] qui ne sauraient peut-être plaire !]

(VII) Calculs dans $\operatorname{Gl}(2, \hat{\mathbb{Z}})' = \operatorname{Gl}(2, \hat{\mathbb{Z}})/\pm 1$
(cocycles sous $\rho$).
Ce groupe est un (tout petit!) quotient de
$\operatorname{Sl}(2, \mathbb{Z})^\wedge$, qui a l'avantage de se prêter
à des calculs algébriques explicites commodes, en termes
de matrices. Désignons, par abus de langage,
par les mêmes lettres $g, g_0, h_0, \xi, \lambda_0, \lambda$ des éléments de
(ou de $\hat{B}_{0,3}$ plus généralement)
$\hat{\pi}_{0,3}$, et leurs images dans $\pi(2)' \subset \operatorname{Sl}(2, \hat{\mathbb{Z}})'$
(... dans $\operatorname{Gl}(2, \hat{\mathbb{Z}})' \simeq \operatorname{Gl}(2, \hat{\mathbb{Z}})/\pm 1$).

Posons de façon générale, considérons une
matrice
(76) $$g = \begin{pmatrix} a & b \\ c & d \end{pmatrix} \in \operatorname{Gl}(2, \hat{\mathbb{Z}})$$

On a alors
(77) $$\rho(g) = \rho g \rho^{-1} = \begin{pmatrix} c+d & -c \\ c+d-a-b & a-c \end{pmatrix} , \det(g) \rho(g^{-1}) = \begin{pmatrix} a-c & c \\ a+b-c-d & c+d \end{pmatrix}$$
$$\rho^2(g) = \rho^{-1} g \rho = \begin{pmatrix} d-b & a+b-c-d \\ -b & a+b \end{pmatrix} , \det(g) \rho^2(g^{-1}) = \begin{pmatrix} a+b & c+d-a-b \\ b & d-b \end{pmatrix}$$

(78) $$[g, \rho] = g \rho g^{-1} \rho^{-1} = \begin{pmatrix} A & B \\ C & D \end{pmatrix}$$
[MARGIN: NB $C+D-B = [unclear]$ $\delta = \det g = ad-bc$]
$$\delta A = (a^2+b^2+ab) - (ac+bd+bc) \quad \delta B = ac+bd+bc$$
$$\delta C = -(c^2+d^2+cd) + (ac+bd+ad) \quad \delta D = c^2+d^2+cd$$

Nous écrirons l'équation
[unclear: (79) $g \rho(g) \rho^2(g) = 1$]
sous la forme équivalente
(79) $$g \rho(g) = \rho^{-1}(g^{-1}) \quad (= \rho^2(g^{-1})) \quad \text{dans } \operatorname{Gl}(2, \hat{\mathbb{Z}})'$$
où
(80) $$g \rho(g) = \begin{pmatrix} P & Q \\ R & S \end{pmatrix}$$

\newpage

%% ===== Page 107 =====

(81)
$$P = (a+b)(c+d-b)$$
$$Q = -a(c-b) - bc = \delta - a(c+d-b)$$
$$R = (c+d)^2 - d(a+b)$$
$$= (c+d)(c+d-b) - \delta$$
$$S = ad - c(c+d) = \delta - c(c+d-b)$$
où on a posé $\delta = ad-bc = \det g$

Comparant (81) et l'expression (77) de $\rho^2(g^{-1})$, l'équation (79) - qui s'écrit aussi
$$g \rho(g) = \pm \rho^2(g^{-1})$$

[MARGIN: Cherche à [unclear] des paramètres de l'équation]
[MARGIN: NB Plus utilitaire de chercher les solutions dans $\mathbb{Z}_l$ (pour $\det g = 1$ i.e. $\delta=1$) evid...]

(82) $g \rho(g) \rho^2(g) = \pm 1$ i.e. $g \rho(g) = \pm \rho^{-1}(g^{-1})$ dans $\operatorname{Gl}(2,\hat{\mathbb{Z}})'$
On distinguera les solutions suivant les facteurs du produit
$$\operatorname{Gl}(2,\hat{\mathbb{Z}})' \simeq \prod_{l \text{ premier}} \operatorname{Gl}(2,\mathbb{Z}_l)'$$

La relation
$$\frac{1}{\delta}(a+b) = P = (a+b)(c+d-b) \quad \text{dans } \mathbb{Z}_l$$
donne les deux solutions $a+b=0$ ou $c+d-b = \frac{1}{\delta}$
(car $\mathbb{Z}_l$ est intègre).

1º) $a+b \neq 0$ (dans $\mathbb{Z}_l$) / d'où $c+d-b = \frac{1}{\delta}$
(83) $\boxed{c+d-b = \frac{1}{\delta}}$

Moyennant (83), on vérifie aussitôt que la formule $g \rho(g) = \rho^{-1}(g^{-1})$ - où il reste à vérifier
$$P = \frac{a+b}{\delta} \quad Q = \frac{c+d-b}{\delta} (= \frac{1}{\delta^2} - \frac{a}{\delta})$$
$$R = \frac{c+d}{\delta} \quad S = \frac{d-b}{\delta} (= \frac{1}{\delta^2} - \frac{c}{\delta})$$

[MARGIN: La nécessité de telles formules était évidente a priori.]

équivaut à

2º) $a+b = 0$ i.e. $a = -b$, i.e. $g = \begin{pmatrix} -b & b \\ c & d \end{pmatrix}$, et
(84)
$$\det g = -b(c+d) = \delta \quad (\text{donc } \delta+b(c+d)=0) \Rightarrow d = -c - \frac{\delta}{b}$$
$$P = 0 \quad Q = \delta+b(c+d-b) = -b^2$$
$$R = (c+d)^2 = \frac{\delta^2}{b^2} \quad S = \delta - c(c+d-b) = \delta + \frac{c}{b} \delta + cb$$

\newpage

%% ===== Page 108 =====

343

1°) $a+b \neq 0$ (dans $\mathbb{Z}_l$) d'où

(83) $$c+d-b = \frac{1}{\delta}$$

[MARGIN: NB les conditions $c+d-b = \frac{1}{\delta}$, $\delta^3 = 1$ sont suffisantes pour la validité de (81) (indépendamment de la condition $a+b \neq 0$ ou $a+b=0$) ; elles sont nécessaires dès lors que $a+b \neq 0$. Voir [unclear: un peu de façon ci-dessous] le cas $g^3=1$ (?)]

Moyennant (83), on vérifie que la formule (82) -
où il reste à vérifier 
$$Q = \frac{c+d-(a+b)}{\delta} (= \frac{1}{\delta^2} - \frac{a}{\delta}) \quad , \quad R = \frac{b}{\delta^2}$$
$$S = \frac{d-b}{\delta} \cdot (= \frac{1}{\delta^2} - \frac{c}{\delta})$$
équivaut à :

(84) $$\delta^3 = 1$$
(dont la nécessité était évidente a priori).

2°) $a+b = 0$ (dans $\mathbb{Z}_l$) d'où

(85) $$a = -b \text{ i.e. } g = \begin{pmatrix} -b & b \\ c & d \end{pmatrix}$$

(86) $$\begin{cases} \text{det } g = -b(c+d) = \delta \\ d = -\frac{\delta}{b} - c \end{cases}$$

i.e.
$$g = \begin{pmatrix} -b & b \\ c & -\frac{\delta}{b} - c \end{pmatrix} ,$$

et
$$\begin{cases} P = 0 \quad , \quad Q = \delta + b(c+d-b) = -b^2 \\ R = \frac{\delta^2}{b^2} \quad , \quad S = \delta + c(\frac{\delta}{b} + b) , \end{cases}$$

L'égalité (82) s'écrit alors, 
$$\rho^2(g^{-1}) = \begin{pmatrix} 0 & -1/b \\ \delta/b & -\frac{1}{b} (\frac{\delta}{b} + b + c) \end{pmatrix} ,$$

et l'équation (81) s'écrit alors
(87)
a) $$b^3 = \delta^3 = 1$$
b) $$c( (b\delta)^2 + (b\delta) + 1 ) + (\delta^2 + b\delta + b) = 0$$

\newpage

%% ===== Page 109 =====

On distingue encore deux cas :

a) $b\delta = 1$ i.e. $b = \frac{1}{\delta}$, alors l'équation (87 b)

$3c + 3\frac{1}{\delta} = 0$ i.e.

[unclear: crossed out (88) $c = -\frac{1}{\delta} = -\delta^2 \dots$]

Or $\delta \in \hat{\mathbb{Z}}^* \implies \delta^2 \equiv 1\ (3)$ donc $2\delta^2+1 \equiv 0\ (3)$,

donc l'équation donne solution dans $\hat{\mathbb{Z}}$

(89) $$c = -\frac{1}{3}(2\delta^2+1)$$

correspond à :

(90) $$g = \begin{pmatrix} -\delta^2 & \delta^2 \\ c & 0 \end{pmatrix}$$

où $\delta' = b = \delta^2$, et

où, je le rappelle, $\delta$ est assujetti à la seule condition $\delta^3=1$.

b) $b\delta \neq 1$ donc (comme $(b\delta)^3=1$) $(b\delta)^2+b\delta+1=0$,

et l'équation (87 b) devient

$(*)$ $$\delta^2 + b^2\delta + b = 0, \quad \text{soit}$$

On aura

$$b = u\frac{1}{\delta} = u\delta^2 \quad \text{avec} \quad u^3=1, u \neq 1 \text{ i.e.}$$

$$u^2+u+1=0$$

et l'équation $(*)$ se réduit à

[unclear: crossed out lines]

$$\delta+u+1=0$$

les $u$ ne sont autres que les solutions distinctes de l'équation cyclotomique $X^2+X+1=0$. Or $u \neq 1$ signifie que $b \neq \frac{1}{\delta}$ i.e. $b \neq \delta^2$ ce qui exige i.e. $b = \delta$.

(91) $$g = \begin{pmatrix} -\delta & \delta \\ c & -1-c \end{pmatrix}$$

\newpage

%% ===== Page 110 =====

Proposition 344 Soit $g \in \operatorname{Gl}(2, \hat{\mathbb{Z}})$, $g = \begin{pmatrix} a & b \\ c & d \end{pmatrix}$, $a,b,c,d \in \hat{\mathbb{Z}}$,
$ad - bc = \delta \in \hat{\mathbb{Z}}^*$. Pour qu'on ait
$$g p(g) p^2(g) = 1 \quad (81)$$
il f. et il s. qu'on ait $\delta^3 = 1$, et qu'on
soit dans l'un des deux cas suivants, qui
s'excluent mutuellement :

1º) $c + d - b = 1$

2º) $g = -\delta' \begin{pmatrix} -1 & 1 \\ 1 & 0 \end{pmatrix} = -\delta' \rho^{-1} \quad (\text{où } \delta' = \frac{1}{\delta} = \delta^2)$

De plus, dans le cas 1º), on a [unclear: crossed out]

Corollaire Les $g$ pour lesquels on a $1^\circ)$ et
$a+b \neq 0$ sont ceux de la forme

$$g = \begin{pmatrix} a & b \\ c & b-c+\frac{1}{\delta} \end{pmatrix}$$

avec $c = \frac{1}{a+b} (a(b+\delta')-\delta)$,

ceux pour lesquels on a $1^\circ)$ et $a+b=0$
sont de la forme

$$g = \begin{pmatrix} -u\delta' & u\delta' \\ c & -\delta' u' - c \end{pmatrix}$$

où $\delta' = \frac{1}{\delta}$, $u' = \frac{1}{u} = u^2$, $u$ étant une solution
de l'équation cyclotomique $u^2+u+1=0$ i.e.
racine primitive $3^e$ de l'unité.

\newpage

%% ===== Page 111 =====

345

et on a de fait alors
[unclear: crossed out equation]

Remarques
1) Evidemment $-g$ est aussi solution de (87)
(tout comme $g^{-1}$) [unclear: crossed out text]
(tout comme $\rho^{-1}$), mais ils sont justiciables
respectivement des cas $1^\circ)$ et [unclear]

2) Pour déterminer les solutions dans $\operatorname{Gl}(2, \hat{\mathbb{Z}})$
de $g \rho g \rho^2 g = 1$, il suffit de choisir pour
$\delta \in \hat{\mathbb{Z}}^*$ une telle solution dans $\operatorname{Gl}(2, \hat{\mathbb{Z}})$, en
utilisant les résultats (86) ; on notera que
les "$a$" obtenus $1^\circ)$ à $3^\circ)$ peuvent varier d'un
$\delta$ à un autre. Même remarque pour
les solutions de l'équation $g \rho g \rho^2 g = -1$.
On fera attention que les cas $1^\circ), 2^\circ), 3^\circ)$
correspondent resp. à des familles de
solutions dépendant de 2, 1 et 0 paramètres.

3) J'ignore si une $\bar{g} \in \operatorname{Gl}(2, \hat{\mathbb{Z}})$ ($\bar{g} = \begin{pmatrix} A & B \\ C & D \end{pmatrix}$) satisfaisant
$\bar{g} \rho \bar{g} \rho^2 \bar{g} = 1$ admet une représentation

[MARGIN: NB $\bar{g} = \lambda_0^{-1} \dots$ relation de cette équation avec $\bar{g} = -\lambda_0 \dots$]

paramétrique de la forme (88) 
[unclear: crossed out equation] (il suffit
de regarder l'équation dans $\operatorname{Sl}(2, \hat{\mathbb{Z}})$).
Or l'équation en $\bar{g}$ implique $\bar{g} \equiv id$ (2)
(ce qui est une condition sur la 2-composante
dans $\operatorname{Sl}(2, \hat{\mathbb{Z}}_2)$ seulement). On voit que

\newpage

%% ===== Page 112 =====

Je lis sur les formules (78) que pour $\xi$
de la forme $\lambda_0\rho(\lambda_0)^{-1}$ (où on a bien sûr
$\xi\rho(\xi)\rho^2(\xi)=1$, bien sûr) on aura $C+D-B=1$ - i.e. on est
dans le
cas "général" à deux paramètres. Je présume

que pour
$$g = \lambda_0\rho(\lambda_0)^{-1} ,$$
on sera dans le cas B), avec $c+d-b=-1$

[MARGIN: ok, cf. (98) ci-dessous -]
i.e. le cas 1°) encore à deux paramètres. Je

laisse de côté ceci -
(sans pouvoir néanmoins exclure le cas
où on aurait en même temps
$A+B=0$ , ce qui signifie en fait que -
$$ \xi = \begin{pmatrix} A & B \\ C & D \end{pmatrix} = -\rho^{-1} = \begin{pmatrix} -1 & 1 \\ -1 & 0 \end{pmatrix} $$
i.e. le cas à un paramètre
$$ \xi = \begin{pmatrix} -u & u \\ C & -u'-C \end{pmatrix} $$
avec $u$ racine primitive 3ème de $1$ ,
$u'=u^{-1}=u^2$ l'autre racine primitive.

\newpage

%% ===== Page 113 =====

345

## (VIII) - Calculs dans $S\ell(2, \hat{\mathbb{Z}})$ : cocycles sous $\sigma_\infty$.

Reprenons un $g = \left(\begin{matrix} a & b \\ c & d \end{matrix}\right) \in S\ell(2, \hat{\mathbb{Z}})$,

et calculons $\sigma_\infty g \sigma_\infty^{-1}$, on trouve, grace à 
$$ \sigma_\infty = \left(\begin{matrix} 0 & -1 \\ 1 & 0 \end{matrix}\right) \quad \quad (48) $$

[MARGIN: NB on a $\sigma_\infty^2 = \left(\begin{matrix} -1 & 0 \\ 0 & -1 \end{matrix}\right) = -I$, $\sigma_\infty^{-1} = \left(\begin{matrix} 0 & 1 \\ -1 & 0 \end{matrix}\right) = -\sigma_\infty \neq \sigma_\infty$ et [unclear] $S\ell(2, \hat{\mathbb{Z}})$ !]

[MARGIN: NB si $g \in G\ell(2, A)$, $A$ anneau quelconque (p. ex $A = \hat{\mathbb{Z}}$), on a $\sigma_\infty g \sigma_\infty^{-1} = \left(\begin{matrix} d & -c \\ -b & a \end{matrix}\right) = (\det g) {}^t(g^{-1})$]

les formules
$$ \sigma_\infty g \sigma_\infty^{-1} = \left(\begin{matrix} d & -c \\ -b & a \end{matrix}\right) = {}^t(g^{-1}) \quad \text{la contragrédiente de } g \text{ !} $$

$$ (90) \quad \sigma_\infty g^{-1} \sigma_\infty^{-1} = {}^t g $$

$$ \underset{\sigma_\infty g^{-1} \sigma_\infty^{-1} g}{[\sigma_\infty, g^{-1}]} = ({}^t g)(g) = \left(\begin{matrix} a^2+b^2 & ac+bd \\ ac+bd & c^2+d^2 \end{matrix}\right) $$

donc, en prenant l'inverse des deux membres dans la dernière formule

$$ (91) \quad \underset{({}^t g . g)^{-1}}{g^{-1} \sigma_\infty(g)} = \left(\begin{matrix} c^2+d^2 & -(ac+bd) \\ -(ac+bd) & a^2+b^2 \end{matrix}\right) $$

D'autre part, soit
[crossed out: ... $h_0 \in S\ell(2, \hat{\mathbb{Z}})$]

on a
$$ h_\infty(h_0) = h_0 ({}^t h_0^{-1}) $$

donc
$$ (*) \quad h_\infty(h_0) = 1 \iff h_0 = {}^t h_0 $$

Soit enfin
$$ (92) \quad \xi = \left(\begin{matrix} A & B \\ C & D \end{matrix}\right) $$

et considérons l'équation (68)
$$ (**) \quad \sigma_\infty(g) = l_1 \xi l_0 $$

qui équivaut au système (*), où $h_0$ donné par (6) [crossed out: ...] et $l_0, l_1$ par (44)

\newpage

%% ===== Page 114 =====

(92) 
$$h_0 = l_0^{-r} \xi = \begin{pmatrix} 1 & +2r \\ 0 & 1 \end{pmatrix} \begin{pmatrix} A & B \\ C & D \end{pmatrix}$$
$$= \begin{pmatrix} A+2rC & B+2rD \\ C & D \end{pmatrix}$$

donc "l'équation en $\xi$" $(**)$ (écrite dans le quotient $\operatorname{Sl}(2, \hat{\mathbb{Z}})'$ de $\underline{C}_{0,3}^+$ , s'écrit $B+2rD=C$ i.e.

[MARGIN: on a rappelé (93) que $2r+1=p$]

(93) $C-B = 2rD$ i.e. $\boxed{pD=1}$ [MARGIN: $\sim \infty$]

(On fera attention que l'équation $(**)$ dans $\operatorname{Sl}(2, \hat{\mathbb{Z}})'$ signifie $h_0 \sigma_{\infty}(h_0) = \pm 1$ i.e. $h_0 = \pm {}^t h_0$, mais le signe $-$ impliquerait qu'on a $D = -D$ i.e. $2D=0$ donc $D=0$, ce qui est impossible si on suppose que $\xi = \begin{pmatrix} A & B \\ C & D \end{pmatrix}$ provient de $\overline{\Pi}_{0,3}$ donc est dans $\overline{\Gamma}'(2)$ i.e. $\equiv 1 \pmod 2$ (de sorte que $D \equiv 1 \pmod 2$).
Donc on devra bien avoir (93).

Prenons maintenant la représentation paramétrique (70) de $\xi$ (tenant compte de l'équation $\{ \rho(s)\rho(s^{-1}) = 1 \}$ (64))

(94) $\xi = \lambda \rho(\lambda)^{-1}$ [unclear: crossed out text] avec $\det \lambda = \pm 1$, ainsi $\lambda \in \operatorname{Gl}(2, \hat{\mathbb{Z}})$ (en fait $\lambda \equiv 1 \ (2)$ dans le cas qui nous intéresse, mais peu importe...)

[unclear: crossed out line at the bottom]

\newpage

%% ===== Page 115 =====

alors la formule $(78^{bis})$ donne $A,B,C,D$ comme
fonctions quadratiques de $a,b,c,d$ [unclear: (en fait] on a même
[crossed out: illegible] que $\delta=1$ dans le cas $2)$)

(95) $\begin{cases} B+C = [crossed out: illegible] D-1 \\ D = \delta(c^2+d^2+cd) \end{cases}$

donc (93) s'écrit

$[crossed out: 1-D] = 2\gamma D \quad$ i.e.

(96) $\quad (2\gamma+1)D = 1 \quad$ i.e. $\quad \rho D = 1 \qquad | \quad D = \delta(c^2+d^2+cd)$

i.e. le multiplicateur $\rho$ s'obtient en fonction
de $\xi \in \hat{\Pi}_{0,3}$ en prenant l'image
de $\xi$ dans $S\ell(2, \hat{\mathbb{Z}})'$, sous forme $\begin{pmatrix} A & B \\ C & D \end{pmatrix}$
(il y a une ambiguité de signe, levée par le
choix de $\lambda$ tel que $\xi = \lambda \rho(\lambda)^{-1}$, l'image
de $\lambda$ dans $G\ell(2, \hat{\mathbb{Z}})'$ ne détermine une
matrice [unclear: (à une unité scalaire près)] de $G\ell(2, \hat{\mathbb{Z}})$ qu'à un
signe près, mais "$\lambda \rho(\lambda)^{-1}$" ne dépend plus du
choix de ce signe ...), par la formule

(97) $\quad \rho = \frac{1}{D} \qquad \qquad (D = \delta(c^2+d^2+cd))$

Notons maintenant qu'en réduction mod 3,
la forme quadratique $c^2+d^2+cd = (c-d)^2-cd$
ne prend que les valeurs $0$ et $1$ (clair si $c=0$,
ou $d=0$, [crossed out: on a $c^2=d^2=1$] [crossed out: $f(c,d)=c^2+d^2$] - dans ces
deux cas $f(c,d)$ est un carré parfait, sinon
$c,d = \pm 1$, et $f(c,d) = 1+1 \pm 1 \equiv 1 \text{ ou } 0 \pmod 3$). Or ici
on doit trouver une unité dans $\hat{\mathbb{Z}}$, [unclear: à savoir] $D$, ceci
[unclear]

\newpage

%% ===== Page 116 =====

[MARGIN: NB On a
$D \equiv \rho \equiv \delta (3)$ (98)
car d'après ce qui précède les congruences mod 3 ne dépendent ni des cas ni de repr. part. prise.]

(car $\delta = \det \lambda \in \{\pm 1\}$, $\xi = \lambda \rho(\lambda)^{-1}$)
Cas $\xi = \lambda \rho(\lambda)^{-1} (\lambda \in \operatorname{Sl}(2,\hat{\mathbb{Z}})) \iff \rho \equiv 1 (3) \implies \gamma \equiv 0 (3)$
Cas $\xi = \lambda' \tau_1 \rho(\lambda')^{-1} (\lambda' \in \operatorname{Sl}(2,\hat{\mathbb{Z}})) \iff \rho \equiv -1 (3) \iff \gamma \equiv -1 (3)$ .

B) Cas $\xi = \lambda' \tau_1 \rho(\lambda')^{-1} (= (\lambda' \tau_1 \rho)(\lambda' \rho)^{-1})$.

[MARGIN: On voit de suite sur les formules ci-dessous mod 3
que ça implique $D \equiv -1$ (car on a $\delta=1$ puis qu'on part de $\lambda'$ [unclear], $\lambda = \lambda' \tau_0 = \lambda' \rho(\tau_1)$, det $\lambda = -\lambda'$ [unclear])]

Utilisant la deuxième formule (87) on trouve
$$ \xi = \begin{pmatrix} a & b \\ c & d \end{pmatrix} \begin{pmatrix} 1 & 2 \\ 0 & 1 \end{pmatrix} \begin{pmatrix} a-c & c \\ a+b-c-d & c-d \end{pmatrix} = \begin{pmatrix} A & B \\ C & D \end{pmatrix} $$

avec

[MARGIN: NB $C+D-B = -4a$]

(99)
$A = 3a^2 + b^2 + 3ab - 3ac - 2ad - bc - bd \qquad B = 3ac + 2ad + bc + bd$
$C = -(3c^2 + 3cd + d^2) + 3ac + 2ad + bc + bd \qquad D = 3c^2 + 3cd + d^2$

Comme tantôt la formule (93) (qui exprime "l'égalité [unclear]")

[MARGIN: On a vu tantôt [unclear]
$r_0 = -r h_0$, car $r_0$ [unclear]
d'où $\xi \equiv 1$ (2) pour $D \equiv 1$ (2).]

(14) où elle ne l'a pas s'écrit , de façon compte tenu de 
$B-C = -1 + D$
sous la forme
$2\gamma D = -1 - D$ i.e.

(100)
$$ (2\gamma + 1)D = -1 \qquad i.e. \qquad \rho D = -1 $$

donc ici le multiplicateur $\rho$ s'exprime en fonction de l'image de $\xi$ dans $\operatorname{Sl}(2,\hat{\mathbb{Z}})'$, mis sous forme $\begin{pmatrix} A & B \\ C & D \end{pmatrix}$ où $A,B,C,D \in \hat{\mathbb{Z}}$ -- et en levant l'ambigüité de signe comme tantôt -- par la formule

(101) $\rho = -\frac{1}{D}$ \qquad ($D = 3c^2 + 3cd + d^2$) .

\newpage

%% ===== Page 117 =====

Ici il est clair que les seules valeurs mod 3
représentées par la forme quadratique
$3c^2+3cd+d^2 \equiv d^2 \quad (3)$ sont $0$ et $1$, d'où

(102) $$p \equiv -1 \quad (3) \quad (\text{i.e. } \gamma \equiv -1 \quad (3)) \text{ si } \xi = \lambda l_0^{-1} p(\lambda)^{-1}$$
$$(\Leftrightarrow g = \lambda_0 l_0^{-1} p(\lambda_0)^{-1})$$

Cette étude montre donc bien que l'on est
dans le cas $\xi = \lambda \rho(\lambda)^{-1}$, resp. $\xi = \lambda l_0^{-1} \rho(\lambda)^{-1}$,
suivant que $p = \chi(u) \equiv 1 \text{ ou } \equiv -1 \quad (3)$,
i.e. suivant que $\gamma \equiv 0 \text{ ou } \gamma \equiv -1 \quad (3)$.
(ou plutôt, dans [unclear])

Remarques. 1) Se plaçant dans $\operatorname{Sl}(2, \hat{\mathbb{Z}})'$ (ou bien
dans le sous-groupe $\pi(2)'= \operatorname{Ker}(\operatorname{Sl}(2, \hat{\mathbb{Z}})' \to \operatorname{Sl}(2, \mathbb{Z}/2\mathbb{Z}))$
de l'un de $\overline{\Pi}_{0,3}$), pour déterminer les
solutions de "l'équation en $\xi$" (**) (cf. (68)),
quand $\xi$ est mis sous forme paramétrique
(94), on voit qu'en tant que condition sur
$\lambda = \begin{pmatrix} a & b \\ c & d \end{pmatrix}$, cette équation est relativement
anodine et très simple — elle dit
simplement que $pD=1$, où
$D = c^2+d^2+cd$ dans le premier cas,
$D = 3c^2+3cd+d^2$ dans le deuxième. Si on
se fixe un $p \in \hat{\mathbb{Z}}^*$ d'avance, cette
condition dit simplement que l'élément

\newpage

%% ===== Page 118 =====

$D$ de $\hat{\mathbb{Z}}$ soit une unité, i.e.

(103) $D \in \hat{\mathbb{Z}}^*$

[MARGIN: NB On [unclear] $p$]

- et on détermine alors de façon unique un terme de $D$, donc de $c$ et de $d$ et du signe $\delta$ par $D = \delta(c^2+d^2+cd)$ (à [unclear] inversibles) - donc aussi en termes de $\lambda \in \hat{\Pi}_{0,3}^{\sim, (1,2)}$, lui-même uniquement déterminé par $\xi$.

Ainsi la connaissance de $\xi \in \hat{\Pi}_{0,3}$ suffit à déterminer le multiplicateur $p = 2\gamma+1 \in \hat{\mathbb{Z}}^*$, donc aussi $\gamma \in \hat{\mathbb{Z}}$, et $h_0 = l_0^{-\gamma}\xi$, donc aussi $g_0$ (par la formule (66) $g_0 \sigma_\alpha(g_0) = h_0$), et par suite $g$ par la formule (62) $g = \sigma_\alpha(g_0)l_1^\gamma \rho(g_0)^{-1}$, et $u$ par (50). Il est [unclear: oiseux] de s'assurer que cela ne change pas quand on fait le changement (69)
$$\xi \longmapsto \xi' = l_0^{-a} \xi l_1^a$$

en termes matriciels, [unclear]

(104) $l_0^{-\alpha} \xi l_1^\alpha = \begin{pmatrix} 1 & 2\alpha \\ 0 & 1 \end{pmatrix} \begin{pmatrix} A & B \\ C & D \end{pmatrix} \begin{pmatrix} 1 & 0 \\ 2\alpha & 1 \end{pmatrix} = \begin{pmatrix} A+2\alpha(B+C)+4\alpha^2D & B+2\alpha D \\ C+2\alpha D & D \end{pmatrix}$

2) Mais on fera attention qu'en toute rigueur la seule connaissance de l'image de $\xi$ dans $\operatorname{Sl}(2, \hat{\mathbb{Z}})'$ ne permet de déterminer le multiplicateur $p$ qu'au signe près - par contre la connaissance de l'image de $\lambda$ dans $\operatorname{Gl}(2, \hat{\mathbb{Z}})'$ c-à-d d'une matrice $\begin{pmatrix} a & b \\ c & d \end{pmatrix} \in \operatorname{Gl}(2, \hat{\mathbb{Z}})$ définie au signe près - permet de déterminer

\newpage

%% ===== Page 119 =====

349

$p$ comme l'inverse de $D$, donné suivant le cas
- où $\delta \in \{\pm 1\}$ est déterminé par la
[crossed out: par $D = \delta(c^2+cd+d^2)$] par $D = c^2+cd+d^2$ ou $D = 3(c^2+cd)+d^2$
condition $D \equiv \delta \ (3)$.

3) On sait par voie arithmétique (car
$$ \Gamma_{\mathbb{Q}} \to \hat{\mathbb{Z}}^* \quad \text{est surjectif)} $$
qu'on doit trouver n'importe quel $p \in \hat{\mathbb{Z}}^*$ sous
la forme $1/D$ (ou encore, sous la forme $D$),
avec $D = c^2+cd+d^2$ si $p \equiv 1 \ (3)$, et
$D = -3(c^2+cd) + d^2$ si $p \equiv -1 \ (3)$, avec $c,d \in \hat{\mathbb{Z}}$,
[unclear: sous la réserve d'imparité dirais-je]
et on a $c \equiv 0 \ (2)$, $d \equiv 1 \ (2)$. (On est ramené à le
vérifier séparément pour les différents $\mathbb{Z}_{\ell}$, etc.
et pour $\mathbb{R}$).

[MARGIN: ou m'en fiche de $-3(c^2+cd)+d^2)$!]
Corollaire. Si $\ell$ est premier $\neq 3$, alors tout
$p \in \mathbb{Z}_{\ell}^*$ se met indifféremment sous la forme
$c^2+cd+d^2$, ou $-3(c^2+cd)+d^2$ (avec $c,d \in \mathbb{Z}_{\ell}$). Si
$\ell = 2$, on peut supposer $c \equiv 0$, $d \equiv 1 \ (2)$. Pour
$p \in \mathbb{Z}_3^*$, $\ell=3$, il se met sous la $1^{ere}$ forme
$c^2+cd+d^2$ (resp. $-3(c^2+cd) + d^2$) si $p \equiv 1 \ (3)$
(resp. $p \equiv -1 \ (3)$), ce qui est mutuellement
exclusif.

4) Revenant à la réduction "anodine" dans
$\operatorname{Sl}(2, \hat{\mathbb{Z}})'$ de l'équation en $\lambda"$ comme
signifiant simplement que $c^2+cd+d^2 \in \hat{\mathbb{Z}}^*$
(resp. $-3(c^2+cd)+d^2 \in \hat{\mathbb{Z}}^*$), on pourrait espérer

\newpage

%% ===== Page 120 =====

tenir une méthode pour construire dans
$\hat{\Pi}_{0,3}$ des gros paquets de solutions de l'équa-
tion en question. Or je n'y suis absolu-
ment pas arrivé - pas à en construire
explicitement une seule, en dehors de celles
qui correspondent à $u=1$ et $u=t_0$.
Une question : Soit $\lambda_0 \in S\ell(2, \hat{\mathbb{Z}})'$, $\lambda_0 \equiv 1 (l)$, tel que
$\xi_0 = \lambda_0 \rho(\lambda_0)^{-1}$ (disons, $-\xi_0 = \lambda_0 \rho(-\lambda_0)^{-1}$) soit tel que
satisfasse dans $S\ell(2, \hat{\mathbb{Z}})'$ à une "équation en $\xi$"
pour $\gamma = p^{-1}$ convenable ($p \in \hat{\mathbb{Z}}^*$) - i.e.
$c^2+d^2+cd \in \hat{\mathbb{Z}}^*$ (resp. $3(c^2+cd)+d^2 \in \hat{\mathbb{Z}}^*$).
Existe-t-il alors un $\tilde{\lambda} \in \hat{\Pi}_{0,3}$ au-dessus de
$\lambda_0$, tel que $\tilde{\xi} = \tilde{\lambda} \rho(\tilde{\lambda})^{-1}$ (resp. ...) satisfasse
aussi une équation en $\tilde{\xi}$ (certainement
avec le même $\gamma$ )?

\newpage

%% ===== Page 121 =====

43 Retour sur les [unclear: relations] 350 - Digression sur les 1-cocycles des groupes cycliques.

Soit $\Pi$ un groupe, $\rho$ un autom. de $\Pi$ tel que
(1) $$ \rho^n = id $$
où $n \in \mathbb{N}^*$. On a donc une opération du groupe cyclique $\mathbb{Z}/n\mathbb{Z}$ sur $\Pi$, d'où un produit semi-direct
(2) $$ \tilde{\Pi} = \Pi \cdot \mathbb{Z}/n\mathbb{Z} $$
et une structure d'extension
(3) $$ 1 \to \Pi \to \tilde{\Pi} \to \mathbb{Z}/n\mathbb{Z} \to 1 . $$
On identifie $\rho$ au générateur de $\mathbb{Z}/n\mathbb{Z}$, et on l'identifie à son image dans $\tilde{\Pi}$, par le scindage canonique
(4) $$ \mathbb{Z}/n\mathbb{Z} \to \tilde{\Pi} \quad , \quad 1 \bmod n \mapsto \rho $$
On a donc
(5) $$ \rho(g) = \rho g \rho^{-1} \qquad (g \in \Pi) . $$
Les scindages de l'extension (3), correspondent biunivoquement aux éléments $\rho'$ de $\tilde{\Pi}$ sur $\rho$, tels que
(6) $$ {\rho'}^n = 1 $$
Dire que $\rho'$ est sur $\rho$ signifie que
(7) $$ \rho' = g \rho \qquad \text{pour } g \in \Pi , $$
et en termes d'un tel $g \in \Pi$ l'équation (6) se réécrit
(8) $$ g \rho(g) \rho^2(g) \dots \rho^{n-1}(g) = 1 $$

\newpage

%% ===== Page 122 =====

[Posons, pour $g \in \pi$, $\alpha \in \pi$,
$\alpha \rho'_g \alpha^{-1} = \dots$ [crossed out]]

Le groupe $\pi$ opère par conjugaison sur l'ensemble des scindages de l'extension (3), i.e. ce qui revient au même, sur l'ensemble des $\rho'$ ci-dessus de $\rho$, satisfaisant (6). On a d'ailleurs, pour $\alpha \in \pi$

(9) $$ \alpha \rho'_g \alpha^{-1} = \alpha g \rho \alpha^{-1} = \alpha g \rho \alpha^{-1} \rho^{-1} \rho = (\alpha g \rho(\alpha)^{-1}) \rho = \rho'_{\alpha g \rho(\alpha)^{-1}} $$

i.e. l'opération par conjugaison de $\pi$ sur l'ensemble des scindages de (3), s'interprète comme l'opération $\Theta_\rho$ de $\pi$ sur l'ensemble des solutions de (8) donnée par :

(10) $$ g \mapsto \alpha g \rho(\alpha)^{-1} $$

(on dira que $\alpha g \rho(\alpha)^{-1}$ est le $\rho$-conjugué de $g$ par l'élément $\alpha \in \pi$). Les classes de $\pi$-conjugaison de scindages de (3) correspondent aux classes de $\pi$-$\rho$-conjugaison de solutions de (8).

Parmi ces classes, la classe "triviale" (du scindage canonique de (3)), i.e. la classe de $g=1$, est donnée paramétriquement comme l'ensemble des

(11) $$ g = \alpha \rho(\alpha)^{-1} \quad (\alpha \in \pi) $$

Soient $g, g'$ deux solutions conjuguées de

\newpage

%% ===== Page 123 =====

l'équation (8), de sorte qu'on a

(12) $$g' = \alpha g \rho(\alpha)^{-1}$$ i.e. $$g' = \Theta(\alpha).g$$

pour un $\alpha \in \pi$ convenable. Quand a-t-on
unicité de $\alpha$ — plus généralement, que peut-
-on dire, pour $g, g'$ fixés, de l'ensemble des $\alpha \in \pi$
qui tels qu'on ait (12) (l'équa- par laquel-
la $\rho$-transformée de $g$ donne $g'$, dans le groupe $\pi$) ?
C'est évidemment une classe à droite modulo le
sous-groupe $\pi(\rho, g)$ de $\pi$ qui stabilise $g$,

(13) $$\pi(\rho, g) = \{ \beta \in \pi \mid \Theta_{\rho}(\beta).g = g \}$$
$$= \{ \beta \in \pi \mid \beta g \rho(\beta)^{-1} = g \}$$
$$= \{ \beta \in \pi \mid int(\beta)(g\rho) = g\rho \}$$
$$= \pi^{g\rho} = \{ \beta \in \pi \mid int(g\rho)\beta = \beta \}$$
$$= \text{Centralisateur de } g\rho \text{ dans } \pi$$

En d'autres termes, si $\alpha$ satisfait (12), l'ensemble des
$\alpha'$ qui satisfont la même relation n'est autre que
l'ensemble $\alpha \pi(\rho, g)$ :

(14) $$\alpha' g \rho(\alpha')^{-1} = \alpha g \rho(\alpha)^{-1} \iff \alpha^{-1}\alpha' \in \pi(\rho, g) = \pi^{g\rho}$$
i.e. $$\alpha^{-1}\alpha' \in \pi^{g\rho}$$

Donc unicité de $\alpha$ satisfaisant (12) ssi

(15) $$\pi^{g\rho} = 1$$

condition qui ne dépend que de la classe de
$\pi$-conjugaison du scindage envisagé de (3) défini
par $g$, i.e. par $g\rho$.

\newpage

%% ===== Page 124 =====

En fait, non seulement le groupe $\pi$, mais le groupe $\tilde{\pi}$ lui-même opère par conjugaison sur l'ensemble des scindages de l'extension (3). Il en est ainsi en particulier de $\rho$, et on a :

$$ (16) \quad ^\rho(g\rho) = \rho (g\rho) \rho^{-1} = (\rho g \rho^{-1}) \cdot \rho = (^\rho g) \rho , $$

i.e. l'opération de $\rho$ sur l'un des scindages en question, s'identifie à l'opération induite par $\rho$ la permutation de $\pi$ sur l'ensemble des solutions de (8) ($g \in \pi$). Notons d'ailleurs que $g$ et $\rho(g)$ sont $\pi$-conjugués, on a

$$ (17) \quad \rho(g) = \alpha g \rho(\alpha)^{-1} \quad \text{où} \quad \alpha = g^{-1} \ . $$

Notons que l'opération de $\tilde{\pi}$ sur l'un des scindages de (3), ou encore sur l'une des solutions de (8), est donc donnée par

$$ \Theta_\alpha : g \mapsto \alpha g \rho(\alpha)^{-1} \quad (1u) $$

où maintenant $\alpha \in \tilde{\pi}$ (et pas nécessairement $\alpha \in \pi$). Les classes de $\tilde{\pi}$-conjugaison (comme il résulte immédiatement de la validité de (17)) coïncident. En particulier, si $\alpha \in \tilde{\pi}$, l'élém.

$$ (18) \quad \alpha \rho(\alpha)^{-1} $$

est un él. de $\pi$, et il satisfait (8) et

\newpage

%% ===== Page 125 =====

se trouve dans la classe principale, i.e. on aura

$$ \alpha \rho(\alpha)^{-1} = \alpha_0 \rho(\alpha_0)^{-1} \quad \text{pour } \alpha_0 \in \pi \text{ conv.} $$
(où $\alpha_0 \in \pi, i \in \mathbb{Z}$)

En fait, on peut écrire $\alpha = \alpha_0 \rho^i$ (d'où $\rho(\alpha)^{-1} = \rho^{-i} \rho(\alpha_0)^{-1}$)
d'où (19). Plus généralement, on aura, pour $\alpha = \alpha_0 \rho^i$ ($\alpha \in \tilde{\pi}, \alpha_0 \in \pi, i \in \mathbb{Z}$)
la relation
$$ \alpha g \rho(\alpha)^{-1} = \alpha_0 \rho^i g \rho^{-i} \rho(\alpha_0)^{-1} = \alpha_0 \rho^i(g) \rho(\alpha_0)^{-1} $$
avec
$$ \rho^i(g) = \beta g \rho(\beta)^{-1} \quad \beta = (g \rho(g) \dots \rho^{i-1}(g))^{-1} $$
d'où

(19)
$$ \alpha g \rho(\alpha)^{-1} = (\alpha_0 \beta) g \rho(\alpha_0 \beta)^{-1} $$
$$ \beta = (g \rho(g) \dots \rho^{i-1}(g))^{-1} $$
(pour $\alpha = \alpha_0 \rho^i$, $i \in \mathbb{N}$).

Bien sûr, la représentation (18) d'une solution de l'équation (8), où l'on suppose non plus $\alpha \in \pi$ mais seulement $\alpha \in \tilde{\pi}$, n'a plus rien d'univoque, puisque pour $\alpha = \rho$ p.ex. on trouve la solution $g = \rho \rho^{-1} = 1$.

Supposons plus généralement qu'on dispose d'un groupe
$$ \mathcal{E} \supset \tilde{\pi} $$
(comme pour le groupe $\pi = \pi_{0,3}$, $\tilde{\pi}$ image inverse de $\{1, \rho, \rho^2\} \subset \mathfrak{S}_3$ dans $\mathfrak{S}_{0,3}$ extension de $\mathfrak{S}_3$ par $\pi_{0,3}$), dans lequel $\pi$ soit invariant, de sorte que $\mathcal{E}$ opère par conjugaison

\newpage

%% ===== Page 126 =====

sur l'un des $\rho'$ dans $\pi$ tel que

(20) $$1 \longrightarrow \pi \longrightarrow \mathcal{E} \longrightarrow G \longrightarrow 1$$

[MARGIN: NB L'extension (20) est l'image inv. par l'inclusion $\mathbb{Z}/n\mathbb{Z} \hookrightarrow G$ de [unclear]]

(21) $$\mathbb{Z}/n\mathbb{Z} \hookrightarrow G, \quad 1 \bmod n \mapsto \dot{\rho} \in G$$

$\pi$ s'identifie à l'image inv. de $\mathbb{Z}/n\mathbb{Z}$ dans $\mathcal{E}$. (élém d'ordre $n$)
On ne peut dire en général, si $\rho'$ est un
élément de $\mathcal{E}$ au dessus de $\dot{\rho}$ (donc $\rho' \in \mathcal{E}$)
($= s'il satisfait \rho'^n=1$), que pour tout
$\alpha \in \mathcal{E}$, $int(\alpha)(\rho')$ est encore un au dessus de $\dot{\rho}$ :
pour qu'il en soit ainsi, il faut e.s. que
$\dot{\alpha} \in G$ centralise $\dot{\rho}$. Soit donc $\mathcal{E}_{\dot{\rho}}$

(22) $$\mathcal{E}_{\dot{\rho}} = \{ \alpha \in \mathcal{E} \mid [\dot{\alpha}, \dot{\rho}] = 1 \}$$

c'est un sous-groupe de $\mathcal{E}$ contenant $\pi$, qui
opère naturellement sur l'ensemble des $\rho'$ au dessus
de $\dot{\rho}$, satisfaisant (8), et ce par les opérations

(23) $$\alpha \rho'_x \alpha^{-1} = \rho'_{\alpha x \alpha^{-1}}$$ [unclear: best guess]

de sorte qu'on peut aussi considérer cette op.
donnée par

(24) $$\Theta_{\dot{\rho}, \mathcal{E}_{\dot{\rho}}} (\alpha) = \Theta(\alpha) : \quad g \mapsto \alpha g \rho(\alpha)^{-1} \quad (\alpha \in \mathcal{E}_{\dot{\rho}})$$

Cette [unclear], une opération $\Theta(\alpha)$ pourra permettre
les classes de $\pi$-conjugaison de telles solutions,
de (8) dans $\pi$. Sur $\mathcal{E}_{\dot{\rho}}$ on aura [unclear]

\newpage

%% ===== Page 127 =====

353

une structure d'extension

(25) $$1 \longrightarrow \pi \longrightarrow \mathcal{E}_\rho \longrightarrow G^\rho \longrightarrow 1$$
$$Cent_G(\rho) = \{g \in G \mid \rho g \rho^{-1} = g\}$$
$$= \{g \in G \mid [\rho, g] = 1\}$$

et $\mathcal{E}_\rho$ opère sur l'ens des
classes de $\pi$-conjugaison de scindages de (21),
(ou des classes de $\pi$-conjug. de solutions de (8))
via sa projection dans $G^\rho$ dessus.

Soit $\alpha \in \mathcal{E}$, considérons l'élément
$$\alpha \rho \alpha^{-1} \rho^{-1} = \alpha (\rho \alpha \rho^{-1})^{-1} = \alpha \rho \alpha^{-1} \rho^{-1} = [\alpha, \rho]$$
Il est dans $\pi$ ssi [unclear]
$$[\dot{\alpha}, \rho] = 1 \quad \text{i.e.} \quad \dot{\alpha} \in G^\rho \quad \text{i.e.} \quad \alpha \in \mathcal{E}_\rho.$$

Pour une solution $g$ de l'équation (8),
il pourra arriver qu'on ne puisse la mettre
sous la forme standard $\alpha \rho \alpha^{-1}$ avec $\alpha \in \pi$, mais qu'on
y arrive avec $\alpha \in \mathcal{E}_\rho$. Pour qu'il existe un
tel $\alpha$, il f. et s. que la classe de $\pi$-conjug.
de $g$ (parmi les solutions de (8)) soit dans
l'orbite de la classe de $\pi$-conjugaison triviale $[\rho]$
par l'action considérée de $G^\rho$ sur cet ens. des
classes de conjugaison ; alors les $\alpha \in \mathcal{E}_\rho$ qui
satisfont

(26) $$g = \alpha \rho \alpha^{-1}$$

\newpage

%% ===== Page 128 =====

satisfaisant : la relation
[crossed out text]
$$ \dot{\alpha} [1]_\rho = [g]_\rho $$
(27)

et pour tout $\dot{\alpha} \in G^\rho$ satisfaisant : cette relation
(27), il existe au moins un représentant $\alpha \in \mathcal{E}_\rho$
satisfaisant (6) (ce qui est pure tautologie).
A quelle condition a-t-on unicité des
$\alpha \in \mathcal{E}_\rho$ satisfaisant : (26) ? L'ens. de ces
$\alpha$ est une classe à droite sous le ss-groupe
de $\mathcal{E}_\rho$ formé des $\alpha$ tel que $\alpha \rho(\alpha)^{-1} = 1$
i.e. $\alpha = \rho(\alpha)$ i.e. $[\alpha, \rho] = 1$
ce groupe n'est autre que $\mathcal{E}^\rho$

(28) $$ \mathcal{E}^\rho = Cent_{\mathcal{E}}(\rho) \subset \mathcal{E}_\rho $$

Ce groupe contient $\rho$ (donc le ss-groupe de $\mathcal{E}$
$\simeq \mathbb{Z}/n\mathbb{Z}$ engendré par $\rho$) - il n'est donc
jamais trivial - i.e. on n'a jamais unicité
(sauf dans le cas trivial où $n=1$ i.e. $\rho=1$,
mais alors $\mathcal{E}^\rho = \mathcal{E}$ et il faut $\mathcal{E}=1$ ... !)

Mais si on suppose [crossed out: $\pi^\rho = \{1\}$] i.e.

(29) $$ \mathcal{E}^\rho \cap \pi = \{1\} $$

alors
$$ \mathcal{E}^\rho \hookrightarrow G^\rho \subset G $$
[MARGIN: injectif]
alors les $\alpha \in \mathcal{E}_\rho$ solutions de (26) [unclear: sont univoquement]

\newpage

%% ===== Page 129 =====

354

donc bi-univoque avec [unclear: les classes de] conjugaison dans $G$, i.e.
avec la transformée de $\Pi_1$ sous $[g_p]$ : si on
choisit un élément de ce transformé, on en déduit
par là-même un $\alpha$, solution de (26).

Dans le cas de $\Pi = \Pi_{0,3}$ ( $= \Pi_1(\mathcal{M}_{0,3})$ ), $\mathcal{E} = \mathcal{E}_{0,3}$ ( $= \Pi_1(U_{0,3}, C_{0,3})$ )
$G = C_{0,3} \simeq \hat{\mathfrak{S}}_3 \times \widehat{\mathbb{Z}/2\mathbb{Z}}$, avec $\rho, \rho'$ comme d'accou-
tumée, l'ens des classes de conjugaison des
[crossed out: sous-groupes de] (3) i.e. des solutions de (8)
s'identifie à $U_{0,3}$, par les deux éléments,
les deux "pôles" $-J$, $-j$ -- dont l'un, choisi
comme point base, correspond à la classe triviale.

[crossed out: Ici le groupe $G \simeq \hat{\mathbb{Z}}$ est le groupe]

$G_P = G_{\bar{P}} \simeq \widehat{\mathbb{Z}/3\mathbb{Z}} \times \widehat{\mathbb{Z}/2\mathbb{Z}}$
est le groupe engendré par $\rho$ et $\tau$, et $\rho$ opère
trivialement sur l'un des deux demi-plans, tandis
que $\tau$ les permute (i.e. $\widehat{\mathbb{Z}/2\mathbb{Z}}$ opère par
façon simplement transitive). Donc si on pose

(30) $\mathcal{E}' = (\Pi_{0,3} . \{1, \tau\}) = $ image inverse de $\{1, \tau\}$ dans $\mathcal{E} = \mathcal{E}_{0,3}$
[MARGIN: (groupe cartographique [unclear: pondéré non orienté...] )]

on voit que

(31) $\mathcal{E}'_\rho = \{1\}$
Conséquence immédiate de la relation $\Pi^\rho = (1)$
et l'ens des solutions dans $\Pi_{0,3}$ de l'équation

(32) $g \rho(g) \rho^2(g) = 1$

\newpage

%% ===== Page 130 =====

est en corr. $1-1$ avec l'ens des éléments
de $\tilde{\Pi}_{0,3}^\sim = \mathcal{E}'$ (le groupe cartographique triangulé pondéré non
orienté, ayant comme générateurs les éléments
$\tau_0, \tau_1, \tau_\infty$ reliés par les relations $\tau_0^2 = \tau_1^2 = \tau_\infty^2 = 1$ sans
plus, avec $l_0 = \tau_\infty \tau_1$, $l_1 = \tau_0 \tau_\infty$, $l_\infty = \tau_1 \tau_0$ ...). Il doit
encore en être de même quand on se place
dans les groupes $\hat{\Pi}_{0,3}$ et $\hat{\Pi}_{0,3}^\sim = \hat{\Pi}_{0,3} \cdot \{ 1, \tau_\infty \}$,
si on admet (chose à vérifier ultérieurement)

(33) $\quad (\hat{\Pi}_{0,3}^\sim)^\rho = \{ 1 \}$

[MARGIN: et que les éléments d'ordre fini de $\mathcal{B}_{0,3}$ sont conjugués des éléments d'ordre fini de $\mathcal{E}_{0,3}$... ici]

Notons que l'on a ici

(34) $\quad \tau_\infty \rho(\tau_\infty)^{-1} = \tau_\infty \rho(\tau_\infty) = l_1^{-1}$, vu

en vertu des formules (23) du § 42

$\rho_\infty \overset{\text{déf}}{=} \tau_\infty \rho \tau_\infty^{-1} = l_1^{-1} \rho = \rho l_0^{-1}$

La situation est encore plus simple pour
l'élément $\sigma_{\vec{\infty}}$ de $\mathcal{G}_{0,3}$, opérant sur $\pi$ par
un automorphisme d'ordre $n=2$, car ici toute
solution de l'équation correspondante

(35) $\quad \beta \sigma_{\vec{\infty}} (\beta) = 1 \quad (\beta \in \pi)$

se met sous la forme évidente

(36) $\quad \beta = \beta' \sigma_{\vec{\infty}} (\beta')^{-1} \quad (\beta' \in \pi)$

et ceci de façon unique. Cela résulte
en effet de la théorie [unclear: générale] ..., compte
tenu que

(37) $\quad U_{0,3}^{\sigma_{\vec{\infty}}} = \{ 1 \}$

\newpage

%% ===== Page 131 =====

l'ens. des pts fixes de $U_{0,3}$ sous $\sigma_{\bar{\infty}}$ est
réduit à un seul élément, ce qui implique
ici par :
$$ \Pi_{0,3}^{\sigma_{\bar{\infty}}} = \{1\} . $$

Nous admettrons encore que ces faits s'étendent
aux éléments de $\widehat{\Pi}_{0,3}$ - ce qui
sur le [unclear] n'a au plus de ce qui en a été
admis à propos de $\rho$) se traduit par l'equation
(38) $$ (\widehat{\Pi}_{0,3})^{\sigma_{\bar{\infty}}} = \{1\} $$

Notons qu' ici
(39) $$ G^{\sigma_{\bar{\infty}}} = \mathcal{C}_{0,3}^{\sigma_{\bar{\infty}}} = \mathbb{Z}/2\mathbb{Z} \times \mathbb{Z}/2\mathbb{Z} \ ,$$
sous-groupe engendré par $\sigma_{\bar{\infty}}, \tau_{\bar{\infty}}$, dans le
[unclear] avec $\sigma_{\bar{\infty}}^2 = \tau_{\bar{\infty}}^2 = 1$ pour ne
mentionner que $\sigma_{\bar{\infty}}, \tau_{\bar{\infty}}$ ; et on a
(40) $$ \mathcal{E}^{\sigma_{\bar{\infty}}} = \widehat{\mathcal{C}}_{0,3}^{\sigma_{\bar{\infty}}} \simeq \hat{\mathbb{Z}}/2\mathbb{Z} \times \hat{\mathbb{Z}}/2\mathbb{Z} \ , $$
groupe engendré par $\sigma_{\bar{\infty}}, \tau_{\bar{\infty}}$ (l'un d'eux
un autom d'ordre 2). Si donc on écrit
une solution $h \in \widehat{\Pi}$ de (35) sous la forme (36),
avec $\beta \in \widehat{\mathcal{C}}_{0,3}$ au lieu de $\beta \in \pi$ (resp $\beta \in \pi$) , alors l'indétermination est dans
ce groupe (40) d'ordre 4 ; de façon plus
précise, si $\beta_0$ est une solution dans $\pi$, alors
les autres solutions sont de la forme
(41) $$ \beta = \beta_0 \sigma_{\bar{\infty}}^\varepsilon \tau_{\bar{\infty}}^{\varepsilon'} \quad \varepsilon, \varepsilon' \in \mathbb{Z}/2\mathbb{Z} $$

\newpage

%% ===== Page 132 =====

(avec $\varepsilon, \varepsilon'$ univoquement déterminés). Pour
qu'on ait $\beta \in \overline{\mathfrak{S}}_{0,3}^+$, il f. et s. que $\varepsilon' = 0$.

\newpage

%% ===== Page 133 =====

[MARGIN: 356]
[MARGIN: Ces notations seront revues encore une fois au § 46]
[MARGIN: autom. extérieurs : $\mathfrak{S}_{3}$]

## 44 Retour sur les notations. Mise en équations dans $\operatorname{Sl}(2, \hat{\mathbb{Z}})$.

Soit $u$ un autom. de $\widehat{\Pi}_{0,3}$ ( — ou ce qui revient au même, un autom. de $\widehat{\mathcal{B}}_{0,3}^{+}$ qui induit l'identité dans $\overline{\mathcal{C}}_{0,3} = \mathfrak{S}_{3}$. Alors $u$ est connu quand on connait ses valeurs sur les deux générateurs $\sigma_{\infty}, \rho$ de $\widehat{\mathcal{B}}_{0,3}^{+}$, qui seront resp. deux éléts de $\widehat{\mathcal{B}}_{0,3}^{+}$ au dessus de $\bar{\sigma}_{\infty}$ et de $\bar{\rho}$ , soit

(1) $$u(\sigma_{\infty}) = h\sigma_{\infty} \quad , \quad u(\rho) = g\rho \quad (h,g \in \widehat{\Pi}_{0,3})$$

où $h, g$ sont assujettis aux conditions

(2) $$h \sigma_{\infty}(h) = 1 \quad , \quad g \rho(g) \rho^{2}(g) = 1 ,$$

exprimant respectivement les relations

(3) $$(h \sigma_{\infty})^{2} = (g\rho)^{3} = 1$$

(Il y a un peu à vérifier la condition liminaire — savoir que (3) sont des relations génératrices de $\widehat{\mathcal{B}}_{0,3}^{+}$ ...).

Pour le moment, il n'y a pas lieu d'imposer aucune relation entre $g$ et $h$ — sauf les relations (2) sur chacun de ces éléments — puisque $\widehat{\mathcal{B}}_{0,3}^{+}$ est le groupe ayant les relations génératrices

$$\sigma_{\infty}^{2} = \rho^{3} = 1 .$$

\newpage

%% ===== Page 134 =====

Les éléments $g$ et $h$, du point de vue
des autom. de $\hat{\Pi}_{0,3}$, expriment les
relations de commutativité de $u$ avec
$\rho$ et $\sigma_{\infty}$ respectivement :

(4) $[u,\rho] = u \rho u^{-1} \rho^{-1} = int(g) \quad , \quad [u,\sigma_{\infty}] = u \sigma_{\infty} u^{-1} \sigma_{\infty}^{-1} = int(h)$

- c'est une chose remarquable que la
connaissance des éléments $g,h$ de $\Pi$
qui vérifient (4), implique déjà la
connaissance de $u$ (grâce au fait que
$\sigma_{\infty}$, $\rho$ engendrent $\widehat{C}_{0,3}^+$ ...).

Écrivons maintenant la compatibilité de
$u$ avec la structure [unclear] : On l'avait
formalisée prudemment par une relation
de la forme

$u(\sigma_0) = int(g_0) (\sigma_0^\rho)$, ou $u(\sigma_0) = int(g_0) \sigma_0^\rho$ ,

pour

$\rho \in \hat{\mathbb{Z}}^*$, $g_0 \in \hat{\Pi}_{0,3}$

convenables. La notation $g_0$ n'est pas telle-
ment heureuse, puisque $g$ et $g_0$ jouent
des rôles tout à fait dissemblables. L'indice
$0$ dans $g_0$ était destiné à pouvoir écrire,
plus généralement

$u(\sigma_i) = int(g_i) \sigma_i^\rho \quad i \in \{0,1,\infty\}$ .

\newpage

%% ===== Page 135 =====

Je préfère utiliser cette lettre $f_0$ - quitte à suppri-
mer l'indice s'il n'y a pas de confusion
possible. Mon choix de la lettre $p$ pour ce multiplicateur
n'est pas très heureux, car par la suite je
aurai besoin de cette lettre pour désigner
une caractéristique. Je préfère la lettre $\mu$
(initiale grecque de "multiplicateur"), et
la relation de lacets prendrait la forme

$$u(\varepsilon_0) = int(f_0) \varepsilon_0^\mu$$

En fait, je préfère y remplacer $f_0$ par $f_0^{-1}$,
pour obtenir par la suite des formules où la
symétrie des rôles joués par $\rho$ et par $\sigma_\infty$
apparaisse plus nettement, - d'autant que la
formule de définition de
$h_0$ en termes de $g_0$ (cf §43, formule (66))

$$h_0 = g_0^{-1} \sigma_\infty(g_0)$$

devienne plutôt
$$h_0 = g_0 \sigma_\infty(g_0)^{-1}$$

en harmonie avec
$$\xi = x \rho(x)^{-1} \dots$$

Donc la "relation de lacets" prend la forme

(5) $$u(\varepsilon_0) = int(f_0^{-1}) \varepsilon_0^\mu \qquad f_0 \in \hat{\Pi}_{0,3}, \mu \in \hat{\mathbb{Z}}^*,$$

[unclear crossed out line]

(6) $$u(l_1) = int(f_1^{-1}) l_1^\mu$$

et on en déduit aussitôt

\newpage

%% ===== Page 136 =====

(7) $$ \begin{cases} u(\varepsilon_1) = int(f_1^{-1})(\varepsilon_1^\mu), & u(h_1) = int(f_1^{-1})(h_1^\mu) \\ u(\varepsilon_\infty) = int(f_\infty^{-1})(\varepsilon_\infty^\mu), & u(h_\infty) = int(f_\infty^{-1})(h_\infty^\mu) \end{cases} $$

où
(8) $$ f_1 = p(f_0) g^{-1} \, , \quad f_\infty = p(f_1) g^{-1} = p^2(f_0) p(g) g^{-1} . $$

Dans la suite, on écrira simplement $f$ au lieu de $f_0$, et on va transformer l'équation
(8 bis) $$ int(f) u(\varepsilon_0) = \varepsilon_0^\mu $$
en explicitant le premier membre

$$ int(f) u(\varepsilon_0) = f u(\sigma_\infty p) f^{-1} = f u(\sigma_\infty) u(p) f^{-1} = f h_\infty g p f^{-1} $$
$$ = f h_\infty \sigma_\infty g \sigma_\infty^{-1} \underbrace{\sigma_\infty p}_{\varepsilon_0} f^{-1} = f h_\infty \sigma_\infty(g) \underbrace{\varepsilon_0 f^{-1} \varepsilon_0^{-1}}_{\dots} \varepsilon_0 $$

i.e.
(9) $$ int(f) u(\varepsilon_0) = \left[ f h_\infty \sigma_\infty(g) (\sigma_\infty p)(f)^{-1} \right] \varepsilon_0 = \left[ f h_\infty \sigma_\infty(g p(f)^{-1}) \right] \varepsilon_0 $$
où les trois facteurs $f, h_\infty, \sigma_\infty(g p(f)^{-1})$ dans $[\,]$ sont $\in \hat{\Pi}_{0,3}$. Ainsi $(8_{bis})$ [unclear], pour le [unclear]

$$ f h_\infty \sigma_\infty(g p(f)^{-1}) = \varepsilon_0^{\mu-1} $$

ce qui donne, en posant
(10) $$ \mu-1 = 2 \nu \qquad (\nu \in \hat{\mathbb{Z}}) $$
(on note donc $\nu$ ce qui était désigné par $\tau$ précédemment) on trouve la relation

$$ f h_\infty \sigma_\infty(g p(f)^{-1}) = l_0^\nu $$

ou encore, en transformant par $\sigma_\infty$
$$ \sigma_\infty(f) \sigma_\infty(h_\infty) g p(f)^{-1} = l_1^\nu \quad , \text{ soit} $$

\newpage

%% ===== Page 137 =====

(11) $$\overline{\sigma}_\infty(h) \cdot g = \overline{\sigma}_\infty(f)^{-1} l_1 \rho(f)$$

Cette relation, qui relie $h$ et $g$ via $f$, peut
à juste titre s'appeler la "[unclear: relation de lacets]".
On voit, par cette relation, que non seulement
[MARGIN: $f$ déterminé mod le s.g. $l_1^m$ ($m \in \hat{\mathbb{Z}}$) sic.]
la connaissance de $h$ et $g$ détermine $f$ et $\nu$
sic. $f$ et $\nu$, mais aussi inversement, une
fois $f$ et $\nu$ connus, $g$ et $h$ se détermi-
nent univoquement via (11).

Notons que pour $u$ correspondant à
$\xi, g, h, f$, si $\gamma \in \hat{\Pi}_{0,3}$, l'automorphisme intérieur $u'=$
$int(\gamma) \circ u$ correspond à
[MARGIN: Cela nous permet de calculer modulo $Int(\hat{\Pi}_{0,3})$ sur $u$, agissant par $\gamma \in \hat{\Pi}_{0,3}$]
(12) $$g' = \gamma g \rho(\gamma)^{-1}, \quad h' = \gamma h \overline{\sigma}_\infty(\gamma)^{-1}, \quad f' = f \gamma_1$$

Les équations (2) "en $h$" et "en $g$" se
résolvent paramétriquement de façon unique,
par

(13) $$g = \alpha \rho(\alpha)^{-1}, \quad h = \beta \overline{\sigma}_\infty(\beta)^{-1}$$

où

(14) $$\begin{cases} \alpha \in \hat{\Pi}_{0,3}^\sim = \hat{\Pi}_{0,3} \cdot \{1, \varepsilon_0\} \\ \beta \in \hat{\Pi}_{0,3}^\sim \end{cases}$$

de façon précise, la donnée de $g, h$ satisfai-
sant à (2) équivaut resp. à celle de $\alpha, \beta$
dans $\hat{\Pi}_{0,3}^\sim$ et $\hat{\Pi}_{0,3}^\sim$, via les formules (13).

\newpage

%% ===== Page 138 =====

les éléments
$g$ et $h$ étant ainsi explicités en termes des
$\alpha$ et $\beta$, la relation (11), peut être considérée
comme une relation entre $\alpha, \beta, f, \nu$

(15) $$ \varpi_{\infty}(\beta \varpi_{\infty}(\beta)^{-1}) (\alpha \rho(\alpha)^{-1}) = \varpi_{\infty}(f)^{-1} l_1^{\nu} \rho(f) $$

où $\nu$ est déterminé de façon unique
par $\alpha, \beta$, et $f$ est déterminé modulo $l_0^{\hat{\mathbb{Z}}}$.
(On notera que le 2e membre de (15) ne
change pas par le changement de $f$ en
$l_0^n f$ , avec $n \in \hat{\mathbb{Z}}$).

La dissymétrie principale, me semble-t-il,
entre $\varpi_{\infty}$ et $\rho$ dans cette théorie, c'est que
l'on a toujours $\beta \in \widehat{\Pi}_{0,3}$, tandis que
$\alpha$ n'est pas nécessairement dans $\widehat{\Pi}_{0,3}$ - il
l'est, ssi $\nu \equiv 0 \ (3)$ i.e. $\mu \equiv 1 \ (3)$.

L'automorphisme
$$ u' = \operatorname{int}(\beta^{-1}) u $$
aura donc pour composantes :

(16) $$ \begin{cases} g' = \beta^{-1} g \rho(\beta) = (\beta^{-1} \alpha) \rho(\beta^{-1} \alpha)^{-1} \\ h' = \beta^{-1} h \varpi_{\infty}(\beta) = 1 \\ f' = f \beta \end{cases} $$

\newpage

%% ===== Page 139 =====

à un [unclear] près que $u'(\sigma_\infty) = \sigma_\infty$ i.e.
$u'$ qui commute à $\sigma_\infty$. Mais
par des considérations symétriques - prendra
$$u'' = int(\alpha^{-1})u$$
correspondant à
$$ (17) \quad \begin{cases} g'' = \alpha^{-1}g \rho(\alpha) = 1 \\ h'' = \alpha^{-1}h \sigma_\infty(\alpha) = (\alpha^{-1}\beta) \sigma_\infty(\alpha^{-1}\beta)^{-1} \\ f'' = f\alpha \end{cases} $$

mais $u''$ sur $\hat{\Pi}_{0,3}$ - conjugué de $u$ - il
sera conjugué à un élément de [unclear] qui
commute à $\rho$ ssi $g$ est $\hat{\Pi}_{0,3}-\rho$-conjugué
à $1$ i.e. ssi $\alpha \in \hat{\Pi}_{0,3}$ i.e. ssi $\mu \equiv 1 \quad (3)$
ou encore $\nu \equiv 0 \quad (3)$.

Dans tous les calculs des §§ précédents
étaient relatifs au cas où $h=1$, i.e. $u$
commute à $\sigma_\infty$, où il reste les
quantités $g$ et $f$, plus le multiplicateur $\mu$.
[unclear crossed out line]
$\nu = \frac{\mu - 1}{3}$, est déterminé sans
ambiguité par $u$. une fois $\mu$ connu i.e. $\nu$ connu,

[MARGIN: On verra que $\mu$ est déterminé [unclear], de même $g$ et $f$ de façon rigoureuse]
$g$ et $f$ se déterminent mutuellement par
(11) ($f$ étant dans $L \rtimes \hat{\Pi}_{0,3}$).
Pour $f$ fixé, l'équation en $g$ peut s'interpréter
comme une équation en $f$, en posant, en

[MARGIN: Cette notation n'est utile que si $\beta = 1$]
(12) $g = f \sigma_\infty(f)^{-1}$

[MARGIN: c'est essentiellement l'élément noté $h_0$ précédemment]

\newpage

%% ===== Page 140 =====

une "équation en $\varphi$"

(13) $$(l_0^\nu \varphi) p(l_0^\nu \varphi) p^2(l_0^\nu \varphi) = 1 \, ,$$

ou encore, posant

(14) $$\boxed{ \zeta = l_0^\nu \varphi \quad \text{i.e.} \quad \varphi = l_0^{-\nu} \zeta }$$

ce qu'on peut interpréter comme "équation en $\zeta$"

(15) $$\zeta p(\zeta) p^2(\zeta) = 1 \, .$$

Ainsi, les autom. [unclear: de l'ext.] de $\widehat{\Pi}_{0,3}$ qui
avec multiplicateur $\mu = 2\nu+1$,
conservent $\sigma_\infty$, correspondent aux
couples de solutions $(\zeta, f)$ des
équations

(16) $$\boxed{ f \sigma_\infty(f) = 1 \qquad \zeta p(\zeta) p^2(\zeta) = 1 }$$

[MARGIN: On verra que $\nu$ est déterminé dès qu'on se donne $\varphi$ diff. de 1 (ou $\zeta$ diff. de 1)]

satisfaisant de plus la condition de
compatibilité (14) (particulièrement simple
si $\mu=1$ i.e. $\nu=0$, puisqu'alors $\varphi = \zeta$ !). On
récupère $g$ (donc $u$) en fonction
de $(\zeta, f)$ par la formule [unclear: du bas] (11), qui
ici prend la forme

(17) $$g = V_{\sigma_\infty}(f)^{-1} l_1^{-\nu} p(f) = (l_0^\nu f)^{-1} (l_0^{-\nu} \varphi) p(l_0^\nu f)$$

$f \in \widehat{\Pi}_{0,3}$ étant déterminé de façon unique en
fonction de $\varphi$ par la formule (12).

De façon symétrique, on va s'intéresser
aux autom. à [unclear: l'ext.] $u$ de $\widehat{\Pi}_{0,3}$ qui conser-

139

\newpage

%% ===== Page 141 =====

tent à $p$, i.e. ceux pour lesquels [unclear] $g=1$,
décrits par suite en termes de $h,f$,

(18) $$u(\sigma_\infty) = h\sigma_\infty, \quad u(p)=1$$
(19) $$u(\varepsilon_0) = [unclear](\varepsilon_0^H)$$

cette dernière relation prouvant la form. (11),

(20) $$\sigma_\infty(h) = \sigma_\infty f^{-1} l_1^{-\nu} p(f)$$

qui s'écrit aussi, [unclear] $\sigma_\infty$

(21) $$h = f^{-1} \cdot l_0^{-\nu} \cdot (\sigma_\infty p)(f)$$

"L'équation en $h$" (2) [unclear]

(*) $$h \sigma_\infty(h) = 1$$

s'écrit lors comme une "équation en $f$", avec

$$h \sigma_\infty(h) = f^{-1} l_0^{-\nu} \underbrace{(\sigma_\infty p)(f) \sigma_\infty(f)^{-1}}_{\sigma_\infty(p(f) f^{-1}) \\ = \psi^{-1}} l_1^{-\nu} p(f) = [unclear]$$

donc ça s'écrit, en posant $\psi = f p(f)^{-1}$

$$l_0^{-\nu} \sigma_\infty(\psi^{-1}) l_1^{-\nu} \psi^{-1} = 1$$

soit (en passant à l'inverse)

$$\psi l_1^\nu \underbrace{\sigma_\infty(\psi) l_0^\nu}_{\sigma_\infty(\psi l_1^{-\nu})} = 1 \quad \text{soit} \quad l_1^{-\nu} \psi \underbrace{l_1^\nu \sigma_\infty(\psi)}_{\sigma_\infty(l_0^{-\nu} \psi)} = 1$$

Donc posant

(22) $$\psi \overset{df}{=} f p(f)^{-1}$$
(23) $$\eta = l_0^{-\nu} \psi \quad \text{i.e.} \quad \psi = l_0^\nu \eta$$

on trouve que les automorphismes à droite de $\hat{\Pi}_{0,3}$ qui commutent à $p$
s'identifient aux triplets $(\eta, f, \nu)$,
($\eta, f \in \hat{\Pi}_{0,3}, \nu \in \hat{\mathbb{Z}}$ (avec $2\nu+1 \in \hat{\mathbb{Z}}^*$)), satisfaisant

\newpage

%% ===== Page 142 =====

[MARGIN: On trouve, en éliminant $h$ et $\eta$ que]

(24) $$ \psi \rho(\psi) \rho^2(\psi) = 1 \quad , \quad \eta \sigma_{\bar{\infty}}(\eta) = 1 $$

et en plus la relation de compatibilité (23)...
particulièrement simple si $\mu=1$ i.e. si $\nu=0$,
auquel cas elle se simplifie $\psi=\eta$ ! Mais en
plus de ces relations, il faut supposer que
la solution $\psi$ de l'équation (24) "en $\psi$"
"provient de la feuille principale", i.e. que,
quand on écrit $\psi$ sous forme paramétrique

(25) $$ \psi = f \rho(f)^{-1} $$

avec $f \in \hat{\overline{\Pi}}_{0,\bar{s}} = \hat{\Pi}_{0,\bar{s}} \cdot \{1,\tau_{\infty}\}$, on ait bien
$f \in \hat{\Pi}_{0,\bar{s}}$ (donc que l'image de $f$ dans
$\hat{\overline{\Pi}}_{0,\bar{s}} / \hat{\Pi}_{0,\bar{s}} \simeq \{1,\tau_{\infty}\} \simeq \mathbb{Z}/2$ soit nulle...).
Ayant ainsi déterminé $f$ en fonction de $\psi$, on
détermine $h$ par (21), d'où [unclear].
Remarques 1) Si
cette condition supplémentaire sur $\psi$ n'est pas
satisfaite, la "dite" formule (21) garde un
sens et définit $h \in \hat{\overline{\Pi}}_{0,\bar{s}}$, satisfaisant
"l'équation en $h$" $h \sigma_{\bar{\infty}}(h) = 1$ donc définissant
un autom. de $\hat{\Pi}_{0,\bar{s}}$, commutant à $\rho$.

On aura dans
ce cas
$$ u(l_0) = int(f^{-1}) (l_0^x) $$
ou encore, posant
$$ f = f_0 \tau_{\infty} \text{ avec } f_0 \in \hat{\Pi}_{0,\bar{s}} $$
$$ u(l_0) = int(f_0^{-1}) (\tau_{\infty}(l_0^x)) = int(f_0^{-1})(l_0^{-x}) $$

\newpage

%% ===== Page 143 =====

i.e. $u$ est un automorphisme extérieur, de multi-
plicateur $-\mu = -(2\nu + 1)$.

2) Notons que l'indétermination de $f$ (pour [unclear] fixé,
i.e. $g,h$ fixés) est dans
(25) $$f \mapsto f' = l_0^n f \quad (n \in \hat{\mathbb{Z}})$$
ce qui donne

$$ \varphi \mapsto \varphi' = l_0^n \varphi l_1^{-n}, \quad \psi \mapsto \psi' = l_0^n \psi l_1^{-n} $$
$$ \xi \mapsto \xi' = l_0^{-n} \xi l_1^{-n}, \quad \eta \mapsto \eta' = l_0^n \eta l_1^{-n} $$

[MARGIN: NB on a de même pour
$\alpha \mapsto \alpha' = l_0^n \alpha l_0^{-n}$
$\beta_0 \mapsto \beta_0' = l_0^n \beta_0 l_0^{-n} \quad (26)$]

(NB -- $l_1 = \sigma_\infty(l_0) = \rho(l_0)$, donc $\xi'$ et $\eta'$ sont les
$\rho$-transformés de $\xi, \eta$ par $l_0^{-n}$, et $\varphi', \psi'$ les
$\sigma_\infty$-transformés de $\varphi, \psi$ par $l_0^n$ -- en particulier
les équations (16) (en $\varphi$ et $\psi$) et (24)
(en $\xi$ et $\eta$) sont bien invariantes par
ces transformations, ainsi d'ailleurs que les
relations de compatibilité (14) resp. (23)!
(réécriture de (11))

3) On a les formules suivantes

(27) $$\sigma_\infty(h) g = (l_0^\nu f)^{-1} \xi \rho(l_0^\nu f) = \sigma_\infty(f^{-1} l_0^{-\nu}) \eta (f^{-1} l_0^{-\nu})^{-1}$$

(en $h=1$) expriment $g$ comme le $\rho$-conjugué
[unclear] de $\xi$ par $(l_0^\nu f)^{-1}$, et de façon symétrique
dans le cas $g=1$, comme le $\sigma_\infty$-conjugué
de $\eta$ par $\sigma_\infty(f^{-1} l_0^{-\nu})$. Si $h=1$, c'est la
première expression qui est [unclear: retenue]
c'est-à-dire l'équation en $g$), pour $g=1$, c'est
la seconde, qui permet de réécrire (aux 2 membres)
(28) $$h = (f^{-1} l_0^{-\nu}) (\eta^{-1}) \sigma_\infty(f^{-1} l_0^{-\nu})^{-1}$$

\newpage

%% ===== Page 144 =====

exprimant $h$ comme $\sigma_\infty$-transformé de $\eta^{-1}$
par $f^{-1} l_0^\sim$. Si on veut l'exprimer comme
$\sigma_\infty$-transformé [MARGIN: $\sigma_\infty^{(kg)}=$] de $\eta$ (au lieu de $\eta^{-1}$), on écrira
$\eta^{-1} \sigma_\infty(\eta) = \eta^{-1} . \eta . \sigma_\infty(\eta)$

d'où
(28 bis) $h = (f^{-1} l_0^\sim \eta^{-1}) \eta \sigma_\infty(f^{-1} l_0^\sim \eta^{-1})^{-1}$.

Les formules permettent de mettre en
relation, dans le cas $h=1$ les solutions des deux
équations
(29) $g = \alpha \rho(\alpha)^{-1}$ et $\xi = \alpha_0 \rho(\alpha_0)^{-1}$
par
(30) $\alpha = (l_0^\sim f)^{-1} \alpha_0 = f^{-1} l_0^{\sim -1} \alpha_0$

et dans le cas où $g=1$, les solutions des
deux équations
(31) $h = \beta \sigma_\infty(\beta)^{-1}$ , $\eta = \beta_0 \sigma_\infty(\beta_0)^{-1}$
par
(32) $\beta = (f^{-1} l_0^\sim \eta^{-1}) \beta_0$ .

La dissymétrie entre ces deux formules
( (32) étant un peu plus compliquée que
(30) par le facteur supplémentaire $\eta^{-1}$
au coefficient de $\beta_0$ ) provient du
fait que $f$ a été défini via la première
des [unclear] (6) , première forme , où $l_0$ est défini
comme la [unclear] de $\sigma_{0\vec{1}} = \varepsilon_0$, et non de $\rho \sigma_\infty = \varepsilon_1$.

\newpage

%% ===== Page 145 =====

Si on identifie maintenant $\widehat{S}_{0,3}$ à $\operatorname{Gl}(2, \hat{\mathbb{Z}})^{\wedge}$,
et qu'on regarde les images des éléments
$g, h, f$ etc dans le quotient
$$\operatorname{Gl}(2, \hat{\mathbb{Z}})^{\pm} = \operatorname{Gl}(2, \hat{\mathbb{Z}})^{\wedge} / \pm 1$$
- qu'on désignera par les m[ême]s lettres, pour sim-
plifier le langage - on sera amené à calculer
avec des éléments de $\operatorname{Gl}(2, \hat{\mathbb{Z}})$, plutôt que
de $\operatorname{Gl}(2, \hat{\mathbb{Z}})^{\pm}$. On choisira $\alpha, \beta$ dans
$\hat{\Pi}_{0,3}$ et $\hat{\Pi}_{0,3}'$ conformément à (13), d'où
$$(33) \qquad \alpha, \beta \in \operatorname{Gl}(2, \hat{\mathbb{Z}})^{\pm 1} \quad \text{modulo signes},$$
mais les éléments correspondants
$$(34) \qquad g = \alpha p(\alpha)^{-1}, \quad h = \beta \sigma_{\infty} p(\beta)^{-1} \quad (\in \operatorname{Sl}(2, \hat{\mathbb{Z}}))$$
sont déterminés dans $\operatorname{Gl}(2, \hat{\mathbb{Z}})$ sans ambiguïté
de signe. De même
$$(35) \qquad \alpha_0 \text{ (cas } h=1), \quad \beta_0 \text{ (cas } g=1) \in \operatorname{Gl}(2, \hat{\mathbb{Z}})^{\pm 1} \quad \text{modulo signe},$$
et alors
$$(36) \qquad \xi = \alpha_0 p(\alpha_0)^{-1}, \quad \eta = \beta_0 \sigma_{\infty} p(\beta_0)^{-1} \in \operatorname{Sl}(2, \hat{\mathbb{Z}})$$
sont déterminés sans ambiguïté de signe. Ainsi
[unclear: crossed out] les éléments
$$(37) \qquad f \in \operatorname{Sl}(2, \hat{\mathbb{Z}}) \text{ mod signes } (\text{pour } f \in \widehat{S}_{0,3} \text{ fixé})$$
et on a alors
$$(38) \qquad g = f \sigma_{\infty} f^{-1}, \quad \eta = f p(f)^{-1} \in \operatorname{Sl}(2, \hat{\mathbb{Z}})$$
sans ambiguïté de signe.
[unclear: crossed out] Dans le choix des représentants

\newpage

%% ===== Page 146 =====

de $\varepsilon_0 \in \operatorname{Gl}(2, \hat{\mathbb{Z}})'$ dans $\operatorname{Sl}(2, \hat{\mathbb{Z}})$ il y a
ambiguité de signes
$$ \varepsilon_0 = \pm \begin{pmatrix} 1 & -1 \\ 0 & -1 \end{pmatrix} \quad , $$
[MARGIN: encore que le signe ne serait qu'un choix secondaire, on aimera mieux un élément unipotent, engendrant un [unclear: pro-p-groupe]...]

mais pour
(39) $$ l_0 = \varepsilon_0^2 = \begin{pmatrix} 1 & -2 \\ 0 & 1 \end{pmatrix} $$
il n'y a pas d'ambiguité de signes. Ceci
étant, les relations
(40) $$ \xi = l_0^\nu \varphi \qquad (\text{cas } h=1) $$
et
(41) $$ \eta = l_0^{-\nu} \psi \qquad (\text{cas } g=1) $$
sont elles vraies sans signes correcteurs ??
Le cas $h=1$ correspond aux données
(dans $\hat{\Pi}_{0,3}$) de $\xi$ et $f$, satisfaisant
$$ \xi \rho(\xi) \rho^2(\xi) = 1, \quad \xi = l_0^\nu \left( f \sigma_{\infty}(f)^{-1} \right) $$
et l'équation en $f$ s'écrit
ou mieux aux données $\alpha_0 \in \hat{\Pi}_{0,3}$, $f \in \hat{\Pi}_{0,3}$, $\nu \in \hat{\mathbb{Z}}$
satisfaisant les relations
(42) $$ \alpha_0 \rho(\alpha_0) = l_0^\nu f \sigma_{\infty}(f)^{-1} $$
en termes de ces données on définit
(43)
$$ \begin{cases} \xi = \alpha_0 \rho(\alpha_0) \\ \varphi = f \sigma_{\infty}(f)^{-1} \\ g = \sigma_{\infty}(f)^{-1} l_1^{-\nu} \rho(f) = (l_0^\nu f)^{-1} (l_0^\nu \varphi) \rho(l_0^\nu f) \end{cases} $$
Ceci en supposant qu'on remonte $\alpha_0, f$ dans $\operatorname{Gl}(2, \hat{\mathbb{Z}})^{\pm 1}$

\newpage

%% ===== Page 147 =====

353

sans ambiguïté de signes, [crossed out: et] [crossed out]
sont bien déterminées dans $\operatorname{Sl}(2, \hat{\mathbb{Z}})$, indépendamment
des choix du signes, [crossed out] et de $g$.

Au moins donc dans $\operatorname{Sl}(2, \hat{\mathbb{Z}})$

[crossed out equation]

(44) $$ \xi = \varepsilon l_0^\nu \varphi $$ i.e. $$ \alpha_0 \rho(\alpha_0) = \varepsilon l_0^\nu f \sigma_\infty(f)^{-1} $$
dans $\operatorname{Sl}(2, \hat{\mathbb{Z}})$

un signe $\varepsilon \in \{\pm 1\}$ bien déterminé, et
la question qui se pose est de savoir si l'on a $\varepsilon = +1$.
Notons que la formule j'espère
[crossed out: Définitive (?)]
$$ \rho(l_0) = l_1 $$
est vraie dans $\operatorname{Sl}(2, \hat{\mathbb{Z}})$ (sans signe correcteur)
donc l'expression de $g$ dans (43) est valable dans
$\operatorname{Sl}(2, \hat{\mathbb{Z}})$ (aux signes près), ce qui montre que
$g \in \operatorname{Sl}(2, \hat{\mathbb{Z}})$ [crossed out] est un $\rho$-transformé
de $l_0^\nu \varphi$, donc cette dernière est un $\rho$-transformé
de $g_0$. On sait que $g$ satisfait [crossed out] la
relation
$$ g \rho(g) \rho^2(g) = \varepsilon_0 \cdot 1 \qquad \varepsilon_0 \in \{\pm 1\} $$

[crossed out: $(l_0 \varphi) \rho(l_0 \varphi) \dots$]

et de même on aura d'ailleurs, si $g' = \alpha g \rho(\alpha)^{-1}$
$(\alpha \in \operatorname{Gl}(2, \hat{\mathbb{Z}}))$, que
$$ g' \rho(g') \rho^2(g') = \varepsilon_0 \cdot 1 $$
donc en particulier pour $g' = l_0^\nu \varphi = \varepsilon \xi$ (44)
on a $\xi$ donné par (42) satisfait à
$$ \xi \rho(\xi) \rho^2(\xi) = 1 $$

(NB on n'a pas $\rho^3 = 1$ dans $\operatorname{Sl}(2, \hat{\mathbb{Z}})$, mais
$\rho^3 = -1$, ce qui implique que l'automorphisme int.
[unclear: défini par $\rho$ est d'ordre 3] ... ), d'où
$$ \varepsilon_0 = \varepsilon $$

\newpage

%% ===== Page 148 =====

i.e. la question de la validité de la formule
$$ \xi = l_0^\vee \varphi $$
dans $\operatorname{Sl}(2, \hat{\mathbb{Z}})$ équivaut à celle de la validité
de la formule
$$ g \rho(g) \rho^2(g) = 1 \quad , $$
que nous laissons en suspens pour le
moment. Notons pour mémoire les
formules suivantes dans $\operatorname{Gl}(2, \hat{\mathbb{Z}})$
($\det \alpha_0 = \pm 1$, $\det f = 1$)
$\alpha_0, f \in \operatorname{Gl}(2, \hat{\mathbb{Z}})$ (à des signes arbitraires
près dans les choix)

[MARGIN: Il me semble maintenant plus sain de définir $g$ par $g = (l_0^\vee fj)^{-1} l_0^\vee \varphi \rho(l_0^\vee f)$ ce qui donne la relation $g \rho(g) \rho^2(g) = \varepsilon^3 = \varepsilon_0$.
NB on aura [unclear: pas] $\varepsilon = \varepsilon_0$]

(45)
(Cas $h=1$)
$$ \xi \overset{df}{=} \alpha_0 \rho(\alpha_0)^{-1} $$
$$ \varphi \overset{df}{=} f \sigma_\infty (fj)^{-1} \quad | \quad \in \operatorname{Sl}(2, \hat{\mathbb{Z}}) $$
$$ [ g \overset{df}{=} \sigma_\infty (fj)^{-1} l_0^\vee \rho(f) = (l_0^\vee fj)^{-1} (l_0^\vee \varphi) \rho(l_0^\vee f) ] $$
$$ \xi = \varepsilon l_0^\vee \varphi \quad \varepsilon \in \{\pm 1\} $$
$$ [ g \rho(g) \rho^2(g) = \varepsilon . 1 $$
$$ g = \varepsilon . \alpha \rho(\alpha)^{-1} \quad \alpha = (l_0^\vee fj)^{-1} \alpha_0 = f^{-1} l_0^{-\vee} \alpha_0 ] $$

De façon symétrique
$\beta_0, f \in \operatorname{Sl}(2, \hat{\mathbb{Z}})$ (à des signes près, i.e. les choix)

(46)
(Cas $g=1$)
$$ \eta \overset{df}{=} \beta_0 \sigma_\infty (\beta_0)^{-1} $$
$$ \psi \overset{df}{=} f \rho(fj)^{-1} $$
$$ h \overset{df}{=} f^{-1} l_0^{-\vee} \sigma_\infty \rho(fj) = $$

\newpage

%% ===== Page 149 =====

364

Je vais reprendre le cas $h=1$ ($cf (45)$), en posant

[MARGIN: NB Dans le cas "arithmétique" loc. cit. provenant de par ex. $\Pi_{0,3} \dots$ j, $\alpha_0, f \equiv 1 (2)$ i.e. $a,d,r,u \equiv 1 (2)$ $b,c,s,t \equiv 0 (2)$]
(47)
$$ \alpha_0 = \begin{pmatrix} a & b \\ c & d \end{pmatrix} $$
$$ f = \begin{pmatrix} r & s \\ t & u \end{pmatrix} $$
[MARGIN: [unclear: changements] de signes dus aux [unclear: lignes] co-invariants]

où $a,b,c,d, r,s,t,u \in \hat{\mathbb{Z}}$, avec

(48) $$ ad-bc = \delta \in \{\pm 1\} $$
(49) $$ ru-st = 1 $$

Ceci pose

(50) $$ \xi = \alpha_0 \rho(\alpha_0)^{-1} = \begin{pmatrix} A & B \\ C & D \end{pmatrix} $$
(51) $$ \varphi = f \sigma_0(f)^{-1} = \begin{pmatrix} R & S \\ T & U \end{pmatrix} = \begin{pmatrix} r^2+s^2 & rt+su \\ rt+su & t^2+u^2 \end{pmatrix} $$

de sorte qu'on aura (cf § 42, formules 78)

[MARGIN: NB on a $\delta = \pm 1$]
(52)
$$ \delta A = (a^2+ab+b^2) - (ac+bd+bc) \qquad \delta B = ac+bd+bc $$
$$ \delta C = ac+bd+ad - (c^2+d^2+cd) \qquad \delta D = c^2+d^2+cd $$

et ( § 42, (90) et 91)

(53) $$ R = r^2+s^2 , \quad S = T = rt+su , \quad U = t^2+u^2 $$

On aura les formules

(54) $$ C+D-B = 1 $$
(55) $$ S = T $$

La relation essentielle reliant $\alpha_0$ et $f$, i.e. reliant les "variables" $a,b,c,d$ aux $r,s,t,u$, est la relation

[MARGIN: "relation de base"]
(56) $$ \xi = \varepsilon l_0^\nu \varphi \quad , \quad \varepsilon \in \{\pm 1\} $$

qui s'écrit aussi $l_0^{-\nu} \xi = \varepsilon \varphi$, i.e.

148

\newpage

%% ===== Page 150 =====

$$(57) \quad \left( \begin{matrix} A+2\nu C & B+2\nu D \\ C & D \end{matrix} \right) = \varepsilon \left( \begin{matrix} R & S \\ T & U \end{matrix} \right) .$$

Mais cette condition implique d'abord que
la matrice du premier membre est symétrique
(puisque $S=T$) i.e.
$$2\nu D = C - B$$
condition qui équivaut [unclear], par (54), à
$$2\nu D = 1 - D \quad \text{i.e.} \quad D = 1/\mu \,,\, \mu^{-1} = D$$

$$(58) \quad \mu D = 1 \quad \text{où} \quad \mu = 2\nu + 1 ,$$

ce qui implique déjà $\mu, D$ inversibles.

si on écrit (56) sous la
forme
$$\varepsilon \left( \begin{matrix} A & B \\ C & D \end{matrix} \right) = \left( \begin{matrix} R-2\nu T & S-2\nu U \\ T & U \end{matrix} \right) ,$$

la condition $C+D-B=1$ devient, sur
les éléments de la matrice $\varphi$
$$T + U - (S - 2\nu U) = \varepsilon \quad ;$$

i.e. $\quad (2\nu + 1)U = \varepsilon \quad$ i.e.
$$(59) \quad \mu U = \varepsilon \quad \text{i.e.} \quad U = \varepsilon/\mu \,,\, \mu = \varepsilon/U$$

ce qui équivaut aussi, tenant compte de
(58), à la relation
$$(60) \quad D = \varepsilon U$$

qui n'est autre que l'égalité des [unclear] des
[unclear] des matrices (57).

\newpage

%% ===== Page 151 =====

365

On peut expliciter les relations (57) comme
quatre relations sur les $9$ variables $a,b,c,d, r,s,t,u, \mu$
en plus des relations (48), (49) ;

(61)
$$D = \varepsilon U$$
$$\mu D = 1 \quad \text{i.e.} \quad \mu U = \varepsilon$$
$$C = \varepsilon T$$
$$A = \varepsilon R - \frac{(\mu-1)}{2\nu} C$$
[MARGIN: $A, B$ [unclear: calculés en fonction des autres par] $C+D-B=1$ et $AD-BC=1$]

(62)
[MARGIN: Conséquences des précédentes]
soit
[crossed out: $\mu(c^2+d^2+cd) = \delta$]
$$b^2+d^2+cd = \varepsilon\delta(t^2+u^2)$$
$$ac+bd+ad-(c^2+d^2+cd) = \delta\varepsilon(rt+su)$$
$$a^2+b^2+ab-(ac+bd+bc) = \delta \varepsilon (r^2+s^2 - (\mu-1)(rt+su))$$

Quand les signes $\varepsilon, \delta$ sont fixés, on trouve
avec (48), (49)
un système de $6$ équations en $9$ variables
$a,b,c,d,r,s,t,u,\mu$ - on doit pouvoir vérifier
que le schéma (sur $Spec \mathbb{Z}$) des solutions
de ces équations est de dim. relative $3$,
[MARGIN: non, il va être de dim. relative $4$]
et de dim relative $1$ sur les "facteurs"
$(a,b,c,d)$ resp. $(r,s,t,u)$. On a donc un
[unclear: point à valeurs dans] un schéma absolu -
une fois [unclear: passés] à coefficients dans $\hat{\mathbb{Z}}$...

Théorème Le multiplicateur $\mu$ est déterminé en fonction
de $\xi = \begin{pmatrix} A & B \\ C & D \end{pmatrix}$ (considéré comme élément de $\operatorname{Sl}(2,\hat{\mathbb{Z}})$) par la
relation
$$\mu D = 1 \quad (58)$$
et il est déterminé au signe près par $\mathcal{U} = \begin{pmatrix} R & S \\ T & U \end{pmatrix}$
par la relation
$$\mu U = \varepsilon \quad (59) \quad \text{où } \varepsilon \in \{ \pm 1 \},$$

\newpage

%% ===== Page 152 =====

et où $\varepsilon$ est le signe qui intervient dans les formules

$$ \begin{cases} [unclear] = \varepsilon l_0 \varphi & (56) \\ g \rho g_1 \rho g_1 = \varepsilon & (45) \end{cases} $$
$g = \varepsilon \alpha \rho \alpha^{-1}$

On a $\varepsilon = 1$ ssi $\mu \equiv 1 \pmod 4$, $\varepsilon = -1$ ssi $\mu \equiv -1 \pmod 4$.
(i.e. $\nu \equiv 0 \pmod 2$) \hfill (i.e. $\nu \equiv 1 \pmod 2$)

[MARGIN: en d'autres termes on a $\mu \equiv \varepsilon \pmod 4$
[crossed out: $\mu \equiv 1 \pmod 4$]
donc $\mu = \varepsilon/U \equiv \varepsilon \pmod 4$.
Ce serait le cas id. si on avait $\mu \equiv \delta \pmod 3$,
relation d'ailleurs démontrée plus bas...]

Pour le prouver, on se rappelle que

$$ U = t^2 + u^2 $$

on note que les seules valeurs mod 4 prises
par $t^2$ et $u^2$ sont $0$ et $1$, donc celles prises par
$t^2+u^2$ sont $0, 1$ et $2$, et la seule unité parmi celles-
ci est $1$. Donc [unclear] mod 4 [unclear] soit $1/\mu \equiv \varepsilon \pmod 4$
ou encore $\mu \equiv \varepsilon \pmod 4$ CQFD.

Ainsi on aurait vu de façon toute [unclear: analogue],
grâce à la formule $p D = 1 \quad (58)$ et
$$ D = \delta(c^2+d^2+cd) $$
que $\mu \equiv \delta \pmod 3$ (i.e. $\delta = \pm 1$ (resp. $-1$) ssi
$\mu \equiv 1 \pmod 3$ (resp. $\mu \equiv -1 \pmod 3$)). On profitera d'ailleurs

notons, pour la commodité, par $\varepsilon_p$ la valeur
du "caractère quadratique" relatif à $p$,

[MARGIN: $(\mathbb{Z}/p\mathbb{Z})^* / ((\mathbb{Z}/p\mathbb{Z})^*)^2 \simeq \{ \pm 1 \}$]

égal à : $1$ ssi $p \in (\mathbb{Z}_p^*)^2$. On aura donc

(63) $$ \begin{cases} \delta = ad - bc = \left(\frac{p}{3}\right) = \varepsilon_3 \\ \varepsilon = D/U = \left(\frac{-1}{p}\right) = \varepsilon_2 \end{cases} $$

et la connaissance de $\varepsilon_3$ (resp. $\varepsilon_2$, resp. des deux)
équivaut à celle de $p \bmod 3$ (resp. mod 4, resp. mod 12).

\newpage

%% ===== Page 153 =====

Nous allons aussi calculer (cf formules (45))
$$ \alpha = (l_0^v f)^{-1} \alpha_0 = f^{-1} l_0^{-v} \alpha_0 $$
$$ g = \varepsilon \alpha \rho(\alpha)^{-1} = \underbrace{\sigma_0(f^{-1}) l_1^v}_{t_f} \rho(f) = (l_0^v f)^{-1} \xi\ \rho(l_0^v f) $$

On trouve
$$ (64)\quad \alpha = \begin{pmatrix} u & -s \\ -t & r \end{pmatrix} \begin{pmatrix} 1 & 2v \\ 0 & 1 \end{pmatrix} \begin{pmatrix} a & b \\ c & d \end{pmatrix} = \begin{pmatrix} u(a+2vc)-cs & u(b+2vd)-ds \\ -t(a+2vc)+cr & -t(b+2vd)+dr \end{pmatrix} $$

[MARGIN: NB Contrairement à (64) et à ce qui en est dit au début du 3ème, il n'est pas du tout évident de l'identité des deux membres de (64), d'où l'égalité d'un des $g$, à moins de le faire à la main. Mais on peut aussi le vérifier à partir de (65)]

$$ f^{-1} \begin{pmatrix} a+2vc & b+2vd \\ c & d \end{pmatrix} $$

$$ g = \begin{pmatrix} r & t \\ s & u \end{pmatrix} \begin{pmatrix} 1 & 0 \\ 2v & 1 \end{pmatrix} \begin{pmatrix} t+bl & -t \\ ... & ... \end{pmatrix} = \begin{pmatrix} r & t \\ s & u \end{pmatrix} \begin{pmatrix} t+bl & -t \\ \mu(t+bl)+tas & -\mu t + s \end{pmatrix} $$

$$ \sigma_0(f^{-1}) l_1^v\ \rho(f) \quad \text{(cf § 43 formule (77))} $$
$$ (65)\quad = \begin{pmatrix} 1 + \mu t(t+u) & -\mu t^2 \\ -1 + \mu u(t+u) & 1 - \mu u t \end{pmatrix} $$

Il apparait que $g$ s'exprime simplement en termes de $f$ i.e. de $r, s, t, u$ directement, et non en termes de $\alpha_0$ ou $\xi$ i.e. via les $a, b, c, d$ — A,B,C,D.
Par contre, pour exprimer $\alpha$ (64), il faut faire intervenir à la fois $r, s, t, u$ et les $a, b, c, d$.

Il était naturel, en relation avec l'homomorphisme
$$ \mathcal{N}_{1,1} \to \operatorname{Gl}(2, \hat{\mathbb{Z}}) $$
(cf § 39, (1)) de chercher une matrice qui pre-
serve $\varepsilon - g$, je vais désigner par la même lettre,
$u$, par l'automorphisme envisagé de $\hat{\Pi}_{1,3}$ dont il
était parti (tout comme je désigne par les mêmes
lettres les éléments $g, f$ etc de $\hat{\Pi}_{0,3}$ ou de $\hat{B}_3$,
et leurs images dans $\operatorname{Gl}(2, \hat{\mathbb{Z}})$, voire même des éléments
correspondants dans $\operatorname{Gl}(2, \mathbb{Z}/l)$ lui-même),

\newpage

%% ===== Page 154 =====

$$u \in \operatorname{Gl}(2, \hat{\mathbb{Z}})$$
autom. de déterminant égal à $\mu$, et induisant sur
$\operatorname{Sl}(2, \hat{\mathbb{Z}})$, par autom. intérieurs, le même autom.
[unclear: que $u_{\bar{f}}$ pour $g$ en question], ce qui s'explicite
par les formules
(66) $$det\ u = \mu$$
(67) $$u \sigma_\infty u^{-1} = \varepsilon' \sigma_\infty \qquad \varepsilon' \in \{\pm 1\}$$
(68) $$u \rho u^{-1} = \varepsilon'' g \rho \qquad \varepsilon'' \in \{\pm 1\}$$

En fait, [unclear: la relation] (68) s'écrit
$$[u, \rho] = u \rho u^{-1} \rho^{-1} = \varepsilon'' g$$
et appliquant aux deux membres l'opération
$x \mapsto x \rho(x) \rho^2(x)$, et tenant compte que $g \rho(g) \rho^2(g) = \varepsilon_2$,
et $x \rho(x) \rho^2(x) = 1$ pour $x = u \rho u^{-1}$, on trouve
(68 bis) $$\varepsilon'' = \varepsilon_2$$
et (68) équivaut donc à
$$u \rho u^{-1} = \varepsilon_2 g \rho = \alpha \rho \alpha^{-1}$$
d'où i.e. à
(69) $$u = \alpha u_0 \qquad \text{avec} \quad [u_0, \rho] = 1 \qquad (\text{i.e. } u_0 \rho u_0^{-1} = \rho).$$
Nous avons déjà remarqué que (67), (68) déter-
minent $u$ seulement modulo multiplication par
un scalaire de $\hat{\mathbb{Z}}^*$, et compte tenu de (66), $u$
est déterminé seulement modulo $\lambda$ satisfaisant
$\lambda^2 = 1$ i.e. modulo $\hat{\mathbb{Z}}^*_2 \simeq \{\pm 1\}^P$ (où $P$ est l'ens.
des nbs premiers), mais on espère néanmoins
trouver un choix canonique.

Dans mes premiers calculs, j'avais pris comme
départ (69), où la condition $u_0$ commutant à $\rho$
s'exprimait sous la forme
(70) $$u_0 = p + q \rho = \begin{pmatrix} p & -q \\ q & p-q \end{pmatrix} \quad (p,q \in \hat{\mathbb{Z}})$$

\newpage

%% ===== Page 155 =====

357

de sorte que

(71) $$det u_0 = p^2 + p q + q^2$$

et que (66) s'écrit (compte tenu de (67), où
[unclear: det barré] $\delta = \varepsilon_3$)

(72) $$p^2 + pq + q^2 = \mu \varepsilon_3$$

ce qui, compte tenu de la relation
$c^2 - cd + d^2 = \varepsilon_3 / \mu$ (52 et 58)
suggérait des relations de la forme
$$p = \mu c, q = \mu d, \text{ ou } p = \mu d, q = \mu c$$

En fait, exprimant (67) comme l'équation en $p, q$
(s'ajoutant à 72), je trouvais (aux signes près bien sûr,
que j'ai vite immolé)

(73) $$p = \mu d, q = \mu c \quad \text{i.e.} \quad u_0 = \mu(c \rho + d)$$

mais en me plaçant dans le cas $\varepsilon' = +1$ exclusivement.
(J'avais cru conclure, un peu trop hâtivement, que le
signe $\varepsilon'$ était toujours $+1$, avant de m'apercevoir
que ce n'était guère compatible avec le corps des
classes, et de subodorer que l'on devait avoir $\varepsilon' = +1$
si $p \equiv 1$ (4) .) En faisant ce détour assez pesant
par les $(a,b,c,d)$, je suis quand même arrivé, après
bien des péripéties (erreurs de calcul etc) à la
formule où il n'y avait plus de $a,b,c,d$ du tout,
et d'une simplicité déroutante :

[MARGIN: Je souligne ici la lettre $\underline{u}$ qui désigne un des deux [unclear] d'intégration, pour la distinguer de $u$ qui correspond manifestement au cas $\varepsilon'_2 = 1$, i.e. celui où $u$ commute à $\sigma_\infty$.]

(74) $$u = \mu(t \sigma_\infty + \underline{u}) = \mu \begin{pmatrix} \underline{u} & t \\ -t & \underline{u} \end{pmatrix} , \quad (\iff \text{sans doute, } \varepsilon_2 = 1)$$

\newpage

%% ===== Page 156 =====

$$ \det(t\sigma_\infty+u) = t^2+u^2 = \varepsilon_2 / \mu (= 1/\mu \text{ si } \varepsilon_2=+1) ,$$
on retrouve bien, dans le cas $\varepsilon_2=1$, que
$$ \det \underline{u} = \mu^2 \det(t\sigma_\infty+u) = \mu^2 1/\mu = \mu $$
i.e. (66). D'autre part, on a ($\mu$)

[MARGIN: on pourrait [unclear] les formules § 93 (78)]

$$ [u,\rho] = [t\sigma_\infty+u, \rho] = \begin{pmatrix} u & t \\ -t & u \end{pmatrix} \begin{pmatrix} 0 & 1 \\ -1 & 1 \end{pmatrix} \begin{pmatrix} u & -t \\ t & u \end{pmatrix} \begin{pmatrix} 1 & -1 \\ 1 & 0 \end{pmatrix} $$
$$ = \mu \begin{pmatrix} -t & u+t \\ -u & u-t \end{pmatrix} \begin{pmatrix} u-t & -u \\ u+t & -t \end{pmatrix} $$
$$ = \mu \begin{pmatrix} (u+t)^2-ut+t^2 & -t^2 \\ ut-t^2 & u^2+t^2-ut \end{pmatrix} = \mu \begin{pmatrix} \frac{1}{\mu} + t(u+t) & -t^2 \\ t(u-t) & \frac{1}{\mu}-ut \end{pmatrix} $$
$$ = \begin{pmatrix} 1+\mu t(u+t) & -\mu t^2 \\ \mu t(u-t) & 1-\mu u t \end{pmatrix} ,$$

et en comparant avec (65) on trouve bien
$$ [u,\rho] = g $$
compte tenu de
$$ \mu(u^2+t^2) = \varepsilon_2 = 1 $$
donc
$$ -1 + \mu u(t+u) = \underbrace{-1 + \mu u^2}_{-\mu t^2} + \mu u t = \mu t(u-t) ! $$

J'aimerais maintenant, sans un retour à ces
nouveaux calculs, trouver "au pif" une expression
pour $\underline{u}$, dans le cas où $\varepsilon_2 = -1$, et où on s'attend
à trouver dans (67) le signe $\varepsilon' = -1$ - i.e. que
$\underline{u}$ soit une anti-similitude [unclear] de la
$$ \underline{u}^t \underline{u} = - \det \underline{u} $$
Mais la matrice
$$ \tau_\infty = \begin{pmatrix} -1 & 0 \\ 0 & 1 \end{pmatrix} \qquad \text{(c'est une reflexion orthog.)} $$
est le cas-type d'une telle antisimilitude, et je
prévois que la matrice
$$ \underline{u} = \mu \tau_\infty (t\sigma_\infty+u) = \mu \begin{pmatrix} -1 & 0 \\ 0 & 1 \end{pmatrix} \begin{pmatrix} u & t \\ -t & u \end{pmatrix} = \mu \begin{pmatrix} -u & -t \\ -t & u \end{pmatrix} $$

\newpage

%% ===== Page 157 =====

conviendra. On aura, bien sûr
$$u^t u = \mu^2 \begin{pmatrix} u & t \\ t & -u \end{pmatrix} \begin{pmatrix} u & t \\ t & -u \end{pmatrix} = \mu^2 \begin{pmatrix} u^2+t^2 & 0 \\ 0 & u^2+t^2 \end{pmatrix} = \mu^2 \begin{pmatrix} -\frac{1}{\mu} & 0 \\ 0 & -\frac{1}{\mu} \end{pmatrix}$$
i.e.
$$u^t u = - \mu . Id$$
et
$$\det u = \mu^2 (-u^2 - t^2) = \mu^2 ( \frac{1}{\mu} ) = \mu$$
compte tenu que maintenant
$$u^2+t^2 = \varepsilon_2 / \mu = - \frac{1}{\mu} \qquad \text{puisque } \varepsilon_2 = -1 .$$
Enfin pour vérifier (68) calculons
$$[u, \rho] = [ \begin{pmatrix} u & t \\ t & -u \end{pmatrix}, \rho ] = \begin{pmatrix} u & t \\ t & -u \end{pmatrix} \begin{pmatrix} 0 & 1 \\ -1 & 1 \end{pmatrix} \mu \begin{pmatrix} -u & -t \\ -t & u \end{pmatrix} \begin{pmatrix} 1 & -1 \\ 1 & 0 \end{pmatrix}$$
$$= \mu \begin{pmatrix} -t & u+t \\ -u & -u+t \end{pmatrix} \begin{pmatrix} -u-t & -u \\ u-t & t \end{pmatrix}$$
$$= \mu \begin{pmatrix} u(u+t) & t^2 \\ \text{[unclear]} & u^2+t^2-ut \end{pmatrix} = \begin{pmatrix} \mu u(u+t) & \mu t^2 \\ 1-\mu u(u+t) & -1-\mu u t \end{pmatrix}$$
(tenant compte que $\mu(u^2+t^2) = -1$), a-t-on bien
$$[u, \rho] = -g \qquad i.e.$$
compte de la formule (65) qui explicite $g$ en
termes de c, s, t, u, on trouve que cette formule
prévue équivaut à
$$2 \mu u t = 0 \qquad i.e. \qquad ut = 0$$
décidément, ça déconne ! Mais je vois que si on
remplace dans la définition de $u$ $t$ par $-t$,
i.e. si on définit
$$u = \mu (t \tau_{\infty} + u) \tau_{\infty} = \mu \begin{pmatrix} u & t \\ -t & u \end{pmatrix} \begin{pmatrix} 1 & 0 \\ 0 & -1 \end{pmatrix} = \mu \begin{pmatrix} u & -t \\ -t & -u \end{pmatrix}$$
(au lieu de $u = \mu \begin{pmatrix} u & t \\ t & -u \end{pmatrix}$), alors $u$ est tjrs
une antisimilitude de déterminant $\mu$, et de
plus maintenant
$$[u, \rho] = \begin{pmatrix} \mu u(u-t) & \mu t^2 \\ 1-\mu u(u-t) & -1+\mu u t \end{pmatrix} = -g$$

\newpage

%% ===== Page 158 =====

ce qui se réduit en fait à vérifier
$$ \mu(u-t) = - (1 + \mu t(t+u)) $$
i.e.
$$ \mu[ u - \mu t + t^2 + tu ] = -1 \qquad O.K. $$

En résumé, on trouve

Théorème Soit $\alpha_2 \in \mathbb{Z}/2\mathbb{Z}$ défini par $(-1)^{\alpha_2} = \varepsilon_2$
(donc $\alpha_2 \equiv 0 \pmod 2$ ssi $\mu \equiv 1 \pmod 4$, $\alpha_2 \equiv 1 \pmod 2$ ssi
$\mu \equiv -1 \pmod 4$). Posons
$$ (75) \qquad \underline{u} = \mu(t\sigma_\infty + u)\tau_\infty^{\alpha_2} = \mu \begin{pmatrix} u & t \\ -t & u \end{pmatrix} \begin{pmatrix} 1 & 0 \\ 0 & -1 \end{pmatrix}^{\alpha_2} \qquad \left( = \begin{pmatrix} u & \varepsilon_2 t \\ -t & \varepsilon_2 u \end{pmatrix} \right) $$

[MARGIN: 
$$ \begin{cases} \mu \begin{pmatrix} u & t \\ -t & u \end{pmatrix} & \alpha_2 = 0 \\ \mu \begin{pmatrix} u & -t \\ -t & -u \end{pmatrix} & \alpha_2 = 1 \end{cases} $$
]

On a alors les formules
$$ (76) \qquad \det \underline{u} = \mu $$
$$ (77) \qquad [\underline{u}, \sigma_\infty] = -\varepsilon_2 \cdot id \qquad \text{i.e. } \underline{u}(\sigma_\infty) = \varepsilon_2 \sigma_\infty $$
$$ (78) \qquad [\underline{u}, \rho] = \varepsilon_2 \cdot g \qquad \text{i.e. } \underline{u}(\rho) = (\varepsilon_2 g) \cdot \rho $$

Corollaire Soit $u$ un autom. du groupe à lacets $\hat{\Pi}_{0,3}$, conservant $\mathfrak{S}_3$ (i.e. $\sigma_\infty$ et $\rho$), $\mu$ son multiplicateur, désignons encore par $\underline{u}$ l'extension canonique de $u$ à $\widehat{G}_{0,3}^+ \simeq \operatorname{Sl}(2, \hat{\mathbb{Z}})^\sim$, induisant l'identité sur $\mathfrak{S}_3 \simeq \widehat{G}_{0,3}^+ / \hat{\Pi}_{0,3}$, et considérons le quotient $\operatorname{Sl}(2, \hat{\mathbb{Z}})$ de $\operatorname{Sl}(2, \hat{\mathbb{Z}})^\sim$. Alors l'action de $u$ passe en un autom. dans le quotient $\operatorname{Sl}(2, \hat{\mathbb{Z}})$, et cet autom. est induit par [unclear: un autom.] intérieur de $\operatorname{Gl}(2, \hat{\mathbb{Z}})$, défini par une matrice $\underline{u} \in \operatorname{Gl}(2, \hat{\mathbb{Z}})$ telle que $\det \underline{u} = \mu$.

\newpage

%% ===== Page 159 =====

(qui décrit l'action extérieure de $u$ sur $\operatorname{Sl}(2,\hat{\mathbb{Z}})'$ en termes du multiplicateur $\mu$).

[MARGIN: NB Pour vérifier le corollaire, il faut voir l'action de $u$ sur $\sigma_\infty$ et $\rho$, ce qui a été fait page précédente (oublié le signe $\varepsilon_2$ !)]

En fait, on a beaucoup mieux encore, modulo signes
puisqu'on a pu définir canoniquement
en termes de $u$, la matrice $\underline{u}$ par la
formule (75), en termes $t, u \in \hat{\mathbb{Z}}$
définis ainsi : la donnée de l'autom $u$
définit $f \in \hat{\Pi}_{0,3}$, modulo opération $f \mapsto l_0^n f$
(pour $n \in \hat{\mathbb{Z}}$), et on désigne encore par $f$
l'image de $f$ dans $\operatorname{Sl}(2,\hat{\mathbb{Z}})/\pm 1$, qui est donc
une matrice $\begin{pmatrix} r & s \\ t & u \end{pmatrix}$ définie mod signes,
On note que dans $\operatorname{Sl}(2,\hat{\mathbb{Z}})$ on a

(79) $$l_0^n \begin{pmatrix} r & s \\ t & u \end{pmatrix} = \begin{pmatrix} 1 & -2n \\ 0 & 1 \end{pmatrix} \begin{pmatrix} r & s \\ t & u \end{pmatrix} = \begin{pmatrix} r-2nt & s-2nu \\ t & u \end{pmatrix}$$

($n \in \hat{\mathbb{Z}}$), donc $(t,u)$ ne change pas (mod signes)
quand on remplace $f$ par $l_0^n f$ - donc $(t,u)$
est déterminé mod signes par l'autom. $u$
de $\hat{\Pi}_{0,3}$. L'observation analogue s'applique
d'ailleurs à $\alpha_0 = \begin{pmatrix} a & b \\ c & d \end{pmatrix}$, qui est déterminée
mod signes et mod multiplication à gauche
par un $l_0^n$ - le vecteur $(c,d)$ est
déterminé mod signes, donc l'élément
$c \rho + d$

\newpage

%% ===== Page 160 =====

est déterminé mod signes, et [unclear: rayé]
on doit avoir une relation
(80) $\quad u \overset{?}{=} \mu \alpha(cp+d) = \mu f^{-1} l_0^{-\nu} \alpha_0 (cp+d)$

[MARGIN: NB cette formule ne colle pas avec les formules (cf (69) (73)), dont au moins $\varepsilon_2=1$ - sinon il y a lieu peut-être de corriger par un facteur $\tau_\alpha$ quelque part... cf plus bas...]
sinon il y a lieu peut-être de corriger par un facteur $\tau_\alpha$ quelque part... [unclear: rayé]

$$l_0^{-\nu} \alpha_0(cp+d) = \begin{pmatrix} 1 & 2\nu=\mu-1 \\ 0 & 1 \end{pmatrix} \begin{pmatrix} a & b \\ c & d \end{pmatrix} \begin{pmatrix} d & c \\ -c & c+d \end{pmatrix}$$

$$= \begin{pmatrix} ad-bc=\delta & ac+bc+bd = \delta B \\ 0 & c^2+cd+d^2 = \delta D \end{pmatrix}$$

$$= \delta \begin{pmatrix} 1 & B \\ 0 & D \end{pmatrix}$$

$$l_0^{-\nu} \alpha_0(cp+d) = \delta \begin{pmatrix} 1 & B+2\nu D = C \\ 0 & D \end{pmatrix}$$

i.e.
(81) $\quad l_0^{-\nu} \alpha_0(cp+d) = \delta \begin{pmatrix} 1 & C \\ 0 & D \end{pmatrix}$

d'où
$$\mu \alpha(cp+d) = (\mu f^{-1}) (l_0^{-\nu} \alpha_0(cp+d)) = \mu \delta \begin{pmatrix} u & -s \\ -t & r \end{pmatrix} \begin{pmatrix} 1 & C \\ 0 & D \end{pmatrix}$$

$$= \mu \delta \begin{pmatrix} u & uC-sD \\ -t & -tC+rD \end{pmatrix}$$

$$= \delta \begin{pmatrix} \mu u & -s + \mu u C \\ -\mu t & r - \mu t C \end{pmatrix}$$

or (61)
$$C = \varepsilon_2 T = \varepsilon_2(rt+su) \quad \text{d'où}$$
$$-s+\mu u C = -s + \varepsilon_2 \mu u r t + \varepsilon_2 \mu u^2 s = s( \underbrace{-1+\varepsilon_2 \mu u^2}_{-\varepsilon_2 \mu t^2} ) + \varepsilon_2 \mu u r t = \varepsilon_2 \mu t \underbrace{(-st+ur)}_{+1}$$

\newpage

%% ===== Page 161 =====

$$r - \mu t C = r - \varepsilon_2 \mu r t^2 - \varepsilon_2 \mu s t u = \varepsilon_2 \mu u (ru-st) = \varepsilon_2 \mu u \quad , \quad d\text{'où}$$
[MARGIN: $\overbrace{r(1-\varepsilon_2 \mu t^2)}^{r(\varepsilon_2 \mu s u)}$]

$$ \mu \alpha(cp+d) = \delta \mu \begin{pmatrix} u & -\varepsilon_2 t \\ -t & \varepsilon_2 u \end{pmatrix} = \delta \underline{u} \quad , $$

i.e.
(82) $\qquad \underline{u} = \delta \mu \alpha(cp+d) = (\delta \mu) (f^{-1} l_0^{-2} \alpha_0) (cp+d)$ [MARGIN: où $\varepsilon_2=1$, et $u \equiv \mu (4)$]

Formule qui s'écrit aussi
i.e.
(83) $\qquad f (t_{0\infty} * \underline{u}) = \delta \alpha (cp+d)$
où $\alpha = f^{-1} l_0^{-2} \alpha_0 \quad , \quad \delta = \varepsilon_3 \in \{\pm 1\}, \delta = ad-bc, \delta \equiv \mu \pmod 3$

ou encore
(84) $\qquad f (t_{0 \infty} * \underline{u}) \tilde{t}_{\infty}^{\alpha_2} = \delta l_0^{-2} \alpha_0 (cp+d)$

$$ \begin{pmatrix} r & s \\ t & u \end{pmatrix} \begin{pmatrix} 1 & 0 \\ 0 & \varepsilon_2 \end{pmatrix} \bigg/ \begin{pmatrix} 1 & 2\nu \\ 0 & 1 \end{pmatrix} \begin{pmatrix} a & b \\ c & d \end{pmatrix} \bigg/ \begin{pmatrix} d & c \\ -c & c+d \end{pmatrix} $$

[unclear crossed out equations: $\begin{pmatrix} 1 & rt+su \\ 0 & \varepsilon_2 t^2 + u^2 \end{pmatrix} \dots = \begin{pmatrix} 1 & C \\ 0 & D \end{pmatrix}$ cf (81)]

et sous cette forme la formule [unclear: extraite] en fait
conserve, [unclear: comme] le contenu essentiel de
la fameuse relation des lacets, sans fioritures
matricielles (56) ou (57)

$$ \varepsilon_2 u = D \quad , \quad \varepsilon_2 t = C \quad ! $$

Remarque. En fait, sous une forme apparemment
d'arithmétique, on a fait par la bande des
[unclear: formules] d'algèbre linéaire, valables sur des
[unclear]

\newpage

%% ===== Page 162 =====

anneaux quelconques et formulables sous
[unclear: dans d'autres] conditions beaucoup plus
générales et plus "compréhensives" -
concernant des équations de la forme
$$[\alpha, \rho] = \xi$$
où $\rho \in \operatorname{Gl}(2, A)$ est donné, ainsi que $\xi$, et
$\alpha \in \operatorname{Gl}(2, A)$ est l'inconnue (assujettie
éventuellement à des conditions du
type $\det \alpha = \delta$, $\delta$ fixé), ainsi que
des équations simultanées en $\alpha$ et $f$
(*) $$[\alpha, \rho] = \varepsilon . \mathcal{E}_0(\mu - 1)[f, \sigma] \quad (cf (56))$$
où $\rho, \sigma \in \operatorname{Gl}(2, A)$ sont donnés ainsi qu'un
scalaire $\varepsilon \in A^*$, et $\lambda \mapsto \mathcal{E}_0(\lambda)$ est un "sous-
groupe à un paramètre", p. ex. on peut
supposer $\mathcal{E}_0(\lambda) = \begin{pmatrix} 1 & \lambda \\ 0 & 1 \end{pmatrix} = \varepsilon_0^\lambda$,
$\varepsilon_0 = \begin{pmatrix} 1 & 1 \\ 0 & 1 \end{pmatrix}$ — en supposant au
besoin une relation de la forme $\sigma f = \varepsilon_0$.
Modulo des ajustements convenables, toutes les
formules du présent § devraient pouvoir
se développer dans ce contexte plus géné-
(ainsi que des compléments substantiels)
ral — mais je ne sais si je vais expliciter

\newpage

%% ===== Page 163 =====

380

ici ces gammes délectables d'algèbre linéaire, voire d'interprétations schématiques du genre que le schéma des solutions de (*)
(pour $f, \alpha_0, \mu$ variables) est un
schéma relatif lisse de dim relative 4, qui etc.
Je vais seulement expliciter, dans le
contexte présent adélique, trop particulier manifestement, la

Proposition Soient $a, b, c, d, r, s, t, u \in \hat{Z}$,
pour qu'ils correspondent à des
données $\alpha_0 = \begin{pmatrix} a & b \\ c & d \end{pmatrix}$, $f = \begin{pmatrix} r & s \\ t & u \end{pmatrix}$ satisfai-
sant
$ad-bc = \varepsilon_3 \in \{\pm 1\}$, $ru-st = 1$, $\alpha_0 \rho(\alpha_0)^{-1} = \varepsilon_2 \varepsilon_0^\mu (f \sigma_0 f^{-1})$
avec $\varepsilon_2 \in \{\pm 1\}$ et $\mu \in \hat{Z}$ convenables (cf (48)(49)(53)),
il faut & il suffit qu'on ait
(85) $c^2 + d^2 = \pm (t^2+u^2) \in \hat{Z}^*$

[texte raturé : Alors $(\varepsilon_2, \varepsilon_3, \mu) \in \{\pm 1\} \times \{\pm 1\} \times \hat{Z}^*$ est déterminé modulo le signe (i.e. modulo passage à $(-\varepsilon_2, -\varepsilon_3, -\mu)$) ... alors déterminés de façon unique ... par les relations]
(où $\mu = 2\nu+1$)

(86) $\mu = \varepsilon_2(t^2+u^2) = \varepsilon_3(c^2+d^2)$
(ce qui détermine $\varepsilon_3, \mu$ quand on se donne $\varepsilon_2$)
et on a

(87) $$ \begin{cases} \varepsilon_2 \equiv \mu \ (4) \\ \varepsilon_3 \equiv \mu \ (3) \\ \nu = (\mu-1)/2 \end{cases} $$

\newpage

%% ===== Page 164 =====

Les solutions s'obtiennent alors ainsi : on choisit
$\varepsilon_2, \varepsilon_3, \mu$ comme dessus (deux possibilités), puis
arbitrairement $r, s \in \hat{\mathbb{Z}}$ satisfaisant
$$ur - ts = 1$$
il y en a toujours (car $t^2+u^2 \in \hat{\mathbb{Z}}^*$ - donc
$\forall p \in \mathcal{P}$, $t_p$ et $u_p$ ne sont pas tous deux dans $p\mathbb{Z}_p$)
et si $r_0, s_0$ est une solution, les autres
s'obtiennent [unclear: sous] formes paramétriques :
(87) $(r, s) = (r_0, s_0) + \lambda (t, u) = (r_0 + \lambda t, s_0 + \lambda u)$
Ayant choisi $r, s$, donc $f = \binom{r \quad s}{t \quad u} \in \operatorname{Sl}(2, \hat{\mathbb{Z}})$,
on en déduit $f \sigma_0 \tilde{f} = f \tilde{f} = \binom{R \quad S}{T \quad U}$
($R = r^2+s^2, S = T = rt+su, U = t^2+u^2$), et on détermine
successivement les quantités
(88)
$$D = \frac{1}{\mu} = \varepsilon_2(t^2+u^2) = \varepsilon_2 U$$
$$C = \varepsilon_2 T (= \varepsilon_2 S)$$
$$B = C + D - 1 = \varepsilon_2(S+U) - 1$$
$$A = \frac{1}{D}(1+BC) = \frac{1}{\varepsilon_2 U}(1 + S(S+U) - \varepsilon_2 S)$$
d'où une matrice $\binom{A \quad B}{C \quad D}$ satisfaisant
$C+D-B=1$ et $D \in \hat{\mathbb{Z}}^*$ et satisfaisant la
relation (57) :
$$\binom{A+2\nu C \quad B+2\nu D}{C \quad \quad \quad D} = \varepsilon_2 \binom{R \quad S}{T \quad U}$$
qui se réduit ici à la relation du coin
sup. gauche $A+2\nu C = \varepsilon_2 R$ i.e.

\newpage

%% ===== Page 165 =====

381

il reste à choisir $a,b$ de façon à satisfaire aux
deux relations $\det \begin{pmatrix} a & b \\ c & d \end{pmatrix} = \varepsilon_3$ et $C = \varepsilon_2 T$ (i.e. (61) (62))

(88) $$ \begin{cases} d a - c b = \varepsilon_3 \\ (c+d)a + db = \varepsilon_2 \varepsilon_3 \underbrace{(rt+su)}_{S} + \underbrace{(c^2+cd+d^2)}_{1/\mu \varepsilon_3} \end{cases} $$

qui est un système de deux équations linéaires
en $a,b$, dont le déterminant est comme
par hasard $c^2+cd+d^2 = 1/\mu\varepsilon_3 \in \hat{\mathbb{Z}}^\times$ - ce
qui assure une solution unique !

Il semble plus naturel peut être de
partir du multiplicateur $\mu$, et de chercher
les systèmes $\begin{pmatrix} a & b \\ c & d \end{pmatrix}$ $\begin{pmatrix} r & s \\ t & u \end{pmatrix}$ qui y
correspondent, mais en mettant l'accent
sur $(c,d)$ $(t,u)$ en premier lieu.

A l'aide de $\mu$ on détermine $\varepsilon_2, \varepsilon_3$ par

$$ \mu \equiv \varepsilon_2 (2), \quad \mu \equiv \varepsilon_3 (3), $$

puis ce qui permet de le représenter
$1/\mu$ sous la forme

$$ 1/\mu = \varepsilon_2(t^2+u^2) = \varepsilon_3 (c^2+cd+d^2) $$

(avec $t,u, c,d \in \hat{\mathbb{Z}}$, les systèmes $(t,u)$ comme
$(c,d)$ "dépendent de 1 paramètre", je ne
suis très convaincu qu'on puisse en déduire
[unclear: des formules à y traduire avec des]
paramètres $\sigma, \tau$ explicites... Une
fois $c,d$, $t,u$ choisis, on choisit $r,s$ comme il a

\newpage

%% ===== Page 166 =====

été dit (avec un paramètre (linéairement)
ce qui détermine ensuite $r,s$ de façon
unique, d'une façon qui fait intervenir
[unclear] de façon linéaire-affine.
En fait, pour $c,d,t,u$ fixés, et $\varepsilon_2, \varepsilon_3$,
satisfaisant : $\varepsilon_2(t^2+u^2) = \varepsilon_3(c^2+cd+d^2)$, les trois éq.s
qui relient $a, b, r, s$ sont linéaires :

[MARGIN: on a souligné les "coefficients" inconnus $a,b,r,s$ : ]
(89)
$$da - cb = \varepsilon_3$$
$$ur - ts = 1$$
$$(c+d)\underline{a} + d\underline{b} - \varepsilon_2 \varepsilon_3 (u\underline{s}+t\underline{r}) = c^2+cd+d^2$$

et le déterminant de ce système est "d"
vaut d ?
l'ens. de ses solutions est
un espace affine de dim 1 sur $\hat{\mathbb{Z}}$...

Remarques On n'a pas tenu compte
des conditions de congruences
$$ \begin{pmatrix} a & b \\ c & d \end{pmatrix} \equiv 1 \ (2) \quad \begin{pmatrix} r & s \\ t & u \end{pmatrix} \equiv 1 \ (\text{[unclear]}) $$
qui expriment que ces matrices proviennent
aux signes près "réductions" d'éléments de
$\hat{\Pi}_{0,3}$, $\hat{\Pi}_{0,2}$. Pour les quantités $(c,d)$, par ex.
on trouve les conséquences

$$ d \equiv \pm 1 \ (4) \text{ donc } d^2 \equiv 1 \ (4) \implies c^2+d^2+cd \equiv 1+cd \ (4) \equiv 1+c \ (4)$$
$$ c \equiv 0,2 \ (4) \text{ donc } c^2 \equiv 0 \ (4) $$
donc $$ \mu = \varepsilon_3(c^2+cd+d^2) \equiv \varepsilon_3(1+c) \ (4) $$

(NB $d$ est inv. mod 4 puisque $d \equiv 1 (2)$ )

\newpage

%% ===== Page 167 =====

382

On voit donc que la connaissance de $c$ [MARGIN: mod 4]
et celle des signes $ad-bc = \varepsilon_3$, permet de
trouver le signe $\varepsilon_2$ par 

(90) $\quad \varepsilon_2 \equiv \varepsilon_3(1+c) \ (4) \quad \text{i.e.} \quad \varepsilon_2\varepsilon_3 \equiv 1+c \ (4)$

Si on se propose, pour $c,d,t,u$ fixés,
satisfaisant
$\quad c \equiv t \equiv 0 \ (2) \quad d \equiv u \equiv 1 \ (2)$
(91) $\left\{ \begin{array}{l} (c^2+cd+d^2) = \varepsilon(t^2+tu+u^2) \in \hat{\mathbb{Z}}^* \\ \text{où } \varepsilon \in \{\pm 1\} \text{ est le signe défini par} \end{array} \right.$
$\quad \varepsilon \equiv 1+c \ (4)$
($\varepsilon=1$ si $c \equiv 0 \ (4)$, $\varepsilon=-1$ si $c \equiv 2 \ (4)$), de trouver
les $a,b,r,s$ correspondants, satisfaisant
les conditions (89), et de plus les
conditions de congruence
(92) $\quad b \equiv s \equiv 0 \ (2) \quad a \equiv r \equiv 1 \ (2)$,
on est amené à relever une solution déjà
donnée mod 2 de ce système d'équations
en une solution dans $\hat{\mathbb{Z}}$. En termes des torseurs
$T$ des solutions, cela signifie qu'on a
donné un élément de $T/2T$ - trouver un point de
ce torseur - l'un de ses relèvements s'identifie
à un torseur sous $2\hat{\mathbb{Z}} \simeq \hat{\mathbb{Z}}$... Aucun
mystère !

\newpage

%% ===== Page 168 =====

Remarque Le groupe additif $\hat{Z}$ opère sur
l'ens des solutions $(\alpha_0, f) = \left( \begin{pmatrix} a & b \\ c & d \end{pmatrix}, \begin{pmatrix} r & s \\ t & u \end{pmatrix} \right)$
des équations envisagées (48)(49)(56)
(93) $\det \alpha_0 \in \{\pm 1\}$, $\det f = 1$, $\alpha_0 f(\alpha_0)^{-1} = \varepsilon l_0^b (f \circ \alpha_0(f)^{-1})$, $\varepsilon \in \{\pm 1\}$,

par
(94) $(\alpha_0, f) \mapsto (\alpha'_0, f') = (l_0^n \alpha_0, l_0^n f)$
$= \left( \begin{pmatrix} a - 2nc & b - 2nd \\ c & d \end{pmatrix} , \begin{pmatrix} r - 2nt & s - 2nu \\ t & u \end{pmatrix} \right)$
($n \in \hat{Z}$).

Ce facteur 2 près dans $2n$, c'est justement l'
indétermination dans l'ensemble des solutions de (93)
lorsque $c, d, t, u$ fixés (cf (87) et proposition)
[unclear] On peut donc 
se demander ce qui est déterminé intrinsèquement (aux
termes de l'action de $\widehat{\Pi}_{0,3}$ comme
[unclear] , car on n'a plus un moyen pour

a) [unclear]
(95) b) les quantités [unclear]

les $c, d, t, u \in \hat{Z}$
satisfont la surprenante relation (86)
(96) $\varepsilon_2(t^2+u^2) = \varepsilon_2(c^2+cd+d^2) = 1/p$
et les relations de congruence
(97) $c \equiv t \equiv 0 \quad (2) \qquad d \equiv u \equiv 1 \quad (2)$

Il est vrai que l'on a vu qu'il y a un
moyen naturel (dans $\widehat{\Pi}_{0,3}$) de normaliser le choix
de $f$, de façon à lever l'ambiguïté des
$l_0^n f$ (cf §91 où c'était effleuré - j'y reviendrai
sans doute par la suite...) mais il me semble [unclear]

\newpage

%% ===== Page 169 =====

389

pour le moment d'ignorer cette possibilité
de normalisation - ce qui met mieux en
évidence, en effet, le caractère crucial,
dans tous ces calculs $l$-adiques, du rôle
joué par les quatre quantités $c,d,t,u \in \hat{Z}$.

\newpage

%% ===== Page 170 =====

384

## 45 L'homomorphisme $\tilde{M}_{0,3} \to \operatorname{Gl}(2, \hat{\mathbb{Z}})'$, construction de l'extension $\tilde{N}_{1,1}$ de $\hat{\Gamma}$ par $\widehat{Sl}(2,\mathbb{Z})^\wedge$.

[unclear: crossed out paragraph starting with "On désigne par $\tilde{M}_{0,3}$..."]

[MARGIN: On désignera (4) par $\tilde{M}_{0,3}$ le groupe des classes d'automorphismes de $\hat{\mathfrak{G}}_{0,3}^+$ (ou ce qui revient au même, des aut. de $\hat{\Pi}_{0,3}$ ext. à bouts) qui opèrent trivialement sur les classes de conjugaison des inerties (i.e. qui opèrent par l'identité sur $\hat{\mathfrak{S}}_3$)]

(1) $$1 \to \hat{\Pi}_{0,3} \to \tilde{M}_{0,3} \to \tilde{N}_{0,3} \to 1,$$

où $\tilde{N}_{0,3}$ désigne le groupe des automorphismes extérieurs de $\hat{\mathfrak{G}}_{0,3}^+$ qui induisent l'identité des [unclear].

On notera que
$$\Gamma = \Gamma_Q \hookrightarrow \hat{\Gamma}$$
a pour l'image inverse de $\Gamma_Q$ dans l'ext (1) s'identifie avec $M_{0,3}$ - 

mais pour le moment je n'ai pas de caractérisation bien plausible à proposer de $\Gamma_Q$ comme sous-groupe de $\hat{\Gamma}$. On a un scindage canonique de l'extension (1), via le sous-groupe $\Gamma_Q$ de $\tilde{M}_{0,3}$ des éléments de $\tilde{M}_{0,3}$ qui commutent à $\tau_{\infty}$.

\newpage

%% ===== Page 171 =====

de sorte que $\tilde{M}_{0,3}$ se présente comme

un produit 1/2 direct

(3) $$ \tilde{M}_{0,3} \simeq \hat{\mathfrak{S}}_{0,3}^+ \cdot \tilde{\Gamma}_{\sigma} \quad . $$
[MARGIN: sous $\hat{\mathfrak{S}}_{0,3}^+$, $S\ell(2,\hat{\Z})^\wedge$ est rayé]

Considérons l'homomorphisme canonique

(4) $$ \hat{\mathfrak{S}}_{0,3}^+ \simeq S\ell(2,\hat{\Z})^\wedge \xrightarrow{\varphi'_0} S\ell(2,\hat{\Z})' \subset G\ell(2,\hat{\Z})' $$

et l'homomorphisme

(5) $$ \tilde{\Gamma}_{\sigma} \xrightarrow{\varphi'_0} G\ell(2,\hat{\Z})' $$
$$ u \longmapsto \text{[unclear]} = \mu \begin{pmatrix} v & \varepsilon t \\ -t & \varepsilon v \end{pmatrix} , $$

où $t, v \in \hat{\Z}$, $\mu \in \hat{\Z}^*$ ($\varepsilon \in \{\pm 1\}$), $\alpha = \alpha_2 \in \hat{\Z}/2\hat{\Z}$

sont les quantités dépendantes de $u \in \tilde{\Gamma}_{\sigma}$

envisagées au § précédent, avec

$$ \varepsilon = (-1)^\alpha . $$

S'il y a lieu, on écrira $t(u), v(u), \varepsilon(u), \mu(u)$

pour indiquer la dépendance par rapport à $u$.

Considérons l'application (continue)

(6) $$ \varphi' \colon \tilde{M}_{0,3} \longrightarrow G\ell(2,\hat{\Z})' $$

définie par

(7) $$ \varphi'(a.u) = \varphi'_0(a) \varphi'_0(u) \quad \begin{cases} a \in \hat{\mathfrak{S}}_{0,3}^+ \\ u \in \tilde{\Gamma}_{\sigma} \end{cases} $$

\newpage

%% ===== Page 172 =====

Théorème. L'application $\varphi'$ est un homomorphis-
me de groupes.

Cela signifie, compte tenu que $\varphi_0'$ est
déjà un hom. de groupes, le théorème
signifie que

a) $\varphi_\sigma'$ est un hom. de groupes, i.e.
(8) $$\varphi_\sigma'(\underline{u}' \underline{u}) = \varphi_\sigma'(\underline{u}') \varphi_\sigma'(\underline{u}) \quad (\underline{u}, \underline{u}' \in \tilde{\Gamma}_{\sigma})$$
b) on a , pour $a \in \hat{\mathfrak{G}}_{0,3}^+$, $\underline{u} \in \tilde{\Gamma}_{\sigma}$

(9) $$\varphi_\sigma'(\underbrace{a \underline{u} a^{-1}}_{\underline{u}(a)}) = \varphi_0'(a) \varphi_\sigma'(\underline{u}) \varphi_0'(a)^{-1}$$

Mais la formule (9) est en fait
déjà établie puisque le choix
par construction même de la matrice
$\varphi_\sigma'(\underline{u}) = \mu(t \sigma_\infty + u)\tau_\infty^\varepsilon$ aux paragraphes
précédents était fait de façon à
ce que cette formule soit vérifiée
(cf (67) (68), et le cor. au th. (77)... ). Pour
vérifier (8), nous allons d'abord expliciter
les quantités $g, h, f$ pour un
produit de deux opérateurs $\underline{u}', \underline{u} \in \tilde{M}_{0,3}$
en termes de leurs valeurs pour $\underline{u}', \underline{u}$
eux-mêmes. Un calcul immédiat
donne, en ce sens, des "relations de cocycle"

(10) 
$$g(\underline{u}' \underline{u}) = \underline{u}'(g(\underline{u})) . g(\underline{u}')$$
$$h(\underline{u}' \underline{u}) = \underline{u}'(h(\underline{u})) . h(\underline{u}')$$
$$f(\underline{u}' \underline{u}) = f(\underline{u}') . \underline{u}'(f(\underline{u}))$$

\newpage

%% ===== Page 173 =====

Remarques 1) On aura d'ailleurs de même, avec les notations du $\S$ précédent

(11)
$$ \begin{cases} \alpha(\underline{u}' \underline{u}) = \underline{u}'(\alpha(\underline{u})) . \alpha(\underline{u}') \\ \beta(\underline{u}'\underline{u}) = \underline{u}' (\beta(\underline{u})) . \beta(\underline{u}') \end{cases} $$

2) La formule (10) a un air un peu plus sympa pour $f$, que pour $g$ et $h$, (on aurait envie de changer encore de notation, et de remplacer également $g$ et $h$ par $g^{-1}, h^{-1}$ (comme précédemment on a [unclear] devenu $f^{-1}$...).

Désignons encore par les mêmes lettres $g(\underline{u}), h(\underline{u})$ etc les images de ces quantités dans $\operatorname{Sl}(2, \hat{\mathbb{Z}})'$, les formules analogues des (11) dans $\operatorname{Sl}(2, \hat{\mathbb{Z}})'$ (ou $(\operatorname{Gl}(2, \hat{\mathbb{Z}}))'$ pour les $\alpha$, $\beta$ de (11) ) qui s'écrivent

(12)
$$ \begin{cases} g(\underline{u}'\underline{u}) = int(\varphi_0'(\underline{u}')) (g(\underline{u})) . g(\underline{u}') \\ h(\underline{u}'\underline{u}) = int(\varphi_0'(\underline{u}')) (h(\underline{u})) . h(\underline{u}') \\ f(\underline{u}'\underline{u}) = f(\underline{u}') . int(\varphi_0'(\underline{u}')) (f(\underline{u})) \\ \alpha(\underline{u}'\underline{u}) = int(\varphi_0'(\underline{u}')) (\alpha(\underline{u})) . \alpha(\underline{u}') \\ \beta(\underline{u}'\underline{u}) = int(\varphi_0'(\underline{u}')) (\beta(\underline{u})) . \beta(\underline{u}') \end{cases} $$

dans $\operatorname{Gl}(2,\hat{\mathbb{Z}})'$.

Ces formules dans $\operatorname{Gl}(2,\hat{\mathbb{Z}})'$ ont besoin [unclear] une vérification entière [unclear] expliciter des preuves de mon calcul hâtif [unclear]

\newpage

%% ===== Page 174 =====

385

au cas des formules initiales dans $\hat{\Pi}_{0,3}$ ou $\tilde{\Pi}_{0,3}$,
où des quantités comme $\underline{w}(g(\underline{u}))$ ne sont pas
explicitables [unclear: par des formules alg.] simples
en fonction des données $(f, h, j \dots)$ relatives
à $g(\underline{u})$, qui déterminent $\underline{u}$ et $\underline{w}'$.
Ces formules se complètent, pour mémoire,
par

$$ (13) \quad \mu(\underline{u}' \underline{u}) = \mu(\underline{u}') \mu(\underline{u}) $$

impliquant notamment

$$ (14) \quad \varepsilon_2(\underline{u}' \underline{u}) = \varepsilon_2(\underline{u}') \varepsilon_2(\underline{u}) $$

Pour faire le calcul de $\varphi_{\infty}(\underline{w}, \underline{u})$ et
vérifier qu'il s'agit de $\varphi_{\infty}(\underline{w}) \varphi_{\infty}(\underline{u})$, on
introduira les coefficients $r, s, t, u$
de $f = f(\underline{u})$, $r', s', t', u'$ de $f' = f(\underline{u}')$, et
$r'', s'', t'', u''$ de $f'' = f(\underline{u}' \underline{u})$, de même $\mu, \mu', \mu''$, $\varepsilon, \varepsilon', \varepsilon''$
d'où

$$ (15) \quad \mu'' = \mu' \mu \, , \, \varepsilon'' = \varepsilon \varepsilon' $$

et la formule en $f$ dans (12) s'écrit alors

$$ (16) \quad f'' = f' \, \underline{w}' \, f \, {\underline{w}'}^{-1} $$

écrit, comme au § précédent,
$\underline{w}'$ au lieu de $\varphi_{\infty}(\underline{w}')$, En termes
explicites

$$ \begin{pmatrix} r'' & s'' \\ t'' & u'' \end{pmatrix} = \begin{pmatrix} r' & s' \\ t' & u' \end{pmatrix} \begin{pmatrix} \mu' & \varepsilon' \\ -\varepsilon' & \mu' \end{pmatrix} \begin{pmatrix} r & s \\ t & u \end{pmatrix} \begin{pmatrix} \mu' & -\varepsilon' \\ \varepsilon' & \mu' \end{pmatrix} $$

[MARGIN: NB $\det \underline{w}' = 1$]

\newpage

%% ===== Page 175 =====

Le produit des trois derniers termes du
second membre est [unclear] :

$$ \underline{u}' f \underline{u}'^{-1} = \mu' \begin{pmatrix} \varepsilon' r u'^2 + s t' u' + t t' u' + \varepsilon' u t'^2 & -\varepsilon' r t' u' + s u'^2 - t t'^2 + \varepsilon' u u' t' \\ -\varepsilon' s u' t' + s t'^2 + t u'^2 + \varepsilon' u u' t' & - \varepsilon' r t'^2 - s u' t' - t u' t' + \varepsilon' u u'^2 \end{pmatrix} $$

mais le produit par $f'$ se simplifie
miraculeusement, grâce à la relation
$$ t'^2 + u'^2 = \varepsilon' / \mu' $$
et on trouve

(17) $$ f'' = \begin{pmatrix} r'' & s'' \\ t'' & u'' \end{pmatrix} = \begin{pmatrix} \text{bordel} & \text{bordel} \\ \varepsilon' t u' + u t' & -t t' + u u' \end{pmatrix} $$

d'où

(18) $$ t'' = \frac{1}{\mu'} (\varepsilon' t u' + u t') \qquad u'' = \frac{1}{\mu'} (-t t' + u u') $$

et par suite (tenant compte de $\mu'' = \mu \mu'$, $\varepsilon'' = \varepsilon \varepsilon'$)

$$ \underline{u}'' = \mu' \mu \begin{pmatrix} -t t' + u u' & \varepsilon t u' + \varepsilon \varepsilon' u t' \\ -\varepsilon' t u' - u t' & - \varepsilon \varepsilon' t t' + \varepsilon \varepsilon' u u' \end{pmatrix} $$

et la formule à vérifier
$$ \underline{u}'' = \underline{u}' \underline{u} $$
se réduit à
$$ \mu' \mu \begin{pmatrix} -t t' + u u' & \varepsilon t u' + \varepsilon \varepsilon' u t' \\ -\varepsilon' t u' - u t' & - \varepsilon \varepsilon' t t' + \varepsilon \varepsilon' u u' \end{pmatrix} = \mu' \begin{pmatrix} u' & \varepsilon' t' \\ -t' & \varepsilon' u' \end{pmatrix} \mu \begin{pmatrix} u & \varepsilon t \\ -t & \varepsilon u \end{pmatrix} $$

qui est bien vraie, OK !

Ici encore, on sent que ces calculs
devraient s'insérer dans un tapis d'algèbre
linéaire qui [unclear] des [unclear]
in extenso. Il faudrait paraphraser la
description de [unclear: $\tilde{M}_{0,3}$] décrite par ces systèmes [unclear]
(ou plutôt [unclear]) des réductions de [unclear]

\newpage

%% ===== Page 176 =====

dans $\operatorname{Gl}(2, \hat{\mathbb{Z}})/_{\pm 1}$, puis on remplacera $\hat{\mathbb{Z}}$ par
un anneau quelconque, et bien sûr
$\operatorname{Gl}(2, \hat{\mathbb{Z}})'$ par $\operatorname{Gl}(2, \hat{\mathbb{Z}})$ lui-même. Peut-
être y reviendrai-je finalement ....

Ayant ainsi construit
$$ \varphi : \widetilde{M}_{0,3} \longrightarrow \operatorname{Gl}(2, \hat{\mathbb{Z}})' = \operatorname{Gl}(2, \hat{\mathbb{Z}})/\pm 1 $$
on déduit, par image inverse de l'extension
centrale $\operatorname{Gl}(2, \hat{\mathbb{Z}})$ de $\operatorname{Gl}(2, \hat{\mathbb{Z}})'$ par $\{\pm 1\}$, une
extension centrale de $\widetilde{M}_{0,3}$ par $\{\pm 1\}$,
que je note $\widetilde{N}_{1,1}$, et on a
un hom d'extensions

(19)
\begin{tikzcd}
1 \ar[r] & \{\pm 1\} \ar[r] \ar[d, equal] & \widetilde{N}_{1,1} \ar[r] \ar[d, "\varphi"] & \widetilde{M}_{0,3} \ar[r] \ar[d, "\varphi'"] & 1 \\
1 \ar[r] & \{\pm 1\} \ar[r] & \operatorname{Gl}(2, \hat{\mathbb{Z}}) \ar[r] & \operatorname{Gl}(2, \hat{\mathbb{Z}})' \ar[r] & 1
\end{tikzcd}

[Paragraphe raturé au bas de la page]

\newpage

%% ===== Page 177 =====

Il faudrait vérifier que l'hom $\varphi'$
qu'on vient de construire ci-dessus sur
$\mathcal{M}_{0,3} \subset \tilde{\mathcal{M}}_{0,3}$ ... est l'hom. "modulaire"
bien connu du § 39 (?) - il suffit
de le vérifier sur le sous-groupe
$\Gamma_{Q,\sigma}$ de $\tilde{\Gamma}_Q$ correspondant à $\Gamma_Q$ -
ce ne doit pas être bien difficile, et
je vais l'admettre ici. Ceci étant
admis, on trouve que l'image
inverse de $\mathcal{M}_{0,3}$ dans $\tilde{\mathcal{N}}_{1,1}$ n'est
autre (à isom. can. près ) que
$\mathcal{N}_{1,1}$ lui-m. - ou tout au moins l'
image inverse du [unclear: crossed out] n'est
autre que $\widehat{C}_{1,1}^+ \simeq \operatorname{Sl}(2,\mathbb{Z})^\wedge$. On aura alors
(avec cette identification)

(20)
\begin{tikzcd}
1 \ar[r] & \widehat{C}_{1,1}^+ \ar[r] \ar[d] & \tilde{\mathcal{N}}_{1,1} \ar[r] \ar[d, "\varphi"] & \tilde{\Gamma} \ar[r] \ar[d, "det"] & 1 \\
1 \ar[r] & \operatorname{Sl}(2,\hat{\mathbb{Z}}) \ar[r] & \operatorname{Gl}(2,\hat{\mathbb{Z}}) \ar[r] & \hat{\mathbb{Z}}^\times \ar[r] & 1
\end{tikzcd}

\newpage

%% ===== Page 178 =====

donc on a reconstruit en termes d'autom[orphisme]s
à l'aide de $\hat{\Pi}_{0,3}$ le diagramme (1) du
§ 39, - plutôt un diagramme un
peu plus général, dans lequel se plonge
celui du § 39 - il resterait encore à
caractériser $\Gamma_Q$ comme sous-groupe de $\tilde{\Gamma}$.

Le diagramme déduit directement du
th. précédent

(21)
\begin{tikzcd}
1 \ar[r] & \hat{C}_{0,3}^+ \ar[r] \ar[d] & \tilde{\mathcal{M}}_{0,3} \ar[r] \ar[d, "\phi'"] & \tilde{\Gamma} \ar[r] \ar[d, "\bar{\chi}"] & 1 \\
1 \ar[r] & \operatorname{Sl}(2,\hat{\mathbb{Z}})' \ar[r] & \operatorname{Gl}(2,\hat{\mathbb{Z}})' \ar[r, "\det"'] & \hat{\mathbb{Z}}^\times \ar[r] & 1
\end{tikzcd}

peut être considéré comme étant déduit
du précédent par réduction mod les
[unclear: sous-groupes centraux] $\{\pm 1\}$ dans les deux lignes.

Que devient dans (20) le scindage
(13) de $\tilde{\mathcal{M}}_{0,3}$ ? Y a-t-il une façon
privilégiée, pour les éléments de
$\tilde{\mathcal{N}}_{1,1}$ - i.e. un couple d'un
$u = (g,h,f) \in \tilde{\mathcal{M}}_{0,3}$, et d'un $u_0 \in \operatorname{Gl}(2,\hat{\mathbb{Z}})$

\newpage

%% ===== Page 179 =====

qui relève l'élt $v = \varphi_e'(u)$ de $\operatorname{Gl}(2,\hat{Z})'$,
de le "normaliser" mod composition
avec un élém. de $\operatorname{Sl}(2, \hat{Z})^\sim$, de
façon à ce qu'il ait une (des...) [unclear]
propriétés de commutation avec $\sigma_\infty$, du
type $u(\sigma_\infty) = \sigma_\infty$, ou $u(\sigma_\infty) = \pm \sigma_\infty$ ?

Une autre façon de poser la question,
en termes du st-groupe $\tilde{\Gamma}_Q \subset M_{0,3}^\sim$,
est la suivante : y-a-t-il un relève-
ment canonique $\tilde{\Gamma}_Q \to \mathcal{N}_{1,1}^\sim$,
associant donc à tout $u = (g, f) \in \tilde{\Gamma}_Q$
(d'où un $\underline{u} = \mu(t \sigma_\infty + n) \tau_\infty^\alpha \in \operatorname{Gl}(2,\hat{Z}) \widehat{\otimes} \hat{Z}$
calculé ici, [unclear: modulo] un $\pm$ [unclear])
un représentant $u_0 \in \operatorname{Gl}(2,\hat{Z})$ — i.e.
une façon de lever cette ambiguïté
de signe ? En fait, si on remplace
$\tilde{\Gamma}_Q$ par $\Gamma_{Q,\infty}$, le yoga
standard suggère que les relèvements en

\newpage

%% ===== Page 180 =====

389

question doivent correspondre aux courbes
elliptiques sur $\mathbb{Q}$, munies d'une rigidifica-
tion de [unclear: niveau] d'échelon 2, correspondant
à l'isomorphisme $\mathcal{M}_{1,1 \mathbb{Q}} \simeq \mathcal{M}_{0,3 \mathbb{Q}}$
au point $1/2$ de $M_{0,3}(\mathbb{Q})$. L'indétermina-
tion pour le choix d'une telle courbe
elliptique est justement dans $H^1(\mathbb{Q}, \pm 1)$
(c'est un [unclear: don] d'ext. quadratique de $\mathbb{Q}$)
qui s'identifie : $\operatorname{Hom}(\Gamma_\mathbb{Q}, \pm 1)$ -
ce qui donne bien l'indétermina-
tion pour un relèvement. Ceci prouve
essentiellement que l'idée suggérée par
le yoga général est bel et bien
correcte - et cela montre en [unclear: même temps]
qu'il n'y a pas (même sur $\tilde{\Gamma}_{\mathbb{Q}, \dots}$) de
choix vraiment privilégié. Qu'il y ait
de tels relèvements (sur $\tilde{\Gamma}_{\mathbb{Q}, \dots}$) provient du fait
qu'il y a bel et bien des courbes,

\newpage

%% ===== Page 181 =====

elliptiques répondant à la question,
la plus simple étant donnée par
l'équation

(22) $y^2 - x(x-1)(x-1/2) = 0$

- et voilà le choix fait ! Il resterait à
essayer de l'expliciter dans un
formulaire calculatoire, si ce n'est
pas sans espoir. Il faudrait que
j'examine cette question de
plus près, à l'occasion.

\newpage

%% ===== Page 182 =====

45 390

Notations définition (?) - nouvelle récapitulation

Pour les autom. à lacets de $\Pi_{0,3}^\wedge$ et de $\widehat{G}_{0,3}^+ \simeq \operatorname{Sl}(2, \mathbb{Z})^\wedge$.

I Les groupes $M_{0,3}^\sim$ et $N_{0,3}^\sim$ :
On pose

(1)
$$M_{0,3}^\sim = \text{groupe des autom de } \widehat{G}_{0,3}^+ \simeq \operatorname{Sl}(2, \mathbb{Z})^\wedge$$
$$\text{qui invariant } \Pi_{0,3}^\wedge, \text{ et qui respectent la}$$
$$\text{"structure à lacets" i.e. satisfont } u(\xi_i) = int(f_i)(\xi_i^{k_i}),$$
$$\text{avec } f_i \in \widehat{G}_{0,3}^+ , i \in \{0,1,\infty\}, k_i \in \hat{\mathbb{Z}}^* \text{ (où } f_i \in \Pi_{0,3}^\wedge,$$
$$\text{et } f_i \text{ est déterminé alors dans } Z_i \backslash \Pi_{0,3}^\wedge )$$

[MARGIN: NB $M_{0,3}^\sim$ est aussi le normalisateur de $\widehat{G}_{0,3}^+$ dans $\widehat{G}_{0,3}$, et $N_{0,3}^\sim$ le normalisat. dans $\widehat{G}_{0,3}$ de $\Pi_{0,3}^\wedge$ ]

resp.
$$\overset{\sim}{\longrightarrow} \text{groupe des autom ext. à lacets de } \Pi_{0,3}^\wedge$$
$$\text{qui commutent extérieurement à}$$
$$\text{l'action extérieure naturelle de } \mathfrak{S}_3.$$

On a une structure d'extension

(2)
$$1 \longrightarrow \Pi_{0,3}^\wedge \longrightarrow M_{0,3}^\sim \longrightarrow N_{0,3}^\sim \longrightarrow 1$$

(3)
$$N_{0,3}^\sim \overset{dfn}{=} M_{0,3}^\sim / \Pi_{0,3}^\wedge$$
$$\simeq \text{groupe des autom ext à lacets}$$
$$\text{de } \Pi_{0,3}^\wedge \text{, commutant à } \mathfrak{S}_3$$

L'hom. naturel $\mathfrak{S}_3 \longrightarrow \operatorname{Aut}(\mathfrak{S}_3)$ est un iso

(4)
$$\mathfrak{S}_3 \overset{x \mapsto int(x)}{\longrightarrow} \operatorname{Aut}(\mathfrak{S}_3)$$

et l'hom. naturel
$$M_{0,3}^\sim \longrightarrow \operatorname{Aut}(\mathfrak{S}_3 = \widehat{G}_{0,3}^+ / \Pi_{0,3}^\wedge)$$

définit alors un homomorphisme

(5)
$$M_{0,3}^\sim \longrightarrow \mathfrak{S}_3$$

qui peut aussi s'interpréter comme la composition

(6)
$$M_{0,3}^\sim \hookrightarrow \underset{\text{tous les autom.}}{\underset{\text{à lacets de } \Pi_{0,3}^\wedge}{\widehat{G}_{0,3}}} \xrightarrow{\text{action sur les choix de lacets}} \mathfrak{S}_3$$

\newpage

%% ===== Page 183 =====

et qui est trivial sur $\Pi_{0,3}^\wedge$, donc se factorise en

$$(7) \quad \mathcal{N}_{0,3}^\sim \longrightarrow \mathfrak{S}_3$$

Les noyaux de (5), (7) sont notés $\mathcal{M}_{0,3}^{\sim !}$, $\mathcal{N}_{0,3}^{\sim !}$ respectivement, de sorte qu'on a un diagramme de suites exactes :

\begin{tikzcd}
& & \Pi_{0,3}^\wedge \ar[d] & & & \\
(8) & 1 \ar[r] & \mathcal{M}_{0,3}^{\sim !} \ar[r] \ar[d] & \mathcal{M}_{0,3}^\sim \ar[r] \ar[d] & \mathfrak{S}_3 \ar[r] \ar[d, equal] & 1 \\
(9) & 1 \ar[r] & \mathcal{N}_{0,3}^{\sim !} \ar[r] \ar[d] & \mathcal{N}_{0,3}^\sim \ar[r] & \mathfrak{S}_3 \ar[r] & 1 \\
& & 1 & & &
\end{tikzcd}

(9) pourra à son tour être considéré comme quotient de (8) par $\Pi_{0,3}^\wedge$ \quad Notons aussi la suite exacte correspondante

$$(9 \text{ bis}) \quad 1 \longrightarrow \Pi_{0,3}^\wedge \longrightarrow \mathcal{M}_{0,3}^\sim \longrightarrow \mathcal{N}_{0,3}^\sim \longrightarrow 1$$

On a un hom. canonique

$$(10) \quad \mathfrak{S}_3 \simeq \widehat{G_{0,3}^+} / \Pi_{0,3}^\wedge \longrightarrow \mathcal{N}_{0,3}^\sim$$

dont le composé avec (7) est l'identité, d'où une scission

$$(11) \quad \mathcal{N}_{0,3}^\sim \simeq \mathfrak{S}_3 \times \mathcal{N}_{0,3}^{\sim !}$$

D'autre part, on a un isomorphisme

\begin{tikzcd}
(12) & \Gamma_{0,3} \ar[r, "\sim"] \ar[d, "\simeq"'] & \mathcal{N}_{0,3}^\sim \\
& \mathfrak{S}_3 \times \mathbb{Z}/2\mathbb{Z} & 
\end{tikzcd}

qui sur le facteur $\mathfrak{S}_3$ ( $= \Gamma_{0,3}^+$ ) coïncide avec (10), et qui en outre applique l'élément non trivial $\tau$ de $\mathbb{Z}/2\mathbb{Z}$ sur l'élément, également noté $\tau$, du deuxième facteur $\mathcal{N}_{0,3}^{\sim !}$ de (11).

\newpage

%% ===== Page 184 =====

L'homomorphisme composé
(13) $$ \mathcal{M}_{0,3}^\sim \to \hat{\mathfrak{S}}_{0,3} \xrightarrow{\text{multiplicateur } \chi} \hat{\mathbb{Z}}^* $$
donne un hom.
(14) $$ \mathcal{M}_{0,3}^\sim \xrightarrow{\chi} \hat{\mathbb{Z}}^* $$
surjectif, se factorise par $\mathcal{N}_{0,3}^\sim$ et même $\mathcal{N}_{0,3}^\sim / \mathfrak{S}_3$
(15) 
\begin{tikzcd}
\mathcal{N}_{0,3}^\sim \ar[rr, "\chi"] \ar[dr] & & \hat{\mathbb{Z}}^* \\
& \mathcal{N}_{0,3}^\sim / \mathfrak{S}_3 \ar[ur] & \\
& = \mathcal{N}_{0,3}^1 &
\end{tikzcd}

et donne lieu à des structures d'extensions :
(16)
\begin{tikzcd}
1 \ar[r] & \mathcal{M}_{0,3}^{\sim +} \ar[r] \ar[d] & \mathcal{M}_{0,3}^\sim \ar[r, "\chi"] \ar[d] & \hat{\mathbb{Z}}^* \ar[r] \ar[d, equal] & 1 \\
1 \ar[r] & \mathcal{N}_{0,3}^{\sim +} \ar[r] & \mathcal{N}_{0,3}^\sim \ar[r, "\chi"] & \hat{\mathbb{Z}}^* \ar[r] & 1
\end{tikzcd}

où on a posé :
(17) $$ \begin{cases} \mathcal{M}_{0,3}^{\sim +} = \operatorname{Ker}(\mathcal{M}_{0,3}^\sim \to \hat{\mathbb{Z}}^*) \\ \mathcal{N}_{0,3}^{\sim +} = \operatorname{Ker}(\mathcal{N}_{0,3}^\sim \to \hat{\mathbb{Z}}^*) \end{cases} $$

- la deuxième ligne de (16) s'interprète comme quotient de la première ligne par $\hat{\Pi}_{0,3}$ .

On a un plongement canonique
(18) $$ \Gamma_{\mathbb{Q}} \stackrel{df}{=} \operatorname{Gal}(\bar{\mathbb{Q}}/\mathbb{Q}) \hookrightarrow \mathcal{N}_{0,3}^\sim $$
[MARGIN: clôture alg. de $\mathbb{Q}$ dans $\mathbb{C}$]
unique à [unclear: un autom. int. de conjugaison près]

\newpage

%% ===== Page 185 =====

complexe $\tau$ ou l'élément déjà noté $(\tau)$
(cf (11)). L'image inverse du sous-
groupe $\Gamma_{\mathbb{Q}}$ par $\mathcal{M}_{0,3}^{\sim} \to \mathcal{N}_{0,3}^{\sim} \simeq \mathcal{N}_{0,3}^{1 \sim} \rtimes \hat{\mathfrak{S}}_3$
i.e l'image inverse de $\hat{\mathfrak{S}}_3 \times \Gamma_{\mathbb{Q}} \subset$
(19) $$ \mathcal{N}_{0,3} \overset{dfn}{=} \hat{\mathfrak{S}}_3 \times \Gamma_{\mathbb{Q}} \subset \mathcal{N}_{0,3}^{\sim} = \hat{\mathfrak{S}}_3 \times \mathcal{N}_{0,3}^{1 \sim} $$
est [unclear] notée $\mathcal{M}_{0,3}$ par ailleurs :

[MARGIN: On a une [unclear] de $\mathcal{M}_{0,3}^{\sim}$ de (7)]
(20)
\begin{tikzcd}
1 \ar[r] & \Pi_{0,3}^{\sim} \ar[r] \ar[d] & \mathcal{M}_{0,3} \ar[r] \ar[d] & \hat{\mathfrak{S}}_3 \times \Gamma_{\mathbb{Q}} \ar[r] \ar[d, "(\alpha, \beta) \mapsto \beta"] & 1 \\
1 \ar[r] & \hat{\mathfrak{S}}_{0,3}^+ \ar[r] & \mathcal{M}_{0,3} \ar[r] & \Gamma_{\mathbb{Q}} \ar[r] & 1
\end{tikzcd}

Bien entendu, c'est par $\mathcal{M}_{0,3}$ et $\mathcal{N}_{0,3}$ plutôt que
par $\mathcal{M}_{0,3}^{\sim}$ et $\mathcal{N}_{0,3}^{\sim}$ que nous serons intéressés -
donc par le sous-groupe remarquable $\Gamma_{\mathbb{Q}}$ de
$\mathcal{N}_{0,3}^{1 \sim}$. Nous allons "débroussailler" un peu
le groupe $\mathcal{N}_{0,3}^{1 \sim}$ - et il s'agira dans la suite
du travail de trouver un maximum
de conditions nécessaires sur des éléments
de $\mathcal{N}_{0,3}^{1 \sim}$ pour provenir de $\Gamma_{\mathbb{Q}}$ - jusqu'à
trouver, si possible, une caractérisation (au
moins plausible, sinon prouvée?) des éléments
de $\Gamma_{\mathbb{Q}}$ - i.e. de $\Gamma_{\mathbb{Q}}$ lui-même.

\newpage

%% ===== Page 186 =====

392

II Détermination combinatoire des éléments de $\mathcal{N}_{0,3}^{!\sim}$.

1) Détermination par $h,g,f$.

Rappelons que

(21) $\mathcal{N}_{0,3}^{!\sim} =$ autom. de $\widehat{\mathcal{G}}_{0,3}^+ \simeq \widehat{Sl}(2,\mathbb{Z})^{+}$ qui stabilisent $\Pi_{0,3}$, induisent l'identité sur $\mathcal{G}_3 = \widehat{\mathcal{G}}_{0,3}^+/\Pi_{0,3}$ et respectent la structure à bouts (définie par la classe de $\epsilon_0$)

Un $u \in \mathcal{N}_{0,3}^{!\sim}$ est déterminé par un système

(22) $(h, g, f)$ $(h, g, f \in \Pi_{0,3})$

[MARGIN: $f$ déterminé mod $f \mapsto h_0^n f$ $n \in \hat{\mathbb{Z}}$]

satisfaisant les formules

(23) $\begin{cases} u(\sigma_\infty) = h \sigma_\infty \\ u(p) = g p \end{cases}$

La condition que ces formules définissent un automorphisme (ou du moins un endomorphisme) de $\widehat{\mathcal{G}}_{0,3}^+$ (qui induise néc. l'identité sur $\mathcal{G}_3$) est exprimée par les formules

(24) $(h \sigma_\infty)^2 = (g p)^3 = 1$

qui s'écrivent aussi

(25) $\begin{cases} h \sigma_\infty(h) = 1 \\ g p(g) p^2(g) = 1 \end{cases}$

et la condition que $u$ respecte les str. à bouts :

[MARGIN: Relations des bouts dans $\widehat{\mathcal{G}}_{0,3}^+$ (forme brute)]

(26) $u(\epsilon_0) = f^{-1} (\epsilon_0^\mu) f$ \hspace{2cm} $\mu \in \hat{\mathbb{Z}}^*$ inversible,

équivaut à : la condition

\newpage

%% ===== Page 187 =====

[MARGIN: relations des bouts dans $\Pi_{0,3}$ (forme brute)]

(27) $$u(l_0) = int(f^{-1})(l_0^\mu) ,$$

elle s'explicite aisément en termes de $h,g,f$ (sa
forme (26) de préférence) compte tenu que $\varepsilon_0 = \sigma_\infty p$

$$(h \sigma_\infty)(g p) = int(f^{-1}) \varepsilon_0^\mu \qquad \text{, i.e.}$$

$$h \sigma_\infty(g) \underbrace{\sigma_\infty p}_{\varepsilon_0} = f^{-1} \varepsilon_0^\mu f$$

$$f h \sigma_\infty(g) \varepsilon_0 f^{-1} = \varepsilon_0^\mu \qquad \text{soit}$$

[unclear: $\varepsilon_0(f) \varepsilon_0 = \sigma_\infty(p(f^{-1})) \varepsilon_0$]

$$f h \sigma_\infty(g p f^{-1}) = \varepsilon_0^{\mu-1} \qquad \text{soit, posant}$$

(28) $$\nu =^{df} \frac{\mu-1}{2} \in \hat{\mathbb{Z}}$$ [MARGIN: possible, car $\mu \in \hat{\mathbb{Z}}^*$ donc $\mu \equiv 1 \pmod 2$]

la relation (compte tenu que $\varepsilon_0^2 = l_0$)

$$f h \sigma_\infty(g p f^{-1}) = l_0^\nu ,$$

ou encore, appliquant $\sigma_\infty$ , comme $\sigma_\infty(l_0) = l_1$

$$\sigma_\infty(f) \sigma_\infty(h) g p f^{-1} = l_1^\nu$$

soit enfin

[MARGIN: $\nu \in \hat{\mathbb{Z}}$ NB $\nu$ est unique-ment déterminé [unclear: par $f, h, g$, et] $\mu = 2\nu+1 \in \hat{\mathbb{Z}}^*$ [unclear: et c'est le multiplicateur...]] 
(29) $$| \sigma_\infty(h) g = \sigma_\infty(f)^{-1} l_1^\nu p(f) |$$

(relation des bouts en $h, g, f$ ).

NB Il faut de plus exiger, pour que $h, g, f$ définisse
bien un autom. , que (24) est un préautom.
générateur de $\widehat{\mathcal{G}}_{0,3}^+$ , par la suite, une ou
[unclear: deux] [unclear: autres] [unclear: conditions]
liminaires".

\newpage

%% ===== Page 188 =====

2) Détermination par $\alpha_0, \beta_0, f$.

Il y a correspondance biunivoque entre les
couples $(g,h)$ satisfaisant (25), et les couples

(31) $(\alpha_0, \beta_0)$ \quad ($\alpha_0 \in \widehat{\Pi}_{0,3}^{\sim} = \Pi_{0,3}^{\sim} . \{1, \tau_\infty\}$, $\beta_0 \in \widehat{\Pi}_{0,3}^{\sim}$)

via les relations

(32) $\begin{cases} g = \alpha_0 p(\alpha_0)^{-1} = [\alpha_0, p] \\ h = \beta_0 \sigma_\infty(\beta_0)^{-1} = [\beta_0, \sigma_\infty] \end{cases}$

Ainsi, les $u \in \hat{M}_{0,3}^{\sim}$ correspondent aux systèmes

(33) $(\alpha_0, \beta_0, f)$ \quad $\alpha_0 \in \widehat{\Pi}_{0,3}^{\sim}$, $\beta_0, f \in \widehat{\Pi}_{0,3}^{\sim}$, $f$ déter-
[unclear: miné modulo un lacet multiplicatif $\in l_0^{\hat{\mathbb{Z}}}$]

satisfaisant la condition (29), considérée comme
condition en $\alpha_0, \beta_0, f$, où on y substitue à $g, h$
les valeurs données par (32). La relation s'écrit

$$ \sigma_\infty(\beta_0) \beta_0^{-1} \alpha_0 p(\alpha_0)^{-1} = \sigma_\infty(f^{-1}) l_1^\nu p(f) $$

ou encore

$$ \sigma_\infty(f \beta_0) \underbrace{\beta_0^{-1} \alpha_0}_{(f \beta_0)^{-1} (f \alpha_0)} p(f \alpha_0)^{-1} = l_1^\nu $$

ou encore

(34) \quad \fbox{$(f \alpha_0) p(f \alpha_0)^{-1} = ((f\beta_0) \sigma_\infty(f\beta_0)^{-1}) l_1^\nu$}

(relation des lacets en $\alpha_0, \beta_0, f$). Comme pour
(29), $\nu \in \hat{\mathbb{Z}}$ est uniquement déterminé en termes
de $\alpha_0, \beta_0$, par la possibilité de trouver $f$ satisfaisant
à la relation (34).

L'avantage de l'introduction de $\alpha_0, \beta_0$ au lieu

\newpage

%% ===== Page 189 =====

de $g, h$ est de rendre inutiles (car automatiques)
les relations (25), - il reste la seule relation
(34) à vérifier. Les relations (23) s'écrivent évidemment :
$$ (35) \qquad u(\sigma_\infty) = int(\beta_0)(\sigma_\infty), \quad u(\rho) = int(\alpha_0)(\rho) $$

3) Détermination par $(\alpha, \beta, f)$.

La forme particulière de (34), où le triple $(\alpha, \beta, f)$
n'intervient que par l'intermédiaire du
couple correspondant $(f\alpha_0, f\beta_0)$, conduit à poser

$$ (35) \qquad \alpha = f \alpha_0, \quad \beta = f \beta_0 \quad , $$

on trouve que les $u \in \tilde{M}_{0,3}^{1 \sim}$ sont déterminés

par les triples
$$ (36) \qquad (\alpha, \beta, f) \qquad \alpha \in \hat{\Pi}_{0,3}^\sim, \beta \in \hat{\Pi}_{0,3}^\sim, f \text{ déf modulo} $$
$$ \alpha \mapsto l\alpha, \beta \mapsto l\beta, f \mapsto l \circ f \quad (l \in \hat{Z}) $$
(même $u$ pour $\alpha, \beta, f$)

[MARGIN: Une disconvenance est l'indétermination du couple $(\alpha, \beta) \mapsto (l\alpha, l\beta)$ $(l \in \hat{Z})$, mais qui va d'ailleurs jouer un rôle plus bas]

satisfaisant à la condition

$$ (37) \qquad (\alpha \rho(\alpha)^{-1}) = (\beta \sigma_\infty(\beta)^{-1}) \cdot h_1^\nu $$

(relation des lacets en $\alpha, \beta$). L'avantage de
cette relation est que $f$ n'y figure plus.
On retrouve $\alpha_0, \beta_0, g, h$ en fonction de $\alpha, \beta, f$
par les formules
$$ (38) \qquad \alpha_0 = f^{-1} \alpha, \quad \beta_0 = f^{-1} \beta $$
$$ (39) \qquad g = \alpha_0 \rho(\alpha_0)^{-1} = f^{-1}(\alpha \rho(\alpha)^{-1}) \rho(f) = f^{-1} \xi \rho(f) $$
$$ (40) \qquad h = \beta_0 \sigma_\infty(\beta_0)^{-1} = f^{-1}(\beta \sigma_\infty(\beta)^{-1}) \sigma_\infty(f) = f^{-1} \eta \sigma_\infty(f) , $$
où on a posé
$$ (41) \qquad \xi = \alpha \rho(\alpha)^{-1}, \quad \eta = \beta \sigma_\infty(\beta)^{-1} \qquad (\xi, \eta \in \hat{\Pi}_{0,3}^\sim) $$

\newpage

%% ===== Page 190 =====

[MARGIN: 394]
Les formules (23) - (35) déterminent en termes
des données calculatoires, deviennent ici

(42) \quad $u(\sigma_\infty) = int(f^{-1}\beta)(\sigma_\infty) \quad u(p) = int(f^{-1}\alpha)(p)$ \quad i.e.
(42 bis) \quad $u = int(f^{-1}) u_{\alpha, \beta}$, où $u_{\alpha, \beta}$ ne dépend que de $\alpha, \beta$ :
(42 ter) \quad $\boxed{u_{\alpha, \beta}(\sigma_\infty) = int(\beta)(\sigma_\infty) = \eta \sigma_\infty, \quad u_{\alpha, \beta}(p) = int(\alpha)(p) = \xi p}$

La formule (37)
devient la
relation en $\xi, \eta$
(43) \quad $\xi = \eta h^\nu$

4) Détermination par $(\xi, \eta, f)$.

La donnée du couple $(\alpha, \beta) \in \hat{\Pi}_{0,3} \times \hat{\Pi}_{0,3}$ ,
équivaut à : celle du couple $(\xi, \eta)$ qui s'en
déduit par (41), assujetti aux équations

(44) \quad $$ \boxed{ \begin{cases} \xi p(\xi) p^2(\xi) = 1 \\ \eta \sigma_\infty(\eta) = 1 \end{cases} } $$

Ainsi, la donnée d'un élément de $\tilde{M}_{0,3}$ équivaut à celle
d'un triplet

(45) \quad $(\xi, \eta, f) \qquad \begin{cases} \xi, \eta, f \in \hat{\Pi}_{0,3} \quad modulo \\ \xi \mapsto l \xi p(l)^{-1}, \eta \mapsto m \eta \sigma_\infty(m)^{-1}, f \mapsto m f l^{-1} \end{cases}$

satisfaisant les relations (44), et la relation

(46)=(43) \quad $$ \boxed{\xi = \eta h^\nu} $$

qui joue le rôle de "relation de centre" en
$\xi, \eta$. L'avantage de cette présentation est

\newpage

%% ===== Page 191 =====

la simplicité maximale des équations (44), (46) -
et le fait que $f$ n'y intervienne pas :
les deux premières équations sont des équations
en $\xi$ seul et en $\eta$ seul, et la troisième relie
$\xi$ et $\eta$ de façon particulièrement simple - en
fait si $\mu=1$, i.e. $\nu=0$, c'est la relation $\xi=\eta$ !
L'inconvénient est qu'il y a trois équations

[MARGIN: mais n'y a-t-il pas une telle expression licite (cf (42 ter) : $u(\sigma_\infty) = int(f^{-1})(\sigma_\infty)$ $u(\rho) = int(f^{-1})(\rho)$ !?)]

au lieu d'une, et qu'on n'a pas d'expression
explicite de $u(\sigma_\infty)$ et de $u(\rho)$ en termes
de $\xi, \eta, f$ : il faut d'abord, pour cela,
résoudre les équations

$$ \alpha \rho(\alpha)^{-1} = \xi \quad (\alpha \in \hat{\Pi}_{0, \tilde{3}}^\wedge) $$
(47)
$$ \beta \sigma_\infty(\beta)^{-1} = \eta \quad (\beta \in \hat{\Pi}_{0, \tilde{3}}^\wedge) $$

en $\alpha, \beta$, pour pouvoir ensuite appliquer (42). !!)
Un autre inconvénient, d'importance, c'est
que lorsqu'on travaille avec les images des
éléments envisagés dans $\widehat{Gl}(2, \hat{\mathbb{Z}})^{\pm \prime}$ la connaiss.
(même dans $\widehat{Sl}(2, \hat{\mathbb{Z}})$ lui-même)
de ces images de $\xi, \eta$ est une donnée
assez faible, qui ne permet nullement de
retrouver $\alpha, \beta$ comme éléments de $\operatorname{Gl}(2, \hat{\mathbb{Z}})^{\pm \prime}$,
cf plus bas.

\newpage

%% ===== Page 192 =====

395

5) Composition aux autom. extérieurs. Relations de cocycles
pour un composé $u'' = u'u$.

Soit $\gamma \in \hat{\Pi}_{0,3}$, et considérons le compo-

$$ (48) \quad u' = int(\gamma) \circ u $$

il correspond à des invariants

$$ f', g', h', \alpha'_0, \beta'_0, \alpha', \beta', \xi', \eta' $$

donnés par

$$ (49) \quad \left\{ \begin{array}{l} 
f' = f \gamma^{-1}, \text{[unclear crossed out]} \\ 
g' = \gamma g \rho(\gamma)^{-1}, h' = \gamma h \sigma_\infty(\gamma)^{-1} \\ 
\alpha'_0 = \gamma \alpha_0, \beta'_0 = \gamma \beta_0 \\ 
\alpha' = \alpha, \beta' = \beta \\ 
\xi' = \xi, \eta' = \eta 
\end{array} \right. $$

Ainsi $\alpha, \beta$ ne changent pas quand on
remplace $u$ par $int(\gamma) u$, ils ne dépendent
que de l'autom. ext. $\bar{u}$ induit par $u$, i.e.
de l'image de $u$ dans $\underline{M}_{0,3}^{\sim}$. C'est là un
avantage substantiel de l'utilisation de

$$ (\alpha, \beta, f) \quad \text{-- ou } (\xi, \eta, f), \quad \text{de préférence à} $$

$$ (\alpha_0, \beta_0, f) \quad \text{-- ou } (g, h, f). $$

Soit maintenant

$$ (50) \quad u, u' \in \underline{M}_{0,3}^{\sim} \quad , \quad \text{soit } u'' = u'u $$

soient

$$ (51) \quad \begin{array}{l} 
f, g, h, \alpha_0, \beta_0, \alpha, \beta, \xi, \eta \\ 
f', g', h', \alpha'_0, \beta'_0, \alpha', \beta', \xi', \eta' \quad \text{et} \\ 
f'', g'', h'', \alpha''_0, \beta''_0, \alpha'', \beta'', \xi'', \eta'' 
\end{array} $$

\newpage

%% ===== Page 193 =====

les invariants relatifs à $u, u', u''$ respectivement, on aura
exprimé (dans la mesure du possible) ceux de $u''$
en fonction de ceux de $u, u'$, et on trouve

(52) 
$g'' = u'(g)g'$
$h'' = u'(h)h'$
[unclear: scribbled out text]
$\alpha''_0 = u'(\alpha_0) \alpha'_0 \quad *$
$\beta''_0 = u'(\beta_0) \beta'_0$
$\alpha'' = (f'u')(\alpha) \alpha' \quad *$
$\beta'' = (f'u')(\beta) \beta'$

[MARGIN: * $u(\alpha_0)$ et $(f'u')(\alpha)$ n'ont de sens à priori que si $\alpha_0, \alpha \in \hat{\Pi}_{0,3}$ (p. ex. $u \in \Gamma(3)$ $>$ on a pour $u$ une action dans un groupoïde de base $\dots$ d'où il y aura lieu d'écrire pour sauver les meubles $\alpha''_0 = u(\alpha_0) \alpha_0'$, avec $u(\alpha_0)$ ayant un sens clair $\dots$ ]

On voit donc que si on considère (en variables dans $\hat{\mathcal{M}}_{0,3}^{1 \sim}$) $f, g, h, \alpha_0, \beta_0, \alpha, \beta$ comme des
fonctions
$$ \hat{\mathcal{M}}_{0,3}^{1 \sim} \to \hat{\Pi}_{0,3} \text{, ou } \hat{T}_{0,3}^{\sim} $$
alors $f, g^{-1}, h^{-1}, \alpha_0^{-1}, \beta_0^{-1}$ satisfont la relation des
cocycles, et $\alpha^{-1}, \beta^{-1}$ aussi quand on se limite
aux sous-groupes pour lequel $f=1$, i.e. les
autom. de $\mathcal{M}_{0,3}$ qui fixent $L_0 = \{ l_0^n \mid n \in \hat{\mathbb{Z}} \}$. A vrai
dire, cet [unclear] est un peu boiteux,
à cause des indéterminations pour $f$, impli-
quant une indétermination correspondante pour
$\alpha, \beta$, de l'ordre $\alpha \mapsto l \alpha, \beta \mapsto l^{-1} \beta$. Pour
construire des cocycles canoniques correspondants
à $f, \alpha^{-1}, \beta^{-1}$, il faudrait introduire une grosse
extension de $\hat{\mathcal{M}}_{0,3}^{1 \sim}$ par $\hat{\mathbb{Z}}$, sur le modèle des
groupes de Teichmüller spinoriaux, sur laquelle
$f, \alpha, \beta$ soient définis sans ambiguïté. Je ne m'appesantis
pas sur ces détails...

\newpage

%% ===== Page 194 =====

## III Calculs dans $\operatorname{Gl}(2, \hat{\mathbb{Z}})'$ et dans $\operatorname{Gl}(2, \hat{\mathbb{Z}})$.

Pour $x \in M_{0,3}^{\sim}$ donné, d'où les quantités

(53) $f, g, h, \alpha_0, \beta_0, \alpha, \beta, \xi, \eta$

[MARGIN: NB le fait d'avoir mal déterminé $f, g, h$ (à un signe près), et le fait d'avoir des indéterminations sur $\alpha_0, \alpha, \beta_0, \beta$ ne doit plus nous inquiéter, ces éléments se trouvant de fait bien définis dans $\operatorname{Gl}(2, \hat{\mathbb{Z}})'$, et [unclear: même] (car de carré 1) bien définis dans $\operatorname{Gl}(2, \hat{\mathbb{Z}})$ [unclear] ! — de même pour $\xi, \eta$ définis à un signe près dans $\operatorname{Gl}(2, \hat{\mathbb{Z}})'$.]
dans $\Pi_{0,\hat{3}}^\sim$ , sauf $\alpha_0, \alpha$ qui ne sont que dans $\Pi_{0,\hat{3}}^\sim \rtimes \{1, \tau_0\}$ ) on désigne par les mêmes symboles leurs images dans

$\operatorname{Gl}(2, \hat{\mathbb{Z}})' \overset{def}{=} \operatorname{Gl}(2, \hat{\mathbb{Z}})/\pm 1$ ,

en fait dans

$\operatorname{Sl}(2, \hat{\mathbb{Z}})' = \operatorname{Sl}(2, \hat{\mathbb{Z}})/\pm 1$ ,

[MARGIN: cf p. 387 on utilisera ces formules plus bas]
sauf $\alpha_0, \alpha$ qui sont seulement dans $\operatorname{Gl}(2, \hat{\mathbb{Z}})'$

(54) $det(\alpha_0) = det(\alpha) = \varepsilon_3 \in \{\pm 1\}$

où $\varepsilon_3$ est défini par

(55) $\varepsilon_3 \equiv \mu \quad (3)$ ,

ce point est [unclear: clair].

(56) $f, g, h, \beta_0, \beta, \xi, \eta \in \operatorname{Sl}(2, \hat{\mathbb{Z}})'$ i.e.
$det(f) = det(g) = det(h) = det(\beta_0) = det(\beta) = det(\xi) = det(\eta) = 1$

Car la fonction $det$ a $det(-1) = 1$, d'où $det$ se factorise

$\operatorname{Gl}(2, \hat{\mathbb{Z}})' \xrightarrow{det} \hat{\mathbb{Z}}^\times$ ,

pour donner une suite exacte

(57) $1 \longrightarrow \operatorname{Sl}(2, \hat{\mathbb{Z}})' \longrightarrow \operatorname{Gl}(2, \hat{\mathbb{Z}})' \xrightarrow{det} \hat{\mathbb{Z}}^\times \longrightarrow 1$ )

On choisira des représentants

(58) $f, \alpha, \beta \in \operatorname{Gl}(2, \hat{\mathbb{Z}})$

ce qui équivaut au choix de trois signes ,

\newpage

%% ===== Page 195 =====

et en fonction de ces choix on définit les
autres quantités par

(59)
$$ \xi = \alpha \rho(\alpha)^{-1} \ , \ \eta = \beta \rho(\beta)^{-1} $$
$$ \alpha_0 = f^{-1} \alpha \ , \ \beta_0 = f^{-1} \beta $$
$$ g = \alpha_0 \rho(\alpha_0)^{-1} = f^{-1} \xi \rho(f) \ , \ h = \beta_0 \rho(\beta_0)^{-1} = f^{-1} \eta \rho(f) $$

On notera que $\xi$ et $\eta$ , $g , h$ sont indépendants
des choix des signes qui ont été faits. La
relation essentielle qui relie $(f, \alpha_0, \beta_0)$ , ou
plutôt $\alpha, \beta$ (en plus des relations
(59), qui jouent le rôle de relations de
définition) est la "relation des lacets", sous
forme matricielle, qui prend la forme
correspondant à (46) (ou : (43) )

(60) $$ \xi = \varepsilon \eta h^\nu $$ $$ \varepsilon \in \{\pm 1\} $$

où $\varepsilon$ est un signe qui ne dépend plus
d'aucun choix - on a (cf aussi formules
plus bas)

(61) $$ \varepsilon = \varepsilon_2 \overset{\text{défini}}{\underset{\text{par}}{\iff}} \mu \equiv \varepsilon_2 \ (4) $$ ;

\newpage

%% ===== Page 196 =====

et où $l_1$ est défini matriciellement par
$$l_1 = \begin{pmatrix} 1 & 0 \\ 2 & 1 \end{pmatrix}$$
d'où
(62) $$l_1^\nu = \begin{pmatrix} 1 & 0 \\ 2\nu & 1 \end{pmatrix} = \begin{pmatrix} 1 & 0 \\ \mu-1 & 1 \end{pmatrix} .$$

Remarques
(Relation avec les notations du § 4 (p. 356 ff, notamment
formulaire p. 364). Dans les groupes
$\widehat{\pi}_{0,3}$ et $\widetilde{\pi}_{0,3}$, les lettres $f, g, h$ du loc. cit.
ont le même sens qu'ici, mais non les lettres
$\alpha, \beta$ etc. Je vais distinguer par des soulignés,
dans le cas de [unclear: conflits] entre les deux notations,
les lettres avec la signification attachée dans
loc. cit, on trouve
(i) $$ \underline{\alpha} = \alpha_0, \quad \underline{\beta} = \beta_0 $$
à partir de p. 357 (i.e. $u(\sigma_\infty)^2 = \sigma_\infty^2$!)
on [unclear: relève] quelques typos
(ii) $\underline{\beta} = f \beta_0 = f$
et loc. cit (11) donne
$$ \underline{\varphi} = f \sigma_{\infty}(f)^{-1} = \beta \sigma_{\infty}(\beta)^{-1} = \eta $$
tandis que (loc. cit (64))
$$ \underline{\xi} = l_0^{2\nu} \underline{\varphi} = l_0^{2\nu} \eta $$
en comparant avec
$$ \xi = \eta l_1^{\nu} \qquad (46) $$
on trouve
(iv) $$ \underline{\xi} = l_0^{2\nu} \xi l_1^{-\nu} = \rho_{\infty}^2 \xi \rho(l_0^2)^{-1} $$
i.e. $\underline{\xi}$ est la $\rho$-transformée de $l_0^{2\nu}$.
La relation (loc. cit (29))
$$ \underline{\xi} = \underline{\alpha}_0 \rho(\underline{\alpha}_0)^{-1} $$
comparée à
$$ \xi = \alpha \rho(\alpha)^{-1} $$
et à (iv), donne ...

\newpage

%% ===== Page 197 =====

(v) $\alpha_0 = l_0^\nu \alpha$.

Les formules (i) à (v) permettent de
traduire le formulaire de loc. cit. en termes
du formulaire que je vais expliciter.

On explicite $\alpha$ et $\beta$ par leurs matrices

(63) $\alpha = \begin{pmatrix} a & b \\ c & d \end{pmatrix}$, $\beta = \begin{pmatrix} r & s \\ t & u \end{pmatrix}$

où $a, b, c, d, r, s, t, u \in \hat{\mathbb{Z}}$, avec

(64) $ad-bc = \delta$, $ru-st = 1$.

On aura

(65) $\xi = \alpha \rho(\alpha)^{-1} = \begin{pmatrix} A & B \\ C & D \end{pmatrix}$

avec (cf. p. 342, 08)

(66) $\begin{cases} \delta A = (a^2+ab+b^2) - (ac+bd+bc) & \delta B = ac+bd+bc \\ \delta C = -(c^2+cd+d^2) + (ac+bd+ad) & \delta D = c^2+cd+d^2 \end{cases}$

(67) $\eta = \beta \sigma_0(\beta)^{-1} = \begin{pmatrix} R & S \\ T & U \end{pmatrix}$, avec

(68) $R = r^2+s^2$, $S = T = rt+su$, $U = t^2+u^2$.

La relation des lacets (60) prend la forme

(69) $\begin{pmatrix} A & B \\ C & D \end{pmatrix} = \varepsilon \begin{pmatrix} R & S \\ T & U \end{pmatrix} \begin{pmatrix} 1 & 0 \\ 2\nu & 1 \end{pmatrix}$

$= \varepsilon \begin{pmatrix} R+2\nu S & S \\ T+2\nu U & U \end{pmatrix}$,

\newpage

%% ===== Page 198 =====

398

Notons les relations (conséquences de (66), (67)) :

(70) $$AD-BC = 1 \qquad RU-ST = 1$$

(71) $$C+D-B = 1 \qquad T-S=0 \text{ i.e. } T=S \quad ,$$

Cette relation des [crossed out: lacets] conjugués (69) [crossed out: prend] est équivalente au système de relations

(72)
$$ \begin{cases} D = \varepsilon U \\ \mu D = 1 \\ B = \varepsilon S \end{cases} \qquad (\text{où } \mu = 2\nu+1) $$

i.e. au système d'équations

(73)
$$ \begin{cases} c^2+cd+d^2 = \varepsilon(t^2+u^2) = \delta/\mu \\ ac+bd+bc = \varepsilon(rt+su) \end{cases} $$

[crossed out: illegible equation]

C'est de là par des considérations de congruences [crossed out: qu'on a ...] tire aussitôt

(74) $$\delta \equiv \mu \; (3), \quad \varepsilon \equiv \mu \; (4)$$

qui caractérisent les signes $\delta, \varepsilon$ comme

(75) $$\delta = \varepsilon_3(\mu) = \varepsilon_3, \quad \varepsilon = \varepsilon_4(\mu) = \varepsilon_4 \quad ,$$

cf th. p. 365.

La formule (83) de la page 370 prend ici la forme très simple (valable sans conditions autres que celles précisées jusqu'ici sur $\alpha, \beta, f$ - $f$ n'intervenant d'ailleurs pas dans le second membre de la formule) :

197

\newpage

%% ===== Page 199 =====

(76) $$ \delta(\alpha(c \rho + d)) = \beta(t \sigma_\infty + u) \tau_\infty^{\alpha_2} \qquad (\alpha = \alpha_-) $$

où $\alpha_2$ est défini par la condition
$$ (-1)^{\alpha_2} = \varepsilon \quad (= \varepsilon_2) . $$

Explicitant (76) par le calcul matriciel, on a
$$ \delta \alpha(c\rho+d) = \delta \begin{pmatrix} a & b \\ c & d \end{pmatrix} \begin{pmatrix} d & c \\ -c & c+d \end{pmatrix} = \delta \begin{pmatrix} 1 & B \\ 0 & D \end{pmatrix} $$
$$ \beta(t \sigma_\infty + u) \tau_\infty^{\alpha_2} = \begin{pmatrix} r & s \\ t & u \end{pmatrix} \begin{pmatrix} 1 & \varepsilon t \\ -t & \varepsilon u \end{pmatrix} = \begin{pmatrix} 1 & \varepsilon S \\ 0 & \varepsilon U \end{pmatrix} $$

de sorte que (76) se réduit aux équations
$$ D = \varepsilon U, B = \varepsilon S , $$
donc la relation des lacets [unclear]

(68, 69) est équivalente à la conjonction de
(76) et de la relation
$$ (77) \quad \mu D = 1 . $$

Ceci posé, on posera (pour $(\alpha, \beta)$ fixés de façon à satisfaire la formule des lacets)
$$ u_{\alpha, \beta} \overset{dfn}{=} \mu \delta \cdot \alpha(c\rho+d) = \mu \beta (t \sigma_\infty + u) \tau_\infty^{\alpha_2} = \mu \begin{pmatrix} 1 & B \\ 0 & D \end{pmatrix} = \begin{pmatrix} \mu & \mu B \\ 0 & 1 \end{pmatrix} $$
de façon que $u_{\alpha,\beta}$ [unclear: ne dépend en fait que]
de $\alpha$ et de $\beta$ et alors n'ignorant $\alpha_2$

[MARGIN: en fait, je pense que la décomp. [unclear] de la relat. (76) dans $\operatorname{Sl}(2, \hat{\mathbb{Z}})$ détermine complètement $\beta$ en fct. de $\alpha$. Cf. p. 410 pour la suite]

qui, lui, n'est pas déterminé par $\beta$, mais
la valeur de $\delta(c^2+cd+d^2) \pmod 4$
... ), et on posera, pour un triple
$(\alpha, \beta, f)$ d'éléments de $M(2, \hat{\mathbb{Z}})$, $dét=1$, $\alpha, \beta$
reliés par la formule des lacets

\newpage

%% ===== Page 200 =====

(79) $$u_{\alpha, \beta, f} \overset{dfn}{=} f^{-1} u_{\alpha, \beta} = (\mu \delta) f^{-1} \overset{\alpha_0}{\alpha} (c \rho + d) =$$
$$= \mu \underset{\beta_0}{\underbrace{f^{-1} \beta}} (t \sigma_\infty + u) \tau_\infty^{\alpha_\infty}$$

et on trouve, comme il se doit

[MARGIN LEFT: NB si $u, \sigma$ sont deux éléments d'un groupe $G$, pour [unclear] on pose bien sûr $u(\sigma) = u \sigma u^{-1}$]
(80) $$u_{\alpha, \beta, f} (\sigma_\infty) = int(\beta_0)(\sigma_\infty)$$
$$u_{\alpha, \beta, f} (\rho) = int(\alpha_0)(\rho)$$
[MARGIN RIGHT: ceci s'explique par le fait ( $u_{\alpha,\beta,f} \in \operatorname{Sl}(2, \hat{\mathbb{Z}})$ ) même [unclear] que l'hom. de [unclear] $u$ de $\widehat{B_{0,3}}$ se déduit (par [unclear]) d'un [unclear] partiel de $\operatorname{Sl}(2, \hat{\mathbb{Z}})$ ]

formules qui se réduisent à des [unclear: relations] où $f$ à nouveau disparait :

(81) $$u_{\alpha, \beta} (\sigma_\infty) = int(\beta)(\sigma_\infty)$$
$$u_{\alpha, \beta} (\rho) = int(\alpha)(\rho)$$

qui [crossed out: à leur tour] résultent trivialement des deux expressions de (78) de $u_{\alpha, \beta}$ !

Les formules (80) se complètent de façon naturelle en

(82) $$u_{\alpha, \beta, f} (\varepsilon_0) = \varepsilon_0$$
(avec en marge [unclear: $(\dots \varepsilon = \varepsilon_0^{\pm 1})$])

qui à son tour se réduit à

(83) $$u_{\alpha, \beta} (\varepsilon_0) = \varepsilon_0$$

où $f$ ne figure plus. Comme $\varepsilon_0 = \sigma_\infty \rho$, cette formule se réduit à :

$$u_{\alpha, \beta} (\sigma_\infty \rho) = u_{\alpha, \beta} (\sigma_\infty) u_{\alpha, \beta} (\rho) = \varepsilon_0$$

qui compte tenu de (81) se réduit à :

$$int(\beta)(\sigma_\infty) int(\alpha)(\rho) = \varepsilon_0$$

[crossed out: $int(\beta)(\sigma_\infty) int(\alpha)(\rho) = \dots$]

\newpage

%% ===== Page 201 =====

or le premier membre s'écrit

[unclear: crossed out lines]

$$ \beta \sigma_\infty \rho^{-1} \alpha \rho \alpha^{-1} = (\beta \sigma_\infty \beta^{-1} \sigma_\infty) (\sigma_\infty \rho) (\rho^{-1} \alpha \rho \alpha^{-1}) $$
$$ = \eta \cdot \varepsilon_0 \cdot \rho^{-1}(\xi) = \eta \underbrace{(\varepsilon_0 \rho^{-1})}_{\sigma_\infty}(\xi) \cdot \varepsilon_0 = \eta \sigma_\infty(\xi) \varepsilon_0 $$

donc la relation (83) s'écrit aussi

$$ \eta \sigma_\infty(\xi) = \varepsilon \varepsilon_0^{\mu-1} \quad (=\varepsilon l_0^{\nu'}) $$

or on a $\xi = \varepsilon y l_0^{\nu'}$ (relation des lacets (60))

donc

$$ \eta \sigma_\infty(\xi) = \varepsilon \eta \sigma_\infty(y) l_0^{\nu'} = \varepsilon l_0^{\nu'} \quad \text{, OK .} $$

Remarque. Toutes ces formules sont valables
dans $\widehat{Sl}(2, \mathbb{Z})$, sans ambiguïté de signes,
dès que $\alpha, \beta, f$ sont choisis. Si on change les
signes de $\alpha$ et de $\beta$ tous les deux, la relation
des lacets n'est pas modifiée - et $u_{\alpha, \beta}$ n'est
pas changé non plus, [unclear: partant ni] $u_{\alpha, \beta, f}$ (ni
ses transfo par $f$). Par contre, un
chgmt. de signe sur $f$ implique un chgmt. de
signe sur $u_{\alpha, \beta, f}$. Mais un tel changement de
signe ne change bien sûr rien à la validité
des formules (80), (82). Quand on fait

\newpage

%% ===== Page 202 =====

sur $(\alpha, \beta, f)$ le changement $l_0^n$
(84) $$ \alpha' = l_0^n \alpha, \quad \beta' = l_0^{-n} \beta, \quad f' = l_0^{-n} f \qquad (n \in \mathbb{Z}) $$
alors
$$ c' = c, d' = d \qquad t' = t, u' = u \qquad f'_0 = f_0 $$
(85) $$ \begin{cases} f'^{-1}\alpha' = f^{-1}\alpha \\ \alpha'_0 \qquad \alpha_0 \end{cases} \qquad \begin{cases} f'^{-1}\beta' = f^{-1}\beta \\ \beta'_0 \qquad \beta_0 \end{cases} $$

[unclear: crossed out text]

enfin
$\delta_{f'} (\det \alpha') = \delta_f (\det \alpha)$ (car $\det l_0^n = 1$)
(86) d'où $\mu' (= \delta_{f'}/({c'}^2+c'd'+{d'}^2)) = \mu (= \delta_f/(c^2+cd+d^2))$
d'où $\varepsilon'_2 = \varepsilon_2$, $\varepsilon'_3 = \varepsilon_3$ ... (voir (79))

c'est au-delà de ce qu'il faut pour vérifier que
(87) $u_{\alpha', \beta', f'} = u_{\alpha, \beta, f}$.

Rappelons aussi (cf page préc. 389)

Proposition Soient $\mu \in \hat{\mathbb{Z}}^*$, d'où les signes $\varepsilon_2 = \varepsilon_2(\mu)$ et $\varepsilon_3 = \varepsilon_3(\mu)$ (définis par $\mu \equiv \varepsilon_2 \; (4)$, $\mu \equiv \varepsilon_3 \; (3)$),
soient $c, d, t, u \in \hat{\mathbb{Z}}$. Pour qu'il existe $a, b, r, s \in \hat{\mathbb{Z}}$

[MARGIN: NB Pour le [unclear] des $\mu, c, d, t, u$ il [unclear] se dispenser de la relation (88) ! heureusement...]

donnant $\alpha = \begin{pmatrix} a & b \\ c & d \end{pmatrix}$, $\beta = \begin{pmatrix} r & s \\ t & u \end{pmatrix}$ satisfaisant (64)
$$ ad-bc = \varepsilon_3 \qquad ru-st=1 $$
et la relation des tresses
$$ \alpha \rho l_0^4 \alpha^{-1} = ( \beta \sigma_{\infty} \beta^{-1} ) l_1^{\nu} \qquad \text{où } \nu = (\mu-1)/2 $$
(liée aux relations (73)) il f. et s. qu'on ait
(88) $$ \varepsilon_3(c^2+cd+d^2) = \varepsilon_2(t^2+u^2) = 1/\mu \quad , $$
et l'ensemble des couples $(\alpha, \beta)$ satisfaisant à ces conditions est un torseur sous $\hat{\mathbb{Z}}$, opérant [unclear]

\newpage

%% ===== Page 203 =====

par $(\alpha, \beta) \longmapsto (l_0^n \alpha, l_0^{-n} \beta)$.

Corollaire Pour qu'on puisse trouver $\alpha, \beta$
tels qu'on ait de plus $\alpha \equiv \beta \equiv 1\ (2)$,
il f. et s. qu'on ait
(89) $$c \equiv t \equiv 0\ (2) \quad d \equiv u \equiv 1\ (2)$$
et alors l'ens. des solutions $(\alpha, \beta)$ satisfaisant
à ces conditions supplémentaires est un
torseur sous $\hat{\mathbb{Z}}$, opérant par
$$(\alpha, \beta) \longmapsto (l_0^n \alpha, l_0^{-n} \beta)$$

Scolie On peut donc dire que pour un
élément de $\mathcal{M}_{0,3}^{\wedge \sim}$ qui définit deux éléments $\alpha, \beta$ de $\widehat{\Pi}_{0,3}$, $\widehat{\Pi}_{0,3}$ modulo la transformation $(\alpha, \beta) \mapsto (l_0^n \alpha, l_0^{-n} \beta)$ (satisfaisant la relation
des lacets, il lui est attaché les
"invariants adéliques"
(90) $$\theta = (\mu, c, d, t, u)$$
($c, d, t, u$ déterminés aux ambiguïtés de
signes, sans plus) satisfaisant (88) -
et c'est ce qu'on peut faire de mieux dans le
contexte actuel : les $\alpha, \beta$ n'ont pas un caractère
invariant. Quand on a un élément de
$\mathcal{M}_{0,3}^{\wedge \sim}$ - i.e. un [unclear: homomorphisme] $f(etc)$, pas
seulement cet [unclear: homo.] - on a de plus un
$f \in \widehat{\Pi}_{0,3}$, donnant un $\bar{f} \in \operatorname{Sl}(2, \hat{\mathbb{Z}})'$, d'où
(91) $$f = \begin{pmatrix} g & * \\ l & m \end{pmatrix}$$ défini mod un signe
et $\mod f \mapsto l_0^n f$

\newpage

%% ===== Page 204 =====

401

[unclear: crossed out] Mais cette donnée supplémentaire est un peu trop faible pour reconstituer l'élément $\underline{u} = u_{\alpha, \beta, f}$ de $\operatorname{Gl}(2, \hat{\mathbb{Z}})'$ associé à $u$.

Il vaut mieux dire que pour $\mu, c, d, t, u$ fixés on a un $\hat{\mathbb{Z}}$-torseur des $(\alpha, \beta) \equiv (1, 1) \ (2)$ satisfaisant aux conditions dites dans la proposition et le corollaire, soit $T = T_\Theta$, on fait opérer alors $\hat{\mathbb{Z}}$ diagonalnt sur $T \times \operatorname{Sl}(2, \hat{\mathbb{Z}})$ par
$$(\alpha, \beta, f) \longmapsto (l_0^n \alpha, l_0^{-n} \beta, l_0^{-n} f)$$
et on regarde le quotient - c'est donc le torseur par $\check{T}$ (torseur opposé à $T$) de $\operatorname{Sl}(2, \hat{\mathbb{Z}})$, sur lequel $\hat{\mathbb{Z}}$ opère par $f \mapsto l_0^{-n} f$. [unclear: crossed out] d'un élément de
$$\check{T} \wedge^{\hat{\mathbb{Z}}} \operatorname{Sl}(2, \hat{\mathbb{Z}})$$
[MARGIN: c'est un $\operatorname{Sl}(2, \hat{\mathbb{Z}})$-torseur à droite]
qui permet de définir un élément
$$\underline{\underline{u}} = u_{\alpha, \beta, f} ,$$
qui changera de signe si $f$ est remplacé par $(-f)$.

Décidément, ces formulations en l'air sont peu enthousiasmantes, et il devient clair qu'il faudra introduire explicit-

\newpage

%% ===== Page 205 =====

ment et étudier de façon un minimum
de soin le groupe adélique formé par les
systèmes $(\alpha, \beta, f)$ satisfaisant les relations des
loc. cit, et dégager l'homomorphisme canoni-
que de $\tilde{M}_{0,3}$ dans ledit groupe.
Mais je voudrais avant cela terminer les
gammes notationnelles.

\newpage

%% ===== Page 206 =====

402

## IV Les trois cas particuliers fondamentaux.

On revient aux éléments $g,h,f$ etc dans $\hat{\Pi}_{0,3}$ ou $\tilde{\Pi}_{0,3}^{\wedge}$ qui correspondent à $u$, pouvant être décrit de quatre façons différentes par l'une ou l'autre des quatre systèmes
(92) $(g,h,f)$, $(\alpha_0, \beta_0, f)$, $(\alpha, \beta, f)$, $(\xi, \eta, f)$.

1er Cas :
(93) $h=1$, i.e. $\beta_0=1$, i.e. $\beta=f$, i.e. $\eta=f\sigma_1(f)^{-1}$,
i.e.
(94) $[u, \sigma_{\infty}]=1$ .

Dans un système $(g,h,f)$ etc quelconque, il y a un automorphisme intérieur et un seul $Int(\gamma)$ tel que $u'=Int(\gamma)(u)$ satisfasse à : la propriété précédente i.e. commute à $\sigma_{\infty}$, ce sera pour
$\gamma = \beta_0^{-1}$ (cf formules (49)).

Les $u$ satisfaisant à (93), exprimés en termes de triples $(\xi, \eta, f)$ satisfaisant les trois relations (44) (46) (ne faisant pas intervenir $f$), correspondent ainsi (par la relation $\eta=f\sigma_1(f)^{-1}$ de (93)) aux couples $(\xi, f)$ d'éléments de $\hat{\Pi}_{0,3}$, satisfaisant :

(95) $\xi\rho(\xi)\rho^2(\xi)=1$
(96) $\xi = (f\sigma_{\infty}(f)^{-1})l_f^2$
avec $l_f \in \hat{\mathbb{Z}}$ comme notée, tel que $2l_f \in \hat{\mathbb{Z}}^*$

[MARGIN: NB on aura $\xi = f g \rho(f)^{-1}$ au lieu de $f$ il a fallu d'écrire $f'$ déduit de $f$ [unclear]]

\newpage

%% ===== Page 207 =====

ou aux couples $(\alpha, \beta)$ d'éléments de
$\hat{\Pi}_{0,3}$ satisfaisant la relation des [unclear: lacets] (37)
[MARGIN: $(\alpha, \beta)$ engendrent $\hat{\Pi}_{0,3}$ (cf. § 4, p. 73)]

(97)=(37) $$(\alpha\rho(\alpha)^{-1}) = (\beta\tau_\infty(\beta)^{-1}) h^k$$
(qu'on complétera dans ce qui suit par $f = \beta$).
Au changement près de $\alpha$ en $l_0^r \alpha = \alpha_0$
(c'étaient les conventions du § 4 (  )
la formule prend encore la forme $$\underbrace{\alpha_0\rho(\alpha_0)^{-1}}_{\xi} = \underbrace{l_0^r (f t_0 l_0^r)^{-1}}_{f \eta}$$

Compte tenu de (79), les conditions
[crossed out: 2ème Cas]
[crossed out: $g=1$, i.e. $x=\beta$, i.e. $f=a$, i.e. $\xi = f\rho(f)^{-1}$ i.e.]
(93) (94) prennent aussi la forme

(98) $$u_{\alpha, \beta, f} = \mu (t_\infty + u)\tau_\infty^{\alpha_2}$$

et expriment que $u = u_{\alpha, \beta, f}$ normalise
la sous-algèbre de l'algèbre des matrices
$M_2(\hat{\mathbb{Z}})$ engendrée par $\sigma_\infty$ (compte tenu
de la relation
(99) $$\tau_\infty(\sigma_\infty) = - \sigma_\infty \text{ dans } M_2(\hat{\mathbb{Z}}) \text{ ).}$$

Si on regarde le sous-groupe de
$\hat{M}_{0,3}^{! \ \sim}$ formé des $u$ qui satisfont aux
conditions (équivalentes) (93) (94) (98)

(100) $$M_{0,3}^{! \ \sim}(1/2) = \{ u \in \hat{M}_{0,3}^{! \ \sim} \ | \ [u, \tau_\infty] = 1 \} \xrightarrow{\sim} \hat{M}_{0,3}^{! \ \sim}$$
on trouve que l'homomorphisme

\newpage

%% ===== Page 208 =====

[MARGIN: NB $\mathcal{M}_{0,3}^{! \sim}$ (101) est stable par $x \mapsto \sigma_{\infty} x \sigma_{\infty}^{-1}$ et par $x \mapsto \tau_{\infty} x \tau_{\infty}^{-1}$ (qui engendrent un groupe $\simeq S_3$)]
[MARGIN: sûrement exact, d'ailleurs aussi les courbes elliptiques d'invariant modulaire $j=0, 1728$ (correspondent à deux orbites $\{1/2, 2, -1\}$ et $\dots$ dans $\mathcal{M}_{0,3}^{\sim}$ sous $\mathcal{U}_{0,3}$)]

(101) $$ \mathcal{N}_{0,3}^{! \sim} \simeq \mathcal{M}_{0,3}(1/2)^\sim \xrightarrow{k_{1/2}} \mathcal{M}_{0,3}^{! \sim} \xrightarrow{u \mapsto \underline{u} = (u_{\alpha,\beta,f})} \operatorname{Gl}(2, \hat{\mathbb{Z}})' $$

factorise par le dit normalisateur

précisant la section en termes du groupe $\operatorname{Gl}(2, \hat{\mathbb{Z}}) \simeq M_2(\hat{\mathbb{Z}})^*$, plutôt que de l'algèbre des matrices, elle se factorise par le normalisateur du tore maximal

(102) $$ T_{\hat{\mathbb{Z}}} = \{ t \sigma_\infty + u \quad (t,u \in \hat{\mathbb{Z}}, \quad u^2-t^2 \in \hat{\mathbb{Z}}^*) \} $$
[unclear: $= \det(t\sigma_\infty+u)$]

Je suspecte qu'il n'y a que le $k_{1/2}$ de (101) de l'extension $\mathcal{M}_{0,3}^{! \sim}$ (de $\mathcal{N}_{0,3}^{! \sim}$ par $\hat{\Pi}_{0,3}$, que dis-je ou de $\mathcal{N}_{0,3}^{! \sim}$ par $\widehat{\mathcal{T}}_{0,3}^{+}$) qui se factorise par ce normalisateur.

2e Cas C'est le cas

(103) $g = 1$ i.e. $\beta \alpha = 1$ i.e. $\alpha = f$ i.e. $\xi = [\rho, f \rho f^{-1}]$, i.e.

(104) $[\alpha, \beta] = 1$.

Pour $u$ quelconque, il y a un $\gamma \in \hat{\Pi}_{0,3} \cdot \{1, \tau_\infty\}$ tel que $u' = i_A(\gamma) \circ u = \gamma u$ commute ainsi à $\rho$, savoir $\gamma = \beta \alpha^{-1}$. Mais attention que $\gamma$ n'est pas toujours dans $\hat{\Pi}_{0,3}$ , en fait

(105) $\gamma = \alpha \beta^{-1} \in \hat{\Pi}_{0,3} \iff \varepsilon_3(u) = 1$ i.e. $f(u) \equiv 1 (3)$

\newpage

%% ===== Page 209 =====

Si on désigne par $\mathcal{N}_{0,3}^{!\sim}(-j)$ le sous-groupe
de $\mathcal{M}_{0,3}^{!\sim}$ formé des autom.
(103)(105)), image de $\mathcal{N}_{0,3}^{!\sim}$ et du
sous-groupe d'indice 2 noyau de
$$u \mapsto \varepsilon_3(\mu(u)) = \varepsilon_3(u)$$
i.e.
(105) $\mathcal{N}_{0,3}^{!\sim}(-j) = \{ u \in \mathcal{N}_{0,3}^{!\sim} \mid \varepsilon_3(u) = 1 \}$
i.e. $\bar{\mu}(u) \equiv 1 (3)$

[MARGIN: NB $\mathcal{M}_{0,3}^{!\sim}(-j)$ est stable par $x \mapsto \rho x \rho^{-1}$]

et l'on a
(107) $\mathcal{M}_{0,3}^\sim[-j] \overset{\sim}{\longrightarrow} \mathcal{N}_{0,3}^{!\sim}[-j]$ .

En réunissant on trouve donc un
hom. composé
(108) $\mathcal{N}_{0,3}^{\Delta \sim}(-j) \overset{\sim}{\longrightarrow} \mathcal{M}_{0,3}^{!\sim}(-j) \hookrightarrow \mathcal{M}_{0,3}^{!\sim}$
$$k(-j)$$
dont le composé avec $\mathcal{M}_{0,3}^{!\sim} \to \dots$ est un
hom. qui se factorise par le tore
maximal $T$ de $\operatorname{Gl}(2, \hat{\mathbb{Z}})'$, formé des
(109) $\{ c\rho + d \mid c,d \in \hat{\mathbb{Z}} , c^2 + cd + d^2 \in \hat{\mathbb{Z}}^\times \}$.

[MARGIN: canular, pour cause, vu que tout à l'heure...]

Je suspecte que c'est le seul sous-groupe
de $\mathcal{M}_{0,3}^{!\sim}$ (voire même de $\mathcal{M}_{0,3}^\sim$)
qui ait cette propriété. Celles-ci résulte
de la forme équivalente des conditions
(103) (101), formées en (79) :

\newpage

%% ===== Page 210 =====

(110) $$u = u_{\alpha, \beta, f} = \mu(cp+d)$$

(NB $\delta = \varepsilon_3(u)$ est égal à $1$ dans ce cas).

Les $u$ en question correspondent aux couples $(\eta, f)$ satisfaisant

(111) $$\eta\sigma_{\varpi}(\eta) = 1$$
(112) $$(f\rho(f))^3 = \eta \varepsilon_1^n$$ [MARGIN: qu'on complétera par $\xi = f \rho(f)^{-1}$]

[MARGIN: On peut écrire la relation (111) en termes de $f$, vu que [unclear] $\alpha, \beta$]

ou encore aux couples $(\alpha, \beta)$ satisfaisant la relation de loc. cit. $(97) = (87)$, qui se complète donc par $f = \alpha$.

3° Cas C'est celui où

(113) $$f \equiv 1\ (L_0) \text{ i.e. } \alpha \equiv \alpha_0\ (L_0 \text{ à g.}) \text{ i.e. } \beta \equiv \beta_0\ (L_0 \text{ à g.})$$
(114) $$u(\varepsilon_0) = \varepsilon_0^\mu\ (\mu \in \hat{\mathbb{Z}}^*) \text{ i.e. } [u, \varepsilon_0] = l_0^\nu\ (\nu \in \hat{\mathbb{Z}} - \nu = \frac{\mu - 1}{2})$$

[MARGIN: NB $L_0 = l_0^{\hat{\mathbb{Z}}}$]

Dans ce cas, modulo l'opération de $L_0 = \{l_0^n \mid n \in \hat{\mathbb{Z}}\}$ - ou encore on normalise $L_0 = \{\varepsilon_0^n \mid n \in \hat{\mathbb{Z}}\}$ - d de même min.
tion $$(\alpha, \beta, f) \longmapsto (l_0^n\alpha, l_0^n\beta, l_0^{-n}f)$$ on exigera

(115) $$f = 1$$

ce qui déterminera $\alpha, \beta$ de façon unique.

[MARGIN: On a renoncé de noter ce sous-groupe par $M_{0,3}^{1\sim}[0]$ car il n'est pas l'ext. par un sect. de $M_{0,3}^{1\sim}$ par $N_{0,3}^{1\sim}$]

Le sous-groupe $M_{0,3}^{1\sim}[0]$ des $u \in M_{0,3}^{1\sim}$ qui normalisent $L_0$ est en corr. bijective avec l'ens des couples
$$(\alpha, \beta) \quad (\alpha \in \hat{\pi}_{0,3}^\sim = \hat{\pi}_{0,3}.\{1, \tau_0\}, \beta \in \hat{\pi}_{0,3}^\sim)$$ satisfaisant
(112')

\newpage

%% ===== Page 211 =====

[MARGIN: Dans tout ceci, je laisse de côté la "condition de lisière" sur les relations qui en découle]

[unclear] relations des boîtes
(11'') $$(\alpha \rho l_0 \alpha^{-1}) = (\beta \rho l_1 \beta^{-1}) l_1^\nu$$ $(l_1^\nu \text{ où } \nu \in \hat{\mathbb{Z}} \text{ conv.})$
(pour $(\alpha, \beta) \mapsto (\alpha \rho(\alpha), \beta \rho(\beta) l_1^\nu)$, et $l_0 l_1 l_2 = 1$)
ou encore aux couples $(g, h)$ d'éléments de $\hat{\Pi}_{0,3}$ satisfaisant

(116) [MARGIN: = (44) + (45)] $$g \rho l_0 g^{-1} = 1, \quad h \rho l_1 h^{-1} = 1, \quad g = h l_1^\nu \quad (\nu \in \hat{\mathbb{Z}} \text{ conv.})$$

au couple $(\alpha, \beta)$ ou $(g, h)$ correspond alors
l'automorphisme $u$ défini par

(117) [MARGIN: = (23)] $$u(v_0) = \alpha v_0 \quad , \quad u(p) = g p$$
ou encore par

(118) $$u(l_0) = int(\alpha) l_0 \quad , \quad u(p) = int(\beta) p$$

Dans le groupe $\widehat{M_{0,3}^{1 \sim}}[0]$, il y a le sous-groupe

(119) $$\widehat{M_{0,3}^{1 \sim}}[0] \cap \hat{\Pi}_{0,3} \simeq L_0$$

l'inclusion

(120) $$\hat{\mathbb{Z}} \simeq L_0 \hookrightarrow \widehat{M_{0,3}^{1 \sim}}$$

étant décrite par

(121) $$\lambda = l_0^n \mapsto (\alpha = l_0^{\frac{n}{2}}, \beta = \lambda^0)$$ [crossed out]

ou encore

(122) $$\lambda = l_0^n \mapsto (f = l_0^n l_1^{-n}, g = l_0^{-n} l_1^{-n})$$
i.e. $f = g = \dots$ [crossed out]

la multiplication par $\lambda \dots$ par $\lambda = l_0^n$ dans $\widehat{M_{0,3}^{1 \sim}}[0]$ étant explicitée par

(123) $$(\alpha, \beta) \mapsto (l_0^n \alpha, l_0^{-n} \beta)$$

(124) $$(f, g) \mapsto (l_0^n f l_1^{-n}, l_0^{-n} g l_1^{-n})$$

\newpage

%% ===== Page 212 =====

405

Comme on a un iso
$$M_{0,3}^{1 \sim}[\vec{0}]/L_0 \xrightarrow{\sim} \mathcal{N}_{0,3}^{1 \sim} \quad ,$$
on trouve le

Corollaire Le groupe $\mathcal{N}_{0,3}^{1 \sim}$ est en bijection
canonique avec l'ens. des couples
$(\alpha, \beta) \in \Pi_{0,\vec{s}}^\sim \times \Pi_{0,\vec{s}}^\sim$ satisfaisant la relation
de Coates (97)
$$\alpha \rho(\alpha)^{-1} = (\beta \tau_0(\beta)^{-1}) l_1^\nu \quad (\text{pour } \nu \in \hat{\mathbb{Z}} \text{ conv.})$$
modulo les transformations $(\alpha, \beta) \mapsto (l_0 \alpha, l_0 \beta)$.
(de $\mathcal{M}_{0,3}^{1 \sim}(\vec{0})$ ou de son quotient $\mathcal{N}_{0,3}^{1 \sim}$)

La loi de groupe s'explicite en termes
de $(g,h)$, ou de $(\alpha,\beta)$, par les "formules
de cocycles" (52) : si $u', u \text{ et } u''=u'u$ sont donnés
par $(g',h'), (g,h), (g'',h'')$, -- par
$(\alpha',\beta'), (\alpha,\beta) \text{ et } (\alpha'',\beta'')$, -- on a :

(125)
$$g'' = u'(g) g' \qquad h'' = u'(h) h'$$
$$\alpha'' = u'(\alpha) \alpha' \qquad \beta'' = u'(\beta) \beta'$$

Remarque sur les notations $\mathcal{M}_{0,3}^{1 \sim}(q)$, pour $q = \frac{1}{2}, -1, 0$.
En fait pour tout point algébrique
(126) $q \in \mathcal{M}_{0,3}(\mathbb{C}) = \mathbb{C} \setminus \{0,1\} = \mathbb{P}^1(\mathbb{C}) \setminus \{0,1,\infty\}$
il est tout à fait naturel d'introduire un scindage partiel

[MARGIN: En fait, il faut [unclear: se donner] plus précisément un [unclear: point] (127) $q^\sim$ au-dessus de $q$, le scindage [unclear: dépendant de ce choix...]]

$$K_q : \mathcal{N}_{0,3}^{1 \sim} \hookrightarrow \mathcal{M}_{0,3}^{1 \sim} \quad , \quad Im K_q = \mathcal{M}_{0,3}^{1 \sim}(q)$$
de l'extension $\mathcal{M}_{0,3}^{1 \sim}$ de $\mathcal{N}_{0,3}^{1 \sim}$ par $\Pi_{0,\vec{s}}^\sim$, --

\newpage

%% ===== Page 213 =====

[MARGIN: La définition implicite [unclear] lorsque $q$ est élément d'une extension cyclotomique de $\mathbb{Q}$, de telle façon que [unclear] coïncide [unclear] sous-groupe de $\Gamma_Q(q)$ et de $\Gamma_Q$]

du moins un scindage partiel au dessus
de $\Gamma_Q(q)$, qu'on s'efforcera de déterminer ;
$\mathcal{M}_{0,3}^{!\sim}(q)$ désigne un sous-groupe
topologique de $\mathcal{M}_{0,3}^{!\sim}$ qu'il
s'agira [unclear: sans doute aussi] : déterminer, si ce n'est
par type le plus fin ... ) , c-à-d.
$\Gamma_Q$ du sous-groupe ouvert $\Gamma_Q(q)$ de $\Gamma_Q$
stabilisateur de $q$. L'analyse des cas

$1^\circ$ et $2^\circ$ ci-dessus reviendrait en fait à
la détermination de tels scindages pour
$q = \frac{1}{2}$ et pour $q = \overline{j}$ - dans le premier
cas $\Gamma_Q(q)$ est $\Gamma_Q$ lui-même, dans le
deuxième c'est le sous-groupe d'indice 2
noyau de $\varepsilon_2 \circ \rho : \Gamma_Q \to \{ \pm 1 \}$, et [unclear: de même] pour
$\mathcal{M}_{0,3}^{!\sim}(q)$. En transformant par [unclear]
on obtiendrait des scindages $K_{1/2}$ les scindages
$K_2$ et $K_{-1}$ - en transformant par une
quelconque des [unclear] on
déduirait du scindage $K_{\overline{j}}$ un scindage de
type $K_{-\overline{j}}$ .
Quant à $\mathcal{M}_{0,3}^{!\sim}[0]$, il correspond (au point

\newpage

%% ===== Page 214 =====

406

0 de $\hat{U}_{0,3} = \mathbb{P}^1_{\mathbb{Q}}$, il y correspond deux sous-groupes
analogues $\tilde{M}_{0,3}[q]$, avec $q=1, q=\infty$, correspon-
dant aux deux autres points "à l'infini" de
$U_{0,3}$ ; et qui se déduisent du précédent en
appliquant $\rho, \rho^2$. J'ai dit des bêtises en disant que
ceci ne correspond pas [unclear] à des scindages partiels
de $\tilde{M}_{0,3}$ sur $\tilde{N}_{0,3}$, puisque l'intersec-
tion n'est pas réduite à $1$, mais à
l'un des trois groupes-lacets type $L_0, L_1, L_{\infty}$.
Mais on sait qu'une [above: le choix d'] uniformisante de
l'anneau local de $\hat{U}_{0,3}$ en $0$ (disons) déter-
mine un scindage de $\tilde{M}_{0,3}(0)$, c-à-d d'une
ext. de $\tilde{N}_{0,3}$ par $L_0$ - et il faudra
dans la suite essayer de saisir un tel scindage
(associé à une uniformisante)
dans l'optique de notre calcul, comme
une façon privilégiée de [unclear: normaliser]
pour [unclear] donc (à un $(\alpha, \beta)$ satisfaisant
aux relations des lacets, de "normaliser"
le choix des $(\alpha, \beta)$ mod $L_0$. J'y revien-
drai par la suite, sûrement.

\newpage

%% ===== Page 215 =====

407

V Récapitulation en termes de $\mathcal{M}_{0,3}^{1 \sim}[0]$

Il apparaît finalement que la façon la plus simple d'étudier les autom. de lacets de $\hat{\Pi}_{0,3}$ (comme on a vu de $\hat{G}_3$, i.e. les éléments de $\hat{N}_{0,3}^{1\sim}$) est de commencer par l'étude dans $M_{0,3}^{1 \sim}$ du sous-groupe $\mathcal{M}_{0,3}^{1\sim}[0]$, normalisateur de $L_0$.

Posant pour $u \in \mathcal{M}_{0,3}^{1\sim}[0]$ :
(128) $$u(\sigma_\infty) = \eta \sigma_\infty \quad , \quad u(\rho) = \xi \rho \quad (\eta, \xi \in \hat{\Pi}_{0,3})$$
i.e.
(129) $$[u, \sigma_\infty] = \eta \quad \quad [u, \rho] = \xi$$
qui expriment la commutation ext. de $u$ à $\mathfrak{S}_3$ (agissant sur $\hat{\Pi}_{0,3}$ par via le $\rho$), la condition de normalisation de $L_0$ s'exprime par une relation de la forme

[MARGIN: (130 bis) (130) équivaut d'ailleurs
vu que $u(E_1) = id$ & $u(E_\rho) = id$
à $u(E_\infty) = 1$ (cf (115') p 260)]
(130) $$u(\varepsilon_0) = \varepsilon_0^\mu \quad \text{ou} \quad \varepsilon(l_0) = l_0^\mu \quad (\mu \in \hat{\mathbb{Z}}^\times).$$

Or on a (compte tenu de $l_0 = \sigma_\infty \rho$, il vient par un calcul immédiat) la form.-
(*31) $$u(\varepsilon_0) = u(\sigma_\infty)u(\rho) = \eta \sigma_\infty \xi \rho = \eta \sigma_\infty \xi (\sigma_\infty \varepsilon_0) = \eta \sigma_\infty(\xi) \varepsilon_0$$
donc (130) s'écrit
$$\eta \sigma_\infty(\xi) = \varepsilon_0^{\mu-1} \quad \text{soit, en posant } \mu-1 = 2\nu \text{ } (\nu \in \hat{\mathbb{Z}})$$
$$\eta \sigma_\infty(\xi) = l_0^\nu \quad \text{soit, appliquant } \sigma_\infty$$
(*31) $$\sigma_\infty(\eta) \xi = l_1^\nu$$

Or les relations
(131) $$(\eta \sigma_\infty)^2 = 1 \quad , \quad (\xi \rho)^3 = 1$$
s'écrivent aussi
(133) $$\eta \sigma_\infty(\eta) = 1 \quad , \quad \xi \rho(\xi) \rho^2(\xi) = 1 \quad ,$$
d'où $\sigma_\infty(\eta) = \eta^{-1}$, de sorte que (*31) s'écrit, de façon plus commode
(134) $$\xi = \eta l_1^\nu \quad \text{(relation des lacets)}$$

214

\newpage

%% ===== Page 216 =====

On trouve ainsi le

Théorème Les éléments $u$ de $M^{\sim \, !}_{0,3}[ \widehat{\mathbb{O}} ]$ i.e. les
autom. ext. de $\widehat{\Pi}_{0,3}$ qui commutent ext. à
$\mathfrak{S}_3$ et qui normalisent $L_0$ sont en corr.
(via les formules (128) ou 129))
biunivoque (avec les couples $(\xi,\eta)$ ($\xi,\eta \in \widehat{\Pi}_{0,3}$) )
satisfaisant les relations (133)(134), et tels que
de plus ( "condition liminaire")
les autom. (132) (équivalents à (133)) soient généra-
-trices. 

[MARGIN: cette condition $2\nu+1 \in \hat{\mathbb{Z}}^*$ est en fait conséquence des rel. (133), (134), cf plus bas th. p. 409!]

Dans (134) $\nu$ est un élém.
convenable de $\hat{\mathbb{Z}}$, évidemment uniquement déterminé,
et assujetti : la condition liminaire sera $2\nu+1 \in \hat{\mathbb{Z}}^*$
(qui est impliquée par la condition liminaire, et
peut-être l'implique...)

L'inclusion évidente
(135) $$L_0 \hookrightarrow M^{\sim \, !}_{0,3}(\omega)$$
$$l_0^\lambda \mapsto i_\sim(\lambda)$$
est donnée par
(136) $$l_0^\lambda \mapsto (\xi = l_0^\lambda \rho(l_0^\lambda)^{-1} = l_0^\lambda l_1^\lambda, \eta = l_0^\lambda \tau_2(l_0^\lambda)^{-1} = l_0^{-\lambda} l_1^{-\lambda})$$
et plus généralement l'action à gauche de $L_0$
sur $M^{\sim \, !}_{0,3}(\omega)$ est donnée, pour $\lambda \in \hat{\mathbb{Z}} \simeq L_0$, par
(137) $$(\xi, \eta) \mapsto (\xi' = l_0^\lambda \xi l_1^{-\lambda}, \eta' = l_0^{-\lambda} \eta l_1^{-\lambda})$$
donc
$$M^{\sim \, !}_{0,3} / L_0 = N^{\sim \, !}_{0,3}$$
est isomorphe au quotient de l'ensemble des
couples $(\xi, \eta)$ satisfaisant (133)(134) [et la condition
liminaire], par les opérations (137) - compte tenu
d'une évidente structure de groupe naturelle; d'ailleurs
sur ces couples $(\xi, \eta)$, donnée par

\newpage

%% ===== Page 217 =====

408

$$ (138) \quad (\xi',\eta') (\xi,\eta) = (\xi'',\eta'') $$

avec

$$ (139) \quad \xi'' = u'(\xi)\xi' \quad , \quad \eta'' = u'(\eta)\eta' $$

$u'$ est [unclear: l'automorphisme] défini par $(\xi',\eta')$.

Corollaire Tout $u \in \tilde{M}_{0,3}^{!\sim}$ s'écrit sous la forme

$$ (140) \quad u = f^{-1} u_{\xi,\eta} \quad , \quad = int(f^{-1}) \circ u_{\xi,\eta} $$

avec $f \in \hat{\Pi}_{0,3}$ et $(\xi,\eta)$ satisfaisant aux conditions des théorèmes précédents, de sorte qu'on aura

$$ (141) \quad u(\sigma_\infty) = h \sigma_\infty \quad , \quad u(\rho) = g \rho \qquad , avec $$
$$ (142) \quad h = f^{-1} \eta \sigma_\infty(f) \quad , \quad g = f^{-1} \xi \rho(f) $$

Pour qu'on ait

$$ f^{'-1} u_{\xi',\eta'} = f^{-1} u_{\xi,\eta} $$

il faut et il suffit qu'il existe $n \in \hat{\mathbb{Z}}$ tel que

$$ (144) \quad \xi' = l_0^n \xi l_1^{-n}, \quad \eta' = l_0^n \eta l_1^{-n}, \quad f' = l_0^n f \quad . $$

NB. Grâce à la formule des lacets, une fois fixé le multiplicateur $\mu$ i.e. $\nu$ (et en fait $\mu$ est déjà déterminé de façon unique par la seule connaissance de $\xi$, comme le montrent les calculs abéliques dans $\operatorname{Gl}(2,\hat{\mathbb{Z}})$ ...), $\xi$ et $\eta$ se déterminent mutuellement, par (134), donc on peut éliminer l'une ou l'autre des [unclear: systèmes] d'équations (133)(134), qui devient soit

$$ (145) \quad \xi \rho(\xi) \rho^2(\xi) = 1 \quad , \quad \xi l_1^{-\nu} \sigma_\infty(\xi l_1^{-\nu}) = 1 \quad i.e. $$
$$ \sigma_\infty(\xi) = l_1^\nu \xi^{-1} l_0^\nu $$

soit

$$ (146) \quad \eta \sigma_\infty(\eta) = 1 \quad , \quad (\eta l_1^\nu) \rho(\eta l_1^\nu) \rho^2(\eta l_1^\nu) = 1 $$

\newpage

%% ===== Page 218 =====

Le calcul adélique. On désigne encore par $\xi, \eta$ les
images de $\xi, \eta$ dans $\operatorname{Sl}(2, \hat{\mathbb{Z}})'$, qui se remontent en
des éléments, notés encore $\xi, \eta$, de $\operatorname{Sl}(2, \hat{\mathbb{Z}})$.
Il y a une ambiguïté de signe, qu'on lève pour
$\xi$ en imposant

$$ (147) \qquad \xi \rho(\xi)\rho^2(\xi) = \underline{\underline{+1}} \qquad \text{(dans } \operatorname{Sl}(2, \hat{\mathbb{Z}}) \text{)} $$

[MARGIN: mais je n'ai pas fait le choix de ces signes [unclear] mais alors y a peut-être des contradictions ! [unclear] conventions [unclear] par ailleurs [unclear]]

il n'y a pas de façon simple de lever l'ambi-
guïté, sauf en admettant la possibilité d'écrire
sous la forme $\beta \varpi_{\infty}(?)$ dont on aura
à priori seulement

$$ (148) \qquad \eta \varpi_{\infty}(\eta) = \pm 1 , $$

sans possibilité de modifier le signe par un
changement $\eta \mapsto -\eta$ ! Posons

$$ (149) \qquad \xi = \begin{pmatrix} A & B \\ C & D \end{pmatrix} , \quad \eta = \begin{pmatrix} R & S \\ T & U \end{pmatrix} $$

On aura

$$ (149') \qquad AD - BC = 1 \qquad RU - ST = 1 $$

et la relation (147) équivaut alors à

$$ (150) \qquad C + D - B = 1 \quad , $$

tandis (du moins dans chaque composante $l$-adique
on devra avoir soit cette relation, soit

$$ \xi_l = \begin{pmatrix} A_l & B_l \\ C_l & D_l \end{pmatrix} = -\rho^{-1} = \begin{pmatrix} 1 & -1 \\ 1 & 0 \end{pmatrix} , $$

cf. prop. de la p. 399 — mais cette solution
est incompatible avec la condition $\xi \equiv 1 \pmod{2}$ ,
d'où $A, D \equiv 1 \pmod{2}$ et $B, C \equiv 0 \pmod{2}$ ... ) dans $\hat{\mathbb{Z}}$
(150) est vraie dans $\mathbb{Z}_l$ pour tout $l$, donc dans $\hat{\mathbb{Z}}$ ... ).

\newpage

%% ===== Page 219 =====

409

La relation (148) s'écrit alors

$$y = \pm \sigma_{\infty}(y^{-1})$$

or $\sigma_{\infty}(y^{-1}) = {}^t y$ , donc (148) s'écrit

$$y = \pm {}^t y \ ,$$

mais si on avait le signe $-$ , on aurait $U = -U$

donc $2U=0$ donc $U=0$, ce qui est incompatible

avec

$$y \equiv 1 \quad (2)$$

qui implique $R, U \equiv 1 \ (2)$, $S, T \equiv 0 \ (2)$. On

aura donc $y = {}^t y$ i.e.

(151) $$y \sigma_{\infty}(y) = +1 \ , \quad i.e.$$

(152) $$S = T$$

La relation [unclear: du cocycle] (138) s'écrit

$$ \begin{pmatrix} A & B \\ C & D \end{pmatrix} = \begin{pmatrix} 1 & 0 \\ 2\nu & 1 \end{pmatrix} \varepsilon \begin{pmatrix} R & S \\ T & U \end{pmatrix} \begin{pmatrix} 1 & 0 \\ 2\nu & 1 \end{pmatrix} $$

$$ = \varepsilon \begin{pmatrix} R+2\nu S & S \\ T+2\nu U & U \end{pmatrix} \qquad \text{avec } \varepsilon \in \{\pm 1\} $$

[MARGIN: Obs ici $\varepsilon = 1$, en choisissant la bonne convention de (150). (La convention d'un coïncider par ex. avec celle ... mais donne [unclear] asymétrique avec les $\alpha, \beta$, elle parait plus symétrique suivant ...)]

et elle équivaut aux relations

(154) $$B = \varepsilon S , \quad D = \varepsilon U \qquad \text{et}$$

(155) $$\mu D = 1 \qquad \text{où } \mu = 2\nu + 1$$

équivalente aussi au vu de (154), à

(155 bis) $$\mu U = \varepsilon \ , \qquad \text{(compte tenu de (152))}$$

et cette dernière au fait qui exprimerait la

relation $$ \overset{A'}{\varepsilon(T+2\nu U)} + \overset{B'}{\varepsilon U} - \overset{C'}{\varepsilon S} = 1 $$ paraphasant

$$ A + B - C = 1 \ , \quad \text{elle ...}$$

\newpage

%% ===== Page 220 =====

elle ne fait qu'exprimer (compte tenu
de (154) et de (150)) la relation
$$C = \varepsilon (\overset{C'}{T+2\nu U})$$

Ainsi les relations (154) (155') expriment le
fait que les deux matrices de (153)
ont trois coefficients égaux à $\varepsilon$ près
$$B = B', C = C', D = D'$$
et l'égalité pour le dernier
$$A = \varepsilon (\overset{A'}{R+2\nu S})$$
s'en déduit alors grâce aux relations (144),
qui donnent d'ailleurs
$$AD-BC=1, \quad A'D'-B'C'=1$$
qui donnent
$$A = \frac{1+BC}{D} \quad , \quad A' = \frac{1+B'C'}{D'}$$
(NB on sait que $D=D'$ est inversible, par (155'))
donc $A=A'$.

On trouve ainsi le
Théorème. Pour une solution $(\xi, \eta)$ dans $\hat{\Pi}_{0,3} \times \hat{\Pi}_{0,3}$
du système d'équations (133)(134), désignant
encore par $\xi, \eta$ des éléments de $\operatorname{Sl}(2, \hat{\mathbb{Z}})$ qui
leur correspondent, l'ambiguité de signe pour $\eta$
étant levée par la relation (147), on aura
les relations (154) et (155'), qui impliquent
les conséquences suivantes :

a) $\mu = 2\nu+1$ est inversible, et il est entièrement déterminé
par $\xi$ : il est déterminé par $\eta$ seulement au
signe près.

\newpage

%% ===== Page 221 =====

Remarques
410

Compte tenu des relations (149) (150)
$$AD - BC = 1 \quad , \quad A + D - B = 1 \quad ,$$
et du fait que $D$ est inversible,
les matrices $\xi$ sont entièrement déterminées par
les coeff $B, D$ — les couples $(\xi,\eta)$
de matrices $\in \operatorname{Sl}(2, \hat{\mathbb{Z}})$, satisfaisant

(156) $\begin{cases} \{ r(\xi,\eta) = 1, s(\xi,\eta) = 1, \eta = \xi l_1^{-\nu} \quad (\nu \in \hat{\mathbb{Z}} \text{ convenable}) \\ \xi \equiv 1 (2), \quad \eta \equiv 1 (2) \end{cases}$

[MARGIN: Bien sûr, l'introduction de $\nu$ congru [...] ici est tout à fait gratuit en [...] nécessaire - éliminé algébriquement [...] indépendante de telle condition...]

sont en correspondance biunivoque avec les
$(B, D) \in \hat{\mathbb{Z}} \times \hat{\mathbb{Z}}^*$ (dont on déduit $B \equiv 0 (2)$)
$\xi = \begin{pmatrix} A & B \\ C & D \end{pmatrix}, \eta = \begin{pmatrix} R & S \\ T & U \end{pmatrix}$, et $\nu \in \hat{\mathbb{Z}}$ par

(157) $C = B - D + 1 \quad , \quad A = \frac{1+BC}{D} = \frac{1}{D}(1+B^2-BD+B)$
$= \mu(1+B+B^2) - B$

(158) $\mu = \frac{1}{D} \quad , \quad \nu = \frac{\mu - 1}{2} \quad ,$

[MARGIN: on a oublié de signaler $\xi \equiv \begin{pmatrix} 1 & 0 \\ 0 & 1 \end{pmatrix}$ (2) dans (156)]

(159) $\eta = \xi l_1^{-\nu} = \begin{pmatrix} A & B \\ C & D \end{pmatrix} \begin{pmatrix} 1 & 0 \\ -2\nu & 1 \end{pmatrix} = \begin{pmatrix} A - 2\nu B & B \\ C - 2\nu D & D \end{pmatrix} \quad c.q.f.d.$

(159) inc. $U = D, S = T = B, R = A - 2\nu B = \mu(1+B^2) = \frac{1+B^2}{D}$

[MARGIN: N.B. On déduit $(\xi,\eta)$ de $u_\xi$ par le
fait que $\mu = \det u_\xi$,
$\nu = \frac{\mu-1}{2}$,
$\xi = \dots$
(161) $\eta = \dots$]

Posons (comme p. 338' (78))
(160) $u_\xi = \mu \begin{pmatrix} 1 & B \\ 0 & D \end{pmatrix} = \begin{pmatrix} \mu & \mu B \\ 0 & 1 \end{pmatrix} = \begin{pmatrix} \frac{1}{D} & B/D \\ 0 & 1 \end{pmatrix} \in \operatorname{Gl}(2, \hat{\mathbb{Z}})$
de sorte que la connaissance de $(B, D)$, i.e. celle de $\xi$ équivaut à celle de $u_\xi \in \operatorname{Gl}(2, \hat{\mathbb{Z}})$, ou encore
à celle des couples $(\xi, \eta)$ (y assimilés par $\eta = \xi l_1^{-\nu}$) . (Dans loc. cit. cet élément était noté $u_{\alpha\beta}$... mais il ne dépend que de $\xi \in \Pi_{0,\hat{3}}$ ,
et le sous-groupe des $\xi \in \operatorname{Sl}(2, \hat{\mathbb{Z}})$ - où sa
connaissance équivaut à celle de $\xi \in \operatorname{Sl}(2, \hat{\mathbb{Z}})$ - )

\newpage

%% ===== Page 222 =====

[MARGIN: Les formules (163) s'écrivent aussi (163 bis)
$(u_\xi \sigma_\infty) = \eta$, $(u_\xi \rho) = \xi \rho$
elles découlent de la façon dont on a calculé
$\xi$ à l'aide de $\eta, \underline{u}_\xi$, i.e. $\xi = \underline{u}_\xi \dots$]

[unclear: Explicitions on aura]

(163) $$u_\xi (\sigma_\infty) = \eta \sigma_\infty = (\xi \rho^{-2 \nu}) \sigma_\infty \quad , \quad u_\xi (\rho) = \xi \rho$$
(164) $$u_\xi (\varepsilon_0) = \varepsilon_0^\mu \quad \text{et} \quad u_\xi (l_0) = l_0^\mu$$

Si on considère une matrice inversible
quelconque
$$\underline{u} = \left(\begin{array}{cc} L & M \\ N & P \end{array}\right) \in \operatorname{Gl}(2, \hat{\mathbb{Z}})$$

la relation
(*) $\underline{u} (\varepsilon_0) = \varepsilon_0^\mu \quad i.e. \quad \underline{u} \varepsilon_0 = \varepsilon_0^\mu \underline{u} \quad (\mu \in \hat{\mathbb{Z}}^*)$
s'écrit équivaut à
$$\left(\begin{array}{cc} L & L+M \\ N & N+P \end{array}\right) = \left(\begin{array}{cc} L+\mu N & M+\mu P \\ N & P \end{array}\right)$$

est équiv. à :
$$N=0 \quad, \quad L = \mu P$$

i.e.
$$\underline{u} = \left(\begin{array}{cc} \mu P & M \\ 0 & P \end{array}\right) \quad , \quad \text{d'où } \det \underline{u} = \mu P^2$$

la condition supplémentaire
(**) $$\det \underline{u} = \mu$$
donne la condition
$$P^2 = 1 \quad ,$$
la solution la plus évidente est $P=1$ ,
ce qui donne juste les matrices
$$\underline{u} = \left(\begin{array}{cc} \mu & M \\ 0 & 1 \end{array}\right) = \underline{u}_\xi$$

\newpage

%% ===== Page 223 =====

où $\xi = \left( \begin{matrix} A & B \\ C & D \end{matrix} \right)$ satisfaisant [unclear] est donné en termes de
$B, D$ par $D = \frac{1}{\mu}, B = DM = M/\mu$.

Théorème L'application

[MARGIN: NB Du fait qu'ils sont naturels ! (168) Dans $! \sim (0)$, on trouve $M_{0,3} \hookrightarrow \operatorname{Gl}(2, \hat{\mathbb{Z}})/\pm 1$ un hom dans $\operatorname{Gl}(2, \hat{\mathbb{Z}})$, pas seulement dans $\operatorname{Gl}(2, \hat{\mathbb{Z}})/\pm 1$]
[MARGIN: NB Cette relation [unclear] est la rel. $B$ ci-dessous (169)]

$$M_{0,3}^{\sim \hat{}} (0) \longrightarrow \operatorname{Gl}(2, \hat{\mathbb{Z}})$$
$$u_{\xi, \eta} \longmapsto \underline{u}_{\xi} \quad \text{, notation où } u \mapsto \underline{u}$$

($\underline{u}_{\xi}$ est donné par (160)) est un hom. de groupes.

Dém
C'est-à-dire $u = u_{\xi, \eta}$, $u' = u_{\xi', \eta'}$, on aura pour $u'' = u_{\xi'', \eta''} = u' u$, $\underline{u} \underline{u}' = \underline{u}''$,
Si $\underline{u}$ image dans $\operatorname{Gl}(2, \hat{\mathbb{Z}})$, et si $x \in \hat{\Pi}_{0,3}$, $\underline{x}$
son image dans $\operatorname{Sl}(2, \hat{\mathbb{Z}})$ :
on a dans $\underline{u}$
(166) $\underline{u(x)} = \underline{u} (\underline{x}) ( \overset{dfn}{=} \underline{u} \underline{x} \underline{u}^{-1} )$

Car il suffit de le voir pour $x = \sigma_{\infty}$ et pour $x = \rho$,
où cette formule se réduit à (161). Ceci noté,
si $u = u_{\xi, \eta}$, $u' = u_{\xi', \eta'}$ deux éléments de
$M_{0,3}^{\sim \hat{}}(0)$, $u'' = u_{\xi'', \eta''} = u' u$ leur produit, on aura (159)

d'où
(167) $\xi'' = u'(\xi) \xi' = \underline{u}'(\xi) \xi'$ , i.e.
(167) $\underline{\xi}'' = int (\underline{u}_{\xi'}) (\underline{\xi}) \cdot \underline{\xi}'$
et il suffit de vérifier que si trois
solutions $\xi, \xi', \xi''$ des équations

$$\xi = \left( \begin{matrix} A & B \\ C & D \end{matrix} \right), \xi' = \left( \begin{matrix} A' & B' \\ C' & D' \end{matrix} \right), \xi'' = \left( \begin{matrix} A'' & B'' \\ C'' & D'' \end{matrix} \right)$$

des équations $AD-BC=1, C+D \cdot B=1, D \in \hat{\mathbb{Z}}^*$,
sont reliées par (167), alors on a
(168) $\underline{u}_{\xi''} = \underline{u}_{\xi'} \cdot \underline{u}_{\xi}$

\newpage

%% ===== Page 224 =====

Comme $\xi$ et $\underline{u} = u_\xi$ se déterminent mutuellement ([unclear]), il revient au m. de dire que si, sur le groupe
$$ \mathcal{B} \subset \operatorname{Gl}(2, \hat{\mathbb{Z}}) $$
$$(169) \quad \{ \begin{pmatrix} \mu & M \\ 0 & 1 \end{pmatrix} \mid \mu, M \in \hat{\mathbb{Z}}, \mu \text{ inv.} \}$$
on définit (cf (162))
$$(170) \quad \left\{ \begin{aligned} & F: \mathcal{B} \longrightarrow \mathfrak{S}l(2, \hat{\mathbb{Z}}) \\ & F_\xi (\underline{u}) = (\det \underline{u})^{-1} \xi [\underline{u}, \tilde{\rho}] \end{aligned} \right.$$
alors on a la relation des 1-cocycles
$$(171) \quad F_\xi ( \underline{u}' \underline{u} ) = \text{int} (\underline{u}') (F_\xi (\underline{u}) ) \cdot F_\xi (\underline{u}') .$$
Par multiplicativité de $\det$, ca se ramène à voir que
$$ \underline{u} \longmapsto [\underline{u}, \tilde{\rho}] = \underline{u} \tilde{\rho} (\underline{u})^{-1} $$
satisfait la relation des cocycles, ce qui est trivial.

Corollaire. L'application
$$ \mathcal{M}_{0,3} \simeq \mathcal{M}_{0,3}^{\sim}[0] \cdot \widehat{\mathfrak{S}_3} \longrightarrow \operatorname{Gl}(2, \hat{\mathbb{Z}})^\# $$
(où le point désigne le produit 1/2 direct)
$$ u_\xi \cdot a \longmapsto \underline{u}_\xi \cdot \underline{a} $$
(où $\underline{u}_\xi$ est la réduction de $u_\xi \bmod \pm 1$, et où $\underline{a}$ désigne l'image de $a$ dans $(\operatorname{Gl}(2, \hat{\mathbb{Z}})^\#)$ ) est
un homomorphisme de groupes.
Résulte aussitôt du th. et des formules (163).

[MARGIN: attention au fait que $\mathcal{M}_{0,3}$ ne se déduit pas de $\mathcal{M}_{0,3}^{\sim}[0]$ par un produit $1/2$ direct! C'est bien $\mathcal{M}_{0,3}^{\sim}[0]$ le bon [unclear]. Faut diviser $\underline{u}_\xi$ par $\pm 1$ pour s'en tirer! Mais ces remarques sont oiseuses, voir pages suivantes!]

\newpage

%% ===== Page 225 =====

470

Dorénavant pour repérer le couple $(\xi,\eta)$ d'un
point de vue adélique (dans $\operatorname{Gl}(2, \hat{\mathbb{Z}})$) à l'aide
de $\underline{u}$, qui est de la forme

(173) $$ \underline{u} = \left( \begin{array}{cc} \mu & M \\ 0 & 1 \end{array} \right) $$ avec $M \equiv 0 (2)$ si $\underline{u}$ provient de $\mathcal{M}_{0,3}^{1\sim}(0)$ (163 bis),

dont on tire $\xi, \eta$ par les formules (163 bis),

(174) $\xi = [\underline{u}, \rho]$, $\eta = [\underline{u}, \sigma_\infty]$ .

Pour $u \in \mathcal{M}_{0,3}^{1\sim}(\infty)$, défini par $(\xi, \eta)$, où on
ne peut plus retenir la possibilité d'écrire
$\xi, \eta$ de façon unique comme $\xi = \alpha \rho (\alpha)^{-1}$, $\eta = \beta \sigma_\infty (\beta)^{-1}$
avec $\alpha \in \hat{\Pi}_{0,3}^\sim$, $\beta \in \hat{\Pi}_{0,3}^\sim$, le seul invariant
adélique qu'on puisse lui associer directement
de façon naturelle est le couple $(\xi, \eta)$,
ou encore (ce qui revient au même, mais est
plus joli) $\underline{u}$. Les éléments
$\underline{\lambda} = l_0^{-n} \in L_0 \hookrightarrow \mathcal{M}_{0,3}^{1\sim}[\vec{\infty}]$
sont repérés par :

[unclear: crossed out] $\xi = \eta = l_0^{-n} l_1^{-n} = \left(\begin{array}{cc} 1 & -2n \\ 0 & 1 \end{array}\right) \left(\begin{array}{cc} 1 & 0 \\ -2n & 1 \end{array}\right) = \left(\begin{array}{cc} A=1+4n^2 & B=-2n \\ C=-2n & D=1 \end{array}\right)$

[MARGIN: Précisons en deça des justifications qu'on tire de la sémantique inhérente au point de vue "adélique", dont il s'agit de se méfier cf point douteux dans (175)]

(175) $$ \underline{\lambda} = \underline{l}_0^{-n} = \left( \begin{array}{cc} 1 & -2n \\ 0 & 1 \end{array} \right) $$

d'ailleurs on a

(177) $$ \left( \begin{array}{cc} 1 & -2n \\ 0 & 1 \end{array} \right) \left( \begin{array}{cc} \mu & M \\ 0 & 1 \end{array} \right) = \left( \begin{array}{cc} \mu & M-2n \\ 0 & 1 \end{array} \right) $$
avec $\underline{l}_0^{-n}$ sous la première matrice et $\underline{u}$ sous la deuxième.

ce qui montre que la description des
sous-groupes $B_1$ de $B$ (169) fournit des

\newpage

%% ===== Page 226 =====

matrices $\underline{u}$ telles que $u \equiv 1$ (2) i.e. $M \equiv 0$ (2)
pour le ss-groupe engendré par $l_0 = \begin{pmatrix} 1 & -2 \\ 0 & 1 \end{pmatrix}$
s'identifie au groupe des matrices $\begin{pmatrix} \mu & 0 \\ 0 & 1 \end{pmatrix}$
un tore maximal $T_0$ de $B_1$, de façon plus
précise on a une structure de produit semi-direct

(178) $$B_1 \simeq T_0 \cdot L_0$$
$$ \left\{ \begin{pmatrix} \mu & M \\ 0 & 1 \end{pmatrix} \middle| \begin{matrix} \mu \in \hat{\mathbb{Z}}^* \\ M \in 2\hat{\mathbb{Z}} \end{matrix} \right\} \quad \left\{ \begin{pmatrix} \mu & 0 \\ 0 & 1 \end{pmatrix} \middle| \mu \in \hat{\mathbb{Z}}^* \right\} \quad \left\{ l_0^n \mid n \in \hat{\mathbb{Z}} \right\} $$
$$ \begin{pmatrix} 1 & -2n \\ 0 & 1 \end{pmatrix} $$

Cela montre que si on s'intéresse aux
quotients $\mathcal{M}_{0,3}^{1\sim} (0) / L_0 \simeq \mathcal{N}_{0,3}^{1\sim}$ , alors
le seul invariant adélique qu'on saura lui
associer dans le contexte présent (quand on
ne travaille pas sur $\alpha, \beta$) est le multiplica-
teur $\mu$. Mais il y a mieux, on trouve grâce
à la décomposition (178) dans $\operatorname{Gl}(2, \hat{\mathbb{Z}})$
un scindage canonique de l'extension
$\mathcal{M}_{0,3}^{1\sim} (0)$ de $\mathcal{N}_{0,3}^{1\sim}$ par $L_0$ :

Proposition. Pour tout $\dot{u} \in \mathcal{N}_{0,3}^{1\sim}$, il existe
un représentant unique $u \in \mathcal{M}_{0,3}^{1\sim} (0)$ tel
qu'on ait $\underline{u} \in T_0$ i.e. $\underline{u}$ de la forme $\begin{pmatrix} \mu & 0 \\ 0 & 1 \end{pmatrix}$

[MARGIN: si $\dot{u} = \overline{\begin{pmatrix} A & B \\ C & D \end{pmatrix}}$ dans $\mathcal{N}_{...}$, $\underline{u}$ de la forme [unclear: $\begin{pmatrix} A & 0 \\ C & D \end{pmatrix} = \begin{pmatrix} \mu & 0 \\ 0 & 1 \end{pmatrix}$] ]

\newpage

%% ===== Page 227 =====

413

$$ (179) \qquad \mathcal{M}_{0,3}^{\sim} (0) \overset{df}{=} \{ u \in \hat{\mathcal{M}}_{0,3}^{\sim} (0) \mid \underline{u} \in T_0 \text{ i.e. } \underline{u} = \begin{pmatrix} \mu & 0 \\ 0 & 1 \end{pmatrix} \} $$

est un ss-groupe section de l'extension $\hat{\mathcal{M}}_{0,3}^{\sim}(0)$ de $\mathcal{N}_{0,3}^{\sim}$ par $L_0$, donc aussi un scindage canonique $\hat{\mathcal{M}}_{0,3}^{\sim}$ de $\mathcal{N}_{0,3}^{\sim}$ par $\hat{\Pi}_{0,3}$, et de l'extension $\tilde{\mathcal{M}}_{0,3}$ de $\tilde{\mathcal{N}}_{0,3}$ par $\tilde{\mathcal{B}}_{0,3}$ déduite.

Cela provient du fait que si $u \in \hat{\mathcal{M}}_{0,3}^{\sim}(0)$, $\underline{u} = \begin{pmatrix} \mu & M \\ 0 & 1 \end{pmatrix}$, avec $M \equiv 0 (2)$, i.e. $M \in 2\hat{\mathbb{Z}}$, il existe un unique $n \in \hat{\mathbb{Z}}$ tel que $l_0^n \underline{u} = \begin{pmatrix} \mu & M-2n \\ 0 & 1 \end{pmatrix}$ soit dans $T_0$ i.e. tel que $l_0^n u \in \mathcal{M}_{0,3}^{\sim}(0)$, savoir $n = \frac{M}{2}$.

Si $(\xi, \eta)$ satisfont aux conditions (133)(134) i.e. définissent un $u = u_{\xi, \eta}$ (cf th. p. 407'), lors on dira que $(\xi, \eta)$ est normalisé si $u_{\xi, \eta} \in \mathcal{M}_{0,3}^{\sim}(0)$. De même, si un élément quelconque de $\hat{\mathcal{M}}_{0,3}^{\sim}$ est écrit dans la forme 
$$ u = f^{-1} \cdot u_{\xi, \eta} = f^{-1} u_{\xi} $$
[MARGIN: (en vertu de l'identification conventionnelle 136)] 
on dit que cette représentation (par $\{\xi, \eta, f \in \hat{\pi}_{0,3} \}$) si $u_{\xi, \eta}$ est normalisé [ i.e. pour $\underline{u}_{\xi, \eta} = \begin{pmatrix} A & B \\ C & D \end{pmatrix}$, on a $B=0$ ]. Ainsi, tout $u$ de $\hat{\mathcal{M}}$ se met de façon unique sous forme

[unclear: (177)] $\quad u = f^{-1} u_{\xi, \eta}$ ,

\newpage

%% ===== Page 228 =====

avec $f$ normalisé. En fait, ceci correspond
à la description en produit $1/2$-direct
$$(180) \quad \mathcal{M}_{0,3}^{!\sim} \simeq \mathcal{M}_{0,3}^{!\sim}(0) \cdot \widehat{\Pi}_{0,3}$$
— et d'ailleurs de même
$$(181) \quad \mathcal{M}_{0,3}^{\sim} \simeq \mathcal{M}_{0,3}^{!\sim}(0) \cdot \widehat{\mathfrak{S}}_{0,3}^{+}$$

[unclear: Rem] Il y aurait à l'occasion moyen de mettre
ces scindages canoniques en relation avec le
choix d'une uniformisante en $0$ de $\mathbb{P}_{\mathbb{Q}}^1 \simeq \mathcal{U}_{0,3\mathbb{Q}}$.
Bien entendu, la même question se pose exac-
tement de même pour les "points à l'infini" $1, \infty$ de $\mathcal{U}_{0,3\mathbb{Q}}$,
relativement aux sous-groupes d'inertie
$$(182) \quad \mathcal{M}_{0,3}^{!\sim}(q) \quad q \in \{1,\infty\}$$
correspondant à ces points à l'infini -
la construction d'une [unclear: quasi-section]
$\mathcal{M}_{0,3}^{!\sim}(0)$ dans le normalisateur de $L_0$,
s'applique évidemment trait pour trait pour $L_1, L_\infty$.
On aura de même
$$(183) \quad \rho(\mathcal{M}_{0,3}^{!\sim}(q)) = \mathcal{M}_{0,3}^{!\sim}(\rho(q)) \quad \text{pour } q \in \{0, 1, \infty\}$$
donc les trois $\mathcal{M}_{0,3}^{!\sim}(q)$ se déduisent de l'un
quelconque d'entre eux en appliquant $\rho, \rho^2$.

\newpage

%% ===== Page 229 =====

47. Retour sur l'homomorphisme $\Pi_{0,3}^\wedge \xrightarrow{\hat{\varphi}} \widehat{G}_{0,3}^+$
et la condition de commutation extérieure de celui-ci. Le groupe $\mathcal{M}_{0,3}^\sim$, $\mathcal{M}_{0,3}^{!\sim}$ et l'hom. $\varphi : \mathcal{M}_{0,3}^\sim \to \mathcal{M}_{0,3}^{!\sim}$.

Rappelons que $\hat{\varphi}$ est décrit par
(1) $$ \hat{\varphi}(l_0) = \varepsilon_0 = \sigma_\infty \rho \ , \ \hat{\varphi}(h) = \sigma_\infty \ , \ \hat{\varphi}(l_\infty) = \rho^{-1} \ . $$

Soit
$$ u = u_{\alpha, \beta, f} \in \mathcal{M}_{0,3}^{!\sim} $$

un autom à lacets de $\Pi_{0,3}^\wedge$ , s'étendant ext$^t$ à $\widehat{G}_{0,3}$ , défini par

(2) $$ u(\sigma_\infty) = h \sigma_\infty = int(\beta_0) \sigma_\infty \quad (h = \beta_0 \sigma_\infty \beta_0^{-1} , \ \beta_0 = f^{-1} \beta) $$

(3) $$ u(\rho) = g \rho = int(\alpha_0) \rho \quad (g = \alpha_0 \rho (\alpha_0)^{-1} , \ \alpha_0 = f^{-1} \alpha) $$

et écrivons que $u$ commute ext$^t$ avec $\hat{\varphi}$ , i.e. qu'il existe

(4) $$ \lambda \in \Pi_{0,3}^\wedge $$

tel que l'on ait

(5) $$ \hat{\varphi}(u(x)) = int(\lambda) u \hat{\varphi}(x) \qquad \forall x \in \Pi_{0,3}^\wedge \ . $$

Il suffit de vérifier (5) pour $x = l_0$ et $x = l_1$,

1) Cas $x = l_0$ , d'où
[unclear: crossed out text, possibly $u(x) = int(f^{-1})(l_0^\mu)$]
$$ \hat{\varphi}(u(x)) = int(\hat{\varphi}(f)^{-1}) \hat{\varphi}(l_0^\mu) = int(\hat{\varphi}(f)^{-1}) \varepsilon_0^\mu $$
$$ u \hat{\varphi}(x) = u(\varepsilon_0) = int(f^{-1}) (\varepsilon_0^\mu) $$
donc la relation (5) s'écrit
$$ int( \lambda \hat{\varphi}(f) f^{-1} ) (\varepsilon_0^\mu) = \varepsilon_0^\mu $$

\newpage

%% ===== Page 230 =====

[MARGIN: NB on n'a point vu en detail [unclear] $\hat{\varphi}(f) \lambda f^{-1} \in \hat{\Pi}_{0,3}$ i.e. que c'est [unclear] $\in \hat{\Pi}_{0,3}$ i.e.]

équivaut à $\hat{\varphi}(f) \lambda f^{-1} \in L_0^\wedge = Zentr_{\hat{\Pi}_{0,3}}(\varepsilon_0^M)$
: l'existence de $r \in \hat{\mathbb{Z}}$ tel que l'on ait
$\hat{\varphi}(f) \lambda f^{-1} = \varepsilon_0^r$, i.e.

(6) $\lambda = \hat{\varphi}(f)^{-1} \varepsilon_0^r f$ .

Ceci permet donc d'éliminer $\lambda$ de la deuxième relation :

2) Cas $x = l_1$, d'où
$u(x) = int(f_1^{-1})(l_1^M) \quad (f_1 = \text{[unclear]} p(f^{-1}))$
$\hat{\varphi}(u(x)) = int(\hat{\varphi}(f_1)^{-1}) \sigma_\infty^M = int(\hat{\varphi}(f_1)^{-1}) \sigma_\infty$
(compte tenu que $\sigma_\infty^M = \sigma_\infty$ car $M \equiv 1 \pmod 2$).

[unclear crossed out line]
$u \hat{\varphi}(x) = u(\sigma_\infty) = int(\beta_0)(\sigma_\infty)$ ,

Donc la relation (5) s'écrit
$int((\hat{\varphi}(f_1))^{-1} \lambda \beta_0) \sigma_\infty = \sigma_\infty$
ce qui signifie qu'il existe $s \in \hat{\mathbb{Z}}/2\hat{\mathbb{Z}}$ tel qu'on ait
$\hat{\varphi}(f_1)^{-1} \lambda \beta_0 = \sigma_\infty^s \quad \text{soit} \quad \hat{\varphi}(f_1)^{-1} = \lambda \beta_0 \sigma_\infty^{-s}$

et remplaçant $\lambda$ par sa valeur (6), et $f_1^{-1}$ par sa valeur $g p(f)^{-1}$, on trouve la condition

[unclear crossed out line]
$\hat{\varphi}(g p(f)^{-1})^{-1} = \hat{\varphi}(f)^{-1} \varepsilon_0^r f \beta_0 \sigma_\infty^s \quad \text{soit}$

$\hat{\varphi}(f g p f^{-1}) = \varepsilon_0^r \beta_0 \sigma_\infty^{-\alpha}$
$\xi = \alpha p(\alpha)^{-1}$

\newpage

%% ===== Page 231 =====

4d5

Ainsi :
Théorème Pour que
$$ \hat{\phi} : \Pi_{0, \tilde{3}} \longrightarrow \hat{S}_{0, 3}^+ $$
commute extérieurement à l'élément
$u \in \mathcal{M}_{0, \tilde{3}}^\sim$, correspondant aux invariants
$g, h, f, \alpha, \beta, \alpha, \beta_0, \xi, \eta$ (cf §4c), il f. et s. qu'on
ait

(7) $$ \hat{\phi}(\xi) = \xi_0^r \beta \sigma_\infty^\alpha $$
pour $r \in \hat{\mathbb{Z}}$
$\alpha \in \mathbb{Z}/2\mathbb{Z}$
convenables.

[MARGIN: NB $r, \alpha$ sont bien déterminés par [unclear: $\hat{\phi}, \alpha$ et $\beta_0$]. Car aucun $\xi_0^n$ n'est divisible.]

Notons que $\lambda$ donne par (6) : $\lambda = \hat{\phi}(f)^{-1} \xi_0^r f$.

Notons que $u$ ne définit le triple
$(\alpha, \beta, f)$ que modulo une transformation
de la forme

(8) $$ (\alpha, \beta, f) \longmapsto (\alpha' = l_0^n \alpha, \beta' = l_0^n \beta, f' = l_0^{-n} f) $$
où $n \in \hat{\mathbb{Z}}$. On aura alors (pour $\xi' = \alpha' \beta (\alpha')^{-1} = l_0^n \xi l_0^{-n}$)

*) $$ \hat{\phi}(\xi') = \xi_0^{r'} \beta' \sigma_\infty^{\alpha'} $$
avec $r' \in \hat{\mathbb{Z}}$, $\alpha' \in \mathbb{Z}/2\mathbb{Z}$ également bien déterminés.

Or on aura
$$ \hat{\phi}(\xi') = \hat{\phi}(l_0^n \xi l_0^{-n}) = \hat{\phi}(l_0^n) \hat{\phi}(\xi) \hat{\phi}(l_0^{-n}) = \varepsilon_0^n \hat{\phi}(\xi) \varepsilon_0^{-n} $$

et la relation (*) s'écrit
$$ \hat{\phi}(\xi) = \xi_0^{r'-2n} \beta' \sigma_\infty^{\alpha'-n} $$

où encore, comme $\beta' = l_0^n \beta = \varepsilon_0^{2n} \beta$
$$ \hat{\phi}(\xi) = \xi_0^{r'+n} \beta \sigma_\infty^{\alpha'-n} \ , $$

ce qui, comparé à (7), donne

(9) $$ r' = r - n \ , \ \alpha' = \alpha + n \mod 2 $$

\newpage

%% ===== Page 232 =====

Cela implique le

Corollaire 1 Supposons que $u$ commute exté-
rieurement à $\hat{\phi} : \Pi_{0,3} \longrightarrow \widehat{\mathcal{G}}_{0,3}^+$. Alors
il est possible de façon unique de
décrire $u$ par un triple $(\alpha, \beta, f)$
(a priori déterminé modulo une transfor-
mation du type (8)), de façon à ce qu'on
ait (posant $\xi = \alpha \rho(\alpha)^{-1}$, comme à l'accou-
tumée) la relation

[MARGIN: Il faudrait examiner si cette définition coïncide avec celles de la p. 40', il faud. formuler [unclear] un tel choix sera appelé normalisé.]
(10) $$ \hat{\phi}(\xi) = \beta \sigma_\infty^i $$ pour $i \in \Z/3\Z$
convenable.

Ainsi les $u \in \mathcal{M}_{0,3}^\sim$ qui commutent
extérieurement à $\hat{\phi}$ correspondent
biunivoquement aux simples triples [unclear: crossed out]
$(\alpha, \beta, f)$, $\alpha \in \Pi_{0,3}^\sim$, $\beta, f \in \Pi_{0,3}$
satisfaisant d'une part la relation des
lacets
(11) $$ \xi = \eta h_1^r $$ où $\xi = \alpha \rho(\alpha)^{-1}$, $\eta = \beta \sigma_0(\beta)^{-1}$,
et $r \in \hat{\Z}$ convenable
(tel que $2r+1 \in \hat{\Z}^*$)
d'autre part la relation (10) (pour $i \in \Z/3\Z$ conv.)
— l'une ni l'autre relation ne faisant
d'ailleurs pas intervenir $f$.

\newpage

%% ===== Page 233 =====

Notons que $i \in \mathbb{Z}/2\mathbb{Z}$ est défini en termes de
$\xi$ compte tenu de $\beta \in \Pi_{0,\hat{3}} = \operatorname{Ker}(\widehat{\mathfrak{S}}_{0,3}^+ \to \mathfrak{S}_3)$
[unclear: en vertu de] (10) mod $\Pi_{0,\hat{3}}$ . Désignons par

(12) $\hat{\varphi}^{\cdot} : \Pi_{0,\hat{3}} \to \mathfrak{S}_3$

la composée $\Pi_{0,\hat{3}} \overset{\hat{\varphi}}{\to} \widehat{\mathfrak{S}}_{0,3}^+ \to \mathfrak{S}_3$ ,

la formule (10) donne donc

(13) $\hat{\varphi}^{\cdot}(\xi) = \sigma_{\infty}^i$

ce qui détermine bien $i \in \mathbb{Z}/2\mathbb{Z}$ en termes de
$\xi \in \Pi_{0,\hat{3}}$, et implique une condition
intrinsèque sur l'élément $\xi$ de $\Pi_{0,\hat{3}}$ ,
à savoir que son image dans $\mathfrak{S}_3$ par $\hat{\varphi}^{\cdot}$
soit dans le sous-groupe $\{1, \sigma_{\infty}\}$ de
$\mathfrak{S}_3$. En fait, le noyau de $\hat{\varphi}^{\cdot}$ est
le sous-groupe inv. de $\Pi_{0,\hat{3}}$ engendré
par $l_0^2, l_1^2, l_{\infty}^3$ , et la condition
envisagée signifie donc que modulo ce
sous-groupe, $\xi$ soit de la forme $l_{\infty}^i$ .
Par la formule (10) , on peut éliminer
$\beta$ de la donnée $(\alpha, \beta, f)$ , et on
peut remplacer la donnée $\alpha$ par la donnée
$\xi$. Ainsi :

\newpage

%% ===== Page 234 =====

Corollaire 2 Les éléments du s/groupe de $M_{0,3}^\sim$ formé des $u$ qui commutent extérieurement à $\hat{\rho}$, sont en corr. biunivoque avec l'ens des couples $(\xi, f)$ ( $\xi, f \in \Pi_{0,\hat{3}}$) satisfaisant les conditions (ne faisant intervenir que $\xi$)

[MARGIN: il faut de plus la "condition éliminatoire"...]

(14) $$ \xi \rho(\xi) \rho^2(\xi) = 1 $$

(15) $$ \xi = (\hat{\varphi}(\xi) \sigma_\infty (\hat{\varphi}(\xi))^{-1}) l_1^\nu $$
$$ (\nu \in \hat{\mathbb{Z}} \text{ convenable, tel que } 2\nu+1 \in \hat{\mathbb{Z}}^\times). $$

Pour un tel $\xi$, $\hat{\varphi}^\circ(\xi)$ ds $\mathfrak{S}_3$ est de la forme $\sigma_\infty^i$ , $i \in \hat{\mathbb{Z}} / 2\mathbb{Z}$, et posant

[MARGIN: en fait, $\sigma_\infty^i$ est vrai s-groupe par def. de $\alpha, \beta$ (16) bien sûr.]

$$ \beta = \hat{\varphi}(\xi) \sigma_\infty^{-i} \ (\in \Pi_{0,\hat{3}}) $$

$$ \alpha \in \Pi_{0,\hat{3}} \text{ l'unique él. de } \Pi_{0,\hat{3}}^\sim \simeq \dots \text{ tel que } \alpha \rho(\alpha)^{-1} = \xi $$

$u$ se déduit de $\xi, f$ par les formules

(17)
$$ \begin{cases} u(\sigma_\infty) = int(f^{-1} \beta)(\sigma_\infty) = int(f^{-1}) (int(\beta)(\sigma_\infty)) = int(f^{-1})(\xi l_1^i \sigma_\infty) \\ u(\rho) = int(f^{-1} \alpha)(\rho) = int(f^{-1}) (int(\alpha)(\rho)) = int(f^{-1})(\xi \rho) \end{cases} $$

On notera que ds ces formules, $f$ n'intervient que par le fait de composer avec $int(f)$. Si on pose d'ailleurs, pour $\xi \in \Pi_{0,\hat{3}}$

\newpage

%% ===== Page 235 =====

[MARGIN: cette notation a déjà été utili-
sée dans des condi-
tions plus géné-
rales différant
un peu de celles
ici. Voy. plus
bas.]

faisant (14)(15), par $u_\xi$ l'autom.
(de $\widehat{\Gamma}_{0,3}$) correspondant, aux termes du
Corollaire 2, au couple $(\xi,1)$, de
sorte que

$$ (18) \quad \begin{cases} u_\xi(\sigma_\infty) = (\xi h^{-\nu}) \sigma_\infty \\ u_\xi(\rho) = \xi \rho \end{cases} $$

les formules (17) s'écrivent aussi
simplement

$$ (19) \quad u_{\xi,f} = int(f^{-1}) u_\xi $$

D'autre part les formules (5), (6)
donnent :

$$ (20) \quad \hat{\phi} \circ u_\xi = u_\xi \circ \hat{\phi} $$

où, bien sûr, dans le premier membre
$u_\xi$ désigne la restriction de $u_\xi$ à $\Pi_{0,\hat{3}}$ .

Notons

Corollaire 3 Les automorphismes (géométriques$^\ast$)
de $\widehat{\Gamma}_{0,3}$ de la forme $u_\xi$ sont
exactement ceux pour lesquels on
a la propriété formule (20) (i.e. qui
commutent à $\hat{\phi}$) et qui de plus
normalisent $L_0$.

[MARGIN: $\ast$ au sens de p. 416]

\newpage

%% ===== Page 236 =====

En effet, en termes des triplets "spéciaux" $(\alpha, \beta, f)$ ces
propriétés s'expriment par les conditions
que l'on ait $f$ de la forme $l_0^n$
(qui exprime la condition de normaliser $L_0$)
et que $\lambda$ donné par (6) soit $= 1$, i.e.
(21) $$ \lambda = \hat{\varphi}(f)^{-1} f $$
soit égal à $1$, i.e. la condition
(22) $$ \hat{\varphi}(f) = f $$ ,
qui en l'occurrence s'écrit
$$ \varepsilon_0^n l_0^n = l_0^n $$ i.e. $$ \varepsilon_0^n = 1 $$
ce qui implique $f = l_0^0 = 1$, donc
$u$ de la forme $u_{(\alpha, \beta, 1)} = u_\xi$, où $\xi = \alpha p(\alpha)^{-1}$.

Corollaire 4. Les $u_\xi$ forment un
sous-groupe $\operatorname{Norm}_{M_{0,3}^{\sim}} (L_0)$,
soit $M_{0,3}^{1\sim}$, dont l'intersection avec
$\Pi_{0,3}^\wedge$ se réduit à $\{1\}$, et dont l'image dans
[MARGIN: Remarque $N_{0,3}$] $N_{0,3}$ est formé des automorphismes
ext. [unclear] de $\widehat{\mathcal{C}}_{0,3}^+$, stabilisant
$\Pi_{0,3}^\wedge$ (i.e. $u \in M_{0,3}^\wedge$) tels que $\hat{\varphi} : \Pi_{0,3}^\wedge \xrightarrow{\sim}$
commute extérieurement à $u$, d'où
un isomorphisme
(23) $$ M_{0,3}^{1\sim} \simeq \{ u_\xi \mid \xi \text{ solution de } (14), (15) \} \xrightarrow{\sim} N_{0,3}^{1\sim} $$

\newpage

%% ===== Page 237 =====

419

Il reste seulement à vérifier qu'un
$u_{\xi}$ ne peut être un autom intérieur que
si $u_{\xi} = 1$ i.e. $\xi = 1$ (donc $\nu = 0$). Or
$u_{\xi}$ est un $u_{\alpha, \beta} = u_{\alpha, \beta, 1}$ ($\alpha, \beta$ définis
par $\alpha \rho(\alpha)^{-1} = \xi$, [unclear: $\rho \alpha (\rho \beta)^{-1} = \xi \ell_1^{-\nu}$]),
qui est un autom. int. ssi on a
$\alpha = \ell_0^{-n}, \beta = \ell_0^{-n}$ (cf prop. de la page 404), i.e.
$\xi = \ell_0^{-n} \ell_1^n$, [unclear: $y = \rho(\alpha) \beta^{-1} = \ell_0^{-n} \ell_1^{-n}$], et la relation
[unclear] $\xi = y \ell_1^{\nu}$ s'écrit $\ell_1^{\nu} = 1$ i.e. $\nu = 0$. La
relation (15) donne donc (pour
$\varphi(\xi) = \hat{\varphi}(\ell_0^{-n} \ell_1^n) = \varepsilon_0^{-n} \sigma_{\infty}^{-n}$
$\sigma_{\infty}(\hat{\varphi}(\xi)) = \varepsilon_1^{-n} \sigma_{\infty}^{-n}$ ) donne la relation
$\ell_0^{-n} \ell_1^{-n} = \varepsilon_0^{-n} \sigma_{\infty}^{-n} \varepsilon_1^{-n} \sigma_{\infty}^{-n}$
i.e. (comme $\ell_0 = \varepsilon_0^2, \ell_1 = \varepsilon_1^2$) si $n$ pair
$\varepsilon_0^n = \varepsilon_1^{3n}$ d'où $n=0$,
et si $n$ impair
$\varepsilon_0^{-n} = \varepsilon_0^{-2} \sigma_{\infty} \varepsilon_1^{-n} \sigma_{\infty} \varepsilon_1^{-2}$
$\sigma_{\infty}(\varepsilon_1)^{-n} = \varepsilon_0^{-n}$
i.e. $\varepsilon_1^{2n} = 0$, i.e. encore $n=0$, cqfd.

[MARGIN: N.B. Il y a un choix de signes des générateurs de [unclear] de $\hat{N}_{0,3}^{1 \simeq}$ fait au tirage [unclear] pour voir si sa restriction à $N_{0,3}^{1 \simeq}$ ...]

On a donc trouvé une section
partielle canonique de $\hat{M}_{0,3}^{1 \simeq}$ au dessus
des sous-groupes $N_{0,3}^{1 \simeq}$ de $\hat{N}_{0,3}^{1 \simeq}$. Comme

\newpage

%% ===== Page 238 =====

ce sous-groupe central $\Gamma_{\mathbf{Q}}$, on trouve
de [crossed out] une section de l'extension $\mathcal{M}_{0,3}^1$ de
$\Gamma_{\mathbf{Q}}$ par $\hat{\Pi}_3$, [crossed out] qui en fait une
section du sous-groupe de $\mathcal{M}_{0,3}^1$ normalisateur
de $L_0$. Il faudrait vérifier que
cette section canonique est associée
au choix d'une uniformisante pour la
pointe $0$ de $U_{0,3 \mathbf{Q}}$, [crossed out] i.e. d'un isomorphisme
[crossed out: $\mathcal{M}_0 / \mathcal{M}_0^2$] à un $\mathbf{Q}$-isomorphisme , où $\mathcal{M}_0$ désigne
l'idéal maximal de l'anneau local en $0$
de $U_{0,3 \mathbf{Q}} = \mathbf{P}_{\mathbf{Q}}^1 \setminus \dots$ - il faudrait
déterminer cet élément de $\mathcal{M}_0 / \mathcal{M}_0^2$ -
On a envie de prendre un élément qui
soit, non seulement un bon sur $\mathbf{Q}$,
mais (vu la situation sur les entiers)
un bon sur $\mathbf{Z}$ - condition qui détermine
un élément de $\mathcal{M}_0 / \mathcal{M}_0^2$ [crossed out] à un
signe près - ce sera en fait le
choix correspondant à l'uniformisante
$z$, ou $-z$. Mais est-ce bien le bon ?
Il faudrait y revenir !

\newpage

%% ===== Page 239 =====

Travaillons désormais dans le sous-groupe
$\mathcal{M}_{0,3}^{!\simeq}$ de $\mathcal{M}_{0,3}^{!\sim}$, formé des $u$ qui
commutent rationnellement à $\hat{\rho}$ ,
i.e. l'image inverse de $\mathcal{N}_{0,3}^{!\simeq}$, de sorte

qu'on a
(24) $1 \to \Pi_{0,\hat{3}} \to \mathcal{M}_{0,3}^{!\simeq} \to \mathcal{N}_{0,3}^{!\simeq} \to 1$

On a donc trouvé un sous-$\mathbb{Q}$-groupe section
(25) $\mathcal{M}_{0,3}^{!\simeq}(0) \overset{\sim}{\to} \mathcal{N}_{0,3}^{!\simeq}$
$\{ u_f \mid f \in \Pi_{0,\hat{3}} \text{ sol. de } (14), \text{ et } (15) \} \subset \mathcal{M}_{0,3}^{!\simeq}$

En termes de celui-ci on peut donc
écrire $\mathcal{M}_{0,3}^{!\simeq}$ comme un produit 1/2 direct
(26) $\mathcal{M}_{0,3}^{!\simeq} \simeq \mathcal{M}_{0,3}^{!\simeq}(0) \cdot \Pi_{0,\hat{3}}$ ,

en écrivant tout élément de $\mathcal{M}_{0,3}^{!\simeq}$ de
façon unique comme

(27) $u = u_0 f \quad \begin{cases} u_0 \in \mathcal{M}_{0,3}^{!\simeq}(0) \text{ i.e. } \\ u_0 = u_{\dots} \text{ sol. de } (14)(15) ; \\ f \in \Pi_{0,\hat{3}} \text{ bien déterminé} \end{cases}$

( $u_f$ du (19), soit , )

Ce sera justement le 
(28) $u_f = u_{\alpha, \beta} \cdot f$

avec $\alpha, \beta$ déterminés en termes de $f$ par
(29) $\alpha \rho(\alpha)^{-1} = \{ f \}$ , $\beta \varpi(\beta)^{-1} = \{ f \}^{-1}$ ,

\newpage

%% ===== Page 240 =====

c'est ainsi qu'on pu interpréter les
triples $(\alpha, \beta, f)$ normalisés du cor. 1.
(récapitulation)

Rappelons que le groupe $\mathcal{N}_{0,3}^{!\simeq} \simeq \mathcal{M}_{0,3}^{!\simeq}(0)$ est
en corr. 1-1 avec l'ens des $\xi \in \widehat{\Pi}_{0,3}$
satisfaisant les deux relations

(30) $$ \begin{cases} \xi \rho(\xi)\rho^2(\xi)=1 \\ \xi = (\hat{\varphi}(\xi) \sigma_\infty (\hat{\varphi}(\xi))^{-1}) l_1^\nu \end{cases} \quad \text{pour } \nu \in \hat{\mathbb{Z}} \text{ conv.} $$

ou encore (posant $\xi = \alpha \rho(\alpha)^{-1}$, $\alpha \in \widehat{\Pi}_{0,3}$)
avec l'ens des $\alpha \in \widehat{\Pi}_{0,3}$ satisfaisant la
deuxième relation (30), dans laquelle on remplace
$\xi$ par $\alpha \rho(\alpha)^{-1}$. L'él. $u_\xi$ de $\mathcal{M}_{0,3}^{!\simeq}(0)$
défini par $\xi$ est défini par les formules
(18), i.e. par

(31) $$ u_\xi(\sigma_\infty) = (\xi l_1^{-\nu}) \sigma_\infty \quad , \quad u_\xi(\rho) = \xi \rho $$

Enfin, si on a $u_\xi$ et $u_{\xi'}$, alors

(32) $$ u_{\xi'} u_\xi = u_{\xi''} $$

avec $\xi''$ donné grâce aux formules
p. 39[unclear] (52) — en notant qu'ici $\xi = g, \xi' = g', \xi''=g''$ :
$\xi'' = \alpha'' \rho(\alpha'')^{-1}$ avec $\alpha'' = u_{\xi'}(\alpha) \cdot \alpha'$ d'où
[unclear crossed out equation: $\xi'' = u_{\xi'}(\alpha) \alpha' \rho(\alpha')^{-1} \rho(u_{\xi'}(\alpha))^{-1}$]

\newpage

%% ===== Page 241 =====

(33) $$\xi'' = u'(\xi) . \xi'$$

Posons maintenant

(34) $$\mathcal{N}_{0,3}^{\approx} = \mathfrak{S}_3 \times \mathcal{N}_{0,3}^{! \approx} \subset \mathcal{N}_{0,3}^{\sim} = \mathfrak{S}_3 \times \mathcal{N}_{0,3}^{! \sim}$$

et soit $\mathcal{M}_{0,3}^{\approx}$ l'image inverse

dans $\mathcal{M}_{0,3}^{\sim}$ de $\mathcal{N}_{0,3}^{\approx} \subset \mathcal{N}_{0,3}^{\sim}$, de sorte

qu'on a une structure d'extension

(35) $$1 \longrightarrow \widehat{\Pi}_{0,3} \longrightarrow \mathcal{M}_{0,3}^{\approx} \longrightarrow \mathcal{N}_{0,3}^{\approx} \longrightarrow 1$$

contenant l'extension (24) comme

[unclear: crossed out lines]

(36) $$1 \longrightarrow \widehat{C}_{0,3}^+ \longrightarrow \mathcal{M}_{0,3}^{\approx} \longrightarrow \mathcal{N}_{0,3}^{! \approx} \longrightarrow 1,$$

contenant (24) comme sous-extension ou

quotient par un sous-groupe plus petit.

Le sous-groupe $\mathcal{M}_{0,3}^{! \approx}(\infty)$, qui est une

section de (24), en donne ainsi

section de (36), d'où un iso can.

(37) $$\mathcal{M}_{0,3}^{\approx} \simeq \mathcal{M}_{0,3}^{! \approx}(\infty) \ltimes \widehat{C}_{0,3}^+$$

(produit $\frac{1}{2}$-direct).

\newpage

%% ===== Page 242 =====

Considérons maintenant l'application

(38)
$$ \underline{\Phi} : M_{0,3}^{!\approx} \overset{\simeq}{\longrightarrow} M_{0,3}^{\approx} $$
$$ \simeq \qquad \qquad \simeq $$
$$ N_{0,3}^{!\approx}(\omega) \cdot \hat{\Pi}_{0,3} \qquad M_{0,3}^{\approx}(\omega) \cdot \widehat{\mathfrak{S}}_{0,3}^+ $$

qui est donné par

(39) $$ \underline{\Phi}(u.a) = u \hat{\varphi}(a) $$

($u \in N_{0,3}^{!\approx}(\omega)$, $a \in \hat{\Pi}_{0,3}$). Il résulte de la caractérisation de $M_{0,3}^{\approx}(\omega)$ (cor. 3) - plus particulièrement du fait que les éléments de $M_{0,3}^{\approx}(\omega)$ opèrent effectivement (pas seulement à autom. int. près ...) ; que $\underline{\Phi}$ est un homomorphisme de groupes. Ce n'est autre bien sûr que l'hom. $\underline{\Phi}$ envisagé dans le (à restriction près aux sous-groupes $M_{0,3}^! = \mathcal{E}_{0,3}$ et $M_{0,3}$) § 41 (p. 316') - dans loc. cit. cet hom. pouvait se définir d'ailleurs de façon effective - par générateurs : $\hat{\Pi}_{0,3} - conj.$ puis à $\widehat{\mathfrak{S}}_{0,3}^+ - conj.$ puis - et avec un peu de soin devrait pouvoir vérifier l'égalité de ces homomorphismes $\underline{\Phi}$ du § 41 et d'ici, de façon effective.

\newpage

%% ===== Page 243 =====

Calcul de $x=l_\infty$. 422

Je vais revenir sur la détermination des
$u_f \in \mathcal{M}_{0,3}^{! \sim}[0] = \operatorname{Norm}_{\mathcal{M}_{0,3}^{! \sim}}(L_0)$ qui sont normalisés
i.e. qui commutent à $\hat{\varphi}$ i.e. satisfont (5)
avec $\lambda=1$. Pour $x=l_0$ cette formule, en vertu
de la forme particulière de $u_f$, est tautologique
($\hat{\varphi} u(l_0) = \hat{\varphi}(l_0^u) = \xi_0^u$, $u\hat{\varphi}(l_0) = u(l_0) = \xi_0^u$), donc on
voit qu'il reste = à la la poser pour
$x=l_1$, - pour $x=l_\infty$. Pour $x=l_1$, la condition
prend la forme (1D) — remarquable du fait
qu'elle donne un calcul de la "délicate"
"délicate" $\beta$ (dont l'existence en principe,
en termes de $\gamma$ repose sur des fondements
qui n'ont pas encore été établis...), en termes
des paramètres "grossiers" $\gamma$ — compte tenu
du fait que $\epsilon \in \hat{\mathbb{Z}}^*$ lui-même est déterminé...
aisément en termes de $\gamma$.

Nous allons donner cette condition sous
forme [unclear] un peu différente, en
écrivant la relation de commutation (5)
pour $x=l_\infty$, au lieu de $x=l_1$. Notons d'abord
qu'on a la relation
(40) $$ \tau'_\infty(\rho) = \rho^{-1} $$
car $$ \tau'_\infty(\rho) = \sigma_\infty(\tau_\infty(\rho)) = \sigma_\infty \rho \sigma_\infty^{-1} = \xi_0^{-1} \sigma_\infty^{-1} = \rho^{-1} $$

\newpage

%% ===== Page 244 =====

il en résulte que si $u \in \tilde{\mathcal{M}}_{0,3}^\sim$ satisfait
(41) $u(\rho) = \rho^\mu \quad \mu \in \hat{\mathbb{Z}}^*$ convenable
alors on a nécessairement $\mu \equiv 1 \text{ ou } -1 \quad (3)$
et si $\alpha_3 = \alpha_3(\mu) \in \mathbb{Z}/2\mathbb{Z}$ est l'entier
(42) $\alpha_3(\mu) \in \mathbb{Z}/2\mathbb{Z}$ caractérisé par $\mu \equiv (-1)^{\alpha_3(\mu)} \ (3)$
on aura donc $u(\rho) = \rho$ ou $u(\rho) = \rho^{-1}$ suivant
que $\alpha_3(\mu) = 0$ ou $\alpha_3(\mu) = 1$, et en tous cas
on aura
(43) $u(\rho) = \tau_\infty^{\alpha_3}(\rho)$.

Ceci posé, faisons $x = \ell_\infty$ dans $(5)$ :
$\hat{\varphi} u(\ell_\infty) = \hat{\varphi}(\text{int}(f_0^{-1}) \ell_\infty^\mu) = \text{int}(\hat{\varphi}(f_0)^{-1}) \rho^{-\mu}$
(car $\hat{\varphi}(\ell_\infty) = \rho^{-1}$), et
$u(\hat{\varphi}(\ell_\infty)) = u(\rho^{-1}) = u(\rho)^{-1} = $ [unclear crossed out]

donc la formule $(5)$ pour $\lambda = 1$ devient
$\text{int}( \hat{\varphi}(f_0) \alpha) (\rho^{-1}) = \rho^{-\mu}$ soit
$\text{int}( \hat{\varphi}(f_0) \alpha) \rho = \rho^\mu$
et en vertu de $(43)$ [unclear crossed out]
$\rho^\mu = \tau_\infty^{\alpha_3}(\rho) \stackrel{\text{déf}}{=} \text{int}( (\tau_\infty')^{\alpha_3} ) (\rho)$
donc la condition devient
$\text{int}( (\tau_\infty')^{-\alpha_3} \hat{\varphi}(f_0) \alpha ) \rho = \rho$
i.e. que l'élément $(\tau_\infty')^{-\alpha_3} \hat{\varphi}(f_0) \alpha$ de
$\widehat{G}_{0,3}^\sim = \widehat{G}_{0,3}^+ \cdot \{ 1, \tau_\infty \} $ commute à $\rho$.

\newpage

%% ===== Page 245 =====

[MARGIN: il faudrait voir les formules précises !!!]

d'après des principes généraux (pas bien
stables, à vrai dire, dans le cadre profini)
le commutant de $\rho$ dans le groupe en
question est réduit au groupe engendré par
$\rho$, $\{1, \rho, \rho^2\}$, de sorte que la relation *)
prend la forme

$$ \tau_\infty^{\alpha_3} \hat{\varphi}(f_\infty) \alpha = \rho^j \quad \text{pour } j \in \hat{\mathbb{Z}}/3\hat{\mathbb{Z}} \text{ conv.} $$
(bien déterminé).

On l'explicite en utilisant la valeur de $f_\infty$
$$ f_\infty^{-1} = g \rho(g) \rho^2(f_0^{-1}) = g \rho(g) = \rho^2(g)^{-1} = \rho^{-1}(\xi^{-1}) $$
NB ici $f_0^{-1} = 1$

donc la relation devient
$$ \tau_\infty^{\alpha_3} \hat{\varphi}(\rho^{-1}(\xi)) \alpha = \rho^j $$

soit
$$ \hat{\varphi}(\rho^{-1}(\xi)) = \tau_\infty^{-\alpha_3} \rho^j \alpha^{-1} $$

ou aussi
[unclear: crossed out text]

comme $\tau_\infty'(\rho) = \rho^{-1}$ et $\tau_\infty^{'\alpha_3}(\rho^j) = \rho^{(-1)^{\alpha_3} j}$

$$ (44) \quad \hat{\varphi}(\rho^{-1}(\xi)) = \rho^{j'} \tau_\infty^{'\alpha_3} \alpha^{-1} $$
où $\alpha_3 = \alpha_3(u) \in \hat{\mathbb{Z}}/2\hat{\mathbb{Z}}$
défini par $(-1)^{\alpha_3} = \mu$ (3)
et $j' \in \hat{\mathbb{Z}}/3\hat{\mathbb{Z}}$ (bien déterminé)

Nous allons encore transformer un peu cette
relation, en écrivant
$$ \tau_\infty^{'\alpha_3} = (\sigma_\infty \tau_\infty)^{\alpha_3} = \sigma_\infty^{\alpha_3} \tau_\infty^{\alpha_3} \quad \text{d'où} $$
$$ \tau_\infty^{'\alpha_3} \alpha^{-1} = \sigma_\infty^{\alpha_3} \underbrace{(\alpha \tau_\infty^{-\alpha_3})^{-1}}_{\alpha} $$

\newpage

%% ===== Page 246 =====

où

(45) $$\tilde{\alpha} = \alpha \tau_{\infty}^{\alpha_3}$$

a l'avantage sur $\alpha$ lui-même d'être dans
$\hat{\pi}_{0,\hat{1}}$ dans tous les cas (NB on a
$\alpha \in \hat{\pi}_{0,\hat{1}} . \{1, \tau_{\infty}\}$, et $\alpha \in \hat{\pi}_{0,\hat{1}}$ ssi $\alpha_3 = 0$).
Il est égal à $\alpha$ ssi $\alpha_3 = 0$ et dans le
cas $\alpha_3 = 1$ il est relié directement à $\xi$
par la formule

(46) $$\xi = \tilde{\alpha} l_1^{-1} \rho(\tilde{\alpha})^{-1}$$

(déduite de (45), de $\xi = \alpha \rho(\alpha)^{-1}$ et de
(**) $$\tau_{\infty} \rho(\tau_{\infty})^{-1} = l_1^{-1}$$ ).

[MARGIN: tous ces calculs faits au signe près $\pm 1$ en exposant qui peut s'introduire dans [unclear] i.e. à multiplicité près]

Ainsi (44) prend la forme définitive

(47) $$\hat{\phi}(\rho^{-1}(\xi)) = \rho^{j'} \tilde{\alpha}^{-1}$$ \quad ($\tilde{\alpha} = \alpha \tau_{\infty}^{\alpha_3}$)

qui peut s'écrire, en passant aux inverses

(47 bis) $$\hat{\phi}(\rho^{-1}(\xi^{-1})) = \tilde{\alpha} \rho^{-j'}$$

ou encore, en isolant $\tilde{\alpha}$

(48) $$\tilde{\alpha} = \hat{\phi}(\rho^{-1}(\xi^{-1})) \rho^{j'}$$
$$_{\|} $$
$$\alpha \tau_{\infty}^{\alpha_3}$$

On rapproche cette formule de la
formule (10), qui se résoud en $\beta$ sous la forme

(49)=(10) $$\hat{\phi}(\xi) = \beta \sigma_{\infty}^{-1}$$
(50) $$\beta = \hat{\phi}(\xi) \sigma_{\infty}$$

\newpage

%% ===== Page 247 =====

424

Le fait que la formule (47) est bien équivalente à (10), ou encore à (49), lorsque $\alpha \in \hat{\Pi}_{0,1}^\wedge$ et $\beta \in \hat{\Pi}_{0,1}$ sont liés par $\xi = y l_1^i$ (où $\xi = \alpha \rho(\alpha)^{-1}$, $y = \beta \sigma_{\infty}(y)^{-1}$), mériterait une vérification directe. L'intérêt de (47), - ou de façon résolue en $\alpha$ (48), tout comme de (10) $\overset{(49)}{\leftrightarrow}$ (50), consiste en ce que ces formules donnent les paramètres "délicats" $\alpha, \beta$ en termes des "paramètres grossiers" $\xi$. L'une et l'autre sont des formules dans $\hat{\mathcal{B}}_{0,s}^+$ - redonnées dans $\hat{\Pi}_{0,1}$, elles donnent des formules dans $\hat{\mathcal{B}}_1$, savoir (compte tenu que $\tilde{\alpha}, \beta \in \hat{\Pi}_{0,1}$) :

(51) $$ \begin{cases} \hat{\varphi}(p^{-1}(q_1)) = \rho^{j'} \sigma_\infty^{\alpha_3} \\ \hat{\varphi}(\xi) = \sigma_\infty^i \end{cases} $$

ce qui redonne exactement, en termes des éléments $\hat{\varphi}(y)$ et $\hat{\varphi}(p^{-1}(y))$ de $\hat{\mathcal{B}}_1$, les valeurs de $i, \alpha_3 \in \mathbb{Z}/2\mathbb{Z}$ et de $j' \in \mathbb{Z}/3\mathbb{Z}$.

[MARGIN: A vrai dire [unclear] $(i, \alpha_3) \in \mathbb{Z}/2\mathbb{Z}$ [unclear] multiplicateur $\mu$]
-- (Il faudrait voir si $i, j'$ peuvent, comme $\alpha_3$, s'expliciter en termes du multiplicateur $\mu$, voire en termes de [unclear: eq])

\newpage

%% ===== Page 248 =====

restes quadratiques $\varepsilon_l(\mu)$ ) On peut les
compléter par une formule déduite de
l'une et de l'autre, compte tenu que
$$ \{ p(\xi) p'(\xi) = 1 \quad \text{ie.} \quad p^{-1}(\xi) \xi = p'(\xi)^{-1} $$
[unclear: en tirant de,] en multipliant (47) et (10),
$$ \hat{\phi}( p^{-1}(\xi) . \xi ) = p' j' \sigma_\infty^{\alpha_3} \tilde{\alpha}^{-1} \beta \sigma_\infty^i $$
$$ (p'(\xi)^{-1}) $$
soit encore

(52) $$ \hat{\phi}(p(\xi)) = \sigma_\infty^{-i} (\beta^{-1} \tilde{\alpha}) \sigma_\infty^{-\alpha_3} p j' $$

soit, sous forme résolue en $\beta^{-1} \tilde{\alpha}$

(53) $$ \beta^{-1} \tilde{\alpha} = \sigma_\infty^i \hat{\phi}(p(\xi)) \sigma_\infty^{\alpha_3} p j' \ . $$

On peut considérer que les trois
"paramètres délicats" principaux sont
les quantités $\beta, \tilde{\alpha}$ et $\beta^{-1} \tilde{\alpha}$, qui s'expri-
ment directement, de façon remarqua-
blement simple, en termes des valeurs
de $\hat{\phi}$ pour $\xi, p^{-1}(\xi)$ et $p'^{-1}(\xi) = p(\xi)$, ce
qui dispenserait [unclear: faisant] sans doute (en

247

\newpage

%% ===== Page 249 =====

sachant les calculs en termes de $\xi,\eta$
exclusivement, ce qui demanderait un
peu plus de soins...) de faire appel à des
résultats délicats non encore démontrés,
du type $\xi = \alpha \rho(\alpha)^{-1}$, $\eta = \beta \sigma_{\infty}(\beta)^{-1}$,
pour justifier l'introduction de ces paramè-
tres $\alpha$, $\beta$.

[MARGIN: (54) récapitulation]
$$ \hat{\varphi}(q) = \beta \sigma_{\infty}^{-i} $$
$$ \hat{\varphi}(\rho^{-1}(q)) = \rho' \sigma_{\infty}^{\alpha_3} \tilde{\alpha}^{-1} \qquad (\tilde{\alpha} = \alpha \tau_{\infty}^{\alpha_3}) $$
$$ \hat{\varphi}(\rho(q)) = \sigma_{\infty}^{-i} (\beta^{-1} \tilde{\alpha}) \sigma_{\infty}^{\alpha_3} \rho'^{-1} $$
$$ \rho(q) \qquad \sigma_{\infty}^{-i} (\beta^{-1} \tilde{\alpha}) \sigma_{\infty}^{\alpha_3+i} \rho'^{-1} $$

Supposons pour simplifier

(55) $\alpha_3 = 0 \quad$ i.e. $\quad \alpha = \tilde{\alpha} \quad$ i.e. $\quad \alpha \in \Pi_{0, \hat{S}}$

Alors on tire de la deuxième formule
(54), en prenant les inverses
$$ \hat{\varphi}(\rho^{-1}(q^{-1})) = \alpha \rho'^{-1} $$
et le $\times \rho(X)^{-1}$ des deux membres
(56) $$ \xi = \hat{\varphi}(\rho^{-1}(q^{-1})) \rho' \hat{\varphi}(\rho^{-1}(q))^{-1} $$
qui implique déjà la première formule
(14) dans les formules fond. (14) (15).

\newpage

%% ===== Page 250 =====

425

On trouve donc dans ce cas ($\mu\equiv 1 (3)$ i.e. $\nu\equiv 0 (3)$) que les deux relations (14) (15) équivalent en fait à l'une des deux relations

[MARGIN: valable seulement si $\nu\equiv 0 (3)$ i.e. $\mu\equiv 1 (3)$ i.e. $\varepsilon_3(u)=1$]

(57) $$ \xi = \hat{\varphi}(\rho'(q^{-1})) \rho(\hat{\varphi}(\rho'(q^{-1})))^{-1} $$
(58) $$ \xi = \hat{\varphi}(q) \sigma_\infty(\hat{\varphi}(q))^{-1} h^\nu $$

qui caractérisent ainsi les $u = u_q \in M_{0,3}^{1 \approx}(0)$, dont le multiplicateur $\mu$ est $\equiv 1 (3)$.

Rmq Le seul élément de $M_{0,3}^{1 \approx}(0) \simeq \widehat{N_{0,3}}^\approx$ qui se relève explicitement, en dehors de l'élément unité (et donc l'unique aussi des $N_{0,3}^{1 \approx} \simeq \widehat{N_{0,3}}^\approx$) est l'unique $\tau \in C_{0,3}^{1 \approx} \subset \widehat{N_{0,3}}^\approx$ qui dans $M_{0,3}^{1 \approx}(0)$ correspond à l'élément $\tau_\infty$, on a $u = u_y$ [unclear].

[MARGIN: Attention, ici $\alpha_3=1$ donc (54) n'est pas valable]

(59) $h=\beta_0=1$, $f=1$, $\alpha=\alpha_0=\tau_\infty$, $\xi=g=l_1^{-1}$, $\eta=1$

$\nu=-1$, $\mu=-1$

satisfaisant bien à $\xi \rho(\xi) \rho^2(\xi) = 1$ [unclear]

la $1^{\underline{e}}$ relation [unclear]

la $2^{\underline{e}}$ relation (54) [unclear]

\newpage

%% ===== Page 251 =====

avec

(60) $i=1, \alpha_3=1$ i.e. $\varepsilon_3=-1, j'=-1$

Rmq On peut se poser la question, de façon générale, de déterminer les
éléments $u_{\alpha,\beta} = u_{\xi,\eta}$ de $\mathcal{M}_{0,3}^{1 \sim}[0]$

(pas nécessairement normalisés dans un sens ou dans l'autre) tels qu'on
ait non seulement par hyp., $f=1$;
mais aussi $\beta=1$ i.e. $\eta=1$
(NB On sait déjà que si $\alpha=1$
i.e. $\xi=1$, alors $u=1$) On aura ainsi
$\xi = \eta h^\nu = \rho_1^{2\nu}$, et il reste à vérifier

[MARGIN: en refaisant le calcul de la page 424 je m'aperçois que les "multiplicateurs" sont bien plus compliqués à cause des "coefficients" ... [unclear]]

$\{\rho(\xi)\rho(\eta)\}=1$ avec $\xi=h^\nu$,

dans $\operatorname{Sl}(2,\hat{\mathbb{Z}})'$, cela dit on note que

$$ \xi = \begin{pmatrix} 1 & 0 \\ 2\nu & 1 \end{pmatrix} $$

(p. 248)

ceci signifie que l'on ait soit

$\xi = \pm f^{-1}$ [unclear crossed out math] - ce

qui manifestera rien pour ces - soit

$C+D - B = \pm 1$ i.e. $2\nu + 1 = \mu = \pm 1$ i.e.

$\nu=0$ ou $-1$ ce qui donne les deux cas

triviaux $u_\xi = id$ et $u_\xi = \tau_\infty$.

\newpage

%% ===== Page 252 =====

Notons la

[MARGIN: NB. Cette proposition revient à dire que dans $W \simeq$ [unclear]]
Proposition. Le commutant de $\tau$ (élément non trivial de $\Gamma_{0,3} \simeq \mathbb{Z}/2\mathbb{Z} \subset M_{0,3}^{\sim !}$) dans $M_{0,3}^{\sim !}$ est réduit à $\{1, \tau\}$.

On fera le calcul dans le groupe isomorphe $M_{0,3}^{\sim !}(0)$, où $\tau$ correspond à $\tau_\infty$, en fait il se trouve que le centralisateur de $\tau$ dans $M_{0,3}^{\sim !}[0]$ tout entier (le groupe normalisateur au lieu de $M_{0,3}^{\sim !}$, formé des $u_\xi$ sans conditions de normalisation) est réduit à $\{1, \tau\}$. En effet, il faut exprimer la commutation d'un $u = u_\xi$ avec $u_{\xi_1} = \tau_\infty$, $\xi_1 = l_1^{-1}$, on a
$$
\begin{cases}
u_{\xi_1} u_{\xi} = u_{\xi''} & \text{avec } \xi'' = u_{\xi_1}(\xi)\xi_1 \\
u_{\xi} u_{\xi_1} = u_{\xi'''} & \text{avec } \xi''' = u_{\xi}(\xi_1)\xi
\end{cases}
$$
et il faut exprimer $\xi'' = \xi'''$, i.e.
(61) $\tau_\infty(\xi) l_1^{-1} = u_\xi(l_1^{-1})\xi$

Or on a 
$u_\xi(\varepsilon_1) = int(\xi)(\varepsilon_1^H)$ d'où
$u_\xi(\underbrace{l_1^{-1}}_{\varepsilon_1^{-2}}) = int(\xi)(l_1^{-H})$ et (61) s'écrit
$\tau_\infty(\xi) l_1^{-1} = \xi l_1^{-H}$ i.e.
(62) $\xi^{-1} \tau_\infty(\xi) = l_1^{1-H} = l_1^{-2\nu}$

\newpage

%% ===== Page 253 =====

427

Cette équation a la solution
(63) $$ \xi = l_1^\nu $$
dans $\Pi_{0,\hat{3}}$ (car $\tau_\infty(l_1) = l_1^{-1}$ donc $\tau_\infty(\xi) = l_1^{-\nu}$)
et cette solution est unique dès qu'on admet
que les résultats discrets se prolongent au
cas profini, de façon à avoir
(64) $$ (\Pi_{0,\hat{3}})^{\tau_\infty} = \{1\} $$
Or il faut de plus que $\xi$ satisfasse la
relation
$$ \xi \rho(\xi) \rho^2(\xi) = 1 $$
et on vient de voir que cela implique
$\nu = \pm 1$ , i.e. $u \in \{1, \tau_\infty\}$.

\newpage

%% ===== Page 254 =====

428

48. Construction d'un schéma en groupes.

I) Étude de l'application $\alpha \longmapsto \alpha \rho(\alpha)^{-1} = [\alpha, \rho]$.

Soit $G$ un schéma en groupes sur un anneau $A$ (p.ex. $A = \mathbb{Z}$, $G = \operatorname{Sl}(2)_{\mathbb{Z}}$ ou $\operatorname{Gl}(2)_{\mathbb{Z}}$), $\rho$ une section de $G$, et considérons le morphisme
(1) $F_\rho : \alpha \longmapsto \alpha \rho(\alpha)^{-1} = [\alpha, \rho] : G \to G$

Soit
$T_\rho = Centr_G(\rho) = G^\rho \subset G$

alors
(2) $F_\rho(\alpha') = F_\rho(\alpha) \iff \alpha^{-1}\alpha'$ section de $T_\rho$

donc on a un monomorphisme canonique induit par $F_\rho$
(3) $G/T_\rho \hookrightarrow G$ ,

dont il s'agit de caractériser l'image.

Dans le cas qui nous occupe, le groupe quotient $G_{ab}$ est représentable (c'est le groupe unité pour $G = \operatorname{Sl}(2)_{\mathbb{Z}}$, et c'est $\mathbb{G}_m$ pour $G = \operatorname{Gl}(2)_{\mathbb{Z}}$), et [unclear: on cherche] une condition pour une section $\xi \in G(S)$, au dessus d'un $A$-schéma $S$, pour qu'elle soit dans l'image, c'est à dire [unclear: qu'on] ait
(4) $\delta(\xi) = 1$

\newpage

%% ===== Page 255 =====

(5) $\delta : G \to G_{ab}$

est l'hom. canonique. Une deuxième condition, c'est que $\xi p$, qui doit être (localement) de la forme $\alpha p \alpha^{-1}$, soit conjugué localement à $p$ - cette condition est d'ailleurs nécessaire et suffisante, tautologiquement, pour que $\xi p$ appartienne à l'image de $F_p$, et implique donc la condition (4). Mais cette dernière peut servir parfois d'ingrédient, [unclear: car dans un] jeu de conditions de nature "calculatoire", qui soient néc. et suff. pour qu'on ait $\xi p \sim p$. Dans le cas d'un groupe réductif $G$, la condition de conjugaison locale de sections est d'ailleurs assez bien comprise, depuis notamment les travaux de Steinberg. On a un "schéma des invariants" $I_G$, et un morphisme canonique

(6) $G \to I_G$

(sur $A$, et comm. [unclear: après tt chg de base]) qui soit universel pour les morphismes de $G$ dans des schémas, qui soient inv. par autom. int. - i.e. $I_G$ est l'enveloppe représentable de $G/Int(G)$ (faisceau des classes [unclear])

\newpage

%% ===== Page 256 =====

429

de conjugaison
$$ (7) \quad G/int(G) \longrightarrow I_G $$

A vrai dire, ce morphisme est un épimorphisme (pour la top fppf), mais pas un iso i.e. pas un mono. En d'autres termes, pour deux sections (en l'occurence $\rho, \xi\rho$) de $G$, il est bien néc, mais pas tjrs suffisant, que leurs images dans $I_G$ soient égales, pour qu'elles soient conjuguées loc.t fppf. Mais dans le cas où les deux sections sont régulières au sens de Steinberg (ce qui est une condition sur chaque fibre) [unclear: le résultat est vrai].

Elle l'est aussi si les deux sections sont 1/2 simples (ce qui est également une condition sur chaque fibre).

Dans le cas d'un groupe loc. isom. à $\operatorname{Sl}(2)$, [unclear], ou $\operatorname{Gl}(2)$, les sections régulières au sens de Steinberg sont d'ailleurs celles qui en aucun point ne sont centrales - ce qui est bien le cas pour le $\rho$ ou le $\xi \rho$ des $\S$ précédents. D'autre part, $I_G$ est alors bien connu, dans le cas d'un groupe

\newpage

%% ===== Page 257 =====

déployés, plus généralement de $\operatorname{Sl}(E)$ ou $\operatorname{Gl}(E)$,
$E$ loc. libre de rang 2, il est can. isomorphe
à $E^1$ ou à $E^{1*} \times E^1$ (suivant qu'on a $Sl$ ou $Gl$)
dans le cas d'une forme de $\operatorname{Sl}(2)$
c'est donné par la trace
(7) $$G \longrightarrow I_G \simeq E^1$$
$$x \longmapsto \text{Tr} x$$
dans le cas d'une forme de $\operatorname{Gl}(2)$ c'est
donné par
(8) $$G \longrightarrow I_G = E^{1*} \times E^1$$
$$x \longmapsto (\det x, \text{tr} x)$$
On trouve donc (pour mémoire)

Proposition. Soit $G$ une forme de $\operatorname{Sl}(2)$ ou
$\operatorname{Gl}(2)$ sur un schéma $S$, soit $\rho$ une section
de $G$ qui n'est centrale en aucune fibre
(localement).
Pour qu'une section $\xi$ de $G$ soit (dans
l'image de $F_\rho$ ((1) ou (3))), il faut et
il suffit qu'elle satisfasse les conditions suivantes :

a) $\xi/\rho$ n'est centrale en aucune fibre, i.e.
en aucune fibre on n'a $\xi = \lambda \rho$ , où
$\lambda$ est élément central

b) $\text{tr} \xi = \text{tr} \rho$

c) (si $G$ forme de $\operatorname{Gl}(2)$) $\det \xi/\rho = 1$
i.e. $\xi/\rho$ section de $\operatorname{Sl}(2)$.
[MARGIN: équivaut à : $\det(\xi) = \det \rho$]

\newpage

%% ===== Page 258 =====

930

Corollaire Quand ces conditions sont satisfaites, le schéma des solutions (de l'équation $[\alpha, \rho] = \xi$)
est un torseur (à droite) sous $T_\rho$, qui est un schéma en groupes lisse commutatif sur $S$, de dimension relative 1 (si $G$ forme de $\operatorname{Sl}(2)$) ou 2 (si $G$ forme de $\operatorname{Gl}(2)$).

Exemples $G = \operatorname{Gl}(2)_\mathbb{Z}$, $\rho = \begin{pmatrix} 0 & 1 \\ -1 & 1 \end{pmatrix}$ est l'élément de $\operatorname{Sl}(2, \mathbb{Z})$ déjà noté ainsi (à partir de p. 335). Il est bien régulier, et de même l'élément $\sigma = \sigma_\infty = \begin{pmatrix} 0 & -1 \\ 1 & 0 \end{pmatrix}$.

De façon générale, la donnée d'une forme de $\operatorname{Gl}(2)$ équivaut à celle d'une forme de $\operatorname{Sl}(2)$, ou de $PGl(2)$, ou encore à celle d'une Algèbre d'Azumaya de rang 4 sur $S$, soit $\mathcal{M}$ - i.e. une Alg. loc. (fppf) isom. à $M_2(\mathcal{O}_S)$ - par

(9)
$$G(\mathcal{M}) = \text{Schéma des unités de } \mathcal{M}$$
$$SG(\mathcal{M}) = \text{......... de dét. } 1$$
$$PG(\mathcal{M}) = G(\mathcal{M}) / \text{Centre} \simeq G(\mathcal{M}) / \mathbb{G}_m$$

[unclear: Dans ce dictionnaire la donnée d'un $\rho$ correspond] à celle d'une section $\rho \in \mathcal{M}$, et

\newpage

%% ===== Page 259 =====

la condition de régularité signifie que la
sous-algèbre engendrée sur chaque fibre est
de rang $2$ (ce qui implique que $\rho$ engendre
une sous-algèbre quadratique

(10) $$ \mathcal{O}_S + \mathcal{O}_S \rho = A_\rho $$

qui est en fait un ordre dans $M$ comme
sous-algèbre, i.e. de formation commutante au
changement de base. Et on a :

(11) $$ T_\rho = \text{schéma des unités de } A_\rho $$
$$ ( \text{cas de } G = \mathbb{G}(M) ) $$

(11 bis) $$ T_\rho = \text{schéma des unités de } A_\rho $$
de norme égale à $1$
$$ (\text{cas de } G = S\mathbb{G}(M) ) $$

On voit que $T_\rho$ est un tore [unclear: (au moins sur une ext.)]
$\rho$ est semi-simple, i.e. semi-simple sur chaque fibre
ce qui équivaut, puisqu'il est régulier, au
fait qu'il a sur chaque fibre ses deux val. propres
bien distinctes. Pour

(12) $$ \rho = \begin{pmatrix} 0 & 1 \\ -1 & 1 \end{pmatrix} \quad (\rho^6 = 1, \rho^3 = -1) $$

ceci est le cas en car. $\neq 3$ — dans le
sens strict $T_\rho$ en tant que schéma en
groupes sur $\operatorname{Spec}(\mathbb{Z})$ n'est pas un tore — il

\newpage

%% ===== Page 260 =====

devient [unclear: produit] d'un gr. unipotent par un
autre $\mathbb{G}_m$ en car. 3. De même,
pour
$$ (13) \qquad \sigma = \sigma_\infty = \begin{pmatrix} 0 & -1 \\ 1 & 0 \end{pmatrix} \quad , \quad \sigma^4=1, \sigma^2=-1 $$
$\sigma$ est semi-simple en toutes car. sauf en car. 2,
où il est unipotent - et $T_\sigma$ est un
tore sauf en car. 2, où il devient un
produit d'un $\mathbb{G}_m$ par un unipotent.

[MARGIN: Remarque en général]
Quand on prend la [unclear: variante] dans
$S\mathbb{G}(M)$, on trouve
$$ S T_\rho \subset T_\rho \quad , \quad S T_\rho = \operatorname{Ker}(T_\rho \xrightarrow[\text{induite par le déterminant}]{\text{norme}} \mathbb{G}_m) $$

L'avantage de $T_\rho$ sur $ST_\rho$, c'est que
le torseur sous $T_\rho = \text{schéma des autos de } A_\rho$
est localement trivial au sens de Zariski, tandis
qu'il n'en est pas de même, a priori, pour
un torseur sous $ST_\rho$. Ainsi, quand
on se propose de résoudre l'équation en $\alpha, \beta$
$$ [\alpha, \beta] = \rho \quad, $$
pour $\rho$ donné satisfaisant les conditions de la
proposition, c'est possible (loc. Zari) dans
$\mathbb{G}(M)$, mais peut-être pas dans $S\mathbb{G}(M)$.

\newpage

%% ===== Page 261 =====

Revenons au cas particulier

$$ \rho = \begin{pmatrix} 0 & 1 \\ -1 & 1 \end{pmatrix} \quad , \quad A = \mathbb{Z} $$

alors

(14) $A_{\rho} =$ anneau des entiers du corps $\mathbb{Q}(\sqrt{-3})$, $\rho$ correspondant à
$$ \exp(2i\pi/6) = \frac{1+i\sqrt{3}}{2} $$

Quand on se donne $\xi = \begin{pmatrix} A & B \\ C & D \end{pmatrix}$, d'où $\xi \rho = \begin{pmatrix} -B & A+B \\ -D & C+D \end{pmatrix}$

les conditions b), c) de la proposition deviennent

$$ \det \xi = 1, \quad Tr(\xi \rho) = Tr \rho $$

(16) $$ AD-BC=1 \qquad C+D-B=1 $$

- conditions bien familières ! Elles impliquent la possibilité de trouver (sur une base quelconque) loc. Zar. un $\alpha \in \operatorname{Gl}(2)$
[MARGIN: [unclear] (*)]
qui ait

$$ [\alpha, \rho] = \xi , $$

et sur $\mathbb{Z}$ c'est possible globalement,
car l'obstruction est un $A_{\rho}$-Module inversible, or
$A_{\rho}$ est principal i.e.
(17) $$ Pic(A_{\rho}) = 1 $$

\newpage

%% ===== Page 262 =====

932

Mais si $\alpha$ est une solution, on a a priori

(18) $$det \alpha \in \mathbb{Z}^* = \{\pm 1\}$$

et toute autre solution sera de la forme

$\alpha' = \alpha f$, avec $f \in A_\rho^*$, d'où

$$det \alpha' = det \alpha . det f = det \alpha$$

car $$det f = f \bar{f} > 0 \quad \text{donc} \quad det f = + 1 \quad,$$

ainsi le signe dans (18) ne dépend pas
du choix de $\alpha$ ; s'il est $-1$, on ne peut
trouver de solution de l'équation (*) dans
$\operatorname{Sl}(2, \mathbb{Z})$. Ce signe ne dépend que de $\xi$ i.e. que
de $A, B, C, D$ - et en fait bien sûr que de $B, D$
(qui déterminent à eux deux $\xi$, grâce aux
formules (16) ) - les calculs faits par ailleurs
suggèrent que ce signe est aussi égal au
reste de $D \pmod 3$ - $1$ si $D \equiv 1 \ (3)$, $-1$ si
$D \equiv -1 \ (3)$.

Revenant au cas général, si on a trouvé
une solution

(19) $$[\alpha, \rho] = \alpha \rho(\alpha)^{-1} = \xi \quad \quad \alpha \in \Gamma G(M) \quad,$$

pour pouvoir en trouver une solution $\alpha'$ qui
soit $\in \Gamma S G(M)$ i.e. $det \alpha' = 1$, on note que
les solutions $\alpha' \in \Gamma G(M)$ sont données par
$$\alpha' = \alpha f \quad \quad f \in \Gamma \mathcal{O}_S^* \quad,$$

\newpage

%% ===== Page 263 =====

et la condition devient
$$det \alpha' = 1 \quad i.e. \quad det(\alpha f) = 1 \quad i.e.$$
$$det(f) = det \alpha$$
donc la condition néc. et suffisante pour l'existence
d'une solution, est que la section $det \alpha$ de
$\underline{O}_S^*$ soit ds l'image de l'homomorphisme

[MARGIN: NB $S=Spec A$]
(20) $$N_{A_\rho / A} : A_\rho^* \longrightarrow A^*$$
i.e. l'obstruction est dans le groupe

(21) $$A^* / N_{A_\rho/A} A_\rho^*$$
qui est un invariant particulièrement bien
connu dans la théorie des corps de classes !

[MARGIN: NB $N(x+y\rho)=...$]
Sauf erreur, il résulte de la théorie des corps
de classes que pour $\rho$ donné par (10),
et pour $A = \mathbb{Z}_l \quad (l \text{ nb premier})$, ce
(entiers $l$-adiques)
groupe quotient est nul dans tous les cas,
sauf pour $l=3$, où il est isomorphe à
$\mathbb{Z}/2\mathbb{Z}$. -- A [unclear] --

[MARGIN: [unclear: base] quelconque]
Notons que (tjs ds le cas particulier (12))

[MARGIN: ici]
(22) $$\alpha = \begin{pmatrix} a & b \\ c & d \end{pmatrix}$$
des calculs immédiats déjà faits montrent que

\newpage

%% ===== Page 264 =====

(23) $$(ad-bc) D = (c^2+cd+d^2)$$

donc $D$ est inversible ssi $c^2+cd+d^2 = \mathcal{N}_{A_\rho/A}(c\rho+d)$
l'est, i.e. ssi $c\rho+d$ est inversible , dans $A$.

alors 
$$ad-bc = (c^2+cd+d^2)/D$$
[MARGIN: cas $c^2+cd+d^2 \equiv 1 (3)$]
ce qui montre que le signe de $\det \alpha = ad-bc$
est alors égal à la valeur de $D \pmod 3$.
(Mais est-ce que $D$ peut être congru à $0 \pmod 3$ ?
alors mod 3 la matrice $\xi$ dépendrait
des seuls paramètres $A,B$
$\xi = \begin{pmatrix} A & B \\ 1+B & 0 \end{pmatrix} \implies -B(1+B) = 1$ i.e. $B^2+B+1 = 0$ donc $B \equiv 1$ (car $3|B-1$)
donc $\xi = \begin{pmatrix} A & 1 \\ -1 & 0 \end{pmatrix}$,
$\xi\rho = \begin{pmatrix} -1 & A+1 \\ 0 & -1 \end{pmatrix}$, cette matrice est régulière ssi
$A \neq -1$ , ok.) Donc a priori on ne peut
exclure le cas où $D$ peut prendre une valeur
nulle mod 3 ...)

Lorgnons pour la section $\sigma$ de $\operatorname{Gl}(2)_{\mathbb{Z}}$,
donnée par

(24) $$\sigma = \sigma_\infty = \begin{pmatrix} 0 & -1 \\ 1 & 0 \end{pmatrix} \qquad \sigma_\infty^4 = 1, \sigma_\infty^2 = -1$$

on a encore (sauf erreur)
[MARGIN: cas $A = \mathbb{Z}$]
(25) $$Pic(\Lambda_\sigma) = 0$$

La condition numérique sur 

(26) $$\eta = \begin{pmatrix} R & S \\ T & U \end{pmatrix}$$
pour [unclear] d'être locale dans l'image

\newpage

%% ===== Page 265 =====

de $F_\sigma$ (en plus de la condition de régularité
sur $\eta\sigma$)

(27) $\eta = \beta \sigma \beta \sigma^{-1} = [\beta,\sigma]$

c'est qu'on ait

(28) $det \, \eta = 1, \ Tr(\eta\sigma) = Tr\sigma$

or

$$ \eta\sigma = \begin{pmatrix} R & S \\ T & U \end{pmatrix} \begin{pmatrix} 0 & -1 \\ 1 & 0 \end{pmatrix} = \begin{pmatrix} S & -R \\ U & -T \end{pmatrix} $$

et (28) devient

(29) $RU-ST=1, \ S-T=0 \quad i.e. \ S=T$

i.e. $\eta$ doit être symétrique. [crossed out: Pour] Si cette
condition est satisfaite, il s'ensuit par (15)
- quand la base est $\mathbb{Z}$ - que l'on peut
trouver une solution $\beta \in \operatorname{Gl}(2, \mathbb{Z})$ de (27),
l'obstruction à trouver une solution dans
$\operatorname{Sl}(2,\mathbb{Z})$ est [unclear: mesurée par son signe], savoir
sa classe dans

$$ \mathbb{Z}^* / N(A^*_\sigma) \simeq \pm 1 $$

- comme on a, pour

(30) $\beta = \begin{pmatrix} r & s \\ t & u \end{pmatrix}$

[unclear: une matrice] satisfaisant (27), la relation

(31) $(ru-st) U = r^2+u^2 = N(r\sigma+u)$

\newpage

%% ===== Page 266 =====

434 \hfill (bien [unclear: entendu])

on sait comme tantôt que $U$ est [unclear: impair],
toute base [unclear: in.] peut l'être, et que dans ce
[MARGIN: car $N(k) \equiv 1 \; (4)$]
cas, l'obstruction $\rightarrow$ (dans $\mathbb{Z}$, ou $\mathbb{Z}/2\mathbb{Z}$) est
le reste de $U$ mod 4 (qui est égal à
$\pm 1$)...

Résumons pour mémoire

Proposition. Soit $A = \mathbb{Z}$, ou $\mathbb{Z}_{\ell}$, ou $\hat{\mathbb{Z}}$, et
considère les équations dans
$$ [\alpha, \rho] \overset{def}{=} \alpha \rho (\alpha)^{-1} = \xi $$
$$ [\beta, \sigma] = \beta \sigma (\beta)^{-1} = \eta $$
où $\rho = \begin{pmatrix} 0 & 1 \\ -1 & 1 \end{pmatrix}$, $\sigma = \begin{pmatrix} 0 & -1 \\ 1 & 0 \end{pmatrix}$,
$\xi$ et $\eta \in \operatorname{Gl}(2,A)$ donnés, & on cherche $\alpha$ (resp $\beta$)
dans $\operatorname{Gl}(2,A)$. On suppose que
$\xi = \begin{pmatrix} A & B \\ C & D \end{pmatrix}$ resp $\eta = \begin{pmatrix} R & S \\ T & U \end{pmatrix}$
satisfont les conditions nécessaires pour
l'existence d'une solution
$$ AD - BC = 1 \qquad B + D - C = 1 \qquad (16) $$
$$ RU - ST = 1 \qquad S - T = 0 \text{ i.e. } S = T \qquad (29) $$

Alors :
1) Pour qu'il existe une solution $\alpha$ (resp $\beta$) dans $\operatorname{Gl}(2, A)$
il faut & il suffit que $\xi \rho = \begin{pmatrix} -B & A+B \\ -D & C+D \end{pmatrix}$ soit
régulière, i.e. un [unclear]
(resp. que $\eta \sigma = \begin{pmatrix} S & -R \\ U & -T=-S \end{pmatrix}$ le soit) - et pour

\newpage

%% ===== Page 267 =====

ceci est en fait [unclear] que sur un [unclear] $\xi$ (resp $\eta$) soit un multiple scalaire de $\rho^{-1}$ (resp. $\sigma^{-1}$)
- et pour ceci il suffit que $D$ (resp. $U$) soit inversible, car le coeff. correspondant dans 
$$\rho^{-1} = \begin{pmatrix} -1 & -1 \\ 1 & 0 \end{pmatrix} \quad \text{resp} \quad \sigma^{-1} = \begin{pmatrix} 0 & 1 \\ -1 & 0 \end{pmatrix} \quad \text{est zéro} ).$$

2) Lorsque $D$ (resp $U$) est inversible, il y a tjs une solution $\alpha$ resp $\beta$, avec $\det \alpha \in \pm 1$ (resp $\det \beta \in \pm 1$).

3) Pour qu'il y ait une solution dans $\operatorname{Sl}(2, A)$, il faut et il suffit qu'on ait $D \equiv 1 \ (3)$ (resp. $U \equiv 1 \ (4)$) - condition qui est automatique si $A = \mathbb{Z}_{\ell}$ avec $\ell \neq 3$ (resp avec $\ell \neq 2$).

dans $\operatorname{Gl}(2)$ [unclear]

[MARGIN: Etude d'abord de l'équation]

## V) Etude de l'équation
(32) $$[\alpha, \rho] = [\beta, \sigma]$$

On suppose pour ça que $G$ soit une forme de $\operatorname{Gl}(2)$, associée à un alg. d'Azumaya $M$ sur la base $S = Spec \ A$, et qu'on a deux sections $\rho, \sigma$ régulières de $G$ i.e. deux sections $1 \to \dots$ de $M$. On se donne de plus un "sous-groupe à 1 paramètre" (additif) de $M$, i.e. hom.
(33) $$G: \mathbb{G}_a \to S\mathcal{G}, \quad \lambda \mapsto \xi(\lambda)$$

\newpage

%% ===== Page 268 =====

435

Le cas qui nous intéresse plus particulièrement
est celui où $G = \operatorname{Gl}(2)_{\mathbb{Z}}$, $\rho$ et $\sigma$ comme
dessus, et

(34) $\varepsilon_1(\lambda) = \varepsilon_1^\lambda = \begin{pmatrix} 1 & 0 \\ \lambda & 1 \end{pmatrix}$ .

On notera qu'on a

(35) $\varepsilon_1 = \varepsilon_1(1) = \begin{pmatrix} 1 & 0 \\ 1 & 1 \end{pmatrix} = $ [unclear]

Nous considérons (32) comme des
équations reliant les quantités
toutes variables
$\alpha, \beta, \mu$ - $\alpha, \beta$ sections de $G$,
$\mu$ section de [unclear] - les hypothèses que
nous allons faire impliqueront d'ailleurs
que $\mu$ sera en fait une section de
$\mathbb{G}_m$ (s'il satisfait $\alpha, \beta$ : la
relation (32)).

On commencera par étudier
les syst. de solutions $(\xi, \eta, \mu)$ satisfaisant

(33) $\xi = \eta \varepsilon_1(\mu-1)$
(34) $\det \xi = \det \eta = 1$
[MARGIN: NB On aura $\det \varepsilon_1(\mu-1) = 1$ donc $\det \xi = \det \eta$ [unclear]]
(35) $Tr \, \xi = Tr \, \rho \quad Tr \, \eta = Tr \, \sigma$

On peut éliminer $\eta$ de ces équations
grâce à (33), et trouve

\newpage

%% ===== Page 269 =====

(36) $$det \xi = 1, \quad Tr \xi \rho = 0, \quad Tr((\xi \varepsilon_1 (1-\mu) - 1)\sigma) = 0$$.

posons pour simplifier
$$i = Tr \rho \qquad j = Tr \sigma$$

La troisième relation (36), qui en développant par rapport aux coefficients de $\xi$, peut s'écrire
(37) $$P\mu - Q = 0$$

Les coefficients $P, Q$ sont des fonctions linéaires affines par rapport aux composantes de $\xi$, et données par l'identité (en $\xi, \mu$)
(38) $$Tr((\xi \varepsilon_1(1-\mu) - 1)\sigma) = P(\xi)\mu - Q(\xi)$$

Le terme constant $-Q(\xi)$ est égal (s'obtient en prenant la valeur du premier membre de (38) pour $\mu = 0$)
i.e.
(39) $$-Q(\xi) = Tr((\xi \varepsilon_1 - 1)\sigma)$$
$$= Tr(\xi \varepsilon_1 \sigma) - Tr \sigma$$

On suppose qu'on a [MARGIN: $\lambda$ section centrale de $G$ i.e. $\lambda$ scalaire inversible]
(40) $$\varepsilon_1 \sigma = - \lambda \rho$$
de sorte que
(40 bis) $$\varepsilon_1 = - \lambda \rho \sigma^{-1}$$ [MARGIN: NB Si $det \rho = det \sigma = 1$, alors on doit avoir $\lambda^2 = 1$ donc si $A$ intègre, $\lambda \in \{ \pm 1 \}$]
et (39) devient,
$$Tr(\xi \varepsilon_1 \sigma) = Tr(\xi(-\lambda \rho)) = - \lambda Tr \xi \rho$$ [unclear]

\newpage

%% ===== Page 270 =====

435

on trouve

(41) $$Q(\xi) = \underbrace{\lambda_0 \text{Tr} \rho + \text{Tr} \sigma}_{c} \quad (\text{constante, indépendante de } \xi)$$

et on supposera

(42) $$c = \lambda_0 \text{Tr} \rho + \text{Tr} \sigma \quad \text{scalaire inversible dans } \underline{O}_S$$

(Dans le cas qui nous intéresse, on a
$\sigma^{-1} = -\sigma$, $\varepsilon_1 = \rho \sigma = - \rho \sigma^{-1}$, i.e. $\lambda_0 = 1$
et $c = \lambda_0 \text{Tr} \rho + \text{Tr} \sigma = \underset{1}{\text{Tr} \rho} + \underset{0}{\text{Tr} \sigma} = 1 \qquad ok. \quad)$

Ceci posé, la troisième équation (36) devient (compte
tenu des deux premières)

(43) $$P(\xi) \mu = c$$

ce qui signifie aussi que $P(\xi)$ est
inversible, et

(44) $$\mu = c / P(\xi) \quad.$$

Ainsi, $\mu$ est déterminé de façon unique
par $\xi$, par l'équation (44) - et
$\mu$ sera nécessairement inversible.

En résumé :

Proposition Si on suppose satisfaite (40 bis)
pour l'élément unipotent $\varepsilon_1 = \varepsilon_1(\underline{1})$, et que l'on
ait (42) ($\lambda_0 \text{Tr} \rho + \text{Tr} \sigma$ inversible) , alors

\newpage

%% ===== Page 271 =====

l'ens des solutions $(\xi,\eta,\mu)$ des

équations simultanées (33)(34)(35)

est corr. 1-1 avec l'ens des $\xi$

sections de $G=G(M)$ satisfaisant les

conditions

(45) $$ \begin{cases} \det \xi = 1, \text{Tr } \xi\rho = \text{Tr}\rho \ , \text{ et} \\ P(\xi) \text{ section inversible de } \underline{O}_S \ . \end{cases} $$

En termes de $\xi$ , la solution corr. du

syst. (33 à 35) est donnée par

$$ (\xi, \eta = \xi \varepsilon_1(1-\mu), \mu = c/P(\xi) ) $$

- et $\mu$ (le "multiplicateur") est

nécessairement inversible.

NB Si on avait résolu en $\eta$ , on aurait

trouvé une relation de la forme

$$ P(\eta)\mu - Q'(\eta) = 0 $$

où

$$ Q'(\eta) = \text{Tr } \eta \varepsilon_0^{-1} \rho - \text{Tr}\rho $$

et la condition naturelle serait

$$ \varepsilon_0^{-1}\rho = - \lambda_0' \sigma^- \quad (\lambda_0' \text{ scalaire inv.}) \ , $$

qui est équivalente à (40 bis) avec $\lambda_0' = 1/\lambda_0$ ,

et implique que

$$ Q'(\eta) = \text{cte} = \lambda_0' \text{Tr }\sigma - \text{Tr}\rho = c' \ , $$

\newpage

%% ===== Page 272 =====

435

on voit tout de suite que
$$ c' = \lambda'_0 c $$
donc la condition $c$ inv. [unclear: équivaut] à $c'$
inv., et on a dans le système d'équations
en $\eta$
$$ det \eta = 1, Tr \eta \sigma = Tr \sigma , P'(\eta) \mu = c' \text{ i.e. } \mu = c' / P'(\eta) $$
[MARGIN: $P'(\eta)$ inv.]

On voit donc que les hypothèses faites sont en
fait symétriques en $\xi, \eta$, et permettent
un traitement symétrique de la résolution,
soit via $\xi$, soit via $\eta$.

Détermination de $P(\xi)$. On aura
$$ \varepsilon_1(\lambda) = 1 + N_1 \lambda = (1+N_1)^{\lambda} $$
$N_1$ une matrice de carré nul
$$ N_1^2 = 0 , $$
et
$$ Tr(\xi \varepsilon_1(1-\mu)\sigma) - Tr \sigma = Tr \xi \sigma - Tr \sigma + (1-\mu) Tr \xi N_1 \sigma $$
donc
(47) $$ P(\xi) = Tr \xi (- N_1 \sigma) $$
[MARGIN: sous $- N_1 \sigma$ il est écrit: $N_1 \rho = \rho + \sigma$]

D'ailleurs la formule [unclear: donne]
(48) $$ \varepsilon_1 = 1 + N_1 $$
et la formule (40 bis) s'écrit
$$ \sigma + N_1 \sigma = - \lambda_0 \rho^{-1} $$
(49) $$ N_1 = - ( \lambda_0 \rho^{-1} \sigma^{-1} + 1 ) $$
soit en [unclear]

\newpage

%% ===== Page 273 =====

$$N_1\sigma = - (\lambda_0 \rho + \sigma)$$

donc on a

$$(50) \quad P(\xi) = \mathrm{Tr} \, \xi (\lambda_0 \rho + \sigma)$$

Corollaire de la Proposition

Le schéma des [MARGIN: en $\xi, y, \mu$] solutions simultanées des équ. (33), (35) est un schéma relatif lisse, de dimension relative 2 - car l'équation $\mathrm{Tr}(\xi\rho) = \text{cte}$ définit toujours un sous-schéma lisse de dim 1 de $\mathcal{SL}(\mathcal{M})$ (dès que $\rho$ est une section qui n'est nulle part [unclear]), - dont il faut prendre ensuite l'ouvert défini par "$P(\xi)$ inversible".

Dans le cas qui nous intéresse,

$$ \rho = \begin{pmatrix} 0 & 1 \\ -1 & 1 \end{pmatrix} \quad , \quad \sigma = \begin{pmatrix} -1 & -1 \\ 1 & 0 \end{pmatrix} , $$

la condition $P(\xi)$ inv. s'écrit

posant

$$ \xi = \begin{pmatrix} A & B \\ C & D \end{pmatrix} \quad , \quad y = \begin{pmatrix} R & S \\ T & U \end{pmatrix} $$

on a

$$ P(\xi) = D \quad , \quad P(y) = U $$

$$(51) \quad y \varepsilon_1(\mu-1) = \begin{pmatrix} R + (\mu-1)B & S \\ T + (\mu-1)A & U \end{pmatrix}$$

\newpage

%% ===== Page 274 =====

338

et l'équation
$$ \xi = \eta \varepsilon_1(\mu-1) $$
implique
(52) $$ B = S \ , \quad D = U \ , $$
et en fait elle équivaut à l'une de ces
deux équations et à
(53) $$ \mu D = 1 \quad (\implies \mu U = 1) $$
quand on suppose déjà que $\xi, \eta$ satisfont
à l'équation des déterminants et des traces
(54) $$ \begin{matrix} AD - BC = 1 \\ A+D - B = 1 \end{matrix} \qquad \begin{matrix} RU - ST = 1 \\ R+U - S = 1 \end{matrix} $$

Le sous-schéma des solutions des
équations simultanées (33), (35),
est donc celui des $(D,B)$ avec $D$
inv., ou encore des $(S,U)$ avec
$U$ inv.. On peut l'identifier
au schéma des matrices
(55) $$ \begin{pmatrix} \mu & M \\ 0 & 1 \end{pmatrix} = \begin{pmatrix} 1/D & B/D \\ 0 & 1 \end{pmatrix} $$

(cf p. 410 - ), qui normalisent le
$$ \begin{pmatrix} 1 & -\lambda \\ 0 & 1 \end{pmatrix} = \rho^{-1} \varepsilon_1(\lambda) \rho \ . $$

\newpage

%% ===== Page 275 =====

Revenons au cas général, où on suppose néanmoins que $\varepsilon_1$ donné par (40 bis),
et posons

$$ (56) \qquad \varepsilon_0 = \rho^{-1} \varepsilon_1 \rho = -\lambda_0 \sigma^{-1} \rho = 1 + \mathcal{N}_0 $$
$$ (57) \qquad \mathcal{N}_0 = \rho^{-1} \mathcal{N}_1 \rho = -(\lambda_0 \sigma^{-1} \rho + 1) $$

de sorte que $\mathcal{N}_0^2 = 0$, et

$$ (58) \qquad \lambda \mapsto \varepsilon_0(\lambda) \overset{df}{=} 1 + \lambda \mathcal{N}_0 $$

est un sous-groupe à 1 paramètre de $St_{\mathcal{B}}(M)$, soit $L_0$, dont le normalisateur dans $\operatorname{Gl}(M)$ sera noté $B_0$. On a deux hom.

$$ B_0 \xrightarrow{\chi(\underset{\sim}{u}) = \det(\underset{\sim}{u})} \mathbb{G}_m $$
$$ (59) \qquad B_0 \xrightarrow{\chi'(\underset{\sim}{u}) \text{ donné par }} \underset{\sim}{u}(\varepsilon_0(\lambda)) = \varepsilon_0(\chi'(\underset{\sim}{u})\lambda) $$

[MARGIN: je souligne $\underset{\sim}{u}$ par un trait sinueux, pour le distinguer des scalaires (t.q. $\lambda$ !)]

(en termes linéaires dans $M$)
$$ (60) \qquad \underset{\sim}{u} \mathcal{N}_0 \underset{\sim}{u}^{-1} = \chi'(\underset{\sim}{u}) \mathcal{N}_0 $$

Et on considère le sous-groupe
$$ B_0^* \subset B_0 \quad \text{ défini par } \chi(\underset{\sim}{u}) = \chi'(\underset{\sim}{u}) $$
$$ (61) \qquad \text{ie par } \quad \underset{\sim}{u} \mathcal{N}_0 \underset{\sim}{u}^{-1} = \det(\underset{\sim}{u}) \mathcal{N}_0 $$

[MARGIN: il faudrait of course vérifier [unclear], cf plus bas]

$L_0 \subset B_0^*$ et une suite exacte de sch. en gr.
$$ 1 \to L_0 \to B_0^* \xrightarrow{\det} \mathbb{G}_m \to 1 $$

\newpage

%% ===== Page 276 =====

439

Soit $\underline{u}$ une section de $\tilde{B}_0^*$, j'ai envie de prouver que si on pose

$$ (63) \qquad \xi = [\underline{u}, \rho], \quad \eta = [\underline{u}, \sigma] $$

on aura automatiquement la relation

$$ (64) \qquad \xi = \eta \varepsilon_1(\mu-1) \qquad \text{où } \mu = \det \underline{u} $$

Pour nous en rassurer, prenons d'abord $\underline{u}$ dans $L_0$, par ex. 

$$ \underline{u} = \varepsilon_0(\lambda) $$

d'où 

$$ [\underline{u}, \rho] = \varepsilon_0(\lambda) \rho \varepsilon_0(\lambda)^{-1} \rho^{-1} = \varepsilon_0(\lambda) \varepsilon_1(-\lambda) $$

$$ [\underline{u}, \sigma] = \varepsilon_0(\lambda) \sigma \varepsilon_0(\lambda)^{-1} \sigma^{-1} $$

Or

$$ \sigma \varepsilon_0(-\lambda) \sigma^{-1} = 1 + (\sigma N_0 \sigma^{-1}) \lambda $$

[unclear: Et compte tenu] de (57) on a

$$ \sigma N_0 \sigma^{-1} = -(\lambda_0 \rho \sigma^{-1} + 1) \stackrel{(49)}{=} N_1 $$

donc

$$ (65) \qquad \begin{cases} \sigma \varepsilon_0(\lambda) \sigma^{-1} = \varepsilon_1(\lambda) \\ \rho \varepsilon_0(\lambda) \rho^{-1} = \varepsilon_1(\lambda) \end{cases} $$
[MARGIN: et par suite]

donc on trouve bien

$$ [\underline{u}, \rho] = [\underline{u}, \sigma] = \varepsilon_0(\lambda) \varepsilon_1(-\lambda) $$

i.e. (64), puisque ici $\mu = \det \underline{u} = 1$, donc $\mu - 1 = 0$.

275

\newpage

%% ===== Page 277 =====

Nous avons vu que $\rho$ et $\sigma$ tous les
deux transforment $N_0$ sur $N_1$, donc le
sous-groupe à 1 paramètre additif
$$ L_0 = \{ \underset{= 1 + \lambda N_0}{\varepsilon_0(\lambda)} \mid \lambda \text{ dans } \mathbb{G}_a \} $$
dans le sous-groupe analogue
$$ L_1 = \{ \underset{= 1 + \lambda N_1}{\varepsilon_1(\lambda)} \mid \lambda \text{ dans } \mathbb{G}_a \} , $$
et donc $\tilde{B}_0^*$ dans $\tilde{B}_1^*$ (défini au
moyen de $L_1$, ou de $N_1$, comme
$\tilde{B}_0^*$ en termes de $L_0$). Donc la
formule (64) s'écrit alors
$$ \rho \underline{u} \rho^{-1} = \sigma \underline{u} \sigma^{-1} \varepsilon_1(\mu-1) $$
ou aussi, en remplaçant $\underline{u}$ par $\underline{u}^{-1}$
$$ (65) \quad \rho \underline{u} \rho^{-1} = \sigma \underline{u} \sigma^{-1} \varepsilon_1(\frac{1}{\mu}-1) $$
Or
$$ \rho \underline{u} \rho^{-1} = int(\lambda) (\sigma \underline{u} \sigma^{-1}) $$
où $\lambda = \rho \sigma^{-1}$, et la formule (65) s'écrit
$$ (**) \quad int(\lambda)(U) = U . \varepsilon_1(\frac{1}{\mu}-1) $$
[MARGIN: pour $U$ dans $\tilde{B}_1^*$ et $\mu = \det U$ ;]

et en vertu de (40 bis) $int(\lambda) = int(\varepsilon_1)$, donc
la formule se réduit à
$$ int(\varepsilon_1)(U) = U . \varepsilon_1(\frac{1}{\mu}-1) $$

\newpage

%% ===== Page 278 =====

440

on aura

[crossed out: $(1+N_1)v(1-N_1) = v(1+N_1(\frac{1}{\mu}-1))$]

soit 
$$v + N_1v - vN_1 - N_1vN_1 = v + vN_1(\frac{1}{\mu}-1)$$
$$\varepsilon_1 v \varepsilon_1(-1) = v \varepsilon_1(\frac{1}{\mu}-1) \text{ soit}$$
$$\underset{1+N_1}{\varepsilon_1} v = v \underset{1+\frac{1}{\mu}N_1}{\varepsilon_1(\frac{1}{\mu})} \text{ soit } v + N_1 v = v + \frac{1}{\mu} v N_1$$

soit enfin
$$v N_1 v^{-1} = \mu N_1 \quad (= \det(v).N_1)$$

ce qui n'est autre que la définition de
$\tilde{B}_1^\gamma$ !

On trouve donc un morphisme
$$ (66) \quad \left\{ \begin{array}{l} \tilde{B}_1^\gamma \longrightarrow \mathcal{F}_{\rho,\sigma} \\ \underline{u} \longmapsto ([\underline{u}, \rho], [\underline{u}, \sigma], \det \underline{u}) \end{array} \right. $$

de $\tilde{B}_1^\gamma$ dans l'ensemble des solutions du
système (33) à (35) , est-ce un isomor-
phisme ?
Notons que l'égalité
$$[\underline{u}, \rho] = [\underline{u}', \rho] \text{ signifie que } \underline{u}' \underline{u}^{-1} \in \Gamma_\rho$$
[crossed out: donc que (66) est injectif]
signifie simplement que
$$\Gamma_\rho \cap \tilde{B}_1^\gamma = \{1\}$$
en fait l'analyse faite [unclear: justifie notre]

\newpage

%% ===== Page 279 =====

que cette intersection est égale à
: $T_0 \cap \tilde{B}_0^{\gamma}$, donc contenue dans

$T_0 \cap T_{\rho}$. Je vais supposer

(68) $$T_0 \cap T_{\rho} = \mathbb{G}_m (= Cent \, \mathfrak{S}(\mathcal{M})$$
[unclear: est $\simeq$ il omet \dots ] $)$

- il suffira pour ceci de supposer qu'en
degrés [unclear: fixés], $s, \sigma, \rho$ sont linéairement indé-
pendants - c'est le cas bien sûr dans le
cas qui nous intéresse. Mais donc
on a

$$T_0 \cap T_{\rho} \cap \tilde{B}_0^{\gamma} \subset \mathbb{G}_m \cap \tilde{B}_0^{\gamma},$$

mais

$$\mathbb{G}_m \cap \tilde{B}_0^{\gamma} = \mu_2$$

car si $\lambda$ est un [unclear: scalaire] local
dans $\mathbb{G}_m$, alors

$$\lambda N_0 \lambda^{-1} = N_0 = 1 . N_0$$

mais $\det \lambda = \lambda^2$, et la condition que
$\lambda$ [crossed out: satisfasse \dots] [unclear: \dots relative aux] $(B_0)$
soit dans $\tilde{B}_0^{\gamma}$, c'est bien $\lambda^2 = 1$.

On trouve donc ainsi

(69) $$T_0 \cap \tilde{B}_0^{\gamma} = T_{\rho} \cap \tilde{B}_0^{\gamma} = \mu_2$$

et le morphisme (66) n'est pas un
monomorphisme, mais se factorise

\newpage

%% ===== Page 280 =====

441

en un monomorphisme

(70) $$\tilde{B}_0 / \mu_2 \hookrightarrow \tilde{X}_{p,\sigma}$$

A vrai dire, on voit que $\tilde{B}_0$ n'est pas le
"bon" groupe à introduire ici - dans le
cas particulier $\varepsilon_0 = \begin{pmatrix} 1 & -1 \\ 0 & 1 \end{pmatrix}$, on a

$$B_0 = \{ \begin{pmatrix} p & e \\ 0 & q \end{pmatrix} \mid p, q \text{ inv.} \}$$

et la relation $u \varepsilon_0 u^{-1} = \varepsilon_0^{\det u}$ s'écrit, par un
calcul immédiat

(71) $$p/q = p q \quad \text{i.e.} \quad q^2 = 1$$

tandis que le "bon" groupe qui nous intéresse
est défini par

$$B_0' = \{ \begin{pmatrix} p & e \\ 0 & 1 \end{pmatrix} \mid p \text{ inv.} \} ,$$

donc avec $q=1$. Pour donner la description
dans le cas général, on note que la
fonction déterminant sur $\tilde{B}_0$ se décompose
intrinsèquement en deux facteurs

(72) $$\det \underline{u} = p(\underline{u}) q(\underline{u}) ,$$

qu'on peut distinguer l'un de l'autre
justement par la relation

(73) $$\underline{u} N_0 \underline{u}^{-1} = \frac{p(\underline{u})}{q(\underline{u})} N_0$$

et la condition qui va nous dicter le
"bon" groupe n'est pas (71), mais bien

(74) $$q = q(\underline{u}) = 1 .$$

\newpage

%% ===== Page 281 =====

qui définit donc un groupe qui est libre
(sur un car. 2 !) et s'insère cette fois-ci
dans une suite exacte canonique

$$(74) \quad 1 \longrightarrow L_0 \longrightarrow B'_0 \xrightarrow{det} \mathbb{G}_m \longrightarrow 1$$

On a [unclear]

$$(75) \quad B'_0 \cap G_m = \{ 1 \}$$

ce qui implique que le morphisme
induit par (56)

$$(76) \quad B'_0 \longrightarrow \mathcal{X}_{p,\sigma}$$

est un monomorphisme.

Pour montrer que c'est un isomorphisme [MARGIN: iso, il suffit], il suffit
maintenant, en vertu des principes généraux,
de montrer que c'est bijectif sur les
pts : i.e. sur un corps alg clos (on se
pourra simplifier la vie). Je laisse
tomber le détail des vérifications.

Rem. On peut considérer, pour la condition
(40 bis), que l'image d'une section
$\varepsilon_1 = 1 + t \bar{N}_1 \quad (N_1^2 = 0)$
du morphisme de $StA(M)$, et d'une section
$\rho$, le groupe $\sigma$ ... [unclear] déduit par [unclear]

\newpage

%% ===== Page 282 =====

442

(77) $\sigma = - \varepsilon_1^{-1} \rho ( = -(1-N_1)\rho$ i.e. $\sigma = -\rho + N_1 \rho$, i.e. $\rho + \sigma = N_1 \rho)$

(je laisse tomber la constante $\lambda_0$, ou cela revient à dire que changer $\rho, \sigma$ par des multiples scalaires, car $F_\rho$ et $F_\sigma$ n'est pas changé...) La condition (48) devient alors
$$c = Tr \rho + Tr \sigma = Tr (1-\varepsilon_1^{-1})\rho = Tr N_1 \rho \in \Gamma \underline{O}_S^*$$
(section inversible de $\underline{O}_S$). On suppose les $\rho, \sigma$ sont linéairement indépendantes dans chaque fibre, il vient au vu des dim. que $1, \rho$, et $N_1 \rho$ le sont, ou encore que $\rho^{-1}, 1$ et $N_1$ le sont, ou que $1, \sigma, N_1 \sigma$ le sont, ou $\sigma^{-1}, 1, N_1$. En résumé -

Proposition. Soit $M$ une alg. d'Azumaya sur le schéma de base $S$, $N_1$ une section de trace nulle, i.e. $\varepsilon_1 = 1+N_1$ une section unipotente, $\rho, \sigma$ des sections reliées par

(78) $\varepsilon_1 = -\rho \sigma^{-1}$ i.e. $\sigma = - \varepsilon_1^{-1} \rho$ i.e. $\rho + \sigma = N_1 \rho$ i.e. $\sigma + \rho = - N_1 \sigma$

ce qui détermine donc chacune des trois sections $\varepsilon_1, \rho, \sigma$ (et $\varepsilon_1 \iff N_1$) en fonction des deux autres, et par des [unclear: conditions] arbitraires $\rho$ et $\sigma$ (i.e. $\rho$ et $N_1$ de trace nulle) -

[MARGIN: Je prends la constante qui intervient dans la condition [unclear] égale à $-1$ (au lieu de $-\lambda_0$), ce qui donne $\varepsilon_1 = -\rho \sigma^{-1}$]

[MARGIN: N.B. $\rho + \sigma = N_1 \rho = N_0$
$= -N_1 \sigma = -\sigma N_0$
(78 bis)
Quitte à multiplier l'une des sections (p. ex $\sigma$) par $-1$, on a $\rho - \sigma = 1$ (c.a.d. $\rho + \sigma' = 1$), on ramène ça à la situation des projecteurs...]

\newpage

%% ===== Page 283 =====

supposer de plus (29) $1, \rho, \sigma$ [unclear: sys. de générateurs] en
chaque fibre i.e. $1, \rho^{-1}, N_1$ lin.
indép. en chaque fibre i.e. $1, \sigma^{-1}, N_0$ lin. indép.
en chaque fibre, et (30) $Tr N_1 \rho$ inversible.
(N.B. $N_1 \rho = - N_0 \sigma = \rho + \sigma$). Ceci posé, on aura
$int(\rho^{-1}) \varepsilon_1 = int(\sigma^{-1}) \varepsilon_1$, soit $\varepsilon_0$
$\rho \varepsilon_0 \rho^{-1} = \sigma \varepsilon_0 \sigma^{-1} = \varepsilon_1$

Soit $B_0$ le normalisateur de $\varepsilon_0$
$B'_0 \subset B_0$ le centralisateur de $\varepsilon_0$
défini par (74), (on a $\varepsilon_0 = 1 + N_0$)

$\varepsilon_0(\lambda) = 1 + \lambda N_0 \quad , \quad \varepsilon_1(\lambda) = 1 + \lambda N_1$
(Sous-groupes à $1$-paramètres additifs
dans $\operatorname{Gl}(M)$ (ou $\operatorname{Aut}(M)$))
et soit $L_0 =$ sous-schéma des groupes $\varepsilon_0(\mathbb{G}_a)$
$L_1 = \varepsilon_1(\mathbb{G}_a)$

(donc $\rho(L_0) = \sigma(L_0) = L_1$)

Soit $B'_0$ le sous-schéma en groupes de
de $B_0$ défini par (74), de sorte que $B_0$
s'insère dans une suite exacte
$$1 \longrightarrow L_0 \longrightarrow B_0 \xrightarrow{det} \mathbb{G}_m \longrightarrow 1$$
et que $u \in \Gamma(B_0)$ implique
$u(N_0) = det(u) N_0$, i.e.
$u(\varepsilon_0(\lambda)) = \varepsilon_0(det(u) \lambda)$.

Ceci posé, considérons le morphisme
$$B'_0 \longrightarrow G(M) \times G(M) \times E'$$
$$u \longmapsto ([u, \rho], [u, \sigma], det(u)).$$

[MARGIN: NB 29) est équivalent à dire que (30) $Tr N_1 \rho$ est inversible. * (car $\varepsilon_1 = \begin{pmatrix} 1 & 1 \\ 0 & 1 \end{pmatrix}$, $\rho = \begin{pmatrix} a & b \\ c & d \end{pmatrix}$ ... $N_1 = \begin{pmatrix} 0 & 1 \\ 0 & 0 \end{pmatrix}$ $N_1 \rho = \begin{pmatrix} c & d \\ 0 & 0 \end{pmatrix}$ donc $Tr N_1 \rho = c$ et un unipotent $\neq 1$ ... $Tr N_1 \rho = 1$]

\newpage

%% ===== Page 284 =====

443

Celui-ci prend ses valeurs dans $\mathcal{X}_{\rho, \sigma}$, le schéma des solutions des équations simultanées (33) à (35), et induit un iso

(79)
$$B'_0 \overset{\sim}{\longrightarrow} \mathcal{X}_{\rho, \sigma} \ .$$

[unclear: crossed out] Rappelons-le pour mémoire

Corollaire  La projection $\mathcal{X}_\rho \to \mathbb{G}(M)$
$$(\xi, \eta, \mu) \mapsto \xi$$
est un mono, et l'image est formée des $\xi$ satisfaisant
$$\det \xi = 1, \ Tr \xi \rho = Tr \rho, \ Tr \xi (\rho + \sigma) \ \text{inv.}$$
[unclear: $-M_\rho = N_{\rho \xi} = \rho N_0 - \sigma N_0$]

De même la projection $\mathcal{X}_\rho \to \mathbb{G}(M)$
$$(\xi, \eta, \mu) \mapsto \eta$$
est un mono, et l'image est formée des $\eta$ satisfaisant
$$\det \eta = 1, \ Tr \eta \sigma = Tr \sigma, \ Tr \eta (\rho + \sigma) \ \text{inv.}$$

Enfin, on a
(80)
$$Tr \xi (\rho + \sigma) = Tr \eta (\rho + \sigma) = 1/\mu \ .$$

Scholie  Par transport de structures, on trouve une structure de groupe sur $\mathcal{X}_{\rho, \sigma}$.

\newpage

%% ===== Page 285 =====

En l'explicitant, on trouve

si $\underline{u}$ correspond à $\xi, \eta, \mu$

$\underline{u}'$ — $\xi', \eta', \mu'$

$\underline{u}'' = \underline{u}' \underline{u}$ — $\xi'', \eta'', \mu''$

on aura

(81) $\mu'' = \mu' \mu$, $\xi'' = u'(\xi)\xi'$, $\eta'' = u'(\eta)\eta'$

où on a posé, comme d'habitude
$$u'(x) = \underline{u}' x \underline{u}'^{-1}$$
pour deux éléments $x, \underline{u}'$ d'un groupe.

III) Etude de l'équation [MARGIN: (Suite)]

(82) $\alpha \rho(\alpha)^{-1} = \beta \sigma(\beta)^{-1} \varepsilon_1(\mu-1)$
$(\alpha, \beta \in \mathbb{G}(\mathcal{M}), \mu \in \Gamma \mathbb{S}'_S)$

i.e.
(82bis) $[\alpha, \rho] = [\beta, \sigma] \varepsilon_1(\mu-1)$ .

Soit $G_{\rho, \sigma}$ le schéma des solutions, on a
donc un morphisme

(83) $G_{\rho, \sigma} \longrightarrow X_{\rho, \sigma}$
$(\alpha, \beta, \mu) \longmapsto (\alpha \rho(\alpha)^{-1}, \beta \sigma(\beta)^{-1}, \mu)$

et il résulte du th. précédent que c'est
un épimorphisme, plus précisément

\newpage

%% ===== Page 286 =====

On a une section canonique

(84)
\begin{tikzcd}
\mathcal{G}_{\rho,\sigma} \ar[r] & \mathcal{X}_{\rho,\sigma} \\
& B'_0 \ar[u]
\end{tikzcd}
$$(\underline{u}, \underline{u}, \det \underline{u}) \longleftarrow \underline{u} \longmapsto ([\underline{u},\rho], [\underline{u},\sigma], \det \underline{u})$$

Cette section prend ses valeurs dans la « diagonale », [unclear: formée par ce qui] correspond aux couples $(\underline{u}, \mu)$ (avec $\underline{u} \in \Gamma \mathcal{M}^*, \mu \in \Gamma E'$) tels que

(85) $$[\underline{u}, \rho] = [\underline{u}, \sigma] \varepsilon_1(\mu-1) \quad ,$$

qu'on tâchera de déterminer ainsi.

Conjuguant les résultats de (I) et (II) on voit que la solution générale de (82) sera fournie par

$$(\underline{u}\underline{a}, \underline{u}\underline{b}, \mu = \det \underline{u})$$

où $\underline{u}$ est une section de $B'_0$, $\underline{a}$ une section locale de $T_\rho$, $\underline{b}$ une section de $T_\sigma$.

De façon précise

Th. Considérons

(86)
\begin{tikzcd}
B'_0 \times T_\rho \times T_\sigma \ar[r] & G \times G \times E' \\
(\underline{u}, \underline{a}, \underline{b}) \ar[r, maps to] & (\underline{u}\underline{a}^{-1}, \underline{u}\underline{b}^{-1}, \det \underline{u})
\end{tikzcd}

\newpage

%% ===== Page 287 =====

[unclear] prend ses valeurs dans $\mathcal{G}_{p, \sigma}$ ,
et est même un isomorphisme

$$ (87) \quad B'_0 \times T_p \times T_\sigma \overset{\sim}{\longrightarrow} \mathcal{G}_{p, \sigma} $$

De façon plus précise, soit
$(\alpha, \beta, \mu)$ une section de $\mathcal{G}_{p, \sigma}$ , alors
il existe d'uniques sections
$\underline{a}$ de $T_p$ , $\underline{b}$ de $T_\sigma$ , telles que l'on ait

$$ (88) \quad \alpha \underline{a} , \beta \underline{b} \in \Gamma B'_0 $$

et on aura alors nécessairement
$$ \alpha \underline{a} = \beta \underline{b} \ (= \underline{u}) $$
d'où bien sûr
$$ ([\alpha, p], [\beta, \sigma]) = ([\underline{u}, p], [\underline{u}, \sigma]) $$
et l'unique $\mu$ qui satisfait à (82)
est donc $\det \underline{u}$ , d'où
$$ (89) \quad (\alpha, \beta, \mu) = (\underline{u} \, \underline{a}^{-1}, \underline{u} \, \underline{b}^{-1}, \det \underline{u}) \text{ (de façon unique)} $$

Dém. Il reste seulement à expliciter
l'unicité de $\underline{a}, \underline{b}$ satisfaisant
(88), ce qui est conséquence de
$$ (90) \quad T_p \cap B'_0 = T_\sigma \cap B'_0 = \{1\} $$
qu'on avait omis d'expliciter dans (II).

\newpage

%% ===== Page 288 =====

445 suite

En termes de la donnée de $\alpha$, on aura donc un $\underline{a}$ uniquement déterminé par la condition $\alpha \underline{a} \in \Gamma B'_0$, et on aura

(91) $$ \underline{a} = c_0 \rho + d_0 $$ 
($c_0, d_0 \in \Gamma(\underline{\mathcal{O}_S})$ déterminés par $\alpha$)

de même on aura

(92) $$ \underline{b} = t_0 \sigma + u_0 $$ 
($t_0, u_0 \in \Gamma(\underline{\mathcal{O}_S})$ déterminés par $\beta$)

de telle façon donc qu'on aura

(93) $$ \alpha \underline{a} = \beta \underline{b} $$ 
($\overset{\text{déf}}{=} \underline{u}$)

Dans le cas numérique
$$ \rho = \begin{pmatrix} 0 & 1 \\ -1 & 1 \end{pmatrix} , \quad \sigma = \begin{pmatrix} 0 & -1 \\ 1 & 0 \end{pmatrix} $$
posant
$$ \alpha = \begin{pmatrix} a & b \\ c & d \end{pmatrix} , \quad \beta = \begin{pmatrix} r & s \\ t & u \end{pmatrix} $$
on aura (cf p. 398' (75)(78) sauf erreur)

[MARGIN: $\delta = \det \alpha = ad - bc$]
[MARGIN: $\delta' = \det \beta = ru - st$]

(94) $$ \underline{a} = \frac{1}{\delta D}(c\rho + d) $$ 
où $\delta D = c^2+cd+d^2$, [unclear] Dét $\alpha$
$$ = \frac{1}{c^2+cd+d^2}(c\rho + d) $$
on avait posé $\xi = \alpha \rho \alpha^{-1} = \begin{pmatrix} A & B \\ C & D \end{pmatrix}$

[MARGIN: Il faudra [unclear] dans le cas général]

(95) $$ \underline{b} = \frac{1}{\delta' U}(t\sigma + u) $$
où $\delta' U = t^2+u^2$
$$ = \frac{1}{t^2+u^2}(t\sigma + u) $$
où l'on a posé $\eta = \beta \sigma \beta^{-1} = \begin{pmatrix} R & S \\ T & U \end{pmatrix}$

Le seul ennui, pour faire le lien avec les calculs "bourrins" de loc. cit., est dans le facteur $\tau_\infty^\alpha$ de p. 398' (78), où $\alpha = [unclear: \varphi_2](\mu) \in \mathbb{Z}/2\mathbb{Z}$, dans le cas où $\alpha = 1$ --

\newpage

%% ===== Page 289 =====

à première vue il semble y avoir incom-
patibilité entre la théorie (plus simple)
développée ici, et ce facteur $\tau_{\infty}^2$ - dont je
suis persuadé qu'il a de très bonnes
raisons d'être là ! Et vrai - dans loc.
cit. on écrivait l'équation de base avec
la possibilité d'un signe

[MARGIN: [unclear: Vu le chgt de $\xi, \eta$ dans nouveaux $\alpha \rho(\alpha)^{-1}$ désormais $\beta \sigma \rho(\beta)^{-1}$]]

$$ \xi = \varepsilon \eta h_1 \qquad \varepsilon \in \{ \pm 1 \} $$
($\xi, \mu_{-1}$ avec les
notations actuelles)

c'est donc une généralisation p.r. à l'équation
posée au début, où on supposait $\varepsilon = +1$.
Il faudra donc tantôt reprendre la
théorie précédente, dans le contexte plus
général d'une équation

(95) $\qquad \xi = \varepsilon . \eta \varepsilon_1(\mu_{-1}) \qquad \varepsilon \text{ section de } \underline{O}_S$
$\qquad \quad || \qquad \quad ||$
$\qquad \alpha \rho(\alpha)^{-1} \quad \beta \sigma \rho(\beta)^{-1}$

où on devra avoir (en prenant les
déts. des deux membres) 
$$ \varepsilon^2 = 1 $$
- mais il n'y a pas lieu sans doute de
supposer qu'on ait $\varepsilon \in \{ \pm 1 \}$. Mais je vais
continuer d'abord les calculs dans le cas
présent $\varepsilon = +1$.

\newpage

%% ===== Page 290 =====

## IV Un épinglage canonique
[MARGIN: 445]

Rappelons que la donnée (sur un schéma de base $S$) d'une forme $G$ de $\operatorname{Gl}(2)_S$, ou $SG$ de $\operatorname{Sl}(2)_S$, ou $PG$ de $P\operatorname{Gl}(2)_S$, d'une algèbre d'Azumaya $\mathcal{M}$ (forme de $M_2(\mathcal{O}_S)$), revient à la donnée d'un fibré en droites projectives $\mathcal{P}$ sur $S$. $\mathcal{P}$ peut s'interpréter tantôt comme le schéma des sous-groupes de Borel $B$ de $G$, ou $SB$ de $SG$, ou $PB$ de $PG$, (ou comme le schéma des "sous-groupes additifs à un paramètre" "boréliens" i.e. identiques à la partie unipotente d'un Borel convenable) $L$ de $G$, ou de $SG$, ou de $PG$ (NB. Les morphismes canoniques $SG \to G \to PG$ induisent des bijections entre les ens. des sous-groupes de Borel des trois groupes en question ; et de même entre les ens. des sous-groupes à $1$ paramètre "boréliens" de ces groupes). Également, $\mathcal{P}$ peut s'interpréter en termes des sections $N$ de $\mathcal{M}$ qui sont non nulles sur tte fibre et telles que

(1) $\det N = 0$, $Tr N = 0$ d'où $N^2=0$ (les sections nilpotentes et régulières, c-a-d non nulles sur ch. pt. géom.), "strictement nilpotentes" ), à toute telle section correspond un sous-groupe borélien de $SG$, param. isom. de $\mathbb{G}_a$ sur lui,

(2) $\lambda \mapsto \exp(\lambda N) = 1 + \lambda N$

Mais il est évident que si $N$ et $N'$ st de telles
[MARGIN: 289]

\newpage

%% ===== Page 291 =====

sections qui sont homothétiques i.e. $N' = a N$
[unclear: en fait, avec] $a \in \Gamma(S, \mathcal{O}_S)^*$, alors elles définissent le \^m
sous-groupe à 1 paramètre, - donc ces dernières
correspondent en fait à des sections du
fibré projectif (de dim relative 3) $\mathbb{P}(M)$, -
si on tient compte de $Tr N = 0$, soit à des
sections du fibré projectif (de dim relative
2 - un fibré en plans projectifs)
$\mathbb{P}(M_0)$ (où $M_0 \subset M$ est défini justemt
par la condition $Tr N = 0$) - à savoir les
sections de la quadrique conique "définie"
par l'équation
$$det N = 0$$

Les $N$ correspondent d'ailleurs aussi aux
opérateurs
$$e = e_N = \text{[unclear: exp(N)]} \ e_N(1) = 1 + N$$
satisfaisant les conditions
(3) \ \ \ $det e = 1$, $Tr e = 2$
ce qui équivaut en effet, en posant
$e = 1 + N$, i.e. en définissant $N$ par
$$N = e - 1$$
aux conditions (1) [car $Tr$ et $det$ de
deux opérateurs $N$ et $e$ reliés par $e = \lambda + N$
sont [unclear: liés] mutuellement par
$$
\begin{cases}
Tr e = 2 \lambda + Tr N, & det e = \lambda^2 + \lambda Tr N + det N \\
Tr N = -2 \lambda + Tr e, & det N = \lambda^2 - \lambda Tr e + det e
\end{cases}
$$
]

\newpage

%% ===== Page 292 =====

447

[MARGIN: terminologie de quasi-involutions]
Les quasi-involutions sections de $G$ (en fait de $SG$) satisfaisant à (3), et qui ne sont pas $=1$, s'appellent régulières. Deux telles sections $e, e'$ sont considérées "équivalentes" si :

(4) 
$$e' = e^\lambda \overset{def}{=} e(\lambda) \quad (\lambda \in \Gamma_{G_a} = \Gamma_{\mathcal{O}_S})$$
$$\overset{def}{=} (1+\lambda N) \quad \text{on aura plutot intérêt} \quad \lambda \in \Gamma_{G_m}$$
$$= 1 + \lambda (e-1) = (1-\lambda) + \lambda e$$

On voit donc que les sections de $\mathcal{M}$ (strictement) nilpotentes régulières correspondent 1-1 aux sections unipotentes régulières, et la relation d'équivalence d'homothétie devient celle d'équivalence par $e \mapsto e^\lambda$ - et les classes d'équivalence de ces dernières correspondent aux sous-groupes additifs de $G$ (ou $SG$, ou $PG$ via $SG \to PG$) - ou encore (en prenant les normalisateurs de ceux-ci dans $G, SG$ ou $PG$) aux Borels de ces groupes, i.e. aux sections de $\mathcal{P}$ sur $S$.

Supposons donnée une telle section de $\mathcal{P}$, elle permet donc de déduire de [unclear] un fibré affine

(5) 
$$E (= P \setminus \{ \infty \})$$ 
(fibré vectoriel [unclear: des translations étant] $\mathcal{T}$)

\newpage

%% ===== Page 293 =====

Il $x_\infty$ est donné [unclear: rigidement] par un sous-
groupe à 1-paramètre borélien
(6) $$L \subset SG \subset G$$
(puis, par $SG \to PG$, donnera encore un
sous-groupe additif borélien $L \subset PG$ dans
$PG$...). Localement, $L$ sera engendré par
une section unipotente régulière
(7) $$e = 1 + N \quad \text{($N$ nilpotente régulière)}$$
comme le sous-groupe [unclear]
(8) $$e^\lambda = e(\lambda) = e_N(\lambda) = 1 + \lambda N$$,
mais comme on a dit, $e, N$ ne sont pas
déterminés intrinsèquement en termes de
$L \subset SG$. De façon précise, $L$ est munie
canoniquement d'une structure de fibré
vectoriel (respectée par ses [unclear: automorphismes]
dans $G$), et il devient justement celle
déduite par transport de structure de
celle de $\mathbb{G}_a$ par
$$e : \mathbb{G}_a \xrightarrow{\sim} L \qquad \lambda \mapsto e^\lambda = 1 + \lambda N$$,
et le choix de $e$ s'identifie [unclear] fibré
vectoriel canonique [unclear] fibré en
droites $\mathcal{N} \subset \mathcal{M}$ sous-[unclear: fibré] défini par les $N$
(9) $$L \simeq \check{V}(\mathcal{N})$$

\newpage

%% ===== Page 294 =====

448

le choix de $e$, ou de $N$, correspond simple-
ment ; celui d'une base de $\mathcal{N}$, ou
d'une section * de $L$ partout $\neq 0$
(* i.e. d'une section $e$ de $L$
qui soit $\neq 0$ sur les fibres).

[MARGIN: Notons que $L$ est aussi caractérisé comme fibré vectoriel des translations du fibré en droites affines $P \setminus \{0\}$.]

Revenons à l'équation (3b), et notons
que $\rho$ et $\sigma$ n'y interviennent qu'au module
des transformations $\rho \mapsto \lambda \rho$, $\sigma \mapsto \lambda \sigma$ (qui ne
changent pas les commutateurs
$[\alpha, \rho]$ et $[\beta, \sigma]$) ; en d'autres termes, les [unclear: formes]
[unclear: qu'on] [unclear: a vu] s'interprètent comme sections
du [unclear: sous-groupe] additif [unclear]
de $G$ dans $\Gamma(PG)$, d'
ailleurs

(4) $$PG \simeq \underline{Aut}_{S-gr}(G) \simeq \underline{Aut}_{S-gr}(SG) \simeq \underline{Aut}_S(P_G)$$
$$ \simeq \underline{Aut}_{S-Alg}(M) \simeq \underline{Aut}_S(P)$$

Il est à remarquer [unclear: ici], à dire que
$\rho$ et $\sigma$ n'interviennent que par leur action
sur $G$ (ou sur $SG$), ce qui est clair
[unclear: vu] sur ces formules
$$[\alpha, \rho] = \alpha \rho(\alpha)^{-1}, \quad [\beta, \sigma] = \beta \sigma(\beta)^{-1}.$$

D'autre part, la [unclear: section] $\varepsilon_1(\mathcal{H}-1)$ de (32)
peut être considérée comme une section
d'un [unclear: sous-fibré] sous-groupe additif $L_1$ de
$G$ (en fait, de $SG$), qui s'exprime bien

\newpage

%% ===== Page 295 =====

mutatis mutandis - l'équation (32),
s'écrit alors sous la forme (où je n'écris plus l'indice $\alpha, \beta$)

$$ (11) \qquad [\beta, \sigma]^{-1} [\alpha, \rho] \in \Gamma L_1 $$
ou
$$ (12) \qquad [\alpha, \rho] = [\beta, \sigma] \cdot \gamma \qquad \text{avec} \quad \gamma \in \Gamma L_1 . $$

Sous la forme (12), on peut la regarder comme une équation reliant $\alpha, \beta, \gamma$.

Considérons l'image $\tilde{L}_1$
$$ (13) \qquad \tilde{L}_1 \subset P G $$
de $L_1$ par $G \to P G$, la relation (40 bis)
s'écrit alors sous la forme
$$ (14) \qquad \tilde{\xi}_1 = \tilde{\rho} \tilde{\sigma}^{-1} \in \Gamma \tilde{L}_1 \qquad ( [unclear] \gamma \in \Gamma L_1 ) . $$

Cette section provient donc de une section bien déterminée indép. de $\tilde{\rho}, \tilde{\sigma}$

$$ (15) \qquad \xi_1 \in \Gamma L_1 $$
décrite en termes de $\rho, \sigma$ d'ailleurs par la relation qui est une relation de
forme

[MARGIN: NB $\lambda_0$ est uniquement déterminé p. $\rho$, si 2 [unclear] ne p. se recoller qu'en [unclear] ]

$$ (16) = \text{(40 bis)} \qquad \xi_1 = - \lambda_0 \rho \sigma^{-1} \qquad ( \lambda_0 \in \Gamma G_{ab} ) $$

par commodité
(à cause de l'application au cas abélien qui nous intéresse le plus ... ) .

\newpage

%% ===== Page 296 =====

443

Revenons sur la situation générale d'une forme $G$ de $PGL_2$, et supposons qu'elle puisse s'écrire sous la forme

$$ (17) \quad G = \operatorname{Aut}_{\mathcal{O}_S-Mod}(F) $$

où $F$ est un $\mathcal{O}_S$-Module loc. libre de rg 2.
C'est le cas dès qu'une section $x$ de $P$ sur $S$, i.e. un sous-groupe [unclear: additif?] borélien $L$ de $G$ est [unclear] (en fait on peut prendre alors pour $F$ l'extension de $\mathcal{O}_S$ par $\mathcal{N}$ qui détermine la fibré en droites affine $P \setminus \{x\}$). Alors le $P$ se [unclear]

$$ (18) \quad P = P(F) \simeq \mathbb{P}(\check{F}) $$

le fibré projectif associé à $F$, et les sections de $P$ s'identifient aux "droites" $D \subset F$.

On a alors

$$ (19) \quad \mathcal{M} = End_{\mathcal{O}_S-Mod}(F) \simeq \check{F} \otimes_{\mathcal{O}_S} F $$

et le sous-Module $\mathcal{N}$ associé à la section $x$ de $P$ n'est autre que

$$ (20) \quad \mathcal{N} = D^\circ \otimes D \quad (D^\circ = (F/D)\check{} \subset \check{F}, \text{ l'orthogonal de } D \text{ dans } \check{F}) $$

ses sections sont les $u : F \to F$ telles que

$$ (21) \quad u(F) \subset D, \quad u(D) = \{0\} $$

295

\newpage

%% ===== Page 297 =====

Proposition. Soit $G$ forme de $\operatorname{Gl}(n_S)$, $\rho, \sigma$ deux sections de $G$ telles que $\rho \sigma^{-1} = \varepsilon_1$ soit sect. unipot. regulière de $P(G)$, i.e. les conditions suiv. st. équivalentes :

(i) $\varepsilon_1$ est sect. regulière (i.e. [unclear] des fibres)
$\varepsilon_1 = 1 + \mathcal{N}_1$, avec $\mathcal{N}_1$ une [unclear] de $\dots$) et

[MARGIN: NB On a désigné par $X_1$ la section de $\check{P}$ qui correspond à $L_1$]

$\check{\rho}^{-1}(X_1)$ est une section de $\check{P}$ partout $\neq X_1$ dans chaque fibre.

(ii) Idem, avec $\check{\sigma}^{-1}$ au lieu de $\check{\rho}$.

(iii) On a (42), i.e. $c = \lambda_0 Tr \rho + Tr \sigma$ est inversible

Dém. On a, grâce au fait que [unclear]
normalise $L_1$, (la relation
$\check{\rho}^{-1}(L_1) = \check{\sigma}^{-1}(L_1)$ )
ce qui revient au m, on a

(22) $\check{\rho}^{-1}(X_1) = \check{\sigma}^{-1}(X_1)$,

ce qui montre l'équivalence de (i) et (ii).

Montrons que (i) $\iff$ (iii). OPS que $S$ est le
spectre d'un corps alg. clos, et OPS $\lambda_0 = 1$ (quitte à remplacer
$\rho$ par $\lambda_0 \rho$). Si on avait $\varepsilon_1 = 1$ i.e.
$\rho = \sigma$, on aurait $Tr(\rho+\sigma) = Tr \rho + Tr \sigma = 2$, ce qui n'implique pas la
régularité de $\varepsilon_1$, i.e. de $\mathcal{N}_1$, OPS donc que
$\varepsilon_1$ est régulière. Il faut montrer que si $\varepsilon_1 = 1 + \mathcal{N}_1$
est regulière, [unclear]
$D \subset F$,

\newpage

%% ===== Page 298 =====

et une involution $\sigma$ de $F$, d'où
$\rho = - \varepsilon_1 \sigma$, on a $c = Tr \sigma^{-1} + Tr \rho = Tr \sigma - Tr \varepsilon_1 \sigma$
$= - Tr N_1 \sigma$ est inversible ssi $D + \sigma \cdot D = F$
i.e. $\sigma \cdot D + D = F$, Prenons une base
$e_0, e_1$ de $F$ telle que $e_1$ engendre $D$, et
que $N_1$ s'écrive par la matrice $\begin{pmatrix} 0 & 0 \\ 1 & 0 \end{pmatrix}$,
soit $\sigma = \begin{pmatrix} a & b \\ c & d \end{pmatrix}$, on a $N_1 \sigma = \begin{pmatrix} 0 & 0 \\ a & b \end{pmatrix}$,
$Tr N_1 \sigma = b$, $\sigma \cdot e_1 = a e_0 + b e_1$, et $D + \sigma \cdot D = D + k \sigma \cdot e_1$
est égal à $F = k^2$ ssi $b$ inv., cqfd.

Supposons satisfaites ces conditions
équivalentes, qui s'expriment
en disant que $L_1$ et

(23) $$L_0 = \tilde{\rho}^{-1}(U_{-1}) = \tilde{\sigma}^{-1}(U_{-1})$$

sont distincts sur chaque fibre, i.e. définissent
des Borels $B_1, B_0$ de $G$ tels que

(24) $$T = B_1 \cap B_0$$

soit un tore (maximal) de $G$.
Ce choix [unclear: équivaut] à [unclear: se donner] l'intersection
des stabilisateurs $B_0, B_1$ dans $G$ des sections
$x_0, x_1$ de $\tilde{P}$ qui correspondent à $L_0, L_1$
(à l'origine $x_1 = \tilde{\rho} x_0 = \tilde{\sigma} x_0$),
(25) et à [unclear: un choix entre] ces deux sous-groupes (Borel) ...

\newpage

%% ===== Page 299 =====

dont $\tilde{T}$ dans $\underline{PGL} = \underline{Aut}_S(\mathbb{P}_1)$. On peut dire
aussi que la donnée de deux sections
partout distinctes $x_0, x_1$ du fibré projectif
$\mathbb{P}$ revient à la donnée : celle d'un Borel
$B_1$ (correspondant au divis. de $x_1$), et d'un
tore maximal $T$ de $B_1$ (qui s'écrit
donc

[MARGIN: le couple $(B_1, T)$ s'appelle un couple de Killing.]

(26) $$B_1 = T \cdot L_1 \quad \text{(prod. semi direct)}$$

qui [unclear: par la donnée] de $x_0$ cor-
respond à un unique Borel $B_0$ tel que
$T = B_0 \cap B_1$. (les Borels $B_0$ tels que
$B_0 \cap B_1$ soit un tore, i.e. tels que $B_0 \neq B_1$ sur chaque fibre, corresp. en bijection
aux tores maximaux $T$ de $B_1$,
par la formule précédente $T = B_0 \cap B_1$ .)

Cela revient à ceci : la donnée, en
plus de $B_1$ (i.e. de $x_1$), d'une section $x_0$ du
fibré en droites affines $U_1 = \mathbb{P} \setminus \{x_1\}$, équi-
vaut à se donner par translation
un isomorphisme du fibré $U_1$ sur le
fibré vectoriel $L_1$ qui opère par translation

(27) $$U_1 = \mathbb{P} \setminus \{x_1\} \xrightarrow{\sim} L_1$$

Pour arriver "à épingler" la situation [unclear]

\newpage

%% ===== Page 300 =====

de réduire $\mathcal{P}$ à un fibré projectif de type $\mathbb{P}^1_S$ )
il faut avoir en plus une base de
$L_1$ i.e. une section $x_\infty$ de

(28) $$L_1^* \simeq U_1 \setminus \{x_0\} \simeq \mathcal{P} \setminus \{x_1, x_0\} .$$

Or la section $e_1$ de $L_1$ est justement
une section de $L_1^*$ - on a donc un
épinglage complet de la situation. En
termes de $(F, D=D_1)$ correspondant à $L_1$,

on a $$F = D_1 \oplus \sigma^{-1}(D_1) = D_1 \oplus \wp^*(D_1) ,$$

[crossed out text]

(29)
$$ \begin{cases} \wp^{-1}(D_1) = \sigma^{-1}(D_1) \overset{df}{=} D_0 \\ F = D_0 \oplus D_1 \end{cases} $$

(30)
$$N_1 | D_0 : D_0 \overset{\sim}{\longrightarrow} D_1$$

et (quitte à tensoriser par $D_1^{\otimes -1}$ ) OPS
$D_1 \simeq \mathcal{O}_S$, d'où

(31)
$$ \begin{cases} F \simeq \mathcal{O}_S^2 = \mathcal{O}_S e_0 + \mathcal{O}_S e_1 \\ D_0 = \mathcal{O}_S e_0 , \quad D_1 = \mathcal{O}_S e_1 \\ N_1 e_1 = 0, N_1 e_0 = e_1 \quad \text{i.e. } \text{[unclear crossed-out matrix]} \end{cases} $$

\newpage

%% ===== Page 301 =====

i.e.

(32) $$N_1 = \begin{pmatrix} 0 & 0 \\ 1 & 0 \end{pmatrix} \quad , \quad \varepsilon_1 = \begin{pmatrix} 1 & 0 \\ 0 & 1 \end{pmatrix} \quad .$$

D'autre part nous choisirons $\rho, \sigma$ relèvent les [unclear: aut(2)] ci-dessus en $\bar{\rho}, \bar{\sigma}$ (sections de $GP(1) \simeq \operatorname{Gl}(2)/\mathbb{G}_m$) satisfaisant $\bar{\rho} \bar{\sigma}^{-1} = \bar{\varepsilon}_1$, de façon à avoir

(33) $$\rho \sigma^{-1} = -\varepsilon_1$$

ce qui détermine chacune des deux sect. $\rho, \sigma$ en fonction de l'autre (déterminées elles-mêmes modulo un scalaire près, une fonction de $\mathcal{O}_{S}^\ast$). Comme

$$\rho D_0 = D_1 \quad ,$$

on aura

(*) $$\rho = \begin{pmatrix} 0 & u \\ v & w \end{pmatrix} \quad , \quad \det \rho = -uv \text{ inv.}$$

et on peut lever l'ambiguité scalaire sur $\rho$ en [unclear] par

(34) $$\rho = \begin{pmatrix} 0 & 1 \\ v & w \end{pmatrix} \quad v \in \mathbb{G}_m \quad .$$

\newpage

%% ===== Page 302 =====

Notons que (33) s'écrit $\rho = -(1+N_1)\sigma$ i.e.

(35) $\rho + \sigma = -N_1 \sigma$

multipliant par $N_1$ à g. on trouve aussi

$N_1 \rho + N_1 \sigma = 0$

(puisque $N_1^2 = 0$), d'ailleurs

$N_1 \rho = \rho N_0$ , $N_1 \sigma = \sigma N_0$

(car $N_1 = \rho N_0 \rho^{-1} = \sigma N_0 \sigma^{-1}$), donc (35) s'écrit

(36) $\rho + \sigma = -N_1 \sigma = +N_1 \rho = -\sigma N_0 = +\rho N_0 \overset{\text{df}}{=} N$

d'ailleurs quand $\rho$ est sous la forme

générale $\rho = \begin{pmatrix} 0 & u \\ v & w \end{pmatrix}$ on trouve

$N = \begin{pmatrix} 0 & 0 \\ 0 & uv \end{pmatrix}$ , $N^2 = \begin{pmatrix} 0 & 0 \\ 0 & u^2 v^2 \end{pmatrix}$

dans la normalisation (34) $u=1$ signifie

(37) $\rho + \sigma = \begin{pmatrix} 0 & 0 \\ 0 & 1 \end{pmatrix}$

ou encore 

(37 bis) $N^2 = N$ (où $N = \rho + \sigma$)

ou encore que $\sigma$ est donné en termes

de $\rho$ par

[MARGIN: NB On voit aussitôt que les conditions envisagées ($tr \sigma = - tr \rho$) équivalent à dire que $N = \rho + \sigma = \rho N_0 \rho^{-1}$ satisfait $N^2 = N$, i.e. $Tr N = 1$, i.e. $uv = 1$]

$\sigma = \begin{pmatrix} 0 & -1 \\ -v & 1-w \end{pmatrix}$

$N_0$ sera de la forme $\begin{pmatrix} 0 & \alpha \\ 0 & 0 \end{pmatrix}$

un calcul immédiat montre que

$\alpha$ est donné par la formule $u = \alpha v$, i.e. $\alpha = v^{-1}$

en normalisant par $u=1$)

\newpage

%% ===== Page 303 =====

(39) $$N_0 = \begin{pmatrix} 0 & 1/w \\ 0 & 0 \end{pmatrix}, \quad \varepsilon_0 = \begin{pmatrix} 1 & 1/w \\ 0 & 1 \end{pmatrix}$$

[MARGIN: N.B. On aura
(39bis) $\sigma^{-1} \rho = -\varepsilon_0$
Car $\sigma^{-1} \rho = \sigma^{-1} (\rho \sigma^{-1}) \sigma = \sigma^{-1}(-\varepsilon_1) \sigma = - \sigma^{-1}(\varepsilon_1) \sigma = -\varepsilon_0$]

Ces normalisations canoniques étant faites, revenons à la proposition de la page 436 ; comme $\varepsilon_1$ a été choisi (de façon tt ce qu'il y a de canonique) il n'y a plus d'inconvénient à laisser l'équation (32), ( mais sous la forme un peu plus intrinsèque (41)... (43) ) sous la forme primitive, où a priori on aurait à introduire le paramètre $\gamma$ avec

$$[\alpha, \rho] = [unclear: [\rho, \beta] ] \varepsilon_1(\gamma)$$

mais il apparait que c'est $\gamma+1=\mu$ qui a une signification particulièrement remarquable dans la théorie - d'où on gardera la forme (32), avec le facteur $\varepsilon_1(\mu-1)$.

L'expression (37) de $P(\xi)$, grâce à 
(37), s'écrit

(40) $$P(\xi) = D \quad \text{pour} \quad \xi = \begin{pmatrix} A & B \\ C & D \end{pmatrix},$$

\newpage

%% ===== Page 304 =====

453

Donc la proposition nous dit que le schéma $\widetilde{X}_{\rho, \sigma}$
des solutions simultanées $(\xi, \eta, \mu)$ de
(41) $$ \det \xi = \det \eta = 1, \quad \xi = \eta \varepsilon_1(\mu-1), \quad \text{Tr}(\xi \rho) = \text{Tr} \rho, \quad \text{Tr} \eta \sigma = \text{Tr} \sigma $$
est isomorphe, via $(\xi, \eta, \mu) \mapsto \xi$, à celui
des $\xi = \begin{pmatrix} A & B \\ C & D \end{pmatrix}$ satisfaisant les
conditions,
(42) $$ \det \xi = 1, \quad \text{Tr}(\xi \rho) = \text{Tr} \rho, \quad \text{et } D \text{ section inversible de } \mathcal{O}_S $$
[MARGIN: $AD-BC$] [MARGIN: $Cw+vB$]
et on retrouve $\mu$ en termes de $\xi$ par
(43) $$ \mu = 1/D $$
(ici $c = \text{Tr}(\rho \varepsilon_1 \rho^{-1}) = \text{Tr} \begin{pmatrix} 0 & 0 \\ 0 & 1 \end{pmatrix} = 1$ ), d'où
$\eta$ par $\eta = \xi \varepsilon_1(1-\mu)$.

Le sous-schéma en groupes $B_0'$ de
la prop 442' n'est autre que le schéma
des matrices
(44) $$ B_0' = \left\{ \begin{pmatrix} \mu & \lambda \\ 0 & 1 \end{pmatrix} \mid \mu \in \mathbb{G}_m, \lambda \in \mathbb{G}_a \right\} $$

Nous allons expliciter l'isomorphisme (39)
(45) $$ u \longmapsto ( [u, \rho], [u, \sigma], \det u ) $$
en termes des coefficients $\mu, \lambda$ de $B'_0$,
à l'aide des paramètres $v, w$ dans $\rho = \begin{pmatrix} 0 & 1 \\ v & w \end{pmatrix}$

\newpage

%% ===== Page 305 =====

On trouve

$$ [\underline{u},\rho] = \begin{pmatrix} \mu & \lambda \\ 0 & 1 \end{pmatrix} \begin{pmatrix} 0 & 1 \\ v & \omega \end{pmatrix} \frac{1}{\mu} \begin{pmatrix} 1 & -\lambda \\ 0 & \mu \end{pmatrix} \begin{pmatrix} \omega & -1 \\ -v & 0 \end{pmatrix} $$
$$ = - \frac{1}{\mu v} \begin{pmatrix} \lambda v & \mu+\lambda\omega \\ v & \omega \end{pmatrix} \begin{pmatrix} \omega+\lambda v & -1 \\ -\mu v & 0 \end{pmatrix} $$
$$ = \frac{1}{\mu} \begin{pmatrix} -v\lambda^2 + \omega\lambda\mu + \mu^2 - \omega\lambda & \lambda \\ -v\lambda + \omega\mu - \omega & 1 \end{pmatrix} $$

[MARGIN: on pose $\sigma = \begin{pmatrix} 0 & -1 \\ v' & \omega' \end{pmatrix}$ avec $v' = v$, $\omega' = 1-\omega$]

$$ [\underline{u},\sigma] = \begin{pmatrix} \mu & \lambda \\ 0 & 1 \end{pmatrix} \begin{pmatrix} 0 & -1 \\ v' & \omega' \end{pmatrix} \frac{1}{\mu} \begin{pmatrix} 1 & -\lambda \\ 0 & \mu \end{pmatrix} \begin{pmatrix} \omega' & 1 \\ -v' & 0 \end{pmatrix} $$
$$ = \frac{1}{\mu v'} \begin{pmatrix} \lambda v' & -\mu+\lambda\omega' \\ v' & \omega' \end{pmatrix} \begin{pmatrix} \omega'+\lambda v' & 1 \\ -\mu v' & 0 \end{pmatrix} $$
$$ = \frac{1}{\mu v'} \begin{pmatrix} v'\lambda^2 + v'\mu^2 - v'\omega'\lambda\mu + v'\omega'\lambda & \lambda v' \\ v'\lambda - v'\omega'\mu + v'\omega' & v' \end{pmatrix} $$
$$ = \frac{1}{\mu} \begin{pmatrix} v'\lambda^2 - v'\omega'\lambda\mu + \mu^2 + \omega'\lambda & \lambda \\ v'\lambda - \omega'\mu + \omega' & 1 \end{pmatrix} $$
$$ = \frac{1}{\mu} \begin{pmatrix} -v\lambda^2 + (\omega-1)\lambda\mu + \mu^2 - (\omega-1)\lambda & \lambda \\ -v\lambda + (\omega-1)\mu - (\omega-1) & 1 \end{pmatrix} $$

En résumé

[MARGIN: [unclear: abréviation]]
$$ (46) \begin{cases} \xi = [\underline{u},\rho] = \frac{1}{\mu} \begin{pmatrix} \mu^2 + (-v\lambda + \omega(\mu-1))\lambda & \lambda \\ -v\lambda + \omega(\mu-1) & 1 \end{pmatrix} \\ \eta = [\underline{u},\sigma] = \frac{1}{\mu} \begin{pmatrix} \mu^2 + (-v\lambda + (\omega-1)(\mu-1))\lambda & \lambda \\ -v\lambda + (\omega-1)(\mu-1) & 1 \end{pmatrix} \end{cases} $$

$u = \det \underline{u}$

donc, si $\xi = \begin{pmatrix} A & B \\ C & D \end{pmatrix}$
(47) $B = \lambda / \mu$, $D = 1 / \mu$ i.e.
(48) $\mu = 1 / D$, $\lambda = B / D$, donc $\underline{u} = \begin{pmatrix} 1/D & B/D \\ 0 & 1 \end{pmatrix}$

\newpage

%% ===== Page 306 =====

459 bis

De même, si $y = \begin{pmatrix} R & S \\ T & U \end{pmatrix}$, on a $S=B, U=D$

d'où

(47 bis) $S = B = \lambda / \mu$, $U = D = 1 / \mu \quad$ d'où

et $\mu = 1/U$, $\lambda = S/U$

donc

(48 bis) $\underline{u} = \begin{pmatrix} 1/U & S/U \\ 0 & 1 \end{pmatrix}$

Considérons maintenant l'équation

(49) $[\alpha, \rho] = [\beta, \sigma] \varepsilon_1(\mu_1) , \quad \text{avec}$

(50) $\alpha = \begin{pmatrix} a & b \\ c & d \end{pmatrix} , \beta = \begin{pmatrix} r & s \\ t & u \end{pmatrix} ,$

et posons, pour des calculs plus élégants

(51) $\rho = \begin{pmatrix} 0 & v_1 \\ v & w \end{pmatrix} , \sigma = \begin{pmatrix} 0 & v_1' \\ v' & w' \end{pmatrix}$

[MARGIN: on aura pour $v_1=1$, $v'_1=1$ [unclear] le calcul pour $[\beta, \sigma]$ [unclear: symétrique de] pour $[\alpha, \rho]$ !]
(au lieu de $u, u'$, pour ne pas interférer avec la notation $u$ pour un des coefficients de $\beta$ ), avec

(52) $v' = -v, v_1' = -v_1, w + w' = v_1 \quad , \text{ posons}$

(53) $\delta_\alpha = \det \alpha = ad-bc, \delta_\beta = \det \beta = ru-st ,$

On a alors

$y = [\alpha, \rho] = \begin{pmatrix} a & b \\ c & d \end{pmatrix} \begin{pmatrix} 0 & v_1 \\ v & w \end{pmatrix} \frac{1}{\delta_\alpha} \begin{pmatrix} d & -b \\ -c & a \end{pmatrix} \frac{1}{-v_1 v} \begin{pmatrix} w & -v_1 \\ -v & 0 \end{pmatrix}$

$= \frac{1}{-\delta_\alpha v_1 v} \begin{pmatrix} bv & a v_1 + bw \\ dv & c v_1 + dw \end{pmatrix} \begin{pmatrix} d w + b v & -d v_1 \\ -c w - a v & c v_1 \end{pmatrix} = \begin{pmatrix} A & B \\ C & D \end{pmatrix}$

[MARGIN: NB les formules pour B/D s'implifient (!)]
(54) $\begin{cases} (-\delta_\alpha v v_1) A = \dots \text{(banale)} , & (-\delta_\alpha v v_1) B = -v v_1 b d + v_1^2 a c + v_1 w b c \\ (-\delta_\alpha v v_1) C = \dots \text{(banale)} , & (-\delta_\alpha v v_1) D = -v v_1 d^2 + v_1 w c d + v_1^2 c^2 \end{cases}$

305

\newpage

%% ===== Page 307 =====

et on trouve de même

$$y = (\beta, \sigma) = \begin{pmatrix} R & S \\ T & U \end{pmatrix} \quad avec$$

(55) $\begin{cases} (-v v_1 \delta_\beta) R = \dots & (-v v_1 \delta_\beta) S = (v v_1) su + v_1^2 tr - v_1 w' st \\ (-v v_1 \delta_\beta) T = \dots & (-v v_1 \delta_\beta) U = -(v v_1) u^2 - v_1 w' tu + v_1^2 t^2 \end{cases}$
[MARGIN: NB on peut diviser par $v_1$]

On en tire en particulier

(56) $\begin{cases} D = \frac{-v v_1 d^2 + v_1 w cd + v_1^2 c^2}{-v v_1 (ad-bc)} \\ \mu = 1/D = \frac{-v v_1 (ad-bc)}{-v v_1 d^2 + v_1 w cd + v_1^2 c^2} \end{cases}$

[crossed out] $\frac{B}{D} = \frac{-v v_1 bd + v_1^2 ac + v_1 w bc}{-v v_1 d^2 + v_1 w cd + v_1^2 c^2}$

ce qui permet donc d'expliciter

$$y = \begin{pmatrix} 1/D = \mu & B/D = \lambda \\ 0 & 1 \end{pmatrix}$$

en termes de $a, b, c, d$ (et des paramètres $v, v_1, w$ qui déterminent $\rho$).

Pour $v_1 = 1$, la formule pour $D$ de (56) s'écrit

(57) $$D = \det (d\rho+c)/-v \delta_\alpha = \det (d\rho+c)/\det\rho\det\alpha$$

la condition [crossed out] sur $\alpha$ "D invers."

(58) $$D(\alpha) \text{ inversible } \iff d\rho+c \text{ inversible.}$$

mais je m'aperçois que c'est une fausse piste,

\newpage

%% ===== Page 308 =====

[MARGIN: (de l'équation 45?) [crossed out]]

V $\underline{\text{Étude dans } PGl(1), \text{ et la condition } Tr \sigma = 0.}$

Jusqu'à présent nous avons étudié l'équation (32)

(*) $$[\alpha,\rho] = [\beta,\sigma] \varepsilon_1(\mu_{-1})$$

dans $G (\simeq \operatorname{Gl}(2))$, ou ce qui revient au même, en tant que relation dans $\operatorname{Sl}(2)$. Ici $\alpha, \beta$ sont des sections variables dans $G (\simeq \operatorname{Gl}(2))$, mais $[\alpha,\rho]$ et $[\beta,\sigma]$ ne dépendent que des sections correspondantes $\tilde{\alpha}, \tilde{\beta}$ de $PG (\simeq PGl(1)) \simeq G/\mathbb{G}_m \simeq \operatorname{Aut}(\mathbb{P}_1)$ (en ce qui [unclear: concerne] la dépendance de $\rho, \sigma$ qui est $\tilde{\rho}, \tilde{\sigma}$). On peut donc considèrer que c'est une équation liant $\tilde{\alpha}, \tilde{\beta}, \rho, \sigma$.
On peut se proposer d'étudier l'équation correspondante dans $PG \simeq PGl(1)$,

(1) $$[\tilde{\alpha},\tilde{\rho}] = [\tilde{\beta},\tilde{\sigma}] \cdot \tilde{\varepsilon}_1(\mu_{-1})$$

ce qui s'écrit aussi, dans $\operatorname{Gl}(2)$, sous la forme équivalente

(2) $$[\alpha,\rho] = z [\beta,\sigma] \tilde{\varepsilon}_1(\mu_{-1})$$

où $z$ est une section de $\mathbb{G}_m$, et en fait d'ailleurs uniquement déterminée par $\tilde{\alpha}, \tilde{\beta}, \mu$. On [unclear].

\newpage

%% ===== Page 309 =====

(3)
$$ \varepsilon^2 = 1 $$

(prendre le dét. des deux membres). On
considère (2) comme une équation
reliant les variables $\alpha, \beta, \mu, \varepsilon$ ($\alpha, \beta$ dans
$\operatorname{Gl}(2)$, $\mu$ dans $\hat{\mathbb{Z}}'$ - et en fait [unclear]
dans $\mathbb{G}_m$, et $\varepsilon$ pris dans $\mu_2 = (\subset \mathbb{G}_m)$),
ou en $\tilde{\alpha}, \tilde{\beta}, \mu, \varepsilon$. Posons, par analogie avec (II),

(4)
$$ \xi = [\alpha, \rho] , \quad \eta = \varepsilon [\beta, \sigma] $$

de sorte que (2) prenne la forme

(5)
$$ \xi = \eta \varepsilon_1(\mu-1) $$

comme dans (II). Pour pouvoir appliquer
(II) à cette situation, il faudra s'assurer
que $\eta = \varepsilon [\beta, \sigma]$ satisfait aux conditions
préliminaires que

$$ \eta_0 = [\beta, \sigma] $$

lui-même, savoir $\det \eta = 1$ ($\iff \varepsilon^2 = 1$) et

$$ Tr \eta = Tr \sigma $$
$$ \varepsilon Tr \eta_0 = \varepsilon Tr \sigma $$

i.e.

$$ \varepsilon Tr \sigma = Tr \sigma $$

\newpage

%% ===== Page 310 =====

Lemme : Soit $\sigma$ une section de $G$ (fibré
des $\operatorname{Gl}(2)$, associé à une Alg d'Azumaya $M$)
— i.e. section de $M$ telle que $det \sigma$ in-
versible. Soit $A_\sigma$ le schéma des sections
$y$ de $G$ satisfaisant les conditions
$det y = 1$ (i.e. $y$ dans $SG$), $Tr y\sigma = Tr \sigma$.

Conditions équivalentes a) b) c) :

a) $A_\sigma$ est stable par translations par $\mu_2 \subset \mathbb{G}_m$ :
$$ y \mapsto \varepsilon y \quad (\varepsilon^2 = 1) $$

b) $Tr \sigma = 0$
[unclear: texte raturé] III (p. 451),

d) (Dans le contexte des n° précédents, avec les
normalisations (34) (38) pour $\rho, \sigma$) on a .

$$ \sigma = \begin{pmatrix} 0 & -1 \\ v & 0 \end{pmatrix} = \begin{pmatrix} 0 & -1 \\ 1 & 0 \end{pmatrix} \quad (6) $$

i.e. (avec les notations de loc. cit.) $v = 1$,

re.

$$ \rho = \begin{pmatrix} 0 & 1 \\ 1 & 1 \end{pmatrix} \quad (7) $$

[unclear: paragraphe raturé et illisible contenant "Enfin ces conditions impliquent..." et "a) $\iff$ b)"]

que $E' = (0)$ , ce qui est d'ailleurs
de vérification immédiate. L'équivalence de
b) et d) est immédiate, reste à prouver
b) $\iff$ c) . Or cette condition est satisfaite en vertu du th. de Hamilton-Cayley
$$ \sigma^2 - (Tr \sigma) \sigma + det \sigma = 0 . $$

[MARGIN: c) On a $\sigma^2 \in \Gamma \mathbb{G}_m$, i.e. $\sigma^2 = 1$]

\newpage

%% ===== Page 311 =====

qui montre que $Tr \sigma = 0 \implies \sigma^2 = - det \sigma \cdot 1$,
et inversement si $\sigma^2$ est un scalaire,
on aura que $(Tr \sigma) \sigma = \sigma^2 + det \sigma \cdot 1$ est
un scalaire, ce qui n'est possible (si $\sigma$ n'est
pas un scalaire sur aucune fibre) que
si $Tr \sigma = 0$.

Supposons par la suite ces conditions
satisfaites. Soit alors 

(8) $$G_{\rho, \sigma}^* \subset G \times G \times \underline{E}' \times \underline{E}'$$

le schéma des solutions $(\alpha, \beta, \mu, \varepsilon)$
de l'équation (2), qui est aussi l'image
inverse dans $G \times G \times \underline{E}' \times \underline{E}'$ du sous-
schéma

(9) $$\tilde{G}_{\rho, \sigma}^* \subset PG \times PG \times \underline{E}'$$

formé des solutions $(\alpha, \beta, \mu)$ de (4). La
projection sur le dernier facteur $\underline{E}'$
fournit un morphisme

(10) $$\tilde{G}_{\rho, \sigma}^* \xrightarrow{\tilde{pr}} \underline{E}'$$ ,

\newpage

%% ===== Page 312 =====

456

qui se factorise d'ailleurs par $\tilde{\mathcal{G}}^*_{\rho, \sigma}$ :
(11) $$ \tilde{\varepsilon} : \tilde{\mathcal{G}}^*_{\rho, \sigma} \to \mu_2 $$
Notons qu'on a
$$ \mathcal{G}_{\rho, \sigma} \subset \mathcal{G}^*_{\rho, \sigma} $$
plus précisément
(12) $$ \mathcal{G}_{\rho, \sigma} = \varepsilon^{-1}(1) $$
ce qui conduit à introduire aussi
(13) $$ \begin{matrix} \tilde{\mathcal{G}}_{\rho, \sigma} \subset \tilde{\mathcal{G}}^*_{\rho, \sigma} \\ \parallel dfn \\ \tilde{\varepsilon}^{-1}(1) \end{matrix} $$
de sorte que $\mathcal{G}_{\rho, \sigma}$ est l'image inverse dans $G \times G \times E'$ de $\tilde{\mathcal{G}}_{\rho, \sigma}$.

On a un morphisme
$$ (\alpha, \beta, \mu, \varepsilon) \longmapsto ( (\xi = \alpha \rho(\alpha)^{-1}, \eta = \varepsilon \beta \sigma(\beta)^{-1}, \mu), \varepsilon ) $$
$$ \in \Gamma \mathcal{H}_{\rho, \sigma} \times \mu_2 $$
(14) de 
$$ \mathcal{G}^*_{\rho, \sigma} \to \mathcal{H}_{\rho, \sigma} \times \mu_2 $$
et l'image inverse d'une section $((\xi, \eta, \mu), \varepsilon)$ du second membre est le schéma des solutions en $\alpha, \beta$ des équations
(15) (a) $\alpha \rho(\alpha)^{-1} = \xi$, (b) $\beta \sigma(\beta)^{-1} = \varepsilon \eta$

\newpage

%% ===== Page 313 =====

Or on a, vu l'isom.
$$B_0 \simeq \mathcal{H}_{\rho,\sigma}$$
une solution privilégiée $\alpha = u$ de la
première équation (15), qui est
caractérisée comme l'unique solution
qui soit une section de $B_0$ (remarquons
que $T_0 \cap B_0 = \{1\}$).

Mais on a, pour cet $u$ :
$$u \sigma(u)^{-1} = \eta \qquad (\text{et non } \varepsilon \eta \text{ !})$$
$$[u, \sigma]$$

pour trouver en terme de $u$ une solution
de l'équation (15 b)
$\beta \sigma(\beta)^{-1} = \varepsilon \eta$, on est
conduit à l'essayer sous la forme
$$(16) \qquad \beta = u \tau \quad \text{d'où} \quad \beta \sigma(\beta)^{-1} = u \tau \sigma(\tau)^{-1} \sigma(u)^{-1} ,$$
en essayant de déterminer $\tau$ de façon
qu'on ait
$$(17) \qquad \tau \sigma(\tau)^{-1} = \varepsilon$$
$$[\tau, \sigma]$$
ce qui implique bien
$$\beta \sigma(\beta)^{-1} = u (\varepsilon \sigma(u)^{-1}) = \varepsilon (u \sigma(u)^{-1}) = \varepsilon \eta$$
i.e. (15 b).

\newpage

%% ===== Page 314 =====

La relation (17) [MARGIN: 457] s'écrit aussi

[crossed out: illegible]

(18) $$ \tau \sigma \tau^{-1} \sigma^{-1} = \varepsilon \quad \text{ou} \quad \sigma \tau \sigma^{-1} \tau^{-1} = \varepsilon \quad \text{ou} $$

(19) $$ \sigma^\tau = \varepsilon \sigma \quad \text{ou} \quad \tau(\sigma) = \varepsilon \sigma $$

ce qui implique que $\tau$ normalise $T_\sigma$

(20) $$ \tau(T_\sigma) = T_\sigma \quad \text{i.e.} \quad \tau \in \mathcal{N}_\sigma $$

où
(21) $$ \mathcal{N}_\sigma \overset{\text{déf}}{=} \operatorname{Norm}_G(T_\sigma) . $$

Corollaire au lemme (p 455) [MARGIN: algébriques] les conditions algébriques suivantes [unclear: sont équivalentes pour les] sections $\sigma$
(e) $\forall$ section $\varepsilon$ de $\mu_n$, existe section $\tau$ de $G$ telle qu'on ait (17) (ou encore (19))
(f) (si 2 inv. dans $S$) idem, avec $\varepsilon = -1$.

Si on a (e), alors on voit que pour [unclear: toute] $\varepsilon$, on doit avoir 
$$ Tr \varepsilon \sigma = Tr \sigma $$
(puisque $\varepsilon \sigma$ est conj. à $\sigma$)
Or $Tr \varepsilon \sigma = \varepsilon Tr \sigma$, donc $(Tr \sigma (\varepsilon - 1) = 0)$.
Si on suppose [unclear: alors] que $\varepsilon = -1$, on trouve $Tr \sigma = - Tr \sigma$ i.e.
$$ 2 Tr \sigma = 0 \quad \text{i.e.} \quad Tr \sigma = 0 \text{ si 2 inv.} $$
Inversement, si $Tr \sigma = 0$, alors $Tr \varepsilon \sigma = 0$, d'ailleurs $det(\varepsilon \sigma) = \varepsilon^2 = 1$, donc $\varepsilon \sigma$ est [unclear: conj. à $\sigma$]

[unclear: d'où] "conditions nécessaires" pour être de la forme $[\tau, \sigma]$ pour $\tau$ convenable. Puis [unclear]

\newpage

%% ===== Page 315 =====

assurera l'existence locale de $\tau$, il reste
à s'assurer que $\varepsilon \sigma$ est section régulière
au sens de Steinberg i.e. n'est nulle
part un scalaire, ce qui provient du fait
qu'il en est ainsi de $\sigma$.

Ainsi, sous les conditions actuelles,
soit $\mathcal{N}'_\sigma \subset \mathcal{N}_\sigma$ le sous-schéma de [unclear]
défini par
(22) $$ \mathcal{N}'_\sigma = \{ \tau \mid [\tau, \sigma] \in \Gamma \mu_r \} $$
alors $\mathcal{N}'_\sigma$ est stable par translations à droite
par $T_\sigma$ - mieux, c'est un sous-schéma
(contenant $T_\sigma$) en groupes de $\mathcal{N}_\sigma$, car si
$\tau, \tau'$ sont tels que
$$ \tau(\sigma) = \varepsilon \sigma \quad , \quad \tau'(\sigma) = \varepsilon' \sigma , $$
alors $\tau' \tau(\sigma) = \varepsilon' \varepsilon \cdot \sigma$,
[unclear: ce qui prouve en particulier que]
$\tau \mapsto \varepsilon$
donne un hom. de groupes $\mathcal{N}'_\sigma \to \mu_r$, et on a
une suite exacte de schémas en
groupes
(23) $$ 1 \to T_\sigma \hookrightarrow \mathcal{N}'_\sigma \xrightarrow{\varepsilon_\sigma} \mu_r \to 1 $$ .

On notera que si $\tau \in \Gamma G$ normalise $T_\sigma$,
on a $\tau(\sigma) = a \sigma + b \quad a,b \in \Gamma \underline{O}_S$,

\newpage

%% ===== Page 316 =====

et prenant les traces des deux membres on
trouve (compte tenu que $Tr \tau(\sigma) = Tr \sigma = 0$)
$Tr\ b.1 = 0 \quad$ i.e. $\quad 2b = 0$, donc si $2$ est inv.
$b = 0$, donc $\tau(\sigma) = a\sigma$ et nécessairement
aussi $a^2 = 1 \quad$ i.e. $\quad a$ est dans $\varepsilon$. Donc

(24) $\qquad \mathcal{N}_\sigma^\prime = \mathcal{N}_\sigma$ en car. résiduelle $\neq 2$.

[MARGIN: crossed out text]
Plus généralement, sur une base $S$ telle que
$2$ ne soit pas nilpotent (i.e. sans div. de
zéro) dans $\Gamma(S, \mathcal{O}_S)$, on voit que toute
section de $\mathcal{N}_\sigma^\prime$ est en fait une section de
$\mathcal{N}_\sigma$. Mais il n'en est plus ainsi en
car. $2$, car
$$ \sigma^\prime = a\sigma + b \qquad (a \in \Gamma\mathbb{G}_m, b \in \Gamma\mathcal{O}_S) $$
dès qu'il satisfait
$$ Tr\ \sigma^\prime = 0, \quad \det\ \sigma^\prime = \det\ \sigma $$
est conjugué loc^t. à $\sigma$ - et il suffit pour
s'en assurer (comme $\det(a\sigma+b) = a^2 \det\ \sigma + [unclear: 0] + b^2$
$$ \begin{cases} - \nu a^2 + b^2 = - \nu \\ 2b = 0 \qquad \text{(automatique en car. 2)} \end{cases} $$
qui n'impliquent nullement $b = 0$ ....

Considérons maintenant le
groupe produit
$$ B_0^\prime \times A_p \times \mathcal{N}_\sigma^\prime $$

\newpage

%% ===== Page 317 =====

et le morphisme de schémas

(25)
$$ \left\{ \begin{array}{l} B'_0 \times A_p \times N'_0 \longrightarrow G \times G \times E' \times E' \\ (\underline{u}, \underline{a}, \underline{b}) \longmapsto (\alpha=\underline{u}\underline{a}^{-1}, \beta=\underline{u}\underline{b}^{-1}, \mu=\det \underline{u}, \varepsilon=\varepsilon_{\sigma}(\underline{b})) \end{array} \right. $$ ,

je dis que l'image est dans $\mathcal{G}_{p,\sigma}^\ast$, ie.

qu'on aura bien

$$ \alpha \rho(\alpha)^{-1} = \varepsilon (\beta \sigma(\beta)^{-1}) \varepsilon_1(\mu - 1) $$

En effet

$$ \alpha \rho(\alpha)^{-1} = \underline{u} \rho(\underline{u})^{-1} $$

$$ \beta \sigma(\beta)^{-1} = \underline{u} \underline{b}^{-1} \sigma(\underline{b}) \sigma(\underline{u})^{-1} = \varepsilon \underline{u} \sigma(\underline{u})^{-1} $$
(car $\varepsilon_{\sigma}(\underline{b})^{-1} = \varepsilon_{\sigma}(\underline{b})$)

et on sait (par II) que

$$ \underline{u} \rho(\underline{u})^{-1} = \underline{u} \sigma(\underline{u})^{-1} \varepsilon_1(\mu - 1) $$

dans ce cas, puisque $a^2=1$. Mais de

plus le morphisme (25) est mono,

car en termes de $(\alpha, \beta, \mu, \varepsilon)$, $\underline{u}$ est déterminé
(en [unclear: termes] de $\alpha$ déjà), comme l'unique section de $B'_0$ qui

soit de la forme $\alpha \underline{a}^{-1}$, où $\underline{a} \in \Gamma T_p$

(car $T_p \cap B'_0 = \{1\}$), et $\underline{b}$ est alors déterminé

par la relation $\beta = \underline{u} \underline{b}^{-1}$. On trouve

donc que (25) est un iso

(26)
$$ \left\{ \begin{array}{l} B'_0 \times A_p \times N'_0 \overset{\sim}{\longrightarrow} \mathcal{G}_{p,\sigma}^\ast \\ (\underline{u}, \underline{a}, \underline{b}) \longmapsto (\alpha=\underline{u}\underline{a}^{-1}, \beta=\underline{u}\underline{b}^{-1}, \mu=\det \underline{u}, \varepsilon=\varepsilon_{\sigma}(\underline{b})) \end{array} \right. $$

\newpage

%% ===== Page 318 =====

459

Si on a deux sections $(\underline{u}, \underline{a}, \underline{b})$ et
$(\underline{u}', \underline{a}', \underline{b}')$ de $B'_0 \times \tilde{A}_P \times N'_0$, à quelle
condition a-t-on (27) $\tilde{\alpha} = \tilde{\alpha}'$, $\tilde{\beta} = \tilde{\beta}'$ ? Il
est évidemment [unclear: suffisant] qu'on ait

(28) $\underline{a}' = \lambda_1 \underline{a} \quad , \quad \underline{b}' = \lambda_2 \underline{b} \qquad \lambda_1, \lambda_2 \in \Gamma \tilde{G}$

et je dis que ces conditions sont
nécessaires. Cela provient du fait que
$B'_0 \hookrightarrow P_G$ [unclear: crossed out text] est mono (sur $B'_0$
lire $\tilde{B}'_0$), que $A_P \to P_G$ et
$N'_0 \to P_G$ sont [unclear: comme couronnes]
(soient $\tilde{T}_P$ et $\tilde{N}'_0$ les images), et que
dans $P_G$

(29) $\tilde{T}_P \cap \tilde{B}'_0 = \{1\}$ \quad ([unclear: $\tilde{T}_P \tilde{B}'_0 = \dots$])

de sorte que $\tilde{\underline{u}}$ est défini de façon
unique en termes de $\tilde{\alpha}$ par la relation

$\tilde{\alpha} = \tilde{\underline{u}} \tilde{\underline{a}}^{-1}$ \qquad [unclear: un terme de $\tilde{T}_P$]

— et qu'alors $\tilde{\underline{b}}$ est défini par la relation
$\tilde{\beta} = \tilde{\underline{u}} \tilde{\underline{b}}^{-1}$.

Ceci montre que l'on a un isomorphisme

(30) $\tilde{B}'_0 \times \tilde{A}_P \times \tilde{N}'_0 \xrightarrow{\sim} G_{P, \sigma}^* \subset P_G \times P_G \times E'$
$(\tilde{\underline{u}}, \tilde{\underline{a}}, \tilde{\underline{b}}) \longmapsto (\tilde{\underline{u}} \tilde{\underline{a}}^{-1}, \tilde{\underline{u}} \tilde{\underline{b}}^{-1}, \delta(\tilde{\underline{u}}))$

\newpage

%% ===== Page 319 =====

où
(31) $$ \delta(\tilde{u}) = det(\underline{u}) $$
quand $\tilde{u}$ provient de $\underline{u} \in \Gamma B'_0$ - ce qui
revient au même (cf p. 441) $\delta(\tilde{u})$ est
défini par
(32) $$ \tilde{u}(N_0) = \delta(\tilde{u}) N_0 $$ .

Rmq. Notons qu'une section $(\tilde{\alpha}, \tilde{\beta}, \mu)$ de [unclear: $g_{P,\sigma}$]

[MARGIN: NB le centralis de $P$ dans $G = PL_2 \ltimes P_2$ est $T_{0'} \times N_{0'}$, d'où [unclear]]

(mod op. de $N'_{0}$ sur $P$)

[MARGIN: NB $\tilde{\alpha}$ dépend du trio paramétrique de telle sorte que $\tilde{\gamma} = \tilde{\alpha} p \tilde{\alpha}^{-1}$ redonne $\gamma \in \Gamma P$ (cas $g_{P,\sigma} \simeq$ [unclear])]

est entièrement déterminée, par
ce qui précède, en termes de $\tilde{\alpha}$ -
la condition à satisfaire sur $\tilde{\alpha}$ pour
qu'il existe $\tilde{\beta}, \mu$ qui la complètent
en une section de $g_{P,\sigma}$, étant la
condition locale : 

[MARGIN: NB On a alg. locale : $\alpha \in \Gamma G$ de classe dans $\bar{g}$ donnée s'écrit $g = \left( \begin{matrix} A & B \\ C & D \end{matrix} \right)$ d'un choix du relèvement]

le coefficient $D$ de $\xi = \alpha p(\alpha)^{-1}$
est inversible (ce qui permet
de déterminer $\mu$ par $\mu = 1/D$), puis
on détermine $\tilde{y}$ par $\tilde{y} = \tilde{\xi} \tilde{\varepsilon}_1 (\mu-1)$, et
enfin on trouve $\tilde{\beta}$ par la condition
$\tilde{\beta} \tilde{\sigma}(\tilde{\beta})^{-1} = \tilde{y}$. De même, une section
$(\tilde{\alpha}, \tilde{\beta}, \mu, \varepsilon)$ de $\mathcal{G}^*_{P,\sigma}$ est déterminé par $\tilde{\alpha}$ seul
(mod $N_{0'}$ -- on trouve d'abord $\underline{u}, \underline{a}$ en termes

\newpage

%% ===== Page 320 =====

460

de $\alpha$ comme dessus, d'où $\mu = \det \underline{u}$, d'où
[unclear: crossed out line]
$\beta = \underline{u} \underline{b}^{-1}$, où $\underline{b} \in \Gamma \mathcal{N}'_0$ arbitraire, $\varepsilon$ étant
défini alors en termes des choix de $\underline{b}$
comme $\varepsilon = \varepsilon(\underline{b})$. La condition pour
que $\alpha$ fixé puisse se compléter par
$(\beta, \mu, \varepsilon) \leadsto (\alpha, \beta, \mu, \varepsilon) \in \Gamma \hat{\mathcal{G}}_{\vec{p}, \vec{o}}$ , est toujours
la condition que le coeff. $D$ de $\xi = \alpha \rho(\alpha)^{-1}$
$$(33) \quad \xi = \alpha \rho(\alpha)^{-1} = \left(\begin{array}{cc} A & B \\ C & D \end{array}\right)$$
soit inversible, ce qui permet de déterminer $\mu$ par
$$(34) \quad \mu = 1/D$$

Posons
$$(35) \quad \alpha = \left(\begin{array}{cc} a & b \\ c & d \end{array}\right)$$
et explicitons $\alpha \rho(\alpha)^{-1}$, avec $\rho = \left(\begin{array}{cc} 0 & 1 \\ v & 0 \end{array}\right)$

[MARGIN: $\delta = \det \alpha = a d - b c$]

$$ \alpha \rho(\alpha)^{-1} = \left(\begin{array}{cc} a & b \\ c & d \end{array}\right) \left(\begin{array}{cc} 0 & 1 \\ v & 0 \end{array}\right) (\frac{1}{\delta}) \left(\begin{array}{cc} d & -b \\ -c & a \end{array}\right) \left(\begin{array}{cc} 0 & \frac{1}{v} \\ 1 & 0 \end{array}\right) $$
$$ = - \frac{1}{\delta v} \left(\begin{array}{cc} b v & a \\ d v & c \end{array}\right) \left(\begin{array}{cc} b v & -d \\ -a v & c \end{array}\right) \quad , \quad d'où $$

[unclear: crossed out terms in the matrix calculation]

$$(36) \quad \begin{matrix} (-\delta v) A = v b(b v) - a(a v) \\ (-\delta v) B = -v b d + c a \\ (-\delta v) C = v d(b v) - c(a v) \\ (-\delta v) D = -v d^2 + c^2 \end{matrix}$$

la quantité la plus essentielle étant donc
$$(37) \quad (-\delta v) D = -v d^2 + c^2 = [unclear: crossed out text] \quad ,$$
et d'autre part on a $\delta = \det \alpha = ad-bc$
[unclear: et la convention] qui définit $\rho = \left(\begin{array}{cc} 0 & 1 \\ v & 0 \end{matrix}\right)$

\newpage

%% ===== Page 321 =====

et ...

(38) $$(-\delta v) B = a c + b c - v b d$$

ce qui redonne $u$ en termes de $a, b, c, d$,
grâce à la formule $\mathbb{IV}(48)$ (p. 453')

(39) $$u = \begin{pmatrix} 1 & B/D \\ 0 & 1 \end{pmatrix} = \begin{pmatrix} \frac{-v(ad-bc)}{-vd^2+cd+d^2} & \frac{ac+bc-vbd}{-vd^2+cd+d^2} \\ 0 & 1 \end{pmatrix}$$

Remarque. Par contre, la donnée de $\beta$ seule
(sans la "figure" $\varepsilon$)
ne permet pas de récupérer $\mu$. Mais
la donnée de $(\beta, \varepsilon)$ le permet, car on
en déduit le comm. associé

(40) $$\eta = \varepsilon \beta \varepsilon \beta^{-1} \overset{def}{=} \begin{pmatrix} R & S \\ T & U \end{pmatrix} ,$$

qui sera relié à $\zeta = \alpha \rho \alpha^{-1}$ par
$$\zeta = \eta \varepsilon(\mu^{-1})$$
et par le théorème de (II), on aura

(41) $$\mu = 1/U$$

(mais en termes de $\beta \rho \beta^{-1} \overset{def}{=} \begin{pmatrix} R_0 & S_0 \\ T_0 & U_0 \end{pmatrix} ,$
on aura

(42) $$\mu = \frac{\varepsilon}{U_0}$$ )

i.e. il n'est [unclear: plus] donné par $U_0$ seul, mais
seulement par $U_0$ et $\varepsilon$.

Pour préciser ces points, on va regarder
quand est-ce que, pour $(\beta, \varepsilon)$ fixés, on peut
compléter par $\alpha, \mu$ de façon à ...

\newpage

%% ===== Page 322 =====

461

section $(\alpha, \beta, \mu, c)$ de $\mathcal{G}_{\rho, \sigma}$. Une condition
nécessaire est que l'on ait $U_0$ inversible
(cela ne dépend que de $\beta$, $\rho$ et $\varepsilon$) —
et alors [unclear: raturé] la donnée de $\varepsilon$ en détermine uniquement
un $\eta = \varepsilon \beta \sigma \gamma \beta^{-1}$ sera dans une
section de $A_{\sigma}$ dans $G$ des relations
de $\det \eta = 1, \text{Tr } \eta \sigma = \text{Tr } \sigma$ (cf lemme,
p. 455), donc par le théorème des (II), pour
[unclear: raturé] on prend $\mu = 1/\mathcal{U} = \frac{\varepsilon}{U_0}$ [unclear: raturé]
et on définit
$$ \xi = \eta \ \varepsilon_1(\mu-1) \quad , \quad \text{loc. cit.} $$
$\xi$ pourra par la base s'écrire dans la
forme $\alpha \rho \alpha^{-1}$, l'indétermination étant
d'ailleurs dans $T_{\rho}$ via $\alpha \mapsto \alpha' = \alpha \underline{b} , \underline{b} \in T_{\rho}$.

Ceci montre bien que $\mu$ ne peut prendre,
pour un $\beta$ fixé, que les valeurs
$$ \frac{\varepsilon}{U_0} \quad , \quad \text{où } \beta \sigma \gamma \beta^{-1} = \begin{pmatrix} R_0 & S_0 \\ T_0 & U_0 \end{pmatrix} \dots $$

\newpage

%% ===== Page 323 =====

462

## (VI) Caractérisation des triplets $(E_{1},\rho, \sigma)$ modulaires.

1) Soient $G, SG, PG$ des formes associées
de $\operatorname{Gl}(2), \operatorname{Sl}(2), PGl(2)$ sur une base $S$ -
correspondant donc à un choix d'un
fibré en droites projectives $P$ sur $S$, ou une Algèbre
d'Azumaya $M$ de rang 4, et considérons
deux sections $\tilde{\rho}, \tilde{\sigma}$ de $PG$, telles que
(1) $\tilde{\varepsilon}_1 = \tilde{\rho} \tilde{\sigma}^{-1}$
soit une section unipotente régulière, i.e.
la section
(2) $\tilde{\varepsilon}_0 \overset{df}{=} \tilde{\rho}^{-1}(\tilde{\varepsilon}_1) = \tilde{\sigma}^{-1}(\tilde{\varepsilon}_1)$
s'en déduit soit [unclear]
(équivalente à : $\tilde{\varepsilon}_0$ t.q.c. que les sections
$\tilde{x}_0, x_1$ de $P$ qui correspondent à $\varepsilon_0, \varepsilon_1$
soient propres et distinctes en chaque fibre).

[MARGIN: NB si $\varepsilon_0, \varepsilon_1 \in G$ des relèvements de $\tilde{\varepsilon}_0, \tilde{\varepsilon}_1$, et $\varepsilon_1 = \rho \sigma^{-1}$, la condition $\rho \varepsilon_0 = \varepsilon_1$ (i.e. $\varepsilon_0 = \rho^{-1} \varepsilon_1$) entraîne $\sigma \varepsilon_0 = \sigma \rho^{-1} \varepsilon_1 = \varepsilon_1^{-1} \varepsilon_1 = e$ (?) (donc $\varepsilon_0 = \sigma^{-1}$). L'unicité de $(\rho, \sigma)$ s'en déduit sauf s'il y a des points de $S$ où il y a symétrie exacte entre $\varepsilon_0$ et $\varepsilon_1$ i.e. $x_0 \leftrightarrow x_1$ est induit par un autom. du "triangle".]

Alors il existe un isomorphisme unique
(3) $\operatorname{Gl}(2)_S \overset{\sim}{\to} G$
et des sections $u, v, w, w'$ de $\mathcal{O}_S$, vérifiant
(4) $u+v=0$, $w w' = 1$, $v \in \Gamma(\mathcal{O}_S^*)$
telles que $\tilde{\rho}, \tilde{\sigma}$ soient les images, par la
projection $\operatorname{Gl}(2)_S \overset{\sim}{\to} G \to PG$, des sections
(5) $\rho = \begin{pmatrix} 0 & w \\ v & u \end{pmatrix}$, $\sigma = \begin{pmatrix} 0 & -1 \\ w' & v \end{pmatrix}$

\newpage

%% ===== Page 324 =====

de $\operatorname{Gl}(2)_S$ - et inversement, pour tout
épinglage $\gamma$ et [crossed out text] $v, v', w, w'$
satisfaisant (4) - i.e.
(4bis) $$v' = -v, w' = 1-w, \quad v \in \Gamma(S, \mathcal{O}_S^\times)$$
$$w \in \Gamma(S, \mathcal{O}_S)$$
les images $\tilde{\rho}, \tilde{\sigma}$ par la composition $\operatorname{Gl}(2)_S \to PG \simeq$ [unclear]
des sections $\rho, \sigma$ de $\operatorname{Gl}(2)_S$ données par
(5) satisfont les conditions plus haut
($\tilde{\rho} \tilde{\sigma}^{-1} \stackrel{df}{=} \tilde{\varepsilon}_1$ est unipotente régulière, et
ponctue [unclear: son] équivalence : $\tilde{\sigma}^{-1}(\tilde{\varepsilon}_1) = \tilde{\rho}^{-1}(\tilde{\varepsilon}_1) \stackrel{df}{=} \tilde{\varepsilon}_0$).

On a alors
(6) $$\varepsilon_1 = \left( \begin{array}{cc} 1 & 0 \\ 1 & 1 \end{array} \right) \quad , \quad \varepsilon_0 = \left( \begin{array}{cc} 1 & 1/v \\ 0 & 1 \end{array} \right)$$

(7) $$\rho \sigma^{-1} = - \varepsilon_1 \quad , \quad \sigma^{-1} \rho = - \varepsilon_0$$
(7bis) $$\varepsilon_0 = \rho^{-1}(\varepsilon_1) = \sigma^{-1}(\varepsilon_1), \quad \varepsilon_1 = \rho(\varepsilon_0) = \sigma(\varepsilon_0)$$
dont les images par $\operatorname{Gl}(2)_S \to PG$ des
sections $\varepsilon_1, \varepsilon_0$ sont $\tilde{\varepsilon}_1, \tilde{\varepsilon}_0$. On a
les formules
(8) $$Tr \sigma = w', Tr \rho = w, Tr(\sigma + \rho) = 1$$
(9) $$\sigma + \rho = \left( \begin{array}{cc} 0 & 0 \\ 0 & 1 \end{array} \right) \stackrel{df}{=} N$$
(10) $$Tr \xi N = D \quad \text{si} \quad \xi = \left( \begin{array}{cc} A & B \\ C & D \end{array} \right)$$
(11) $$\det \sigma = \det \rho = -v = v'$$

\newpage

%% ===== Page 325 =====

2°) Les conditions suivantes sont équivalentes

(12)
\begin{tikzcd}
\boxed{Tr \sigma = 0} \ar[r, "\Leftrightarrow"] \ar[d, "\Updownarrow"] & Tr \rho = 1 \ar[r, "\Leftrightarrow"] \ar[d] & \boxed{\tilde{\sigma}^2 = 1} \ar[d] \\
\boxed{w' = 0} & \boxed{w = 1} & \sigma^2 \in \mathbb{G}_m \ar[d] \\
& & \boxed{\sigma^2 = v \text{ id.}}
\end{tikzcd}

Elles équivalent encore à ceci : désignant par $N'_\sigma$ le sous-schéma en groupes de $\operatorname{Gl}(2)_S$ défini par la condition
(13) $\tau(\sigma) = \varepsilon \sigma$ i.e. $\tau(\sigma) = \varepsilon \sigma$, $\varepsilon \in \mathbb{G}_m$
(avec nécessairement $\varepsilon^2 = 1$ i.e. $\varepsilon \in \mu_2$)
le morphisme

(14)
\begin{tikzcd}
N'_\sigma \ar[r] & \mu_2 \\
\tau \ar[r, |->] & \varepsilon
\end{tikzcd}

est un épimorphisme, de façon qu'on a une suite exacte

(15) $$1 \longrightarrow T_\sigma \longrightarrow N'_\sigma \longrightarrow \mu_2 \longrightarrow 1$$

où
(16) $T_\sigma = Cent_{\operatorname{Gl}(2)_S}(\sigma) = \{ \dots \}$
($N_\sigma \subset N'_\sigma$ est $\operatorname{Norm}_{\operatorname{Gl}(2)_S}(T_\sigma)$, les
[unclear]
une section de $N'_\sigma$ [unclear] vers $\mu_2$ est une section de
$N_\sigma$ [unclear] en fait. On a [unclear]
(18) $N'_\sigma = N_\sigma \quad (\text{car } \dots)$

La condition ci-dessus signifie que $\sigma, \rho$ prennent la forme

\newpage

%% ===== Page 326 =====

(19) $$ \sigma^- = \begin{pmatrix} 0 & -1 \\ v' & 0 \end{pmatrix} , \quad \rho = \begin{pmatrix} 0 & 1 \\ v & 1 \end{pmatrix} $$

avec $v + v' = 0$.

Ces résultats de 1°) 2°) ne sont qu'une récapitulation d'une partie des résultats des sections IV, V précédentes. [unclear: Il convient d'illustrer des conditions] qui reviennent à la détermination [unclear: des valeurs de $v, v'$ (sans avoir $w$)] de façon très explicite. On a à [unclear: se tourner] vers des conditions qui sont [unclear: relatives] [unclear: pour $v$ (ou encore $v'$)] de façon explicite.

3°) Conditions équivalentes
(20) $$ \begin{cases} \det \sigma = 1 \iff \det \rho = 1 \iff v = -1 \iff v' = 1 \\ \iff \varepsilon_0 = \begin{pmatrix} 1 & -1 \\ 0 & 1 \end{pmatrix} \end{cases} $$

Ces équivalences résultent aussitôt de (6) et (11) ! Les conditions signifient aussi que $\rho$ et $\sigma^-$ ont la forme

(21) $$ \sigma^- = \begin{pmatrix} 0 & -1 \\ 1 & w' \end{pmatrix} , \quad \rho = \begin{pmatrix} 0 & 1 \\ -1 & w \end{pmatrix} $$
$$ (w + w' = 1) $$

4°) Conjonction de 2°) et 3°).
Conditions équivalentes :

\newpage

%% ===== Page 327 =====

Conjonction des conditions de $2^\circ$ et des conditions de $3^\circ$
$$ \iff \sigma^2 = -1 \iff \tilde{\sigma}^2 = \tilde{\rho}^3 = 1 \text{ [unclear, crossed out box: } \implies \rho^2-\rho+1=0 \text{]} $$
i.e. $\sigma^2+1=0$
cf cons. plus bas

(22)
$$ v=-1, w=1 \iff \rho = \begin{pmatrix} 0 & 1 \\ -1 & 1 \end{pmatrix} \iff Tr \rho = \det \rho = 1 $$
$$ v'=1, w'=0 \iff \sigma = \begin{pmatrix} 0 & -1 \\ 1 & 0 \end{pmatrix} \iff (Tr \sigma = 0, \det \sigma = 1) $$

Dim. Le fait que la conjonction des conditions de $2^\circ$) et $3^\circ$) soit équivalente à chacune des conditions en $2^e$ et $3^e$ ligne, est trivial. L'équivalence avec $\sigma^2 = -1$ provient du fait que au [unclear: facteur] scalaire près, la condition $2^\circ$ signifie que $\sigma^2 = v id$, tandis que $3^\circ$ signifie $v=-1$, d'où $\sigma^2 = -1$ équivaut bien à la conjonction.

[MARGIN: Pour avoir une démonstration de l'équivalence de $2^\circ)+3^\circ)$ et de $\tilde{\sigma}^2=\tilde{\rho}^3=1$, cf ci-dessous]

Notons que $2^\circ$) s'écrit aussi $\tilde{\sigma}^2 = 1$, et équivaut à : $w'=0$, soit $w=1$, et
dans ce cas, en écrivant
$$ \rho = \begin{pmatrix} 0 & 1 \\ v & 1 \end{pmatrix} $$

(23)
$$ \rho^3 = \begin{pmatrix} v & v+1 \\ v(v+1) & 2v+1 \end{pmatrix} $$
donc $\rho^3$ ne peut être un scalaire i.e. $\tilde{\rho}^3 = 1$ que si les deux termes diagonaux sont égaux, or leur diff est $(2v+1) - v = v+1$,

\newpage

%% ===== Page 328 =====

donc la condition s'écrit $v=-1$, i.e. c'est
la condition 3°), mais alors l'expression
(23) de $\rho^3$ nous montre que $\rho^3=-1$.

Remarques Des conditions qui précèdent
impliquent les formules
(24) $\varepsilon_0 = \sigma \rho$
(25) $\varepsilon_1 = \rho \sigma$
(26) $\rho \varepsilon_1 = \varepsilon_0^{-1}$
(27) $\rho^2 \sigma = \varepsilon_0^{-1}$

[MARGIN: (sauf la dernière (27))]
dont chacune / les impliquent à son tour
- p.ex. on a $\varepsilon_0 = -\sigma^{-1}\rho$, donc (24)
équivaut à ($\sigma = -\sigma^{-1}$, i.e.) $\sigma^2=-1$, et (26) équivaut à
(25) en multipliant par $\rho$ aux deux mbres.

On a d'ailleurs
$$ \rho \varepsilon_1 = \begin{pmatrix} 0 & 1 \\ v & w \end{pmatrix} \begin{pmatrix} 1 & 1 \\ 0 & 1 \end{pmatrix} = \begin{pmatrix} 0 & 1 \\ v & v+w \end{pmatrix} , \quad \varepsilon_0^{-1} = \begin{pmatrix} 1 & -1 \\ 0 & 1 \end{pmatrix} $$
qui montre que $\rho \varepsilon_1 = \varepsilon_0^{-1}$ équivaut à
$w=1$, $v=-w \quad$ i.e. $\quad w=1, v=-1 \quad$ i.e.
(sur la forme $\rho \varepsilon_1 = \rho^2 \sigma$) $\rho = \begin{pmatrix} 0 & 1 \\ -1 & 1 \end{pmatrix}$,
OK. Comme (25) (et (26) - par suite)
sont conséquences des conditions
2°) + 3°), l'implication inverse résulte
d'un calcul brutal
$$ \rho^2 \sigma = \begin{pmatrix} -w & w-(v+w^2) \\ -v(v+w) & -w(v+w)+v(-1) \end{pmatrix} , \quad \varepsilon_0^{-1} = \begin{pmatrix} 1 & -1 \\ 0 & 1 \end{pmatrix} $$
c'est impossible on ne peut avoir
l'égalité $\rho^2 \sigma = \varepsilon_0^{-1}$ qui lie $-v(v+w)=0$ i.e. $v=-w^2$
qui implique (car $w=-1$ !) $v+w=0$, [unclear]

\newpage

%% ===== Page 329 =====

$$\rho^2 \sigma = \left( \begin{matrix} w^3 & w \\ 0 & \text{[unclear]} \end{matrix} \right)$$

[MARGIN: 265]

[MARGIN: On a vu plus haut que $v+w^2=0$, $w^3=1 \Rightarrow w=1, v=-1$ OK]

cela montre que :
(27 bis) $\quad \rho^2 \sigma = \varepsilon_0^{-1} \iff v+w^2=0, w^3=1 \quad (\Rightarrow v=-w^2=-1/w)$

Mais $(27')$, multipliant par $\rho$, équivaut à

$$\rho^3 \sigma = \rho \varepsilon_0^{-1} \quad (7')$$
or on sait que $\varepsilon_0 = - \sigma^{-1} \rho$ i.e. $\varepsilon_0^{-1} = - \rho^{-1} \sigma$
i.e. $\rho \varepsilon_0^{-1} = - \sigma$, donc la relation précédente
s'écrit aussi $\rho^3 \sigma = - \sigma$ ou encore
(28) $\quad \boxed{\rho^3 = -1}$
de (28) on tire immédiatement (27 bis).

On peut voir algébriquement directement
en calculant $\rho^3$ en fonction de $v$ et $w$, on
trouve
(29) $\quad \rho^3 = \left( \begin{matrix} v w & v+w^2 \\ v(v+w^2) & w(2v+w^2) \end{matrix} \right)$

dit qu'on ne peut avoir $\tilde{\rho}^3 = 1$ i.e. $\rho^3$ un
scalaire que si $v+w^2=0$
i.e. $v=-w^2$
ce qui implique
$$\rho^3 = \left( \begin{matrix} -w^3 & 0 \\ 0 & -w^3 \end{matrix} \right)$$
i.e. on trouve
Lemme. On a $\tilde{\rho}^3 = 1$ ssi $v+w^2=0$ i.e.
$v=-w^2$, ie. on a
(30) $\quad \rho = \left( \begin{matrix} 0 & 1 \\ -w^2 & w \end{matrix} \right)$
et alors on a
(31) $\quad \rho^3 = -w^3 \cdot 1 \quad \text{et}$
donc on a
(32) $\quad \rho^3 = -1 \iff w^3 = 1 \text{ i.e. } w \in \mu_3 \text{ et } v = -w^2 (= -1/w)$

328

\newpage

%% ===== Page 330 =====

Corollaire 1 Les conditions équivalentes $(A2)$
équivalent aussi à

(33) $\rho^3 = -1$ et $det \rho = 1$

ou encore :

(34) $\rho^3 = -1$ et $Tr \rho = 1$

en d'autres termes, considérant les
trois relations

(35) $det \rho = 1$, $Tr \rho = 1$, $\rho^3 = -1$

la conjonction de deux d'entre elles-ci impliquent la troisième,
et équivaut aux conditions équivalentes
de $(A2)$.

Notons une autre façon de trouver ce
corollaire, en notant [unclear: que]

(35) $\rho^3+1 = (\rho+1)(\rho^2-\rho+1)$ (identité algébrique)

or $\rho$ n'est pas un scalaire sur aucune fibre,
[MARGIN: faux] $\rho+1$ est inversible, donc $\rho^3=-1$ i.e. $\rho^3+1=0$
équivaut à la relation $\rho^2-\rho+1=0$ :

(36) $\rho^2-\rho+1=0$, ($\iff \rho^3=-1$),

qu'on compare avec Hamilton-Cayley

(37) $\rho^2 - (Tr \rho)\rho + det \rho = 0$

donc $Tr \rho = det \rho = 1$ implique (36),
d'autre part si on a (36), on conclut
via 37

$(Tr \rho - 1)\rho = (det \rho - 1)1$

\newpage

%% ===== Page 331 =====

Une autre façon de trouver certaines
équivalences est via les équations de Hamilton-
Cayley pour $\sigma$ et $\rho$

(36) $$ \sigma^2 - (\underbrace{Tr \sigma}_{w'=1-w}) \sigma + \underbrace{det \sigma}_{v'=-v} = 0 $$

(37) $$ \rho^2 - (\underbrace{Tr \rho}_{w}) \rho + \underbrace{det \rho}_{v} = 0 $$

L'équation (36) redonne que la conjonction
de $Tr \sigma = 0$, $det \sigma = 1$ implique $\sigma^2+1=0$ i.e.
$\sigma^2=-1$, et inversement, compte tenu que
$1$ et $\sigma$ sont [unclear: lin. indépendantes] dans $M_2(\mathcal{O}_S)$.

De même (37) montre

(38) $(Tr \rho = 1 \text{ et } det \rho = 1) \iff (\rho^2 - \rho + 1 = 0)$.

et l'on retiendra l'implication

(39) $\rho^2 - \rho + 1 = 0 \implies \rho^3+1=0 \text{ i.e. } \rho^3=-1$

qui résulte de l'identité algébrique

$$ \rho^3+1 = (\rho+1)(\rho^2 - \rho + 1) \quad . $$

Donc
Corollaire 2. Les conditions équivalentes de
(22) équivalent aussi à la relation
(40) $\rho^2 - \rho + 1 = 0$.

\newpage

%% ===== Page 332 =====

5°) Théorème d'épi-génèse modulaire.

Théorème Soient $G, SG, PG$ des formes [unclear: (sans relation)] sur un schéma $S$ de $\operatorname{Gl}(2), \operatorname{Sl}(2), P\operatorname{Gl}(2)$ (associées), et $\tilde{\sigma}, \tilde{\rho}$ deux sections de $PG$.

5°) Un retour sur 1°), relatif à la situation
(41) $$ \begin{cases} 1^o) \quad \tilde{\varepsilon}_0 = \tilde{\sigma}^{-1} \tilde{\rho} \text{ unipotents réguliers } (i.e. \\ \tilde{\varepsilon}_1 = \tilde{\rho} \tilde{\sigma}^{-1} = \tilde{\rho}(\tilde{\varepsilon}_0) = \tilde{\sigma}(\tilde{\varepsilon}_0) \text{ unipotents réguliers}) \text{ et} \\ 2^o) \quad \tilde{\varepsilon}_0 \text{ et } \tilde{\varepsilon}_1 \text{ sont équivalentes } \text{[unclear: en d. gen. fibré]} \end{cases} $$

Soient $\varepsilon_0, \varepsilon_1$ les relèvements unipotents réguliers dans $G$ s'identifiant à des sections de $\mathcal{M}$

[MARGIN: $N_0$ et $N_1$ strictement nilp. i.e. d'une det. et de trace nuls (42)]
(43) $\varepsilon_0 = 1 + N_0$, $\varepsilon_1 = 1 + N_1$
$$ \tilde{\rho}(\varepsilon_0) = \varepsilon_1 \iff \tilde{\sigma}(N_0) = \tilde{\rho}(N_0) = N_1 $$

On peut alors relever $\tilde{\sigma}$ et $\tilde{\rho}$ de façons canoniques en des sections $\sigma, \rho$ de $G$ (donc de $\mathcal{M}$) satisfaisant les conditions suivantes :

(44) $$ \begin{cases} 1^o) \quad -\sigma^{-1} \rho = \varepsilon_0 \quad (\iff -\rho \sigma^{-1} = \varepsilon_1) \\ \quad \iff \rho + \sigma = -N_0 \sigma (= - \sigma N_0) \\ \quad \iff \rho + \sigma = N_1 \rho (= \rho N_0) \\ 2^o) \quad (\rho + \sigma)^2 = u(\rho + \sigma) \text{ avec } u \in \Gamma \mathbb{G}_m \end{cases} $$

De plus, il existe un unique morphisme (épi-génèse)
$$ G_0 = \operatorname{Gl}(2) \xrightarrow{\phi} G $$
(45) $$ \begin{cases} \rho_0, \sigma_0 \in \Gamma(G_0), \rho_0 = \begin{pmatrix} 0 & 1 \\ -v & w \end{pmatrix}, \sigma_0 = \begin{pmatrix} 0 & -1 \\ v' & w' \end{pmatrix} \\ (\text{avec } v' = -v, w' = 1 - w), v \in \Gamma \mathbb{G}_m \end{cases} $$
tel qu'on ait
(46) $\phi(\rho_0) = \rho$, $\phi(\sigma_0) = \sigma$
qui impliquent
(47) $\phi( (\varepsilon_0)_0 ) = \varepsilon_0$, $\phi( (\varepsilon_1)_0 ) = \varepsilon_1$

\newpage

%% ===== Page 333 =====

457
		
où [crossed out text] $(\varepsilon_1)_0 = -\rho_0 \sigma_0^{-1} =$

(48) $$ \varepsilon_0 = -\sigma_0^{-1} \rho_0 = \begin{pmatrix} 1 & \frac{1}{\nu} \\ 0 & 1 \end{pmatrix} , \quad \varepsilon_{10} = \begin{pmatrix} 1 & 0 \\ \nu & 1 \end{pmatrix} $$

— et inversement, pour $\rho_0, \sigma_0$ comme dessus,
et tout hom. $\varphi : G_0 \simeq \operatorname{Gl}(2)_{\hat{\mathbb{Z}}} \to G$, les images
$\tilde{\rho}, \tilde{\sigma}$ de $\rho_0, \sigma_0$ dans $PG$ satisfont (41), et les
images $\rho, \sigma$ de $\rho_0, \sigma_0$ par $\varphi$ sont des relèvements
de $\tilde{\rho}, \tilde{\sigma}$ satisfaisant (44).

Théorème d'épinglage modulaire. Soient $G, SG, PG$
des formes (aux schémas $S$) de $\operatorname{Gl}(2), \operatorname{Sl}(2), PGl(2)$
associées à une Alg. d'Azumaya $M$ de rg. 4.

1) Appelons [crossed out text] triplet modulaire de $PG$ un
triplet de sections $(\tilde{\rho}, \tilde{\sigma}, \tilde{\varepsilon}_0)$ de $PG$,
satisfaisant les conditions (49) (50) (51).

[crossed out text]

(49) $$ \tilde{\varepsilon}_0 = \tilde{\sigma}^{-1} \tilde{\rho} \quad i.e. \quad \tilde{\rho}^{-1} \tilde{\sigma} \tilde{\varepsilon}_0 = 1 $$

[ de sorte que
(49 bis) $$ \tilde{\varepsilon}_1 \overset{def}{=} \tilde{\rho} \tilde{\sigma}^{-1} = \tilde{\rho}(\tilde{\varepsilon}_0) = \tilde{\sigma}(\tilde{\varepsilon}_0) \quad ] \ ; $$

(50) $$ \tilde{\rho}^3 = \tilde{\sigma}^2 = 1 , \quad \tilde{\varepsilon}_0 \text{ est unipotente régulière} $$
$$ (\iff \tilde{\varepsilon}_1 \text{ est unipotente régulière}) \ ; $$

(51) $\tilde{\varepsilon}_0$ et $\tilde{\varepsilon}_1$ [unclear: possèdent un] équivalent d. gen. fibr.
(ce qui signifie, en introduisant les relèvements
$$ \varepsilon_0 = 1 + N_0 , \quad \varepsilon_1 = 1 + N_1 \quad (N_0 \text{ et } N_1 \text{ strictement} $$
$$ \text{nilpotents... cf (42)}) $$
dans $G$, que $N_0$ et $N_1$ sont homothétiques sur les
fibres).

2) Appelons triplet modulaire de $G$ un

[MARGIN: 332]

\newpage

%% ===== Page 334 =====

triplet de sections $(\rho, \sigma, \varepsilon_0)$ satisfaisant les
[unclear: cinq] conditions suivantes

(52) $-\sigma^{-1} \rho = \varepsilon_0 \qquad (\iff \rho^{-1} \sigma \varepsilon_0 = -1 .)$
$\left[ \iff - \rho \sigma^{-1} = \varepsilon_1 \text{ où } \varepsilon_1 = \rho \varepsilon_0 \rho^{-1} \right.$
$\left. \iff - \rho \sigma^{-1} = \varepsilon_1 \text{ et } \varepsilon_1 = \sigma \varepsilon_0 \sigma^{-1} \right]$

(53) $(\rho + \sigma)^2 = \rho + \sigma$

[MARGIN: $\det \varepsilon_0 = 1$
$Tr \varepsilon_0 = 2$
$\varepsilon_0 \neq 1$ en ch fibre]
(54) $\varepsilon_0$ est unipotente régulière i.e. $\varepsilon_0 = 1 + N_0$ avec
$N_0$ nilpotent régulier ($Tr N_0 = \det N_0 = 0$, $N_0$ non
nul en chaque fibre).

$(\iff \varepsilon_1 = \rho \varepsilon_0 \rho^{-1} = \sigma \varepsilon_0 \sigma^{-1} \text{ est unipot. rég})$

(55) $\varepsilon_0$ et $\varepsilon_1$ sont inéquivalentes en chaque fibre
$(\iff N_0 \text{ et } N_1 \text{ non homoth en ch fibre})$

(56) On a l'une des quatre conditions équivalentes
[ $\sigma^2+1=0$ ] [ $\rho^2-\rho+1=0$ ]
[ $Tr \sigma=0, \det \sigma=1$ ] [ $Tr \rho = \det \rho = 1$ ]

3) Appelons triplet modulaire standard de
$\operatorname{Gl}(2)_S$ (resp. dans $PGl(2)_S$) le triplet

(57) $(\rho_0 = \begin{pmatrix} 0 & 1 \\ -1 & 1 \end{pmatrix}, \sigma_0 = \begin{pmatrix} 0 & -1 \\ 1 & 0 \end{pmatrix}, \varepsilon_0 = \begin{pmatrix} 1 & -1 \\ 0 & 1 \end{pmatrix})$

(resp. son image $(\tilde{\rho}_0, \tilde{\sigma}_0, \tilde{\varepsilon}_0)$ dans $PGl_S$) -
qui est bel et bien un triplet modulaire.
[unclear: le désignerons comme d'hab]

Ceci posé, on a des applications
canoniques

\newpage

%% ===== Page 335 =====

(58)
$$
\begin{tikzcd}
Isom(\operatorname{Gl}(2)_S, G) \ar[r, "\varphi \mapsto (\varphi(\rho_0), \varphi(\sigma_0), \varphi(\varepsilon_0))"] \ar[d, "\simeq"'] & TripMod(G) \ar[d] & (\rho, \sigma, \varepsilon_0) \ar[d, mapsto] \\
Isom(PGl(2)_S, PG) \ar[r, "\psi \mapsto (\psi(\tilde{\rho}_0), \psi(\tilde{\sigma}_0), \psi(\tilde{\varepsilon}_0))"'] & Tripmod(PG) & (\tilde{\rho}, \tilde{\sigma}, \tilde{\varepsilon}_0)
\end{tikzcd}
$$

et on a le

Théorème d'épinglage. Toutes les flèches
de (58) sont bijectives.

Ce théorème ne fait que résumer certains
des résultats principaux des reflexions
du présent § VI. Notons que l'injectivité
des flèches horizontales signifie qu'un
hom. $\operatorname{Gl}(2)_S = G_0 \to G$ (resp. $PGl(2)_S \to PG$)
est connu quand on connait ses effets sur
$\sigma$ et $\rho$ (donc aussi sur $\varepsilon_0$) (resp. quand on
connait ses effets sur $\tilde{\sigma}$ et $\tilde{\rho}$) - ce qui
signifie aussi, [unclear: à le dire] pour les
automorphismes de $\operatorname{Gl}(2)_S$ (resp. de $PGl(2)_S$)
ce qui signifie aussi que
(59) $Z_\rho \cap Z_\sigma = Cent \ G_0$
(l'intersection des centralisateurs de $\rho, \sigma$ dans
$\operatorname{Gl}(2)_S$ est [crossed out] au centre $Cent \ G_0$) -
resp. la relation plus forte
(60) $N_\rho \cap N_\sigma = Cent \ G_0$

\newpage

%% ===== Page 336 =====

où $N'_{\tilde{\rho}}$, $N'_{\tilde{\sigma}}$ sont les sous-schémas en groupes de $\operatorname{Gl}(2)_{\mathbb{Z}}$ définis [unclear: comme] stabilisateurs de $\tilde{\rho}$, $\tilde{\sigma}$ dans $\operatorname{Gl}(2)_{\mathbb{Z}}$, i.e.

$$ (61) \begin{cases} N'_{\tilde{\rho}} = \{ g \mid g(\tilde{\rho}) = \varepsilon \tilde{\rho}, \varepsilon \text{ scalaire} \} \\ N'_{\tilde{\sigma}} = \{ g \mid g(\tilde{\sigma}) = \varepsilon \tilde{\sigma}, \varepsilon \text{ scalaire} \} \end{cases} $$
[MARGIN: N.B. [unclear] avec $\varepsilon^2=1$]

Rmq Il est plausible, [unclear: au vu de ces résultats], que $GP(1)_{\mathbb{Z}}$ en tant que faisceau en groupes "absolu" (sur $\operatorname{Spec}\ \mathbb{Z}$) pour f.p.q.f., est engendré par les hom.
$$ \mathbb{Z}/2\mathbb{Z} \xrightarrow{\alpha \mapsto \tilde{\sigma}^\alpha} GP(1)_{\mathbb{Z}} $$
$$ \mathbb{Z}/3\mathbb{Z} \xrightarrow{\beta \mapsto \tilde{\rho}^\beta} GP(1)_{\mathbb{Z}} $$
$$ G_m \xrightarrow{\lambda \mapsto \tilde{\varepsilon}_0(\lambda)} GP(1)_{\mathbb{Z}} $$
satisfaisant seuls la relation ($\Omega$)
$$ (62) \quad \tilde{\rho}^{-1} \tilde{\sigma} \tilde{\varepsilon}_0(1) = 1 $$

- ce vaudrait bien la peine de tirer cette question au clair, ainsi que la question analogue pour $\operatorname{Gl}(2)_{\mathbb{Z}}$, où les [unclear: dits] groupes $\mathbb{Z}/2\mathbb{Z}, \mathbb{Z}/3\mathbb{Z}$ doivent se remplacer par $\mathbb{Z}/4\mathbb{Z}, \mathbb{Z}/6\mathbb{Z}$, et la relation (62) par
$$ (63) \quad \tilde{\rho}^{-1} \tilde{\sigma} \tilde{\varepsilon}_0(1) = \tilde{\rho}^3 = \tilde{\sigma}^2 $$

\newpage

%% ===== Page 337 =====

469

Rmq J'ai la nette impression que le paquet de conditions (51) à (56) est surabondant, et se reduit à p.ex. si l'on se borne aux conditions algébriques simples

[MARGIN: NB Il faut peut-être préciser que les conditions $Tr \sigma = 0$, $det \sigma = 1$, $Tr \rho = \pm 1$, $det \rho = 1$ sont vigoureusement satisfaites fibre par fibre.]

$$ (64) \qquad \sigma^2 + 1 = \rho^2 - \rho + 1 = 0 , \quad (\rho + \sigma)^2 = \rho + \sigma $$
(qui définit $\varepsilon_0$ par (51)), [unclear: s'impliquent par les autres conditions].

La condition (65) signifie que
$$ N = \rho + \sigma $$

$$ (66) \qquad N^2 = N , $$

et en termes de $N$ on aura $\sigma^{-1} N = \sigma^{-1} \rho + 1$ i.e.

$$ (67) \left\{ \begin{array}{ll} \varepsilon_0 \overset{df}{=} -\sigma^{-1} \rho = 1 - \sigma^{-1} N = 1 + N_0 & N_0 = - \sigma^{-1} N = \sigma N \\ \varepsilon_1 \overset{df}{=} -\rho \sigma^{-1} = 1 - N \sigma^{-1} = 1 + N_1 & N_1 = - N \sigma^{-1} = N \sigma \end{array} \right. $$

et la condition $\varepsilon_0$ stricte s'exprimera par les conditions

$$ (68) \qquad det N = Tr N = 0 , \quad \sigma N \text{ nul en toute fibre,} $$
i.e. $N$ nul en toute fibre.

Or la [unclear: troisième] condition (68) résulte des autres, car si on avait $N=0$ sur une fibre, on pourrait trouver un $S$ spectre d'un corps, $N=0$, i.e. $\rho = -\sigma$, d'où l'on aurait à la fois

$$ \sigma^2 + 1 = 0 \qquad \sigma^2 + \sigma + 1 = 0 \qquad \text{d'où } \sigma = 0, \text{ absurde} $$

Je dis aussi qu'on n'aura aucune fibre où $N=1$, sinon on aurait en cette fibre

$$ \rho + \sigma = 1 \text{ et en substituant dans} $$
$$ \sigma^2 + 1 = 0 \implies (1-\rho)^2 + 1 = 0 \implies \rho^2 - 2\rho + 2 = 0 $$

\newpage

%% ===== Page 338 =====

comparant avec $\rho^2-\rho+1=0$, on trouve par
différence $\rho-1=0$ i.e. $\rho=1$, mais
la relation $\rho^2-\rho+1=0$ devient $1=0$, absurde.

Donc l'opérateur idempotent $N$ est partout
sur chaque fibre
de rang $1$, donc un [unclear: thm] (sauf erreur)
dit que $N$ est décomposable (un tenseur
de $M \simeq \check{F} \otimes F$) loc^t décomposable sous la forme
$x' \otimes x$, avec $\langle x, x' \rangle = 1$, et par
suite $Tr \ N = 1$, $det \ N = 0$
donc la relation $det (\sigma^{-1} N) = 0$ dans (68) est
également automatique, puisque $det \sigma N$
$= det \sigma \ det N = 0$. Il reste la seule relation
$$Tr \ \sigma N = 0$$
(que j'écris aussi
(69) $$Tr \ \sigma \rho = -2 \quad (\Leftrightarrow Tr \ \sigma^{-1} \rho = 2)$$
pour laquelle il faudrait voir si
elle résulte des autres relations.

Prenons
(70) $$ \begin{cases} \rho = \rho_0 = \begin{pmatrix} 0 & 1 \\ -1 & 1 \end{pmatrix} \\ \sigma = u \ \sigma_0 \ u^{-1} = \begin{pmatrix} a & b \\ c & d \end{pmatrix} \begin{pmatrix} 0 & -1 \\ 1 & 0 \end{pmatrix} \begin{pmatrix} d & -b \\ -c & a \end{pmatrix} \\ \text{où } u = \begin{pmatrix} a & b \\ c & d \end{pmatrix} \in \operatorname{Sl}(2), \text{i.e. } ad-bc = 1 \end{cases} $$
on trouve

\newpage

%% ===== Page 339 =====

(70bis) $$ \sigma = \left(\begin{array}{cc} bd+ac & -(a^2+b^2) \\ c^2+d^2 & -(bd+ac) \end{array}\right) $$

d'où
$$ N = \rho + \sigma = \left(\begin{array}{cc} (bd+ac) & 1-(a^2+b^2) \\ -1+(c^2+d^2) & 1-(bd+ac) \end{array}\right) $$

$$ N-1 = \left(\begin{array}{cc} (bd+ac)-1 & 1-(a^2+b^2) \\ -1+(c^2+d^2) & -(bd+ac) \end{array}\right) = -det(N) N^{-1} \text{ si } det N \text{ inv.} $$

et on trouve
(71) $$ N^2 - N = N(N-1) = - det(N) . 1 $$
donc la condition
(72) $$ N^2 = N \iff det N = 0 \quad , \quad (\Longrightarrow Tr N = 1) $$

or on a
$$ det N = bd+ac - (bd+ac)^2 $$
$$ - \left[ -1 - (a^2+b^2)(c^2+d^2) + a^2+b^2 + c^2+d^2 \right] $$
$$ = bd+ac - b^2d^2 - a^2c^2 - 2abcd + 1 + a^2c^2+a^2d^2+b^2c^2+b^2d^2 $$
$$ - a^2-b^2-c^2-d^2 $$
$$ = (bd+ac) - (a^2+b^2+c^2+d^2) + \underbrace{(ad-bc)^2}_{1} + 1 $$
$$ = 2 + (bd+ac) - (a^2+b^2+c^2+d^2) $$

donc la condition (72) équivaut à
(73) $$ a^2+b^2+c^2+d^2 = 2 + (bd+ac) $$

par ailleurs on a par 70), (70bis)
$$ \rho \sigma = \left( \begin{array}{cc} c^2+d^2 & -(bd+ac) \\ (c^2+d^2)-(bd+ac) & a^2+b^2-(bd+ac) \end{array} \right) $$

donc
(74) $$ Tr \rho \sigma = a^2+b^2+c^2+d^2 - (bd+ac) = 2 $$

\newpage

%% ===== Page 340 =====

Ce calcul suggère fortement que
les relations (64) impliquent déjà (68)..
c'est à dire pas de choses pires aux dénominateurs....
Si on admet ce point, il
reste seulement à vérifier que les
conditions précédentes (64) (i.e. essentiellement déjà
(61) à (56) sauf (55)) impliquent en sus (55),
i.e. l'indépendance de $\varepsilon_0, \varepsilon_1$ en toute fibre,
i.e. que $N_0, N_1$ sont non-homothétiques en
tt fibre. Or (67)

(*) $$N_1 = \tilde{\sigma}(N_0)$$ [unclear crossed out text]

Si on avait sur une fibre
$$N_1 = \lambda N_0$$
on en concluerait, en [unclear] utilisant
et $\tilde{\rho}^3=1$ les relations
[unclear crossed out text]

$\tilde{\sigma}^2 = 1$, 
$$ \begin{cases} N_1 = \lambda N_0 \\ \lambda^2=1 \end{cases} $$ 
donc [unclear crossed out text] $N_0 = \lambda^2 N_0$ d'où

Mais un calcul [unclear: immédiat] donne aussi
(**) $$N_0 = \tilde{\rho}^{-1} N_1$$
$$\lambda^2 = 1, \quad \lambda^3 = 1 \quad \text{d'où} \quad \lambda = 1$$
i.e. $N_0 = N_1$ i.e. $N_0$ commute à $\sigma, \rho$
donc aussi $N = - \sigma N_0$ [with $= \tilde{\rho}^{-1} N_0$ above] commute à $\sigma, \rho$,
et comme $N = \sigma + \rho$, $\sigma$ et $\rho$ commutent.
Mais alors

\newpage

%% ===== Page 341 =====

471

$$(\sigma+\rho)^2 = \underbrace{\sigma^2}_{-1} + \underbrace{\rho^2}_{\rho-1} + 2\sigma\rho = \rho - 2 + 2\sigma\rho =$$

donc la relation $(\sigma+\rho)^2 = \sigma+\rho$ s'écrit

$\rho - 2 + 2\sigma\rho = \sigma + \rho$ i.e.

$2\sigma\rho = \sigma + 2$ soit $2\sigma^2\rho = \sigma^2 + 2\sigma$ i.e. [unclear: $2(-\rho) = -1 + 2\sigma$]

[Le passage suivant est barré :]
Prenons les traces des deux mbres, en utilisant
$Tr \sigma\rho = 2 (\neq 4)$, $Tr \sigma = 0$, $Tr 1 = 2$, on trouve
$4=4$, pas de contradiction, mais prenant les
déterminants des deux mbres (et en admettant
que $\det \sigma = \det \rho = 1$, ce qui est une conséquence
de (64) pour peu qu'on exige de plus
que $\sigma, \rho$ soient réguliers i.e. ne soient des
scalaires sur aucune fibre) on trouve
(**) [unclear] $= 2^2 + 2 Tr \sigma + (\det \sigma)^2$ i.e. [unclear] $= 0$
ce qui est absurde.
[Fin du passage barré]

(**) $2 N = 1 ,$

d'où, en prenant les déterminants
$4 \det N = 1$ i.e. $0 = 1$ , absurde.

Je ne serais pas étonné
qu'on ait ainsi des conditions équivalentes
à (51) - (56) sous la forme

(75) $\sigma^2+1 = \rho^2-\rho+1 = 0$, $\varepsilon_0 \overset{dfn}{=} -\sigma\rho$ est unipotent régulier

340

\newpage

%% ===== Page 342 =====

en exigeant de plus, au besoin, que $\sigma, \rho$ soient
également réguliers i.e. ne soient des scalaires
sur aucune fibre. Il faudrait en
conséquence que cela implique que
$$N \stackrel{dfn}{=} \rho + \sigma$$
satisfait bien la relation d'idempotence
(75) $$N^2 - N = 0$$
Ce qui équivaut, si $N$ régulier, à :
(77) $$Tr N = 1, det N = 0 \quad (?)$$
Mais $N$ est régulier, car si sur une fibre
[unclear: valant] un scalaire on aurait que (en
se restreignant à cette fibre) $\rho$ et $\sigma$
commutent, on aurait
$$-\sigma \rho = \sigma \rho = \varepsilon_0 = 1 + N_0$$
[unclear crossed out line]
avec
$$(1-\sigma \rho)^2 = 0 \quad i.e. \quad 1 - 2\sigma \rho + \underbrace{\sigma^2 \rho^2}_{-1 \ -1} = 0 \quad i.e.$$
$$1 - 2\sigma \rho + 1 = 0 \quad i.e. \quad 2 = 2\sigma \rho$$
soit, en multipliant par $\sigma$
$$2\sigma = \rho \sigma + 2\underbrace{\sigma^2}_{-1} \rho \quad i.e. \quad 2(\sigma + \rho) = \underbrace{\rho \sigma}_{N}$$
Mais par hyp. on a $N = \lambda \cdot 1$, $\lambda$ scalaire,
d'où $\rho \sigma = 2\lambda$ , et prenant les
déterminants $\det(\rho \sigma) = 4 \lambda^2$
i.e. $\underbrace{1}_{1} \cdot 2 \quad i.e \quad \lambda = \frac{\varepsilon}{2} , \text{ avec } \varepsilon^2 = 1,$
i.e. $\rho \sigma = \varepsilon \quad , \quad \varepsilon^2 = 1$
$\underbrace{\rho^2 \sigma^2}_{1} = 1 \quad , \quad \underbrace{\rho^2}_{-1} = -1 \quad , \quad i.e. \quad 2=0, \text{ absurde}$

\newpage

%% ===== Page 343 =====

472

Mais $Tr N = 0$ [unclear: pas]
Donc $N$ est régulier, il reste à prouver (77).
Mais $Tr N = \underbrace{Tr \sigma}_{0} + \underbrace{Tr \rho}_{1} = 1$, OK, reste
à prouver $det(\rho + \sigma) = 0$ mais
$det(\rho + \sigma) =$ [crossed out: $\underbrace{det \rho}_{1} + \underbrace{det \sigma}_{1} + Tr(\sigma^{-1}\rho)$] est la relation
$\underbrace{det(\sigma(1 + \sigma^{-1}\rho))}_{\text{qui équivaut à}} = det \sigma \cdot det(1 + \sigma^{-1}\rho) \quad = \underbrace{det \sigma^{-1}\rho}_{(= det \sigma det \rho)} + \underbrace{Tr(\sigma^{-1}\rho)}_{- Tr \sigma \rho} + 1$

[MARGIN: calculs déjà faits ? i.e. v. au dos (cf formule (69))]

$det(\rho+\sigma) = 2 + Tr(\sigma^{-1}\rho) = 2 - Tr \sigma \rho \quad,$ donc
la condition $det N = 0$ équivaut à : (69)
$$Tr \sigma \rho = 2$$
(c'est du moins juste).

Pour la prouver (cette relation), on peut
songer à procéder comme tantôt, sauf que je
suis parti de (68) au lieu de (75). Ce
n'est pas entièrement élucidé - à revoir
: le danger!

\newpage

%% ===== Page 344 =====

473

VII Récapitulation sur les [unclear: épi-gèbres].

On désigne par $G, SG, PG$ des formes de $\operatorname{Gl}(2_S), \operatorname{Sl}(2_S), P\operatorname{Gl}(2_S)$ associées à une algèbre d'Azumaya $M$, i.e. à un fibré en droites projectives $P$ sur $S$.

1) Sections régulières de $G$. (pas nec. inv.)

Soit $\rho$ une section de $M$, posons
(1) $$v = \det \rho, \quad w = \text{Tr} \rho$$
de sorte qu'on a l'équation d'Hamilton-Cayley
(2) $$\rho^2 - w\rho + v = 0 \quad \text{i.e. } \rho^2 = w\rho - v$$

Proposition. Conditions équivalentes :

a) $\rho$ n'est un scalaire sur aucune fibre.

[MARGIN: NB À cause de (2), $A_\rho$ est une sous-algèbre de $M$.]
a') $A_\rho = \underline{O}_S + \underline{O}_S \rho$ est un sous-module localement facteur direct libre de rang 2 de $M$.

b) (Si $M = \underline{\text{End}}_{\mathcal{O}_S}(F)$) Il existe loc. (Zar) une section $e_0$ de $F$ telle que $(e_0, \rho e_0 \overset{\text{def}}{=} \rho(e_0))$ forment une base de $F$

b_1) (Si $M = \underline{\text{End}}_{\mathcal{O}_S}(F)$) Il existe loc. Zar un "sous-fibré en droites" $D_0 \subset F$, tel que $F = D_0 \oplus D_1$, où $D_1 = \rho(D_0)$

[Crossed out section]
b') (Si $M = \underline{\text{End}}_{\mathcal{O}_S}(F)$) Le $\mathcal{O}_S[T]$-module $F_\rho$ ($= F$ avec indéterminée, qui opère sur $F$ par $\rho$) est [unclear: monogène] loc (Zar).
[End crossed out section]

c) (Si $M = \underline{\text{End}}_{\mathcal{O}_S}(F)$) Le $\mathcal{O}_S[T]$-module $F_\rho \simeq \mathcal{O}_S[T]/(T^2 - wT + v)$

c') (Si $M = \underline{\text{End}}_{\mathcal{O}_S}(F)$) Il existe loc Zar une base $e_0, e_1$ de $F$, telle que pour la matrice de $\rho$ p.r. à cette base... [cut off]

\newpage

%% ===== Page 345 =====

soit
(3) $\rho_{e_0,e_1} = \begin{pmatrix} 0 & -v \\ 1 & w \end{pmatrix}$
(i.e. $\rho e_0 = e_1$, d'où $\rho e_1 = \rho^2 e_0 = -v e_0 + w e_1$)

[MARGIN: (Cas $M = \underline{End}(F)$)]
c'') $\exists$ loc Zar une base $(e_0, e_1)$ de $F$ telle que la matrice de $\rho$ p.r. à celle-ci soit
(4) $\rho_{e_0, e_1} = \begin{pmatrix} 0 & 1 \\ -v & w \end{pmatrix}$

d) Le commutant de $\rho$ dans $M$ est réduit à $A_\rho = \mathcal{O}_S + \mathcal{O}_S \rho$, et il en reste ainsi après changement de base $S' \to S$.

d') Idem dans chaque fibre.

Dim
$a \iff a'$ trivial (algèbre linéaire des Modules $M$ et des sous-modules $\mathcal{O}_S \oplus \mathcal{O}_S \rho$ loc. facteurs directs)

$b \iff b' \iff b''$ trivial, $b) \implies a)$ imm.

$a \implies b$ il suffit de choisir $e_0$ en chaque point de façon à ne pas être vecteur propre sous $\rho$ (possible sur un corps, pour un op. non scalaire) puis on relève aux voisinages.

$c) \iff c')$ trivial, les bases pour lesquelles $\rho_{e_0, e_1}$ soit de la forme $\begin{pmatrix} 0 & -v \\ 1 & w \end{pmatrix}$ étant justement les bases de la forme $(e_0, e_1 = \rho e_0)$.

$c') \iff c'')$ signifie que si $\rho$ satisfait aux conditions envisagées, $^t \rho$ aussi. C'est [unclear: vrai] et incidemment il en résulte que la condition pour $^t \rho$ l'implique pour $\rho = ^t(^t \rho)$, i.e. $c'') \implies c')$ - mais c'est immédiat sous la forme a).

Ainsi les conditions a) à c'') sont équivalentes, d'ailleurs $d) \implies d') \implies a)$ est trivial, reste à prouver que $b$ (p.ex) $\implies d)$. [unclear: Comme c'est stable] par chgt de base, il suffit de voir que

\newpage

%% ===== Page 346 =====

974

NB De ce fait, $G_{\rho} = \underline{M}^*$ opère sur $\underline{M}$ (par autom intérieurs) ou sur $G$, $S G$, via $P G = \underline{M}^*/\mathbb{G}_m$, et $P G$ est quotient ffpf de $S G$ (et $=$ pour la top. (ét), si $2$ inv.).
Donc on voit que la notion de conjugaison précédente, dans la formule
(6) $\rho' = int(\alpha) \rho = \alpha \rho \alpha^{-1}$,
où $\alpha \in \Gamma(S, \underline{M}^*)$.

Cor. 2 Soient $\rho, \xi$ deux sections de $G$, $\rho$ étant régulière. Pour qu'il existe loc. sur $S$ un $\alpha \in \Gamma G$ telles que
(7) $int(\alpha) \rho = \xi$ i.e. $\alpha \rho \alpha^{-1} = \xi$
il f. et s. qu'on ait
a) $\xi$ est régulière, i.e. [unclear: $\xi \ne$ scalaire dans chaque fibre]
b) $\det \xi = 1$, $Tr \xi = Tr \rho$
(On peut alors prendre $\alpha$ dans $S G$...)
En effet par les conditions du cor. 1 pour la conjugaison de $\xi$ à $\rho$, la condition
$\det \xi \rho = \det \rho$ s'écrit ici $\dots$ $\det \xi \det \rho = \det \rho$
i.e. $\det \xi = 1$.

Corollaire 3 L'indétermination de $\alpha$ dans la résolution de (7) est de la forme $\alpha \mapsto \alpha t$, où $t$ section de $A_{\rho}^*$ i.e. $t = c \rho + d$, avec $c, d \in \Gamma(\mathcal{O}_S)$, $\det(c \rho + d)$ inv. (NB $\iff c^2 v + c d u + d^2$ inv. où $v = \det \rho$, $w = Tr \rho$, cf (1) ).

\newpage

%% ===== Page 347 =====

si $u$ est une section de $M$ telle qu'elle commute
à $\rho$, alors $u = a + b\rho$ ($a, b \in \Gamma\mathcal{O}_S$). Pour $a) \Rightarrow b)$
est [unclear: immédiat], et $b)$ résulte que la question est
locale fppf. OPS alors $M = End(F)$, $F$ muni d'une
base $(e_0, e_1 = \rho(e_0))$, on a,

$u(e_0) = a e_0 + b e_1$, posant $u' = a + b\rho$,
[crossed out: $u(e_1) = u\rho(e_0) = \rho u(e_0) = \rho(a e_0 + b e_1) = a \rho(e_0) + b \rho(e_1)$]

et si $u$ commute à $\rho$, on a
$u(e_1) = u\rho(e_0) = \rho u(e_0) = \rho(a e_0 +$
$u'(e_1) = u'\rho(e_0) = \rho u'(e_0)$ car $u'$ commute à $\rho$
donc
$u(e_0) = u'(e_0)$, $u(e_1) = u'(e_1)$ donc $u=u'$
cqfd.

Définition - Une section de $M$ qui satisfait aux conditions
équivalentes de la prop. est dite régulière. Idem
pour sections de $G, Sl_2$, interprétées comme
sections de $M$ - [crossed out] telle section
est régulière ssi [crossed out] son polynôme caract.
n'est pas un carré, i.e. n'est dans le centre. Ceci est
stable par multiplication par un scalaire inver-
sible - donc par passage au quotient on définit
la notion de section régulière de $PG$ -
ce sont les sections qui ne sont pas $=1$ en
aucune fibre.

Corollaire 1. Soient $\rho, \rho'$ deux sections régulières de
$M$. Pour qu'elles soient conjuguées (loc Zar)
via $M^*$, il faut et il suffit a) que $\rho'$ soit régulière
[MARGIN: (Condition de régularité)] et b) qu'on ait
(5) $det \rho = det \rho'$, $Tr \rho = Tr \rho'$ (Conditions numériques)

\newpage

%% ===== Page 348 =====

Ceci résulte de l'équivalence $a) \Leftrightarrow d)$ de la prop.

Proposition — Soit
(9) $$\varepsilon = 1 + N$$
une section de $M$. Conditions équivalentes :

[MARGIN: NB $\det(1+N) = 1 + Tr N + \det N$
$Tr(1+N) = 2 + Tr N$]

a) $\det \varepsilon = \det 1 \quad (=1)$
(10) $Tr\ \varepsilon = Tr\ 1 \quad (=2)$
$\varepsilon$ régulière
$\Leftrightarrow \varepsilon$ régulière et $N$ régulière

b) $\det N = Tr\ N = 0$
(11)

c) (Si $M = \underline{End}_{\mathcal{O}_S}(F)$) $\exists$ sous-fibré en droites $D \subset F$ tel
que l'on ait
$$N(F) = D, \ N(D) = 0$$
(Alors $D$ est unique, déterminé par $\varepsilon$ i.e. par $N$,
$D = Im N = \ker N$)
c') [crossed out: et si $M = \underline{End}_{\mathcal{O}_S}(\mathcal{O}_S^2)$]. $\exists$ loc. une base $e_1, e_2$ tel que $\varepsilon = \begin{pmatrix} 1 & 1 \\ 0 & 1 \end{pmatrix}$ i.e. $N = \begin{pmatrix} 0 & 1 \\ 0 & 0 \end{pmatrix}$

d) $\exists$ Borel $B \subset G$, tel que $\varepsilon$ soit une section
partout $\neq 1$ de $B_u$.

d') $\exists$ Borel $SB \subset SG$, tel que $\varepsilon$ soit une
section partout $\neq 1$ de $(SB)_u$.

De plus, $D, B, SB$ sont bien déterminés par
$\varepsilon$ i.e. par $N$, et déterminent $N$ à homothétie près.

Définition : on dit que $\varepsilon$ est (strictement) unipotente
si on a $\det \varepsilon = 1, Tr\ \varepsilon = 2$, que $N$ est strictement
nilpotente $\Leftrightarrow \det N = Tr\ N = 0$ (d'où $N^2 = 0$ par
Hamilton-Cayley). On dit que $\varepsilon$ est unipotente
régulière resp. $N$ nilpotente régulière (en laissant
tomber le qualificatif "strictement") si de plus
il est régulier.

\newpage

%% ===== Page 349 =====

Proposition. Soit $P$ une section de $M$. Conditions équivalentes

a) $P$ est régulier, et $P^2=P$

b) $\det P=0$, $Tr P=1$

c) (Si $M=End_{\mathcal{O}_S}(F)$) $\exists$ base $e_0, e_1$ de $F$ telle que la matrice de $P$ est $P_{e_0, e_1} = \begin{pmatrix} 0 & 0 \\ 0 & 1 \end{pmatrix}$ i.e. $P$ est la projection de $\mathcal{O}_S e_0 \oplus \mathcal{O}_S e_1$ sur le facteur $\mathcal{O}_S e_1$.

d) $\exists_{loc. Zar.} \varphi : M \simeq M_2(\mathcal{O}_S)$, transformant $\begin{pmatrix} 0 & 0 \\ 0 & 1 \end{pmatrix}$ en $P$.

c') (Si $M=End_{\mathcal{O}_S}(F)$) $P$ est un projecteur dont le noyau a rg $1$.

Pour la démonstration, on notera que la condition b) implique déjà que $P$ est régulier, car si $P$ était un scalaire $\lambda$ (sur une fibre) on aurait
$$ \det P = \lambda^2 = 0, \quad Tr P = 2\lambda = 1 $$
donc $\lambda$ serait inversible et serait de carré nul, absurde.

Noter par cet argument un peu plus général, que si $P$ est une section de $M$ telle que
(12) $\delta(P) = (Tr P)^2 - 4 \det P = w^2 - 4v$ (= discriminant du polynôme de Hamilton-Cayley de $P$)
soit une section inversible de $\mathcal{O}_S$, alors $P$ est une section régulière de $M$.

Les $P$ satisfaisant les conditions de la proposition précédente sont appelés les projecteurs élémentaires de $M$.

Corollaire. On a une corr. biunivoque entre les projecteurs élémentaires de $M$, et les couples de Borel $(B,T)$ dans $G$ ($= \mathcal{S}B, \mathcal{S}T$ des $\mathcal{S}\ell_2$, $= \mathcal{P}B, \mathcal{P}T$ dans $PG$), en ...

\newpage

%% ===== Page 350 =====

associant à $P$ un tore $T = \underline{A}_P^*$, et le
Borel $B_0$ formé des $u \in G$ qui
stabilisent le drapeau $D_0$, et induisent
l'identité sur $D_0$ et sur $F/D_0 \simeq \underline{D}_1$, [unclear: pour $e_0, e_1$]
$$ (13) \quad T = \left\{ \begin{pmatrix} H_0 & 0 \\ 0 & H_1 \end{pmatrix} \mid H_0, H_1 \in \Gamma \underline{G}_m \right\} $$
$$ B_0 = \left\{ \begin{pmatrix} H_0 & \lambda \\ 0 & H_1 \end{pmatrix} \mid H_0, H_1 \in \Gamma \underline{G}_m, \lambda \in \Gamma \underline{G}_a \right\} $$
Le couple de Killing opposé à $(T, B_0)$ est
$(T, B_1)$, avec
$$ (14) \quad B_1 = \left\{ \begin{pmatrix} H_0 & 0 \\ \lambda & H_1 \end{pmatrix} \mid H_0, H_1 \in \Gamma \underline{G}_m, \lambda \in \Gamma \underline{G}_a \right\} $$
de façon un peu plus intrinsèque, $B_1$
est le stabilisateur de $\underline{D}_1$. Le proj.
él. associé au couple de Killing [unclear: en conservant les notations] opposé $(T, B_1)$
est $1 - P$.

Remarque. Le schéma des $\underline{P}$ est un fibré principal
(i.e. un $\underline{G}_m$-torseur)
sur le schéma des Borels (de $G$, de $SG$ ou
de $PG$, cela revient au $=$), il est
lisse de dimension relative 2. Le schéma des
projecteurs élémentaires $P$ est isom.
à celui des couples de Killing, fibré
lisse de dim. relative 1 sur le schéma
des Borels. Le produit fibré de
ces deux schémas relatifs au dessus de $\underline{B} =$ schéma
des Borels, est isomorphe au schéma
des épinglages de $G$ (ou de $SG$, ou de $PG$)
(i.e. des isomorph. $G_0 \simeq G$ (resp. $S\ell(2)_S \simeq SG$),

\newpage

%% ===== Page 351 =====

resp! $G(\mathbb{P}^1_S) \simeq P G$ ), qui est un torseur
sous $GP(1)_S = \underline{Aut}_{S-gr} \mathbb{G}_{a|S} \simeq \underline{Aut}_{S-gr} \hat{\mathbb{G}}_{m|S}$
$\simeq \underline{Aut}_{S-gr} (\hat{\mathbb{G}}_{a|S})$. On dit que $\varepsilon$ et
$P, N$ associés [MARGIN: (... $N = \varepsilon - 1$ et $P$ en associés)] sont dits [unclear] un Borel,
i.e. proviennent d' un épi-plongem.
On vérifie que ces séries ont relations
suivantes
(15) $$P N = 0 \quad , \quad N P = N$$
ou encore, en termes de $P$ et $\varepsilon = 1 + N$
(16) $$P \varepsilon = P \quad , \quad \varepsilon P = P + \varepsilon - 1,$$ [MARGIN: d'où $N$ nilp !]

Nous considérons maintenant la
question de la conjugaison de sections
$\tilde{\rho}, \tilde{\rho}'$ de $P G$, définies par des sections $\rho, \rho'$
de $G$. On suppose $\rho$ régulière i.e. $\neq 1$ en
ch. pt. géom. [unclear] $\rho$ régulière - La
conjugaison par une section $\alpha$ de
$G$ (div...) signifie qu'on a
(17) $$\rho' = \lambda \rho \lambda^{-1} \quad , \quad \lambda \in \Gamma G_m$$
d'où on déduit
$$\det \rho' = \lambda^2 \det \rho \quad , \quad Tr \rho' = \lambda Tr \rho$$
et par suite
(18) $$\frac{Tr^2 \rho'}{\det \rho'} = \frac{Tr^2 \rho}{\det \rho}$$
Posons, pour une section $\rho$ de $G$
(19) $$\Delta(\rho) = \frac{Tr^2 \rho}{\det \rho}$$

\newpage

%% ===== Page 352 =====

[unclear: On a] - fait ce qu'il fallait pour que [unclear: le morph.]
$$G \longrightarrow E'$$ passe au
quotient en
(20) $$PG \overset{\Delta}{\longrightarrow} E'$$
$$\Delta(\tilde{\rho}) = \Delta(\rho) = \frac{(Tr \rho)^2}{\det \rho}$$
et que cela remonte à un autom. int.

Proposition. Pour que deux sections $\tilde{\rho}, \tilde{\rho}'$ de $PG$ soient conjuguées loc. fppf [unclear: et $\tilde{\rho}$ est régulière], il faut et il suffit que $\tilde{\rho}'$ soit régulière, et qu'on ait
(21) $$\Delta(\tilde{\rho}) = \Delta(\tilde{\rho}')$$
(où au dessus de $\rho, \rho'$ relevant $\tilde{\rho}, \tilde{\rho}'$) que $\rho'$ soit régulière et satisfasse (21). Ces conditions sont satisfaites quand $\Delta(\tilde{\rho})$ inv. i.e. $Tr \rho$ inv.

Il reste à prouver la suffisance. Notons que loc. fppf, existe $\lambda \in \Gamma \mathbb{G}_m$ tel que
$$\det(\rho') = \det(\lambda \rho) \quad \text{i.e.} \quad \lambda^2 = \frac{\det \rho'}{\det \rho}$$
[unclear: Quitte à des ops. : on remplace] $\rho$ par $\lambda \rho$) pour avoir $\det \rho = \det \rho'$. Alors (18) s'écrit :
(22) $$Tr^2 \rho' = Tr^2 \rho$$
d'où la question revient à voir que $\exists \varepsilon \in \Gamma \mu_2$, telle que
$$Tr \rho' = \varepsilon Tr \rho \quad ( = Tr(\varepsilon \rho) ),$$
ainsi
$$\det \rho' = \det \varepsilon \rho \quad ( = \varepsilon^2 \det \rho = \det \rho ),$$
d'où la conjugaison de $\rho'$ et $\varepsilon \rho$, qui termine. Mais si $\tau, \tau'$ sont deux sections de $E_s$, telles que $\tau'^2 = \tau^2$, on aura bien
$\tau' = \varepsilon \tau$ ($\varepsilon = \tau' / \tau$), et $\varepsilon^2 = 1$, OK.

\newpage

%% ===== Page 353 =====

Rmq Si $\Delta(\rho)$ non inv. il y a des
contre-exemples évidents, p.ex. avec $Tr \rho'=0$,
$(Tr \rho')^2=0$, $det \rho = det \rho'=1$, $Tr \rho$ et $Tr \rho'$ n'étant
pas homothétiques. On peut prendre
$$ \rho = \begin{pmatrix} 0 & 1 \\ -1 & \omega \end{pmatrix} , \rho' = \begin{pmatrix} 0 & 1 \\ -1 & \omega' \end{pmatrix} $$
$\omega^2 = \omega'^2 = 0$, $\omega, \omega'$ non homothétiques, p.ex.
sur un [unclear] de base $A = A_0[\omega, \omega']/(\omega^2, \omega \omega', \omega'^2)$
([unclear: l'idéal d'augm. est le module libre de rg 2]
engendré par $\omega, \omega'$).

2) Position relative de deux sections $\rho, \sigma$
de $G$, ou $\tilde{\rho}, \tilde{\sigma}$ deux sections de $PG$.

On s'intéresse dans $PG$ à
(23) $\tilde{\varepsilon}_0 = \tilde{\sigma}^{-1} \tilde{\rho}$ et $\tilde{\varepsilon}_1 = \tilde{\rho} \tilde{\sigma}^{-1} = \tilde{\rho}(\tilde{\varepsilon}_0) = \tilde{\sigma}(\tilde{\varepsilon}_0)$
et à la condition que $\tilde{\varepsilon}_0$ ([unclear: ou] $\tilde{\varepsilon}_1$)
soit unipotente régulière - ce qui proviendra d'un
couple canonique de $\varepsilon_0, \varepsilon_1$ unipotentes régulières de $SL_2$,
satisfaisant
(24) $\varepsilon_1 = \tilde{\rho}(\varepsilon_0) = \tilde{\sigma}(\varepsilon_0)$
[MARGIN: d'après [unclear] condition sur $\tilde{\rho}, \tilde{\sigma}$]
(25) $\begin{cases} \varepsilon_0 = (\text{Scalaire inv.}) \sigma^{-1} \rho \text{ soit } \varepsilon_0 = z_0 \sigma^{-1} \rho \\ \varepsilon_1 = ( \text{ - - - - } ) \rho \sigma^{-1} \quad \varepsilon_1 = z_1 \rho \sigma^{-1} \end{cases}$
et (20) $\tilde{\varepsilon}_0, \tilde{\varepsilon}_1$ soient [unclear: purement] "équivalentes",
i.e. posons
(26) $\varepsilon_0 = 1 + N_0, \quad \varepsilon_1 = 1 + N_1$

\newpage

%% ===== Page 354 =====

que $N_0$ et $N_1$ soient partout non-homothétiques.

Alors $\tilde{\varepsilon}_0(G_a)$, $\tilde{\varepsilon}_1(G_a)$, $\tilde{\varepsilon}_0 \in \Gamma \tilde{\varepsilon}_0(G_a)$ [unclear] définissent
un ~~épi-gloge~~ épinglage canonique de $PG$ dans $G$ ...
On aura alors, pour cet épinglage

(26) $\rho = \begin{pmatrix} 0 & u \\ v & w \end{pmatrix} \qquad \sigma = \begin{pmatrix} 0 & u' \\ v' & w' \end{pmatrix}$

avec $u, v, u', v' \in \Gamma(G_m)$, $w, w' \in \Gamma G_a$

et on aura

$$ (27) \begin{cases} - \sigma^{-1} \rho = \begin{pmatrix} - \frac{v}{v'} & \frac{u w' - w u'}{u' v'} \\ 0 & - \frac{u}{u'} \end{pmatrix} = \varepsilon_0 = \begin{pmatrix} 1 & \ast \\ 0 & 1 \end{pmatrix} \\ - \rho \sigma^{-1} = \begin{pmatrix} - \frac{u}{u'} & 0 \\ \frac{v' w - v w'}{u v'} & - \frac{v}{v'} \end{pmatrix} = \varepsilon_1 = \begin{pmatrix} 1 & 0 \\ \ast & 1 \end{pmatrix} \end{cases} $$
[MARGIN: bien définies au niveau épinglage]

donc il faudra

(28) $u/u' = -1, v/v' = -1$, i.e. $u' = -u, v' = -v$

~~(29) ... donc~~
$$ - \rho \sigma^{-1} = \begin{pmatrix} 1 & 0 \\ \frac{1}{u}(w+w') & 1 \end{pmatrix} $$
et il faut

(29) $w' + w = u$ i.e. $w' = u - w$

en d'autres termes (vu (26)(28))

(30) $\rho + \sigma = u P \qquad$ où $\quad P = \begin{pmatrix} 0 & 0 \\ 0 & 1 \end{pmatrix}$

est le projecteur évident associé à l'épinglage envisagé.
ce qui détermine $\sigma$ en termes de $\rho$ et du coeff. $u$ de $\rho$, et on a

(31) $- \sigma^{-1} \rho = \varepsilon_0 = \begin{pmatrix} 1 & u/v \\ 0 & 1 \end{pmatrix}, \quad N_0 = \begin{pmatrix} 0 & u/v \\ 0 & 0 \end{pmatrix}$

\newpage

%% ===== Page 355 =====

On aimerait donner des "critères numériques" intrinsèques pour que $\rho$, $\sigma$ aient les propriétés précédentes ... [unclear]

Proposition Soient $\rho$, $\sigma$ deux sections [unclear: de] $G_g$
[unclear: Posons]
$$ \varepsilon_0 = -\sigma^{-1}\rho \quad , \quad \varepsilon_1 = -\rho\sigma^{-1} $$
de sorte qu'on ait ...
$$ \varepsilon_1 = \sigma(\varepsilon_0) = \rho(\varepsilon_0) .$$

Conditions équivalentes

a) 1) $\varepsilon_0$ est unipotente régulière (partie unipotente régulière) et 2) $\varepsilon_0$ et $\varepsilon_1$ [unclear] posant
$$ \varepsilon_0 = 1 + N_0 \quad , \quad \varepsilon_1 = 1 + N_1 $$
$N_0$ ($\Leftrightarrow N_1$) est nilpotent régulier, et $N_0$ et $N_1$ sont partout sans homothétiques.

[MARGIN: cf prop. p. 497]
b) On a
$$ \alpha) \text{det } \sigma = \text{det } \rho $$
$$ \beta) \text{Tr } \sigma^{-1}\rho = -2 $$ $\left\} \begin{array}{l} \text{car d'une relation canonique} \\ \text{det}(\varepsilon_0)=1, \text{Tr}\varepsilon_0=2 \text{ i.e. } \varepsilon_0 \\ \text{est (strictement) unipotente} \end{array} \right.$
$$ \gamma) \quad u \stackrel{\text{df}}{=} \text{Tr}(\sigma + \rho) \text{ est inversible} \left\{ \begin{array}{l} \text{ceci exprime} \\ \text{que } \varepsilon_0 \text{ est régulière} \\ \text{et indépendante} \\ \text{de } \varepsilon_1, \text{ cf} \\ \text{loc. cit.} \end{array} \right. $$

c) On a
(33) $\alpha) \text{ det}(\sigma+\rho)=0$
$\beta) \text{ Tr}(\sigma+\rho)$ inversible

[MARGIN: NB $P$ et $u$ uniquement déterminés, $u = \text{Tr}(\sigma+\rho)$, $P = \frac{1}{u}(\sigma+\rho)$]
d) On a $\text{det } \sigma = \text{det } \rho$, et il existe une projection $P$ et un scalaire inv. $u \in \Gamma \mathbb{G}_m$, tels que l'on ait
(34) $\sigma + \rho = u P$

\newpage

%% ===== Page 356 =====

[MARGIN: NB cet isomorph. est univoque]

[crossed out: e) (Si $M= \underline{End}(F)$) $\exists$ une base de $F$, telle que par rapport à cette base on ait]

e) $\exists$ isomorphisme $\operatorname{Gl}(2) \xrightarrow{\varphi} G$, tel que l'on ait
(35)
$$ \varphi^{-1}(\rho) = \begin{pmatrix} 1 & 0 \\ 1 & 1 \end{pmatrix} $$
$$ \varphi^{-1}(\varepsilon_0) = \begin{pmatrix} 1 & [unclear] \\ 0 & 1 \end{pmatrix} $$

[MARGIN: et $\varphi$ est unique; avec 
$u = Tr(\sigma + \rho)$ ($= -u'$)
$w = Tr \rho$ ($w' = Tr \sigma$)
$-uw = det \rho = det \sigma$
ce qui dicte $u,w,w'$
en termes de $\rho, \sigma$,
ainsi $u',v'$...]

f) $\exists$ isomorphisme $\operatorname{Gl}(2) \simeq G$ tel que l'on ait
(36)
$$ \varphi^{-1}(\rho) = \begin{pmatrix} 0 & 0 \\ 1 & w \end{pmatrix} \quad \varphi^{-1}(\sigma) = \begin{pmatrix} v & w' \\ 0 & u' \end{pmatrix} $$
avec $uu' = vw' = 0$, $w+w' = u$

[MARGIN: NB On aura $P = \begin{pmatrix} 0 & 0 \\ 0 & 1 \end{pmatrix}$]

g) On aura dans (35), avec
(38) $$ \varphi^{-1}(\varepsilon_0) = \begin{pmatrix} 1 & u/v \\ 0 & 1 \end{pmatrix} $$
et on aura
(39) $$ \sigma + \rho = u P , \quad \text{où } \varphi^{-1}(P) = \begin{pmatrix} 0 & 0 \\ 0 & 1 \end{pmatrix} . $$

Démonstration facile. L'équivalence a) $\Leftrightarrow$ b) est dans prop. p. 449. Pour b) $\Leftrightarrow$ c), on écrit

$$ det(\sigma + \rho) = det(\sigma(1 + \sigma^{-1}\rho)) = det \sigma \cdot det(1 + \sigma^{-1}\rho) $$
$$ = det \sigma + det \rho + (det \sigma) Tr(\sigma^{-1}\rho) $$
[MARGIN: $1 + Tr \sigma^{-1}\rho + det \sigma^{-1}\rho$]
[MARGIN: $(det \sigma^{-1}det\rho)$]

donc au vu de l'hypothèse
$det \sigma = det \rho$
commune à b) et c), on a

(40) $$ det(\sigma + \rho) = det \sigma (2 + Tr(\sigma^{-1}\rho)) $$

donc la première n'est [unclear: autre] que $Tr \sigma^{-1}\rho = -2$, cqfd.

\newpage

%% ===== Page 357 =====

Posant $Q = \sigma \varepsilon \rho$, l'équivalence $b \iff c$ prouve alors l'équivalence (valable pour tte section $Q$ de $M$)

$$(4a) \quad ( \det Q = 0, \text{Tr } Q \text{ inv.} ) \Longleftrightarrow ( \exists \text{ projecteur idém. } P, \text{ et } u \in \Gamma(G_m) \text{ tels que } Q = u P )$$

qui résulte facilement des résultats du n$^{\circ}$ précédent c à d du [unclear: corollaire de le prop. 115], que j'ai oublié d'expliciter.

L'équivalence avec e), f) est aussi immédiate, et est un fait qui résume certains calculs précédents.

[MARGIN: Corollaire Sous ces conditions $1, N_0, N_1, P$ forment une base de $\mathcal{M}$, et [unclear: $\rho, \sigma$ sont de la forme $\rho = N_0 + N_1 + wP$, $\sigma = -u(N_0+N_1)+(u-w)P$ qui forment une base... ]]

Notre point de vue ici sera que nous partons de $\bar{\rho}, \bar{\sigma}$, que nous voulons relever - l'épinglage est en effet défini déjà en termes de $\bar{\rho}, \bar{\sigma}$, mais il reste la question de normaliser le relèvement de $\bar{\rho}, \bar{\sigma}$ en $\rho, \sigma$. Ceci a été déjà fait une première restriction par la formule

$$\varepsilon_1 = -\sigma^{-1} \rho$$

qui élimine le scalaire $\lambda_0$ - on peut normaliser plus en multipliant $\rho, \sigma$ par un facteur $u^{-1}$, de façon à obtenir la relation simple.

$$(4b) \quad Tr(\rho + \sigma) = 1 \quad \text{i.e. } Tr \rho + Tr \sigma = 1 \quad \text{i.e. } u = 1,$$

qui signifie donc que $\rho, \sigma$ sont de la

\newpage

%% ===== Page 358 =====

480

formes
(41) $\rho = \begin{pmatrix} 0 & 1 \\ v & w \end{pmatrix} , \quad \sigma = \begin{pmatrix} 0 & -1 \\ v' & w' \end{pmatrix}$

$v+v'=0, w+w'=1 \quad \text{i.e.} \quad v'=-v, w'=1-w$ ,

donc
(42) $\varepsilon_0 = \begin{pmatrix} 1 & 1/v \\ 0 & 1 \end{pmatrix}$ ,

où on a identifié, par $\varphi$, $\operatorname{Gl}(2)_{\mathbb{Z}}$ et $G$.
Y a-t-il unicité de $\rho, \sigma$ au dessus de $\bar{\rho}, \bar{\sigma}$,
par les conditions $\varepsilon_1 = -\rho \sigma$ (i.e. $\varepsilon_0 = -\sigma \rho^{-1}$)
et (40)? Je voudrais d'abord vérifier avec
soin une assertion faite par la bande à
certains endroits de ces notes, savoir que l'application
$G \to PG \simeq G/\{\pm 1\}$
induit une bijection

(43) unipotents réguliers de $G \to$ unipotents réguliers de $PG$

(au moins localement [unclear])

[MARGIN: Donc, si l'on se donne une section de $PG$ unipotente régulière elle se relève loc. [unclear] d'une section régulière de $G$.]

la surjectivité est en effet une consé-
quence des définitions des unipotents réguliers
dans $PG$. Donc il s'agit de tester l'injecti-
vité, ce qui signifie aussi ceci : si
$\varepsilon, \varepsilon'$ sont deux sections unipot. régulières de
$M$, telles que $\bar{\varepsilon}' = \bar{\varepsilon}$ i.e. $\varepsilon' = \lambda \varepsilon \quad (\lambda \in \Gamma(\dots))$, alors
$\varepsilon = \varepsilon'$, i.e. $\lambda=1$. Prenant les déterminants des
deux membres, on trouve

(44) $\lambda^2=1$ ,

\newpage

%% ===== Page 359 =====

prenant les traces on trouve

(45) $2\lambda = 2$

ce sont les conditions néc. et suff. sur $\lambda$ pour que $\forall$ ($\varepsilon$ étant unipot. régulier) $\lambda \varepsilon$ soit encore unipot. (donc néc. unipot. régulier). Si $2$ est inv. sur $S$, plus gén. si $2$ est régulier dans $\mathcal{O}_S$, la condition implique $\lambda=1$, OK, mais en l'absence de cette hypothèse (p.ex. en caractéristique $2$ sur un schéma non réduit) il n'en est plus néc. ainsi. Ecrivons

$\lambda = 1 + \nu$

les conditions (44) (45) s'écrivent

(45') $2\nu = 0$, $\nu^2 = 0$

donc pour que sur $S$ l'application (43) soit injective, il f et s. que la seule section de $S$ qui satisfasse (45') soit la section nulle.

En général, désignons par $\mu'_2$ le sous-schéma en groupes de $\mu_2$ défini par (45') (c'est en fait un schéma en groupes fini sur $\mathbb{Z}$, plat sur $\mathbb{Z}$ [crossed out], groupe infinitésimal absolu, son anneau en spectre

$\mathbb{Z}[T]/(T^2-1, 2T-2) \simeq \mathbb{Z}[S]/(S^2, 2S)$,

isomorphe au schéma des nbs duaux sur $\mathbb{Z}$, relatif au module $\mathbb{Z}/2\mathbb{Z}$ sur $\mathbb{Z}$,

donc $\mu'_2 \simeq \operatorname{Spec}( \mathbb{Z} \oplus \mathbb{Z}/2\mathbb{Z} )$ ...
(idéal de carré nul)

358

\newpage

%% ===== Page 360 =====

Il se trouve que $\mu'_{2, S}$ opère [MARGIN: librement par translations] (sur le schéma des unipotents réguliers de $G$, et on a

(47) $$ Unip.rég(G)/\mu'_{2, S} \xrightarrow{\sim} Unip.rég(PG) $$ .

Ceci semble à première vue en contradiction avec la bijection des Borels

(48)
\begin{tikzcd}
B \ar[d] & Borels(G) \ar[r, "\sim", "B \mapsto B/G_{...}"] \ar[d, "\wr"] & Borels(PG) \ar[d, "\wr"] & \tilde{B} \ar[d] \\
B_u & \text{"Sous-groupes unipots. à 1 param." de } G \ar[r, "\sim"] & \text{"S.-groupes unipots. à 1 par." de } PG & \tilde{B}_u \\
L \ar[rr, "|->"] & & p(L) & 
\end{tikzcd}

(où $p : G \to PG$ est le morphisme can. !)

La difficulté proviendrait-elle du fait que les flèches verticales (surjectives a priori par définition des s-groupes unipot. à 1 paramètre de $G$ resp de $PG$) ne seraient pas injectives, en dehors de l'hyp. 2 inv. (i.e. car. rés. $\neq 2$) ?

J'ai envie de dire que

(49) $$ B = \operatorname{Norm}(B_u) \text{ , } \tilde{B} = \operatorname{Norm}(\tilde{B}_u) $$

sont des formules qui permettent de retrouver $B, \tilde{B}$ à l'aide de $B_u, \tilde{B}_u$ - il suffirait bien sûr de vérifier la première de ces formules pour un groupe d'alg. lin. $M$ (i.e. le genre :

les sous-groupes unipot. à 1 paramètre de $G$ correspondent 1-1 aux sections de $\mathbb{P}$ sur $S$, i.e. aux s-fibrés en droites $D$ de $F$ ...) -

\newpage

%% ===== Page 361 =====

mais ça marche pareil dans $PG$, c'est à
dire $u \in \Gamma(M)^*$ est telle que $u N u^{-1} = \mu N + \lambda,$
[unclear: pour $\mu=1$ trivial.]

La difficulté doit plutôt provenir du
fait qu'une section unipotente régulière
$\varepsilon$ de $PG$ ne définit pas (sa restriction
sur $S$) un sous-groupe unipotent à $1$
paramètre - alors qu'une section $\varepsilon$ de
unipotents réguliers de $G$
$$\varepsilon = 1 + N$$
en définit par la formule
$$\varepsilon(t) = 1 + t N$$

On notera que si on remplace $\varepsilon$ par $\varepsilon' = \lambda \varepsilon$
(avec $\lambda$ satisfaisant (44)(45) $\lambda^2=1, 2\lambda=2$)

on a
$$\varepsilon' = (1+\nu)(1+N) = 1 + N' \quad N' = (1+\nu)N + \nu$$
et
$$\varepsilon'(t) = 1 + t N'$$
je dis que si $\nu \neq 0$, alors les
hom. composés
$$\tilde{\varepsilon}, \tilde{\varepsilon}' : \mathbb{G}_a \to G \to PG$$
[MARGIN: | |] associés à $\varepsilon, \varepsilon'$ seront distincts,
bien qu'ils aient même valeur pour $t=1$.
En effet, on aura les dérivées à l'origine
de $\tilde{\varepsilon}, \tilde{\varepsilon}'$ s'identifient aux composés
$$\mathcal{O}_S \to \mathcal{M}_S \to \mathcal{M}_S/\mathcal{O}_S \simeq P\mathcal{G}$$

\newpage

%% ===== Page 362 =====

482

qui transforment 1 en $\tilde{N}$ resp. en $\tilde{N}'$,
( i.e. la différence dans le p-groupe quotient
additif ) on a
$$ \tilde{N}' = (1+\nu) \tilde{N} $$
et $\tilde{N}$ est une section de $P\mathcal{G}$ partout non
nulle donc on ne peut avoir $\tilde{N}' = \tilde{N}$
i.e. $(1+\nu)\tilde{N} = \tilde{N}$ i.e. $\nu\tilde{N} = 0$
que si $\nu = 0$, i.e. $\lambda = 1+\nu = 1$, cqfd.

En conclusion, il faut considérer que
l'on se donne dans $P\mathcal{G}$, non seulement
$\tilde{\rho}$ et $\tilde{\sigma}$ ( d'où $\tilde{\varepsilon}_0 = \tilde{\sigma}^{-1}\tilde{\rho}$, $\tilde{\varepsilon}_1 = \tilde{\rho}\tilde{\sigma}^{-1}$ ) , mais
aussi des "sous-groupes à 1 paramètre"
(paramétrés) $\tilde{\varepsilon}_0, \tilde{\varepsilon}_1$ de $P\mathcal{G}$,

$$ (49) \left\{ \begin{array}{l} \tilde{\varepsilon}_0 : \mathbb{G}_a \to P\mathcal{G} \\ \tilde{\varepsilon}_1 : \mathbb{G}_a \to P\mathcal{G} \end{array} \right. $$

reliés bien sûr par

$$ (50) \left\{ \begin{array}{l} \tilde{\varepsilon}_1 = \tilde{\rho}(\tilde{\varepsilon}_0) = \tilde{\sigma}(\tilde{\varepsilon}_0) \quad i.e. \\ \tilde{\varepsilon}_1(t) = \tilde{\rho}(\tilde{\varepsilon}_0(t)) = \tilde{\sigma}(\tilde{\varepsilon}_0(t)) \end{array} \right. $$

ce qui signifie justement qu'on se donne,
en plus de $\tilde{\rho}, \tilde{\sigma}$, [unclear: des] (relèvements unipotents
$\varepsilon_0$ de $\tilde{\varepsilon}_0$ ( dont le choix n'est donc pas
a priori unique ) , et on en déduit $\varepsilon_1$
par
$$ \varepsilon_1 = \tilde{\rho}(\varepsilon_0) = \tilde{\sigma}(\varepsilon_0) $$

[MARGIN: on aura reconnu une relation $\sigma^{-1}\rho = \tilde{\varepsilon}_0$ donc $int(\rho)(\tilde{\varepsilon}_0) = \dots$]

361

\newpage

%% ===== Page 363 =====

[crossed out: Eett] Ainsi, on a ajouté la [crossed out: pb] description des
données, pour [unclear] pouvoir donner des
choix de normalisation précis - savoir

Proposition. Soient $\tilde{\rho}, \tilde{\sigma}$ des sections additives de $PG$,
et $\tilde{\varepsilon}_0$ un sous-groupe à 1-paramètre "strict"
[MARGIN: (i.e. mono et dont l'image est un $B_u$)]
$\tilde{\varepsilon} : \mathbb{G}_a \to PG$
tels que l'on ait
(51) $\tilde{\sigma}^{-1} \tilde{\rho} = \tilde{\varepsilon}_0(1) \stackrel{déf}{=} \tilde{\varepsilon}_0$
d'où
$\tilde{\sigma}(\tilde{\varepsilon}_0) = \tilde{\rho}(\tilde{\varepsilon}_0) \stackrel{déf}{=} \tilde{\varepsilon}_1$, et
(52) $\tilde{\rho} \tilde{\sigma}^{-1} = \tilde{\varepsilon}_1(1) \stackrel{déf}{=} \tilde{\varepsilon}_1 = \rho(\tilde{\varepsilon}_0) = \sigma(\tilde{\varepsilon}_0)$

Alors :
1) [crossed out: Il] (unicité) $\exists$ uniques $\varepsilon_0, \varepsilon_1$ sections
[crossed out: unipotentes de] $G$ telles que (posant $\varepsilon_i(1) = 1+N_i$,
$\varepsilon_i(t) = 1+t N_i$) on ait
(53) $\tilde{\varepsilon}_i(t) = \varepsilon_i(t)^{\sim}$.

2) [crossed out: On aura alors]
(54) $\varepsilon_1 = \tilde{\rho}(\varepsilon_0) = \tilde{\sigma}(\varepsilon_0)$
[crossed out: $\varepsilon_0 = -\lambda_0 \sigma^{-1} \rho$, $\varepsilon_1 = -\lambda_0 \rho \sigma^{-1}$]
Pour que $\varepsilon_0$ et $\varepsilon_1$ soient [unclear] indépendantes
[unclear], i.e. $N$

3) Soient $\rho, \sigma$ des relèvements de $\tilde{\rho}, \tilde{\sigma}$, on aura
alors
(55) $\varepsilon_0 = -\lambda_0 \sigma^{-1} \rho$, $\varepsilon_1 = -\lambda_0 \rho \sigma^{-1}$
pour $\lambda_0 \in \Gamma$ bien uniquement déterminé.
Pour un choix convenable de $\rho$ (quand $\sigma$ est déjà
choisi) on a [unclear] $\lambda_0 = 1$ (ce qui détermine $\rho$ ...)

\newpage

%% ===== Page 364 =====

483

4) de sorte qu'on a
(56) $\varepsilon_0 = - \sigma^{-1} \rho$ (d'où $\varepsilon_1 = - \rho \sigma^{-1}$)

4) $\rho$ et $\sigma$ étant choisis de façon à satisfaire à (en chaque fibre)
(56), pour que $\varepsilon_0, \varepsilon_1$ soient pointées et indépendantes
i.e. pour que l'on ait $N_0$ et $N_1$ non homothé-
tiques sur les fibres, il f. et i.s. qu'on ait
(57) $Tr(\rho + \sigma)$ inv.

Cf. au surplus p. 478' pour d'autres conditions
équivalentes ). (en chaque fibre).

5) Pour que $\varepsilon_0, \varepsilon_1$ soient indépendantes (loc. f. p. p.), il f. et
i.s. que l'on puisse trouver $\rho, \sigma$ ci-dessus
relevant $\tilde{\rho}, \tilde{\sigma}$ de façon à satisfaire à (56) et
à
(5?) $Tr(\rho \circ \sigma) = 1$ i.e. $Tr \rho \cdot Tr \sigma = 1$
[MARGIN: par un argument préliminaire] et cela détermine $\rho, \sigma$ de façon unique
(en termes de $\tilde{\rho}, \tilde{\sigma}, \tilde{\varepsilon_0}$).

Remarques 1) Les conditions sur $(\rho, \sigma)$ envisagées dans
le par. p. 478', avec la normalisation $n=1$ i.e.
$Tr(\sigma \rho) = 1$ i.e. $\sigma + \rho$ proj. elim., sont symétriques en
$(\rho, \sigma)$ : elles sont satisfaites pour $(\rho, \sigma)$ ssi elles le sont
pour $(\sigma, \rho)$ — l'échange de $\rho$ et $\sigma$ a pour effet de changer
$\varepsilon_0$ et $\varepsilon_1$ en leurs inverses, posons
(58) $\rho' = \sigma, \sigma' = \rho$ , $\varepsilon'_0 = - \sigma'^{-1} \rho' = - \rho^{-1} \sigma$, $\varepsilon'_1 = - \rho' \sigma'^{-1} = - \sigma \rho^{-1}$
donc
(59) $\varepsilon'_0 = \varepsilon_0^{-1}$, $\varepsilon'_1 = \varepsilon_1^{-1}$ de la p. 478'
De même, $(\rho^{-1}, \sigma^{-1})$ satisfait aux mêmes conditions, et permet

\newpage

%% ===== Page 365 =====

(60) $\rho'' = \rho^{-1}, \sigma'' = \sigma^{-1}, \varepsilon_0'' = \underbrace{-\sigma''^{-1}\rho''}_{-\sigma\rho^{-1}}, \varepsilon_1'' = \underbrace{-\rho''\sigma''^{-1}}_{-\rho^{-1}\sigma}$

(61) $\varepsilon_0'' = \varepsilon_1^{-1} \quad , \quad \varepsilon_1'' = \varepsilon_0^{-1} \quad ,$

mais il faudra faire attention qu'on n'aura plus
en général $Tr(\rho''+\sigma'')=1$, mais (par un calcul
matriciel immédiat)

(62) $Tr(\rho''+\sigma'') = Tr(\rho^{-1}+\sigma^{-1}) = \underbrace{\frac{1}{det \rho}}_{det \rho^{-1}} = (det \rho)^{-1} = (det \sigma)^{-1}$
(en supposant bien sûr $u_1=1$ i.e. $Tr(\rho\sigma)=1$)

On a ainsi égal. à 1 que si

(63) $det \rho \; (= det \sigma) = 1$

- c'est la condition de la relation 3a) de (A') (p. 463'),
qui signifie donc en outre que la normalisation
$u_1=1$ est stable par la transformation $\rho \mapsto \rho^{-1}, \sigma \mapsto \sigma^{-1}$.

Par contre, le choix de l'"épi-graphe" associé à
$(\rho,\sigma)$ (tiré $\tilde{\rho},\tilde{\sigma}$) [dans lequel l'accent a été mis sur
$\varepsilon_1 = -\rho\sigma^{-1}$, qui reçoit la forme simple
$\varepsilon_1 = \left( \begin{smallmatrix} 1 & 1 \\ 0 & 1 \end{smallmatrix} \right) \quad ,$
alors que $\varepsilon_0 \underbrace{= -\sigma^{-1}\rho}_{\text{[unclear]}}$ la forme plus compliquée
$\varepsilon_0 = \left( \begin{smallmatrix} 1 & u_1/u \\ 0 & 1 \end{smallmatrix} \right) \quad \left(= \left( \begin{smallmatrix} 1 & -1 \\ 0 & 1 \end{smallmatrix} \right) \text{ dans le} \atop \text{cas modulaire} \right)$]
fait intervenir implicitement l'ordre dans lequel
on a donné $\rho,\sigma$ - rien ne change, [unclear: à d... près d'échanger]
l'épi-graphe pour la [unclear] et $\varepsilon_1 \mapsto \varepsilon_1^{-1}$.
En fait, la situation offre le choix entre
quatre "épi-graphes" "invariants", associés à la paire de

\newpage

%% ===== Page 366 =====

demi-groupes ; aux paramètres $\{L_0, L_1\}$, correspondent
quatre sections $\varepsilon_0, \varepsilon_0^{-1}, \varepsilon_1, \varepsilon_1^{-1}$ de l'un [unclear: et de] l'autre.

On a préféré normaliser à l'aide de $L_1$, à cause
de la forme la plus commode de l'équation
des lacets (à laquelle j'ai fini par parvenir d'
ailleurs après des pérégrinations assez lourdes, où
c'est $L_0$ que je m'obstinais à privilégier dans
les formules - cf § 44, la notation définitive
n'apparait qu'au § 45 - ). La préférence donnée
à $-p\sigma^{-1}$ (noté $\varepsilon_1$) sur son inverse $-\sigma p^{-1}$ (noté $\varepsilon_1^{-1}$)
provient de la préférence (justifiée par les faits)
donnée à $p$ sur $\sigma$, dans la formule des
lacets écrite sous la forme
$$[\alpha, p] = [\beta, \sigma] \varepsilon_1(p)$$
au lieu de l'autre
$$[\beta, \sigma] = [\alpha, p] \varepsilon'_1(p) \qquad (\varepsilon'_1 = \varepsilon_1^{-1})$$
Donc c'est $\varepsilon_1$ qui méritait alors d'avoir la forme
matricielle la plus simple possible - on a
pris $\varepsilon_1 = \begin{pmatrix} 1 & 0 \\ 1 & 1 \end{pmatrix}$ ...

2) En termes du système
(64)
$$ \tilde{l}_\infty, \tilde{l}_1, \tilde{l}_0 \quad \text{satisfaisant} \quad \tilde{l}_\infty . \tilde{l}_1 . \tilde{l}_0 = 1 $$
$$ || \quad || \quad || $$
$$ \tilde{p} \quad \tilde{\sigma} \quad \tilde{\varepsilon}_0 $$
(i.e. $\tilde{p} . \tilde{\sigma} . \tilde{\varepsilon}_0 = 1$ i.e. $\tilde{\varepsilon}_0 = \tilde{\sigma}^{-1} \tilde{p}^{-1}$ ), le groupe
$\tilde{\mathfrak{S}}_3$, opère (cf § 42 I) par permutations circulaires
et par [unclear] $(\tilde{\varepsilon}_0^{-1}, \tilde{\sigma}, [\tilde{p}^{-1}]^{-1} = \tilde{p})$, qui ramène à son
tour par permutation circulaire donne $(\tilde{\sigma}^{-1}, \tilde{p}, \tilde{\varepsilon}_0^{-1})$ -

\newpage

%% ===== Page 367 =====

le passage de $(\tilde{\rho}^{-1}, \tilde{\sigma}, \tilde{\varepsilon}_0)$ à $(\tilde{\rho}^{-1}, \tilde{\rho}, \tilde{\varepsilon}_0^{-1})$ corr-
pond alors à l'échange $(\rho, \sigma) \mapsto (\sigma, \rho)$. Comment
[unclear: s'applique - crossed out] interpréter la transformation
$(\rho, \sigma) \mapsto (\rho^{-1}, \sigma^{-1})$, i.e. $(\tilde{\rho}^{-1}, \tilde{\sigma}, \tilde{\varepsilon}_0) \mapsto (\tilde{\rho}, \tilde{\sigma}^{-1}, \tilde{\varepsilon}_1)$ ?
Interprétons ces triples comme des hom.
$$ \Pi_{0,3} \xrightarrow{\varphi} T\Gamma G $$
$$ l_0 \mapsto l'_0 $$
$$ l_1 \mapsto l'_1 $$
$$ l_\infty \mapsto l'_\infty $$
on voit que l'échange $(\rho, \sigma) \mapsto (\sigma, \rho)$ correspond
à la composition de $\varphi$ avec l'autom de $\Pi_{0,3}$ tel que
$l_0 \mapsto l_1^{-1}$, $l_1 \mapsto l_0^{-1}$ donc $l_\infty \mapsto l_\infty^{-1}$
i.e. avec $\tau'_\infty$. De même la transformation
$(\rho, \sigma) \mapsto (\rho^{-1}, \sigma^{-1})$ correspond à la composition avec
l'autom
$$ l_0 \mapsto l_0^{-1} \quad \text{d'où} \quad l_\infty \mapsto l_0 l_1 $$
$$ l_1 \mapsto l_1^{-1} $$
qui n'est autre que $\tau_\Delta = \sigma_{01} \tau'_\infty$. Donc le groupe
de quatre transformations qui opère sur l'ens.
des couples $(\rho, \sigma)$ qu'on a envisagé ci-dessus
10) d'autre part provient du sous-groupe $\simeq \mathbb{Z}/2\mathbb{Z} \times \mathbb{Z}/2\mathbb{Z}$
engendré par $\tau'_\infty, \tau_\Delta$ dans $\operatorname{Aut}_{ext}(\Pi_{0,3}) = \mathfrak{S}_{0,3}$ ,
qui stabilise $L_\infty \subset \Pi_{0,3}$ (donc son action
sur le troisième facteur $\tilde{\varepsilon}_0 = \dots$ sera
de le [unclear: remplacer] par son inverse [unclear: $\tilde{\varepsilon}_0^{-1}$],
ce qui aura la vertu de le conserver [unclear: 2-continu]
unipotent, auquel nous tenons !).

\newpage

%% ===== Page 368 =====

VIII Récapitulation sur l'équation $^{(1)} [\alpha,\rho]=[\beta,\sigma]\varepsilon_1(\nu)$ .

Le résultat essentiel auquel nous sommes parvenus dans (VII), c'est la possibilité de paramétriser les solutions d'une telle équation $\alpha,\beta,\nu$ (sous des hypothèses convenables sur $\rho,\sigma$) par un groupe $B'_0 \times \Gamma_\rho \times \Gamma_\sigma$, où $B'_0$ est un "pré-Borel" dans $G$ (précisément, pour une base convenable, la forme $B'_0 = \left\{ \begin{pmatrix} \mu & \tau \\ 0 & 1 \end{pmatrix} \mid \begin{matrix} \mu \in \Gamma_{G_0} \\ \tau \in \Gamma_{G_0} \end{matrix} \right\}$) et $\Gamma_\rho, \Gamma_\sigma$ sont les centralisateurs des sections régulières $\rho,\sigma$ de $G$, par le morphisme

(2) $\begin{cases} B'_0 \times \Gamma_\rho \times \Gamma_\sigma \longrightarrow \mathcal{G}_{\rho,\sigma} \\ (\underline{u}, \underline{a}, \underline{b}) \longmapsto (\underline{u}\underline{a}^{-1}, \underline{u}\underline{b}^{-1}, \underline{\delta}(\underline{u})) \end{cases}$

(Ceci le fait que $\underline{\nu}(\underline{u}) = \det(\underline{u})-1$ n'a pas pour le moment d'importance capitale).

[MARGIN: Nous partons maintenant de $\rho,\sigma, B'_0, \varepsilon_1$ (donc $L_1$) et nous posons la question : pour quelles [unclear] les paramétrisations (2) ...]

Pour qu'un tel morphisme se prolonge en un morphisme de $G \times G \times E'$ dans l'espace des solutions de (1), il faut et il suffit ainsi que l'on ait identiquement

$$\underline{u} \rho \underline{u}^{-1} = \underline{u} \sigma \underline{u}^{-1} \varepsilon_1(\underline{\nu}(\underline{u}))$$

ce qui s'écrit aussi, en simplifiant par $\underline{u}$ et prenant les inverses

(3) $\rho(\underline{u}) = \varepsilon_1(-\underline{\nu}(\underline{u}))\sigma(\underline{u})$ .

La possibilité de trouver $\underline{\nu} : B'_0 \to G_m$ (morphisme de schémas) donnant, via (2), un isomorphisme $B'_0 \times \Gamma_\rho \times \Gamma_\sigma \to \mathcal{G}_{\rho,\sigma} = \dots$

\newpage

%% ===== Page 369 =====

des solutions de (1), équivaut donc à celle
que [crossed out: pour tout $\sigma$ et $\rho$] $\rho$ admet des sections
sur $B'_0$, modulo multiplication [unclear: à gauche] par
le sous-groupe [unclear] $L_1$, qui
intervient dans l'écriture [unclear: univoque] de l'équation
(1) que nous étudions. Cela se conçoit
difficilement si [unclear: ce n'est] que
$L_1 \subset \rho(B'_0) \overset{df}{=} B'_1$, ce qui implique d'ailleurs
$\sigma(B'_0) \subset B'_1$ (et par un argument connu $\sigma(B'_0) = B'_1$)
et $L_1 = (B'_1)_m$. Inversement, si
$\rho(B'_0) = \sigma(B'_0)$, i.e. [crossed out: alors] $\sigma^{-1} \rho$ normalise $B'_0$,
donc i.e. $\sigma^{-1} \rho$ section de $B_0$, alors
$\sigma(u) = \sigma u \sigma^{-1}$ et $\rho(u) = \rho u \rho^{-1}$ sont des sections
de $B'_1$ qui ont [unclear: un = déterminant - sont]
donc qui coïncident bien modulo $(B'_1)_m$.

Ainsi on voit que la propriété essentielle
pour obtenir des solutions de (1) via (2),
est d'avoir
$$(\sigma^{-1} \rho) (B'_0) = B'_0 \quad \text{i.e.} \quad \sigma^{-1} \rho \in \Gamma B_0$$
L'unicité de la représentation de ces
solutions $(\alpha, \beta, \gamma)$ sous la forme $(\underline{u} a^{-1}, \underline{u} b^{-1}, \underline{v}(u))$
fait supposer de plus que l'on ait
[Il suffit pour ça que l'on ait $T_\rho \cap T_\sigma = \{1\}$ et $G_m \cap B'_0 = \{1\}$]
$$T_\sigma T_\rho \cap B'_0 = \{1\},$$
[crossed out: qui implique d'ailleurs $T_\rho \cap B'_0 = \{1\}$]

[MARGIN: NB Cette condition n'est pas rigoureuse ici. (Cf la [unclear] de la page qui suit). Pour assurer l'unicité de la représentation des solutions $(\alpha, \beta, \gamma)$ sous la forme [unclear] pratiquement il faut supposer (5) [unclear] (6)]

\newpage

%% ===== Page 370 =====

486

donc la flèche $(\underline{u}, \underline{a}, \underline{b}) \mapsto (U\natural, a\natural, b\natural)$, a une section de $T_\rho \cap T_\sigma \cap B'_0$. Reste la question de la surjectivité de (2) - toute solution de (4) s'écrit-elle bien sous la forme $(u a', v b', \nu(u))$ ? Je renonce à [unclear] de dire que l'on a surjectivité de

(6) 
$$ B'_0 \longrightarrow \mathcal{H}_{\rho, \sigma} $$
$$ u \longmapsto ([\underline{u}, \rho], [\underline{u}, \sigma], \underline{\nu}(u)) $$

où $\mathcal{H}_{\rho, \sigma}$ est le schéma des solutions simultanées en $\underline{\xi}, \underline{\eta}, \nu$ des

(7)
$$ \det \xi = \det \eta = 1, \text{Tr } \xi \rho = \text{Tr } \rho, \text{Tr } \eta \sigma = \text{Tr } \sigma, \xi = \eta \varepsilon_1(\nu). $$

C'est essentiel pour obtenir cette surjectivité que nous avions été amenés dans la section (II) à faire les hypothèses draconiennes (p. 435, 436) qui s'expriment, au prix d'un miracle ad-hoc, par les conditions que

(8)
$\rho, \sigma$ sections de $(B_0)_u = L_0$ (et par suite de $B_0$)
$B_0$ et $B_1 = \rho(B_0) = \sigma(B_0)$ en "position générale"
i.e. $\rho_0$ et $L_1 = \rho(L_0) = \sigma(L_0)$ indépendants ;
([unclear])

mais il n'est pas clair si ces conditions (suggérées par des calculs antérieurs, et donnant lieu en effet à de très jolis calculs) soient vraiment indispensables pour la validité de [unclear].

\newpage

%% ===== Page 371 =====

Dans notre contexte il est naturel de supposer $\varepsilon_0 = \delta^{-1} \rho$ unipotent (ou au moins quasi-unipotent, i.e. une puissance entière $\varepsilon_0^n$ unipotente), mais la condition b) de (8), plus subtile, apparait comme moins impérative. La suite montrera si cette restriction est judicieuse.

Dans le cas où on considère une équation

(9) $$[\alpha,\rho] = \varepsilon [\beta,\sigma] \varepsilon_1(v)$$

avec un paramètre $\varepsilon$ (pouvant notamment prendre la valeur importante -1) $\in [unclear]$, on a été amené pour la première fois dans ces calculs détruit la symétrie entre $\rho$ et $\sigma$, en introduisant la condition $Tr \sigma = 0$ comme condition de stabilité des solutions $\mathcal{H}_{\rho,\sigma}$ en $\eta$ satisfaisant

$$\det \eta = 1, \quad Tr \eta \sigma = Tr \sigma$$

pour $\mathcal{fl}_\sim$. Le résultat essentiel auquel nous tendions, était l'isomorphisme

(10) $$B'_0 \times T_{\rho} \times \mathcal{N}'_\sigma \overset{\sim}{\longrightarrow} \mathcal{G}_{\rho,\sigma}^*$$
$$(\underline{u}, \underline{a}, \underline{b}) \longmapsto (\alpha = \underline{a}^{-1}, \beta = \underline{u} \underline{b}^{-1}, \mu = \det \eta, \varepsilon = \varepsilon(\underline{b}))$$

où $\mathcal{N}'_\sigma \subset \mathcal{N}_\sigma$ est naturellement le normalisateur de $T_\sigma$ defini par la condition

(11) $$\underline{b} \sigma (\underline{b})^\sim = \varepsilon(\underline{b}) \sigma \quad \text{i.e.} \quad \underline{b}(\sigma) = \varepsilon(\underline{b}) \sigma$$

\newpage

%% ===== Page 372 =====

Que le morphisme (10) s'envoie bien dans $G_{\rho, \sigma}^*$,
est trivial par définition de $N_\sigma'$ et via le fait
analogue pour (2) - sans hypothèse du genre
$Tr \sigma = 0$. Mais la possibilité d'obtenir (localement)
n'importe quel $\varepsilon$ via un $b \in \Gamma N_\sigma'$ - ou
simplement $\varepsilon = -1$, si 2 invers.) exige justement
$Tr \sigma = 0$.

Bien entendu, nous considérons
$G_{\rho, \sigma}^*$ comme un schéma en groupes,
par transport de structure
de groupe provenant de $B_0' \times T_p \times N_\sigma'$.
On y distingue le sous-groupe

(12) $G_{\rho, \sigma}^{* (0)} \text{ correspondant à : } T_0' \times T_p \times N_\sigma'$

où

(13) $T_0' = B_0' \cap B_1'$

est le tore maximal privilégié de $B_0'$, qui
pour notre épinglage canonique s'écrit
comme

(14) $T_0' = \left\{ \begin{pmatrix} \mu & 0 \\ 0 & \mu^{-1} \end{pmatrix} \mid \mu \in \Gamma G_m \right\}$

On pose de même

(15) $G_{\rho, \sigma}^{\#} = G_{\rho, \sigma} \cap G_{\rho, \sigma}^{* (0)} \text{, correspondant à : } T_0' \times T_p \times T_\sigma$

\newpage

%% ===== Page 373 =====

Bien entendu donne des structures de produit
semi-direct

(16) $$G_{\rho,\sigma} \simeq G_{\rho,\sigma}^* \ltimes \Lambda_0 , \quad G_{\rho,\sigma}^{* \natural} \simeq G_{\rho,\sigma,(0)}^{* \natural} \ltimes \Lambda_0 ,$$

où
(17) $$\Lambda_0 \subset G_{\rho,\sigma} , \Lambda_0 = \left\{ (\alpha = \lambda, \beta = \lambda, \mu = 1) \mid \lambda \in \Gamma L_0 \right\}$$
[unclear: $\varepsilon_0^n \quad \varepsilon_0^n \quad ( \lambda \in \Gamma G_0 )$]
correspondant à $L_0 \times \{1\} \times \{1\} \subset B'_0 \times T_\rho \times T_\sigma$

Nous allons enfin considérer un autre
schéma en
sous-groupes de $G_{\rho,\sigma}^*$ par les
conditions

(18) $$ \begin{cases} \det \alpha \in \{ \pm 1 \} \text{ i.e. } (\det \alpha)^2 = 1 \\ \det \beta = 1 \end{cases} \quad \text{d'où } SG_{\rho,\sigma}^* $$

[unclear: crossed out text] $SG_{\rho,\sigma}^*$ peut s'interpréter
comme le schéma des systèmes
$(\alpha, \beta, \mu, \varepsilon, \varepsilon')$
satisfaisant $\quad \det \beta = 1$,

(19) $[\alpha, \beta] = \varepsilon [\beta, \sigma] \varepsilon (\mu - 1) , \det \alpha = \varepsilon' , (\varepsilon^2 = {\varepsilon'}^2 = 1$
[MARGIN: NB $\varepsilon^2 = 1$ résulte déjà de la première relation.]

En termes des systèmes
$(\underline{u}, \underline{a}, \underline{b}) \in \Gamma B'_0 \times T_\rho \times N'_0$
les équations (19) (dont la première est
automatique) se réduisent à

(20) $\det \underline{u} = \det \underline{b} = \varepsilon' \det \underline{a} , \quad {\varepsilon'}^2 = 1 ,$

\newpage

%% ===== Page 374 =====

qui montrent que $Sg_{p,\sigma}^*$ est [unclear: bien] un sous-groupe
de $G_{p,\sigma}^*$. Ce sous-groupe contient $\Lambda_0$
(correspondant à $L_0 \times \{1\} \subset B'_0 \times T_p \times N'_0$ ) - qui
correspond au cas où
(21) $$det \underline{v} = det \underline{b} = det \underline{a} = 1 \quad (\varepsilon' = 1) \quad ,$$
et on a une structure de produit semi-direct
(22) $$Sg_{p,\sigma}^* \simeq Sg_{p,\sigma}^{*(0)} \cdot \Lambda_0$$
où
(23) $$Sg_{p,\sigma}^{*(0)} = Sg_{p,\sigma}^* \cap g_{p,\sigma}^{*(0)}$$
correspond au sous-groupe de $B'_0 \times T_p \times N'_0$
défini par les deux équations
(20). Bien sûr, on a des
(24) $$Sg_{p,\sigma} \quad , \quad Sg_{p,\sigma}^{(0)} \quad (\text{ss-}[unclear])$$
d'intersections de $Sg_{p,\sigma}^* , Sg_{p,\sigma}^{*(0)}$ avec
$Sg_{p,\sigma}$ - se déduisent de l'expression
de ceux-ci en termes de $(B'_0, T_0), T_p, N'_0$
en remplaçant $N'_0$ par $T_\sigma$.

Pour expliciter la structure de $Sg_{p,\sigma}^*$,
on est ramené vu (22) à étudier le
groupe (25) $$Sg_{p,\sigma}^{*(0)} \hookrightarrow T_0 \times T_p \times N'_0$$

\newpage

%% ===== Page 375 =====

décrit par les équations (20). Mais on a

$$T_0 \stackrel{\sim}{\longrightarrow} G_m$$

$$u \longmapsto \det u$$

ce qui implique que la projection

$$S G_{p, \sigma}^{*(0)} \longrightarrow T_p \times N_0'$$

$$(\underline{u} \underline{a}'', \underline{u} \underline{b}'', \det \underline{u}, \varepsilon = \varepsilon(\underline{b})) \longmapsto (\underline{a}, \underline{b})$$
$$\simeq$$
$$(\underline{u}, \underline{a}, \underline{b})$$

est [unclear: stricte], et induit un iso. de

$S G_{p, \sigma}^{*(0)}$ sur le sous-groupe de $T_p \times N_0'$

défini par l'équation

(21) $$\det \underline{b} = \varepsilon' \det \underline{a} \quad , \quad \varepsilon'^2 = 1$$

ou encore $(\det \underline{b})^2 = (\det \underline{a})^2$ .

En termes des homomorphismes ([unclear: déduits]

[unclear: naturellement] )

$$T_p \stackrel{\delta_p}{\longrightarrow} G_m \quad , \quad N_0' \stackrel{\delta_{0'}}{\longrightarrow} G_m$$

cela revient à prendre l'image inverse de la
diagonale de $G_m \times G_m$ par $\delta_p^2, \delta_{0'}^2$.

On voit tout de suite (comme $\delta_p, \delta_{0'}$

sont épi) que l'hom.

$$S G_{p, \sigma}^{*(0)} \longrightarrow \mu_2 \times \mu_2$$

$$([unclear], \varepsilon, \varepsilon') \longmapsto (\varepsilon, \varepsilon')$$
$$\text{i.e.} \quad (\underline{u}, \underline{a}, \underline{b}) \longmapsto (\varepsilon(\underline{b}), \det \underline{b} / \det \underline{a} )$$

est épi, et le noyau [unclear: donne lieu à] d'une

suite exacte

374

\newpage

%% ===== Page 376 =====

989

(23) $$1 \to SG_{\rho, \sigma^{(\omega)}}^{* \natural} \to SG_{\rho, \sigma}^{* \natural} \to \mu_2 \times \mu_2 \to 1$$

où $SG_{\rho, \sigma^{(0)}}^{* \natural}$ est isom. au sous-schéma de

$T_\rho \times T_\sigma$ défini par l'équation

(24) $$\det \underline{a} = \det \underline{b}$$

et se dévisse donc : [unclear] par un

(25) $$1 \to S T_\rho^\natural \times S T_\sigma^\natural \to SG_{\rho, \sigma^{(0)}}^{* \natural} \to \mathbb{G}_m \to 1$$
$$(\underline{a}, \underline{b}) \mapsto \det \underline{a} = \det \underline{b}$$

où on a posé

(26) 
$$S T_\rho^\natural = \operatorname{Ker} (T_\rho \xrightarrow{\delta_\rho} \mathbb{G}_m)$$
$$S T_\sigma^\natural = \operatorname{Ker} (T_\sigma \xrightarrow{\delta_\sigma} \mathbb{G}_m)$$

Si par exemple $\rho$ est [unclear], donc $T_\rho$ un tore maximal, alors $S T_\rho^\natural$ est lui-m. un tore [MARGIN: de dim $\le 1$]

Ainsi, dans le "cas modulaire"

$$\rho = \begin{pmatrix} 0 & 1 \\ 1 & 0 \end{pmatrix}, \quad \sigma = \begin{pmatrix} 0 & -1 \\ 1 & -1 \end{pmatrix}, \quad S T_\rho^\natural \text{ est un}$$

tore de dim 1 en car. 3, $S T_\sigma^\natural$ tous de dim 1 en car. 2. Dans ce cas, $SG_{\rho, \sigma^{(0)}}^{* \natural}$ est un tore de dim. 3.

Pour terminer de déterminer la structure de $SG_{\rho, \sigma}^{* \natural}$ en termes de $SG_{\rho, \sigma^{(0)}}^{* \natural}$ et de la décomposition (22) en produit semi-direct, il faut expliciter l'opération de

\newpage

%% ===== Page 377 =====

$SG_{\rho,\sigma}^{*(0)}$ sur $\Lambda_0$, en identifiant
$SG_{\rho,\sigma}^{*(0)}$ comme sous-groupe de $T_\rho \times N'_\sigma$,
et $\Lambda_0$ comme $\simeq \mathbb{G}_m$ par $\lambda \mapsto \varepsilon_0(\lambda)$.

Or la relation

$$ \underline{u} (\varepsilon_0(\lambda)) \underline{u}^{-1} = \varepsilon_0( \det(\underline{y}) \lambda ) \quad , \quad \det(\underline{y}) = \det(b) = \delta_\sigma(b) $$

montre que l'opération se fait via
l'application composée

(27)
$$ SG_{\rho,\sigma}^{*(0)} \hookrightarrow N'_\sigma \xrightarrow{\delta_\sigma} \mathbb{G}_m \quad , \quad OK. $$

Pour la suite, nous aurons besoin
d'un hom. canonique

(28)
\begin{tikzcd}
G_{\rho,\sigma}^* \ar[r, "t"] & G \\
B'_0 \times T_\rho \times N'_\sigma \ar[ur, dashed] &
\end{tikzcd}

qui s'identifie à la première projection
$$ (y, \underline{a}, \underline{b}) \mapsto y $$
du produit, mais considérée comme
un morphisme dans $G$ grâce à l'inclusion
$B'_0 \hookrightarrow G$.

On pose aussi
$$ u = \underline{\Theta}(\alpha, \beta, \mu, \varepsilon), $$
ce dont on aura [unclear]

\newpage

%% ===== Page 378 =====

$$ (29) \quad \begin{cases} \underline{u}(\rho) = \xi \rho \\ \underline{u}(p) \sim \text{int}(h) p \\ \det \underline{u} = \mu \end{cases} \quad , \quad \begin{matrix} \underline{u}(\sigma) = \gamma \sigma \\ \sim \varepsilon \text{ int}(p)(\sigma) \end{matrix} \quad \left\{ \begin{matrix} \xi = \alpha \rho(\alpha)^{-1}, \gamma = \beta \sigma \beta^{-1} \\ \text{à un facteur } \varepsilon ! \end{cases} \right. $$

Proposition. — Le morphisme de schémas $\Theta$ est caractérisé par les formules (29), mod un morphisme $G_{\rho, \sigma} \xrightarrow{\lambda} \mu_2$.
En effet, les premières de ces formules, qui déterminent l'action de $\underline{u}$ sur $\rho$ resp. $\sigma$, déterminent (comme on sait) $\underline{u}$ mod $T_\rho$ resp. mod $T_\sigma$, donc (soit $T_\rho \cap T_\sigma = \mu_2$), la deuxième la détermine mod $SG_1$, — dans le cas précédent, mod $\mu_2$ — donc un autre morphisme $\Theta'$ ayant ces propriétés sera de la forme
$$\Theta'(g) = \lambda(g) \Theta(g), \quad \text{où } \lambda: G_{\rho, \sigma} \to \mu_2.$$

Rem. d'ailleurs Notons que $\Theta$ est normalisé par
$$\Theta(1,1,1,1) = 1$$
ce qui impose une normalisation supplémentaire $\lambda(1) = 1$, sur $\lambda$ — mais il y en aura une [unclear] (par exemple $\rho, \sigma$ [unclear] dans $T_\rho, T_\sigma$ des torus...) qui donne lieu...

[MARGIN: Si on travaillait dans $SG_{\rho, \sigma}$ [unclear] sur la composante neutre [unclear] de se multiplier sur $B'_0 \times T_\rho \times T_\sigma$ (de dimension $n+2$) par des termes [unclear] essentiels comme [unclear] indétermination possible la multiplication de $u$ par $\varepsilon$ (seule correspondant à un hom. non trivial de $N_0 / T_\sigma \simeq \mu_2$ dans $\mu_2$).]

\newpage

%% ===== Page 379 =====

(IX) Les groupes $\mathcal{U}_{\rho,\sigma}$ et $S\mathcal{U}_{\rho,\sigma}$.

Soient tjrs $\rho,\sigma$ deux sections de $\underline{G}$, satisfaisant
les conditions générales (cf p.ex. pag p. 478),
plus la condition $Tr \sigma=0$ (cf V - p.ex.
lemme p. 415 et son cor. p. 457) - nous pourrions
faire leur construction sans cette condition
supplémentaire (en nous contentant de
travailler avec $\underline{G}_{\rho,\sigma}$ au lieu de $\underline{G}_{\rho,\sigma}^*$) -
mais peu importera au reste pour la suite.
Nous pourrons tjrs supposer $\rho,\sigma$ normalisés
de façon qu'on ait $Tr(\rho+\sigma)=1$.

On désigne par $S\mathcal{M}_{\rho,\sigma}$ le schéma

[MARGIN: Par la suite, il "apparaîtra" préférable de noter ces groupes $S\mathcal{M}_{\rho,\sigma}$, cf p. 559 ff]

produit
(1) $$S\mathcal{M}_{\rho,\sigma} = S\underline{G}_{\rho,\sigma}^* \times S\underline{G}$$
ses sections sont notées
$$(\alpha,\beta,\mu,\varepsilon,f)$$
$f$ étant donc une section arbitraire de $S\underline{G}$
$$(= \operatorname{Sl}(2)_S)$$.

Dans le cas où $S=\operatorname{Spec}(A)$, $A=\hat{\mathbb{Z}}$,
$\underline{G}=\operatorname{Gl}(2)_S$, $\rho = \begin{pmatrix} 0 & 1 \\ -1 & 1 \end{pmatrix}$, $\sigma = \begin{pmatrix} 0 & -1 \\ 1 & 0 \end{pmatrix}$,
et on a donc un homomorphisme
canonique
(2) $$ \widehat{\underline{G}}_{\rho,\sigma} \to G(\hat{\mathbb{Z}})^{\pm 1} \simeq \operatorname{Gl}(2,\hat{\mathbb{Z}})^{\pm 1}/(\pm 1) $$
à valeurs dans le sous-groupe des éléments de det $\pm 1$,

\newpage

%% ===== Page 380 =====

envoyant $\hat{\mathcal{C}}_{0,3}$ dans $\underline{S}G(\hat{\mathbb{Z}})' = \operatorname{Sl}(2,\hat{\mathbb{Z}})' = \operatorname{Sl}(2,\hat{\mathbb{Z}})/_{\pm 1}$
on va construire une variante du groupe
(cf p. [unclear: 374]),
$\mathcal{M}_{0,3}^\sim$, qui s'envoie canoniquement dans le
produit (1), ce qui nous amènera à définir
une structure de groupe sur $S\mathcal{M}_{\rho,\sigma}$, de
façon à rendre ledit morphisme multiplicatif.

Considérons d'abord la variante "spéciale"
$S\mathcal{M}_{0,3}^\sim$ de $\mathcal{M}_{0,3}^\sim$, extension de $\mathcal{M}_{0,3}^\sim$
par $\hat{\mathbb{Z}}^3$, déduite d'une extension de
$\overline{\mathcal{M}}_{0,3} = \mathcal{M}_{0,3}^\sim / \hat{\Pi}_{0,3} \simeq \hat{\mathcal{C}}_{0,3}$ , par [unclear]
cette extension de $\hat{\mathcal{C}}_{0,3}$ (via $\mathcal{M}_{0,3}^\sim \to \hat{\mathcal{C}}_{0,3}$)
formée des

(3) $(u, f_0, f_1, f_\infty) \quad (u \in \hat{\mathcal{C}}_{0,3}, f_0, f_1, f_\infty \in \hat{\mathbb{Z}})$

satisfaisant

(4) $u(l_i) = int(f_i^{-1})(l_{u(i)})$

où [unclear] la multiplication de u, et
où l'action de u sur $\{0, 1, \infty\}$ est
donnée [unclear]. Le produit est défini par

(5) $(u', f_0', f_1', f_\infty') (u, f_0, f_1, f_\infty) = (u'u, f_0' u'(f_0), f_1' u'(f_1), f_\infty' u'(f_\infty))$

bien entendu (disons $f_0 g_0$ !) - et on a
un $u \in S\mathcal{M}_{0,3}^\sim$ si $\prod f_i = 1$.

\newpage

%% ===== Page 381 =====

On a un hom. surjectif $S\mathcal{M}_{0,3}^\sim \to \mathcal{M}_{0,3}^\sim$ ,

dont le noyau est formé des

(6) $$ (1, f_0, f_1, f_\infty) \quad f_i \in L_i $$

d'où

(7) $$ 1 \longrightarrow \hat{\mathbb{Z}}^3 \longrightarrow S\mathcal{M}_{0,3}^\sim \longrightarrow \mathcal{M}_{0,3}^\sim \longrightarrow 1 $$

Nous allons définir un sous-groupe intéressant

(8) $$ S_0\mathcal{M}_{0,3}^\sim \subset \mathcal{M}_{0,3}^\sim \ , $$

en utilisant maintenant la définition de

$\mathcal{M}_{0,3}^\sim$ , savoir la propriété des $u \in \mathcal{M}_{0,3}^\sim$

de commuter extérieurement à $\rho$ , de sorte qu'on a

(9) $$ u \rho u^{-1} = g \rho \quad \text{i.e.} \quad [u, \rho] = g , \quad g \in \hat{\Pi}_{0,3} $$

Ceci permet, quand on connait un des $f_i \in \hat{\Pi}_{0,3}$

satisfaisant

(10) $$ u(l_i) = int(f_i^{-1})(l_i^u) $$

(en supposant pour simplifier $u \in S\mathcal{M}_{0,3}^\sim$ ),

d'en déduire les autres $f_j$, via

la formule

(11) $$ f_{i+1} = g \rho(f_i) \quad \left[ \text{et } f_{i-1} = g \rho(f_{i+1}) = g \rho(g) \rho^2(f_i) \right] $$

(où $g = [u, \rho] \in \hat{\Pi}_{0,3}$ )

(on pose $\rho(i) = i+1$ pour les indices, d'où

$0+1=1, 1+1=\infty, \infty+1=0 \quad ! \quad$ ) on a

$$ f_i = f_{(i+2)+1} = g \rho(f_{i+2}) = \underline{ g \rho(g \rho(f_{i+1})) } = \underline{ g \rho(g) \rho^2(g) \underbrace{\rho^3(f_i)}_{f_i} } \ . $$

\newpage

%% ===== Page 382 =====

492

les $(u, f_0, f_1, f_\infty) \in \mathcal{S}M_{0,3}^\sim$ qui satisfont
(11) pour tout $i \in \{0,1,\infty\}$ (il suffit qu'il le
vérifie pour deux indices, pour que ce soit
vérifié pour tous les trois) forment un
sous-groupe de $\mathcal{S}M_{0,3}^\sim$, que je note
$\mathcal{S}_0M_{0,3}^\sim$ - il s'identifie aussi au groupe
des couples $(u, f_0)$ ($u \in M$, $f_0 \in \hat{\Pi}_{0,3}^\sim$)
satisfaisant
(12) $$u(l_0) = int(f_0^{-1})(l_0^u) \quad (l_0^u = \chi(u)) .$$

On a une structure d'extensions
(13) $$1 \to \hat{\mathbb{Z}} \to \mathcal{S}_0M_{0,3}^{!\sim} \to M_{0,3}^{!\sim} \to 1 ,$$

et un hom d'extensions
(14) 
\begin{tikzcd}
1 \ar[r] & \hat{\mathbb{Z}} \ar[r] \ar[d, "diagonale"'] & \mathcal{S}_0M_{0,3}^{!\sim} \ar[r] \ar[d] & M_{0,3}^{!\sim} \ar[r] \ar[d, "\simeq"] & 1 \\
1 \ar[r] & \hat{\mathbb{Z}}^3 \ar[r] & \mathcal{S}M_{0,3}^{\sim} \ar[r] & M_{0,3}^{\sim} \ar[r] & 1
\end{tikzcd}

(où $\mathcal{S}_0M_{0,3}^{!\sim}$ est [unclear: par définition] l'image inverse de $M_{0,3}^{!\sim} \subset M_{0,3}^\sim$ dans $\mathcal{S}M_{0,3}^\sim$). On peut donc dire que l'extension $\mathcal{S}M_{0,3}^\sim$ (de $M_{0,3}^\sim$ par $\hat{\mathbb{Z}}^3$)
est déduite, par l'hom. de groupes

(15) $$\hat{\mathbb{Z}} \to \hat{\mathbb{Z}}^3$$
$$n \mapsto (n_0=n, n_1=n, n_\infty=n)$$

d'une extension "plus fine" $\mathcal{S}_0M_{0,3}^{!\sim}$ de

\newpage

%% ===== Page 383 =====

$\mathcal{M}_{0,3}^{\sim}$ par $\hat{\mathbb{Z}}$.

[MARGIN: à $(u,f)$ [unclear] en suite [unclear] par un autom. de [unclear] $u \mapsto p u$]

Les calculs de $§ 46$ V (p 407 ff)
permettent d'autre part d'expliciter les
[unclear] $(x, f_0=f)$ en termes de systèmes

(16) \quad $(\alpha, \beta, f)$ \quad $\left\{ \begin{array}{l} \alpha \in \hat{\Pi}_{0,3}^{\sim} = \hat{\Pi}_{0,3} , \beta \in T_{0,3} = \text{im. inv. de } \{1, \tau\} \\ \text{par } \widehat{\mathfrak{S}}_{0,3} \to \mathfrak{S}_{3} \end{array} \right.$

satisfaisant des relations (ne faisant pas intervenir $f$)

(17) \quad $\alpha \rho(\alpha)^{-1} = \beta \sigma(\beta)^{-1} \varepsilon_1(\mu-1)$

où $\mu \in \hat{\mathbb{Z}}^\times$ est défini sans ambiguïté par
$\alpha$ et ne dépend seulement comme on sait du tirage
$x \mapsto \alpha \rho(\alpha)^{-1}$, c. à d. de l'image de $\alpha$ dans $G(\bar{\mathbb{Q}}/\mathbb{Q})$!

Le cas $f=1$ correspond au sous-groupe de
$S_0 \mathcal{M}_{0,3}^{\sim}$, qui s'identifie au normalisateur
de $f_0$ dans $\mathcal{M}_{0,3}^{\sim}$, que j'avais noté
$\mathcal{M}_{0,3}^{\sim}[0]$ ($§ 46$ V $\dots$) :

(18) \quad $\mathcal{M}_{0,3}^{\sim}[0] \simeq \left\{ (\alpha, \beta, f) \Big| \begin{array}{c} f=1 , \alpha, \beta \\ \text{satisfaisant} \\ (17) \end{array} \right\} \subset S_0 \mathcal{M}_{0,3}^{\sim}$

\quad \quad \quad $||$ def

\quad \quad \quad $\operatorname{Norm}_{\mathcal{M}_{0,3}^{\sim}}(f_0)$

Si on désigne par $u_{\alpha, \beta}$ l'autom de [unclear]
un tel autom défini en termes de

(19) \quad $\xi = \alpha \rho(\alpha)^{-1}$, $\eta = \beta \sigma(\beta)^{-1}$

\newpage

%% ===== Page 384 =====

(20) $$u(\rho) = \xi \rho, \quad u(\sigma) = \eta \sigma$$

(NB $\eta$ est défini en termes de $\xi$ grâce à (14), où $\mu = \mu(\xi)$ est déterminée par $\xi$), on trouve par construction (loc. cit.)

(21) $$u_{\alpha, \beta, f} = \text{int}(y^{-1}) u_{\alpha, \beta}$$

et un triplet $(\alpha, \beta, f)$ correspond à l'élément $(u_{\alpha, \beta, f}, f)$ de $S_0 M_{0,3}^{\sim \wedge}$ :

(22) $$(\alpha, \beta, f) \longmapsto (u_{\alpha, \beta, f}, f) \in S_0 M_{0,3}^{\sim \wedge}$$ .

Ayant ainsi explicité les éléments de $S_0 M_{0,3}^{\sim \wedge}$, on définit une flèche

(23) $$S_0 M_{0,3}^{\sim \wedge} \longrightarrow \underline{S}M_{\rho, \sigma}(\hat{\mathbb{Z}})' \overset{*}{=} S\underline{G}_{\rho, \sigma}(\hat{\mathbb{Z}})' \times Sell_2 \hat{\mathbb{Z}}$$

$$(\alpha, \beta, f) \longmapsto (\underline{\alpha}, \underline{\beta}, \mu (\xi = \alpha p(\underline{\alpha})^{-1}), \varepsilon_2(\mu), \underline{f})$$

où les soulignés $\underline{\alpha}, \underline{\beta}, \underline{f}$ désignent les images de $\alpha, \beta, f$ dans $Sell_2 \hat{\mathbb{Z}}$, et le ' accent ' dans $\underline{S}M_{\rho, \sigma} (\hat{\mathbb{Z}})'$ désigne l'ensemble des systèmes $(\underline{\alpha}, \underline{\beta}, \mu, \varepsilon, \underline{f})$
[crossed out: éléments $\in (Sell_2 \hat{\mathbb{Z}})^{\#} \times S\underline{G}l(2, \hat{\mathbb{Z}}) \times \hat{\mathbb{Z}} \times \{\pm 1\} \times$]
dans $Sell_2 \hat{\mathbb{Z}}$ (pas dans $Sell_1 \hat{\mathbb{Z}}$ !), $(Sell_2 \hat{\mathbb{Z}})$
satisfaisant la relation

(24) $$\underline{\alpha} p(\underline{\alpha})^{-1} = \varepsilon p(\underline{\beta})^{-1} \varepsilon_1(\mu-1)$$

(ce qui ne dépend pas des choix des signes .... )

\newpage

%% ===== Page 385 =====

On a "fait ce qu'il fallait" pour que la
restriction de (23) à $S_0 M^{\sim \prime}_{0,3}[0]$ ((18) identifié
aux $(\alpha, \beta, f)$ pour lesquels dans le groupe $S G^*_{p, \sigma}$ $\underline{f}=1$) soit un
homomorphisme - en définissant justement de
façon ad hoc (astucieuse) la loi de composition dans
$S G^*_{p, \sigma}$. Je vais la réécrire quand même
pour me rassurer, en considérant $u=u_{\alpha,\beta}$, $u'=u_{\alpha',\beta'}$,
$u''=uu'$ sous la forme $u_{\alpha'',\beta''}$, et on se rappelle
la "relation des cocycles" qui donne, à un [unclear] du
dessus près, $(\alpha'',\beta'')$ en termes de $(\alpha',\beta')$, $(\alpha,\beta)$ (§46 II 5°))

(25) $\alpha'' = u'(\alpha)\alpha' \quad \beta'' = u'(\beta)\beta'$ [MARGIN: à un facteur $\mu'' = \mu \mu'$ près dont je]
[unclear: n'ai cure ici]

Soient alors $\underline{\alpha}, \underline{\beta}, \dots$ des relevés dans
$\operatorname{Gl}(2, \hat{\Z})^{\pm}$ (déterminés aux signes près par
ceux des images de $\alpha, \alpha'$, $\beta, \beta'$ de
$\operatorname{Gl}(2, \hat{\Z})'$, etc) avec donc $(\underline{\alpha}, \underline{\beta}, \mu, \varepsilon_1, \varepsilon_2)$ et $(\underline{\alpha}', \underline{\beta}', \mu', \varepsilon'_1, \varepsilon'_2)$
dans $S G^*_{p, \sigma}(\hat{\Z})$, et appliquons $\Theta$ - trouvons
$\underline{u}, \underline{u}' \in \operatorname{Gl}(2, \hat{\Z})'$ (en fait $\in B'_0(\hat{\Z})$) tels que

(26) $\begin{cases} \underline{u}'(\rho) = int(\underline{\alpha}')\rho \quad \underline{u}'(\sigma) = \varepsilon'_2(\rho) int(\underline{\beta}')(\sigma) \\ \det \underline{u}' = \mu' \end{cases}$

(et relations itou pour $\underline{u}$)
et on aura alors, pour tout $x \in \hat{\pi}_{0,3}$

(27) $\underline{u}''(x) = \underline{u}'(\underline{u}(x))$

(cf. [unclear]) Donc, pour des relevés convenables
$\underline{\alpha}'', \underline{\beta}''$ des images de $\alpha'', \beta''$ dans $\operatorname{Gl}(2, \hat{\Z})^{\pm}$, on aura:

(28) $\underline{\alpha}'' = \underline{u}'(\underline{\alpha}) \underline{\alpha}'$, $\underline{\beta}'' = \underline{u}'(\underline{\beta}) \underline{\beta}'$ [MARGIN: à un facteur $\mu'' = \mu\mu'$ près $\varepsilon''_2 = \varepsilon_2 \varepsilon'_2$]

\newpage

%% ===== Page 386 =====

et il reste à vérifier que l'élément
$$g"=(\underline{\alpha}", \underline{\beta}", \mu"=\mu\mu', \varepsilon"_2=\varepsilon'_2\varepsilon_2) \in [unclear: \Gamma(\hat{Z}; ... )]$$
est bien le produit des éléments
$$g' = (\underline{\alpha}', \underline{\beta}', \mu', \varepsilon') \quad \text{et} \quad g = (\underline{\alpha}, \underline{\beta}, \mu, \varepsilon).$$
Ceci est ok pour $\mu"$ et $\varepsilon"$ [unclear: est trivial], il reste
à vérifier que dans $\mathcal{G}_{\rho, \sigma}$, le produit est bien
celui qui convient, les composantes en $\alpha, \beta$ (resp. $\alpha', \beta'$)
données par les formules (28), ce qui
sera valable, bien entendu, [unclear: sans aucun cas]
quelconque, et dans les conditions générales du
V (on [unclear: l'utilisera au III] du reste pour $\mathcal{G}_{\rho,\sigma}$ au lieu
de $\mathcal{G}^*_{\rho,\sigma}$). On aura

(29)
$\underline{\alpha} = \underline{u}\underline{a}^{-1}$, $\underline{\beta} = \underline{u}\underline{b}^{-1}$ \qquad $\underline{\alpha}, \underline{\alpha}' \in \hat{T}_\rho(\hat{Z})$
$\underline{\alpha}' = \underline{u}'\underline{a}'^{-1}$, $\underline{\beta}' = \underline{u}'\underline{b}'^{-1}$ \qquad $\underline{\beta}, \underline{\beta}' \in \hat{N}_\sigma(\hat{Z})$

et par définition du produit dans $\mathcal{G}_{\rho,\sigma}$,
le produit de $g'g$ aura de même
composantes $\underline{\alpha}"$ et $\underline{\beta}"$ :
$$\underline{\alpha}" = (\underline{u}'\underline{u})(\underline{a}'\underline{a})^{-1} \quad , \quad \underline{\beta}" = (\underline{u}'\underline{u})(\underline{b}'\underline{b})^{-1}$$

et il reste à prouver
$$\underline{\alpha}" = \underline{\alpha}" \quad , \quad \underline{\beta}" = \underline{\beta}" \quad , \text{soit les relations}$$

(30)
$$\underline{u}'(\underline{\alpha}) . \underline{\alpha}' = (\underline{u}'\underline{u})(\underline{a}'\underline{a})^{-1} \qquad \underline{u}'(\underline{\beta}) \underline{\beta}' = (\underline{u}'\underline{u})(\underline{b}'\underline{b})^{-1}$$

où $\underline{\alpha}, \underline{\beta}, \underline{\alpha}', \underline{\beta}'$ sont donnés en fonction de
$\underline{u}, \underline{a}, \underline{u}', \underline{a}', \underline{b}, \underline{b}'$ par (29). On a donc
$$\underline{u}'(\underline{\alpha}) \underline{\alpha}' = \underline{u}'(\underline{u}\underline{a}^{-1}) \underline{u}'\underline{a}'^{-1} = (\underline{u}'\underline{u})(\underline{a}^{-1}\underline{a}'^{-1}) = (\underline{u}'\underline{u})(\underline{a}'\underline{a})^{-1},$$
Idem pour la 2ème des relations 30,
ce qui prouve notre assertion.

\newpage

%% ===== Page 387 =====

Notons à cette occasion que les systèmes
(31) $(\alpha, \beta, 1, 1)$ avec $\alpha \in \{\pm 1\}$, $\beta \in \{\pm 1\}$
sont dans $S\mathcal{G}_{\rho, \sigma^\circ}$. [unclear]
par multiplications, d'un $(\alpha, \beta, \mu, \varepsilon)$,
des $\alpha$ et $\beta$ [unclear: donnés] par une section de [unclear]
[unclear] par $\tilde{\alpha}, \tilde{\beta}$ ...] ils correspon-
dent, par l'isomorphisme $\mathcal{G}_{\rho, \sigma^\circ} \simeq B^\times_0 \times T_\rho \times T_{\sigma^\circ}$
aux systèmes $(\underline{u}, \underline{a}, \underline{b})$ avec $\underline{u} = 1$, $\underline{a} \in \Gamma_{...}$,
$\underline{b} \in \Gamma_{...}$ ) et forment un sous-groupe
(avec la multiplication ordinaire [unclear]
comme loi des groupes) - si $\alpha, \beta$
sont sections des [unclear], forment [unclear] des sous-groupes
"de sections triviales" de $S\mathcal{G}_{\rho, \sigma^\circ}$. L'action
par translation de ce sous-groupe $\simeq \Gamma_m \times \Gamma_m$
sur $\mathcal{G}^*_{\rho, \sigma^\circ}$ se fait par multiplication des
deux composantes en question. Cela montre
que le fait de travailler avec des
$\alpha, \beta$ définis seulement "à signes près"
revient simplement à travailler dans les
groupes quotients de $\mathcal{G}^*_{\rho, \sigma^\circ}(\hat{\mathbb{Z}})$ (ou de $S\mathcal{G}^*_{\rho, \sigma^\circ}(\hat{\mathbb{Z}})$)
par le sous-groupe $Z_{\rho, \sigma^\circ} \simeq \mathbb{Z}/2\mathbb{Z} \times \mathbb{Z}/2\mathbb{Z}$ engendré des
éléments $(\alpha, \beta, 1, 1)$ avec $\alpha, \beta \in \{\pm 1\}$. Ceci

\newpage

%% ===== Page 388 =====

495

précisera en même temps le sens de la locution :
" homom. de groupes ... ".

Revenons maintenant à l'application (23)
avec $f$ quelconque, et précisons la loi de
composition dans $S_0 \mathcal{M}_{0,3}^{\sim}$ en termes des
$(\alpha, \beta, f)$ i.e. en termes des $u_{\alpha,\beta}$. Considérons
le sous-groupe formé des $(1,1, f)$, ou encore
ce qui correspond aux couples $(u, f)$ avec
$u = int(f)$, i.e. aux couples $(int(f), f)$
— il forme un groupe isomorphe à $\Pi_{0,\hat{3}}$,
et on a une structure de produit semi-direct :

(31) $$S_0 \mathcal{M}_{0,3}^{\sim} \simeq \mathcal{M}_{0,3}^{\sim} [0] \cdot \Pi_{0,\hat{3}}$$

(32) 
$$ \begin{cases} \Pi_{0,\hat{3}} \longrightarrow S_0 \mathcal{M}_{0,3}^{\sim} \\ f \longmapsto ( [unclear], f ) \leftrightarrow (1,1, f^{-1}) \end{cases} $$
l'opération de $\mathcal{M}_{0,3}^{\sim} [0]$ sur $\Pi_{0,\hat{3}}$ étant
l'opération tautologique (déduite de
l'homom. $\mathcal{M}_{0,3}^{\sim} [0] \subset \mathcal{M}_{0,3}^{\sim} \to \operatorname{Aut}(\Pi_{0,\hat{3}})$).

Le sous-groupe $\hat{\mathbb{Z}}$ de $S_0 \mathcal{M}_{0,3}^{\sim}$ est
donné par les couples $(u,f) = (1, l_0^n)$ ($n \in \hat{\mathbb{Z}}$)
qui sont ici mis sous la forme $(\alpha, \beta, f)$
c.à.d. tel que $u_{\alpha,\beta} = f u = int(l_0^n)$,
donc $\alpha = \beta = l_0^n$, i.e.

(33) $$\hat{\mathbb{Z}} \hookrightarrow S_0 \mathcal{M}_{0,3}^{\sim} \text{ s'identifie à } n \mapsto (l_0^n, \beta=l_0^n, f=l_0^n)$$

387

\newpage

%% ===== Page 389 =====

On aura
$$(f'^{-1} u_{\alpha', \beta'}) (f^{-1} u_{\alpha, \beta}) = \text{[crossed out: } f'^{-1} (f' u_{\alpha', \beta'} f'^{-1}) \dots \text{]}$$
$$f'^{-1} u_{\alpha', \beta'}(f^{-1}) (u_{\alpha', \beta'} u_{\alpha, \beta})$$

i.e. la multiplication dans $S_0 \mathcal{M}_{0, 3}^{\sim}$ se fait,
en termes des composantes $(\alpha', \beta', f')(\alpha, \beta, f)$,
par $(\alpha'', \beta'', f'')$, avec

(34)
$$
\begin{cases}
u_{\alpha'', \beta''} = u_{\alpha', \beta'} u_{\alpha, \beta} \text{ i.e. } \begin{cases} \alpha'' = u_{\alpha', \beta'}(\alpha) \alpha' \\ \beta'' = u_{\alpha', \beta'}(\beta) \beta' \end{cases} \text{ , et} \\
f'' = u_{\alpha', \beta'}(f) f'
\end{cases}
$$

Revenons maintenant à la défini-
-tion de $S \mathcal{M}_{g, \nu}$ (1), nous allons y intro-
-duire la multiplication correspondante, ce
qui revient à définir $S \mathcal{M}_{g, \nu}$ comme
un produit $1/2$ direct

(35)
$$ S \mathcal{M}_{g, \nu} \overset{df}{=} S \mathcal{G}_{g, \nu} \overset{*}{.} G $$
(produit $1/2$ direct)

par l'opération de $S \mathcal{G}_{g, \nu}$ sur $G$
via l'homomorphisme $\Theta$ (p. 488') ,

et on définit un isomorphisme

(36)
$$ S \mathcal{M}_{g, \nu} \longrightarrow S \mathcal{G}_{g, \nu} \times_s G $$
$$ \text{[crossed out: } u_{\dots} \cdot f \text{]} \quad u \cdot f \longmapsto (u, f^{-1}) $$
[MARGIN: * Attention on [unclear] à $f^{-1}$ !]

où $u = (\alpha, \beta, \mu, \varepsilon)$ est dans $S \mathcal{G}_{g, \nu}$ .

\newpage

%% ===== Page 390 =====

495

On a de fait ce qu'il fallait, dans le
cas $S = \operatorname{Spec}(\hat{\mathbb{Z}})$ etc, pour avoir un
homomorphisme de groupes

(37) $$ \widehat{S_0 \mathcal{M}_{0,3}} \stackrel{\sim}{\longrightarrow} S\mathcal{M}_{\bar{p},\sigma}(\hat{\mathbb{Z}}) / \Sigma $$

où $\Sigma_0 \simeq \mathbb{Z}/2 \times \mathbb{Z}/2$ est le s-groupe de $S\mathcal{G}_{\bar{p},\sigma}^*$ introduit p. 494', qui opère
trivialement sur $G$, donc est central dans
$S\mathcal{M}_{\bar{p},\sigma}$, donc dans $S\mathcal{M}_{\bar{p},\sigma}(\hat{\mathbb{Z}})$, et où

(38) $$ \Sigma = \Sigma_0 \times \{\pm 1\} \simeq \mathbb{Z}/2 \times \mathbb{Z}/2 \times \mathbb{Z}/2 $$
[MARGIN: qui s'identifie au centre de $S\mathcal{G}_{\bar{p},\sigma}^*$
$= \{\pm 1\}$, centre de $\operatorname{Gl}(2,\hat{\mathbb{Z}})$
$= G(\hat{\mathbb{Z}})$]

est le s-groupe distingué correspondant
dans le produit 1/2 direct - le fait
de passer au quotient par $\Sigma$ expulse
les trois ambiguïtés de signes, sur $\alpha$, sur $\beta$
et sur $f$.

Considérons maintenant l'hom.

(39) $$ \Gamma_A \longrightarrow \widehat{S_0 \mathcal{M}_{0,3}} $$
$$ A \longmapsto (\alpha = \varepsilon_0(A), \beta = \varepsilon_0(A), \mu=1, \varepsilon=1, f = \varepsilon_0(A)) $$

c'est un hom. de groupes, et fait
de $\Gamma_A$ un s-groupe distingué.
Pour ceci le plus commode est d'identifier
$S\mathcal{M}_{\bar{p},\sigma}$ au produit 1/2-direct (35),

389

\newpage

%% ===== Page 391 =====

où on identifie $S\mathcal{G}_{p,\vec{\sigma}}^\ast$ à un sous-
groupe de $B'_0 \times T_p \times N'_0$, opérant
sur $G$ par autom. ext. via $B'_0$.

Le morphisme précédent correspond
alors à 
[MARGIN: l'hom. du groupe canonique]
$$ (40) \quad \lambda \mapsto (\lambda, 1, 1, \lambda^{-1}) $$

[unclear: crossed out text]
Soit d'autre part $u = (\lambda, \beta, \varepsilon, f)$ un
élément quelconque du produit $1/2$ direct
et calculons $u x(\lambda) u^{-1} -$
$u \cdot x(\lambda) \cdot u^{-1} = u_0 \cdot f x(\lambda) f^{-1} u_0^{-1} =$
[end crossed out text]

On identifie $S\mathcal{G}_{p,\vec{\sigma}}^\ast$ à un sous-groupe
de $\mathcal{G}_{p,\vec{\sigma}}^\ast$, et de ce produit $1/2$ direct
$S\mathcal{M}_{p,\vec{\sigma}}^\ast$ à un sous groupe du
produit $1/2$ direct
$$ (41) \quad \mathcal{G}_{p,\vec{\sigma}}^\ast \cdot G \simeq (B'_0 \times T_p \times N'_0) \cdot G $$
et on voit que l'image du morph. $i$
$(40)$ dans $\mathcal{G}_{p,\vec{\sigma}}^\ast \cdot G$ est distingué
dans ce groupe plus grand. Pour ne pas
nous y perdre dans les notations, je
désigne par 
[MARGIN: DIAGRAM: arrows from $B'_0, T_p, N'_0$ pointing to a bracket leading to $\mathcal{G}_{p,\vec{\sigma}}^\ast \hookrightarrow \mathcal{M}_{p,\vec{\sigma}}^\ast$, and an arrow from $G$ pointing directly to $\mathcal{M}_{p,\vec{\sigma}}^\ast$]
$$ (42) \quad i_B, i_p, i_N, i_G \text{ les morphismes d'inclusion} $$

\newpage

%% ===== Page 392 =====

[unclear: 297]

étant entendu, bien sûr, qu'on fait opérer
$G_{p,\sigma}^*$ sur $G$ intérieurement via le
facteur $B_0'$. Il faut vérifier que
$$i(L_0) = \{ i_B(\lambda) i_G(\lambda^{-1}) \mid \lambda \in \Gamma_{B_0} \}$$
est normalisé par $i_p(x), i_\sigma(x)$,
$i_B(x)$ et $i_G(x)$. Pour les deux premiers,
c'est trivial -- car action triviale sur $L_0$.
Si $x \in \Gamma_{B_0'}$, on aura

[MARGIN: d'où tout de suite que $i(L_0)$ est un sous-groupe invariant de [unclear]]

$$i_B(x)(i(\lambda)) = i_B(x)(i_B(\lambda) i_G(\lambda^{-1})) =$$
$$= i_B(x)(i_B(\lambda)) \cdot i_B(x)(i_G(\lambda^{-1}))$$
$i_B(x \lambda x^{-1})$ \quad $i_G(x)(i_G(\lambda^{-1})) = i_G(x \lambda^{-1} x^{-1})$
$$= i(x \lambda x^{-1})$$
ce qui montre que $L_0 \to G_{p,\sigma}^* \cdot G = \tilde{S} M_{p,\sigma}^*$
commute à l'action naturelle de $B_0'$ sur
les deux cotés (sur le deuxième via $i_B : B_0' \to \tilde{S} M_{p,\sigma}^*$). Enfin
$$i_G(x)(i(\lambda)) = i_G(x) \cdot (i_B(\lambda) i_G(\lambda^{-1})) \cdot i_G(x)^{-1}$$
or par définition, $i_B(\lambda)$ et $i_G(\lambda)$ opèrent
de la $=$ façon sur $G$, donc $i(\lambda) = i_B(\lambda) i_G(\lambda^{-1})$
y opère trivialement, donc le dernier facteur
$i_G(x)^{-1}$ peut sauter sur le facteur $i(\lambda)$, on trouve
$$i_G(x) i_G(x)^{-1} i(\lambda) = i(\lambda)$$
O.K.

\newpage

%% ===== Page 393 =====

\underline{Rmq} $\tilde{S}M_{\rho,\sigma}$ est un groupe plat de dim relative
$\underbrace{2}_{B'_0} + \underbrace{2}_{T_\rho} + \underbrace{2}_{N_0} + \underbrace{4}_{G} = 10$, $SM_{\rho,\sigma}$ est de dim relative
$4+3=7$, en passant au quotient par $i(L_0) \simeq G_a$
[unclear: crossed out text]
produit des groupes croisés $\tilde{M}_{\rho,\sigma}$, $M_{\rho,\sigma}$,
de dim. relative $9$ resp. $6$, de sorte qu'on
a un diagramme
$$
(43) \begin{tikzcd}
1 \ar[r] & G_a \ar[r, "\lambda \mapsto i(\lambda)"] \ar[d, equal] & S\tilde{M}_{\rho,\sigma} \ar[r] \ar[d, "sur"'] & \tilde{M}_{\rho,\sigma} \ar[r] \ar[d, "sur"'] & 1 \\
1 \ar[r] & G_a \ar[r] & SM_{\rho,\sigma} \ar[r] & M_{\rho,\sigma} \ar[r] & 1
\end{tikzcd}
$$
Le passage en quotient ne se fait pas sans
difficulté - les quotients sont en effet
représentables comme des produits 1/2 directs
$$
(44) \begin{cases} 
\tilde{S}M_{\rho,\sigma} \simeq S\tilde{G}_{\rho,\sigma}^* \ltimes SG \\ 
\tilde{M}_{\rho,\sigma} \simeq \tilde{G}_{\rho,\sigma}^* \ltimes G 
\end{cases}
$$
où les sous-schémas en groupes
$S\tilde{G}_{\rho,\sigma}^*$ et $\tilde{G}_{\rho,\sigma}^*$ de $S\tilde{G}_{\rho,\sigma}$, $\tilde{G}_{\rho,\sigma}$
sont définis dans (VIII), $\tilde{G}_{\rho,\sigma}^*$ s'identifiant
à $T_0 \times T_\rho \times N'_0$, et $S\tilde{G}_{\rho,\sigma}^*$ à son intersection
avec $S\tilde{G}_{\rho,\sigma}$ (défini par les conditions
simples (20) de p. 487!).

\newpage

%% ===== Page 394 =====

478

[MARGIN: compte tenu de (13)]
Revenant maintenant à la situation arithmétique (par passage au quotient via B)), on trouve un isomorphisme canonique

(45) $$ \mathcal{M}_{0,3}^{\prime \sim} \overset{\sim}{\longrightarrow} M_{\rho,\sigma}(\hat{\mathbb{Z}})/\Sigma $$
[unclear: crossed out text]
opérant sur $\operatorname{Sl}(2,\hat{\mathbb{Z}})$ via le facteur $T_{\sigma}$.

$\Sigma$ étant le sous-groupe $\simeq \{\pm 1\}^3$, produit de trois facteurs $\mathbb{Z}/2\mathbb{Z}$ inclus resp.t dans $T_{\rho}(\hat{\mathbb{Z}})$, $T_{\sigma}(\hat{\mathbb{Z}})$ et $\operatorname{Sl}(2,\hat{\mathbb{Z}})$.

Remq. Cet homomorphisme a été obtenu [unclear: crossed out text] ici, à l'aide de l'hom. de $\mathcal{M}_{0,3}^{\prime \sim}$ (ou ?)

(46) $$ \mathcal{M}_{0,3}^{\sim}(\infty) = \operatorname{Norm}_{\mathcal{M}_{0,3}^{\prime \sim}} (L_0) $$
[MARGIN: sous-groupe des autom. "unimodulaires"]
[MARGIN: (cf §46, V)]

\begin{tikzcd}
& \searrow \\
& S\mathcal{G}_{\rho,\sigma}^{**}(\hat{\mathbb{Z}})/\Sigma_0 \ar[r, hook] & S\mathcal{G}_{\rho,\sigma}^{*}(\hat{\mathbb{Z}})/\Sigma_0
\end{tikzcd}

et du scindage

$$ \hat{\Pi}_{0,3} \hookrightarrow \hat{\Gamma}_{0,3} \twoheadrightarrow \operatorname{Sl}(2,\hat{\mathbb{Z}})^{\sim} $$
[MARGIN: $\simeq \operatorname{Sl}(2,\mathbb{Z})^{\wedge}$]

en utilisant

(47) $$ \mathcal{M}_{0,3}^{\sim} \simeq \mathcal{M}_{0,3}^{\prime \sim} \cdot \hat{\Pi}_{0,3} $$

\newpage

%% ===== Page 395 =====

Mais on a ici un produit semi direct
$$(48) \quad \mathcal{M}_{0,3}^{\sim} \simeq \mathcal{M}_{0,3}^{! \sim}(L_0) \cdot \hat{\mathfrak{S}}_{0,3}^+$$
qui, combiné directement à
$$ \hat{\mathfrak{S}}_{0,3}^+ \longrightarrow \operatorname{Sl}(2, \hat{\mathbb{Z}})' $$
donne un hom.
$$(49) \quad \mathcal{M}_{0,3}^{\sim} \longrightarrow \left( M_{p,\sigma}(\hat{\mathbb{Z}}) \simeq SG_{p,\sigma}^*(\hat{\mathbb{Z}}) \cdot \operatorname{Sl}(2,\hat{\mathbb{Z}}) \right)_{/\Sigma}$$
qui prolonge l'hom (45) , défini d'abord
sur le sous-groupe distingué $\mathcal{M}_{0,3}^{! \sim}$
d'indice 6.

\newpage

%% ===== Page 396 =====

499

$\overline{\underline{X}}$ Étude des groupes $M_{p,r}$ et de $M_{p,r} \to \operatorname{Gl}(2)$

Par construction des produits semi-directs
$$M_{p,r}^\sim = \mathcal{G}_{p,r}^{* \mathrel{\hat{*}}} . G \simeq (T_0 \times T_p \times N'_0) . G$$
[MARGIN: opère dans $G$ via $T_0 \hookrightarrow G$]
(1)
$$M_{p,r} = S\mathcal{G}_{p,r}^{* \mathrel{\hat{*}}} . SG \hookrightarrow (T_p \times N'_0) . G$$
[MARGIN: opère dans $G$ via $N'_0 \xrightarrow{\text{det}} \hat{\mathbb{Z}}^\times \xrightarrow{\text{def}} T_0 \hookrightarrow G$]

on a un homomorphisme de groupes
(2)
$$ \begin{cases} M_{p,r}^\sim \xrightarrow{\Theta} G \\ \Theta( i_{T_0}(x) i_{T_p}(y) i_{N'_0}(z) i_G(f) ) = x.f \\ (x \in T_0, y \in T_p, z \in N'_0, f \in G). \end{cases} $$

L'homomorphisme composé
(3)
$$ S M_{p,r}^\sim \to M_{p,r}^\sim \xrightarrow{\Theta} G $$
$$ \simeq \mathcal{G}_{p,r}^{*\hat{*}}.G $$
$$ B_0 \times T_p \times N'_0 $$
[MARGIN: (composé noté encore $\Theta$)]
est décrit par
$$ \Theta( i_B(x) i_P(y) i_0'(z) i_G(f) ) = x.f $$
et $\Theta$ [unclear]
$$ i(\lambda) = i_B(\lambda) i_G(\lambda^{-1}) \quad (\lambda \in \Gamma L_0) $$
et passe donc au quotient en $\Theta : M_{p,r} \to G$.

\newpage

%% ===== Page 397 =====

Je m'intéresse surtout à l'hom. induit

(4) $$M_{\rho,\sigma} \overset{\Theta}{\longrightarrow} G \simeq \widehat{Gl}(2)_S$$

il est évident épi, car son image contient
trivialement $SG$, donc il suffit de voir
que le composé avec
$$G/SG \overset{\det}{\underset{\sim}{\longrightarrow}} \widehat{\mathbb{G}}_m$$
est épi, i.e. que le composé sus-dit l'est.
Or ce composé est aussi l'hom. composé

(5)
\begin{tikzcd}
M_{\rho,\sigma} \ar[r, "\text{projection}", "\text{canonique}"'] \ar[d, "\simeq"'] & S\mathcal{G}^{*4}_{\rho,\sigma(\varpi)} \ar[r, "*4 \text{ multiplication}"] & \widehat{\mathbb{G}}_m \\
S\mathcal{G}_{\rho,\sigma(\varpi)}^{*4}. SG & (\alpha, \beta, \mu, \varepsilon) \ar[r, mapsto] & \mu
\end{tikzcd}

qui est épi — en termes de

(6) $$S\mathcal{G}^{*4}_{\rho,\sigma(\varpi)} \simeq \text{[unclear: crossed out text]}$$
ce qui $(\underline{a}, \underline{b}) \mapsto \det \underline{b}$

en termes de

(7) $$S\mathcal{G}^{*4}_{\rho,\sigma(\varpi)} \simeq T_\rho \times N'_0$$
ce sera $(\underline{a}, \underline{b}) \mapsto \det \underline{b}$

(cf VIII, p. 189 en particulier (25)).

Définition

(8) $$M^+_{\rho,\sigma} \overset{df}{=} \operatorname{Ker}(M_{\rho,\sigma} \overset{\Phi}{\longrightarrow} \widehat{\mathbb{G}}_m) \quad ,$$

où $\Phi$ désigne l'hom. de surjection [unclear: considéré].

\newpage

%% ===== Page 398 =====

(9)
\begin{tikzcd}
1 \ar[r] & M_{\rho,\sigma}^+ \ar[r] \ar[d] & M_{\rho,\sigma} \ar[r] \ar[d, "\Theta"] & \mathbb{G}_m \ar[r] \ar[d, equal] & 1 \\
1 \ar[r] & SG \ar[r] & G \ar[r, "\det"] & \mathbb{G}_m \ar[r] & 1
\end{tikzcd}

qu'on rapprochera de §39 p.299 (1)

Comme $M_{\rho,\sigma}$ est plat de dim relative 6, et $G$ lisse de dim relative 4, le noyau de l'épi $\Theta : M_{\rho,\sigma} \to G$ est nécessairement plat de dim relative 2. Notons-le $I_{\rho,\sigma}$, il est contenu dans le noyau $J_{\rho,\sigma}$ (de dim relative 3) du composé
$$M_{\rho,\sigma} \to G \to PG \ ,$$
de sorte qu'on a
(10)
$$1 \to I_{\rho,\sigma} \to J_{\rho,\sigma} \to \mathbb{G}_m \to 1$$
et une suite de composition de
$M_{\rho,\sigma}$

(11)
\begin{tikzcd}
1 \ar[r, phantom, "\subset"] & I_{\rho,\sigma} \ar[r, phantom, "\subset"] & J_{\rho,\sigma} \ar[r, phantom, "\subset"] & M_{\rho,\sigma} \\
& & & PG \simeq SG/\mu_2 \\
& & G &
\end{tikzcd}

Le groupe $J_{\rho,\sigma}$ joue le rôle d'un radical - on prévoit qu'il est commutatif, [unclear: à isogénie près], (de dim relative 3) - et torique si...

\newpage

%% ===== Page 399 =====

$T_\rho^0$, $T_\sigma^0$ sont des tores.
Mais je vois que je me complique
l'existence pour être trop savant dans mes
déductions. En effet, par construction
du produit $\frac{1}{2}$ direct, on a :
$$M_{\rho,\sigma} = I\!I_{\rho,\sigma}^a . SG \subset (T_\sigma \times T_\rho \times N'_\sigma) . SG$$
$$\simeq (T_\rho \times N'_\sigma) \times_G (T_\sigma . SG)$$

ainsi

(12)
\begin{tikzcd}
M_{\rho,\sigma} \ar[r, hook] \ar[dr, "\Theta"'] & T_\rho \times N'_\sigma \times G \ar[d, "pr_3"] \\
& G
\end{tikzcd}

d'ailleurs $M_{\rho,\sigma}$ est défini dans ce
produit par les conditions

(13) $$M_{\rho,\sigma} \subset T_\rho \times N'_\sigma \times G$$
$$ \simeq \{ (x,y,f) \mid \det f = \det y = \varepsilon' \det x , \; {\varepsilon'}^2 = 1 \}$$

Cela montre que

(14) $$I\!I_{\rho,\sigma} \subset T_\rho \times N'_\sigma \times \mathbb{G}_m$$
$$ \simeq \{ (x,y,\lambda) \mid \det y = \varepsilon \det x = \lambda^2 , \; \varepsilon^2 = 1 \}$$

$$J_{\rho,\sigma} \subset T_\rho \times {N'_\sigma}^1$$
$$ \simeq \{ (x,y) \mid \det y = 1 , (\det x)^2 = 1 \}$$

Il me semble plus naturel de
faire le dévissage de $M_{\rho,\sigma}$ en commençant
par descendre à sa composante neutre $M_{\rho,\sigma}^0$,

\newpage

%% ===== Page 400 =====

[MARGIN: 501]

le groupe [unclear: quotient]
(15) $$\mathcal{M}_{p,\sigma} / \mathcal{M}_{p,\sigma}^\circ \simeq g_{p,\sigma}^*(0) / g_{p,\sigma}^*(0)^\circ = \mu_2 \times \mu_2$$

(cf p. 489 (23)), on aura donc

$$\mathcal{M}_{p,\sigma}^\circ = g_{p,\sigma}^{*\circ} \cdot SG \subset (T_\sigma \times T_p \times T_\sigma) \cdot SG$$
$$\simeq T_p \times T_\sigma \times (T_\sigma \cdot SG)$$
$$\simeq T_p \times T_\sigma \times G$$

(16) $$\mathcal{M}_{p,\sigma}^\circ \subset T_p \times T_\sigma \times G$$
$$= \{ (x, y, f) \mid \det x = \det y = \det f \}$$

[unclear: D'où on a]
$$\mathcal{J}_{p,\sigma}^\circ \overset{def}{=} \operatorname{Ker}(\mathcal{M}_{p,\sigma}^\circ \to PG)$$

(17)
$$\simeq \{ (x, y, \lambda) \mid \det x = \det y = \lambda^2 \}$$
$$\simeq T_p \times_{G_m} T_\sigma = \{ (x, y) \mid \det x = \det y \}$$
$$\simeq Sg_{p,\sigma}^*(0)^\circ$$

Enfin, on en déduit que
$$\mathcal{I}_{p,\sigma}^\circ \overset{def}{=} \operatorname{Ker}(\mathcal{M}_{p,\sigma}^\circ \to G)$$
(18)
$$\simeq ST_p \times ST_\sigma = \{ (x, y) \mid \det x = \det y = 1 \}$$

Récapitulons :

[MARGIN: 399]

\newpage

%% ===== Page 401 =====

Théorème. Le groupe
(19) $$M_{\rho,\sigma} \overset{def}{=} G_{\rho,\sigma}^*(0) \cdot SG$$
[MARGIN: $\uparrow$ cf définition p. 487' (30)]
est extension de $\mu_2 \times \mu_2$ par $M_{\rho,\sigma}^\circ$

(20) $$1 \to M_{\rho,\sigma}^\circ \to M_{\rho,\sigma} \to \mu_2 \times \mu_2 \to 1$$
déduit de $G_{\rho,\sigma}^*(0) \to \mu_2 \times \mu_2$
$(\alpha, \beta, \mu, \varepsilon) \mapsto (\det \alpha, \varepsilon \det \beta)$

Dans $M_{\rho,\sigma}^\circ$ on a
[unclear: remarquable : le sous-groupe central (de $M_{\rho,\sigma}^\circ$ + action)]
(21) $$ST_\rho \times ST_\sigma \hookrightarrow G_{\rho,\sigma}^\circ(0) \hookrightarrow M_{\rho,\sigma}^\circ \quad ,$$
et on a une suite exacte
(22)
\begin{tikzcd}
1 \ar[r] & ST_\rho \times ST_\sigma \ar[r] & M_{\rho,\sigma}^\circ \ar[r, "\Theta"] & G \ar[r] & 1 \\
& & SG \ar[u, hook] \ar[ur, "\simeq"'] & & 
\end{tikzcd}

Comme [unclear: dans $M_{\rho,\sigma}^\circ$ on a]
[unclear: $ST_\rho \times ST_\sigma \times SG \subset J_{\rho,\sigma}^\circ \times SG$]
[MARGIN: NB $J_{\rho,\sigma}^\circ$ est encore central dans $M_{\rho,\sigma}^\circ$ et $J_{\rho,\sigma}^\circ \cap SG = 1$]

(23) $$1 \to ST_\rho \times ST_\sigma \times SG \to M_{\rho,\sigma} \to \mathbb{G}_m \times \mu_2 \times \mu_2 \to 1$$
$x \mapsto (\chi(x), \varepsilon(x), \varepsilon'(x))$

(24) $$M_{\rho,\sigma}^{0+} \simeq ST_\rho \times ST_\sigma \times SG$$
provenant de la
section ($G \hookrightarrow M_{\rho,\sigma}^\circ$)
[unclear: $ST_\rho \times ST_\sigma \times G \to M_{\rho,\sigma}^\circ \to 1$]

On a un hom. canonique
(25) $$\mu_2 \times \mu_2 \times \mu_2 \hookrightarrow M_{\rho,\sigma}^{0+} \overset{def}{=} \operatorname{Ker}(M_{\rho,\sigma}^\circ \to \mathbb{G}_m)$$
$\simeq$ (24)
$$ST_\rho \times ST_\sigma \times SG$$

\newpage

%% ===== Page 402 =====

grâce aux inclusions canoniques
(26) $$ \mu_2 \hookrightarrow S T_{\underline{p}}, \mu_2 \hookrightarrow S T_{\underline{\sigma}}, \mu_2 \hookrightarrow S G $$
donnant naissance à l'hom. canonique
déjà envisagé par ailleurs
(27) $$ \Sigma = (\mathbb{Z}/2\mathbb{Z})^3 \hookrightarrow \Gamma \mathcal{M}_{\underline{p},\underline{\sigma}}^{0+} \subset \Gamma \mathcal{M}_{\underline{p},\underline{\sigma}}^* $$

Je vais introduire au sujet des groupes
$\mathcal{M}_{\underline{p},\underline{\sigma}}$ des notations qui paraphrasent les
notations $\mathcal{M}_{0,3}$ etc, dans le contexte des
schémas en groupes. Le sous-groupe
invariant $SG$ de $\mathcal{M}_{\underline{p},\underline{\sigma}}$ sera ainsi noté
$\mathfrak{S}_{\underline{p},\underline{\sigma}}^{0+}$ (l'analogue de $\widehat{\mathfrak{S}}_{0,3}^+$), le

[MARGIN: pour un peu j'aurais préféré le rayer ! Celui (28) est tombé du ciel. même dans la section svie de $\underline{\text{III}}$, c'est $\Gamma_{\underline{p},\underline{\sigma}}$ (et non $\Gamma_{\underline{p},\underline{\sigma}}^{0+}$ ou $\Gamma \mathcal{M}_{\underline{p},\underline{\sigma}}^{0+}$ !) qu'il conviendrait de prendre [unclear] ! (29)]

$$ \mathcal{M}_{\underline{p},\underline{\sigma}} / \mathfrak{S}_{\underline{p},\underline{\sigma}}^{0+} \simeq \mathcal{T}_{\underline{p},\underline{\sigma}}^! $$
($\mathcal{M}_{0,3}^!$, on a d'ailleurs
$$ \mathcal{T}_{\underline{p},\underline{\sigma}}^! \simeq \mathcal{G}_{\underline{p},\underline{\sigma}(0)}^* \simeq T_0 \times T_p \times \mathcal{N}_0^1 $$
- par la suite la notation plus simple et
plus suggestive $\mathcal{T}_{\underline{p},\underline{\sigma}}$ sera utilisée de
préférence à $\mathcal{G}_{\underline{p},\underline{\sigma}(0)}^*$. On a donc

(30) $$ 1 \longrightarrow \underset{\simeq SG \simeq \operatorname{Sl}(2)}{\mathfrak{S}_{\underline{p},\underline{\sigma}}^{0+}} \longrightarrow \mathcal{M}_{\underline{p},\underline{\sigma}} \longrightarrow \mathcal{T}_{\underline{p},\underline{\sigma}}^! \longrightarrow 1 $$

et un scindage canonique en produit semi-direct
(31) $$ \mathcal{M}_{\underline{p},\underline{\sigma}} \simeq \mathcal{M}_{\underline{p},\underline{\sigma}(0)}^* \cdot \mathfrak{S}_{\underline{p},\underline{\sigma}}^{0+} $$

\newpage

%% ===== Page 403 =====

(32)
\begin{tikzcd}
M_{\rho,\sigma}(\infty) \ar[r, hook] \ar[dr, "\sim"'] & M_{\rho,\sigma} \ar[d, "can"] \\
& \tau_{\rho,\sigma}^!
\end{tikzcd}

est un sous-groupe isomorphe à $\tau_{\rho,\sigma}^!$
qui correspond au sous-groupe $M_{0,3}^{!\sim}(\infty)$ dans
$M_{0,3}^\sim$ ($\S$ V).

Rem. On aurait envie d'introduire une
$\tau_{\rho,\sigma}$ extension (un produit) de $\tau_{\rho,\sigma}^!$ par un groupe
tel que $\mathfrak{S}_3$ (qui serait noté $\tau_{\rho,\sigma}^+$), en
retenant un certain sous-groupe $\Pi_{\rho,\sigma}$ de
$\mathfrak{S}_{\rho,\sigma}^+$ tel que $\mathfrak{S}_{\rho,\sigma}^+ / \Pi_{\rho,\sigma} \simeq \mathfrak{S}_3$ -
mais ce n'est pas possible dans le contexte
des schémas en groupes (ni dans le cas
particulier modulaire), vu que $\mathfrak{S}_{\rho,\sigma}^+ \simeq \operatorname{Sl}(2)_{\mathbb{Z}}$
est [unclear: connexe] sur $S$. Par contre on
peut, dans le "cas modulaire" $\rho = \begin{pmatrix} 0 & 1 \\ -1 & 0 \end{pmatrix}, \sigma = \begin{pmatrix} 0 & 1 \\ -1 & 1 \end{pmatrix}$
introduire une section canonique

(33) $$t = \bar{\tau}_\infty \in \Gamma(\tau_{\rho,\sigma}^!)$$

dont le représentant $\tau_\infty$ dans $\Gamma M_{\rho,\sigma} (\infty)$ s'explicite
comme

(34) $$\tau_\infty = \begin{pmatrix} -1 & 0 \\ 0 & +1 \end{pmatrix}$$

[unclear: Il serait intéressant de voir si, en dehors des]

\newpage

%% ===== Page 404 =====

cas modulaire, on peut introduire des "anti-invo-
lutions" qui généralisent les réflexions géométriques
$$
(35) \qquad \left\{ \begin{array}{l}
\Gamma_0 = \tau'_0 = \begin{pmatrix} 0 & 1 \\ 1 & 0 \end{pmatrix} \\
\Gamma_1 = -\tau'_0 = \begin{pmatrix} 0 & -1 \\ -1 & 0 \end{pmatrix} \\
\Gamma_\infty = \tau_\infty = \begin{pmatrix} -1 & 0 \\ 0 & +1 \end{pmatrix}
\end{array} \right. \qquad (\text{cf. p. } 335, \text{ \S } 42, \text{ IV})
$$

Je suis convaincu que oui. D'ailleurs n'avions-
nous pas mis en évidence dans $M_{p,\sigma}$ un
sous-groupe $\mu_2 \times \mu_2 \times \mu_2$, correspondant d'une part à
trois sections remarquables (correspondant respect. aux
trois sections $-1$ des trois facteurs $\mu_2$) ?
Mais je me mélange là les pinceaux - ces
sections étaient des sections de $M_{p,\sigma}^\circ$ - tandis
que les sections $\tau_\infty, \tau'_0, \dots$ de multiplicateur
[unclear: que fait la réflexion $Z_{p,\sigma} \simeq \mu_2 / \text{Centr}(M_{p,\sigma})$]
$-1$, donc correspondaient à $\varepsilon = \varepsilon' = -1$, dont corres-
pondait à la section "la moins triviale" $\underline{(-1,-1)}$ de
$M_{p,\sigma} / M_{p,\sigma}^\circ = \mu_2 \times \mu_2$.

$$
(36) \qquad \mathcal{T}_{p,\sigma}' \simeq \mathcal{G}_{p,\sigma}(0) \simeq \mathcal{T}_0 \times \mathcal{T}_p \times \mathcal{N}_\sigma'
$$
$$
\simeq \{ (\underline{u}, \underline{a}, \underline{b}) \mid \det \underline{u} = \det \underline{b} = \varepsilon' \det \underline{a} , \varepsilon'^2 = 1 \}
$$

et rappelons les homomorphismes
$$ \delta_\sigma : u \mapsto \det u / \mathcal{T}_0 \to \mathbb{G}_m $$
et des hom. canoniques "positifs"
$$ \mathbb{G}_m \overset{j_p}{\longrightarrow} \mathcal{T}_p $$
$$ \mathbb{G}_m \overset{j_\sigma}{\longrightarrow} \mathcal{T}_0 $$
dont le composé avec $\delta_p, \delta_\sigma$ (les hom. cités) est $\lambda \mapsto \lambda^2$.
Ainsi dans la deuxième [unclear] de (36), il y a [unclear]

\newpage

%% ===== Page 405 =====

sous-groupe de Cartan trivial [unclear: crossed out text]
$$ \mathbb{G}_m \times \mathbb{G}_m \times \mathbb{G}_m \hookrightarrow T_0 \times T_p \times \mathcal{N}'_0 $$

Faisons les calculs pour déterminer les invariants $\alpha$, $\beta$ ... correspondants à $\tau_0, \tau_1, \tau_\infty$.

i) $\tau_\infty = \tau_\infty^* = \begin{pmatrix} -1 & 0 \\ 0 & 1 \end{pmatrix} \in \operatorname{Gl}(2, \mathbb{Z}) \subset \operatorname{Gl}(2, \hat{\mathbb{Z}})^{\pm 1}$

On a $\tau_\infty \in \mathcal{T}_0(\mathbb{Z})$, c'est le seul élément [unclear: crossed out text] non trivial de $\mathcal{T}_0(\mathbb{Z}) \simeq \mathbb{G}_m(\mathbb{Z}) \simeq \{\pm 1\}$.
L'automorphisme extérieur qu'il définit dans $\Pi_{0,S}$ ou dans $\operatorname{Sl}(2, \mathbb{Z})$ a une caractéristique [unclear: familière] -
il stabilise le sous-groupe $L_0$ engendré par $\varepsilon_0$.
Ecr. $L_0$ [unclear: symétrique] ($\varepsilon_0^2 = l_0$)
$\tau_\infty(\varepsilon_0) = \varepsilon_0^{-1}$, $\tau_\infty(l_0) = l_0^{-1}$
où $\tau_\infty \in \mathcal{M}_{0,S}^{\sim}[0]$

etc. - je [unclear] il est dû de la forme $u_{\alpha, \beta}$, [unclear]
des calculs déjà faits donnent
$\alpha = \tau_\infty$, $\gamma = \alpha \beta \alpha^{-1} = -l_1^{-1} = \begin{pmatrix} 1 & 0 \\ -2 & 1 \end{pmatrix}$
(attention aux signes - quand on fait les calculs dans $\operatorname{Sl}(2, \mathbb{Z})$ et pas dans $\operatorname{Sl}(2, \mathbb{Z})^\sim$!)

$\beta = 1$, $\nu = l_1 = 1$,
$\varepsilon = -1$, $\varepsilon = \varepsilon_2(l_1) = -1$, $\varepsilon' = \varepsilon_3(l_1) = -1$

on a ici $B=0$, $D=-1$, donc
$$u = \begin{pmatrix} 1/D & B/D \\ 0 & 1 \end{pmatrix} = \begin{pmatrix} -1 & 0 \\ 0 & 1 \end{pmatrix} = \tau_\infty$$

$\underline{a} = \alpha^{-1} u = \tau_\infty^{-1} \tau_\infty = 1$, $\underline{b} = \beta^{-1} u = \tau_\infty$
(élément [unclear: antécédent] en $\mathcal{T}_0(\mathbb{Z}) \simeq _2\mathcal{T}_0(\mathbb{Z})$)

donc
(36) $\Gamma_\infty = \tau_\infty = (\underline{u}_\infty = \tau_\infty, \underline{a}_\infty = 1, \underline{b}_\infty = \tau_\infty) \in T_0 \times T_p \times \mathcal{N}'_0$

NB Soit plus généralement
$$u = \begin{pmatrix} \mu & \lambda \\ 0 & \nu \end{pmatrix} \in \operatorname{Gl}(2, \mathbb{Z}) \subset \operatorname{Gl}(2, \hat{\mathbb{Z}})$$
et considérons l'automorphisme extérieur de $\operatorname{Sl}(2, \mathbb{Z})^\sim$ (et $\mathcal{M}_{0,S}^\sim$) qu'il induit, dont la restriction à $B_{0,s}^+$ opère,
quels en sont les invariants $\alpha, \beta$ etc.? On pourra dévisser tout ce [unclear] ...

\newpage

%% ===== Page 406 =====

569

mieux prendre $f=1$, et on aura
$\alpha = u$, d'où $\underline{a = 1}$,
d'autre part, notons que $\mu \in \{\pm 1\} (= (-1)^K, K \in \mathbb{Z}/2\mathbb{Z})$,
on aura $\varepsilon = \varepsilon' = \mu$, et si $\mu = 1$ ($ie.$ $u \in S\ell(2,\mathbb{Z})$) on aura $\beta = u$
[unclear] donc $\underline{b = 1}$, sinon $\beta = \tau_\infty u \tau_\infty$ donc $\underline{b = \tau_\infty}$,
dans tous les cas on aura $\beta = u \tau_\infty^K$, $b = \tau_\infty^K$
en résumé

(37) $\underline{u} = (\underbrace{u}_{\in \Gamma_{B_0}}, \underbrace{a=1}_{\in \Gamma_p}, \underbrace{b=\tau_\infty^K}_{\in \Gamma_{N_0}}) \in \hat{M}_{p,\sigma \{0\}}^*(\hat{\mathbb{Z}}) = \Gamma_{B_0} \times \Gamma_p \times \Gamma_{N_0}$

[unclear struck out text]
Mais ce n'est un élément [unclear] pour [unclear] $\Gamma_{p, \sigma \{0\}}$,
que si $\lambda=0$. [unclear]

2°) $\tau_{\hat{1}} = -\tau'_0 = \begin{pmatrix} -1 & 0 \\ 0 & 1 \end{pmatrix}$, c'est justement un cas
qui tombe sous la remarque précédente, on

[MARGIN: NB réduir-t-on [unclear] Loc $\hat{G}_{\bar{p}, \sigma}^+$ [unclear] $\Gamma_i$ ont [unclear] fournir $(1, \tau_\infty)$]

a donc (38) $\tau_{\hat{1}} = \tau'_1 = (\underline{u}_{\hat{1}} = -\tau'_0 = \begin{pmatrix} -1 & 0 \\ 0 & 1 \end{pmatrix}, \underline{a=1}, \underline{b=\tau_\infty}) \in M_{p, \sigma \{0\}}(\mathbb{Z})$

3°) $\tau_{\hat{0}} \tau'_{\infty} = \begin{pmatrix} 0 & 1 \\ 1 & 0 \end{pmatrix}$, cette fois-ci c'est un
automorphisme qui [unclear] pas [unclear] pour $L_0$
ça va mieux, n'est ce pas que $\tau_0, \tau_1, \tau_\infty$ engen-
-drent $G\ell(2,\mathbb{Z})$ ( on a "seulement" $S\ell(2,\mathbb{Z})$ ) !
On a $\tau'_{0} = \tau_{\infty} \tau_{\alpha}$, donc son image dans
$\hat{M}_{p, \sigma \{0\}}(\hat{\mathbb{Z}})$ est le produit des
images de $\tau_{\infty}$ et $\tau_{\alpha}$ - la première a été calculée
dans 1°), la deuxième est tautologique
puisque $S G_{p,\sigma} \subset \hat{G}_{p,\sigma}^+$ et $\tau_{\alpha} \in S G(\mathbb{Z})$. On
n'a qu'à [unclear] $\tau_{\hat{0}} = \tau_\infty$ et on a $\dots$ [unclear]

\newpage

%% ===== Page 407 =====

dans $\mathcal{T}_{p,\sigma}(\mathbb{Z})$ par $\tau_{\infty} = \tau_{\infty}$ et $\tau_1 = -\tau_0'$, c'est
toujours le = élément

(39) $$(-1, 1, \tau_{\infty}) \in (\underline{T}_0 \times \underline{T}_p \times N'_{\sigma})(\mathbb{Z})$$

Calculons d'ailleurs
$$\underline{T}_0 \times \underline{T}_p \times N'_{\sigma}(\mathbb{Z}) = T_0(\mathbb{Z}) \times T_p(\mathbb{Z}) \times N'_{\sigma}(\mathbb{Z})$$

On a

(40)
$$T_0 \simeq \mathbb{G}_m, T_0(\mathbb{Z}) \simeq \{\pm 1\}$$
$$T_p(\mathbb{Z}) \simeq (\mathbb{Z} + \mathbb{Z} \rho)^* \simeq \mathbb{Z}/6\mathbb{Z} \quad \text{groupe engendré par } \rho$$
$$N'_{\sigma}(\mathbb{Z}) \simeq \{1, \tau_{\infty}\} \cdot \underline{T}_{\sigma}(\mathbb{Z}) \quad (\text{produit semi-direct})$$
$$T_{\sigma}(\mathbb{Z}) \simeq (\mathbb{Z} + \mathbb{Z} \sigma)^* \simeq \mathbb{Z}/4\mathbb{Z} \quad \text{groupe engendré par } \sigma$$
$$N'_{\sigma}(\mathbb{Z}) \simeq D_4 \quad \text{groupe diédral d'ordre } 4$$
$$(\text{"symétries du carré"})$$

d'où

(41) $$(T_0 \times T_p \times N'_{\sigma})(\mathbb{Z}) \simeq (\underline{\mathbb{Z}/2\mathbb{Z}}) \cdot (\mathbb{Z}/2\mathbb{Z} \times \mathbb{Z}/6\mathbb{Z} \times \mathbb{Z}/4\mathbb{Z})$$
groupe engendré par $\tau_{\infty}$
$\tau_{\infty}$ opère trivialement sur les deux premiers facteurs et par symétrie sur le dernier
$$\simeq \mathbb{Z}/2\mathbb{Z} \times \mathbb{Z}/6\mathbb{Z} \times \underbrace{(\mathbb{Z}/2\mathbb{Z} \cdot \mathbb{Z}/4\mathbb{Z})}_{D_4}$$

[MARGIN: NB la condition det $a = \varepsilon' \det u$, $\varepsilon''=1$ est automatique car $\mathbb{G}_m(\mathbb{Z}) = \{\pm 1\}$]

(42) $$S\mathcal{G}_{p,\sigma}(\mathbb{Z}) = \mathcal{T}_{p,\sigma}(\mathbb{Z}) = \{ (u, a, b) \in (T_0 \times T_p \times N'_{\sigma})(\mathbb{Z}) \mid \det u = \det b \}$$
$$\simeq (T_p \times N'_{\sigma})(\mathbb{Z}) \simeq \mathbb{Z}/6\mathbb{Z} \times D_4$$

et de façon précise
$$D_4 \to M_{p,\sigma}[0](\mathbb{Z})$$
[unclear]

\newpage

%% ===== Page 408 =====

505

l'élément recherché, i.e. qui correspond
à $\tau_0 = \tau_\infty, \tau_1 = -\tau_0', \tau_\infty = \tau_\infty$ [barré : correspond-t-il à]
l'élément $(1, \tau_\infty)$ de $\mathbb{Z}/6\mathbb{Z} \times (D_4 \simeq \{1,\tau_\infty\} \cdot \mathbb{Z}/4\mathbb{Z})$.

Ceci nous donne d'ailleurs une expression
pour $\underline{M}_{p,\sigma}(\mathbb{Z})$ :

(43)
$$ \underline{M}_{p,\sigma}(\mathbb{Z}) \simeq \underline{\mathcal{T}}_{p,\sigma}(\mathbb{Z}) \cdot S\mathcal{L}(2,\mathbb{Z}) $$
$$ (\text{produit semi-direct}) $$
$$ \simeq (\mathbb{Z}/6\mathbb{Z} \times D_4) \cdot S\mathcal{L}(2,\mathbb{Z}) $$

où $\mathbb{Z}/6\mathbb{Z} \times D_4$ opère sur $S\mathcal{L}(2,\mathbb{Z})$ via l'hom.
caractère canonique composé

(44) $$ \mathbb{Z}/6\mathbb{Z} \times D_4 \xrightarrow{\text{projection canonique}} D_4 \xrightarrow{\text{hom. can.}} \{1, \tau_\infty\} $$

$\tau_\infty$ opère sur $S\mathcal{L}(2,\mathbb{Z})$ par conjugaison intérieure avec
la matrice
$$ \tau_\infty = \begin{pmatrix} -1 & 0 \\ 0 & 1 \end{pmatrix} $$

On note que l'épim. canonique
$$ \underline{M}_{p,\sigma} \to \mu_2 \times \mu_2 $$
se factorise par l'hom. $\underline{\mathcal{T}}_{p,\sigma}$, et induit des
hom. sur les points entiers

(45) $$ \underline{M}_{p,\sigma}(\mathbb{Z}) \to \underline{\mathcal{T}}_{p,\sigma}(\mathbb{Z}) \simeq \mathbb{Z}/6\mathbb{Z} \times D_4 \to (\mu_2 \times \mu_2)(\mathbb{Z}) \simeq \{\pm 1\} \times \{\pm 1\} $$

Cet hom. est nul sur $\mathbb{Z}/6\mathbb{Z} \times \mathbb{Z}/4\mathbb{Z}$ (avec $D_4^+$ noté sous $\mathbb{Z}/4\mathbb{Z}$)
(correspondant à des déterminants égaux à $1$), et
sur $\tau_\infty$ prend la valeur $(-1, -1)$. Les fibres
de ce [unclear] est lisses sur les pts entiers [unclear: de la]

\newpage

%% ===== Page 409 =====

pas surjectif.

[MARGIN: faux dès qu'il s'agit d'un produit $1/2$ direct par le centre d'un groupe, qui y opère trivialement] Centre de $M_{g,\sigma} = M_{g,\sigma(0)} . Sl$. C'est le produit semi-direct du sous-groupe centre de $M_{g,\sigma(0)}$ (qui opère trivialement par automorphismes sur $Sl$) par le centre de $Sl$. Ce dernier est bien connu, c'est $\mu_2$, et $M_{g,\sigma}$ y opère trivialement. D'autre part le sous-groupe de $M_{g,\sigma(0)} \simeq \mathcal{T}_{g,\sigma} = B'_0 \times T_{\rho} \times N'_0$ qui opère trivialement sur $Sl_2$, est l'intersection de $\mathcal{T}_{g,\sigma}$ avec $T_{\rho} \times N'_0$, car $B'_0$ opère fidèlement, tandis que
$$ (46) \qquad B'_0 \cap \text{centre de } Sl_2 = \{1\} $$
C'est donc le sous-groupe des
$$ (1, \underline{a}, \underline{b}) \mid (\det \underline{a})^2 = 1, \det \underline{b} = 1 $$
$\in T_{\rho} \ \in N'_0$

[MARGIN: NB! Cela donne [unclear] éléments du centre de $\mathcal{T}_{g,\sigma}$] Reste : déterminer le centre de $\mathcal{T}_{g,\sigma}$, extension de $\mu_2 \times \mu_2$ par $\mathcal{T}_{g,\sigma}^\circ \simeq g_{g,\sigma}^\circ(0)$
$\subset T'_0 \times T_{\rho} \times N'_0$. Comme $\mathcal{T}_{g,\sigma}^\circ$ est commutatif, il suffit de déterminer l'action de $\mu_2 \times \mu_2$ sur $\mathcal{T}_{g,\sigma}^\circ$, et d'en déduire les invariants.

\newpage

%% ===== Page 410 =====

505

Notons que l'on a
$$ \mathcal{T}_{\rho,\sigma} = S G_{\rho,\sigma}^*(0) \subset G_{\rho,\sigma}^*(0) \simeq T_0 \times T_\rho \times N'_{0'} $$
et un hom. de suites exactes
$$ ((u,a,b) \mapsto (\varepsilon(b), \det\underline{a}/\det(\underline{u}))) $$
(48)
\begin{tikzcd}
1 \ar[r] & \mathcal{T}_{\rho,\sigma} \ar[r] \ar[d] & G_{\rho,\sigma}^*(0) \ar[r] \ar[d] & \mu_2 \times \mu_2 \ar[r] \ar[d, "\text{première} \atop \text{proj.}"] & 1 \\
1 \ar[r] & T_0 \times T_\rho \times \bar{N}_0 \ar[r] & T_0 \times T_\rho \times N'_{0'} \ar[r] & \mu_2 \ar[r] & 1
\end{tikzcd}

qui montre que l'action de $\mu_2 \times \mu_2$ sur $\mathcal{T}_{\rho,\sigma}$ se fait par l'intermédiaire de la projection de $\mu_2 \times \mu_2$ sur son premier facteur, et par l'opération de celui-ci sur $T_0 \times T_\rho \times \bar{T}_0$ qui opère trivialement sur $T_0 \times T_\rho$, et qui sur $\bar{T}_0$ est l'opération canonique de $\mu_2$ sur $\bar{T}_0$, $\sigma \mapsto \varepsilon \sigma$, dont les invariants se réduisent aux scalaires $\mathbb{G}_m$. Donc ce raisonnement montre, plus généralement, que la trace du centre de l'ext. $G^*_{\rho,\sigma}$ sur le sous-groupe $\mathcal{T}'_{\rho,\sigma}$ noyau de 
$$ \mathcal{T}_{\rho,\sigma} \rightarrow \mu_2 \times \mu_2 \xrightarrow{1^{\text{ère}} \text{ proj.}} \mu_2 $$
$$ (u,a,b) \mapsto \varepsilon(b) $$
est aussi la trace, sur ce sous-groupe, de $T_0 \times T_\rho \times \mathbb{G}_m$. D'ailleurs, comme $\mu_2$ opère fidèlement sur $\bar{T}_0$ via $\sigma \mapsto \varepsilon \sigma$, on voit que l'on constate qu'il opère aussi fidèlement sur $\mathcal{T}_{\rho,\sigma}$ ([unclear: qui en a en effet s'envoie épimorphiquement] sur $\bar{T}_0$) donc le
$$ \text{Centre}(\mathcal{T}_{\rho,\sigma}) \subset \mathcal{T}'_{\rho,\sigma} = \operatorname{Ker}(\mathcal{T}_{\rho,\sigma} \xrightarrow{((u,a,b) \mapsto \varepsilon(b))} \mu_2) $$

\newpage

%% ===== Page 411 =====

Donc on résume :

(49) 
$$ Centre(\mathbb{G}_{\rho, \sigma}^*(\underline{o})) = \mathbb{T}_0 \times \mathbb{T}_\rho \times \mathbb{G}_m $$
$$ Centre(C_{\rho, \sigma}) = [unclear: (\mathbb{T}_0 \times \mathbb{T}_\rho \times \mathbb{G}_m)] \cap C_{\rho, \sigma} $$
$$ = \{ (\underline{u}, \underline{a}, \lambda) \in \mathbb{T}_0 \times \mathbb{T}_\rho \times \mathbb{G}_m \mid u = \lambda^2, \underline{a} = \varepsilon^{\text{det } \underline{a}}, \varepsilon^{12} = 1 \} $$

Ce sous-groupe du centre, qui opère trivialement sur $S \mathbb{G}$ est donc, en vertu de (47), donné par la restriction $u=1$, il s'identifie à

(50)
$$ 1 \times S \mathbb{T}_\rho \times \mu_2 \subset C_{\rho, \sigma} \subset \mathbb{T}_0 \times \mathbb{T}_\rho \times N_0' $$
$$ \text{où } S \mathbb{T}_\rho \subset \mathbb{T}_\rho $$
$$ = \{ \underline{a} \mid (\text{det } \underline{a})^2 = 1 \} $$

Donc

[MARGIN: Faux. Le centre est plus grand]
(51) $$ Centr. (\mathcal{M}_{\rho, \sigma}) \simeq (1 \times S \mathbb{T}_\rho \times \mu_2) \times \mu_2 $$

Le groupe de points entiers de ce groupe, dans le cas modulaire, est réduit à :

(52) $$ \mathbb{Z}/6\mathbb{Z} \times \mathbb{Z}/2\mathbb{Z} \times \mathbb{Z}/2\mathbb{Z} $$

[MARGIN: NB Ce groupe est en fait contenu dans $\mathcal{M}_{\rho, \sigma}(\mathbb{Z})$.]
$$ \simeq \mathbb{Z}/3\mathbb{Z} \times (\mathbb{Z}/2\mathbb{Z} \times \mathbb{Z}/2\mathbb{Z} \times \mathbb{Z}/2\mathbb{Z}) $$

produit d'un groupe cyclique d'ordre 3 (engendré visiblement par $-\rho$) par un groupe $(\mathbb{Z}/2\mathbb{Z})^3$, qui n'est autre que le groupe $\Sigma$ déjà décrit dans (I).

Plus [unclear: instructif] serait de [unclear: calculer] le

\newpage

%% ===== Page 412 =====

507

centre de $\mathcal{M}_{p,\sigma^\circ}$, qui sera donc le produit
direct du centre de $SG$ (sur lequel le groupe
$\overline{\mathcal{M}_{p,\sigma^\circ}}(0) \simeq \tau_{p,\sigma^\circ}$ opère trivialement) par le ss-
groupe de $\tau_{p,\sigma^\circ} \subset T_0 \times T_p \times T_\sigma$ qui opère
trivialement sur $Sl_1$, savoir 
[MARGIN: Faux]
$$ \tau_{p,\sigma^\circ} \cap 1 \times T_p \times T_\sigma = \{ (1, \underline{a}, \underline{b}) \mid \det \underline{b} = 1, (\det \underline{a})^2 = 1 \} $$
$$ \simeq ST_p \times ST_\sigma $$
(cet [unclear] $(\det \underline{a})^2=1$, car $\underline{a}$ dans $\mathcal{T}_p^\circ \subset \mathcal{T}_{p,\sigma^\circ}^{(0)}$)

Donc
[MARGIN: Faux, le centre est de dim 3]
(53) $Centre (\mathcal{M}_{p,\sigma^\circ}) = (1 \times ST_p \times ST_\sigma) \times \mu_2$

d'où dans le cas modulaire :
(54) $Center (\mathcal{M}_{p,\sigma^\circ})(\mathbb{Z}) \simeq 1 \times \mathbb{Z}/6\mathbb{Z} \times \mathbb{Z}/4\mathbb{Z} \times \mathbb{Z}/2\mathbb{Z}$

qui n'est que deux fois plus gros que
le groupe des pts entiers du centre
de $\mathcal{M}_{p,\sigma^\circ}$. [unclear]

Rmq. La composante neutre du centre de
$\mathcal{M}_{p,\sigma^\circ}$ est bien sûr $ST_p \times ST_\sigma$, comme
le suggérait la discussion du th. p. 501.

\newpage

%% ===== Page 413 =====

508

## $\underline{\text{XI}} \quad$ Récapitulation sur le groupe $\mathcal{M}_{p,\tilde{\sigma}}$.

[MARGIN: Préliminaires]

1°) Etant donnés, sur un schéma $S$, une forme $G$ de $\operatorname{Gl}(2)$, d'où des formes $\mathcal{S}l_G$, $P_G$ de $\mathcal{S}l(2)$, $P(1)$, et deux sections $\tilde{\rho}$, $\tilde{\sigma}$ de $P_G$, et un sous-groupe à 1-paramètre additif fermé $\tilde{\xi}_0 : \mathbb{G}_a \hookrightarrow P_G$ (ce qui revient à une section nilpotente régulière de Lie de $G$, ou une section nilpotente régulière $N_0$ dans la forme $M$ de $M_2(\mathcal{O}_S)$ associée à $G$) - satisfaisant les conditions

(1) $\tilde{\sigma}\tilde{\rho} = \tilde{\xi}_0(1)$, $\tilde{\xi}_0$ et $\tilde{\xi}_1 = \tilde{\rho}(\tilde{\xi}_0) \quad (= \tilde{\sigma}(\tilde{\xi}_0))$ sont indépendantes sur chaque fibre

(2) $Tr \, \sigma = 0$ où $\sigma$ est un relèvement local de $\tilde{\sigma}$

On définit des relèvements canoniques $\rho$, $\sigma$ de $\tilde{\rho}$, $\tilde{\sigma}$ à $G$, satisfaisant les conditions

(3) $\begin{cases} \xi_0 = - \sigma^{-1} \rho & \text{i.e.} \quad \xi_1 = - \rho \sigma^{-1} \\ Tr(\rho + \sigma) = 1 & \text{i.e.} \quad Tr \rho = 1 \end{cases} ,$

et les constructions qui suivent ont un sens.

On construit un épinglage

[crossed out: (4) (5) $Epingl_S(\operatorname{Gl}(2) \to G)$]

par la condition que $\xi_1$ prenne la forme

(6) $\xi_1 = \begin{pmatrix} 1 & 0 \\ 1 & 1 \end{pmatrix}$

d'où

(7) $\begin{cases} \rho = \begin{pmatrix} 0 & 1 \\ -v & 1 \end{pmatrix} , \quad \sigma = \begin{pmatrix} 0 & -1 \\ v & 0 \end{pmatrix} \quad (v \in \Gamma(\mathbb{G}_m)) \\ \det \rho = \det \sigma = -v \end{cases}$

(8) $\xi_0 = \begin{pmatrix} 1 & 1/v \\ 0 & 1 \end{pmatrix}$

\newpage

%% ===== Page 414 =====

on considère les sous-groupes

(9)
$$ \begin{cases} L_0 = \{ \varepsilon_0(\lambda) \} = \left\{ \begin{pmatrix} 1 & \lambda \\ 0 & 1 \end{pmatrix} \right\} \qquad (\lambda \in \Gamma \mathbb{G}_a) \\ L_1 = \sigma(L_0) = \rho(L_0) = \{ \varepsilon_1(\lambda) \} = \left\{ \begin{pmatrix} 1 & 0 \\ \lambda & 1 \end{pmatrix} \right\} \\ B'_0 = \left\{ \begin{pmatrix} \mu & \lambda \\ 0 & 1 \end{pmatrix} \mid \mu \in \Gamma \mathbb{G}_m, \lambda \in \Gamma \mathbb{G}_a \right\} \\ T_0 = \left\{ \begin{pmatrix} \mu & 0 \\ 0 & 1 \end{pmatrix} \mid \mu \in \Gamma \mathbb{G}_m \right\} \end{cases} $$

d'où

(10)
$$ B'_0 = T_0 \cdot L_0 \qquad (\text{produit } 1/2\text{-direct}) $$

et on a des iso.

(11)
$$ \begin{cases} \xi_0 : \mathbb{G}_a \simeq L_0 \qquad \lambda \mapsto \varepsilon_0(\lambda) = \begin{pmatrix} 1 & \lambda \\ 0 & 1 \end{pmatrix} \\ \xi_1 : \mathbb{G}_a \simeq L_1 \qquad \lambda \mapsto \varepsilon_1(\lambda) = \begin{pmatrix} 1 & 0 \\ \lambda & 1 \end{pmatrix} \end{cases} $$
$$ \text{avec } \varepsilon_1(\lambda) = \rho(\varepsilon_0(\lambda)) = \sigma^-(\varepsilon_0(\lambda)) $$

On [unclear: pose] considère

(12)
$$ \begin{cases} T_{\rho} = \text{schém. des unités de } \mathcal{A}_{\rho} = \underline{\mathcal{O}}_S + \underline{\mathcal{O}}_S \rho \simeq \underline{\mathcal{O}}_S[X]/(X^2-X-\nu) \\ T_{\sigma} = \text{schém. des unités de } \mathcal{A}_{\sigma} = \underline{\mathcal{O}}_S + \underline{\mathcal{O}}_S \sigma \simeq \underline{\mathcal{O}}_S[Y]/(Y^2-\nu) \end{cases} $$

(13)
$$ N'_{\sigma} = \{ g \in G \mid g(\sigma) = \varepsilon \sigma, \varepsilon^2 = 1 \} $$
$$ \cap $$
$$ G $$

d'où

(14)
$$ 1 \to T_{\sigma} \to N'_{\sigma} \xrightarrow{\varepsilon_{\sigma}} \mu_2 \to 1 $$

On pose

(15)
$$ \begin{cases} \mathcal{G}_{\rho,\sigma} = B'_0 \times T_{\rho} \times T_{\sigma} \\ \mathcal{G}^*_{\rho,\sigma} = B'_0 \times T_{\rho} \times N'_{\sigma} \end{cases} $$

\newpage

%% ===== Page 415 =====

509

[MARGIN: 
$M_{\rho,\sigma}[0]$, $M_{\rho,\sigma}^{(0)}$
$\mathcal{G}_{\rho,\sigma}$, $\tau_{\rho,\sigma}$
$M_{\rho,\sigma}$, $\mathcal{G}_{\rho,\sigma}^*$
$\tau_{\rho,\sigma}^*$
]

2° On introduit

(16) $$M_{\rho,\sigma}[0] \subset \mathcal{G}_{\rho,\sigma}^* = B'_0 \times T_\rho \times N'_\sigma$$
$$\| dfn$$
$$\{ (u, a, b) \mid \det u = \det b = \varepsilon' \det a , \varepsilon'^2 = 1 \}$$

(17) $M_{\rho,\sigma}(0) \subset M_{\rho,\sigma}[0]$ défini par $u=1$ i.e.
$$M_{\rho,\sigma}(0) = M_{\rho,\sigma}[0] \cap (T_0 \times T_\rho \times N'_\sigma)$$
dans $\mathcal{G}_{\rho,\sigma}^*$

(18) $$L_0 \hookrightarrow M_{\rho,\sigma}[0]$$
$$\lambda \longmapsto (\lambda, 1, 1)$$

donc on a
(19) $$M_{\rho,\sigma}[0] \simeq M_{\rho,\sigma}(0) . L_0 \quad \text{produit } 1/2 \text{ direct}$$

On a un homomorphisme
(20) $$M_{\rho,\sigma}[0] \longrightarrow \mathbb{G}_m \times \mu_2 \times \mu_2$$
noté $g \mapsto \chi(g), \varepsilon_1(g), \varepsilon_2(g)$
pour $g=(u,a,b)$ on a
$\chi(g) = \det u = \det b$
$\varepsilon_1(g) = \varepsilon(b)$ (cf (13)(14))
$\varepsilon_2(g) = \det(a)/\det b$ i.e. $\det a = \varepsilon_2 \mu$
(où $\mu = \chi(g)$)

trivial sur $L_0$ (donc se factorise par le quotient). On pose ainsi

[MARGIN: Cette notation n'est pas très adéquate par un [unclear: choix] précis à la suite de ce qui précède, le [unclear: portrait] [unclear: tiré] d'un brouillon combinant les notations $\Gamma_{\rho,\sigma}$ et $\Gamma_{0,s}$]

(21) $$\tau_{\rho,\sigma} \simeq M_{\rho,\sigma}[0] / L_0 \simeq$$

composé
$$M_{\rho,\sigma}(0) \hookrightarrow M_{\rho,\sigma}[0] \longrightarrow \tau_{\rho,\sigma}$$

est un iso, donc on a

(23) $$M_{\rho,\sigma}(0) \simeq \tau_{\rho,\sigma} \hookrightarrow T_0 \times T_\rho \times N'_\sigma$$

\newpage

%% ===== Page 416 =====

La composition avec la projection can. dans $T_{\rho} \times N'_{\sigma}$ donne un mono

[MARGIN: car $\det$ sur $T_0 \xrightarrow{\sim} \mathbb{G}_m$]
(24) $$ \mathcal{T}_{\rho, \sigma} \hookrightarrow T_{\rho} \times N'_{\sigma} $$
$$ \simeq \left\{ (\underline{a}, \underline{b}) \mid (\det \underline{a} / \det \underline{b})^2 = 1 \right\} $$
$$ \text{i.e. } \det \underline{a} = \varepsilon' \det \underline{b} \text{ avec } {\varepsilon'}^2=1 $$

et on aura dans

(25) $$ \begin{cases} \chi(\underline{a}, \underline{b}) = \det \underline{b} \\ \varepsilon_2(\underline{a}, \underline{b}) = \varepsilon_{\sigma}(\underline{b}) \\ \varepsilon_3(\underline{a}, \underline{b}) = \det(\underline{a}) / \det(\underline{b}) = \det \underline{b} / \det \underline{a} \end{cases} $$

On définit

(26) $$ \underset{\simeq T_{\rho} \times_{\mathbb{G}_m} T_{\sigma}}{\mathcal{T}_{\rho, \sigma}^{\circ}} = \text{Ker } \Big( \mathcal{T}_{\rho, \sigma} \xrightarrow{(\varepsilon_2, \varepsilon_3)} \mu_2 \times \mu_2 \Big) $$

de sorte on a

(27) $$ 1 \to \mathcal{T}_{\rho, \sigma}^{\circ} \to \mathcal{T}_{\rho, \sigma} \to \mu_2 \times \mu_2 \to 1 $$

et une suite exacte canonique

(28) $$ 1 \to S T_{\rho} \times S T_{\sigma} \to \mathcal{T}_{\rho, \sigma}^{\circ} \xrightarrow{\chi} \mathbb{G}_m \to 1 $$

où on définit

(29) $$ \begin{cases} S T_{\rho} = \{ g \in T_{\rho} \mid \det g = 1 \} \\ S T_{\sigma} = \{ g \in T_{\sigma} \mid \det g = 1 \} \end{cases} $$

[MARGIN: L'opération de $\mathcal{T}_{\rho, \sigma} \simeq \mathcal{M}_{\rho, \sigma}(0)$ sur $L_0$ est donnée par $g(\lambda, \mu) = \varepsilon_3(g)(\lambda, \mu)$, où $g \in \mathcal{T}_{\rho, \sigma}$]

On fait opérer $\mathcal{M}_{\rho, \sigma}(0) (\simeq \mathcal{T}_{\rho, \sigma})$ sur $SG$, par projection

(30)
\begin{tikzcd}
\mathcal{M}_{\rho, \sigma}(0) \ar[r, "\text{projection}"] \ar[d, hook] & T_0 \ar[r, "\text{can.}"', "\text{(mono)}"] \ar[d, hook] & P G \ar[r, "\sim"] \ar[d, equal] & \operatorname{Aut}_{gr}(SG) \\
T_0 \times T_\rho \times N'_{\sigma} & G & G / \mathbb{G}_m \simeq \operatorname{Aut}(G) &
\end{tikzcd}

\newpage

%% ===== Page 417 =====

510

et on définit

$$ (32) \qquad \mathcal{M}_{p,\sigma} \simeq \mathcal{M}_{p,\sigma}(0) \cdot SG $$
produit $1/2$ direct

On a donc un hom. can.

$$ (33) \qquad SG \hookrightarrow \mathcal{M}_{p,\sigma} $$

dont l'image est notée $\mathcal{G}_{p,\sigma}^+$, d'où

$$ (34) \qquad 1 \longrightarrow \mathcal{G}_{p,\sigma}^+ \longrightarrow \mathcal{M}_{p,\sigma} \longrightarrow \mathcal{T}_{p,\sigma} \longrightarrow 1 $$

$$ (35) \qquad \mathcal{M}_{p,\sigma} \simeq \mathcal{M}_{p,\sigma}(0) \cdot \mathcal{G}_{p,\sigma}^+ $$
(où $\mathcal{M}_{p,\sigma}(0) \simeq \mathcal{T}_{p,\sigma}^0$, produit $1/2$ direct).

NB. On utilise la notation $\mathcal{T}_{p,\sigma}$ pour le groupe quotient

$$ (36) \qquad \mathcal{T}_{p,\sigma} = \mathcal{M}_{p,\sigma} / \mathcal{G}_{p,\sigma}^+ \ , $$

la notation $\mathcal{M}_{p,\sigma}(0)$ pour le sous-groupe de $\mathcal{M}_{p,\sigma}$.

On définit un plongement canonique

$$ (37) \qquad \mathcal{M}_{p,\sigma}[0] \hookrightarrow \mathcal{M}_{p,\sigma} $$

qui prolonge l'identité du sous-groupe $\mathcal{M}_{p,\sigma}(0)$ commun de $\mathcal{M}_{p,\sigma}[0]$ et de $\mathcal{M}_{p,\sigma}$ , et qui sur $L_0$ est donné par

$$ (38) \qquad i(\lambda) = j(\lambda) \qquad (cf \ (33)) $$

(ce qui est bien compatible avec les opérations de $\mathcal{M}_{p,\sigma}(0)$ sur $L_0$ d'une part, sur $SG \simeq \mathcal{G}_{p,\sigma}^+$ de l'autre).

\newpage

%% ===== Page 418 =====

Une autre façon de le dire, c'est de constater
que $\mathcal{M}_{\rho,\sigma}(0)$ opère sur $SG \simeq \widehat{\mathcal{S}}_{\rho,\sigma}^+$
normalise le sous-groupe $j(L_0)$ - que nous
identifions désormais à $L_0$), de sorte qu'on
aura un sous-groupe
$$\mathcal{M}_{\rho,\sigma}(0) . L_0 \subset \mathcal{M}_{\rho,\sigma}$$
qui est canoniquement isom. à $\mathcal{M}_{\rho,\sigma}[0]$
défini plus haut, grâce à l'écriture (35) de
celui-ci comme produit $1/2$-direct.

Remarque.
$$\mathcal{M}_{\rho,\sigma}[0] = \operatorname{Norm}_{\mathcal{M}_{\rho,\sigma}}(L_0)$$

Par construction, on a $\subset$, l'inclusion
en sens inverse reviendrait à dire que $L_0$
est son propre normalisateur dans $SG \simeq \widehat{\mathcal{S}}_{\rho,\sigma}^+$
ce qui bien sûr n'est pas vrai du tout.

Vérifier que le fait analogue est vrai
dans le groupe profini $\widehat{\mathcal{S}}_{0,\vec{s}}^+$, pour le sous-
groupe engendré par $\tau_0$, ou dans $T_{0,\vec{s}}^+$ pour
celui engendré par $L_0$ (curieuse vérification
d'ailleurs que je n'ai pas pu faire).

\newpage

%% ===== Page 419 =====

511

[MARGIN: $\mathcal{M}_{\rho,\sigma}^\circ, \mathcal{M}_{\rho,\sigma}^{+0}$, $\tau_{\rho,\sigma}^{+0}, \tau_{\rho,\sigma}^\circ \dots$ quelques sous-groupes [unclear] calqués de ... au dessus...]

On désigne par $\mathcal{M}_{\rho,\sigma}^\circ$, $\mathcal{M}_{\rho,\sigma}[0]^\circ$, $\mathcal{M}_{\rho,\sigma}(0)^\circ$
composantes neutres, on aura p.ex

(39)
$$1 \to \mathcal{M}_{\rho,\sigma}^\circ \to \mathcal{M}_{\rho,\sigma} \to \mu_s \times \mu_2 \to 1$$
$$\simeq \mathcal{M}_{\rho,\sigma}(0)^\circ . \widehat{\mathfrak{S}}_{\rho,\sigma}^+$$
produit $1/2$ direct

(40) $\mathcal{M}_{\rho,\sigma}(0)^\circ \simeq \mathcal{T}_{\rho,\sigma}^\circ \simeq T_\rho \times_{\mathbb{G}_m} T_\sigma$

(41) $\mathcal{M}_{\rho,\sigma}^\circ$ est l'image inverse de $\mathcal{T}_{\rho,\sigma}^\circ$
par $\mathcal{M}_{\rho,\sigma} \xrightarrow{can} \mathcal{T}_{\rho,\sigma}$ ,

on a un hom canonique surj.
(42) $\mathcal{M}_{\rho,\sigma} \to \mathcal{T}_{\rho,\sigma} \xrightarrow{\chi} \mathbb{G}_m$
également noté $\chi$, son noyau est noté
$\mathcal{M}_{\rho,\sigma}^+$, d'où extensions $\mathcal{M}_{\rho,\sigma}^{0+}$, etc

(43) 
$$1 \to \mathcal{M}_{\rho,\sigma}^+ \to \mathcal{M}_{\rho,\sigma} \xrightarrow{\chi} \mathbb{G}_m \to 1$$
$$1 \to \mathcal{M}_{\rho,\sigma}^{0+} \to \mathcal{M}_{\rho,\sigma}^\circ \xrightarrow{\chi} \mathbb{G}_m \to 1$$

(44) 
$$\mathcal{M}_{\rho,\sigma}^+ \simeq \underline{\mathcal{M}_{\rho,\sigma}^+(0)} . \widehat{\mathfrak{S}}_{\rho,\sigma}^+$$
$$\simeq \mathcal{T}_{\rho,\sigma}^+$$
$$\mathcal{M}_{\rho,\sigma}^{0+} \simeq \underline{\mathcal{M}_{\rho,\sigma}^{0+}(0)} . \widehat{\mathfrak{S}}_{\rho,\sigma}^+$$
$$\simeq \mathcal{T}_{\rho,\sigma}^{0+} \simeq ST_\rho \times ST_\sigma$$
cf (28)

Donc on a

(45) $\mathcal{M}_{\rho,\sigma}^{0+}(0) \stackrel{dfn}{=} \operatorname{Ker}(\mathcal{M}_{\rho,\sigma}^\circ(0) \xrightarrow{\chi, \varepsilon_1, \varepsilon_2} \mathbb{G}_m \times \mu_s \times \mu_2)$
$\simeq ST_\rho \times ST_\sigma$

et ce sous-groupe est central dans $\mathcal{M}_{\rho,\sigma}^\circ$
(puis dans $\mathcal{M}_{\rho,\sigma}$), de façon précise

\newpage

%% ===== Page 420 =====

[MARGIN: Toute la fin du § devrait venir après la "variante" de la page 426 "variante diédrale" à la fin du §, d'où des répétitions, avec égalités, cf plus bas...]

on a
$$ (45) \quad \mathcal{M}_{\rho, \sigma}^{0+} \overset{dfn}{=} \operatorname{Ker}( \mathcal{M}_{\rho, \sigma}^{0} \xrightarrow{(\chi, \varepsilon_1, \varepsilon_2)} \mathbb{G}_m \times \mu_2 \times \mu_2 ) $$
$$ \simeq ST_\rho \times ST_\sigma \times \widehat{\mathfrak{S}}_{\rho, \sigma}^+ $$
(avec $\widehat{\mathfrak{S}}_{\rho, \sigma}^+ \subset \mathcal{M}_{\rho, \sigma}(0)$)

$$ (47) \quad \text{Centr}(\mathcal{M}_{\rho, \sigma}^0) \supset 1 \times ST_\rho \times ST_\sigma \times \mu_2 $$
(où $\mu_2 \subset \mathcal{M}_{\rho, \sigma}(0)^\circ$, et $\widehat{\mathfrak{S}}_{\rho, \sigma}^+ = SG$)

D'où un sous-groupe
$$ (48) \quad \mu_2 \times \mu_2 \times \mu_2 \subset ST_\rho \times ST_\sigma \times \mu_2 $$
(ce dernier contenu dans $\text{Centr}(\mathcal{M}_{\rho, \sigma}^0)$)

qui est en fait contenu dans le Centre du $\mathcal{M}_{\rho, \sigma}$ tout entier, car
$$ (49) \quad \text{Centr}(\mathcal{M}_{\rho, \sigma}) \supset 1 \times S'T_\rho \times \mu_2 \times \mu_2 $$
(où $\mu_2 \times \mu_2 \subset \mathcal{M}_{\rho, \sigma}(0)$, et $\widehat{\mathfrak{S}}_{\rho, \sigma}^+ = SG$)

avec
$$ (50) \quad S'T_\rho \overset{dfn}{=} \{ g \in \Gamma T_\rho \mid (\text{det } g)^2 = 1 \} $$
de sorte qu'on a
$$ (51) \quad 1 \to ST_\rho \to S'T_\rho \xrightarrow{\varepsilon_3} \mu_2 \to 1 $$

On a donc
$$ (52) \quad \text{Centr}(\mathcal{M}_{\rho, \sigma}) \cap \mathcal{M}_{\rho, \sigma}^0 \supset 1 \times ST_\rho \times \mu_2 \times \mu_2 $$
(les facteurs correspondant à $\mathcal{T}_\rho$, $\mathcal{T}_\rho$, $\mathcal{T}_\sigma$, $\widehat{\mathfrak{S}}_{\rho, \sigma}^+$ respectivement)

Le "cas modulaire" est celui où
$$ (53) \quad \begin{cases} \text{det } \rho = 1 \quad \text{i.e.} \quad \text{det } \sigma = 1 \quad \text{i.e.} \quad v = -1 \\ \text{i.e.} \quad \rho = \begin{pmatrix} 0 & 1 \\ -1 & 1 \end{pmatrix} \quad \text{i.e.} \quad \sigma = \begin{pmatrix} 0 & 1 \\ -1 & 0 \end{pmatrix} \end{cases} $$

(cf VI pour d'autres caractérisations, telles
$\sigma^2 = -1$, ou $\rho^2 - \rho + 1 = 0$, ou $\tilde{\sigma}^2 = \tilde{\rho}^3 = 1$ ... ) ,
cela signifie aussi qui plus est que l'image de $\mathcal{T}_\rho$ est dans $ST_\rho$ donc dans le centre de $\mathcal{M}_{\rho, \sigma}^0$.

\newpage

%% ===== Page 421 =====

512

[unclear: vu] que $\sigma$ [unclear: est dans] l'image de $T_\sigma$ et de
$ST_\sigma$ i.e. dans le centre de $M_{\rho,\sigma}^\circ$. Donc dans ce
cas on trouve un homomorphisme

(54)
$$ \left\{ \begin{array}{ccc} \mathbb{Z}/6\mathbb{Z} \times \mathbb{Z}/4\mathbb{Z} & \to & \underline{Centr}(M_{\rho,\sigma}^\circ) \\ (i,j) & \mapsto & i_\rho(\rho^i) i_\sigma(\sigma^j) \end{array} \right. $$

(où on définit $i_\rho, i_\sigma$ par les trois morphismes canoniques

(55)
$i_\rho : S' T_\rho \to M_{\rho,\sigma}^\circ(0) \subset M_{\rho,\sigma}$
$i_\sigma : S' T_\sigma \to M_{\rho,\sigma}^\circ(0) \subset M_{\rho,\sigma}$
$i_G : SG \simeq \hat{\mathbb{G}}_{\rho,\sigma}^+ \to M_{\rho,\sigma} \quad , \quad )$

dont l'image pour $2$ inversible sur $S$, contient
contient l'image de $\mu_2 \times \mu_2 \times 1$ dans (48).
Pour $2$ inv. sur $S$, l'hom. (54) est un mono
(mais pas en car. $2$, bien sûr). Sur le bon
absolu $\mathbb{Z}$, on a par ailleurs, dans le cas
modulaire

(56)
$$ \left\{ \begin{array}{l} \mathcal{T}_{\rho,\sigma}(\mathbb{Z}) \simeq T_\rho(\mathbb{Z}) \times N_\sigma(\mathbb{Z}) \simeq \mathbb{Z}/6\mathbb{Z} \times D_4 \\ M_{\rho,\sigma}(\mathbb{Z}) \simeq \mathcal{T}_{\rho,\sigma}(\mathbb{Z}) . \operatorname{Sl}(2,\mathbb{Z}) \end{array} \right. $$
produit $\frac{1}{2}$ direct
$\mathcal{T}_{\rho,\sigma} \simeq \mathbb{Z}/6\mathbb{Z} \times D_4$ opérant
via le caractère $D_4 \xrightarrow{pr_2} \mathbb{Z}/2\mathbb{Z} \simeq \{\pm 1\}$
$\mathcal{T}_{\rho,\sigma}$ opérant sur $\operatorname{Sl}(2,\mathbb{Z})$ par
autom. induit par la matrice
$\tau_\infty = \begin{pmatrix} -1 & 0 \\ 0 & 1 \end{pmatrix}$

\newpage

%% ===== Page 422 =====

[MARGIN: Les homs $\Theta, \Theta^\circ$]
[MARGIN: 40]

On définit un hom.
(57) $$M_{\rho,\sigma} \overset{\Theta}{\longrightarrow} G$$
$\Theta | \underline{G}_{\rho,\sigma}^+ \simeq SG$ est l'identité
$\Theta | M_{\rho,\sigma}(0) \simeq T_0 \times \Gamma_\rho \times N_\sigma'$ donné par $\Theta(\underline{u}, \underline{a}, \underline{b}) = \underline{u}$

d'où
(58) $\Theta | M_{\rho,\sigma}[0] \simeq B_0' \times \Gamma_\rho \times N_\sigma'$ donné par $\Theta(\underline{u}, \underline{a}, \underline{b}) = \underline{u}$

On a
(59)
\begin{tikzcd}
M_{\rho,\sigma}(0) \ar[r, equal] & T_0 \times \Gamma_\rho \times N_\sigma' \\
\operatorname{Ker} \Theta \ar[u, hook] \ar[r, "\subset"] & 1 \times \Gamma_\rho \times N_\sigma' \ar[u, hook, "cart."'] \\
& \{ (1, \underline{a}, \underline{b}) \mid \det \underline{b} = 1, (\det \underline{a})^2 = 1 \} \ar[u, "\simeq"'] \\
& \Gamma T_\rho \times \Gamma N_\sigma' \ar[u, "\subset"'] \\
& S T_\rho \times S N_\sigma' \ar[u, "\simeq"']
\end{tikzcd}
où $SN_\sigma' \overset{d\acute{e}f}{=} \operatorname{Ker}(N_\sigma' \to G_m)$

Soit
(60) $\Theta^\circ = \Theta | M_{\rho,\sigma}^\circ$ , on trouve

(61)
\begin{tikzcd}
1 \ar[r] & \operatorname{Ker} \Theta^\circ \ar[r] \ar[d, "\simeq"'] & \operatorname{Ker} \Theta \ar[r, "(\varepsilon_1{,}\varepsilon_2)"] & \mu_2 \times \mu_2 \ar[r] & 1 \\
& S T_\rho \times S T_\sigma \ar[d, "\simeq"'] & & & \\
& \operatorname{Rad}(M_{\rho,\sigma}^\circ) \ar[d, "\simeq"'] & & & \\
& \operatorname{Centr}(M_{\rho,\sigma}^\circ) & & &
\end{tikzcd}

Posons
(62) $SR_{\rho,\sigma}^\circ = S T_\rho \times S T_\sigma = \operatorname{Ker} \Theta^\circ = (\operatorname{Ker} \Theta) \cap M_{\rho,\sigma}^\circ$
qui est un sous-groupe central de $M_{\rho,\sigma}^\circ$ (il est

\newpage

%% ===== Page 423 =====

distingué [crossed out: invariant] mais pas central dans $M_{\rho,\sigma}$ tout entier)

le groupe
(63) $$M_{\rho,\sigma} / S R_{\rho,\sigma}^0 = M_{\rho,\sigma}^\natural$$
est donc une extension de $\mu_2 \times \mu_2$ par
$M_{\rho,\sigma}^0 / S R_{\rho,\sigma}^0 \underset{\text{déf}}{=} M_{\rho,\sigma}^{0\natural}$, isomorphe à $G$ par
$\Theta^0$ (puisque $S R_{\rho,\sigma}^0 = \text{Ker } \Theta^0$), d'où
(64) $$1 \to G \to M_{\rho,\sigma}^\natural \to \mu_2 \times \mu_2 \to 1,$$
d'ailleurs l'hom.
$$\Theta : M_{\rho,\sigma} \to G$$
étant trivial sur $S R_{\rho,\sigma}^0$ passe au quotient en
(65) $$\Theta^\natural : M_{\rho,\sigma}^\natural \to G,$$
dont la restriction à $G$ (via (64)) est l'
identité, donc on trouve, via (64)
(66) $$M_{\rho,\sigma}^\natural \simeq G \times \mu_2 \times \mu_2.$$
On trouve donc
(67)
\begin{tikzcd}
1 \ar[r] & S R_{\rho,\sigma}^0 \ar[r] \ar[d, "\simeq"'] & M_{\rho,\sigma} \ar[r, "((\Theta), \varepsilon_3, \varepsilon_2)"] & M_{\rho,\sigma}^\natural \ar[r] \ar[d, "\simeq"] & 1 \\
& S T_\rho \times S T_\sigma & & G \times \mu_2 \times \mu_2 &
\end{tikzcd}

Pour voir comment $M_{\rho,\sigma}^\natural \simeq G \times \mu_2 \times \mu_2$
opère sur $S T_\rho \times S T_\sigma$, on [unclear: note que le s-groupe]
$M_{\rho,\sigma}^0$ [crossed out: est l'] image inverse de $G \simeq M_{\rho,\sigma}^{0\natural}$
dans l'extension (67), or $S R_{\rho,\sigma}^0$ est central
dans $M_{\rho,\sigma}^0$, donc $G$ y opère trivialement.

\newpage

%% ===== Page 424 =====

D'autre part on a
$$
\begin{array}{ccc}
SR^\circ_{\rho,\sigma} & \subset & M_{\rho,\sigma(0)} \xrightarrow[(\varepsilon_3,\varepsilon_2)]{\text{épi}} \mu_2 \times \mu_2 \\
\wr\parallel & \cap & \\
S\Gamma_\rho \times S\Gamma_\sigma & & \Gamma_0 \times \Gamma_\rho \times \mathcal{N}_\sigma
\end{array}
$$
et on trouve que $\mu_2 \times \mu_2$ opère sur $SR_{\rho,\sigma}$
via son deuxième facteur (correspondant
à $\varepsilon_2(u, \underline{v}, \underline{\underline{k}}) = \varepsilon_\sigma(\underline{\underline{k}})$), un $\varepsilon$ de celui-ci
($\varepsilon^2=1$) opère sur $SR_{\rho,\sigma} \simeq S\Gamma_\rho \times S\Gamma_\sigma$ en laissant
intouché le premier facteur, et en opérant sur
$\Gamma_\sigma$ donc sur $S\Gamma_\sigma$ via
$$ (68) \quad \sigma \longrightarrow \varepsilon \sigma \quad \text{(compatible avec l'équation de Hamilton-Cayley } \sigma^2 - v = 0 \text{ de } \sigma) $$

[unclear crossed out text]

[MARGIN: Il y aurait intérêt peut-être à l'étudier en soi et le relier aux (102) plus bas de $M_{\rho,\sigma}$, en se bornant au $y^{(6)}$.]

Soit
(69) $R^\circ_{\rho,\sigma}$ l'image inverse de $G_m \subset G$
dans $M_{\rho,\sigma}$ par $\Theta : M_{\rho,\sigma} \to G$

donc une structure d'extension
sur $R^\circ_{\rho,\sigma}$:
$$ (70) \quad 1 \longrightarrow SR^\circ_{\rho,\sigma} \longrightarrow R^\circ_{\rho,\sigma} \longrightarrow G_m \longrightarrow 1 $$

On aura
$$ R^\circ_{\rho,\sigma} = \left\{ (\underset{\in \Gamma_0}{u}, \underset{\in \Gamma_\rho}{\underline{v}}, \underset{\in \Gamma_\sigma}{\underline{\underline{k}}}), f \ \left| \begin{array}{l} \det u = \det \underline{v} = \det \underline{\underline{k}}, \ u, f \in \Gamma_m \\ \text{i.e. } f = \lambda y^\alpha \\ \text{où } \lambda \in \Gamma_{unclear}, \ \lambda^2 = \det u \end{array} \right. \right\} $$

\newpage

%% ===== Page 425 =====

514

$$ (71) \quad R_{\rho,\sigma}^0 \simeq \left\{ (\mu, \underset{\in \Gamma_{T_\rho}}{\underline{a}}, \underset{\in \Gamma_{T_\sigma}}{\underline{b}}) \mid \mu \in \mathbb{G}_m, \text{det } a = \text{det } b = \mu^{-2} \right\} $$
[MARGIN: NB. le $\text{det}$ est $1$ car $\text{det } j_{\mathbb{G}}(\mu) = \mu^{-2}$]

où
$$ (72) \quad \begin{cases} j_{+}(\mu) = \begin{pmatrix} \mu & 0 \\ 0 & 1 \end{pmatrix} \in S\mathbb{G} \quad (\mu \in \mathbb{G}_m) \\ j_{T_0} : \mathbb{G}_m \to T_0 \quad \text{défini plus haut} \\ j_{\mathbb{G}} : \mathbb{G}_m \to \mathbb{G} \quad \text{défini plus haut} \end{cases} $$

On trouve ainsi que la composition
$$ R_{\rho,\sigma}^0 \hookrightarrow M_{\rho,\sigma}^0 \twoheadrightarrow \tau_{\rho,\sigma}^0 \simeq T_\rho \times_{\mathbb{G}_m} T_\sigma $$
$$ M_{\rho,\sigma}^0 / S P_{\rho,\sigma}^0 $$

$$ (73) \quad \begin{cases} \text{fait de } R_{\rho,\sigma}^0 \text{ un [unclear: revêtement] de groupes} \\ \text{de } \tau_{\rho,\sigma}^0 \end{cases} $$

et plus précisément qu'on a un carré cartésien

(74)
\begin{tikzcd}
R_{\rho,\sigma}^0 \ar[r, "(73)"] \ar[d, "\text{l'hom. can.}"', "\text{de (70)}"] & T_\rho \times_{\mathbb{G}_m} T_\sigma \simeq \tau_{\rho,\sigma}^0 \ar[d, "\text{can}"] \\
\mathbb{G}_m \ar[r, "\mu \mapsto \mu^{-2}"'] & \mathbb{G}_m
\end{tikzcd}

(75) Proposition : $R_{\rho,\sigma}^0 = \text{Centre} (M_{\rho,\sigma}^0)$

Évidemment le centre de $M_{\rho,\sigma}^0$ est contenu dans $R_{\rho,\sigma}^0$, il s'agit inversement de [unclear: justifier] des assertions...

\newpage

%% ===== Page 426 =====

de $M_{\rho,\sigma}^\natural = M_{\rho,\sigma} / \mathfrak{S}_{\rho,\sigma} \simeq G$, il reste à
prouver que $R_{\rho,\sigma}$ est central dans
$M_{\rho,\sigma}^0$. Mais [unclear: crossed out words] $M_{\rho,\sigma}^0$ est
engendré par $R_{\rho,\sigma}^0$ et par $\mathfrak{S}_{\rho,\sigma}^+ \simeq SG$,
puisque $R_{\rho,\sigma}^0 \to M_{\rho,\sigma}^0 / \mathfrak{S}_{\rho,\sigma}^+$
est épi (73)(74), et comme $R_{\rho,\sigma}$ est
commutatif (74), il reste à montrer
qu'il commute à $\mathfrak{S}_{\rho,\sigma}^+$ i.e. opère
trivialement sur $\mathfrak{S}_{\rho,\sigma}^+$. Or un
élément de $R_{\rho,\sigma}^0$ est de la forme
$(j_{T_0}(\mu^2), a, \underline{b}) . ( \lambda j_{T_0}(\mu^{-2}) )$
le premier facteur opère comme
$$ int \ j_G(\mu^2) $$
le deuxième comme
$$ int \ j_G(\mu^{-2}) $$
donc leur produit opère trivialement, cqfd.

Notons maintenant que
$$ (75) \quad R_{\rho,\sigma}^0 \cap \underset{\overset{\|}{SG}}{\mathfrak{S}_{\rho,\sigma}^+} = \operatorname{Ker}(R_{\rho,\sigma}^0 \to T_{\rho,\sigma}^0) = \underset{\overset{|}{cf (74)}}{\mu_2} $$
(on peut dire aussi plus simplement, que

\newpage

%% ===== Page 427 =====

$R_{p,\sigma}^\circ \cap G_{p,\sigma}^+$ , étant contenu dans le centre de $G_{p,\sigma}^+ \simeq SG$ qui est $\mu_2$ , est $\subset \mu_2$ - mais cela ne montre pas l'égalité.

Rmq. Cela donne une description plus simple de $M_{p,\sigma}^\circ$ que comme produit $\frac{1}{2}$ direct, de la façon suivante. On définit $R_{p,\sigma}^\circ$ par le diagramme (74), d'où une suite exacte

$$(77) \quad 1 \to \mu_2 \xrightarrow{j'} R_{p,\sigma}^\circ \longrightarrow T_p \times_{G_m} T_\sigma \to 1$$

et on définit $M_{p,\sigma}^\circ$ comme somme amalgamée centrale

$$(78) \quad 1 \to \mu_2 \longrightarrow R_{p,\sigma}^\circ \times \underset{SG}{\overset{\parallel}{G_{p,\sigma}^+}} \longrightarrow M_{p,\sigma}^\circ \to 1$$
$$ \varepsilon \mapsto (j'(\varepsilon), j''(\varepsilon^{-1})) $$

(où $j'$ est donné dans (77) et $j'' : \mu_2 \xrightarrow{\sim} SG = G_{p,\sigma}^+$ l'isomorphisme, [unclear: description bien connue]).

Notons qu'on a de même pour $G$ [unclear: en termes] de $SG$

$$(79) \quad 1 \to \mu_2 \longrightarrow G_m \times SG \longrightarrow G \to 1$$
$$ \varepsilon \mapsto (\varepsilon, j''(\varepsilon^{-1})) $$

et on a un hom. de (78) dans (79), grâce à l'hom. can.

$$(80) \quad R_{p,\sigma}^\circ \longrightarrow G_m \quad \text{de (74)}.$$

\newpage

%% ===== Page 428 =====

ce qui par passage aux quotients dans 
$$R_{\rho,\sigma}^o \times SG \xrightarrow{\tilde{\Theta}^o} G_{\sim} \times SG$$
fournit un isomorphisme
$$\Theta^o : M_{\rho,\sigma}^o \longrightarrow G \quad ,$$
tout comme le groupe $M_{\rho,\sigma}^o$ lui-même, [unclear: s'exprime]
en termes des groupes ($\underline{pro}$-toriques) "en
l'air" $T_\rho, T_\sigma$ et des épimorphismes
$$T_\rho \xrightarrow{\delta_\rho} G_{\sim} \quad , \quad T_\sigma \xrightarrow{\delta_\sigma} G_{\sim} \quad ,$$
et de la forme $G$ de $\operatorname{Gl}(2)_{\hat{\mathbb{Z}}}$ bien [unclear: sûr]
sans avoir à utiliser [unclear: ni] des plongements
$$T_\rho \hookrightarrow G, T_\sigma \hookrightarrow G \quad ,$$ et encore moins
des sections $s_\rho, s_\sigma$ de $G...$). On
trouve $C_{\rho,\sigma}^+$ comme l'image de $SG$ dans $M_{\rho,\sigma}^o$,
[crossed out text]

$$T_{\rho,\sigma}^o \overset{df}{=} M_{\rho,\sigma}^o / C_{\rho,\sigma}^+ \simeq R_{\rho,\sigma}^o / \mu_2$$
$$\simeq T_\rho \times_{G_{\sim}} T_\sigma$$

Mais pour notre propos, le
"dévissage" de l'extension ainsi construite
de $T_{\rho,\sigma}^o$ par $C_{\rho,\sigma}^+ \simeq SG$, correspondant
au ss-groupe section $M_{\rho,\sigma}^o(0)$, jouera un
rôle essentiel - en comparant à cette section.

\newpage

%% ===== Page 429 =====

515

$$ \tau_{\rho, \sigma}^\circ \simeq T_\rho \times_{G} T_\sigma \longrightarrow M_{\rho, \sigma}^\circ $$

avec

$$ M_{\rho, \sigma}^\circ \longrightarrow G $$

on trouve un homomorphisme qui se factorise

ainsi :

$$ \tau_{\rho, \sigma}^\circ = T_\rho \times_{G_m} T_\sigma \xrightarrow{\text{épi, can}} G_m \hookrightarrow G $$

dont l'image est le tore $T_0$ ...

Remarque On peut bien sûr construire $M_{\rho, \sigma}$ de façon analogue, en introduisant l'image inverse - notons le $R_{\rho, \sigma}$ - de $G_{ar} \times fl_x \times fl_x$ par $M_{\rho, \sigma} \rightarrow G_{ar} \times fl \times fl$ dans le [unclear: raisonnement] avec $\widetilde{C}_{\rho, \sigma}^+$ évidemment égal à $fl_x$ , puisque

$$ R_{\rho, \sigma}^\circ = R_{\rho, \sigma} \cap M_{\rho, \sigma}^\circ = (\dots) $$

[unclear: (compris nul)]

(il vaudrait peut-être mieux noter $R_{\rho, \sigma}$ plutôt $R_{\rho, \sigma}^\circ$ , et l'ancien $R_{\rho, \sigma}$ deviendrait $R_{\rho, \sigma}'$ , de $S P_{\rho, \sigma}^\circ$ deviendrait $S R_{\rho, \sigma}'$ ... )

Je ne vais pas expliciter ici la construction de $R_{\rho, \sigma}^\circ$ en termes de $T_\rho, N_{\dots}$ etc, n'étant pas sûr qu'il y en aura besoin pour la suite.

[MARGIN: NB $R_{\rho, \sigma}$ est aussi, plus simplement, l'image inv. de $C \hookrightarrow P(G?)$ par $\tilde{\theta}: M_{\rho, \sigma} \rightarrow P(G)$ d'où une suite exacte $1 \rightarrow R_{\rho, \sigma} \rightarrow M_{\rho, \sigma} \rightarrow \dots$]

\newpage

%% ===== Page 430 =====

Dans le cas modulaire, sur la base
absolue $\mathbb{Z}$, on aura

(81) $$S C_{\rho,\sigma}^0(\mathbb{Z}) \stackrel{\sim}{\to} C_{\rho,\sigma}^0(\mathbb{Z}) \simeq T_\rho(\mathbb{Z}) \times T_\sigma(\mathbb{Z}) \simeq \mathbb{Z}/6\mathbb{Z} \times \mathbb{Z}/4\mathbb{Z}$$
$S T_\rho(\mathbb{Z}) \times S T_\sigma(\mathbb{Z})$ car $T_\rho(\mathbb{Z}) = S T_\rho(\mathbb{Z})$, $T_\sigma(\mathbb{Z}) = S T_\sigma(\mathbb{Z})$

d'où, via (74)

(82) $$R_{\rho,\sigma}^0(\mathbb{Z}) \simeq (\mathbb{Z}/6\mathbb{Z} \times \mathbb{Z}/4\mathbb{Z}) \times \{\pm 1\}$$

les deux projections sur les deux facteurs
du second membre correspondant justement
aux hom.

$$R_{\rho,\sigma}^0 \to C_{\rho,\sigma}^0, \quad R_{\rho,\sigma}^0 \to G_m$$

de (74), de sorte que l'inclusion

(83) $$\mu_2(\mathbb{Z}) = \{\pm 1\} \hookrightarrow R_{\rho,\sigma}^0(\mathbb{Z}) = (\mathbb{Z}/6\mathbb{Z} \times \mathbb{Z}/4\mathbb{Z}) \times \{\pm 1\}$$

est l'inclusion évidente $\varepsilon \mapsto (0) \times \varepsilon$.

Une façon plus parlante de dire ceci
est que l'inclusion canonique

(84) $$S R_{\rho,\sigma}^0 \times \mu_2 \hookrightarrow R_{\rho,\sigma}^0$$
(avec $S R_{\rho,\sigma}^0 = S T_\rho \times S T_\sigma$ et $\mu_2 \cap M_{\rho,\sigma}^0(0) = Centr(S G)$)

induit un isomorphisme

(85) $$S R_{\rho,\sigma}^0(\mathbb{Z}) \times \mu_2(\mathbb{Z}) \stackrel{\sim}{\to} R_{\rho,\sigma}^0(\mathbb{Z})$$
(avec $S R_{\rho,\sigma}^0(\mathbb{Z}) = \mathbb{Z}/6\mathbb{Z} \times \mathbb{Z}/4\mathbb{Z}$ et $\mu_2(\mathbb{Z}) = \{\pm 1\}$)

\newpage

%% ===== Page 431 =====

[MARGIN: Commençons l'étude du radical $C R_{\rho,\sigma}$, $S R_{\rho,\sigma}$, $R_{\rho,\sigma}$, $S \dots$]

$6^o$) Je vais revenir sur

(86) $$R_{\rho,\sigma} \overset{\text{df}}{=} \operatorname{Ker}(\mathcal{M}_{\rho,\sigma} \overset{\overline{\Theta}}{\longrightarrow} \mathbb{P} G)$$

donnant lieu à la suite exacte

(87) $$1 \longrightarrow R_{\rho,\sigma} \longrightarrow \mathcal{M}_{\rho,\sigma} \longrightarrow \mathbb{P} G \longrightarrow 1 \ .$$

On a aussi

(88) $$R_{\rho,\sigma}^o \underset{\substack{\text{défini} \\ \text{dans 69}}}{=} R_{\rho,\sigma} \cap \mathcal{M}_{\rho,\sigma}^o = \substack{\text{comp. neutre} \\ \text{de } R_{\rho,\sigma}}$$

D'ailleurs on a vu

(89) $$R_{\rho,\sigma} = \Theta^{-1}(\mathbb{G}_m) \quad (\Theta: \mathcal{M}_{\rho,\sigma} \to G)$$

d'où une suite exacte

(90) 
$$1 \longrightarrow S R_{\rho,\sigma} \longrightarrow R_{\rho,\sigma} \overset{\mathbf{k}}{\longrightarrow} \mathbb{G}_m \longrightarrow 1$$
$$ \Vert \text{déf}$$
$$\text{Ker } \mathbf{k} \, (R_{\rho,\sigma} \overset{\mathbf{k}}{\longrightarrow} \mathbb{G}_m)$$
$$ \Vert$$
$$\text{Ker } \Theta \, (\mathcal{M}_{\rho,\sigma} \overset{\Theta}{\longrightarrow} G)$$

où $\mathbf{k}$ est défini par la condition que

(91)
\begin{tikzcd}
R_{\rho,\sigma} \ar[r, "\mathbf{k}"] \ar[drrr, "\Theta"', bend right] & \mathbb{G}_m \ar[r, "\simeq"] & \text{Centr}(G) \ar[r, hook] & G
\end{tikzcd}

soit commutatif.

Notons en passant que l'hom. $\chi$ de (42) $\mathcal{M}_{\rho,\sigma} \overset{\chi}{\longrightarrow} \mathbb{G}_m$ peut se déduire aussi de $\Theta$ (cf (57) par composition) :

(92)
\begin{tikzcd}
\mathcal{M}_{\rho,\sigma} \ar[r, "\Theta"] \ar[dr, "\chi"'] & G \ar[d, "\text{det}"] \\
& \mathbb{G}_m
\end{tikzcd}

et on a donc

(93) $$\chi(g) = \mathbf{k}(g)^2 \quad \text{si } g \in R_{\rho,\sigma}$$

i.e.

\begin{tikzcd}
R_{\rho,\sigma} \ar[r, "\mathbf{k}"] \ar[dr, "\chi | R_{\rho,\sigma}"'] & \mathbb{G}_m \ar[d, "\lambda \mapsto \lambda^2"] \\
& \mathbb{G}_m
\end{tikzcd}
commutatif.

\newpage

%% ===== Page 432 =====

On a le diagramme d'inclusions
(94)
\begin{tikzcd}
\underset{\text{"Ker }\kappa\text{"}}{SR_{\rho,\sigma}} \ar[r, hook, "\mathbb{G}_m \ (90)"] \ar[d, "\mu_2 \times \mu_2"'] & \underset{\text{"Ker }\Theta\text{"}}{R_{\rho,\sigma}} \ar[r, hook, "PG \ (86)"] \ar[d, "\text{cart et cocart}"'] & M_{\rho,\sigma} \ar[d, "\mu_2 \times \mu_2"] \\
\underset{\text{"Ker }\kappa^\circ\text{" \ (62)}}{\underset{\simeq S T_\rho \times S T_\sigma}{SR_{\rho,\sigma}^\circ}} \ar[r, hook, "\mathbb{G}_m"] & \underset{\text{"Ker }\Theta^\circ\text{" \ (69)}}{R_{\rho,\sigma}^\circ} \ar[r, hook, "PG"] & M_{\rho,\sigma}^\circ
\end{tikzcd}

[MARGIN: NB il y a une vérification à faire, qui se réduit à voir que les flèches $SR_{\rho,\sigma}^+ \to SR_{\rho,\sigma}$, $R_{\rho,\sigma}^+ \to R_{\rho,\sigma}$ sont des isomorph. — oui, cf plus bas en (100) où la flèche est rendue explicite par ex.]

La deuxième ligne est formée des composantes neutres de la première, qu'on peut réécrire plus joliment encore

(95)
\begin{tikzcd}
& R_{\rho,\sigma} \ar[dr, "PG"] \ar[dd] \\
SR_{\rho,\sigma} \ar[ur, "\mathbb{G}_m"] \ar[rr, dashed, "G"] \ar[dd, "\mu_2 \times \mu_2"'] & & M_{\rho,\sigma} \ar[dd, "\mu_2 \times \mu_2"] \\
& R_{\rho,\sigma}^\circ \ar[dr, "PG" near start] \\
\underset{(S T_\rho \times S T_\sigma)^-}{\underset{\parallel}{SR_{\rho,\sigma}^\circ}} \ar[ur, "\mathbb{G}_m"] \ar[rr, dashed, "G"'] & & M_{\rho,\sigma}^\circ
\end{tikzcd}

sous forme d'un prisme dont les trois carrés latéraux sont à la fois cartésiens et cocartésiens — donc donnent lieu à des quotients (sur des cotés parallèles) isomorphes — les quotients qui interviennent étant $\mu_2 \times \mu_2$, $\mathbb{G}_m$, $PG$ et l'extension $G$ de $

\newpage

%% ===== Page 433 =====

510

qui explicite la structure de $R_{p,\sigma}$ en termes
des groupes très simples $ST_p, ST_\sigma, \mathbb{G}_m, \mu_2$.

On a un cube, construit sur la face
verticale postérieure gauche de (95) (qui devient la
face verticale postérieure)

(95)
\begin{tikzcd}
S R_{p,\sigma} \ar[rr, "\mathbb{G}_m", hook] \ar[dd, "\mu_2 \times \mu_2"'] \ar[dr, "\mu_2"] & & R_{p,\sigma} \ar[rr, dashed] \ar[dd, "\mu_2 \times \mu_2"'] \ar[dr, "\mu_2"] & & \mathbb{G}_m \ar[dd, equal] \ar[dr, dashed, "\lambda \mapsto \lambda^2"] \\
& T_{p,\sigma}^+ \ar[rr, "\mathbb{G}_m", hook, crossing over] \ar[dd] & & T_{p,\sigma} \ar[rr, dashed, crossing over] \ar[dd] & & \mathbb{G}_m \ar[dd, equal] \\
S R_{p,\tilde{\sigma}} \ar[rr, "\mathbb{G}_m"', hook] \ar[dr, "\mu_2"] & & R_{p,\tilde{\sigma}} \ar[rr, dashed] \ar[dr, "\mu_2"] & & \mathbb{G}_m \ar[dr, dashed, "\lambda \mapsto \lambda^2"'] \\
& T_{p,\tilde{\sigma}}^+ \ar[rr, "\mathbb{G}_m"', hook] & & T_{p,\tilde{\sigma}} \ar[rr, dashed] & & \mathbb{G}_m
\end{tikzcd}

dont les faces verticales (antérieure et
postérieure) sont cartésiennes et cocartésiennes
et les flèches horizontales qui vont de l'une à
l'autre sont des morphismes de passage
au quotient, ayant tous des noyaux isomorphes
soit à $\mu_2$ (noyaux des flèches), soit à 1
(noyaux des autres flèches), ce qui
explicite les relations :

(97)
$\begin{cases}
R_{p,\sigma} \cap S G = R_{p,\sigma} \cap S h = \mu_2 \\
S R_{p,\sigma} \cap S h = S R_{p,\sigma} \cap S G = \{1\} \quad \text{explicite bien ça}
\end{cases}$

Pour la première ligne il s'agit de voir que
[unclear] est vrai dans la première ligne
ce qui équivaut à dire que l'on a une inclusion
$$\mu_2 = Sh_{p,\sigma} \cap S G \subset R_{p,\sigma} \cap S G$$

\newpage

%% ===== Page 434 =====

519

on a
$$ \underset{Ker \tilde{\Theta}}{R_{\rho,\sigma}} \cap SG = Ker (\tilde{\Theta} | SG) \quad , \quad \underset{Ker \tilde{\Theta}^\circ}{R_{\rho,\sigma}^\circ} \cap SG = Ker (\tilde{\Theta}^\circ | SG) ,$$
$$ \tilde{\Theta} | SG = \tilde{\Theta}^\circ | SG = \text{morphisme can } SG \to PG $$
dont le noyau est bien $\mu_2$.

D'autre part
$$ \underset{Ker \Theta}{S R_{\rho,\sigma} \cap S h} = Ker (\Theta | S h) = \{1\} $$
puisque $\Theta | S h$ est l'inclusion can. $S h \subset C_{\rho}$,

a fortiori $S R_{\rho,\sigma}^\circ \cap S h = \{1\}$. Le fait qu'on a

des isomorphismes

(98)
$$ S R_{\rho,\sigma} \stackrel{\sim}{\longrightarrow} C_{\rho,\sigma}^+ \simeq S \Gamma_{\rho} \times W_{\sigma}' $$
$$ S R_{\rho,\sigma}^\circ \stackrel{\sim}{\longrightarrow} C_{\rho,\sigma}^{+ \circ} \simeq S \Gamma_{\rho} \times S \Gamma_{\sigma} $$

résulte, pour le deuxième, de la comparaison
de (62) et de (28), pour le premier il s'en
déduit par le lemme des cinq, puisque
$$ S R_{\rho,\sigma} / S R_{\rho,\sigma}^\circ \stackrel{\sim}{\longrightarrow} C_{\rho,\sigma}^+ / C_{\rho,\sigma}^{+ \circ} \simeq \mu_2 \times \mu_2 . $$

Le fait que l'on ait des iso (98)
implique que dans
$$ \underset{[unclear: (= \operatorname{Ker}(\dots))]}{M_{\rho,\sigma}^+} \simeq \underset{\simeq C_{\rho,\sigma}^+}{M_{\rho,\sigma(0)}^+} . \underset{SG}{S_{\rho,\sigma}^+} \quad (\text{produit } \frac{1}{2} \text{ direct cf (44)}) $$
regardé comme extension de $C_{\rho,\sigma}^+$ par $SG$,
il y a, à côté du scindage par $M_{\rho,\sigma(0)}^+ \subset M_{\rho,\sigma}^+$
(donné par $M_{\rho,\sigma(0)}$), un autre scindage,

\newpage

%% ===== Page 435 =====

liée de façon beaucoup plus naturelle à la
structure du groupe $M_{\rho,\sigma}$ et à ses hom.
[crossed out: D...s] $G$ et $\gamma = Out \Theta$ dans $\hat{\Gamma}_{n,1}$, savoir le
sous-groupe section $S R_{\rho,\sigma}$, permettant de l'écrire
dans $M_{\rho,\sigma}^+$ comme un produit direct
$$ (99) \quad M_{\rho,\sigma}^+ \simeq \underset{\parallel S' T_{\rho} \times S N'_{\sigma}}{S R_{\rho,\sigma}} \times \underset{\parallel S G}{\mathfrak{S}_{\rho,\sigma}^+} $$

induisant une décomposition analogue

[MARGIN: c'est (45), par le biais d'un autre dévissage immédiat]

$$ (100) \quad M_{\rho,\sigma}^{+0} \simeq \underset{S T_{\rho} \times S T_{\sigma}}{S R_{\rho,\sigma}^0} \times \underset{S G}{\mathfrak{S}_{\rho,\sigma}^{+0}} $$
$$ \simeq S T_{\rho} \times S T_{\sigma} \times S G $$

qui dit mieux !. D'autre part, le fait
que $R_{\rho,\sigma} \to \Gamma_{\rho,\sigma}$ soit un isomorphisme
nous montre que l'on a, à isomorphisme
près, une "section" [crossed out text] de $M_{\rho,\sigma}$ sur $\Gamma_{\rho,\sigma}$ en tant
qu'extension de $\Gamma_{\rho,\sigma}$ par $\mathfrak{S}_{\rho,\sigma}^+ \simeq S G$.

En fait, on a
$$ R_{\rho,\sigma} \cap S G = \{1\} $$
[crossed out text]
$R_{\rho,\sigma}$ commute à $S G$
(trivial par définition de
$R_{\rho,\sigma}$ comme $\ker \tilde{\Theta}$)
donc on a une suite exacte

\newpage

%% ===== Page 436 =====

510

(102)
$$ 1 \longrightarrow \mu_2 \longrightarrow R_{\rho,\sigma} \times S G \longrightarrow M_{\rho,\sigma} \longrightarrow 1 $$
$$ \varepsilon \longmapsto (\varepsilon, \varepsilon^{-1}) $$

et de même

(103)
$$ 1 \longrightarrow \mu_2 \longrightarrow R_{\rho,\sigma}^+ \times S G \longrightarrow M_{\rho,\sigma}^+ \longrightarrow 1 $$

déjà envisagée (un peu timidement) dans (78).
C'est le moment décidément d'expliciter
entièrement la description de $R_{\rho,\sigma}$ en
termes de $T_\rho, N_\sigma'...$ -- on l'obtient, dans
le contexte de la p. 510, comme

(104)
$$ R_{\rho,\sigma} = \left\{ (\underset{\in}{\underset{\Gamma_{T_0}}{u}}, \underset{\in}{\underset{\Gamma_{T_\rho}}{\underline{a}}}, \underset{\in}{\underset{\Gamma_{N'_\sigma}}{\underline{b}}}, \underset{\in}{\underset{SG}{f}}) \left| \begin{array}{l} \det u = \det \underline{b} = \varepsilon' \det \underline{a},\ \varepsilon'^2=1 \\ u f = \lambda \\ (\text{d'où } \lambda^2 = \det u) \end{array} \quad (\varepsilon',\lambda \in \Gamma_{\mathbb{G}_m}) \right. \right\} $$
$$ = \left\{ (\underset{\in}{\underset{T_0}{u}}, \underset{\in}{\underset{T_\rho}{\underline{a}}}, \underset{\in}{\underset{N'_\sigma}{\underline{b}}}, \underset{\in}{\underset{\mathbb{G}_m}{\lambda}}) \left| \det u = \det \underline{b} = \varepsilon' \det \underline{a} = \lambda^2 \right. \right\} $$
$$ \simeq \left\{ (\underset{\in}{\underset{T_\rho}{\underline{a}}}, \underset{\in}{\underset{N'_\sigma}{\underline{b}}}, \underset{\in}{\underset{\mathbb{G}_m}{\lambda}}) \left| \det \underline{b} = \lambda^2,\ \det \underline{a}^2 = \det \underline{b}^2 \right. \right\} $$
$$ \simeq (\widetilde{T}_\rho, \delta_\rho^2) \times_{\mathbb{G}_m} (\widetilde{N}'_\sigma, \delta_\sigma^2) \times_{\mathbb{G}_m (\delta^2)} \mathbb{G}_m (\delta^2) $$

ou plus simplement ,

[raturé]
$$ R_{\rho,\sigma} \simeq \left\{ (\underset{\in}{\underset{T_\rho}{\underline{a}}}, \underset{\in}{\underset{N'_\sigma}{\underline{b}}}, \underset{\in}{\underset{\mathbb{G}_m}{\lambda}}) \left| \det \underline{a} = \det \underline{b} = \lambda^2 \right. \right\} $$

(105)
\begin{tikzcd}
R_{\rho,\sigma} \ar[r] \ar[d] & \widetilde{C}_{\rho,\sigma} \ar[d] \\
\mathbb{G}_m \ar[r, "\lambda \mapsto \lambda^2"'] & \mathbb{G}_m
\end{tikzcd}
diag. cartésien

\newpage

%% ===== Page 437 =====

s'insère dans un diagramme à suites exactes

\begin{tikzcd}
& & & 1 \ar[d] & 1 \ar[d] & \\
& & & S R_{\rho,\sigma} \ar[r, "\sim"] \ar[d] & \mathcal{C}^+_{\rho,\sigma} \ar[d] & \\
(105) & 1 \ar[r] & \mu_2 \ar[r] \ar[d, equal] & R_{\rho,\sigma} \ar[r, "can."] \ar[d, "\pi"] & \mathcal{C}_{\rho,\sigma} \ar[r] \ar[d, "\chi"] & 1 \\
& 1 \ar[r] & \mu_2 \ar[r] & \mathbb{G}_m \ar[r, "\lambda \mapsto \lambda^2"] \ar[d] & \mathbb{G}_m \ar[r] \ar[d] & 1 \\
& & & 1 & 1 & 
\end{tikzcd}

qui n'est qu'une explicitation de la
fin supérieure (*) [MARGIN: (*) plus au-dessus] des diagrammes analogues
(95) -- faire attention que celle-ci s'est
construite [plus bas] sur l'hom.
(107) groupe $\mathcal{C}_{\rho,\sigma}$, et $(\mathcal{C}_{\rho,\sigma} \xrightarrow{\chi} \mathbb{G}_m, \mathcal{C}_{\rho,\sigma} \xrightarrow{(\varepsilon_g, \varepsilon_d)} \mu_2 \times \mu_2)$
[ et par là exprime : (105). - (106) $R_{\rho,\sigma}$,
$S R_{\rho,\sigma} \simeq \mathcal{C}^+_{\rho,\sigma}$, et aussi (en [unclear: remplaçant]
$\mathcal{C}^0_{\rho,\sigma} = \operatorname{Ker}(\mathcal{C}_{\rho,\sigma} \to \mu_2 \times \mu_2)$
$R^0_{\rho,\sigma}, S R^0_{\rho,\sigma} \simeq \mathcal{C}^0_{\rho,\sigma}$ ] -- en termes des
groupes
(108) $T_{\rho}, N'_{\sigma}$ et de $T_{\rho} \xrightarrow{\delta_{\rho}} \mathbb{G}_m, N'_{\sigma} \xrightarrow{\delta_{\sigma}} \mathbb{G}_m$
$N'_{\sigma} \xrightarrow{\varepsilon_{\sigma}} \mu_2$

\newpage

%% ===== Page 438 =====

511

comme sous-groupe

(109)
\begin{tikzcd}
\mathcal{T}_{\rho, \sigma} \ar[r, hook] \ar[d, equal] & T_0 \times T_\rho \times N'_\sigma \\
\{(\underline{u}, \underline{a}, \underline{b}) \mid \delta_0(\underline{u}) = \delta_\rho(\underline{a}) = \varepsilon' \delta_\sigma(\underline{b}) , {\varepsilon'}^2 = 1 \} & T_0 \ar[d, "\delta_0"] \\
\simeq \{(\underline{a}, \underline{b}) \in T_\rho \times N'_\sigma \mid \delta_\rho(\underline{a}) = \varepsilon' \delta_\sigma(\underline{b}) , {\varepsilon'}^2 = 1 \} & \mathcal{G}_m
\end{tikzcd}

ou

(110)
$$
\begin{cases}
\chi(\underline{a}, \underline{b}) = \delta_\rho(\underline{a}) \\
\varepsilon_3(\underline{a}, \underline{b}) = \delta_\rho(\underline{a}) / \delta_\sigma(\underline{b}) \\
\varepsilon_1(\underline{a}, \underline{b}) = \varepsilon_\sigma(\underline{b})
\end{cases}
$$

A l'aide de ceci on peut construire $M_{\rho, \sigma}$
par (102), et bien entendu l'hom.

(111) $$ \underline{\theta} : M_{\rho, \sigma} \longrightarrow \underline{G} $$

se déduit par passage au quotient des hom.
évidents

(112)
\begin{tikzcd}
S\underline{G} \ar[r, hook] & \underline{G} , & R_{\rho, \sigma} \ar[r] \ar[dr, "\kappa"'] & \underline{G} \\
& & & \underline{G}_m \ar[u, "\simeq"'] \ar[r, "\simeq"] & \underline{C}ent(\underline{G})
\end{tikzcd}

On vérifie que on en déduit les compatibilités

(113)
\begin{tikzcd}
M_{\rho, \sigma} \ar[r, "\underline{\theta}"] \ar[dr, "\underline{\tilde{\theta}}"'] & \underline{G} \ar[d] \\
& P\underline{G}
\end{tikzcd}

$$ M_{\rho, \sigma} \overset{\underline{\theta}}{\longrightarrow} \underline{G} \overset{det}{\longrightarrow} \underline{G}_m $$

et on en déduit la suite exacte

(114)
\begin{tikzcd}
1 \ar[r] & S\underline{G} \ar[r] \ar[dr, "can"'] & M_{\rho, \sigma} \ar[r] & \mathcal{T}_{\rho, \sigma} \ar[r] \ar[d, equal] & 1 \\
& & R_{\rho, \sigma} \times S\underline{G} \ar[u] & R_{\rho, \sigma} / \mu_2
\end{tikzcd}

\newpage

%% ===== Page 439 =====

de la description (102), d'où des composés

(115)
\begin{tikzcd}
& & T_\rho \\
M_{\rho, \sigma} \ar[r] & T_{\rho, \sigma} \ar[ur, "pr_1"] \ar[r, "pr_2"'] \ar[dr] & N'_\sigma \ar[r, "det"] & \mathbb{G}_m \\
& & \mu \times \mu_2
\end{tikzcd}

Dans tout ceci, on n'a pas eu besoin de la
donnée de sections $\rho$, $\sigma$ de $G$, ni de
connaître quoi que ce soit sur la façon
dont $T_\rho$, $N'_\sigma$ sont [unclear: obtenus] à partir de telles
sections. Mais des structures supplémentaires
plus délicates dans $M_{\rho,\sigma}$ - telle la
donnée d'un sous-groupe section $M_{\rho,\sigma}(s)$ -
en dépendent, et permettront peut-être de
les retrouver.

Rmq. Le cadre (96) est le "compagnon"
d'un diagramme analogue, que je
ne sais expliciter, dont le [unclear]
complète (95) - qui peut être considéré
comme un diagramme partiel de
celui-ci :

\newpage

%% ===== Page 440 =====

Comme d. (95) (96), 6 faces inf. 
et la face sup. du prisme ne sont pas des
carrés cartésiens (ni cocartésiens), les quatre groupes
quotients s'identifient canoniquement entre eux par les
flèches en jeu, ils sont isomorphes à 
$M_{\rho,\sigma}^0$ --

(116)
\begin{tikzcd}[row sep=1.5em, column sep=1.5em]
G_m \ar[ddd, "\lambda \mapsto \lambda^2"'] \ar[drr, dashed] & & & G_m \ar[ddd, "\lambda \mapsto \lambda^2"'] \ar[drr, dashed] & & & G_m \ar[ddd, "\lambda \mapsto \lambda^2"'] \ar[ddr, dashed] & & \\
& & S R_{\rho,\sigma} \ar[rrr, hook] \ar[ddd, "pr_1 \times pr_2" near start] \ar[dl] & & & R_{\rho,\sigma} \ar[rrr] \ar[ddd, "pr_1 \times pr_2" near start] \ar[dl] & & & P G \ar[ddd, "id"] \ar[dl] \\
& S G \ar[rrr, crossing over, hook] \ar[ddd, "id"'] & & & G \ar[rrr, crossing over] \ar[ddd, "id"'] & & & M_{\rho,\sigma} \ar[ddd, "pr_1"'] & \\
G_m \ar[drr, dashed] & & & G_m \ar[drr, dashed] & & & G_m \ar[ddr, dashed] & & \\
& & S R_{\rho,\sigma}^0 \ar[rrr, hook] \ar[dl] & & & R_{\rho,\sigma}^0 \ar[rrr] \ar[dl] & & & P G \ar[dl] \\
& S G \ar[rrr, hook] & & & G \ar[rrr] & & & M_{\rho,\sigma}^0 & 
\end{tikzcd}

Dans le cube ci-dessous (117), les faces latérales du cube sont
des carrés cartésiens, mais pas les faces
inf. et sup. en général.

\newpage

%% ===== Page 441 =====

(117)
\begin{tikzcd}[row sep=2em, column sep=2em]
& R_{\rho, \sigma} \ar[dl] \ar[d, "\mathbb{G}_m"', dashed] \ar[dr] \ar[dddd, "pr"'] & \\
S R_{\rho, \sigma} \ar[d] \ar[drr, "\simeq" near start] \ar[dddd, "pr"'] & G \ar[dl] \ar[r] \ar[dr] \ar[dddd, "pr"'] & P G \ar[dd] \\
S G \ar[dr] \ar[dddd, "pr"'] & & M_{\rho, \sigma} \ar[dl, "+"] \ar[dr] \ar[dddd, "pr"'] & \\
& \mathcal{M}_{\rho, \sigma}^+ \ar[dddd, "pr"'] & & C_{\rho, \sigma} \ar[dddd, "pr"'] \\
& R_{\rho, \sigma}^0 \ar[dl] \ar[d, "\mathbb{G}_m"', dashed] & \\
S R_{\rho, \sigma}^0 \ar[d] \ar[drr, "\simeq" near start] & G^0 \ar[dl, dashed] \ar[dr] & \\
S G \ar[dr] & & M_{\rho, \sigma}^0 \ar[dl, "+ 0"] \ar[dr] & \\
& \mathcal{M}_{\rho, \sigma}^{+ 0} & & C_{\rho, \sigma}^0
\end{tikzcd}

La face comprise entièrement à gauche de ce cube est de nature particulièrement simple, puisque les termes médians $\mathcal{M}_{\rho, \sigma}^+$, $\mathcal{M}_{\rho, \sigma}^{+ 0}$ s'identifient explicitement aux produits directs des tronçons $S R_{\rho, \sigma} \simeq \mathcal{T}_{\rho, \sigma}^+$ et $S R_{\rho, \sigma}^0 \simeq \mathcal{T}_{\rho, \sigma}^{+ 0}$
$$ \simeq S T_{\rho} \times S N_{\sigma}' \quad \text{et} \quad \simeq S T_{\rho} \times S T_{\sigma} $$
avec le groupe $S G$. C'est pourquoi le [unclear: as-is] pour la face parallèle (face postérieure droite), on gagne [unclear: le prix] - [unclear: à condition] de remplacer les groupes
médians $M_{\rho, \sigma}$, $M_{\rho, \sigma}^0$ par les pr-revêtements

(118) $$ \mathcal{M}_{\rho, \sigma} \simeq R_{\rho, \sigma} \times S G, \quad \mathcal{M}_{\rho, \sigma}^0 \simeq R_{\rho, \sigma}^0 \times S G $$

cf (102), (103)) , ce qui [unclear: assure] la [unclear: suite].

\newpage

%% ===== Page 442 =====

L'utilité de rendre les quatre carrés des
faces sup. et inférieures des deux cubes accolés
cartésiens et cocartésiens, ce qu'ils n'étaient pas
avant, regrettablement.

Pour terminer ces gloses sur le
radiciel, je vais introduire un hom
canonique

$$ (118) \qquad \mathbb{G}_m \overset{i_R}{\hookrightarrow} R_{\rho,\sigma}^\circ $$

via un hom canonique

$$ (119) \qquad \mathbb{G}_m \overset{i_\tau}{\hookrightarrow} \tau_{\rho,\sigma}^\circ \simeq \mathbb{T}_\rho \times_{\mathbb{G}_m} \mathbb{T}_\sigma $$
$$ \lambda \longmapsto (\lambda, \lambda) $$

dont la composée avec

$$ \nu : \tau_{\rho,\sigma}^\circ \longrightarrow \mathbb{G}_m $$

est donnée en fait par $\lambda \longmapsto \lambda^2$, de sorte
que grâce à la description (105) (ou (109)) de $R_{\rho,\sigma}$
comme produit fibré (paraphrasée en une
expression de $R_{\rho,\sigma}^\circ$ en termes de $\tau_{\rho,\sigma}^\circ$)
on a un unique hom (118), rendant
commutatif le diagramme

\newpage

%% ===== Page 443 =====

52-9

(120)
\begin{tikzcd}
\mathbb{G}_m \ar[r, "i_\tau"] \ar[dr, "i_R"'] & \mathcal{T}_{\rho,\sigma}^\circ \\
& R_{\rho,\sigma}^\circ \ar[u, "can"']
\end{tikzcd}

-, plus joliment, le diagramme en
suite exacte

(121) $$1 \longrightarrow \mu_2 \longrightarrow \mathbb{G}_m \xrightarrow{\lambda \mapsto \lambda^2} \mathbb{G}_m \longrightarrow 1$$ .

Par construction, on trouve aussitôt

(122)
\begin{tikzcd}
\mathbb{G}_m \ar[r, "i_R"] \ar[dr, "id"'] & R_{\rho,\sigma}^\circ \ar[d, "k"] \\
& \mathbb{G}_m
\end{tikzcd}

d'où des décompositions en produit direct

(123)
$$
\begin{cases}
R_{\rho,\sigma} \simeq S R_{\rho,\sigma} \times \mathbb{G}_m \simeq S' \mathcal{T}_\rho \times S \mathcal{N}_\sigma' \times \mathbb{G}_m \\
R_{\rho,\sigma}^\circ \simeq S R_{\rho,\sigma}^\circ \times \mathbb{G}_m \\
\quad \simeq S \mathcal{T}_\rho \times S \mathcal{T}_\sigma \times \mathbb{G}_m
\end{cases}
$$

qui précise considérablement (96) dont j'étais
fort peu satisfait !

En termes de cette décomposition, on
explicite l'hom.
$$ \mu_2 \hookrightarrow R_{\rho,\sigma}^\circ \subset R_{\rho,\sigma} $$ (et $R_{\rho,\sigma} \twoheadrightarrow \mathcal{T}_\sigma$, $R_{\rho,\sigma}^\circ \twoheadrightarrow \mathcal{T}_\sigma^\circ$)
qui intervient dans (102), (103), $\mu_2 = R_{\rho,\sigma}^\circ \cap S R_{\rho,\sigma}^\circ = R_{\rho,\sigma} \cap S R_{\rho,\sigma}$
~~[unclear crossed out text]~~

\newpage

%% ===== Page 444 =====

Pour l'hom. 
$$R_{\rho,\sigma} \to T_{\rho,\sigma} \subset T_\rho \times N_\sigma$$
il suffit d'expliciter les composées
(124)
$$
\begin{tikzcd}
R_{\rho,\sigma} \ar[r] \ar[d, "\simeq"'] & T_{\rho,\sigma} \ar[r] \ar[dr] & T_\rho \\
S' T_\rho \times S N'_\sigma \times \mathbb{G}_m & & N_\sigma
\end{tikzcd}
$$
$$(a,b,\lambda) \mapsto (\lambda a, \lambda b)$$
[MARGIN: on a en effet $\det(\lambda a) = \det(\lambda b)^2$, $\lambda^2 = \lambda^2$ car $\det(a)=1, \det \underline{b}=1$]

Le noyau de cet hom. est donc
formé des $(\lambda^{-1}, \lambda^{-1}, \lambda)$, avec $\lambda \in \mathbb{G}_m$ tel que
$\lambda^{-1} \in S'T_\rho$, $\lambda^{-1} \in S N'_\sigma$, c.a.d. i.e. $(\lambda^2)^? = 1$ et $\lambda^2 = 1$,
soit $\lambda \in \mu_2$. Donc

(115) [unclear: (125)]
$$
\begin{tikzcd}
\mu_2 \ar[r, "can"] & R_{\rho,\sigma}^\circ \ar[r, hook] & R_{\rho,\sigma} \\
R_{\rho,\sigma} \cap SG \ar[u, phantom, sloped, "\simeq"] & S T_\rho \times S T_\sigma \times \mathbb{G}_m \ar[u, phantom, sloped, "\simeq"] & S' T_\rho \times S N'_\sigma \times \mathbb{G}_m \ar[u, phantom, sloped, "\simeq"]
\end{tikzcd}
$$
$$\lambda \mapsto (\lambda, \lambda, \lambda)$$

On trouve donc l'explicitation particulièrement
simple
(126)  $$T_{\rho,\sigma} \simeq (S' T_\rho \times S N'_\sigma \times \mathbb{G}_m) / \mu_2$$
par envoi "diagonal" dans le
produit triple

On notera qu'on a des hom. canoniques
$$\mu_2 \to S T_\rho \subset S' T_\rho, \quad \mu_2 \to S T_\sigma \subset S N'_\sigma, \quad \mu_2 \to \mathbb{G}_m$$
(qui sont utilisés dans (125)) d'où
(127)
$$
\begin{tikzcd}
\mu_2 \times \mu_2 \times \mu_2 \ar[r] & S T_\rho \times S T_\sigma \times \mathbb{G}_m \ar[r, hook] & S' T_\rho \times S N'_\sigma \times \mathbb{G}_m \\
& R_{\rho,\sigma}^\circ \ar[u, phantom, sloped, "\simeq"'] & R_{\rho,\sigma} \ar[u, phantom, sloped, "\simeq"']
\end{tikzcd}
$$
$$(\lambda', \lambda'', \lambda) \mapsto (\lambda', \lambda'', \lambda)$$
qui n'est autre que (98)).

\newpage

%% ===== Page 445 =====

Le centre de $M_{p,\sigma}$. [MARGIN: 525] Il s'ensuit par $\widetilde{\Theta}$ dans le cadre du $PG$ que ce trivial, on a

$$ \underline{Cent}(M_{p,\sigma}) \subset Ker \widetilde{\Theta} = R_{p,\sigma} \simeq S'T_p \times SN'_\sigma \times \mathbb{G}_m $$

et évidemment. Comme $R_{p,\sigma}$ commute à $SG$, et que $R_{p,\sigma} . SG = M_{p,\sigma}$, on en conclut

(128) $\underline{Cent}(M_{p,\sigma}) = \underline{Cent}(R_{p,\sigma}) = S'T_p \times \underline{Cent}(SN'_\sigma) \times \mathbb{G}_m$

déterminons donc $\underline{Cent}(SN'_\sigma)$. Notons que

$$ \begin{cases} \mu_2 \simeq N'_\sigma / T_\sigma \text{ opère fidèlement sur } T_\sigma \\ \text{et } T_\sigma = \mathbb{G}_m . ST_\sigma \end{cases} $$
[MARGIN: donc $\mu_2$ opère fidèlement sur $ST_\sigma$]

d'où

(129) $\underline{Cent}(SN'_\sigma) \subset ST_\sigma$ d'où $\underline{Cent}(SN'_\sigma) \simeq (ST_\sigma)^{\mu_2}$

Où on constate que $T_\sigma / \mu_2 \simeq \mathbb{G}_m \hookrightarrow ST_\sigma$ par $x \mapsto x^{-1}$ ce qui montre que

$\underline{Cent}(SN'_\sigma) \subset (ST_\sigma)^{\mu_2}$

On a pour une section de $ST_\sigma$, $g = c\sigma + d$
$det g = - \nu c^2 + d^2 = 1$
$g \bar{g} = \nu c^2 + d^2 + 2cd\sigma = 1$

i.e.
$$ \begin{cases} - \nu c^2 + d^2 = 1 \\ \nu c^2 + d^2 = 1 \\ 2cd = 0 \end{cases} $$

d'où i.e.
$$ \begin{cases} 2\nu c^2 = 0 \quad \text{donc } 2c^2 = 0 \quad \text{(car } \nu \text{ inv.)} \\ 2cd = 0 \\ -\nu c^2 + d^2 = 1_{c^2=0 \Rightarrow d^2=1} \text{ (donc } d \text{ inv.)} \end{cases} $$

2 inv. cela revient à $c=0$, $d=\pm 1$ i.e. $g \in \mu_2$

(131) $\underline{Cent}(SN'_\sigma) = \mu_2$

Dans le cas général on aura simplement

(132) $\mu_2 \subset \underline{Cent}(SN'_\sigma)$

et cette inclusion induit un iso. sur les [unclear] sections dans le cas où 2 est inversible [unclear].

\newpage

%% ===== Page 446 =====

[unclear: $\underline{C}_{es}$. Diagrammes 
(131) $\mu'_2 \subset S T_\sigma$
sous-groupe de $S T_\sigma$ d'où finales, qu'on a :]

(128) $$\underline{Cent} (\underline{M}_{\rho,\sigma}) \simeq S' T_\rho \times \mu'_2 \times \mathbb{G}_m \subset \underline{R}_{\rho,\sigma} \simeq S' T_\rho \times S N'_\sigma \times \mathbb{G}_m$$
d'où
(129) $$\underline{Cent} (\underline{M}_{\rho,\sigma}) \cap \underline{M}_{\rho,\sigma}^\circ = S T_\rho \times \mu'_2 \times \mathbb{G}_m \subset \underline{R}_{\rho,\sigma}^\circ \simeq S T_\rho \times S T_\sigma \times \mathbb{G}_m$$

Revenons aux sections entières dans le cas
modulaire absolu

(130) $$\underline{R}_{\rho,\sigma} (\mathbb{Z}) \simeq \underbrace{S' T_\rho (\mathbb{Z})}_{T_\rho(\mathbb{Z}) \simeq \mathbb{Z}/6\mathbb{Z}} \times \underbrace{S N'_\sigma (\mathbb{Z})}_{T_\sigma(\mathbb{Z}) \simeq \mathbb{Z}/4\mathbb{Z}} \times \underbrace{\mathbb{G}_m (\mathbb{Z})}_{\mu(\mathbb{Z}) \simeq \pm 1}$$

(131) $$\underline{Cent} (\underline{M}_{\rho,\sigma}) (\mathbb{Z}) \simeq S' (T_\rho [\mathbb{Z}]) \times \mu_2 (\mathbb{Z}) \times \mathbb{G}_m (\mathbb{Z})$$
$$\simeq \mathbb{Z}/6\mathbb{Z} \times \{\pm 1\} \times \{\pm 1\}$$
et pour mémoire (cf (75), p. 514)

(132) $$\underline{Cent} (\underline{M}_{\rho,\sigma}^\circ) (\mathbb{Z}) \simeq S T_\rho (\mathbb{Z}) \times S T_\sigma (\mathbb{Z}) \times \mathbb{G}_m (\mathbb{Z})$$
$$\underline{R}_{\rho,\sigma}^\circ \simeq \mathbb{Z}/6\mathbb{Z} \times \mathbb{Z}/4\mathbb{Z} \times \{\pm 1\}$$

En résumé, on a le diagramme d'inclusions (strictes)
de schémas en groupes

(133)
\begin{tikzcd}
\text{dim 3} & \begin{matrix} \underline{Cent} (\underline{M}_{\rho,\sigma}^\circ) = \underline{R}_{\rho,\sigma}^\circ \\ \simeq S T_\rho \times S T_\sigma \times \mathbb{G}_m \end{matrix} \ar[r, hook, "\text{(quotient } \mu_2 \times \mu_2)"'] & \begin{matrix} \underline{R}_{\rho,\sigma} \\ \simeq S' T_\rho \times S N'_\sigma \times \mathbb{G}_m \end{matrix} & \text{dim 3} \\
\text{dim 2} & \begin{matrix} \underline{Cent} (\underline{M}_{\rho,\sigma}) \cap \underline{M}_{\rho,\sigma}^\circ \\ = S T_\rho \times \mu'_2 \times \mathbb{G}_m \end{matrix} \ar[u, hook] \ar[r, hook, "\text{(quotient } \mu_2)"'] & \begin{matrix} \underline{Cent} (\underline{M}_{\rho,\sigma}) \\ \simeq S' T_\rho \times \mu'_2 \times \mathbb{G}_m \end{matrix} \ar[u, hook] & \text{dim 2}
\end{tikzcd}

\newpage

%% ===== Page 447 =====

526

valable dans le cas général - et qui, dans le
cas modulaire absolu, induit sur les
points entiers des isomorphismes par les
deux flèches horizontales, correspondant
resp. aux sous-groupes

(134) $$ \begin{cases} \underline{Cent}(\mathcal{M}_{p,\sigma}^\circ)(\mathbb{Z}) \simeq \mathbb{Z}/6\mathbb{Z} \times \mathbb{Z}/4\mathbb{Z} \times \{\pm 1\} \\ \underline{Cent}(\mathcal{M}_{p,\sigma})(\mathbb{Z}) \simeq \mathbb{Z}/6\mathbb{Z} \times \{1\} \times \{\pm 1\} \end{cases} $$

qui en est un sous-groupe d'indice 2.
Enfin, pour mémoire signalons à nouveau
le groupe

(135) $$ \Sigma = \{\pm 1\} \times \{\pm 1\} \times \{\pm 1\} $$

contenu dans $ \underline{Cent}(\mathcal{M}_{p,\sigma})(\mathbb{Z}) $ ([unclear: càd] d'ordre 8,
d'indice 3), qui joue un rôle particulier
dans l'application arithmétique, puisqu'on
a un homomorphisme canonique

(136) $$ \mathcal{M}_{0,3}^{\sim} \to \mathcal{M}_{p,\sigma}(\hat{\mathbb{Z}}) / \Sigma $$

\newpage

%% ===== Page 448 =====

[MARGIN: Bien décrire $\mathcal{M}_{\rho,\sigma}$]

7° Faisons une liste récapitulative des principaux homomorphismes de groupes envisagés au cours de la section 6°, impliquant le groupe $\mathcal{M}_{\rho,\sigma}$

(134)
\begin{tikzcd}[row sep=1.5em, column sep=3em]
& & & & PG \\
& SG \ar[dr, "j_G"] & & G \ar[ur] \ar[dr, "\xi = \det_G"'] \\
& S'T_\rho \ar[dr, "j'_\rho"] & & & \mathbb{G}_m \\
\mu_2 \ar[uuuur] \ar[uur] \ar[ddr] \ar[ddddr] & & \mathcal{M}_{\rho,\sigma} \ar[uuur, "\Theta = \Theta_G" near start] \ar[ur, "\Theta_\rho" near start] \ar[dr, "\Theta_\sigma" near start] \ar[dddr, "\varepsilon_3" near start] & T_\rho \ar[r, "\delta_\rho = \det_{T_\rho}"] & \mathbb{G}_m \\
& SN'_\sigma \ar[ur, "j'_\sigma"'] & & N'_\sigma \ar[r, "\delta'_\sigma = \det_{N'_\sigma}"] \ar[dr, "\varepsilon_\sigma"'] & \mathbb{G}_m \\
& & & & \mu_2 \\
& \mathbb{G}_m \ar[uuur, "j_\mu"'] & & \mu_2
\end{tikzcd}

Notons que ce diagramme se reconstruit (avec les propriétés essentielles que nous venons de repasser en revue) : à partir des données

$$
(135) \qquad
\begin{cases}
\mathbb{G}_m \xrightarrow{i_\rho} T_\rho \xrightarrow{\delta_\rho} \mathbb{G}_m & (\delta_\rho i_\rho(\lambda) = \lambda^2) \\
\mathbb{G}_m \xrightarrow{i_\sigma} N'_\sigma \xrightarrow{\delta'_\sigma} \mathbb{G}_m & (\delta'_\sigma i_\sigma(\lambda) = \lambda^2) \\
N'_\sigma \xrightarrow{\varepsilon_\sigma} \mu_2 & \varepsilon_\sigma i_\sigma(\lambda) = 1
\end{cases}
$$

[diagramme rayé]
\begin{tikzcd}
1 \ar[r] & \mu_2 \ar[r] \ar[d] & SG \ar[r] \ar[d] & PG \ar[r] \ar[d, equal] & 1 \\
1 \ar[r] & \mathbb{G}_m \ar[r] \ar[d, "\lambda \mapsto \lambda^2"'] & G \ar[r] \ar[d] & PG \ar[r] & 1 \\
& \mathbb{G}_m \ar[r, equal] & \mathbb{G}_m
\end{tikzcd}

plus de la donnée familière

(136)
\begin{tikzcd}
& \mathbb{G}_m \ar[r, equal] & \mathbb{G}_m \\
1 \ar[r] & \mathbb{G}_m \ar[r, "i_G"] \ar[u, "\lambda \mapsto \lambda^2"'] & G \ar[r] \ar[u, "\delta_G"'] & PG \ar[r] & 1 \\
1 \ar[r] & \mu_2 \ar[r] \ar[u, hook] & SG \ar[r] \ar[u, hook] & PG \ar[r] \ar[u, equal] & 1 \\
& 1 \ar[u] & 1 \ar[u]
\end{tikzcd}

\newpage

%% ===== Page 449 =====

517

Rappelons qu'on a défini, en termes de (135) :

(137) 
$$S T_p = \ker \delta_p, S' T_p = \ker \delta_p^2$$
$$T_\sigma = \ker \varepsilon_\sigma, S N'_\sigma = \ker \delta'_\sigma, S T_\sigma = S N'_\sigma \cap T_\sigma = \ker \delta_\sigma$$
$$(\delta_\sigma = \delta'_\sigma | T_\sigma)$$
$$\mu_2 \to S T_p \subset S' T_p \text{ comme comp. de } \mu_2 \subset \mathbb{G}_m \xrightarrow{i_p} T_p$$
$$\mu_2 \to S T_\sigma \subset S N'_\sigma \text{ [unclear: par] } \mu_2 \subset \mathbb{G}_m \xrightarrow{i_\sigma} T_\sigma$$

ce qui détermine la partie gauche du diagramme (134), et $\mathcal{M}_{g,\nu}$ peut être considéré tout bêtement défini comme quotient de la somme amalgamée des quatre facteurs $SG$, $S' T_p$, $S N'_\sigma$, $\mathbb{G}_m$, par $\mu_2 \hookrightarrow$ .
Cela permet en principe de réduire la description des quatre morphismes fondamentaux de $\mathcal{M}_{g,\nu}$ dans (134) — $\Theta_G, \Theta_p, \Theta_\sigma$ et $\varepsilon_3$ — à leurs composés avec les quatre hom. $j_G, j_p, j_\sigma, j_\mu$ — quitte à vérifier que le produit des quatre valeurs sur l'image d'un même $\varepsilon \in \mu_2$ est trivial. On aura :

$$\Theta_G j_G : S G \hookrightarrow G , \text{ inclusion dans (136)}$$
(138)
$$\Theta_G j_p, \Theta_G j_\sigma \quad \text{triviaux}$$
$$\Theta_G j_\mu : \mathbb{G}_m \longrightarrow \mathbb{G}_m \text{ id. can.}$$

(conditions de compatibilité satisfaites !)

\newpage

%% ===== Page 450 =====

(139)
$\begin{cases}
\Theta_P j_G, \Theta_P j_\sigma \text{ triviaux} \\
\Theta_P j_\rho : S' T_P \hookrightarrow T_P \text{ inclusion} \\
\Theta_P j_\mu : G_m \xrightarrow{i_P} T_P
\end{cases}$

compatibilités triviales,

(140)
$\begin{cases}
\Theta_\sigma j_G, \Theta_\sigma j_\rho \text{ triviaux} \\
\Theta_\sigma j_\sigma : S N'_\sigma \hookrightarrow N'_\sigma \text{ inclusion} \\
\Theta_\sigma j_\mu : G_m \xrightarrow{i_\sigma} N'_\sigma
\end{cases}$

compatibilités triviales,

(141)
$\begin{cases}
\varepsilon_3 j_G, \varepsilon_3 j_\sigma, \varepsilon_3 j_\mu \text{ triviaux} \\
\varepsilon_3 j_\rho : S' T_P \xrightarrow{\delta_P} \mu_2
\end{cases}$

On définit donc, en termes de (134)

(142)
$$
\begin{tikzcd}
M_{\rho, \sigma} \ar[r, "\varepsilon_2"] \ar[d, "\Theta_\sigma"'] & \mu_2 \\
N'_\sigma \ar[ur, "\varepsilon_\sigma"'] &
\end{tikzcd}
\quad \varepsilon_2 = \varepsilon_\sigma \circ \Theta_\sigma
$$

$$
\begin{tikzcd}
M_{\rho, \sigma} \ar[r, "\chi"] \ar[d, "\Theta_G"'] & G_m \\
G \ar[ur, "\delta_G=\det"'] &
\end{tikzcd}
\quad \chi = \delta_G \circ \Theta_G
$$

$$
\begin{tikzcd}
M_{\rho, \sigma} \ar[r, "\widetilde{G}"] \ar[d, "\Theta_G"'] & P G \\
G \ar[ur, "can"'] &
\end{tikzcd}
\quad \widetilde{G} = can \circ \Theta_G
$$

on aura donc
(143)
$$
M_{\rho, \sigma} \xrightarrow{(\chi, \varepsilon_3, \varepsilon_2)} G_m \times \mu_2 \times \mu_2
$$
) épimorphisme

\newpage

%% ===== Page 451 =====

et les autres deux composés $\delta_p \Theta_p$ et $\delta'_0 \Theta_\sigma$
de (134) s'explicitent donc par les formules

$$
(144)
\begin{cases}
\delta'_0 \Theta_\sigma = \chi \quad (\overset{dfn}{=} \delta_G \Theta_G) \\
\delta_p \Theta_p(g) = \varepsilon_3(g) \chi(g)
\end{cases}
$$

En termes de (143) on définit

$$
(145)
\begin{cases}
\mathcal{M}_{p,\sigma}^+ = \text{Ker } (\mathcal{M}_{p,\sigma} \xrightarrow{\chi} G_m) \\
\mathcal{M}_{p,\sigma}^\circ = \text{Ker } (\mathcal{M}_{p,\sigma} \xrightarrow{(\varepsilon_3, \varepsilon_2)} \mu_2 \times \mu_2) \\
\mathcal{M}_{p,\sigma}^{+\circ} = \mathcal{M}_{p,\sigma}^+ \cap \mathcal{M}_{p,\sigma}^\circ = \text{Ker } (\mathcal{M}_{p,\sigma} \xrightarrow{(\chi, \varepsilon_3, \varepsilon_2)} G_m \times \mu_2 \times \mu_2)
\end{cases}
$$

on définit

$$
(146)
\begin{cases}
R_{p,\sigma} = \text{Ker } (\mathcal{M}_{p,\sigma} \xrightarrow{\tilde{\Theta}} P G) \\
R_{p,\sigma}^\circ = \mathcal{M}_{p,\sigma}^\circ \cap R_{p,\sigma} = \text{Ker } (\mathcal{M}_{p,\sigma}^\circ \xrightarrow{\tilde{\Theta}^\circ} P G) \\
S R_{p,\sigma} = \text{Ker } (\mathcal{M}_{p,\sigma} \xrightarrow{\Theta} G) \\
S R_{p,\sigma}^{\circ+} = \text{Ker } (\mathcal{M}_{p,\sigma} \xrightarrow{(\Theta, \varepsilon_3, \varepsilon_2)} G \times \mu_2 \times \mu_2) \\
\quad = S R_{p,\sigma} \cap R_{p,\sigma}^\circ
\end{cases}
$$

d'où les carrés cartésiens de caractérisation

(94), et plus [unclear: richement], les diagrammes
[crossed out: paradigmatiques]
analogues (95) (96) (116) (117) , et les

dévissages

$$
(147)
\begin{cases}
R_{p,\sigma} \simeq S' T_p \times S N'_\sigma \times G_m \\
R_{p,\sigma}^\circ \simeq S T_p \times S T_\sigma \times G_m \\
S R_{p,\sigma} \simeq S' T_p \times S N'_\sigma \times 1 \\
S R_{p,\sigma}^\circ \simeq S T_p \times S T_\sigma \times 1
\end{cases}
$$

[MARGIN: dont le sens est évident sur le diagr. (134)]

\newpage

%% ===== Page 452 =====

(148)
\begin{tikzcd}
1 \ar[r] & \mu_2 \ar[r] \ar[d, equal] & R_{\rho,\sigma} \times S G \ar[r] \ar[d, hook, "\mu_2 \times \mu_2"'] & M_{\rho,\sigma} \ar[r] \ar[d, hook, "\mu_2 \times \mu_2"'] & 1 \\
1 \ar[r] & \mu_2 \ar[r] & R_{\rho,\mathring{\sigma}} \times S G \ar[r] & M_{\rho,\mathring{\sigma}} \ar[r] & 1
\end{tikzcd}
[MARGIN: c'est ici la définition de $M_{\rho,\mathring{\sigma}}$]

$$
(149) \qquad
\begin{cases}
M_{\rho,\sigma}^+ \simeq S R_{\rho,\sigma} \times S G \\
M_{\rho,\mathring{\sigma}}^{+0} \simeq S R_{\rho,\mathring{\sigma}} \times S G
\end{cases}
$$

Enfin, définissons

$$
(150) \qquad
\begin{cases}
\mathcal{T}_{\rho,\sigma} \simeq M_{\rho,\sigma} / S G \\
\mathcal{T}_{\rho,\mathring{\sigma}} = M_{\rho,\mathring{\sigma}} / S G = \operatorname{Ker}(\mathcal{T}_{\rho,\sigma} \to \mu_2 \times \mu_2) \\
\mathcal{T}_{\rho,\sigma}^+ = M_{\rho,\sigma}^+ / S G = \operatorname{Ker}(\mathcal{T}_{\rho,\sigma} \to \mathbb{G}_m) \\
\mathcal{T}_{\rho,\mathring{\sigma}}^{+0} = M_{\rho,\mathring{\sigma}}^{+0} / S G = \operatorname{Ker}(\mathcal{T}_{\rho,\sigma} \to \mathbb{G}_m \times \mu_2 \times \mu_2) \\
\quad = \mathcal{T}_{\rho,\sigma}^+ \cap \mathcal{T}_{\rho,\mathring{\sigma}} \quad \text{dans } \mathcal{T}_{\rho,\sigma}
\end{cases}
$$

d'où des structures d'extensions

$$ 1 \to S G \to M_{\rho,\sigma} \to \mathcal{T}_{\rho,\sigma} \to 1 $$

etc (idem avec 
$M_{\rho,\mathring{\sigma}}, \mathcal{T}_{\rho,\mathring{\sigma}}$
$M_{\rho,\sigma}^+, \mathcal{T}_{\rho,\sigma}^+$
$M_{\rho,\mathring{\sigma}}^{+0}, \mathcal{T}_{\rho,\mathring{\sigma}}^{+0}$)

et des suites exactes

$$ 1 \to \mathcal{T}_{\rho,\sigma}^+ \to \mathcal{T}_{\rho,\sigma} \to \mathbb{G}_m \to 1 $$

etc (idem pour 
$\mathcal{T}_{\rho,\mathring{\sigma}} \subset \mathcal{T}_{\rho,\sigma}$
$\mathcal{T}_{\rho,\mathring{\sigma}}^{+0} \subset \mathcal{T}_{\rho,\mathring{\sigma}}$
$\mathcal{T}_{\rho,\sigma}^{+0} \subset \mathcal{T}_{\rho,\sigma}^+$
$\mathcal{T}_{\rho,\mathring{\sigma}}^{+0} \subset \mathcal{T}_{\rho,\sigma}$)

\newpage

%% ===== Page 453 =====

529

enfin des iso
$$ (151) \begin{cases}
R_{\rho,\sigma} / \mu_2 \simeq \tau_{\rho,\sigma} \\
R_{\rho,\sigma}^{\circ} / \mu_2 \simeq \tau_{\rho,\sigma}^{\circ} \\
S R_{\rho,\sigma} \simeq \tau_{\rho,\sigma}^{+} \\
S R_{\rho,\sigma}^{\circ} \simeq \tau_{\rho,\sigma}^{+\circ}
\end{cases} $$

80
Si $\underline{U}$ est un élément de $\mathcal{M}_{\rho,\sigma}$, on lui
associe ses images par $\theta_{\underline{a}}, \theta_{\rho}, \theta_{\sigma}$, et par
$\varepsilon, \varepsilon_1, \varepsilon_2$, on les notera respectivement

[MARGIN: Il faillait expliciter que les conditions (152)' conduisent $\mathcal{M}_{\rho,\sigma}$ comme un p. fibré principal sur $G \times \Gamma_{\underline{a}} \times \mathcal{N}_{\sigma}$, ... il faudrait le vérifier... mais peut on le faire ici ??? (153)' ne suffit pas vrai-]

$$ (152) \begin{cases}
u \in \Gamma_{\underline{a}} \\
\underline{a} \in \Gamma_{\rho} \\
\underline{b} \in \Gamma_{\sigma} \\
(\varepsilon, \varepsilon' \in \Gamma_{ph}) \\
\mu \in \Gamma_{\underline{a}}
\end{cases} \quad \text{satisfaisant} \quad (152)' \begin{cases}
\mu = \delta_{\underline{a}}(u) = \delta_{\rho}(\underline{a}) = \varepsilon \delta_{\sigma}(\underline{b}) \\
\varepsilon' = \varepsilon_{\sigma}(\underline{b})
\end{cases} $$
donc ces quantités sont
déterminées par la seule
conn. de $u, \underline{a}, \underline{b}$

au lieu des notations des § 45 et suivants, on reprendra, pour l'optique de la présente section, des éléments notés

$$ u, \alpha, \beta, \alpha_0, \beta_0, f \quad \text{avec} \quad \begin{cases} u \in \Gamma_{\underline{a}} \subset \Gamma_G \\ \alpha, \alpha_0 \in \Gamma G^{\pm 1} = \Gamma_G \\ \beta, \beta_0, f \in \Gamma S_G \end{cases} $$

[MARGIN: NB des cond (152)' ou (153) on déduit des identités autres... On pourrait introduire $g = \alpha \rho(\alpha_0)^{-1} (= a \rho(a)^{-1})$, $h = \beta \sigma(\beta_0)^{-1}$ mais je ne vois pas à quoi ça sert]

(153) reliés par des relations
$$ \begin{cases}
\alpha = u \underline{a} \\
\beta = u \underline{b}
\end{cases} \text{ i.e. } u = \alpha \underline{a}^{-} = \beta \underline{b}^{-} \qquad \text{[MARGIN: formules dans } G \text{]} $$
$$ \underline{U} = f^{-1} u $$
$$ \alpha_0 = f^{-1} \alpha \text{ i.e. } \alpha = f \alpha_0 $$
$$ \beta_0 = f^{-1} \beta \text{ i.e. } \beta = f \beta_0 $$

\newpage

%% ===== Page 454 =====

formules qui ont un sens seulement si
on suppose donnés des plongements (ou
du moins morphismes)

(155) $$ \begin{cases} \Gamma_p \to G \\ \Gamma_\sigma \to G \end{cases} \quad \text{et un sous-groupe} \quad \Gamma_0 \hookrightarrow G $$

Quand on a de tels plongements, les formules
(154) ont un sens, et définissent les
(pour les quantités (152) connues) les
quantités (153) en termes de l'une d'elles,
disons $\underline{u}$, ou $\underline{f}$, ou $\underline{\alpha}$, ou $\underline{\beta}$ (mais
pas en termes de $\underline{\alpha}_0$ ou de $\underline{\beta}_0$).

En fait, on dispose d'un sous-groupe
(156) $$ \Gamma_0 \subset G $$
tel qu'on ait $$ \begin{cases} \Gamma_0 \xrightarrow{\delta|\Gamma_0} G_m \\ \text{i.e. } G \simeq \Gamma_0 \cdot SG \\ \text{produit 1/2 direct} \end{cases} $$
pour $\underline{u} \in \Gamma_G$,
de sorte que (la [unclear] il $\exists$ un
unique $f \in SG$ tel que
$$ \underline{U} = f^{-1} \underline{u} $$

[MARGIN: 453]
d'où des uniques $\underline{\alpha}, \underline{\beta}, \underline{\alpha}_0, \underline{\beta}_0 \in \Gamma_G$ satisfaisant
à ( [unclear: voir f. précédente] ) i.e. (154) :
(157) $\underline{\alpha} = \underline{u}\underline{a}^{-1}, \quad \underline{\beta} = \underline{u}\underline{b}^{-1}, \quad \underline{\alpha}_0 = f^{-1}\underline{\alpha} = f^{-1}\underline{u}\underline{a}^{-1}, \quad \underline{\beta}_0 = f^{-1}\underline{\beta} = f^{-1}\underline{u}\underline{b}^{-1}$

\newpage

%% ===== Page 455 =====

On aura d'ailleurs
$$det u = det \tilde{u} = \mu$$
$$det \alpha = det (u) det (\underline{a})^{-1} = \varepsilon' \qquad par (152')$$
$$det \beta = det (u) det (\underline{b})^{-1} = 1$$
et $$det \alpha_0 = det \alpha = \varepsilon'$$
$$det \beta_0 = det \beta = 1$$

On aura donc bien $$\alpha, \alpha_0 \in \Gamma G^{\pm 1} \overset{def}{=} \{g \mid (det g)^2 = 1\}$$
$$\beta, \beta_0 \in \Gamma S G$$

- mais on a admis, bien sûr, que les
hom. $\delta_\Gamma, \delta_\sigma$ de (155) sont induits, via les
plongements (155), par le déterminant $\delta_G$.

Pour parvenir à la relation ---
D'équations
(157') $$[\alpha, \beta] = \varepsilon' [ \beta, \sigma] \varepsilon_1 \rho(\mu-1) ,$$
il faut se donner des sections $\rho, \sigma$ de $G$
(ou plus précisément des section resp. de $T_\rho$ et de $T_\sigma$
eux-mêmes), et un unipotent strict $\varepsilon_1 \in \Gamma S G$,
pour que la relation (157) prenne un sens.
Remplaçant $\alpha, \beta$ par leurs valeurs $\underline{u} \underline{a}^{-1}, \underline{u} \underline{b}^{-1}$,
la relation devient (utilisant la commutativité
de $T_\rho$, chose qui n'avait pas même
fait [unclear: necessité] jusqu'ici)
(*) $$[ \underline{u}, \rho] = \varepsilon' \underline{u} ( \underline{b}^{-1} \sigma \underline{b} )^{\underline{u}} \sigma (\underline{u})^{-1} \varepsilon_1 \rho (\mu-1)$$

[MARGIN: 454]

\newpage

%% ===== Page 456 =====

et résultera de l'ensemble des deux conditions

$$ (158) \quad
\begin{cases}
a) & b^{-1} \sigma(b) = \varepsilon_\sigma(b) \text{ identiquement en } b \in \Gamma N_0^1 \\
   & \text{i.e. } \quad b(\sigma) = \varepsilon_\sigma(b) . \sigma \text{ identiquement} \\
b) & [u, \rho] = [u, \sigma] \varepsilon_1(\mu-1) \text{ identiquement en } u \in \Gamma \tilde{T}_0 \\
   & \qquad \qquad \qquad \qquad \qquad \text{pour } \mu = \delta_{\tilde{T}_0}(u) = \det(u) \\
   & \text{i.e. } \rho(u) = \varepsilon_1(1-\mu) \sigma(u) \text{ identiquement en } u
\end{cases}
$$

ce qui s'écrit aussi (en appliquant $\sigma^{-1}$ aux deux mb.)
$$ \sigma^{-1}\rho(u) = \varepsilon_0(1-\mu) . u \qquad \text{où} \qquad \varepsilon_0 = \sigma^{-1}(\varepsilon_1) $$
et n'est guère commode que si $T_0$ normalise $L_0$, et si $\sigma^{-1}\rho \in \Gamma T_0$. $L_0$ est de la forme

$$ (159) \quad 
\begin{cases} 
a) & \sigma^{-1}\rho = \varepsilon_0 u_0 \qquad u_0 \in \Gamma T_0 \\ 
b) & T_0 \text{ opère sur } L_0 \text{ par la connexion induite} \\ 
   & \text{par le déterminant} 
\end{cases} 
$$

[unclear: Ça l'a quand même un air de le voyage par une condition du voisinage ainsi qu'en des VII ?)]

Mais ici l'équation 157) venait un peu comme un cheveu sur la soupe - il doit y avoir des façons plus convaincantes et expéditives de s'en convaincre par un [unclear: retour] sur des formules

[MARGIN: attention au signe $\varepsilon$]
$$ (160) \quad 
\begin{cases}
U(\rho) = \text{int}_{...}(\rho) \quad ( = \rho ( ) \rho^{-1} ) & \text{i.e. } [U, \rho] = \dots \\
U(\sigma) = \varepsilon \, \text{int}_{...}(\sigma) \quad ( = \sigma ( ) \sigma^{-1} ) & \qquad [U, \sigma] = \dots \\
U(\varepsilon_0) = \text{int}_{...} (\varepsilon_0)
\end{cases} 
$$

La troisième relation résulte $(159 \, b)$ de $U= f^{-1} u$, et du fait que par hypothèse $u_0 \in \Gamma T_0$ opère sur $L_0$ (identiquement par la multiplication $\mu = \det(u)$).

\newpage

%% ===== Page 457 =====

Les deux autres formules s'explicitent successive-
ment, en écrivant $\alpha = u a^{-1}$, $\beta = u b^{-1}$, $\eta = f^{-1} u$

$\alpha_0 = f^{-1} \alpha$, $\beta_0 = f^{-1} \beta$

(161) $\begin{cases} u(\rho) = int(\alpha)(\rho) \\ u(\sigma) = \varepsilon_1 int(\beta)(\sigma) \end{cases}$ où $f$ ne figure plus

i.e. $u(\rho) = int(u)(\rho)$ puisque $int(a^{-1})\rho = \rho$
ok ! (car dans $\Gamma_\rho$ comm. )

$u(\sigma) = \varepsilon_1 int(u)(\underbrace{b^{-1}(\sigma)}_{= \varepsilon_0}) = int(u) \sigma$ OK !
cf (158)(a)

Pour prouver (160) on n'a utilisé que (158 a)
et (159) (b). Essayons d'en déduire mainte-
nant la "relation de cocycle" (157) - ce doit être
un calcul formel, que j'ai déjà dû faire une
douzaine de fois, dans des contextes qui ne
différaient au fond en rien ! On l'obtiendra
en écrivant les relations (160) en la forme
équivalente (161) et

(161 bis) $u(\varepsilon_0) = \varepsilon_0(\mu)$ i.e. $\underbrace{[u, \varepsilon_0]}_{u \varepsilon_0 u^{-1}} = \varepsilon_0(\mu^{-1})$ soit

Or $[u, \varepsilon_0] = u \underbrace{\varepsilon_0(u)^{-1}}_{= \sigma^{-1} \rho(u)}$
i.e. $\sigma([u, \varepsilon_0]) = \sigma(u) \rho(u)^{-1}$

Or $\begin{cases} \rho(u^{-1}) = u^{-1} \xi \\ \rho(u) = \xi^{-1} u \\ \sigma(u) = \eta^{-1} u \end{cases}$ $dfn \quad \xi = \alpha \rho(\alpha)^{-1}$, $dfn \quad \eta = \varepsilon \beta \sigma(\beta)^{-1}$

d'où $\sigma([u, \varepsilon_0]) = \eta^{-1} \xi$

et $(*)$ devient, [unclear]

$\eta^{-1} \xi = \varepsilon_1(\mu^{-1})$ i.e. $\xi = \eta \varepsilon_1(\mu^{-1})$ i.e. (157)
OK !

\newpage

%% ===== Page 458 =====

En résumé :

Proposition Supposons donnés, en plus des données
de la section précédente [unclear: $7^\circ$ $§ 26'$] ff 3, notamment
$(135) (136)$, [unclear: une section de $\Gamma_{g, r, \nu}$] ceci

(162)
a) sections $\rho \in \Gamma_{T_{\rho}}$, $\sigma \in \Gamma_{T_{\sigma}}$ sections centrales
b) Plongements $T_{\rho} \hookrightarrow G$, $T_{\sigma} \hookrightarrow G$
c) tore $T_0 \subset G$ et un $\delta_G : T_0 \xrightarrow{\sim} G_m$
i.e. $G \simeq T_0 \cdot SG$ (semi-direct)
d) $\varepsilon_0 \in \Gamma_G$ unipotent strictement régulier, définissant
$\lambda \mapsto \varepsilon_0(\lambda) : G_m \hookrightarrow SG$

et supposons [unclear: les conditions de compatibilité entre $T_{\rho}, T_{\sigma}$] :
(163)
a) $\delta_G | T_{\rho} = \delta_{\rho}$, $\delta_G | T_{\sigma} = \delta_{\sigma}$
b) $int(\underline{\nu})(\sigma) = \varepsilon_0(\underline{\nu})\sigma$ dans $G$ [unclear: identique] pour $\underline{\nu} \in \Gamma_{\tilde{N}}$
c) $T_0$ normalise $L_0$, et on a
[MARGIN: relié $L_0, T_0$] $u(\varepsilon_0(\lambda)) = \varepsilon_0(\lambda^{\mu})$ où $\mu = \delta_G(u)$
[MARGIN: relié $L_0, T_0 / d)$ $\rho, \sigma$] $\sigma^{-1} \rho \in [unclear]$ i.e. $\exists loc.$ scalaire $\lambda \in \Gamma_{G_m}$
[unclear: t.q.] $\lambda_0 \sigma \rho^{-1} = \varepsilon_0 u_0$, avec $u_0 \in \Gamma_{T_0}$

Sous ces conditions, on a pour le système
$(152)$ satisfaisant les relations $(152')$ (i.e. on
fait un [unclear] section de $\mathcal{M}_{g,r}$ !) un
unique syst. $(153)$, satisfaisant les relations
$(154)$ - donnés Un mot de façon unique

sous la forme
(164) $U = f u$, $f \in \Gamma_{SG}$, $u \in \Gamma_{T_0}$
[MARGIN: et on a posé]
(165) $\alpha = u a^{-1}$, $\beta = u b^{-1}$, $\alpha_0 = f^{-1} \alpha = f^{-1} u a^{-1}$, $\beta_0 = f^{-1} \beta = f^{-1} u b^{-1}$.

\newpage

%% ===== Page 459 =====

Tout ceci posé, on a alors les formules
$(160)$, $(161)$, $(161 bis)$, et la formule "des bouts" $(157)$.

Nous nous posons maintenant la question
sous quelles conditions la connaissance des
$(\alpha, \beta, f)$ (pris parmi les quantités $(155)$,
provenant des récepteurs des quantités $(151)$ -
i.e. une section de $M_{g,\nu}$ - et quelles condi-
tions il faut imposer à $(\alpha, \beta, f)$. Il
suffirait (pour l'unicité) de pouvoir construire $\underline{u}$, d'où
$U$, $\underline{a}$, $\underline{b}$ par les trois premières formules
$(154)$. Il faudra
donc trouver $\underline{a} \in \Gamma T_{\rho}$ tel que $\alpha \underline{a}$ soit section
de $T_0$, - indétermination sur $\underline{a}$ est ds
$\Gamma T_0 \cap T_{\rho}$, et $\underline{b} \in \Gamma N'_0$ tel que $\beta \underline{b}$ soit section
de $T_0$ - indétermination ds $\Gamma T_0 \cap N'_0$.
L'indétermination sur $\underline{u}$ est la multiplication
à droite par une section de $T_0 \cap T_{\rho} \cap N'_0$, il
suffira donc qu'on ait

$$(164) \qquad T_0 \cap T_{\rho} \cap N'_0 = \{1\}$$

[MARGIN: en fait, on a $\underline{\underline{même}}$ $T_0 \cap T_{\rho} = \{1\}$ dans les cas envisagés tou... ds II, III et suite ...]

Supposant cette condition satisfaite, quelles
conditions a-t-on sur $(\alpha, \beta, f)$ pour provenir
d'un syst $(152)$ - i.e. d'une section de $M_{g,\nu}$ ?

\newpage

%% ===== Page 460 =====

On voit que si une section $\varphi_0$ de $\underline{M}_{g,\Gamma}$, 
[unclear: se projette] correspond à $\alpha_0, \beta_0, f_0$, et [unclear: une section]
$\varphi_1 = f^* \varphi_0$, correspondant à $\alpha_1, \beta_1, f_1$, alors

(165) $\quad \alpha_1 = \alpha_0$, $\beta_1 = \beta_0$, $f_1 = f_0 f$ ,

On voit que si $\varphi_0$ correspond : $(\underline{U}_0, \underline{a}_0, \underline{b}_0)$,
alors $\varphi_1 : (\underline{U}_1, \underline{a}_1, \underline{b}_1)$, alors on a :

(166) $\quad \underline{U}_1 = f^{-1} \underline{U}_0$, $\underline{a}_1 = \underline{a}_0$, $\underline{b}_1 = \underline{b}_0$

(car $\Theta_{\underline{a}} | S G = $ relation canonique $S e \subset G$,
et $\Theta_{\underline{a}}$ trivialisé sur $S G$). Donc on tire
[crossed out line]
$u_1 = u_0 \quad$ (car $\det \underline{U}_1 = \det \underline{U}_0$)

d'où (165). Cela montre que la condition
cherchée sur $(\alpha, \beta, f)$ ne dépend que de
$\alpha, \beta$ pas de $f$. On voudrait qu'elle
s'exprime par la validité de la formule en
(157), aux paramètres $\varepsilon_j$ près ($\varepsilon_j^2 = 1$)
(qui sera d'ailleurs déterminé - et pour cette
détermination [unclear] ... au niveau local...)

[MARGIN: à vrai dire, cette condition est suffisante - car les automorphismes de [unclear] $\varphi_0 \in T_0$ (qui rev. donc [unclear]) (décrit par $(\alpha, \beta)$) et il faut exiger que $u \in B =$ Aut de la [unclear] ]

Il suffira de cette condition pour
l'objet des relations de (II) (III),
ces conditions (plus draconiennes que celles
prop. précédente (p. 531') qui nous
importent au triplet $(\varepsilon_0, f, \sigma)$. Il

\newpage

%% ===== Page 461 =====

533

est possible que les conditions auxquelles nous
sommes parvenus étaient plus restrictives
que nécessaires, et que je sois obligé
d'y revenir par la suite, dans un contexte
plus général que celui des [unclear] des
[unclear: $\operatorname{Gl}(1), PGl(1), \operatorname{Sl}(2)$], mais pour le moment
l'étude détaillée faite en nous inspirant
des hypothèses dégagées dans II, IV nous
servira de déblayage, pour nous éclairer
les idées sur un cas particulier
important.

Rem. Rappelons que nous avons vu qu'en
fait $p$ est déjà défini en termes de $\alpha$ seuls
il se tire de $\alpha p(\alpha)^{-1} = \xi p$
Mais il n'est déterminé en termes de $p$
— en fait, de $p(\gamma p)^{-1}$ — que modulo
la connaissance de $\varepsilon$ (en pratique, modulo
signe...). Mais si l'on impose (p.ex.) $\varepsilon=1$,
la donnée de $\alpha$ ne détermine $p$ que
modulo $p \mapsto p \gamma_0$ ...

\newpage

%% ===== Page 462 =====

La structure de $\mathcal{M}_{p,\sigma}$ étant bien comprise (cf p. 495),
revenons sur le groupe $\mathcal{S}\mathcal{M}_{p,\sigma}$, considéré comme
produit semi-direct de $\mathcal{M}_{p,\sigma}$ par le sous-groupe
$L_0$ de $SG$, sur lequel en fait opère $\mathcal{M}_{p,\sigma}$
via $\chi$. Désignant par

(165) $$ \mathcal{E}_{L_0} \simeq \operatorname{Ker}(S\mathcal{M}_{p,\sigma} \to \mathcal{M}_{p,\sigma}) $$

correspondant à $\varepsilon_0$, [unclear: est] défini par la condition

(166) $$ u(\mathcal{E}_{L_0}(\lambda)) = \mathcal{E}_{L_0}(\mu\lambda) \quad - \quad \mu = \chi(u). $$

[crossed out: lorsque $u$ ...]

\newpage

%% ===== Page 463 =====

90 Les sous-groupes $M_{g,r}[n]$, et le groupe $S M_{g,r}$, liaison.
Le groupe $M_{g,r}$ est extension de $\bar{\Gamma}_{g,r}$ par
$\pi_{g,r} \simeq \widehat{SG}$, et les trois sous-groupes qu'on a vus
s'explicitent au moy. de sections, partielles --- le
[unclear: grâce à ce que] pour le troisième, on se donne une
section partielle seulement, au dessus du
sous-groupe $\bar{\Gamma}_{g,r}^+$ d'indice 2 de $\bar{\Gamma}_{g,r}$, noyau de
$\varepsilon_2$.

[[ Déterminons de façon générale les sections
de $M_{g,r}$ sur $\bar{\Gamma}_{g,r}$ . Comme

[MARGIN: finalement je ne m'en sers pas dans la suite...]

(1) $$R_{g,r} \xrightarrow{u} \bar{\Gamma}_{g,r}$$ épi de noyau prof. l.
$S' \Gamma_x \times S N_x \times G_m$

il revient au même de se donner une section,
ou un hom.

(2) $$R_{g,r} \xrightarrow{\Psi} M_{g,r}$$

qui soit
a) trivial sur le "noyau" (1)
b) [unclear and crossed out]

Les hom $\Psi$ qui satisfont à a) correspondent
aux scindages de l'extension triviale de
$R_{g,r}$ par $SG$ (extension produit), donc
s'identifient aux homomorphismes

(4) $$R_{g,r} \xrightarrow{\Psi_0} SG$$

\newpage

%% ===== Page 464 =====

en associant à $\psi^0$ l'hom.

(5) $x \mapsto \psi(x) = x \psi_0(x) = \psi^0(x) x$

Donc la condition a) prend la forme

(6) $\psi^0(x) = x$ sur $p_r = R_{\rho, \sigma} \cap SG$.

Les $\psi, \psi^0$ correspondants aux sections
$M_{\rho, \sigma}(0)$, $M_{\rho, \sigma}(h)$, $M_{\rho, \sigma}(-J)$ seront notés resp.t
$(\psi_0, \psi_0^0)$, $(\psi_h, \psi_h^0)$, $(\psi_{-J}, \psi_{-J}^0)$. ]]

10) Cas de $M_{\rho, \sigma}(0)$. Avec les notations
de la p. 529, il correspond au sous-groupe des
$M_{\rho, \sigma}$ satisfaisant à l'une des conditions
équivalentes $y = y_0$,

(7) $f = 1$, ($U \in \Gamma T_0$, $\alpha_0 = \alpha$, $\beta_0 = \beta$)

donc on a

(8) $M_{\rho, \sigma}(0) = \Theta^{-1}(T_0)$
(en termes de la description de $M_{\rho, \sigma}$ de la p. 526)

Vérifions que c'est un sous-groupe-section :
On a
$M_{\rho, \sigma}(0) \cap \widehat{G}_{\rho, \sigma}^+ = \{1\}$
car $\Theta(\widehat{G}_{\rho, \sigma}^+ = SG)$ est l'inclusion
canonique $SG \hookrightarrow G$, $\Rightarrow SG \cap T_0 = \{1\}$.
Reste à vérifier que
$M_{\rho, \sigma}(0) \longrightarrow \overline{G}_{\rho, \sigma}$
(qui est mono par a)) est épi. Mais on

\newpage

%% ===== Page 465 =====

535

termes de la relati. can.

(9) $$ \mathcal{M}_{\hat{p},\sigma} \simeq \underbrace{R_{\hat{p},\sigma} \times SG}_{\widetilde{\mathcal{M}}_{\hat{p},\sigma}} / \underbrace{\mu_2}_{\text{diag.}} $$

et de l'hom.

[MARGIN: $K : R_{\hat{p},\sigma} \to \widehat{G}_{un}$ l'hom. canonique d'après (10) provenant de $\Theta | R_{\hat{p},\sigma}$]

$$ i_{T_0} : \widehat{G}_{un} \buildrel \sim \over \longrightarrow T_0 \subset G $$
l'isom inverse de l'isom
$$ \delta_{T_0} = \det | T_0 : T_0 \buildrel \sim \over \longrightarrow \widehat{G}_{un} $$

on a, pour un élément $(\underset{R_{\hat{p},\sigma}}{\xi}, \underset{\widehat{SG}}{U_0})$ de $\widetilde{\mathcal{M}}_{\hat{p},\sigma}$
$\Theta(\xi, U_0) = K(\xi) U_0$, déterminant $K(\xi)^2$
dire que c'est $\in U$ signifie que l'on a
$$ K(\xi) U_0 = i_{T_0}(K(\xi)^2) \quad \text{i.e.} $$

(11) $$ U_0 = K(\xi)^{-1} i_{T_0}(K(\xi)^2) $$

donc l'image isom. $\mathcal{M}_{\hat{p},\sigma}(\infty)$ de $\mathcal{M}_{\hat{p},\sigma}(\omega)$
dans $\widetilde{\mathcal{M}}_{\hat{p},\sigma}$ est formée des

(12) $$ (\xi, K(\xi)^{-1} i_{T_0}(K(\xi)^2)) \in \Gamma \widetilde{\mathcal{M}}_{\hat{p},\sigma} $$

en d'autres termes cela correspond à une [unclear] section, avec

(13) $$ \psi_0^\circ(\xi) = K(\xi)^{-1} i_{T_0}(K(\xi)^2) . $$

Comme il se doit, c'est bien l'hom.
$$ \psi_0^\circ : R_{\hat{p},\sigma} \longrightarrow S G $$
et sa valeur sur $\mu_2$ est l'identité
car $K(\underline{\epsilon})^2 = 1$ et [unclear] $\psi_0^\circ(\underline{\epsilon}) = K(\underline{\epsilon})^{-1} = K(\underline{\epsilon})$

[unclear: Barré: Pour définir cette]
Remq Pour définir cette section, il suffisait

464

\newpage

%% ===== Page 466 =====

[unclear: C'est d'un pt de vu un peu abusif] de se donner un hom. quelconque, soit $\psi_0$, de $Q_{>0}$ dans $SG$, et de poser (13).

Avant de quitter $\mathcal{M}_{p,\sigma}(0)$, je vais introduire le groupe un peu plus grand (de dim 4, au lieu de 3)

(14) $\mathcal{M}_{p,\sigma}[0] \overset{df}{=} \Theta^{-1}(B'_0)$

où, je le rappelle, $B'_0$ est défini comme le produit 1/2 direct

(15) $B'_0 = T_0 \cdot L_0$

Cette fois-ci, $U = U_0 \kappa(\Gamma) \in \Gamma B'_0$ n'est pas déterminé par la seule connaissance de son déterminant $\kappa(U)^2$, qui [unclear: ne le détermine qu'à un modulo] $L_0$. Cela correspond au fait que $\mathcal{M}_{p,\sigma}[0] \simeq \widetilde{G}_{p,\sigma} \times^{B_0} B'_0 \subset \widetilde{SG}$ n'est plus ident. à (1), mais : $L_0$ y apparaît donc comme un sous-groupe invariant de $\mathcal{M}_{p,\sigma}[0]$, et on a (évidemment compte tenu que $\mathcal{M}_{p,\sigma}(0) \subset \mathcal{M}_{p,\sigma}[0]$ et $\mathcal{M}_{p,\sigma}(0) \overset{\sim}{\to} \widetilde{G}_{p,\sigma}$) par

(16) $\mathcal{M}_{p,\sigma}[0] = \mathcal{M}_{p,\sigma}(0) \cdot L_0$
produit 1/2 direct

On aura une construction qui en sera dès que l'on a $T_0$ comme dessus, et un sous-groupe [unclear] $L_0 \subset SG$,

\newpage

%% ===== Page 467 =====

536

invarié par $T_0$.

L'opération de $\mathcal{M}_{p,\sigma}[0]$ sur $S\hat{G}$ se fait par l'intermédiaire de $\theta$, définie par [unclear] :
$$ \mathcal{M}_{p,\sigma}[0] \overset{\theta | \mathcal{M}_{p,\sigma}[0]}{\longrightarrow} B'_0 \quad ; $$
cette opération étant précisée, on peut introduire

[MARGIN: peut prendre une définition plus symétrique cf p 538, (39)]
(17) $$ S\mathcal{M}_{p,\sigma} \overset{\text{déf}}{=} \mathcal{M}_{p,\sigma}[0] \underset{\text{produit 1/2 direct}}{\cdot} S\hat{G} $$

dans lequel $S\hat{G}$ est invariant, et $\mathcal{M}_{p,\sigma}[0]$ s'identifie à un sous-groupe de $S\mathcal{M}_{p,\sigma}$. Comme $\mathcal{M}_{p,\sigma}(0) \subset \mathcal{M}_{p,\sigma}[0]$, on a une inclusion

(18) $$ \underset{=\mathcal{M}_{p,\sigma}}{\mathcal{M}_{p,\sigma}(0) \cdot S\hat{G}} \subset \underset{=S\mathcal{M}_{p,\sigma}}{\mathcal{M}_{p,\sigma}[0] \cdot S\hat{G}} $$

dont le $1^{er}$ [unclear] $S\mathcal{M}_{p,\sigma}$ s'identifie :

$\mathcal{M}_{p,\sigma}$. [crossed out: Bien dit, le sous-groupe $L_0$ de $\mathcal{M}_{p,\sigma}[0]$ (cf (16)) opère trivialement sur $S\hat{G}$ i.e. $S\mathcal{M}_{p,\sigma}[0]$ qui est trivialise en fait sur $L_0$]

Pourtant je commence à m'égarer : j'y perds, je vais me raccrocher - pour une compréhension de la signification "géom." de $S\mathcal{M}_{p,\sigma}$, $\mathcal{M}_{p,\sigma}[0]$, aux notations de la p. 529, qui méritent d'être explicités par une

\newpage

%% ===== Page 468 =====

Proposition. Considérons l'hom.

$$ (19) \quad \underline{M}_{p,\sigma} \longrightarrow G \times \Gamma_p \times \hat{\Pi}'_\sigma \times \text{[unclear]} $$
$$ x \longmapsto (\Theta_G(x), \Theta_p(x), \Theta_\sigma(x), \mu(x), \varepsilon(x), \varepsilon'(x)) $$

C'est un monomorphisme, dont l'image est
formée des $\mu, \varepsilon, \varepsilon'$ [crossed out: sixtuplets] triplets $(U, \underline{a}, \underline{b})$ satisfaisant à

[crossed out: $\mu = \delta_G(U) = \delta_\sigma(\underline{b}) = \varepsilon' \delta_p(\underline{a})$]
$$ (20) \quad \begin{cases} \mu = \delta_G(U) = \delta_\sigma(\underline{b}) = \varepsilon' \delta_p(\underline{a}) \\ \varepsilon = \varepsilon_\sigma(\underline{b}) \end{cases} $$

Corollaire Le morphisme

$$ (21) \quad \underline{M}_{p,\sigma} \longrightarrow G \times \Gamma_p \times \hat{\Pi}'_\sigma $$
$$ x \longmapsto (\Theta_G(x), \Theta_p(x), \Theta_\sigma(x)) $$

est un monomorphisme, l'image est
formée des triplets $(U, \underline{a}, \underline{b})$ satisfaisant à
[MARGIN: NB on a $\begin{cases} \varepsilon = \varepsilon_\sigma(\underline{b}) \\ \varepsilon' = \delta_p(\underline{a}) / \delta_\sigma(\underline{b}) \\ \mu = \delta_\sigma(\underline{b}) = \delta_G(U) \end{cases}$]

$$ (22) \quad \begin{cases} \delta_G(U) = \delta_\sigma(\underline{b}) \\ (\delta_p(\underline{a}) / \delta_\sigma(\underline{b}))^2 = 1 \end{cases} $$

Dém. de la proposition : [crossed out: On a $\operatorname{Ker}(\dots)$]
$$ \operatorname{Ker}(\Theta_G, \varepsilon, \varepsilon') = SR_{p,\sigma} \simeq ST_p \times ST_\sigma $$
et $\Theta_p, \Theta_\sigma$ vus sur $SR_{p,\sigma}$ s'identifient
aux projections, donc
$$ \operatorname{Ker}((\Theta_p, \Theta_\sigma) | SR_{p,\sigma}) = \{1\}, \text{ donc} $$
$$ \operatorname{Ker}(\Theta_G, \mu, \varepsilon, \varepsilon', \Theta_p, \Theta_\sigma) = (1) $$

\newpage

%% ===== Page 469 =====

537

donc mono. Pour l'épi, on est ramené à la forme
équivalente du corollaire. [unclear: OPS] qu'on a un $\lambda$ tel
que
$$ \lambda^2 = \mu \overset{df}{=} det \underline{U} $$

[MARGIN: NB Où l'épi dans $\mu_{\rho,\sigma} = R_{\rho,\sigma} \times S \ell_2$ que j'écris $G \to \Theta_G$]

on veut écrire
$$ \underline{U} = \Theta(\tau, \underline{U}_0) = k(\tau) \underline{U}_0 $$
avec $\underline{U}_0 \in S \ell_2$, $k(\tau)^2 = det \underline{U}$ i.e. $k(\tau)^2 = \mu$,
on prend $\tau$ avec $k(\tau) = \lambda$,
et on pose donc
$$ \underline{U}_0 = \lambda^{-1} \underline{U}, \underline{a}_0 = \lambda^{-1} \underline{a}, \underline{b}_0 = \lambda^{-1} \underline{b} $$
et on est ramené à voir que $\underline{U}_0, \underline{a}_0, \underline{b}_0$ sont dans
l'image de $M_{\rho,\sigma}$ (on vient de "corriger"
$(\underline{U}, \underline{a}, \underline{b})$ par l'image dans $G \times T_\rho \times T_\sigma$ de
la section $\lambda$ de $G_{m} \hookrightarrow R_{\rho,\sigma}$) [unclear: c'est à dire ramené]
à voir que $(\underline{a}_0, \underline{b}_0)$ est dans l'image de
$$ R_{\rho,\sigma} \to T_\rho \times N'_\sigma. \quad \text{Mais maintenant} $$
$\underline{a}_0$ est dans $S'T_\rho$, $\underline{b}_0$ dans $S N'_\sigma$, et justement
$$ R_{\rho,\sigma} \simeq S'T_\rho \times S N'_\sigma \text{, ce qui gagne....} $$

Remarque On ne peut remplacer $G$ par $P G$ ;
i.e. $$ M_{\rho,\sigma} \to P G \times T_\rho \times N'_\sigma $$ (où on a projeté
encore $G$ [unclear] ) n'est pas [unclear]
le noyau est justement le sous-groupe des
$$ G_{\rho,\sigma}^+ = S \ell_2. $$

\newpage

%% ===== Page 470 =====

Considérons, maintenant que les quantités
$\underline{U}, \underline{a}, \underline{b}$ ($\tilde{\mu}, \varepsilon, \varepsilon'$) de p. 529 (151) sont connues,
les quantités (153) en loc. cit $\underline{u}, f, \alpha, \beta, \alpha_0, \beta_0$
($g, l, s, y$ n'y [unclear: figurent] pas en plus...) Elles sont déterminées,
en termes des $\underline{U}, \underline{a}, \underline{b}$, dès qu'on connait
soit $\underline{u}$, soit $f$, par les formules (154), $\underline{u}$ et
$f$ étant reliés par la seule relation

(23) $$ \underline{U} = f^{-1} \underline{u} $$

avec en plus sur $f$ la condition

(24) $$ f \in \Gamma \frac{S\underline{G}}{\underline{G}_{p,\sigma}^+} \ , $$

plus une condition sur $\underline{u}$ qui, dans loc. cit.,
s'écrivait

(25) $$ \underline{u} \in \Gamma T_0 \ , $$

et était suggérée par les calculs du §46 V
(repris ensuite dans §48 II, III etc). Cette
condition assure l'unicité de $f, \underline{u}$
satisfaisant à (23) - c'est cela, du pt.
de vue du calcul, sa principale utilité ! -
mais dans certains contextes il sera préférable
de la remplacer par la condition plus large

(26) $$ \underline{u} \in \Gamma B'_0 $$ 
[MARGIN: $B'_0 = T_0 . L_0$ (produit [unclear: s. direct])]

\newpage

%% ===== Page 471 =====

On n'a plus de résultat d'unicité pour l'écriture
de $\underline{U}$ sous forme (23) - par contre on a un
théorème d'existence sous forme simple :
Si on part de $(\alpha, \beta ; f)$ sections de $G$ telles
que
(27) $$(\det \alpha)^2 = \det \beta = \det f = 1$$ ,
pour qu'on puisse écrire ce triplet et par conséquent
les quantités qui s'en déduisent : $\alpha_0, \beta_0$,
(28) $$\alpha_0 = f^{-1} \alpha \ , \ \beta_0 = f^{-1} \sigma \beta \sigma$$
$$y = \alpha \rho(\alpha)^{-1} \ , \ y' = \beta \sigma \rho(\beta)^{-1}$$
$$g = \alpha_0 \rho(\alpha_0)^{-1} \ , \ h = \beta_0 \sigma \rho(\beta_0)^{-1}$$
proviennent d'un système $(x, \underline{U})$ ,
avec $x \in \Gamma \underline{M}_{\rho, \sigma}$, $\underline{U} \in \Gamma \underline{B}'_\theta$ tel que
(29) $$f \overset{dfn}{=} \underline{u} \overset{\sigma}{\underline{U}}^{-1} \in \Gamma \underline{S} \underline{G}$$ , i.e. $\det \underline{u} = \det \underline{U}$ ,
par les formules
(30) $$\alpha = \underline{u} \underline{a}^{-1} \ , \ \beta = \underline{u} \underline{b}^{-1} \ , \ f = \underline{u} \underline{U}^{-1}$$ ,
il faut et suffit que l'on ait la relation (qui
ne fait pas intervenir $f$, seulement $\alpha, \beta$) :
(31) $$\alpha \rho(\alpha)^{-1} = \varepsilon \beta \sigma(\beta)^{-1} \varepsilon_1 (\mu_{-1})$$ pour $\varepsilon \in \Gamma \mu_2$, $\mu \in \Gamma \mu_4$ convenables.

[MARGIN: CNB $\varepsilon, \mu$ sont univoquement déterminés par cette formule (i.e. $y=y'$), ou par $\alpha, \beta$).]

Alors le couple $(x, \underline{u})$ est unique. [unclear: Pour le fait]
[unclear: qu'on ait...] correspondance 1-1 entre les [unclear: $\rho, \sigma$-structures]
$\underline{M}_{\rho, \sigma}[\theta]$ (pour $\underline{M}_{\rho, \sigma}(\theta)$ !) sur $\underline{M}_{\rho, \sigma}$, et les [unclear]

\newpage

%% ===== Page 472 =====

des couples $\alpha, \beta$ satisfaisant [unclear] aux deux relations
$$ (\det \alpha)^2 = \det \beta = 1 $$ de (27)
et la relation [unclear] locale. Pour cela, on n'a
d'utilité pour le groupe $S \mathcal{M}_{p,\sigma}$.

Mais comme on a [unclear] [unclear], pour
interpréter la façon agissante sur les
[MARGIN: [unclear] (cf. plus bas (27) (31)),]
triplets $(\alpha, \beta, f)$ (avec la relation de structure de
groupe étrange dessus (sur laquelle il nous
faudra revenir en temps [unclear: utile] (plus généralement...))
où précisément les trois quantités $\alpha, \beta, f$
se mélangent intimement - il y a lieu
d'introduire $S \mathcal{M}_{p,\sigma}$ - qui sera isomorphe
à [unclear] les groupes ) des relations
consistant en (27) (31).

De ce point de vue, $S \mathcal{M}_{p,\sigma}$ se déduit
de $\mathcal{G}$ comme [unclear: lieu] des triplets
[unclear: ex: $(x, u, f)$]

(32) $$ S \mathcal{M}_{p,\sigma} = \text{schéma des triplets } (x, \underline{u}, f) $$
$$ \text{avec } \begin{cases} x \in \mathcal{M}_{p,\sigma} \\ \underline{u} \in B_{\sigma}' \\ f \in \tilde{S}_{\sigma} \end{cases} \quad f^{\underline{u}} \underset{\text{déf}}{=} f^{-1} $$

D'où
(33) $$ S \mathcal{M}_{p,\sigma} \hookrightarrow \mathcal{M}_{p,\sigma} \times B_{\sigma}' \hookrightarrow $$
$$ \simeq \{ x, \underline{u} \mid \gamma(x) = \frac{\text{det } \underline{u}}{\text{det } u_0} \} $$

ou encore, plus simple [unclear]

\newpage

%% ===== Page 473 =====

[MARGIN: 537]

(34)
$$S M_{p, \sigma} \overset{\sim}{\longrightarrow} M_{p, \sigma} \times_{G_m} B'_0$$
où $M_{p, \sigma} \overset{\chi}{\longrightarrow} G_m$
$B'_0 \overset{\delta_{B'} = \delta_G | B'_0}{\longrightarrow} G_m$
, ouf !

Cette formule va nous servir de définition de $S M_{p, \sigma}$ dorénavant - elle définit là-dessus une loi de groupe sur $S M_{p, \sigma}$, et on trouve sur $S M_{p, \sigma}$ deux structures d'extension

(35) $$1 \longrightarrow \underset{\underset{\operatorname{Ker}(B'_0 \overset{\delta_{B'}}{\to} G_m)}{"}}{L_0} \longrightarrow S M_{p, \sigma} \longrightarrow M_{p, \sigma} \longrightarrow 1$$
$$ \tilde{x} = (x, u) \mapsto x $$

(36) $$1 \longrightarrow \underset{\underset{\operatorname{Ker}(M_{p, \sigma} \overset{\chi}{\to} G_m)}{"}}{M_{p, \sigma}^{+}} \longrightarrow S M_{p, \sigma} \longrightarrow B'_0 \longrightarrow 1$$
$$ \underset{\underset{(x, u)}{"}}{\tilde{x}} \mapsto u $$

Le sous-groupe $S M_{p, \sigma}[0]$ de $S M_{p, \sigma}$ corres-
pond à la condition $f = 1$ i.e. $u = \underset{\underset{\theta_G(x)}{"}}{\underline{u}}$ (qui
implique déjà $\underset{\underset{\chi(x)}{"}}{\det \underline{u}} = \det u$ c'est-à-dire le
sous-groupe des $(x, u)$ dans $M_{p, \sigma} \times B'_0$ tels
que l'on ait $u = \theta_G(x)$. Donc $S M_{p, \sigma}[0]$ s'
identifie (cf (14) p. 536) [unclear: isomorphiquement] en un sous-groupe de $M_{p, \sigma}$ par $x \mapsto (x, \theta_G(x))$ :

(37) $$S M_{p, \sigma}[0] = Im (M_{p, \sigma}[0] \hookrightarrow S M_{p, \sigma})$$
$$ x \mapsto (x, \theta_G(x)) $$

$$(\text{où } M_{p, \sigma}[0] = \theta_G^{-1}(B'_0))$$

\newpage

%% ===== Page 474 =====

La donnée d'un scindage de (35) revient
à un scindage de
$$1 \longrightarrow L_0 \longrightarrow B'_0 \longrightarrow \mathbb{G}_m \longrightarrow 1$$
or il y en a un tout fait par le s.s.-groupe
$T_0$. Donc on en déduit un plongement de
$\mathcal{M}_{g, \sigma}$ dans $S\mathcal{M}_{g, \sigma}$ , qui donne un
produit 1/2-direct
(38) $$S\mathcal{M}_{g, \sigma} \simeq \mathcal{M}_{g, \sigma} \cdot L_0$$
$\mathcal{M}_{g, \sigma}$ opérant d'ailleurs sur $L_0$ via le
caractère $\chi : \mathcal{M}_{g, \sigma} \longrightarrow \mathbb{G}_m$ - mais
cette description me semble [unclear: au fond]
beaucoup moins "parlante" que (34) ,
car elle ne met pas en évidence l'homo-
morphisme can.
$$S\mathcal{M}_{g, \sigma} \longrightarrow B'_0$$
qui paraîtra comme la chose naturelle
pour l'introduction de $S\mathcal{M}_{g, \sigma}$. Une fois
ceci posé, on trouve (pour mémoire $^*$) le

Théorème bis. un isomorphisme can.
(39) $$S\mathcal{M}_{g, \sigma} \overset{\sim}{\longrightarrow} S\hat{\Gamma}_{g, \sigma}^*$$

[MARGIN: $^*$ groupe des systèmes $(a,b,f)$ satisfaisant (17) et (31)]

\newpage

%% ===== Page 475 =====

540

donné par

(40) $(x, \underline{u}) \longmapsto (\alpha, \beta, f)$ avec $\alpha = \underline{u} \underline{a}^{-1}, \beta = \underline{u} \underline{b}^{-1}, f = \underline{u} \underline{U}^{-1}$
(comme c'est bizarre !)

où $\underline{a} = \underline{\theta}_P(x), \underline{b} = \underline{\theta}_\sigma(x), \underline{U} = \underline{\theta}_G(x)$

L'opération [unclear: de transformation] $L_0$ dans $\Sigma_{P,\sigma}$ par

(41) $(\alpha, \beta, f) \longmapsto l \alpha, l \beta, l f$

[MARGIN: à gauche] correspond à l'opération de translation à gauche
du sous-groupe $L_0$ de $S\mathcal{M}_{P,\sigma}$ (cf (35)), four-
nie [unclear: par $(l,x) \dots l \in L_0$]. Donc on récupère
$\mathcal{M}_{P,\sigma} \simeq S\mathcal{M}_{P,\sigma}/L_0$ comme la section [unclear: quotient de] $\Sigma_{P,\sigma}^*$
par l'action (41) de $L_0$.

Remarque Dans cette description des $(\alpha, \beta, f)$
interviennent les quantités auxiliaires
$\underline{u}, \underline{a}, \underline{b}, \underline{U}$, ce qui conduit à considérer
le morphisme

(42) $$S\mathcal{M}_{P,\sigma} \longrightarrow B'_0 \times T_P \times B'_0 \times G$$
$$(x,\underline{u}) \longmapsto (\underline{u}, \underline{a}, \underline{b}, \underline{U}) \quad \text{avec } \begin{matrix} \underline{a} = \underline{\theta}_P(x) \\ \underline{b} = \underline{\theta}_\sigma(x) \\ \underline{U} = \underline{\theta}_G(x) \end{matrix}$$

On a à ce sujet :

Proposition L'hom. (42) est un isomorphisme sur
le fermé des systèmes satisfaisant

(43) $\det \underline{u} = \det \underline{b} = \det \underline{U} \text{ soit } \mu, \text{ et } \det \underline{a} / \mu \in \Gamma_{P/ [unclear]}$

C'est une conséquence immédiate des cor.
à la fin de p. 536', et de la définition
de $S\mathcal{M}_{P,\sigma}$ comme produit fibré.

\newpage

%% ===== Page 476 =====

10°) Les sous-groupes-schémas $M_{p,\sigma}(\sigma)$, $M_{p,\sigma}(p)$.

On définit $M_{p,\sigma}(\sigma)$ par la condition
$$ \beta_0 = 1 \quad i.e. \quad \beta = f \quad i.e. \text{ (cf (40)) }$$
(43)
$$ \underline{U} = \underline{b} $$
il est évident par prop. p. 536' que le sous-groupe

(44)
$$ M_{p,\sigma}(\sigma) \subset M_{p,\sigma} $$
$$ "=" $$
$$ \{(\underline{a}, \underline{b}, \underline{U}) \in \Gamma M_{p,\sigma} \mid \underline{U} = \underline{b} \} $$
forme une section - [unclear: pour $\underline{a}, \underline{b}$ d'un el. de $T_{p,\sigma}$]
[unclear: satisfaisant la condition]
$$ \delta_{\sigma}(\underline{a}) / \delta_{\sigma}(\underline{b}) \in \Gamma \mu_r \quad , $$
la façon la plus simple de le relever en
un élém de $M_{p,\sigma}$ n'est-il pas par $(\underline{a}, \underline{b}, \underline{U} = \underline{b})$!
(Et on aura bien $\delta_{\sigma}(\underline{U}) = \delta_{\sigma}(\underline{b})$ ). On définit
de même $M_{p,\sigma}(p)$ par $\alpha_0 = 1$ i.e. $\alpha = f$ i.e.

(45)
$$ M_{p,\sigma}(p) \subset M_{p,\sigma} $$
$$ "=" $$
$$ \{ (\underline{a}, \underline{b}, \underline{U}) \in M_{p,\sigma} \mid \underline{U} = \underline{a} \} $$
Ce qui fait une section partielle [unclear: au dessus du]
sous-groupe $T'_{p,\sigma}$ de $T_{p,\sigma}$ défini comme

(46)
$$ T'_{p,\sigma} \subset T_{p,\sigma} $$
$$ "=" $$
$$ \{ (\underline{a}, \underline{b}) \in \overline{T}_p \times \overline{N}'_{\sigma} \mid \delta_{\sigma}(\underline{a}) = \delta_{\sigma}(\underline{b}) \} $$
$$ "\text{ Ker } ( T_{p,\sigma} \xrightarrow{\varepsilon_g} \mu_r )" $$
donc

(47)
$$ T'_{p,\sigma} = \overline{T}_p \times_{\mathbb{G}_m} \overline{N}'_{\sigma} \quad , \quad \text{et on a} $$

\newpage

%% ===== Page 477 =====

541

donc

(48) $$ \tau_{p,\sigma} / \tau_{p,\sigma}^0 \overset{\varepsilon_2}{\longrightarrow} \mu $$

Rmq. Ces sous-groupes proviennent de
définition, et fournissent une section resp une
section [unclear: au dessus de] / $\tau_{p,\sigma}$, dès qu'on a (dans la
situation générale de 7., p. 526' ff) des
$p$-groupes $T_p \hookrightarrow \hat{G}$, $N\hat{\sigma} \hookrightarrow \hat{G}$, tels que
$\delta_{T_p}$, $\delta_{N\hat{\sigma}}$ soient induits par $\delta_G$.

Dans le cas modulaire, on écrit ainsi

(49) $$ \mathcal{M}_{p,\sigma}(1/2) = \mathcal{M}_{p,\sigma}(\sigma), \quad \mathcal{M}_{p,\sigma}(-1) = \mathcal{M}_{p,\sigma}(\rho) $$

et on définit plus généralement

(50) $$ \mathcal{M}_{p,\sigma}(q), \quad q \in \{ 0, 1, \infty ; 1/2, 2, -1 \} $$

à partir de $\mathcal{M}_{p,\sigma}(0)$ et $\mathcal{M}_{p,\sigma}(1/2)$ par

[MARGIN: Par ext pour tout $q$ utilisé un sous-groupe de $\mathcal{M}_{p,\sigma}$ [unclear]]

(51) $$ int(\rho)(\mathcal{M}_{p,\sigma}(q)) = \mathcal{M}_{p,\sigma}(\rho(q)) $$

$\rho$ agissant par la transformation homographique

(52) $$ \rho : q \longmapsto \frac{1}{1-q} $$

(dont l'itéré troisième est l'identité, tout comme
l'itéré troisième de $\tilde{\rho}$).

Il resterait à définir $\mathcal{M}_{p,\sigma}(-1)$, et comme
on a le choix à l'intérieur d'une classe de
$\widehat{\operatorname{Sl}(2, \mathbb{Z})}$-conjugaison bien déterminée, un choix plus

\newpage

%% ===== Page 478 =====

naturel serait entre

(53) $int(\sigma_i)(M_{p,\sigma}(0)) \qquad i \in \{0,1,\infty\}$

\newpage

%% ===== Page 479 =====

542

## XII Rétrospective.

Par bien des détours inutilement compliqués, j'ai fini de façon bien simple par construire deux schémas en groupes sur $\mathbb{Z}$

(1) $$ \mathcal{M}_{\rho, \sigma} \subset G \times T_\rho \times N'_\sigma \times G $$ [unclear: le dernier $G$ semble raturé]

(2) $$ S\mathcal{M}_{\rho, \sigma} \subset B'_0 \times T_\rho \times N'_\sigma \times \mathbb{G}_m $$

(3) $$ \mathcal{M}_{\rho, \sigma} = \{ (\underline{U}, \underline{a}, \underline{b}) \mid \det \underline{U} = \det \underline{b} = \varepsilon \det \underline{a}, \varepsilon^2 = 1 \} $$

(4) $$ S\mathcal{M}_{\rho, \sigma} = \{ (\underline{U}, \varepsilon, \underline{a}, \underline{b}) \mid \det \underline{U} = \det \underline{b} = \varepsilon \det \underline{a}, \varepsilon^2 = 1 \} $$
dans $S\mathcal{M}_{\rho, \sigma} = B'_0 \times \mathcal{M}_{\rho, \sigma}$

(5) $$ \mathcal{M}_{\rho, \sigma} \simeq S\mathcal{M}_{\rho, \sigma} / L_0 $$
où $L_0 \subset S\mathcal{M}_{\rho, \sigma}$
$$ \lambda \mapsto (\lambda, 1, 1, \lambda) $$

et
(6) $$ S\mathcal{M}_{\rho, \sigma} \simeq \mathcal{M}_{\rho, \sigma} \cdot L_0 $$
produit semi-direct [unclear: écrit "produit 1/2 direct"]

où
(7) $$ \mathcal{M}_{\rho, \sigma} \hookrightarrow S\mathcal{M}_{\rho, \sigma} $$
$$ x=(\underline{U}, \underline{a}, \underline{b}) \mapsto (\underline{U}, \underline{a}, \underline{b}, i_{T_0}(\mu)) \text{ où } \mu = \chi(x) $$

Ici (pour mémoire) dans le cas modulaire

(8) $$ \rho = \begin{pmatrix} 0 & 1 \\ -1 & 1 \end{pmatrix} \quad \sigma = \begin{pmatrix} 0 & 1 \\ -1 & 0 \end{pmatrix} $$

(9) $$ T_\rho = \{ (c\rho + d) \mid \det(c\rho + d) \text{ inv.} \} $$
$$ c^2 + cd + d^2 $$
$$ T_\sigma = \{ (t\sigma + u) \mid \det(t\sigma + u) \text{ inv.} \} $$
$$ t^2 + u^2 $$
[MARGIN:
$$ B'_0 = \left\{ \begin{pmatrix} \lambda & \nu \\ 0 & 1 \end{pmatrix} \mid \text{inv} \right\} $$
$$ T_0 = \left\{ \begin{pmatrix} \mu & 0 \\ 0 & 1/\mu \end{pmatrix} \mid \text{inv} \right\} $$
$$ i_{T_0}(\mu) = \begin{pmatrix} \mu & 0 \\ 0 & 1/\mu \end{pmatrix} $$
$$ L_0 = \left\{ \begin{pmatrix} \lambda & 0 \\ 0 & \lambda \end{pmatrix} \right\} $$
]

(10) $$ N'_\sigma = \{ \underline{b} \text{ dans } G \mid \underline{b}(\sigma) = \varepsilon \sigma, \varepsilon^2 = 1 \} $$
[unclear: $\varepsilon(\underline{b})$ écrit sous le $\varepsilon$]

d'où suite exacte
(11) $$ 1 \to T_\sigma \to N'_\sigma \xrightarrow{\delta_\sigma} \mu_2 \to 1 $$

478

\newpage

%% ===== Page 480 =====

et les trois hom.
(12)
$$ \mathcal{M}_{\rho,\sigma} \xrightarrow{\Theta_G} G , \mathcal{M}_{\rho,\sigma} \xrightarrow{\Theta_\rho} T_\rho , \mathcal{M}_{\rho,\sigma} \xrightarrow{\Theta_\sigma} N'_\sigma $$

donnés dans la définition a) (3) se complètent

par
(12 bis)
$$ \mathcal{M}_{\rho,\sigma} \xrightarrow{\chi = \det \circ \Theta_G} \mathbb{G}_m , \mathcal{M}_{\rho,\sigma} \xrightarrow{\varepsilon_2 = \varepsilon_\sigma \circ \Theta_\sigma} \mu_2 , \mathcal{M}_{\rho,\sigma} \xrightarrow{\varepsilon_3(x) = \frac{\det \underline{a}}{\det \underline{b}}} \mu_2 $$

de sorte qu'on a un [unclear]
(13)
$$ \mathcal{M}_{\rho,\sigma} \xhookrightarrow{(\Theta_G, \Theta_\rho, \Theta_\sigma, \chi, \varepsilon_2, \varepsilon_3)} G \times T_\rho \times N'_\sigma \times \mathbb{G}_m \times \mu_2 \times \mu_2 $$
$$ x \mapsto ( \underline{U}, \underline{a}, \underline{b}, \mu, \varepsilon, \varepsilon' ) $$

dont l'image est formé des systèmes satisfaisant

(14)
$$ \begin{cases} \det \underline{U} = \det \underline{b} = \varepsilon' \det \underline{a} = \mu \\ \varepsilon = \varepsilon_\sigma(\underline{b}) \text{ i.e. } \underline{b}(\sigma) = \varepsilon \sigma \text{ i.e. } \sigma \underset{^t\underline{b}^{-1}}{\underline{U}_\sigma} = \varepsilon \underline{b} \text{ i.e. } {^t\underline{b}}\underline{b} = \varepsilon \end{cases} $$
[MARGIN: dans (14), $\underline{U}_\sigma$ est raturé]

De même on a
(15)
$$ S\mathcal{M}_{\rho,\sigma} \hookrightarrow \underline{B}'_0 \times T_\rho \times N'_\sigma \times \underline{\underline{B}}_0 \times \mathbb{G}_m \times \mu_2 \times \mu_2 $$
$$ x \mapsto (\underline{u}, \underline{a}, \underline{b}, \underline{U}, \mu, \varepsilon, \varepsilon') $$

dont l'image est formé des systèmes satisfaisant

(16)
$$ \begin{cases} \mu = \det \underline{u} = \det \underline{U} = \det \underline{b} = \varepsilon' \det \underline{a} \\ \varepsilon = \varepsilon_\sigma(\underline{b}) \text{ i.e. } \underline{b}(\sigma) = \varepsilon(\sigma) \end{cases} $$

Soit d'autre part

(17)
$$ S\Sigma_{\rho,\sigma}^* \xrightarrow{\sim} G \times G \times G \times \mathbb{G}_a \times \mu_2 $$
$$ \Vert $$
$$ \left\{ (\alpha, \beta, f, \mu, \varepsilon) \mid (\det \alpha)^2 = \det \beta = \det f = 1 , \underbrace{\alpha \rho \alpha^{-1}}_{\xi} = \varepsilon \underbrace{\beta \sigma \beta^{-1}}_{\eta} \varepsilon_1(\mu-1) \right\} $$

où on a définit
(18)
$$ \varepsilon_1 = \begin{pmatrix} 1 & 0 \\ 1 & 1 \end{pmatrix} , \varepsilon_1(\mu) = \varepsilon_1^\mu = \begin{pmatrix} 1 & 0 \\ \mu & 1 \end{pmatrix} $$

\newpage

%% ===== Page 481 =====

543

on fait la projection
$$ (19) \quad \begin{cases} S\Sigma_{p,\sigma}^* \hookrightarrow G \times G \times G \\ y \longmapsto (\alpha, \beta, f) \end{cases} $$
est déjà un mono. On fait opérer $L_0$ [unclear text above]
$$ L_0 = \operatorname{Ker}(B'_0 \xrightarrow{det} \mathbb{G}_m) = \{ l \begin{pmatrix} * & * \\ 0 & 1 \end{pmatrix} \} \text{ sur } S\Sigma_{p,\sigma}^* $$
par
$$ (20) \quad (\alpha, \beta, f) \longmapsto (l\alpha, l\beta, lf) , $$
et on pose
$$ (21) \quad \Sigma_{p,\sigma}^* \overset{def}{=} L_0 \setminus S\Sigma_{p,\sigma}^* \quad . $$

Ceci posé, on a un isomorphisme
$$ (22) \quad S\mathcal{M}_{p,\sigma} \overset{\sim}{\longrightarrow} S\Sigma_{p,\sigma}^* $$
$$ (\underline{U}, \underline{a}, \underline{b}, \underline{V}) \longmapsto (\alpha, \beta, f) $$
$$ (\alpha = \underline{U}\underline{a}^{-1}, \beta = \underline{U}\underline{b}^{-1}, f = \underline{U}\underline{V}^{-1}) $$
transformant l'action de $L_0$ sur $S\mathcal{M}_{p,\sigma}$ portant sur $\underline{U}$ en l'action de $L_0$ décrite dans $(20)$, et déduit un iso
$$ (23) \quad \mathcal{M}_{p,\sigma} \overset{\sim}{\longrightarrow} \Sigma_{p,\sigma}^* $$
par passage au quotient.

On a des hom. canoniques
$$ (24) \quad S_0\mathcal{M}_{0,3}^{\sim} \longrightarrow S\mathcal{M}_{p,\sigma}(\hat{\mathbb{Z}})/\Sigma $$
$$ (25) \quad \mathcal{M}_{0,3}^{\sim} \longrightarrow \mathcal{M}_{p,\sigma}(\hat{\mathbb{Z}})/\Sigma $$
le deuxième déduit du premier par quotient, et coïncide [unclear] sur $\hat{G}_{0,3}$ avec [crossed out text].

\newpage

%% ===== Page 482 =====

l'hom composé § 42 IV

(26) $$ \hat{G}_{0,3}^+ \overset{\sim}{\longrightarrow} S\ell_2(\hat{\mathbb{Z}})' \hookrightarrow M_{\rho,\sigma}(\hat{\mathbb{Z}})/\Sigma $$
$$ \parallel $$
$$ SG(\hat{\mathbb{Z}})/SG(\hat{\mathbb{Z}}) \cap \Sigma $$

grâce à l'hom canonique

(27) $$ SG \hookrightarrow M_{\rho,\sigma} $$
$$ S\ell(2) \cong \left\{ (\underline{u}, \underline{a}, \underline{b}) \mid \underline{a} = \underline{b} = 1 \right\} $$

On a défini

(28) $$ \Sigma = (\mathbb{Z}/2\mathbb{Z})^3 \hookrightarrow M_{\rho,\sigma}(\mathbb{Z}) \subset M_{\rho,\sigma}(\hat{\mathbb{Z}}) $$
$$ \parallel $$
$$ (\mu_2)^3(\mathbb{Z}) $$

via l'hom canonique (central)

(29) $$ \mu_2^3 \longrightarrow M_{\rho,\sigma} \subset T_{\rho} \times N'_{\sigma} \times G $$

défini via les composantes canoniques

(30) 
\begin{tikzcd}
& T_\rho \\
\mu_2 \ar[r] & \mathbb{G}_m \ar[u] \ar[r] \ar[d] & T_\sigma \ar[r, hook] & N'_\sigma \\
& G
\end{tikzcd}

Par passage au quotient, l'hom (25)
induit

(31) $$ \hat{\Gamma}_{0,3} \longrightarrow \Gamma_{\rho,\sigma}(\hat{\mathbb{Z}})/ \underbrace{\{1 \times \pm 1 \}}_{\Sigma_0} $$

où

(32) $$ \Gamma_{\rho,\sigma} \subset T_{\rho} \times N'_{\sigma} $$
$$ = \left\{ (\underline{a}, \underline{b}) \mid \det \underline{a} / \det \underline{b} \in \mu_2 \right\} $$

La donnée - plutôt la construction faite
plus haut - donne un hom de groupes (24)

\newpage

%% ===== Page 483 =====

544

revient à la construction des quatre homomorphismes de groupes

(33) $$S_0 M_{0,3}^\sim \xrightarrow{\Theta_{B'}^\pi} B_0'(\hat{\mathbb{Z}}) = \hat{\mathbb{Z}}^* \cdot \hat{\mathbb{Z}}$$
[MARGIN: donnant $u$]

(34) $$M_{0,3}^\sim \xrightarrow{\Theta_G^\pi} G(\hat{\mathbb{Z}})/\pm 1 = Go(2, \hat{\mathbb{Z}})/\pm 1$$
[MARGIN: donnant $U$ (mod signe) (produit $1/2$ direct)]

(35) $$\Gamma_{0,3}^\sim \xrightarrow{\Theta_\rho^\pi} T_\rho(\hat{\mathbb{Z}})/\pm 1 \quad \text{(unités adéliques (mod $\pm 1$) du corps $\mathbb{Q}(\sqrt{-3})$)}$$
[MARGIN: donnant $a$ (mod signe)]

(36) $$\Gamma_{0,3}^\sim \xrightarrow{\Theta_\sigma^\pi} T_\sigma(\hat{\mathbb{Z}})/\pm 1 \quad \text{(unités adéliques (mod $\pm 1$) du corps $\mathbb{Q}(i)$)}$$
[MARGIN: donnant $b$ (mod signe)]

reliés par les relations (4), que nous allons identifier. On utilise le sous-groupe

(37) $M_{0,3}^{\sim !}[0] =$ sous-groupe de $M_{0,3}^\sim$ qui (invarie les éléms de [unclear: la base] tracés ds $T_{0,3}^\sim$, et qui) normalise $L_{0 \Pi}$ (sous-groupe fermé engendré par $l_0$)

et on note que

(38) $$1 \to L_{0 \Pi} \to M_{0,3}^{\sim !}[0] \to \Gamma_{0,3}^\sim \to 1$$

et que

(39) $$M_{0,3}^\sim = M_{0,3}^{\sim !}[0] \cdot \widehat{C}_{0,3}^+$$
avec $$M_{0,3}^{\sim !}[0] \cap \widehat{C}_{0,3}^+ = L_{0 \Pi}$$
[MARGIN: sous $\widehat{C}_{0,3}^+$: sous-groupe invariant de $M_{0,3}^\sim$]

$-$ bien sûr, (38) est conséquence de (39).

enfin

(40) $$S_0 M_{0,3}^\sim \simeq M_{0,3}^{\sim !}[0] \cdot \widehat{C}_{0,3}^{+ \wedge}$$
de $S_0 M_{0,3}^\sim$ (produit $1/2$ direct)

1º) L'hom. (33) est trivial sur le facteur $\widehat{C}_{0,3}^{+ \wedge}$

\newpage

%% ===== Page 484 =====

[MARGIN: en dépend.
donc la valeur
de (33) resulte de sa
description sur le
$M_{0,3}^\sim[0] \to B_0'(\hat{\mathbb{Z}})$]

$\underline{u}$ en dépend (ne dépend que des couples $(\alpha, \beta)$ pas des triplets $(\alpha, \beta, \gamma)$ qui
décrit un élément de $\mathcal{M}_{0,3}^\sim$ ), et
la [unclear: pire] définition de $\underline{u}$ en termes de
$\alpha, \beta$ - il y a seulement en termes

[MARGIN: en un element de l'image de $\xi = \alpha \rho(\alpha)^{\sim}$ dans $S \ell(2, \hat{\mathbb{Z}})$
en l'équation
$\xi p \xi ! p(\xi) = id !$ ah! l'horreur, par surhaussement dans $S \ell(2, \hat{\mathbb{Z}})$ ! --]

a été faite essentiellement dans le § 45
ou base en vue dans § 45 III, - on aura

(41) $$ \underline{u} = \left(\begin{matrix} 1/D & B/D \\ 0 & 1 \end{matrix}\right) \sim \xi = \left(\begin{matrix} A & B \\ C & D \end{matrix}\right) $$
(dans $S \ell(2, \hat{\mathbb{Z}})$)

ou encore, en termes de
$\underline{\alpha} = \left(\begin{matrix} a & b \\ c & d \end{matrix}\right)$ avec $ad - bc \in \{ \pm 1 \}$

(42) $$ \begin{cases} \varepsilon' D = c^2 + cd + d^2 \\ \varepsilon' B = ac + bd + bc \end{cases} $$

Cela ne dépend pas des fondements non faits
sur ces groupes profinis - les "calculs en
$\xi, \eta$" sont en effet suffisants. C'est grâce
à la description de cet homomorphisme
qu'on peut définir

(43) $$ \mathcal{M}_{0,3}^\sim (0) \subset \mathcal{M}_{0,3}^\sim [0] $$
$"$
$$ \{ \underline{u} \mid \underline{u} \in T_0 (\hat{\mathbb{Z}}) \text{ i.e. } B = 0 \} $$
i.e. $ac+bd+bc=0$
(condition de non-torsion),

\newpage

%% ===== Page 485 =====

545

et obtenir des décompositions canoniques en produits 1/2-directs
$$ (44) \qquad \mathcal{M}_{0,3}^\sim[0] \simeq \mathcal{M}_{0,3}^\sim(0) . L_{0 \pi} $$
$$ (45) \qquad \mathcal{M}_{0,3}^\sim \simeq \mathcal{M}_{0,3}^\sim(0) . \mathfrak{S}_{0,3}^{+\wedge} . $$

2º) Pour décrire (34), comme il est déjà connu sur $\mathfrak{S}_{0,3}^{+\wedge}$ , il suffit de le décrire sur $\mathcal{M}_{0,3}^\sim[0]$ ou = seulement sur $\mathcal{M}_{0,3}^\sim(0)$ (moyennant diverses vérif de compatibilités...). Et sur $\mathcal{M}_{0,3}^\sim[0]$ il coïncide avec (33) - tandis que sur $\mathcal{M}_{0,3}^\sim(0)$ il est pratiquement trivial et se réduit aux multiplicateurs $\mu$

$$ (46) \qquad \Theta_G^\pi = \Theta_{B'_0}^\pi \text{ sur } \mathcal{M}_{0,3}^\sim[0] $$
$$ (47) \qquad \left\{ \begin{array}{l} \Theta_G^\pi(\underline{u}) = i'_{T_0}(\mu) \qquad \text{si } \mu \text{ est le multiplicateur de } \underline{u} \\ \text{si } \underline{u} \in \mathcal{M}_{0,3}^\sim(0) \end{array} \right. $$

Cela = une définition possible de $\mathcal{M}_{0,3}^\sim(0)$ ).

Remarques  Comme $\Theta_G^\pi$ a été construit "à la force du poignet" par l'intermédiaire de $\Theta_{B'_0}^\pi$ i.e. $U$ a été construit par $\underline{u}$ , il y a eu longtemps une espèce de

\newpage

%% ===== Page 486 =====

confusions dans mon esprit entre $\underline{u}$ et $\underline{U}$,
d'autant plus que la plus grande partie
du temps je me réduisais à travailler
dans $\mathcal{M}_{0,3}^{\sim}[0]$ (après avoir pas mal
cafouillé avec des points de vue purement
internes, avant de me rendre compte
que techniquement c'était la façon
la plus élégante de procéder - à
partir du $§ 45$). Cette confusion ne
s'est définitivement clarifiée que hier
et - aujourd'hui, à force de mettre et
remettre les pieds dans mes propres
plats...

3º) Les liens (35) (36) n'ont été obtenus
qu'en travaillant à fond sur les
paramètres $\alpha, \beta$ pour $\underline{u}$ - bien entendu,
il a suffi de travailler dans $\mathcal{M}_{0,3}^{\sim}[0]$,
et de définir 
$\underline{u} \longmapsto \underline{a}$
$\mathcal{M}_{0,3}^{\sim}[0] \longrightarrow \underline{T}_l(\hat{\mathbb{Z}})/\pm 1$
$\longrightarrow \underline{N}_l(\hat{\mathbb{Z}})/\pm 1$
$\underline{u} \longmapsto \underline{b}$

\newpage

%% ===== Page 487 =====

545

de façon que $a$ soit trivial sur $\hat{L_0^\pi}$.
Ces quantités $a, b$ sont apparues dans les
calculs progressivement à partir du
§ 44 (p. 366' ff), le besoin s'en faisant
vite [unclear] à n'en [unclear] seulement
à partir du § 48 III, IV, V - et une compré-
hension complète de leur nature, du point
de vue des schémas en groupes qui inter-
viennent dans la théorie, s'est faite égale-
ment seulement en ces tous derniers
jours. Cette compréhension reste d'
ailleurs pour le moment assez
formelle, et il faudrait approfondir
dans plusieurs directions :

a) Expliciter la signification du point
de vue des corps de classes de
de $\mathbb{Q}(j)$ et de $\mathbb{Q}(i)$.

b) Expliciter la signification pour
la théorie des c-es "congruentielles"
(i.e. sous-entendu : des repr. [unclear]
du groupe $\operatorname{Sl}(2, \hat{\mathbb{Z}})$ ou $\operatorname{Sl}(2, \mathbb{Z})^\wedge$ )
des résultats obtenus,

486

\newpage

%% ===== Page 488 =====

c) Introduire sur $\mathcal{M}_{g,n}$ les éléments de structure caractéristiques des groupes de Galois motiviques — et faire le lien avec les constructions typiques de cet ordre en théorie des corps de classes.

d) Examiner dans quelle mesure dans ce contexte on peut faire intervenir les conditions de compatibilité d'appa-rence dromiciennes [unclear: au dsp] du §47 liées à la compatibilité avec
$$ \varphi : \hat{\Pi}_{0,3} \xrightarrow{\sim} \hat{\mathcal{G}}_{0,3}^+ $$
— je prévois d'ailleurs que le cadre des calculs précédents est trop étroit pour arriver à y faire intervenir ces conditions supplémentaires.

e) Ce qui m'intéresse le plus pour le moment, c'est une systématisation des calculs précédents dans un cadre beaucoup plus vaste que celui du seul groupe $S G$ et de l'hom. standard

\newpage

%% ===== Page 489 =====

547
$$ \widehat{G}_{0,3}^{+ \wedge} \longrightarrow [unclear] \operatorname{Gl}(\hat{\mathbb{Z}})/\pm 1 $$

- j'ai le sentiment d'avoir mis le
doigt sur un procédé systématique
pour construire des représentations
"motiviques" de $\Gamma_{\mathbb{Q}}$ , via des représenta-
tions de la partie "géométrique"
$\widehat{G}_{0,3}^{+ \wedge}$ de $\mathcal{M}_{0,3}^\wedge \simeq \Pi_1(M_{0,3 \mathbb{Q}} , \vec{p}_3)$
$\simeq \Pi_1(M_{1,1 \mathbb{Q}})$ - ou mieux, via des
représentations du groupe $G_{0,3}^+ \simeq \operatorname{Sl}(2, \mathbb{Z})'$
discret (ou de variantes telles
$\operatorname{Sl}(2, \mathbb{Z})$, $\Pi_{0,3}$ ... ) à valeurs dans les
points entiers de schémas en groupes
[crossed out: génériques]
(réductifs ?? ) sur $\mathbb{Z}$ .

f) Il y a pourtant une autre direc-
tion d'investigation qui m'attire
beaucoup, par un cheminement lié
à l'étude des schémas en
groupes, et qui pourrait se pour-
suivre [crossed out: unclear] à l'aide des seuls
"gros" groupes fondamentaux profinis -

\newpage

%% ===== Page 490 =====

savoir l'étude de certains $\mathbb{Q}$-revêtements
particuliers de $M_{0,3 \mathbb{Q}} = \mathbb{P}_{\mathbb{Q}} - \{0,1,\infty\}$,
et des renseignements qu'on doit
en tirer sur des propriétés
supplémentaires les autom. ext.
"arithmétiques" de $\hat{\Pi}_{0,3}$ (du type
des résultats obtenus dans § 4), et
sur des calculs de scindages de
l'extension $M_{0,3}$ de $\hat{\Pi}_{0,3}$ par $\widehat{G}_{0,3}^{+ \wedge}$ en
d'autres points que les points "[unclear: critiques]"
$0, 1, \infty, 1/2, 2, -1, -j, \bar{j}$.

(Sans compter la question d'en revoir
les fondements, via la méthode suggérée
par Deligne - et d'étudier l'équation
des droites dans le compactifié pro-nilp.
et du continu : [unclear: isomorphies] dans le
discret, des conjectures fondamentales,
notamment sur l'anti-involution des
points complexes - etc etc).

\newpage

%% ===== Page 491 =====

548

Revenons pour l'instant à la question d'identi-
fier les $^{quatre}$ hom. obtenus (33) & (36), en termes
de choses classiques. On peut considérer
l'hom. (33), $^{\Theta_{B'_1}^F}$ comme de nature essentielle-
ment tautologique, ou du moins comme
étant subordonné à l'hom. $\Theta_G^F$ de (34).
D'ailleurs, (35) et (36) s'en déduisent,
en les composant avec les deux hom.
$$ \Pi_{0,3}^\sim \to \mathcal{M}_{0,3}^\sim $$
((morphismes de pro-groupes, pour le moment !))
associés aux deux points $1/2$ et $-1$ de
$\mathcal{M}_{0,3\mathbb{Q}}(\mathbb{C})$, et qui correspondent,
au niveau des schémas en groupes,
aux deux $_{(3)}$ scindages $_{(31)}$ $\gamma_p$, $\gamma_\sigma$ de
$\mathcal{M}_{p,\sigma}$ sur $\mathcal{T}_{p,\sigma}$, donnés respectivement par

$$ (38) \begin{cases} \gamma_p : (\underline{a}, \underline{b}) \longmapsto (\underline{a}, \underline{b}, \underline{a}) \\ \gamma_\sigma : (\underline{a}, \underline{b}) \longmapsto (\underline{a}, \underline{b}, \underline{b}) \end{cases} \quad \begin{matrix} \mathcal{T}_{p,\sigma} \to \mathcal{M}_{p,\sigma} \\ \mathcal{T}_{p,\sigma} \to \mathcal{M}_{p,\sigma} \end{matrix} $$

avec un grain de sel que le premier n'est
défini que sur le sous-groupe $\mathcal{T}'_{p,\sigma}$ défini par

$$ (49) \begin{cases} \mathcal{T}'_{p,\sigma} \subset \mathcal{T}_{p,\sigma} \\ "\; \{ (\underline{a}, \underline{b}) \mid \det \underline{a} = \det \underline{b} \} \\ "\; \simeq T_p \times_{\mathbb{G}_m} N'_\sigma \end{cases} $$

\newpage

%% ===== Page 492 =====

et de = [unclear]
$$ \varphi_P^\pi : \Pi_{\mathbb{Q}} \simeq \tilde{C}_{1,1}^{\sim \prime} \longrightarrow \mathcal{T}_P(\hat{\mathbb{Z}})/\pm 1 $$
$$ \Pi_{\mathbb{Q}}^\pi \overset{\sim}{\to} \{ \gamma \in \Pi_{\mathbb{Q}}^\sim \mid \varepsilon_3(H) = 1, \mu = \chi(\gamma) \} $$

[MARGIN: On fera attention cependant que contrairement à $\hat{\Theta}_P^\pi : \Pi_{P} \dots$, l'homomorphisme $\hat{\Theta}_G^\pi$ est défini sans aucune indétermination]

Où sont donc passés les [unclear: inévitables ambiguïtés]
des signes pour $\hat{\Theta}_P^\pi, \hat{\Theta}_G^\pi$ correspondant
aux indéterminations pour $\hat{\Theta}_G$ lui-même - tout va bien ! Cette
dernière est levée [unclear: d'un coup]
en remplaçant $\mathcal{M}_{0,3}^\sim$ par l'
extension [unclear: de groupe] $\tilde{\mathcal{C}}_{1,1}^\sim$ par $\pm 1$,

de sorte qu'on trouve
(50) $$ \tilde{\mathcal{C}}_{1,1}^\sim \longrightarrow \operatorname{Gl}(2, \hat{\mathbb{Z}}) $$
et il faudrait [unclear: pouvoir l'identifier]
à cet homomorphisme fondamental "modulaire" de
la théorie des courbes elliptiques
(que l'image inverse de $\Pi_{\mathbb{Q}} \subset \Pi_{\mathbb{Q}}^\sim$
sur $\tilde{\mathcal{C}}_{1,1}^\sim$ puisse s'interpréter comme
$\tilde{\mathcal{C}}_{1,1}^\sim \simeq \pi_1(M_{1,1/ \mathbb{Q}})$ doit être plus ou
moins tautologique, par construction.)
Il est clair en tous cas que les
deux coïncident (à [unclear: des éléments] de
(51) $$ \mu_2(\hat{\mathbb{Z}}) = \{\pm 1\}^P \quad (P = \text{ens. des nbs premiers}) $$

\newpage

%% ===== Page 493 =====

i.e. ils coïncident modulo un hom.
$$\tilde{\Gamma}_{\mathbb{Q}} \longrightarrow \mu_2(\hat{\mathbb{Z}}) \simeq \{\pm 1\}^{\mathbb{P}}$$
i.e. un hom.
$$G_m(\hat{\mathbb{Z}})_2 = (\hat{\mathbb{Z}}^\ast)_2 \simeq \{\pm 1\}^{\mathbb{P}} \longrightarrow \mu_2 G_m(\hat{\mathbb{Z}}) = \mu_2(\hat{\mathbb{Z}}) \simeq \{\pm 1\}^{\mathbb{P}}$$

je doute en fait [unclear] que
c'est un [unclear] ! C'est moralement sûr que
les deux coïncident - je n'ai pas essayé
de le prouver

Quant à (35) (36) , ils doivent corres-
-dre sûrement à : l'action de $\Gamma_{\mathbb{Q}}$ sur
les modules de Tate des courbes elliptiques
à multiplications complexes correspondantes
justement aux pts $(1/2, -1)$ et $(-j, -j)$
de $U_{0,3 \mathbb{Q}} = M_{1, [2]' \mathbb{Q}}$ . L'indétermination
de signes dont correspond à l'indétermination
à remonter aux points de $M_{1, [2]' \mathbb{Q}}$ en
un point de $M_{1, [2] \mathbb{Q}}$ , i.e. de choisir
une courbe elliptique (avec "rigidification"
d'échelon 2) correspondant aux données
de Legendre $(0, 1, \infty, \rho)$ avec $\rho = 1/2$ ou $\rho = -j$ .

\newpage

%% ===== Page 494 =====

Même une fois que l'ambiguité de signe
pour $\Theta_{\bar{\mathbb{Q}}}^\pi$ est levée en passant à l'exten-
sion $\tilde{C}_{1,1}$ à deux feuillets de $\tilde{M}_{0,3}$, l'in-
détermination de signes pour $\Theta_{\bar{\rho}}^\pi$, $\Theta_{\bar{j}}^\pi$ n'est
pas encor surmontée, car les semi-groupes
[crossed out] ($\psi_{1/2}^\pi, \psi_{-J}^\pi$) ne se prolon-
gent pas en des semi-groupes ([unclear: à p. strictes]) de
[crossed out: cet extension] $\tilde{C}_{1,1}$, mais seulement
de $\tilde{M}_{0,3}$ - et des relevements de tels
semi-groupes sur $\tilde{C}_{1,1}$ revient très précisé-
ment, à choisir justement des courbes
elliptiques à rigidification de Legendre
au dessus des pts $1/2, -J$. Ainsi, des
calculs faits en termes de l'action
du groupe de Galois $\hat{\Pi}_{0,3}$ dans $\widehat{\mathbb{Z}/2 * \mathbb{Z}/3}$,
(si les signes coïncidaient !),
donnent des résultats qui collent
de très près : des réalités arithmético-
géométriques particulièrement cruciales.

\newpage

%% ===== Page 495 =====

050

49. Homomorphismes de $M_{0,3}$, les groupes $M_{g,\nu}$ généraux

I) Considérons un sous-groupe [unclear: fermé] $g$ de $M_{0,3}^\sim$, contenant $\widehat{\mathfrak{G}}_{0,3}^+$, donc image inverse d'un sous-groupe fermé de $\widehat{\Gamma}_{0,3}$.

[MARGIN: Oui on a un [unclear] Tout l'est... section [unclear] un ... le reste, [unclear] a un ... par un sous-groupe d'indice fini de [unclear] débauches p. 553]

On supposera que celui-ci est "d'épaisseur" grosse - disons qu'il contienne un sous-groupe d'indice fini de $\Gamma_{\mathbb{Q}} \subset \widehat{\Gamma}_{0,3}$. Pour le moment, on supposera pour simplifier que

$$g \subset K_{\varepsilon_3} . \quad \text{En résumé}$$

(1) $\widehat{\mathfrak{G}}_{0,3}^+ \subset g \subset M_{0,3}^\sim$, $u \in g \Rightarrow \mu(u) \equiv 1 \pmod 3$

Cette dernière hypothèse signifie donc que pour tout $u \in g$, $\alpha(u) \in \widehat{\mathfrak{G}}_{0,3}^+$, de sorte qu'on a

(2) $g \underset{\beta_0}{\overset{\alpha_0}{\rightrightarrows}} \widehat{\mathfrak{G}}_{0,3}^+$

caractérisés par

(2') $\begin{cases} u(\rho) = \text{int}(\alpha_0)\rho \\ u(\sigma) = \text{int}(\beta_0)\sigma \end{cases} \quad \begin{cases} \alpha_0 = \alpha(u) \\ \beta_0 = \beta(u) \end{cases}$

On considère

(3) $\begin{aligned} g_{(0)} &= g \cap M_{0,3}^\sim(0) \\ g_{[0]} &= g \cap M_{0,3}^\sim[0] = \operatorname{Norm}_g(\mathbb{L}_0^\sim) \end{aligned}$

de sorte que

(4) $g = g_{(0)} \cdot \widehat{\mathfrak{G}}_{0,3}^+$ (produit semi-direct)

(5) $g_{[0]} = g_{(0)} \cdot \mathbb{L}_0^\sim$

(6) $g_{[0]} \cap \widehat{\mathfrak{G}}_{0,3}^+ = \mathbb{L}_0^\sim$

Où on a posé

(7) $\mathbb{L}_0^\sim =$ sous-groupe fermé de $\widehat{\mathfrak{G}}_{0,3}^+$ engendré par $\varepsilon_0$ i.e. [unclear] de $\mathbb{L}_1^\sim , \mathbb{L}_\infty^\sim$.

\newpage

%% ===== Page 496 =====

Considérons un [MARGIN: profini] groupe $M$ , et un hom. (continu) surjectif

(8) $$ \varphi : g \longrightarrow M .$$

On pose
(9) $$ \begin{cases} \varphi( \hat{\mathcal{G}}_{0,3}^+ ) = \mathcal{C} & \text{ss-groupe distingué de } M \\ \varphi( g_{(0)} ) = M_{(0)} \\ \varphi( g_{[0]} ) = M_{[0]} \\ \varphi( \tilde{L_0} ) = L_0 \end{cases} $$

Donc on aura les sous-groupes fermés de $M$, satisfaisant

(10) $$ M_{(0)} \subset M_{[0]} \subset \operatorname{Norm}_M(L_0) $$
(11) $$ L_0 \subset \mathcal{C} \cap M_{[0]} $$

Par abus de langage, on désignera encore par $\rho, \sigma, \varepsilon_0, \varepsilon_1$ ... les images de ces éléments dans $\mathcal{C}$, qui y satisfont donc aux relations connues dans $\hat{\mathcal{G}}_{0,3}^+$ , en particulier

(12) $$ \sigma^2 = \rho^3 = 1 , \quad \sigma \rho = \varepsilon_0 , \quad \rho \sigma = \varepsilon_1 , \quad \varepsilon_1 = \sigma(\varepsilon_0) = \rho(\varepsilon_0) .$$

Les restrictions des opérations $\alpha_0, \beta_0$ de (2, 2') à $\mathcal{M}_{0,3}[0]$ seront plutôt notées $\alpha, \beta$ , où $\alpha, \beta$ sont en fait des opérations

(13) $$ Sg \overset{\alpha}{\underset{\beta}{\rightrightarrows}} \hat{\mathcal{G}}_{0,3}^+ $$
$$ Sg \overset{déf}{=} \text{ ext. par } \tilde{L_0} \text{ s-g de } S\mathcal{M}_{0,3} $$
$$ = \{ (u,f) \in g \times \hat{\mathcal{G}}_{0,3}^+ \mid f u \in g_{[0]} \} $$

dont on regarde les restrictions au sous-groupe ($\simeq \mathcal{M}_{0,3}[0]$) de $Sg$ défini par la relation $f=1$.

Comme $\alpha, \beta : g_{[0]} \rightrightarrows \hat{\mathcal{G}}_{0,3}^+$ ... se factorisent par $\varphi | \hat{\mathcal{G}}_{0,3}^+$ , donnent des opérations, encore notées $\alpha, \beta$ (i.e. $\alpha^M, \beta^M$ si on a une confusion possible)

(14) $$ g_{[0]} \overset{\alpha}{\underset{\beta}{\rightrightarrows}} \mathcal{C} $$

satisfaisant aux relations

(15) $$ \begin{cases} \alpha(p) = \rho \alpha(p) \rho \\ \beta(p) = \rho \beta(p) \sigma \end{cases} \quad \text{pour } u \in M_{[0]} \text{ (} \text{[unclear: image de]} \ \varphi(u_0)\text{),} $$
$$ \text{où } \alpha = \alpha(u_0), \beta = \beta(u_0). $$

\newpage

%% ===== Page 497 =====

A priori, $\alpha$ et $\beta$ dépendent des choix de $u_0$ dans ce procès
pas seulement de $u$ - mais
[MARGIN: (16)]
$$ \begin{cases} \alpha p(\alpha)^{-1} = [u, \rho] \\ \beta p(\beta)^{-1} = [u, \sigma] \end{cases} $$
n'en dépendent pas. Si on pose
[MARGIN: (17)]
$$ T_\rho = Cent_M(\rho) \qquad T_\sigma = Cent_M(\sigma) $$
on voit donc que, modulo $T_\rho$ resp. $T_\sigma$, $\alpha, \beta$ ne dépendent
que de $u$.

Hypothèse $\alpha (= \alpha(u_0))$, $\beta (= \beta(u_0))$ ne dépendent que de $u = \varphi(u_0)$ (pour $u_0 \in G[0]$)
On exprimera cette hypothèse de façon un peu
différente, en introduisant
[MARGIN: (18)]
$$ J = \ker(G[0] \longrightarrow M[0]) , $$
on veut donc que, pour
$u'_0 = u_0 \delta_0$ avec $u_0 \in G[0]$, $\delta_0 \in J$, on ait
$\alpha(u'_0) = \alpha(u_0)$, $\beta(u'_0) = \beta(u_0)$.
Soient $\alpha_0 = \alpha(u_0)$, $\beta_0 = \beta(u_0)$, $\alpha'_0 = \alpha(u'_0)$, $\beta'_0 = \beta(u'_0)$, $\bar{\alpha}_0 = \alpha(\delta_0)$, $\bar{\beta}_0 = \beta(\delta_0)$
dans $G$,
[unclear: d'après ce qu'on a vu], par les formules de cocycles de $G$ :
$$ \alpha'_0 = \delta_0(\alpha_0) \bar{\alpha}_0 , \quad \beta'_0 = \delta_0(\beta_0) \bar{\beta}_0 $$
d'où, en appliquant $\varphi$ et notant que
$$ \varphi(\delta_0(\alpha_0)) = \varphi(\delta_0) (\varphi(\alpha_0)) = \alpha \quad \text{et de même} = $$
$$ \varphi(\delta_0(\beta_0)) = \varphi(\delta_0) (\varphi(\beta_0)) = \beta $$
[MARGIN: (19)]
$$ \alpha' = \alpha \bar{\alpha} , \quad \beta' = \beta \bar{\beta} \quad \text{dans } \mathfrak{G} $$
où $\alpha, \beta, \alpha', \beta', \bar{\alpha}, \bar{\beta}$ sont les images par $\varphi$ des quantités
$\alpha_0$ etc. de $\widehat{\mathfrak{G}_{0,s}}^{+\wedge}$. Donc l'hypothèse
se [unclear: traduit alors] :
[MARGIN: (20)]
$$ \forall \delta_0 \in J = \ker(G[0] \to M[0]) , \quad \bar{\alpha}(\delta_0), \bar{\beta}(\delta_0) \in \ker p $$
Quand cette hypothèse est satisfaite, on a donc des
applications.

\newpage

%% ===== Page 498 =====

(21) 
$$ \mathcal{M}[\sigma_0] \overset{\alpha}{\underset{\beta}{\rightrightarrows}} \hat{G} $$
[unclear: satisfaisant] [crossed out]
précisant (14), et satisfaisant aux

(22)
$$ \begin{cases} u(\rho) = \alpha(u) \rho \\ u(\sigma) = \beta(u) \sigma \end{cases} \quad \text{pour } u \in \mathcal{M}[\sigma_0], \, \alpha = \alpha(u), \, \beta = \beta(u) , $$

qui généralisent (2').

D'ailleurs, posant

(23)
$$ \xi = \alpha \rho \alpha^{-1}, \quad \eta = \beta \sigma \beta^{-1} $$
les relations (22) s'écrivent aussi

(24)
$$ u(\rho) = \xi \rho, \quad u(\sigma) = \eta \sigma $$
où on a

(25)
$$ \xi = [u, \rho], \quad \eta = [u, \sigma] $$
et de plus

(26)
$$ \xi = \eta^\nu \quad \text{où} \quad \nu = \frac{\mu - 1}{2} \in \hat{\mathbb{Z}} \quad \text{(cf le multipli-} $$
[MARGIN: cateur de $u$).]

qui provient de la relation analogue dans $\hat{G}_{0,3}^+$.
(On suppose pour fixer les idées que
l'hom.

$$ g \xrightarrow{\chi^g} \hat{\mathbb{Z}}^\times \quad \text{("multiplicateur")} $$
se factorise par

(27)
$$ g \rightarrow \mathcal{M} \xrightarrow{\chi^\mathcal{M}} \hat{\mathbb{Z}}^\times , $$
de sorte qu'on a une notion de multiplicateur ($\mu \in \hat{\mathbb{Z}}^\times$)
pour les $u \in \mathcal{M}$.)

Considérons l'hom. naturel

(28)
$$ \mathcal{M}[\sigma_0] \rightarrow \operatorname{Aut}(\hat{G}) $$
$$ u \mapsto int(u) | \hat{G} $$
on voit sur (24) qu'il est connu quand on connaît

\newpage

%% ===== Page 499 =====

[MARGIN: à part topologie de $F$ dont] les fonctions $\alpha, \beta : M[T_0] \to \tilde{\mathbb{G}}$, et par suite
en fait seulement $\xi, \eta : M[T_0] \to \tilde{\mathbb{G}}$.
D'ailleurs on pose

(29) $$M'[T_0] = \operatorname{Im}(M[T_0] \to \operatorname{Aut}(\tilde{\mathbb{G}}))$$

ces [unclear: mêmes] fonctions $\xi, \eta$ sont déjà définies sur
$M'[T_0]$ - alors qu'il n'est pas clair qu'il en soit
de même de $\alpha, \beta$, quand $M[T_0]$ n'opère pas fidèlement
sur $\tilde{\mathbb{G}}$. On a ainsi une application [unclear] de
certain sous-groupe $M'[T_0] \subset \operatorname{Norm}_{\operatorname{Aut}(\tilde{\mathbb{G}})}(L_0)$ dans
l'ensemble des solutions de l'équation en $\xi, \eta$
(26) qui s'écrit [unclear]

(30) $$\xi^{-1} \eta \in L_1$$

relation qui n'est que (26) (avec la forme précise de
l'exposant $\nu$) - [unclear] que
$\tilde{\mathbb{G}}_{0,s}^+ \to \operatorname{Sl}(2, \hat{\mathbb{Z}})$
se factorise par $\tilde{\mathbb{G}}$ - ce qu'on
va supposer, pour nous simplifier la vie.

On va faire sur $\tilde{\mathbb{G}}$, muni de $\rho, \sigma$, l'hypo-
thèse [unclear] suivante

$$ \bar{X} = X_{\tilde{\mathbb{G}},\rho,\sigma} \subset \tilde{\mathbb{G}} \times \tilde{\mathbb{G}} $$
(31)
$$ \left\{ \begin{array}{l} (\xi, \eta) \text{ resp. conjugués de } \rho, \sigma \text{ dans } \tilde{\mathbb{G}} \\ \xi \eta^{\nu} \in L_1, \text{ i.e. } \exists l \in L_1 \text{ que } \xi = \eta^{-\nu} l \end{array} \right. $$

[unclear crossed out text]
Pour $u \in \operatorname{Aut}(\tilde{\mathbb{G}})$ [unclear]
[unclear crossed out text] fixant les classes de conjugaison de $\rho$
et celles de $\sigma$, i.e. telles que
$\xi = [u, \rho]$, $\eta = [u, \sigma]$
satisfont aux deux premières conditions de (31), et

\newpage

%% ===== Page 500 =====

et soit $u \in Aut (\tilde{G})$ un automorphisme de $\tilde{G}$,
et posons

(34)
$$ \begin{cases} \xi = [u,\rho] = u \rho u^{-1} \rho^{-1} = u(\rho) \rho^{-1} & \text{i.e. } u(\rho) = \xi \rho \\ \eta = [u,\sigma] = u(\sigma) \sigma^{-1} & \text{i.e. } u(\sigma) = \eta \sigma \end{cases} $$
[MARGIN: cf (24) (25)]

alors dire que $\xi \rho$ $(=u(\rho))$ est conjugué à $\rho$ signifie que
$u$ fixe la cl. de c-j. de $\rho$, dire que $\eta \sigma$ $(=u(\sigma))$
conjugué de $\sigma$ signifie que $u$ fixe la classe de
c-j. de $\sigma$. Enfin la condition $\eta^{-1} \xi \in L_1$ s'écrit
encore, comme

[MARGIN: NB $\sigma^{-1} = \sigma$]

$$ \eta^{-1} \xi = (\sigma u(\sigma^{-1})) (u(\rho) \rho^{-1}) = \sigma u(\sigma^{-1} \rho) \rho^{-1} = \sigma (u(\sigma^{-1} \rho) \underbrace{\rho^{-1} \sigma}_{\xi_0}) \underbrace{\sigma^{-1} \rho^{-1}}_{\xi_0^{-1}} $$
$$ = \sigma (u(\xi_0) \xi_0^{-1}) $$

dans la forme
$$ \sigma (u(\xi_0) \xi_0^{-1}) \in L_1 \quad \text{i.e.} \quad u(\xi_0) \xi_0^{-1} \in L_0 $$
$(L_0 = \sigma^{-1} L_1)$ ) — ou encore (puisque $\xi_0 \in L_0$ [unclear: engendre $L_0$])
$u(\xi_0) \in L_0$, soit (puisque $u$ [unclear: conserve] $L_0$)
$$ u(L_0) = L_0 $$
i.e. que $u$ normalise $L_0$.

[MARGIN: [unclear: Dans notre cas où le $\xi_0$ (qui figure dans $\sigma \rho = \xi_0 \in L_1$) engendre $L_0$, de sorte que la condition $u(\xi_0) \in L_0$ équivaut à $u(L_0) = L_0$ ... $\rho^2=1$ ...]]

Soit donc

(35)
$$ \tilde{B}'_0 = \operatorname{Aut}(\tilde{G}) $$
$$ \| dfn $$
$$ \{u \in \operatorname{Aut}(\tilde{G}) \mid u \text{ fixe classes de c-j. de } \rho \text{ et celle de } \sigma, u \text{ normalise } L_0 \} $$

de sorte qu'on a une application in-jective pourvu que
[MARGIN: pour $\rho,\sigma$ eng. $\tilde{G}$]

(36)
$$ \tilde{B}'_0 \hookrightarrow \Sigma $$
$$ u \longmapsto ([u,\rho] = u(\rho)\rho^{-1}, [u,\sigma] = u(\sigma)\sigma^{-1}) $$

\newpage

%% ===== Page 501 =====

553

Hypothèse fondamentale sur $\overline{\mathcal{G}}$, muni de $\rho, \sigma$, et d'un $\varepsilon_0 = \sigma^{-1}\rho$ ,
$\varepsilon_1 = \rho \sigma^{-1} = \sigma(\varepsilon_0) = \rho(\varepsilon_0)$ :

[MARGIN: (37)]
a) $\overline{\mathcal{G}}$ est engendré topologiquement par $\rho, \sigma$
b) L'application (36) est (injective par a) et bijective.

Cela signifie donc que [unclear: text crossed out] $\Sigma$ est
La condition b) signifie donc que pour tout
$(\alpha, \beta) \in \overline{\mathcal{G}} \times \overline{\mathcal{G}}$, satisfaisant une relation
$$\alpha \rho(\alpha)^{-1} = \beta \sigma(\beta)^{-1} \varepsilon_1^{-1} \quad \text{s'écrivant aussi,}$$
il existe $u \in \operatorname{Aut}(\overline{\mathcal{G}})$ tel qu'on ait
$$[u, \rho] = \alpha \rho(\alpha)^{-1} \quad \text{i.e.} \quad u(\rho) = int(\alpha)(\rho) \quad \text{i.e.} \quad \text{[unclear: } u \rho u^{-1} = \alpha \rho \alpha^{-1}]$$
$$[u, \sigma] = \beta \sigma(\beta)^{-1} \quad \text{i.e.} \quad u(\sigma) = int(\beta)(\sigma) \quad \text{i.e.} \quad \text{[unclear: } u \sigma u^{-1} = \beta \sigma \beta^{-1}]$$
(on aura donc nécessairement $u \in \widetilde{\mathcal{B}}'_0$, et
$u$ sera unique par a) ).

Ainsi, sur l'ens. $R_{\rho, \sigma}$ des $(\xi, \eta)$ satisfaisant
a) b) c) , on aura une structure de groupe,
déduite par transport de structure de celle de $\widetilde{\mathcal{B}}'_0$
[unclear]

\newpage

%% ===== Page 502 =====

[MARGIN: Les groupes $\mathcal{U}_{\rho,\sigma}$ etc]

[MARGIN: Pour une justification des notations qui suivent, cf [unclear: plus loin] p. 564, remarque.]

II Reprenons la situation de 000, en oubliant (au moins pour le moment) $\mathcal{G}, \varphi$ etc.

Soit $\tilde{\mathfrak{S}}$ un groupe - on se placera indifféremment, soit dans le cadre des groupes discrets, soit dans celui des groupes profinis, -- au choix. On ne [unclear: fixera] pas les idées pour se [unclear: restreindre] dans le cadre des groupes profinis. On suppose donnés

(1) $\rho, \sigma \in \tilde{\mathfrak{S}}$

[MARGIN: [unclear: N.B. en outre $\varepsilon_0 = \rho^2, \sigma = \dots$ et $\rho^2 = \rho \varepsilon_0$ ou relation (?) à la place de (2) ]]

(2) $\varepsilon_0 = \sigma^{-1} \rho$ , $\varepsilon_1 = \rho \sigma^{-1} = \sigma(\varepsilon_0) = \rho(\varepsilon_0)$

(3) $L_0 =$ sous-groupe fermé engendré par $\varepsilon_0$
$L_1 =$ -------- '' -------- $\varepsilon_1$
donc $L_1 = \rho(L_0) = \sigma(L_0)$.

On considère

(4) $\operatorname{Aut}(\tilde{\mathfrak{S}}) \xrightarrow{\Phi = (\xi, \eta)} \tilde{\mathfrak{S}} \times \tilde{\mathfrak{S}}$

défini par

(5)
$$ \begin{cases} \xi(u) = [u, \rho] = u \rho u^{-1} \rho^{-1} = u(\rho) \rho^{-1} \quad \text{donc } u(\rho) = \xi \rho \\ \qquad (\xi = \xi(u)) \\ \eta(u) = [u, \sigma] = u \sigma u^{-1} \sigma^{-1} = u(\sigma) \sigma^{-1} \quad \text{donc } u(\sigma) = \eta \sigma \\ \qquad (\eta = \eta(u)) \end{cases} $$

Donc pour $u \in \operatorname{Aut}(\tilde{\mathfrak{S}})$, et $\xi = \xi(u), \eta = \eta(u)$, on a :

(6)
$$ \begin{cases} u \text{ fixe la classe de conjug. de } \rho \iff \xi \rho \sim \rho \text{ par } \tilde{\mathfrak{S}} \\ \qquad \iff \xi = \alpha \rho \alpha^{-1} \quad (\alpha \in \tilde{\mathfrak{S}}) \\ u \text{ fixe la classe de conjug. de } \sigma \iff \dots \\ \qquad \iff \eta = \beta \sigma \beta^{-1} \quad (\beta \in \tilde{\mathfrak{S}}) \\ u \text{ normalise } L_0 \iff \eta^{-1} \xi \in L_1 \iff \xi = \eta \varepsilon_1^\nu \quad \text{avec } \nu \in \hat{\mathbb{Z}} \text{ convenable} \end{cases} $$

\newpage

%% ===== Page 503 =====

554

Soit [crossed out: $B_{\rho,\sigma} \subset$]

(7) $B_{\rho,\sigma} \subset \operatorname{Aut}(\hat{C})$
$=$
$\{ u \mid \text{ satisfaisant les trois conditions (5)} \}$

(8) $\Sigma_{\rho,\sigma} \subset \hat{C} \times \hat{C}$
$=$
$\left\{ (\xi,\eta) \mid \begin{array}{l} \text{i.e. } \xi \rho \sim_{\hat{C}} \rho \text{ , i.e. } \eta \sigma \sim_{\hat{C}} \sigma \\ \xi = \alpha \rho(\alpha)^{-1} \text{ ; } \eta = \beta \sigma(\beta)^{-1} \quad (\alpha, \beta \in \hat{C}) \\ \eta^{-1} \xi \in L_1 \text{ i.e. } \xi = \eta \varepsilon_1^\nu \quad (\nu \in \hat{\mathbb{Z}}) \end{array} \right\}$ ,

de sorte que $B_{\rho,\sigma}$ est l'image inverse de $\Sigma_{\rho,\sigma}$
par (4), et on a un application induite

(9) $B_{\rho,\sigma} \xrightarrow{\varphi = \varphi_{\rho,\sigma}} \Sigma_{\rho,\sigma}$

On supposera
(10) $\begin{cases} a) \text{ } \rho,\sigma \text{ engendrent } \hat{C} \text{ topologiquement} \\ \\ b) \text{ l'application (9) (injective par } a)) \\ \text{ est surjective (donc bijective).} \end{cases}$

[MARGIN: Attention - cette condition a) est peu satisfaite pour $\hat{C} = \widehat{\operatorname{Sl}(2,\mathbb{Z})}$ !! (en tout cas l'analogue si on remplace $\hat{C}$ par [unclear] et $\varepsilon_1$ par [unclear] ! d'où au moins [unclear])]

On transporte alors par $\varphi$ la structure de groupe
de $B_{\rho,\sigma}$ : $\Sigma_{\rho,\sigma}$. On désignera par

[crossed out: $u_{\xi,\eta} = \varphi^{-1}(\xi,\eta)$]
(11) $u_{\xi,\eta} = \varphi^{-1}(\xi,\eta) \in B_{\rho,\sigma} \quad ((\xi,\eta) \in \Sigma_{\rho,\sigma})$

de sorte que $u_{\xi,\eta}$ est caractérisé par

(12) $\begin{cases} u_{\xi,\eta}(\rho) = \xi \rho \\ u_{\xi,\eta}(\sigma) = \eta \sigma \end{cases}$

L'hypothèse b) signifie donc que, dès que $(\xi,\eta) \in \Sigma_{\rho,\sigma}$
i.e. $\xi \rho \sim_{\hat{C}} \rho$, $\eta \sigma \sim_{\hat{C}} \sigma$ et $\eta^{-1} \xi \in L_1$, il existe

\newpage

%% ===== Page 504 =====

un autom. $u_{\xi,\eta}$ de $\underline{C}$ satisfaisant (12)
(ce qui sera assuré par (10) a) ). La loi de
groupe sur $\Sigma_{\rho,\sigma}$ s'explicite ainsi :

$$ (13) \quad (\xi',\eta')(\xi,\eta) = (\xi'',\eta'') \iff \begin{cases} \xi'' = u_{\xi,\eta}(\xi') \\ \eta'' = u_{\xi,\eta}(\eta') \end{cases} $$

Considérons l'hom. [crossed out: tautologique] canonique
$$ \underline{C} \to \operatorname{Aut}(\underline{C}) $$
$$ g \mapsto int(g) $$
et l'image inverse de $B_{\rho,\sigma}$ par cet hom., dans le normalisateur de $L_0$ dans $\underline{C}$, soit $N_0$. On veut
construire alors un groupe $G_{\rho,\sigma}$, contenant $\underline{C}$,
et dans lequel s'envoie $B_{\rho,\sigma}$, de façon que
l'action tautologique de $B_{\rho,\sigma}$ sur $\underline{C}$ soit
induite par cet homomorphisme, et que
$G_{\rho,\sigma}$ soit engendré par $\underline{C}$ et l'image de $B_{\rho,\sigma}$.
ainsi $B'_{\rho,\sigma}$, qui contient $N_0$. Si
le centre de $\underline{C}$ est trivial i.e. $\underline{C} \hookrightarrow \operatorname{Aut}(\underline{C})$,
[crossed out: pour faire cette construction, partons de] 
$B_{\rho,\sigma} \cdot \underline{C}$ (pr. semi/direct) ($\underline{C}$ sous gr. dist)
il suffit de considérer l'application [unclear: de ce groupe]
$N_0 \to B_{\rho,\sigma} \cdot \underline{C}$
$g \mapsto (i(g), g^{-1})$ [unclear]
Mais l'hypothèse
que le centre de $\underline{C}$ soit trivial n'est pas heureuse -
p.ex. elle n'est pas satisfaite si $\underline{C} = \widehat{Sl(2,\mathbb{Z})} / \pm 1$.
Il me semble que l'introduction de $N_0$ (dans
lequel on [unclear: partagera ...] )

\newpage

%% ===== Page 505 =====

n'était pas non plus judicieuse, et que la
condition $B_{\rho,\sigma} \supset N_0$ - il suffisait de $B_{\rho,\sigma} \supset L_0$.

Je vais donc considerer
$$ S_0 G_{\rho,\sigma} \overset{dfn}{=} B_{\rho,\sigma} . \underline{C} \quad \text{le produit 1/2 direct, et} $$
$$ (14) \quad \begin{cases} L_0 \overset{i}{\hookrightarrow} S_0 G_{\rho,\sigma} \\ g \longmapsto i(g).g^{-1} \end{cases} $$
l'image de cet homomorphisme est un sous-groupe
distingué, car a) il commute à $\underline{C}$ (trivial) et
b) il est invariant par $B_{\rho,\sigma}$, car pour $u \in B_{\rho,\sigma}$
$$ int(u)(i(g)) = int(u) (i(g).g^{-1}) = int(u)(i(g)) . \underbrace{int(u)(g^{-1})}_{=u(g)^{-1}} = i(u(g)) $$
(i.e. $i$ commute aux actions de $B_{\rho,\sigma}$ , agissant sur
$L_0$ par restriction à $L_0$ de l'op. de $B_{\rho,\sigma}$ sur
$\underline{C}$, et sur $S_0 G_{\rho,\sigma}$ par autom. intérieurs ).

On peut donc considerer
$$ (15) \quad G_{\rho,\sigma} \overset{dfn}{=} S_0 G_{\rho,\sigma} / i(L_0) $$

Alors les hom. d'inclusion
$$ B_{\rho,\sigma} \hookrightarrow S_0 G_{\rho,\sigma} , \quad \underline{C} \hookrightarrow S_0 G_{\rho,\sigma} $$
donnent, par composition avec $S_0 G_{\rho,\sigma} \twoheadrightarrow G_{\rho,\sigma}$,
des hom.
$$ (16) \quad \underline{C} \to G_{\rho,\sigma} , \quad B_{\rho,\sigma} \to G_{\rho,\sigma} $$
rendant commutatif le diagramme
(17)
\begin{tikzcd}
L_0 \ar[r, hook] \ar[d, hook] & \underline{C} \ar[d] \\
B_{\rho,\sigma} \ar[r] & G_{\rho,\sigma}
\end{tikzcd}

On note :
$$ (18) \quad \operatorname{Ker}(\underline{C} \to G_{\rho,\sigma}) = L_0 \cap \text{Centr}(\underline{C}) = \operatorname{Ker}(L_0 \to B_{\rho,\sigma}) $$

\newpage

%% ===== Page 506 =====

Dans le cas $\underline{G} = S \ell(2, \hat{\mathbb{Z}})'$, (où $\rho,\sigma$ comme à l'ordinaire
l'on a $L_0 \cap Cent(\underline{G}) = \{1\}$. Donc il
ne semble pas a priori qu'il y ait des inconvé-
nients majeurs à supposer que

$$ (18') \quad L_0 \cap Cent(\underline{G}) = \{1\} \text{ i.e. } \underline{G} \hookrightarrow G_{\rho,\sigma} \text{ inj.} $$

[MARGIN: i.e. l'un des 2 diagrammes exacts
(10bis)
\begin{tikzcd}
& & 1 \ar[d] & 1 \ar[d] \\
& & L_0 \ar[d] \ar[r, equal] & L_0 \ar[d] \\
1 \ar[r] & \underline{G} \ar[r] \ar[d, equal] & S G_{\rho,\sigma} \ar[r] \ar[d] & B_{\rho,\sigma} \ar[r] \ar[d] & 1 \\
1 \ar[r] & \underline{G} \ar[r] & G_{\rho,\sigma} \ar[r] \ar[d] & \Gamma_{\rho,\sigma} \ar[r] \ar[d] & 1 \\
& & 1 & 1
\end{tikzcd}
]

Notons que $L_0$ est l'image de
$$ L_0 \to B_{\rho,\sigma} $$
$$ \ell \mapsto i(\ell) $$
Composée de (14) et de la projection $S G_{\rho,\sigma} \to B_{\rho,\sigma}$)
et évidemment $L_0 \cap \underline{G} = 1$, et posant

$$ (19) \quad \Gamma_{\rho,\sigma} \stackrel{dfn}{=} B_{\rho,\sigma} / \text{Im } L_0 $$

on trouve

$$ (20) \quad G_{\rho,\sigma} / \underline{G} \simeq B_{\rho,\sigma} / L_0 \simeq \Gamma_{\rho,\sigma} $$

NB Dans le cas $\underline{G} = S \ell(2, \hat{\mathbb{Z}})'$, $B_{\rho,\sigma}$ est,
(à [unclear: $C P(\hat{\mathbb{Z}})$] près le groupe noté précédemment $\tilde{B}_0'$
ou plutôt un groupe isomorphe $\tilde{B}_0'$), et
l'on a $det \mid \tilde{B}_0' = B_{\rho,\sigma}$ isom. (trivial sur $L_0$)
induisant par passage au quotient un isomorphisme

$$ \Gamma_{\rho,\sigma} \buildrel \sim \over \to G_m $$

et on a donc, pour utiliser l'isomorphisme
$$ \tilde{B}_0' = \tilde{\Gamma}_0 \cdot L_0 \quad (\text{1/2 direct}) \quad \text{où } \tilde{\Gamma}_0 \to \tilde{B}_0' \simeq G_m $$
$$ \Gamma_{\rho,\sigma} = \tilde{\Gamma}_{\rho,\sigma} \cdot L_0 $$
pour donner

$$ (22) \quad G_{\rho,\sigma} = S \ell(2, \hat{\mathbb{Z}})' . $$

[MARGIN: dans le cas de $\underline{G} = \widehat{S \ell}(2, \mathbb{Z}), \bar{\rho}, \bar{\sigma}$ (21)
tout se transpose avec
$\underline{G} = S \ell(2, \hat{\mathbb{Z}})' \dots$]

\newpage

%% ===== Page 507 =====

556

Mais dans des cas plus généraux auxquels nous aspirons, il ne faut pas du tout s'attendre : - que $\gamma_{p,\sigma}$ soit aussi petit - p.ex. qu'il soit commutatif.
Par exemple dans le cas "universel", $\underline{G} = \underline{G}_{0,3}^{+\wedge}$,
on aura

(23) $$B_{p,\sigma} = M_{0,3}^{\sim}[0]' = \operatorname{Ker}(M_{0,3}^{\sim}[0] \xrightarrow{\varepsilon_3} \{ \pm 1 \} )$$

donc

(24) $$\gamma_{p,\sigma} \simeq \Gamma_{0,3}^{\sim '} \quad (\supset \Gamma_Q' !)$$

Par contre, on doit pouvoir dans ce cas universel postuler l'existence d'un scindage (canonique) de l'extension $B_{p,\sigma}$ de $\gamma_{p,\sigma}$ par $L_0$. En fait, si on suppose, dans le cas général, où on a $\sigma^2=p^3=1$, i.e. $\underline{G}$ quotient de $\underline{G}_{0,3}^{+\wedge}$, que le quotient $\underline{G}$ de $\underline{G}_{0,3}^{+\wedge}$ est sans quotient fini s'envoyant dans $CPL(\hat{\mathbb{Z}})$, alors [unclear: l'hom.] $B_{p,\sigma} \simeq \Sigma_{p,\sigma}$ d'où s'envoie dans $\tilde{B}_{0,3}'$, et l'image inverse de $\tilde{T}_0$ est un sous-groupe (qu'on pourrait noter $T_{p,\sigma}$)

(25) $$T_{p,\sigma} \subset B_{p,\sigma} \quad T_{p,\sigma} \simeq \gamma_{p,\sigma}$$

qui est un sous-groupe section, pour $B_{p,\sigma} \to \gamma_{p,\sigma}$, et aussi un sous-groupe de $G_{p,\sigma}$)
idem $G_{p,\sigma} \to \gamma_{p,\sigma}$, de sorte qu'on a

(26) $$B_{p,\sigma} \simeq T_{p,\sigma} . L_0 , \quad G_{p,\sigma} \simeq T_{p,\sigma} . \underline{G} \quad \text{(produits 1/2 directs)}$$

\newpage

%% ===== Page 508 =====

Soit maintenant

(27) 
$$S\mathcal{G}_{\rho,\sigma} \subset \bar{\mathcal{G}} \times \bar{\mathcal{G}}$$
$$\parallel dfn$$
$$\{ (\alpha, \beta) \mid \alpha \rho \alpha^{-1} = \beta \sigma \beta^{-1} \varepsilon_1^\nu \text{ , pour } \nu \in \hat{\mathbb{Z}} \text{ convenable} \}$$
i.e. $\eta^{-1} \xi \in L_1$

On a une application canonique

(28) 
$$S\mathcal{G}_{\rho,\sigma} \longrightarrow \Sigma_{\rho,\sigma}$$
$$(\alpha, \beta) \longmapsto (\xi, \eta) , \quad \xi = \alpha \rho \alpha^{-1} \ , \ \eta = \beta \sigma \beta^{-1}$$

qui est surjective par définition de $\Sigma_{\rho,\sigma}$ ,
et en fait par l'hypothèse (10) pourvue d'une
section canonique

$$\Sigma_{\rho,\sigma} \longrightarrow S\mathcal{G}_{\rho,\sigma}$$
$$(\xi, \eta) \longmapsto u_{\xi, \eta}$$

à cela près que $u_{\xi, \eta}$ n'est pas dans $\bar{\mathcal{G}}$, mais
seulement dans $G_{\rho,\sigma}$. Nous sommes conduits
dans l'ensemble plus gros

(29) 
$$\mathcal{G}_{\rho,\sigma} \subset G_{\rho,\sigma} \times G_{\rho,\sigma}$$
$$\parallel dfn$$
$$\{ (\alpha, \beta) \mid \alpha \rho \alpha^{-1} = \beta \sigma \beta^{-1} \varepsilon_1^\nu \text{ pour } \nu \in \hat{\mathbb{Z}} \text{ conv.} \}$$

[MARGIN: NB on a ici $\xi = \alpha \rho \alpha^{-1}, \eta = \beta \sigma \beta^{-1}$ donc $\xi, \eta \in \bar{\mathcal{G}}$ (puisque $\bar{\mathcal{G}}$ inv. dans $G_{\rho,\sigma}$)]

de sorte que l'on aura

(30) 
$$S\mathcal{G}_{\rho,\sigma} = \mathcal{G}_{\rho,\sigma} \cap (\bar{\mathcal{G}} \times \bar{\mathcal{G}})$$

(sans supposer (18) pour simplifier).

\newpage

%% ===== Page 509 =====

Ainsi l'on trouve une
$$ (31) \quad g_{\rho, \sigma} \longrightarrow \Sigma_{\rho, \sigma} $$
$$ (\alpha, \beta) \longmapsto (\xi, \eta), \quad \begin{matrix} \xi = \alpha\rho\alpha^{-1} = [\alpha, \rho] = \alpha(\rho) \rho^{-1} \\ \eta = \beta\sigma\beta^{-1} = [\beta, \sigma] = \beta(\sigma) \sigma^{-1} \end{matrix} $$

(par définition de $B_{\rho, \sigma}$ on trouve bien que les $(\xi, \eta)$ associés aux $(\alpha, \beta)$ sont dans $\Sigma_{\rho, \sigma}$). Cette fois-ci l'hyp. (10) implique qu'on a une section canonique

$$ (32) \quad \Sigma_{\rho, \sigma} \longrightarrow g_{\rho, \sigma} $$
$$ (\xi, \eta) \longmapsto (u_{\xi, \eta}, u_{\xi, \eta}) $$
et on trouve comme au §10 :

Théorème Considérons les sous-groupes de $G_{\rho, \sigma}$ :
[MARGIN: NB dans le cas $\underline{S} = \text{Ise}(Z, \tilde{Z})$ etc, on a $Z_\rho = T_\rho$, $Z_\sigma = N_\sigma^\circ$]
$$ (33) \quad \begin{cases} Z_\rho = Cent_{G_{\rho, \sigma}}(\rho) \\ Z_\sigma = Cent_{G_{\rho, \sigma}}(\sigma) \end{cases} $$
et le groupe produit
$$ B_{\rho, \sigma} \times Z_\rho \times Z_\sigma $$
alors la flèche est bijective
$$ (34) \quad B_{\rho, \sigma} \times Z_\rho \times Z_\sigma \xrightarrow{\sim} g_{\rho, \sigma} $$
$$ (u, a, b) \longmapsto (u a^{-1}, u b^{-1}) $$

[MARGIN: par définition de $B_{\rho, \sigma}$ il suit que $Z_\rho \subset B_{\rho, \sigma}$, donc on a [unclear] l'hyp (33) ce n'est pas bon !]

Corollaire Considérons les hom. de groupes
$$ \delta_G : G_{\rho, \sigma} \to \bar{G}_{\rho, \sigma} , \delta_\rho : Z_\rho \to \bar{G}_{\rho, \sigma} , \delta_\sigma : Z_\sigma \to \bar{G}_{\rho, \sigma} , $$
où $\delta_G$ l'hom canonique de passage au quotient $\bar{G}_{\rho, \sigma} = \bar{G}$, et $\delta_\rho, \delta_\sigma$ les composés de $\delta_G$ avec l'inclusion $Z_\rho \subset G_{\rho, \sigma}$, $Z_\sigma \subset G_{\rho, \sigma}$, $B_{\rho, \sigma} \subset G_{\rho, \sigma}$

\newpage

%% ===== Page 510 =====

[MARGIN: utile tjrs ds la suite L. §48, IX]

(36) déf. $S\Gamma_{\rho,\sigma} \subset B_{\rho,\sigma} \times Z_{\rho} \times Z_{\sigma}$
$\{ (u, a, b) \mid \delta_G(u) = \delta_{\rho}(a) = \delta_{\sigma}(b) \}$
une bijection dans $S\Gamma_{\rho,\sigma} = B_{\rho,\sigma} \times_{\Gamma_{\rho,\sigma}} Z_{\rho} \times_{\Gamma_{\rho,\sigma}} Z_{\sigma}$.

(37) $S\Gamma_{\rho,\sigma} \longrightarrow S\mathcal{G}_{\rho,\sigma} \quad \text{(défini ds (27))}$
$(u, a, b) \longmapsto (\xi, \eta), \text{ où } \xi = u a^{-1}, \eta = u b^{-1}$

[MARGIN: NB
la lettre S ici est inspirée
de l'usage du $S$ dans $SL_n$, $SP_n$, $SO_n$,
pour dire "gpe sur les ($Z_{\rho}$, $Z_{\sigma}$,...)
spécial", et ici pour $S\Gamma_{\rho,\sigma}$
= sur gpe qui passe par
le centre de $\Gamma_{\rho,\sigma}$
ds un hom. de groupes...
cf L. §48 pour détailler...]

Scolie Par les bijections (34) (38), on
transporte sur $\mathcal{G}_{\rho,\sigma}$, $S\mathcal{G}_{\rho,\sigma}$ les structures
de groupes naturelles de $\Gamma_{\rho,\sigma}$, $S\Gamma_{\rho,\sigma}$.
Expliciter cette structure de groupe. on
trouve notamment pour $S\mathcal{G}_{\rho,\sigma}$

(39) $\begin{cases} (\alpha', \beta') (\alpha, \beta) = (\alpha'' \rho, \beta'' \sigma), \text{ avec } \alpha'' = u'(\alpha)\alpha', \beta'' = u'(\beta)\beta' \\ u' = u_{(\xi,\eta)} : \underline{C} \to \underline{C}, \text{ défini par } \begin{cases} u(\rho) = \xi \rho = \alpha' \alpha \rho \\ u(\sigma) = \eta \sigma = \beta' \beta \sigma \end{cases} \end{cases}$

Dans $S\Gamma_{\rho,\sigma}$ on a ... avec $\xi = \alpha' \alpha^{-1}, \eta = \beta' \beta^{-1}$

[CROSSED OUT:
(39) $L_0 \to S\Gamma_{\rho,\sigma}$
$u \mapsto (u, 1, 1)$
et on a suite exacte
(40) $1 \to L_0 \to S\Gamma_{\rho,\sigma} \to \Gamma_{\rho,\sigma} \to 1$
...
(41) $\Gamma_{\rho,\sigma} = Z_{\rho} \times_{\Gamma_{\rho,\sigma}} Z_{\sigma}$
(42) $1 \to S Z_{\rho} \times S Z_{\sigma} \to S\Gamma_{\rho,\sigma} \to \Gamma_{\rho,\sigma} \to 1$
]

\newpage

%% ===== Page 511 =====

(40) 
$$L_0 \longrightarrow S\underline{\Gamma}_{\rho,\sigma}^{\circ} \subset B_{\rho,\sigma} \times Z_\rho \times Z_\sigma$$
$$l \longmapsto (l, 1, 1)$$

et l'action de $L_0$ par translation à g. sur $S\underline{\Gamma}_{\rho,\sigma}^{\circ}$
devient, par (39), l'action sur $S\mathcal{g}_{\rho,\sigma}$ donnée par

(41) $$(\alpha, \beta) \longmapsto (l\alpha, l\beta)$$

[MARGIN: intr $\underline{\Gamma}_{\rho,\sigma}^{\circ}$ ds §48, $\underline{\underline{IX}}$]

Posons
(42) $$\underline{\Gamma}_{\rho,\sigma}^{\circ} \overset{df}{=} S\underline{\Gamma}_{\rho,\sigma}^{\circ} / Im L_0 \simeq Z_\rho \times_{\Gamma_{\rho,\sigma}} Z_\sigma \ ,$$

on trouve que (38) induit une bijection
canonique
(43) $$\underline{\Gamma}_{\rho,\sigma}^{\circ} \overset{\sim}{\longrightarrow} L_0 \backslash S\mathcal{g}_{\rho,\sigma}$$

[unclear crossed out text]

Posant $\overset{df}{=}$
(44) $$ \begin{cases} S Z_{\sigma} = \operatorname{Ker}(Z_{\sigma} \to \Gamma_{\rho,\sigma}) = Cent_{\underline{\Gamma}_{\hat{\mathbb{Z}}}}(\rho) \\ S Z_{\rho} = \operatorname{Ker}(Z_{\rho} \to \Gamma_{\rho,\sigma}) = Cent_{\underline{\Gamma}_{\hat{\mathbb{Z}}}}(\sigma) \end{cases} $$

[MARGIN: introduire aussi $S\Gamma_{\rho,\sigma} = \operatorname{Ker}(\Gamma_{\rho,\sigma} \to \bar{\Gamma}_{\rho,\sigma})$
$S\mathcal{g}_{\rho,\sigma} = \operatorname{Ker}(\mathcal{g}_{\rho,\sigma} \to \bar{\mathcal{g}}_{\rho,\sigma})$]

on a des suites exactes

(45) $$ \begin{cases} 1 \to S Z_{\sigma} \to Z_{\sigma} \to \Gamma_{\rho,\sigma} \to 1 \\ 1 \to S Z_{\rho} \to Z_{\rho} \to \Gamma_{\rho,\sigma} \to 1 \end{cases} $$

et les expressions (37), (42) de $S\underline{\Gamma}_{\rho,\sigma}^{\circ}$, $\underline{\Gamma}_{\rho,\sigma}^{\circ}$
donnent

(46) $$1 \longrightarrow L_0 \times S Z_\rho \times S Z_\sigma \longrightarrow S\underline{\Gamma}_{\rho,\sigma}^{\circ} \longrightarrow \Gamma_{\rho,\sigma} \longrightarrow 1$$
(47) $$1 \longrightarrow S Z_\rho \times S Z_\sigma \longrightarrow \underline{\Gamma}_{\rho,\sigma}^{\circ} \longrightarrow \Gamma_{\rho,\sigma} \longrightarrow 1 \ .$$

NB Dans le cas $\underline{\mathcal{C}} = \operatorname{Sl}(2,\hat{\mathbb{Z}})$, $\rho$ et $\sigma$ "involutions", on
retrouve les $S_0\underline{\Gamma}_{\rho,\sigma}^{\circ}$ et $\underline{\Gamma}_{\rho,\sigma}^{\circ}$ de §48 $\underline{\underline{IX}}$

510.

\newpage

%% ===== Page 512 =====

559

Soit maintenant

(49)
$$S\mathcal{M}_{\rho,\sigma}^0 \subset S\Gamma_{\rho,\sigma}^0 \times G_{\rho,\sigma} \quad (\subset B_{\rho,\sigma} \times \mathbb{Z}_{\rho} \times \mathbb{Z}_{\sigma} \times G_{\rho,\sigma})$$
$$|| dfn$$
$$\{(u, a, b, U) \mid \delta_B(u) = \delta_{\rho}(a) = \delta_{\sigma}(b) = \delta_G(U)\}$$
(relations dans $\Gamma_{\rho,\sigma}$ cf (19))

(50)
$$\mathcal{M}_{\rho,\sigma}^0 \subset \mathbb{Z}_{\rho} \times \mathbb{Z}_{\sigma} \times G_{\rho,\sigma}$$
$$|| dfn$$
$$\{(a, b, U) \mid \delta_{\rho}(a) = \delta_{\sigma}(b) = \delta_G(U)\}$$
(relations dans $\Gamma_{\rho,\sigma}$)

On a donc un hom. canonique

(51)
$$S\mathcal{M}_{\rho,\sigma}^0 \longrightarrow \mathcal{M}_{\rho,\sigma}^0$$
$$(u, a, b, U) \longmapsto (a, b, U)$$

dont le noyau est formé des $(u,1,1,1)$ avec $u \in B_{\rho,\sigma}$ tel que $\delta_B(u)=1$ i.e. $u \in L_0$, d'où une suite exacte

(52)
$$1 \longrightarrow L_0 \longrightarrow S\mathcal{M}_{\rho,\sigma}^0 \longrightarrow \mathcal{M}_{\rho,\sigma}^0 \longrightarrow 1$$

et des hom. canoniques

(53)
$$ \mathcal{M}_{\rho,\sigma}^0 \begin{matrix} \xrightarrow{\Theta_G} & G_{\rho,\sigma} \\ \xrightarrow{\Theta_{\rho}} & \mathbb{Z}_{\rho} \\ \xrightarrow{\Theta_{\sigma}} & \mathbb{Z}_{\sigma} \end{matrix} $$
$$ (resp. \quad (a,b,U) \begin{matrix} \mapsto & U \\ \mapsto & a \\ \mapsto & b \end{matrix} ) $$

redonnant les hom. analogues sur $S\mathcal{M}_{\rho,\sigma}^0$ par composition avec (51) , et de plus on a un hom. canonique

(54)
$$S\mathcal{M}_{\rho,\sigma}^0 \xrightarrow{\Theta_B} B_{\rho,\sigma}$$
$$(u,a,b,U) \longmapsto u$$

511

\newpage

%% ===== Page 513 =====

On a ainsi des isomorphismes

(55) $$S \mathcal{M}_{\rho, \sigma}^0 \simeq \mathcal{M}_{\rho, \sigma}^0 \times_{\Gamma_{\rho, \sigma}} B_{\rho, \sigma} ;$$

et $S \mathcal{M}_{\rho, \sigma}^0$ [unclear: s'exprime] : l'une de $\mathcal{M}_{\rho, \sigma}^0$ et $B_{\rho, \sigma}$, 
d'où une loi qui s'obtient de $\mathcal{M}_{\rho, \sigma}^0$ par $L_0$ 
ce qui provient de l'extension $B_{\rho, \sigma}$ de 
$\Gamma_{\rho, \sigma}$ par $L_0$. Et on a

(56) $$\mathcal{M}_{\rho, \sigma}^0 \simeq Z_\rho \times_{\Gamma_{\rho, \sigma}} Z_\sigma \times_{\Gamma_{\rho, \sigma}} G_{\rho, \sigma}$$

d'où on déduit (homom $\delta_\rho, \delta_\sigma, \delta_G$ sur $\Gamma_{\rho, \sigma}$)

(57) $$1 \to \underbrace{S Z_\rho \times S Z_\sigma \times \mathcal{C}}_{\overset{df}{=} \mathcal{M}_{\rho, \sigma}^{0+}} \to \mathcal{M}_{\rho, \sigma}^0 \xrightarrow{\delta_M} \Gamma_{\rho, \sigma} \to 1$$

où [unclear crossed out text] (44)
$$ \begin{cases} S Z_\rho = Z_\rho \cap \mathcal{C} = \underline{Cent}_{\mathcal{C}}(\rho) \\ S Z_\sigma = Z_\sigma \cap \mathcal{C} = \underline{Cent}_{\mathcal{C}}(\sigma) \end{cases} $$

De même on a

(49) $$S \mathcal{M}_{\rho, \sigma}^0 \simeq B_{\rho, \sigma} \times_{\Gamma_{\rho, \sigma}} Z_\rho \times_{\Gamma_{\rho, \sigma}} Z_\sigma \times_{\Gamma_{\rho, \sigma}} G_{\rho, \sigma}$$

d'où une suite exacte canonique :

(50) $$1 \to \underbrace{L_0 \times S Z_\rho \times S Z_\sigma \times \mathcal{C}}_{S \mathcal{M}_{\rho, \sigma}^{0+}} \to S \mathcal{M}_{\rho, \sigma}^0 \to \Gamma_{\rho, \sigma} \to 1$$

Considérons l'hom. can.

(51) $$\mathcal{C} \xrightarrow{i_0} \mathcal{M}_{\rho, \sigma}^0$$
$$u \mapsto (1, 1, u)$$

\newpage

%% ===== Page 514 =====

son image est le troisième facteur dans
le noyau de (57), et on a une suite
exacte canonique (déduite de (46) - (47))

(52)
$$1 \longrightarrow \widehat{\mathcal{C}} \longrightarrow M_{\rho, \sigma}^\circ \longrightarrow \Gamma_{\rho, \sigma}^\circ \longrightarrow 1$$
[unclear: par isom. ...]

et de même (via (49) ou 50)

(53)
$$1 \longrightarrow \widehat{\mathcal{C}} \longrightarrow S_0 M_{\rho, \sigma}^\circ \longrightarrow S_0 \Gamma_{\rho, \sigma}^\circ \longrightarrow 1$$

[MARGIN: (53') $1 \longrightarrow S Z_\rho \times S Z_\sigma \longrightarrow M_{\rho, \sigma}^\circ \longrightarrow \Gamma_{\rho, \sigma}^\circ \longrightarrow 1$]

Considérons l'application

(54)
$$S_0 M_{\rho, \sigma}^\circ \longrightarrow \widehat{\mathcal{C}} \times \widehat{\mathcal{C}} \times \widehat{\mathcal{C}}$$
$$\{u, a, b, U\} \longmapsto (u a^\sim, u b^\sim, u U^{-1}) = (\alpha, \beta, f)$$

on voit par le corollaire p. 557 (cf (37))
que cela induit une bijection

(55)
$$S_0 M_{\rho, \sigma}^\circ \simeq (S G_{\rho, \sigma}) \times \widehat{\mathcal{C}}$$

On a ainsi, pour un $x \in S_0 M_{\rho, \sigma}^\circ$,
d'une part les quantités

(56)
$$u, a, b, U$$
$$\in \quad \in \quad \in \quad \in$$
$$B_{\rho, \sigma} \quad Z_\rho \quad Z_\sigma \quad G_{\rho, \sigma}$$

associés, d'autre part les éléments

(57)
$$\alpha, \beta, f, \xi, \eta, \alpha_0, \beta_0, g, h \text{ dans } \widehat{\mathcal{C}}$$

[unclear: les $(\alpha, \beta, f)$] définis en termes de (56), dans (37), et on a

(58)
$$ \begin{cases} \xi = a \rho(a)^{-1}, \eta = b \sigma(b)^{-1}, \alpha_0 = f^{-1} \alpha \xi, \beta_0 = f^{-1} \beta \\ g = \alpha_0 \rho(\alpha_0)^{-1} = f^{-1} \xi \rho(f), h = \beta_0 \sigma(\beta_0)^{-1} = f^{-1} \eta \sigma(f) \end{cases} $$

\newpage

%% ===== Page 515 =====

On aura, en termes de $\alpha, \beta, f$, les expressions
suivantes de l'action de $U$ sur $\widehat{\mathcal{G}}$ :

(59) $$ \left\{ \begin{array}{l} U(\rho) = int(f^{-1})(\xi \rho) \\ U(\sigma) = int(f^{-1})(\eta \sigma) \end{array} \right. $$

Remarques Finalement, dans tous ces
développements, on n'a pas vraiment
utilisé la soi-disante "hyp. fondamentale"
(20), sauf la partie a). En l'absence
de b), il y a lieu d'introduire

(60) $$ \Sigma_{\rho,\sigma}^{eff} \subset \Sigma_{\rho,\sigma} \subset \widehat{\mathcal{G}} \times \widehat{\mathcal{G}} $$
$$ || $$
$$ \{ (\xi,\eta) \in \widehat{\mathcal{G}} \times \widehat{\mathcal{G}} \mid \exists u \in B_{\rho,\sigma}, u(\rho) = \xi\rho, u(\sigma) = \eta\sigma \} $$

qui n'est autre que l'image de l'
application [unclear: (1)] : couples $f \in \widehat{\mathcal{G}}, \xi, \eta$
satisfaisant les conditions définissant $\Sigma_{\rho,\sigma}$ (28)

[unclear: d'ailleurs il n'y a rien à changer, dans]
notons

(61) $$ S\mathcal{G}_{\rho,\sigma}^{eff} = \text{Image rec de } \Sigma_{\rho,\sigma}^{eff} \text{ par } S\mathcal{G}_{\rho,\sigma} \to \Sigma_{\rho,\sigma} \quad (20) $$

l'image inverse
et c'est sur $S\mathcal{G}_{\rho,\sigma}^{eff}$ qu'on aura une [unclear: structure]

\newpage

%% ===== Page 516 =====

du groupe isomorphe : $S_0 \tilde{G}_{\rho,\sigma}$, par l'isomor-
phisme (qui remplace (38))
(62) $$S_0 \tilde{G}_{\rho,\sigma} \overset{\sim}{\longrightarrow} S_0 G^{eff}_{\rho,\sigma}$$
$$(\underline{u}, \underline{a}, \underline{b}) \longmapsto (\underline{u} \underline{a}^{-1}, \underline{u} \underline{b}^{-1})$$

Ceci posé, revenons à $M_{0,3}^\sim$, et
la suite exacte
(63) $$1 \longrightarrow \widehat{\mathfrak{S}}_{0,3}^+ \longrightarrow M_{0,3}^\sim \longrightarrow \Gamma_{0,3}^\sim \longrightarrow 1 .$$

Considérons le sous-groupe
(64) $${\Gamma'}_{0,3}^{\sim} = \operatorname{Ker}(\Gamma_{0,3}^\sim \overset{(43)bis}{\longrightarrow} \text{[unclear]})$$
$$= \{ \gamma \in \Gamma_{0,3}^\sim \mid \mu(\gamma) = 1 \} \text{ [unclear]}$$

et soit
(65) $${M'}_{0,3}^\sim = \text{Image inverse de } {\Gamma'}_{0,3}^{\sim} \text{ dans } M_{0,3}^\sim$$

Soit $\widehat{\mathfrak{S}}$ un groupe profini, et
(66) $$\psi : \widehat{\mathfrak{S}}_{0,3}^+ \longrightarrow \widehat{\mathfrak{S}} .$$
un homomorphisme surjectif - ce qui revient
à se donner dans $\widehat{\mathfrak{S}}$ deux éléments
(67) $$\rho_{\widehat{\mathfrak{S}}}, \sigma_{\widehat{\mathfrak{S}}} \in \widehat{\mathfrak{S}}$$
(les images de $\rho_{0,3}^+, \sigma_{0,3}^+$) satisfaisant
(68) $$\rho_{\widehat{\mathfrak{S}}}^3 = \sigma_{\widehat{\mathfrak{S}}}^2 = 1 .$$

\newpage

%% ===== Page 517 =====

[MARGIN: l'introduire
ces conditions
suivantes équiv. à
a) $M_{0,3}^\sim [0] \to \Sigma_{\rho,\sigma}^{eff}$
b) $\forall u \in M_{0,3}^\sim, \exists y$
tel que $uy \in M_{0,3}^\sim [0]$
b') Idem avec $u\Gamma = y\Gamma$
c) $\exists \Psi_u$ qui prolonge $\psi$
c') $\exists \Psi_G$ qui prolonge $u$
On dira dans
la suite (?)
que $\psi$ est effective
si (70) a) satisfait]

C'est-à-dire que $\psi$ est effectif si pour
tout $u \in M_{0,3}^\sim$, posant $u(\rho) = \xi \rho$, $u(\sigma) = \eta \sigma$,

$$(\xi_{\tilde{\mathbb{G}}}, \eta_{\tilde{\mathbb{G}}}) = (u(\xi), u(\eta)) \in \Sigma_{\rho,\sigma}$$

est effectif (i.e. $\in \Sigma_{\rho,\sigma}^{eff}$, i.e. $\exists\ u_{\tilde{\mathbb{G}}}$ autom.
de $\mathbb{G}$ tel qu'on ait

(69) $$u_{\tilde{\mathbb{G}}}(\rho_{\tilde{\mathbb{G}}}) = \xi_{\tilde{\mathbb{G}}} \rho_{\tilde{\mathbb{G}}} \quad , \quad u_{\tilde{\mathbb{G}}}(\sigma_{\tilde{\mathbb{G}}}) = \eta_{\tilde{\mathbb{G}}} \sigma_{\tilde{\mathbb{G}}}$$

ce qui revient encore, par commutativité dans

(70)
\begin{tikzcd}
\mathbb{G}_{0,3}^{+\wedge} \ar[r, "u"] \ar[d, "\psi"'] & \mathbb{G}_{0,3}^{+\wedge} \ar[d, "\psi"] \\
\mathbb{G} \ar[r, "u_{\tilde{\mathbb{G}}}"] & \mathbb{G}
\end{tikzcd}

(Dans le cas général où on ne l'a pas, le $u \in M_{0,3}^\sim$ effectif signifie que cet $u$ transforme
$\operatorname{Ker}\ \psi$ en lui-même, ou encore que
$M_{0,3}^\sim [0] \quad (= M_{0,3}^\sim [0] \cdot \mathbb{G}_{0,3}^{+\wedge})$ tout entier
stabilise $\operatorname{Ker}\ \psi$. Donc cette condition
est automatiquement remplie
quand $\mathbb{G}, \rho_{\tilde{\mathbb{G}}}, \sigma_{\tilde{\mathbb{G}}}$ satisfait la condition
b) de (20).

On a fait ce qu'il fallait pour avoir...

\newpage

%% ===== Page 518 =====

562

Théorème Soit $\psi : \overline{G}_{0,3}^{+ \wedge} \to \overline{G}$ un
hom surjectif effectif (cf ci-dessus).
Alors il existe un hom. unique

(71) $$ \Psi_\psi : M_{0,3}^{\prime \sim} \to M_\psi $$

[MARGIN: En fait, je crains bien - on aura $S_0 M_{0,3}^\sim \to S_0 M_\psi$]

[unclear: crossed out diagram]
(72)
\begin{tikzcd}
\overline{G}_{0,3}^{+ \wedge} \ar[r, "*"] \ar[d, "\psi"'] & M_{0,3}^{\prime \sim} \ar[d, "\Psi_\psi"] \\
\overline{G} \ar[r, "*"] & M_{\bar{\rho}, \bar{\sigma}} \quad (= M_\psi)
\end{tikzcd}

$\Psi_\psi$ a la propriété suivante: Si 
$u \in M_{0,3}^{\prime \sim}$ est de la forme $u = \alpha \beta \gamma$
($\alpha, \beta, \gamma \in \overline{G}_{0,3}^{+ \wedge}$, satisfaisant aux équations
des lacets [unclear: appelés = polygonaux] ), alors
$\Psi_\psi(u)$ est l'élément de $M_\psi$
défini par les invariants $\psi(\alpha), \psi(\beta), \psi(\gamma)$.

Corollaire On a un diag. de suites exactes

(72)
\begin{tikzcd}
1 \ar[r] & \overline{G}_{0,3}^{+ \wedge} \ar[r] \ar[d, "\psi"'] & M_{0,3}^{\prime \sim} \ar[r] \ar[d, "\Psi_\psi"'] & \Gamma_{0,3}^{\prime \sim} \ar[r] \ar[d, "\overline{\psi}"'] & 1 \\
1 \ar[r] & \overline{G} \ar[r] & M_\psi \ar[r] & \Gamma_\psi \ar[r] & 1
\end{tikzcd}

L'hom. $\Psi_\psi$ s'explicite à l'aide
de trois homomorphismes

517

\newpage

%% ===== Page 519 =====

(73)
$$ \psi_{T_\rho} : \Gamma_Q^{\prime \sim} \to T_\rho $$
$$ \psi_{T_\sigma} : \Gamma_Q^{\prime \sim} \to T_\sigma $$

[MARGIN: NB on écrit pour simplifier $\rho, \sigma$ au lieu de $\rho_{\overline{Q}}, \sigma_{\overline{Q}}$ en indices]

(74)
$$ \psi_G : M_{0,3}^{\prime \sim} \to G_{\rho,\sigma} $$

satisfaisant les conditions : pour $x \in M_{0,3}^{\prime \sim}$, $\dot{x}$ son image dans $\Gamma_Q^{\prime \sim}$,
~~$a = \psi_{T_\rho}(\dot{x}), b = \psi_{T_\sigma}(\dot{x}), U = \psi_G(x)$~~

(75)
$$ a = \psi_{T_\rho}(\dot{x}), \quad b = \psi_{T_\sigma}(\dot{x}), \quad U = \psi_G(x) \quad , $$

(76)
$$ \delta_\rho(a) = \delta_\sigma(b) = \delta_G(U) \quad \text{dans} \quad \Phi_{\rho,\sigma} $$

i.e. les trois homs. composés

(77)
\begin{tikzcd}
& \Gamma_Q^{\prime \sim} \ar[r, "\psi_{T_\rho}"] & T_\rho \ar[dr, "\delta_\rho"] \\
M_{0,3}^{\prime \sim} \ar[ur, "can"] \ar[r, "can"] \ar[dr, "\psi_G"'] & \Gamma_Q^{\prime \sim} \ar[r, "\psi_{T_\sigma}"] & T_\sigma \ar[r, "\delta_\sigma"] & \Phi_{\rho,\sigma} \\
& G_{\rho,\sigma} \ar[urr, "\delta_G"']
\end{tikzcd}

sont égaux. De plus, on aura d'après (73)

(78)
$$ \psi_G | \mathbb{G}_{0,3}^{+ \wedge} = \Psi_{\mathfrak{S}} \quad , $$

et enfin ceci :
(79) Si $x \in M_{0,3}^{\prime \sim}$, $U = \psi_G(x)$ , on a 
commutativité dans

(80)
\begin{tikzcd}
\mathbb{G}_{0,3}^{+ \wedge} \ar[r, "\psi"] \ar[d, "int(x) | \mathbb{G}_{0,3}^{+ \wedge}"'] & \mathfrak{S} \ar[d, "int(U) | \mathfrak{S}"] \\
\mathbb{G}_{0,3}^{+ \wedge} \ar[r, "\psi"'] & \mathfrak{S}
\end{tikzcd}

\newpage

%% ===== Page 520 =====

563

On peut reprendre tout ceci comme
énoncé de fonctorialité des constructions
faites, p. ex. si un hom. surjectif

[MARGIN: 
il aurait
mieux valu
prendre
$\underline{G}' \xrightarrow{\varphi} \underline{G}$
pour
spécialiser
$\underline{G}$
]

(81) $$ \underline{G} \xrightarrow{\varphi = \varphi_{\underline{G}}} \underline{G}' $$

transformant le couple $(\rho, \sigma)$ de générateurs
topologiques de $\underline{G}$ en un couple $(\rho', \sigma')$.
Il y a lors un unique homomorphisme

(82) $$ S_0\mathcal{M}_{\underline{G}} \xrightarrow{\varphi_{S_0\mathcal{M}}} S_0\mathcal{M}_{\underline{G}'} $$

satisfaisant les conditions suivantes :

[MARGIN:
ceci même sans
supposer que
$\varphi_{S_0 \mathcal{M}}$ se déduise
d'un hom.
$\varphi_{\mathcal{M}}$,
pourvu qu'on
se donne les
$\varphi_B, \varphi_\rho, \varphi_\sigma, \varphi_G$
]

a) qui [unclear] sur les quotients [unclear] les hom. :
(83)
$$ \varphi_B : B_{\rho, \sigma} \to B_{\rho', \sigma'} $$
$$ \varphi_\rho : Z \to Z $$
$$ \varphi_\sigma : Z \to Z $$
$$ \varphi_G : G_{\rho, \sigma} \to G_{\rho', \sigma'} ; $$

b) Compatibilité avec $\varphi_G$ (ce qui montre que la connais-
sance de $\varphi_{S_0 \mathcal{M}}$ en revient à celle de $\varphi_G$) ;

c) On a commutativité dans

(85)
\begin{tikzcd}
\underline{G} \ar[r] \ar[d, "\varphi = \varphi_{\underline{G}}"'] & G_{\rho, \sigma} \ar[d, "\varphi_G"] \\
\underline{G}' \ar[r] & G_{\rho', \sigma'}
\end{tikzcd}

(ce qui montre, puisque $G_{\rho, \sigma} = B_{\rho, \sigma} \cdot \underline{G}$, que
$\varphi_G$ est [unclear] quand on connait $\varphi_B$ ;)

d) Pour $u \in B_{\rho, \sigma}$, désignant par $u' = \varphi_B(u)$, on a
commutativité dans
(86)
\begin{tikzcd}
\underline{G} \ar[r, "u"] \ar[d, "\varphi"'] & \underline{G} \ar[d, "\varphi"] \\
\underline{G}' \ar[r, "u'"] & \underline{G}'
\end{tikzcd}

519

\newpage

%% ===== Page 521 =====

-(ce qui détermine $\varphi_B$ en termes de $\varphi$ et $\varphi_G$,
donc achève de déterminer $\varphi_{S\mathcal{M}}$, et par passage aux
quotients, -.-.-

(87) $$ \varphi_M : M_{p,\sigma} \to M_{p',\sigma'} $$

Une autre façon, plus directe, de décrire "paramètres"
$\varphi_{S_0 M_{p,\sigma}}$ ou $\varphi_M$ en termes des coordonnées
$(\alpha, \beta, f)$ au lieu de $(u, a, b, U)$ , est de
dire qu'elle transforme le $u_{\alpha,\beta,f}$ en $u_{\alpha',\beta',f'}$ , où les
$\alpha', \beta', f'$ sont les images de $\alpha, \beta, f$ par $\varphi$ :

(88) $$ \varphi_{S_0 M_{p,\sigma}}(u_{\alpha,\beta,f}) = u_{\alpha',\beta',f'} , \quad \begin{cases} \alpha' = \varphi(\alpha) \\ \beta' = \varphi(\beta) \\ f' = \varphi(f) \end{cases} $$

On retrouve ainsi, sous une forme
un peu plus précise, le th. précédent,
en notant le [unclear] *
[MARGIN: tout au moins / dans le "cas universel" p. ex. / de la fig. / $p'=1, \sigma'=1$]
$\mathcal{C} = \mathcal{C}_{0,3}^{+\wedge}$ (le "cas universel"), $\mathcal{C}'$
quelconque - il vaut mieux dans la situation
(89) inverser - que le rôle de $\mathcal{C}$ et $\mathcal{C}'$
[unclear] les invariants absolus
Il est bon d'expliciter dans le cas "universel" :

pour

(90) $$ \mathcal{C} = \mathcal{C}_{0,3}^{+\wedge} , \quad B_{p,\sigma} = M_{0,3}^{1\sim}[0] , \quad \Gamma_{p,\sigma} = \Gamma_{\mathbb{Q}}^{1\sim} $$
$$ G_{p,\sigma} = M_{0,3}^{1\sim} , \quad Z_p = M_{0,3}^{1\sim}(\bar{\mathbb{F}}_p) , \quad Z_\sigma = M_{0,3}^{1\sim}(\mathbb{Z}/2) $$
$$ \Gamma_{p,\sigma} = \Gamma_{\mathbb{Q}}^{1\sim} , \quad Z'_p \simeq \Gamma_{p,\sigma} , \quad Z'_\sigma \simeq \Gamma_{p,\sigma} $$

\newpage

%% ===== Page 522 =====

[MARGIN: donc $S Z_{\rho'} \simeq S Z_{\rho,\sigma}'$, $G(\hat{\mathbb{Z}}) \simeq S \Gamma_{\rho',\sigma'}$]
$$ \Gamma_{\rho',\sigma'} \simeq \Gamma_{\rho,\sigma}', \quad Z_{\rho'} \simeq Z_{\rho,\sigma}', \quad Z_{\sigma'} \simeq Z_{\rho,\sigma}', \quad \text{donc} $$
(92) $$ S_0 M_{\rho',\sigma'} \simeq S_0 M_{\rho,\sigma}', \quad M_{\rho',\sigma'} \simeq M_{\rho,\sigma}' $$

[MARGIN: Ce n'est pas si simple ! par ex. pour $G = \operatorname{Sl}(2,\hat{\mathbb{Z}})$ $\rho=3$, $\sigma=2$ (ou plus exactement leurs images dans $G/\{\pm 1\}$) on verrait qu'on ne peut se servir des anciennes notations.]
Remarque. En écrivant les relations (80), il devenait évident que les notations utilisées dans cette section II (depuis p. 80 jusqu'à 519) étaient inadéquates, et qu'il convient de les modifier de la façon suivante :
$$ \Sigma_{\rho,\sigma}, B_{\rho,\sigma}, G_{\rho,\sigma}, \gamma_{\rho,\sigma}, Z_{\rho}, Z_{\sigma}, \Gamma_{\rho,\sigma}, S G_{\rho,\sigma} $$
(93) $$ Z_{\rho}, Z_{\sigma}, S_0\Gamma_{\rho,\sigma}, \Gamma_{\rho,\sigma}, S Z_{\rho}, S Z_{\sigma}, S_0 M_{\rho,\sigma}, M_{\rho,\sigma} $$
On met partout un index $'$ pour en faire $\Sigma_{\rho,\sigma}', B_{\rho,\sigma}'$ etc, et on oublie l'exposant $0$ de $S_0\Gamma_{\rho,\sigma}$, $\Gamma_{\rho,\sigma}$, $S_0 M_{\rho,\sigma}$, $M_{\rho,\sigma}$ (qui se trouve donc remplacé par un $'$).

Autres cas particuliers importants explicite-
ment :

2) $G = \operatorname{Sl}(2,\hat{\mathbb{Z}})/\pm 1$, $(\rho, \sigma)$ couple modulaire.

$$ B_{\rho,\sigma}' = \left\{ \begin{pmatrix} \mu & \lambda \\ 0 & 1 \end{pmatrix} \mid \mu \in \hat{\mathbb{Z}}^*, \mu \equiv 1 \, (3), \lambda \in \hat{\mathbb{Z}} \right\} $$
$$ L_0 = \left\{ \begin{pmatrix} 1 & \lambda \\ 0 & 1 \end{pmatrix} \mid \lambda \in \hat{\mathbb{Z}} \right\} $$

[MARGIN: c'est $G(\hat{\mathbb{Z}})/ \mu_2(\hat{\mathbb{Z}})$ avec notations du § 48]
$$ G_{\rho,\sigma}' = Gl'(2, \hat{\mathbb{Z}})/ \pm 1, \quad \text{où } Gl'(2,\hat{\mathbb{Z}}) \overset{\text{def}}{=} \{ U \in \operatorname{Gl}(2,\hat{\mathbb{Z}}) \mid \det(U) \equiv 1 \, (3) \} $$

$$ Z_{\rho}' = \{ c\rho+d \mid c^2+cd+d^2 \in \hat{\mathbb{Z}}^* \} / \pm 1 $$

(94) $$ Z_{\sigma}' = \text{[unclear crossed out text]} \quad \text{avec notation du § 48} $$

$$ S_0 M_{\rho,\sigma}' \simeq S M_{\rho,\sigma}(\hat{\mathbb{Z}})/\Sigma \quad \text{avec notations du § 48} $$

$$ M_{\rho,\sigma}' \simeq M_{\rho,\sigma}(\hat{\mathbb{Z}})/\Sigma \quad id. $$

$$ \Gamma_{\rho,\sigma}' \simeq \Gamma_{\rho,\sigma}(\hat{\mathbb{Z}})/\Sigma_0 \quad \text{avec notations du § 48.} $$

\newpage

%% ===== Page 523 =====

$3^{\circ})$ Prenons $\mathbb{G} = \operatorname{Sl}(2, \hat{\mathbb{Z}})$, ou un couple modulaire

[MARGIN: (95)]
$$B'_{\rho, \sigma} = \left\{ \begin{pmatrix} \mu & \lambda \\ 0 & 1 \end{pmatrix} \mid \mu \in \hat{\mathbb{Z}}^*, \; \mu \equiv 1 \text{ (3)}, \; \lambda \in \hat{\mathbb{Z}} \right\}$$
$$G'_{\rho, \sigma} = GL'(2, \hat{\mathbb{Z}}) = \{ g \in \operatorname{Gl}(2, \hat{\mathbb{Z}}) \mid \det g \equiv 1 \text{ (3)} \}$$
$$Z'_{\rho} = \left\{ c \rho + d \mid \begin{smallmatrix} c, d \in \hat{\mathbb{Z}} \text{ t.q.} \\ c^2 - cd + d^2 \in \hat{\mathbb{Z}}^* \end{smallmatrix} \right\}$$
$$Z'_{\sigma} = T'_{\sigma} \cdot \{ 1, \tau_{\sigma} \}, \text{ où } \begin{cases} T'_{\sigma} \overset{df}{=} \left\{ t + \sigma u \mid \begin{smallmatrix} t, u \in \hat{\mathbb{Z}} \\ t^2 + u^2 \in \hat{\mathbb{Z}}^* \end{smallmatrix} \right\} \\ \tau_{\sigma} = \begin{pmatrix} -1 & 0 \\ 0 & 1 \end{pmatrix} \end{cases}$$

$S_0 \underline{M}'_{\rho, \sigma} = S_0 M'_{\rho, \sigma}(\hat{\mathbb{Z}}) / \Sigma_0$ avec notations de §48
$\underline{M}'_{\rho, \sigma} = M'_{\rho, \sigma}(\hat{\mathbb{Z}}) / \Sigma_0$
$S_0 \Gamma_{\rho, \sigma} = S_0 \Gamma'_{\rho, \sigma}(\hat{\mathbb{Z}}) / \Sigma_0$
$\Gamma_{\rho, \sigma} = \Gamma'_{\rho, \sigma}(\hat{\mathbb{Z}}) / \Sigma_0$

## III Le groupe $M_{\rho, \sigma}$

Considérons l'homomorphisme

[MARGIN: (1)]
$$\hat{\mathbb{Z}} \longrightarrow \mathbb{G}$$
$$n \longmapsto s_0^n \quad \text{ où } s_0 = \sigma^{-1} \rho$$
(d'image $\bar{L}_0$)

et soit $J$ l'intersection des noyaux, de telle façon que la ligne (1) se factorise par

[MARGIN: (2)]
$$(\hat{\mathbb{Z}}/J)^* \simeq \operatorname{Aut}(\bar{L}_0)$$

On a un homomorphisme évident

[MARGIN: (3)]
$$B'_{\rho, \sigma} \longrightarrow \operatorname{Aut}(\bar{L}_0) \simeq (\hat{\mathbb{Z}}/J)^*$$

qui sur $\gamma'_{\rho, \sigma}$ donne

[MARGIN: (4)]
$$\gamma'_{\rho, \sigma} \xrightarrow{\chi_{\rho, \sigma}} (\hat{\mathbb{Z}}/J)^*$$

d'où en composant avec $G'_{\rho, \sigma} \twoheadrightarrow \gamma'_{\rho, \sigma}$

[MARGIN: (4 bis)]
$$G'_{\rho, \sigma} \xrightarrow{\chi_G} (\hat{\mathbb{Z}}/J)^*$$

\newpage

%% ===== Page 524 =====

[III Considérations de...]

Considérons l'extension $G\ell(2, \hat{\mathbb{Z}})$ de $\mathcal{G}_{0,3}^\wedge$ par $\{\pm 1\}$ :

(1)
\begin{tikzcd}
& & 1 \ar[d] & 1 \ar[d] \\
1 \ar[r] & \{\pm 1\} \ar[r] \ar[d, equal] & S\ell(2, \hat{\mathbb{Z}}) \ar[r] \ar[d] & S_0 \mathcal{G}_{0,3}^\wedge \ar[r] \ar[d] & 1 \\
1 \ar[r] & \{\pm 1\} \ar[r] & G\ell(2, \hat{\mathbb{Z}}) \ar[r] \ar[d] & \mathcal{G}_{0,3}^\wedge \ar[r] \ar[d] & 1 \\
& & \{\pm 1\} \ar[r] \ar[d] & \{\pm 1\} \ar[d] \\
& & 1 & 1
\end{tikzcd}

On s'intéresse au $\tilde{N}_{1,1}$, défini comme l'image inverse de $G\ell(2, \hat{\mathbb{Z}})$ dans $G\ell(2, \hat{\mathbb{Z}})/\{\pm 1\}$ [unclear] via l'hom canonique

(2)
$$ \tilde{M}_{0,3} \longrightarrow G\ell(2, \hat{\mathbb{Z}})' = G\ell(2, \hat{\mathbb{Z}})/\{\pm 1\} $$

(cf § 4 [unclear]...), de sorte qu'on trouve un diagramme de suites exactes

(3)
\begin{tikzcd}
1 \ar[r] & \{\pm 1\} \ar[r] \ar[d, equal] & \tilde{N}_{1,1} \ar[r] \ar[d] & \tilde{M}_{0,3} \ar[r] \ar[d] & 1 \\
1 \ar[r] & \{\pm 1\} \ar[r] & G\ell(2, \hat{\mathbb{Z}}) \ar[r] & G\ell(2, \hat{\mathbb{Z}})' \ar[r] & 1
\end{tikzcd}

On trouve, pour l'image inverse dans $S_0 \mathcal{G}_{0,3}^\wedge \subset G\ell(2, \hat{\mathbb{Z}})'$ de $\tilde{M}_{0,3}$ dans $\tilde{N}_{1,1}$ , ... [unclear]

[MARGIN: NB $\tau_{1,1} \simeq G\ell(2, \hat{\mathbb{Z}})$]

\newpage

%% ===== Page 525 =====

donnant un hom.

(4) $\operatorname{Sl}(2,\Z)^\wedge \to \widetilde{\mathcal{N}}_{1,1}$

dont la restriction à $\operatorname{Sl}(2,\Z)^\wedge$ s'insère dans une suite exacte canoni-
que

(5) $$1 \to \underset{\simeq \widetilde{C}_{1,1}^{+ \wedge}}{\operatorname{Sl}(2,\Z)^\wedge} \to \widetilde{\mathcal{N}}_{1,1} \to \underset{\simeq \widetilde{\Gamma}_{0,3}}{\widetilde{\Gamma}_Q} \to 1$$

qui à son tour s'insère dans le diagramme

(6)
\begin{tikzcd}
1 \ar[r] & \operatorname{Sl}(2,\Z)^\wedge \ar[r] \ar[d, "\text{épi}"'] & \widetilde{\mathcal{N}}_{1,1} \ar[r] \ar[d, "\text{épi}"'] & \widetilde{\Gamma}_Q \ar[r] \ar[d, "\text{épi}"'] & 1 \\
1 \ar[r] & \operatorname{Sl}(2,\widehat{\Z}) \ar[r] & \operatorname{Gl}(2,\widehat{\Z}) \ar[r] & \widehat{\Z}^\ast \ar[r] & 1
\end{tikzcd}

On peut écrire 
$\simeq \widetilde{\Gamma}_Q \simeq \operatorname{Sl}(2,\Z)^\wedge \simeq \widehat{Sl(2,\Z)}/\pm 1$
(7) $\widetilde{\mathcal{M}}_{0,3} \simeq \widetilde{\mathcal{M}}_{0,3}(0) \cdot \widetilde{\Gamma}_{0,3}^{+ \wedge}$
(semi-direct)

en termes de cette décomposition, et
compte tenu qu'on a un hom.
canonique
(8) $\widetilde{\mathcal{M}}_{0,3}(0) \to \operatorname{Gl}(2,\widehat{\Z})$

qui relève l'hom. composé de (2)
avec l'iso. can.
$\widetilde{\mathcal{M}}_{0,3}(0) \xrightarrow{\sim} \dots$

\newpage

%% ===== Page 526 =====

525

(9) $$ \mathcal{N}_{1,1}^\sim \simeq \mathcal{N}_{1,1}^\sim(0) . \operatorname{Sl}(2, \mathbb{Z})^\wedge $$
$$ \mathcal{M}_{0,3}^\sim(0) \simeq \Gamma_Q^\sim \qquad \mathcal{N}_{1,1}^\sim(0) \simeq \dots $$

Ainsi on a une opération de $\Gamma_Q^\sim$
sur $\operatorname{Sl}(2, \mathbb{Z})^\wedge$
que je voudrais expliciter,
utilisant le diagramme cartésien

(10)
\begin{tikzcd}
\operatorname{Sl}(2, \mathbb{Z})^\wedge \ar[r] \ar[d, "can"'] & \widehat{\mathfrak{S}}_{0,3}^+ \simeq \operatorname{Sl}(2, \mathbb{Z})^\wedge / \pm 1 \ar[d, "induit \text{ par } (2)"] \\
\operatorname{Sl}(2, \hat{\mathbb{Z}}) \ar[r] & \operatorname{Sl}(2, \hat{\mathbb{Z}}) / \pm 1
\end{tikzcd}

Un élément $\tilde{v}$ de $\Gamma_Q^\sim$ correspond à
un couple $(\alpha, \beta)$ d'élts

(11) $$ (\alpha, \beta) \in \widehat{\Pi}_{0,3} \times \widehat{\Pi}_{0,3} $$

satisfaisant l'équation des tresses

(12) $$ \alpha \rho \alpha^{-1} = \beta \sigma \beta^{-1} l^v \quad (v \in \hat{\mathbb{Z}}) $$

qui opère sur $\widehat{\mathfrak{S}}_{0,3}^+ \simeq \operatorname{Sl}(2, \mathbb{Z})^\wedge$ par

(13)
$$ \begin{cases} u(\rho) = int_\alpha \rho = \xi \rho \quad & \xi = \alpha \rho \alpha^{-1} \\ u(\sigma) = int_\beta \sigma = \eta \sigma \quad & \eta = \beta \sigma \beta^{-1} \end{cases} $$

ce qui est compatible avec une opération
sur $\operatorname{Sl}(2, \hat{\mathbb{Z}})$ induite
par $x \mapsto u x u^{-1}$, $u \in \operatorname{Gl}(2, \hat{\mathbb{Z}})$ où l'élé-

\newpage

%% ===== Page 527 =====

définie comme l'image de $u$ par (8), de façon précise

(14) $\tilde{u} = i_{\Gamma_0} (\mu)$, où $\mu = \chi(u)$ est le multiplicateur, et

(15) $\begin{cases} i_{\Gamma_0} : \hat{G}_m \to Cl(2, \hat{\Z}) \\ \quad \quad \mu \mapsto \begin{pmatrix} \mu & 0 \\ 0 & 1 \end{pmatrix} \end{cases}$

On aura, de façon précise

[MARGIN: les soulignés désignent des éléments de $\operatorname{Sl}(2, \hat{\Z})$ [unclear] des $\alpha, \beta$ de [unclear], pour mémoire $\xi, \eta$ [unclear]]

(16) $\begin{cases} \underline{u}(\underline{\rho}) = \xi(\underline{\rho}) & \quad \text{où } \xi = \underline{\alpha} \rho \underline{\alpha}^{-1} \\ \underline{u}(\underline{\sigma}) = \eta(\underline{\sigma}) & \quad \text{où } \eta = \varepsilon \underline{\beta} \sigma \underline{\beta}^{-1} \end{cases}$

(17) $\varepsilon = \varepsilon_2(\mu) \in \pm 1 \quad \begin{cases} \varepsilon = +1 \text{ si } \mu \equiv 1 (4) \\ \varepsilon = -1 \text{ si } \mu \equiv -1 (4) \end{cases}$

Bien entendu, $\operatorname{Sl}(2, \Z)\hat{}$ est engendré topologiquement par $\rho$ et $\sigma$, et on identifie $\pm 1$ au sous-groupe central de $\operatorname{Sl}(2, \Z)$ donc de $\operatorname{Sl}(2, \Z)\hat{}$, identifiant $-g$ pour le produit de $g \in \operatorname{Sl}(2, \Z)\hat{}$ par $-1$. On aura donc dans $\operatorname{Sl}(2, \Z)\hat{}$

(18) $\begin{cases} u(\rho) = int\ \xi (\rho) = int(\tilde{\alpha})(\rho) = \xi \rho \\ u(\sigma) = \varepsilon\ int(\tilde{\beta})(\sigma) = \eta \sigma \end{cases}$

(19) $\begin{cases} \xi = \tilde{\alpha} \rho(\tilde{\alpha})^{-1}, \quad \eta = \varepsilon \tilde{\beta} \sigma(\tilde{\beta})^{-1} \\ \varepsilon = \varepsilon_2(\mu) \in \{\pm 1\} \subset \operatorname{Sl}(2, \Z)\hat{} \end{cases}$

\newpage

%% ===== Page 528 =====

où $\tilde{\alpha}$ et $\tilde{\beta}$ désignent des
relèvements arbitraires de $\alpha$ et $\beta$ dans $\operatorname{Sl}(2, \mathbb{Z})^\wedge$
- ils ne sont déterminés qu'aux signes près,
$int(\tilde{\alpha})$, $int(\tilde{\beta})$,
mais $\xi$ et $\eta$ sont déterminés sans ambiguïté
par $u$. Notons aussi
(20) $$u(\varepsilon_0) = \varepsilon_0^\mu$$
Si maintenant on a un
(20) $$N_{1,1}^\sim \xrightarrow{\sim \Psi} G$$
déduisant un hom.
(21) $$\operatorname{Sl}(2, \mathbb{Z})^\wedge \xrightarrow{i_T} \tilde{\mathbb{G}} \subset G$$
$Im \Psi$
donc on a des $\rho, \sigma$ dans $\tilde{\mathbb{G}}$ (notées par
les mêmes lettres) satisfaisant
(22) $$\rho^6 = \sigma^4 = 1, \quad \rho^3 = \sigma^2$$ et par suite
$$[\rho^3, \sigma] = [\sigma^2, \rho] = 1$$ i.e. $\rho^3 = \sigma^2$ centraux

[MARGIN: dans donc $\tilde{\mathbb{G}}$ il y a un élément central involutif, noté $-1$]

et pour $u \in N_{1,1}^\sim$, correspon-
dant à des éléments $\tilde{\alpha}, \tilde{\beta}$ de $\operatorname{Sl}(2, \mathbb{Z})^\wedge$
(déterminés aux signes près), d'où $\tilde{\xi}, \tilde{\eta}$
images de ceux-ci dans $\tilde{\mathbb{G}}$ et $\tilde{\mathbb{G}}$,
satisfaisant les relations idoines
(23) $$u(\rho) = \tilde{\xi} \rho \tilde{\xi}^{-1}, \quad u(\sigma) = \tilde{\eta} \sigma \tilde{\eta}^{-1}.$$
et faire jou-jou avec $\tilde{\alpha}$ et $\tilde{\beta}$. Mais on a vu

\newpage

%% ===== Page 529 =====

obligé, comme [unclear: on l'a vu plus haut]
$\tilde{\alpha} \notin \mathcal{S}l(2, \mathbb{Z})^\wedge$, mais $\tilde{\alpha} \in \mathcal{G}l(2,\mathbb{Z})^\wedge$, de
façon à faire [unclear: intervenir] l'élément $\tau_0$ de $\mathcal{G}l(2,\mathbb{Z})^\wedge$,
[crossed out: l'action de] et [unclear] sgr $\tau$ de $G$ sur $\widetilde{\mathfrak{S}}$. On
aura, en posant maintenant

[MARGIN: ce qui m'amène à prendre la convention peu habituelle [unclear: au présent]...]
(24) $\varepsilon_0 \tilde{\alpha} = \tilde{\alpha} \varepsilon_0$, $\quad \varepsilon_1 \tilde{\alpha} = \rho \tilde{\alpha} \sigma \tilde{\alpha} = \sigma (\varepsilon_0 \rho) = \rho (\varepsilon_0 \sigma)$
$= - \sigma \rho \tilde{\alpha} \qquad = - \rho \sigma \tilde{\alpha}$

les relations,

[MARGIN: je laisse tomber les indices $0$]
(25) $\tau(\sigma) = \sigma^{-1}$, $\quad \tau(\rho) = \sigma(\rho^{-1})$

qui impliquent, via (24) et (22)

(26) $\tau(\varepsilon_0) = \varepsilon_0^{-1}$, $\quad \tau(\varepsilon_1) = \varepsilon_1^{-1}$

et on introduit le groupe
(27) [crossed out: illisible] $\quad G \rtimes \{1, \tau\} = \widetilde{\mathfrak{S}}$
[MARGIN: 1/2 direct]

où $\{1, \tau\}$ opère sur $G$ par les formules (25).

Ceci dit, on a une vue d'ensemble
pour plusieurs des constructions dejà faites
par ailleurs.

\newpage

%% ===== Page 530 =====

568

IV Soit

Soit $\overline{G}$ un groupe profini, muni de (générateurs)
$\rho, \sigma \in \overline{G}$ satisfaisant les relations

(1) $$ \rho^6 = \sigma^4 = 1, \quad \rho^3 = \sigma^2 \quad (\overset{dfn}{=} "-1") $$

[MARGIN: Ceci se justifie en [unclear] de la [unclear] de $\operatorname{Sl}(2, \mathbb{Z})^\wedge$]

un cas particulièrement intéressant est
celui où $\rho^3 = \sigma^2 = 1$, auquel [unclear: ramène] un
passage à la limite. Soit $\tau$ un
automorphisme de $\overline{G}$ satisfaisant

[MARGIN: Introduisant tout de suite l'involution standard $\tau$ de $\overline{G}$ (qui a $\tau(\sigma) = \sigma^{-1}$, $\tau(\rho) = \rho^{-1}$, $\tau(\varepsilon_0) = \varepsilon_0^{-1}$, $\tau(\varepsilon_1) = \varepsilon_1^{-1}$)]

(2) $\tau(\sigma) = \sigma^{-1}$, $\tau(\rho) = \sigma(\rho^{-1})$
[unclear: $\downarrow_{= -\sigma}$]
(ce qui signifie que $\tau$ "passe aux quotients"
$\tau$ de $\overline{G}$), ce qui implique

(3) $\tau^2 = 1$

(4) $\tau(\varepsilon_0) = \varepsilon_0^{-1}$, $\tau(\varepsilon_1) = \varepsilon_1^{-1}$, où

[MARGIN: de façon alléchante, les [unclear: équations] ... (5) ... $\varepsilon_0 = \sigma \rho \dots$]

(5) $\varepsilon_0 = \sigma \rho$, $\varepsilon_1 = \rho \sigma = \sigma(\varepsilon_0) = \sigma^{-1}(\varepsilon_0) = \rho(\varepsilon_0)$
$= -\sigma^{-1}\rho$

Posons
$\overline{G}' = \{1, \tau\} \cdot \overline{G}$,

donc on a

(6) $1 \longrightarrow \overline{G} \longrightarrow \overline{G}' \longrightarrow \{1, \tau\} \longrightarrow 1$

Soit

$B_{\rho, \sigma} \subset \operatorname{Aut}(\overline{G}) \times \hat{\mathbb{Z}}^*$

(7) $$ \{ (u, \mu) \mid u(\varepsilon_0) = \varepsilon_0^\mu, u(\sigma) \text{ conjugué dans } \overline{G} \text{ à } \sigma^\mu, u(\rho) \text{ [unclear: conjugué à] } \rho^\mu \} $$

qui se réduit [unclear: visiblement à un sous-groupe],
la troisième condition signifiant d'ailleurs,

529

\newpage

%% ===== Page 531 =====

puisque $\rho^6 = 1$ i.e. $\rho^\mu$ ne dépend que de l'
image de $\mu$ dans $(\mathbb{Z}/6\mathbb{Z})^* \simeq (\mathbb{Z}/3\mathbb{Z})^* \simeq \{\pm 1\}$
[unclear: crossed out text]
$$ (8) \quad \rho^\mu = \rho^{\varepsilon_3(\mu)} $$
que pour $\varepsilon_3(\mu)=1$ on a $u(\rho)$ est conjugué à
$\rho$ dans $\bar{G}$, et pour $\varepsilon_3(\mu)=-1$ il est conjugué
à $\rho^{-1}$ dans $\bar{G}$, c. a d. encore à $\rho$ par un
élément de $\bar{G}'$ i.e. un élém. de la
forme $\alpha, \tau\alpha$ avec $\alpha \in \bar{G}$ (puisque
$\tau(\rho)=\rho^{-1}$ justement...). On a donc,

[crossed out:
(9) $u(\rho) = \xi \rho \xi^{-1}$, $u(\sigma) = \eta \sigma \eta^{-1} \quad (\xi,\eta \in \bar{G})$
des relations
(10) $\xi = \alpha \rho(\alpha)^{-1}$ , $\eta = \beta$
]

On voit d'ailleurs de même que
$$ (9) \quad \sigma^\mu = \sigma^{\varepsilon_2(\mu)} \quad \begin{cases} = \sigma \text{ si } \mu \equiv 1 \text{ (4)} \\ = \sigma^{-1} = -\sigma \text{ si } \mu \equiv -1 \text{ (4)} \end{cases} $$
$$ = \varepsilon_2(\mu) \sigma $$

[crossed out: 
Mais Donc on a ...
Pour $(u,\mu) \in B_{\rho,\sigma}$ ...
Ainsi : $\alpha \in \bar{G}'$, $\beta \in \bar{G}$
]

Pour qu'un élémt $(u,\mu) \in \operatorname{Aut}(\bar{G}) \times \hat{\mathbb{Z}}^*$
soit dans $B_{\rho,\sigma}$ il f. et il s. qu'il
existe $\alpha \in \bar{G}'$, $\beta \in \bar{G}$ tels qu'on ait

\newpage

%% ===== Page 532 =====

(10) $$u(\rho) = int(\alpha)(\rho) \quad , \quad u(\sigma) = \varepsilon_2(\mu) int(\beta)(\sigma)$$

et enfin
(11) $$u(\varepsilon_0) = \varepsilon_0^\mu \quad ,$$

où de plus on a 
[MARGIN: p.e. $\alpha \in \tau$]
(12) [unclear crossed out text] $\mu \equiv 1 \quad (3)$ [unclear crossed out text]

On notera que $(u,\mu)$ est univoquement déterminé
par la connaissance de $\alpha$, de $\beta$, et de $\varepsilon_2(\mu)$
et ce même par par la
(13) $$\xi = \alpha \rho (\alpha)^{-1} \quad , \quad \eta = \varepsilon_2(\mu) \beta \sigma (\beta)^{-1}$$
puisqu'on a alors $u$ par
(14) $$u(\rho) = \xi \rho \quad , \quad u(\sigma) = \eta \sigma \quad .$$

D'ailleurs, dans le cas où
(15) $$\hat{\mathbb{Z}} \longrightarrow \widehat{\mathfrak{G}}$$
$$1 \longmapsto \varepsilon_0$$
est injectif, alors $\mu$ est déterminé [unclear] en
fonction de $u$ comme restriction de
$u$ sur $\hat{\mathbb{Z}} \subset \widehat{\mathfrak{G}}$, image de (15), donc
$(u,\mu)$ est déterminé par $u$ donc par
$(\alpha, \beta, \varepsilon_2(\mu))$. D'ailleurs, si $\widehat{\mathfrak{G}}$ est
assez gros pour s'envoyer dans $\operatorname{Sl}(2,\mathbb{Z})^\wedge$
(de façon à envoyer $\rho, \sigma$ sur les éléments
habituels ) on a que $\alpha$ : [unclear]

\newpage

%% ===== Page 533 =====

détermine déjà $\mu$, donc $\varepsilon$, donc [unclear]
cet élément de $B_{\rho,\sigma}$ est déterminé
à l'aide de $\alpha, \beta$ seuls.

Mais quelles conditions doivent satisfaire
$(\alpha, \beta, \mu)$ pour définir un autom. de $B_{\rho,\sigma}$ ?
Tout d'abord, il faut qu'il existe bien un
homom. (nécess. unique) satisfais. à
(14) — c'est une condition sur $\xi, \eta$ seuls.
Il faut ensuite exprimer (11), qui s'écrit, puisque

$$ \varepsilon_0 = \sigma p $$
$$ u(\varepsilon_0) = u(\sigma p) = [unclear] = \eta \sigma \xi p = \eta \sigma \xi \sigma^{-1} (\overbrace{\sigma p}^{= \varepsilon_0}) $$
$$ \eta (\sigma \xi \sigma^{-1}) \varepsilon_0 = \varepsilon_0^\mu \text{ i.e.} $$
$$ \eta (\sigma \xi \sigma^{-1}) = \varepsilon_0^{\mu-1} [unclear] $$

Or on a,
$$ \sigma(\eta) = [unclear] = \varepsilon \sigma(\beta) \sigma^2(\beta)^{-1} = [unclear] $$
i.e.
(15) $$ \eta \sigma(\eta) = 1 $$

(*)
$$ \varepsilon_0' = \sigma^{-1} p = - \varepsilon_0 $$
d'où
$$ (\varepsilon'_0)^\mu = (-1)^\mu \varepsilon_0^\mu = - \varepsilon_0^\mu \quad (\text{car } \mu \equiv 1 \, (2)) $$
$$ u(\varepsilon'_0) = - u(\varepsilon_0) $$ car $u(-1) = -1$ cf ci-dessous.
donc la relation (11) équivaut à :
(10) $$ u(\varepsilon'_0) = (\varepsilon'_0)^\mu $$
qui, pour la susdite valeur, s'écrit

\newpage

%% ===== Page 534 =====

$$ \tilde{y} \xi = \varepsilon_1'^{2\nu} \quad \nu = (\mu-1)/2, \varepsilon_1' = -\varepsilon_1 = \sigma \sigma^{-1} $$
$$ \varepsilon_1'^{2\nu} = (-1)^{2\nu} \varepsilon_1^{2\nu} = \varepsilon_1^{2\nu} = l_1^\nu, \text{ donc la} $$
relation cherchée prend la forme
$$ (16) \qquad \tilde{y} = y l_1^\nu \qquad \text{où } \nu = (\mu-1)/2 $$

On doit encore prouver le
Lemme. Soit un [unclear: autom.] de $\widehat{\mathbb{Q}}$
tel que $u(\sigma)$ soit conjugué de $\sigma^{-1}$ (i.e. $\sigma' = \alpha \sigma \alpha^{-1}$)
alors $u(-1) = -1$
(17) Dem. $u(-1) = u(\sigma^2) = \alpha (\sigma^{-1})^2 \alpha^{-1} = \alpha (-1) \alpha^{-1} = -1$, OK.

Reste à exprimer que [unclear: crossed out text] $u(\rho)$ est conjugué à $\rho^\mu$, si $u(\sigma) \sim \sigma^{-1}$
et cela est [unclear: crossed out text] puisque
$u(\sigma \tau) = \dots$ (voir 19))
et pour $\rho$ il signifie que l'on a
(17) $\alpha \equiv \tau^{\nu_3(\mu)} \dots$
où $\nu_3(\mu) \in \hat{\mathbb{Z}}$ [unclear: est l'expos.] de $\varepsilon_3(\mu) \dots$
qu'on écrit aussi $\tau^{\nu_3(\mu)} \alpha \in \dots$

Dem. $\dots$

Corollaire L'application
$$ (18) \qquad B_{\hat{\rho}, \sigma} \longrightarrow \mathcal{C} \times \mathcal{C} \times \hat{\mathbb{Z}}^\ast $$
$$ (u, \mu) \longmapsto ([u(\rho)], [u(\sigma)], \mu) $$
est injective, et elle a comme image l'ens.
des triples $([\rho'], [\sigma'], \mu)$ satisfaisant les
conditions suivantes :
1°) la relation $\dots$
2°) $\dots$

\newpage

%% ===== Page 535 =====

3°) On a la relation de tresse (16) .

4°) $\exists$ autom. $u$ de $\widehat{\mathfrak{S}}$, satisfaisant $u(\rho)=\xi \rho$, $u(\sigma)=\eta \sigma$.

[crossed out: Pour ceci, il faut &]

Les conditions 1'), 2'), 3') sont équivalentes
à celle-ci :
[crossed out: elle est équivalente]

4') $\exists (\alpha, \mu) \in \widehat{\mathfrak{S}} \times \hat{\mathbb{Z}}^*$, avec $\mu \equiv 1 (3)$
i.e. $\alpha \bmod \widehat{\mathfrak{S}}' = \tau^{\nu_3(\mu)}$ , et tels qu'on ait
$\xi = \alpha \rho(\alpha)^{-1}$, $\eta = \varepsilon_2(\mu)$ .

Considérons l'hom. de groupes
(19) $$L_0 \xrightarrow{i_{L_0,B}} B_{\rho,\sigma}$$
$$\ell \longmapsto (int(\ell), 1)$$
il est injectif si on suppose tjrs
(20) $$Centr(\widehat{\mathfrak{S}}) = \{1\}$$ ,
et son image est invariante on a
(21) $$u(i_{L_0,B}(\ell)) = i_{L_0,B}(u(\ell)) = i_{L_0,B}(\ell^\mu)$$ ,
où $\mu$ est le multiplicateur. On pose
(22) $$\Gamma_{\rho,\sigma} = B_{\rho,\sigma} / Im \, L_0$$ ,
on construit donc une suite exacte canonique
(23) $$1 \to L_0 \to B_{\rho,\sigma} \to \Gamma_{\rho,\sigma} \to 1$$
et on a un hom. (déduit de
$(u,\mu) \mapsto \mu$ factorisé par $\Gamma_{\rho,\sigma}$ )
(24) $$\Gamma_{\rho,\sigma} \xrightarrow{p_{\Gamma}} \hat{\mathbb{Z}}^*$$
d'où la [unclear] suite exacte :
(25) $$1 \to \Gamma_{\rho,\sigma}^+ \to \Gamma_{\rho,\sigma} \xrightarrow{p_{\Gamma}} \hat{\mathbb{Z}}^* \to 1$$

\newpage

%% ===== Page 536 =====

571

Considérons les hom. composés 
(26)
$$ \gamma_{\rho,\sigma} \xrightarrow{\chi_\gamma} \hat{\mathbb{Z}}^\times \xrightarrow{\varepsilon_2} \pm 1 $$
$$ \gamma_{\rho,\sigma} \xrightarrow{\chi_\gamma} \hat{\mathbb{Z}}^\times \xrightarrow{\varepsilon_3} \pm 1 $$
et soient 
(27)
$$ \gamma_{\rho,\sigma}', \gamma_{\rho,\sigma}'' \subset \gamma_{\rho,\sigma} \text{ sous-groupes d'indice 2, définis par} $$
$$ \gamma_{\rho,\sigma}' = \operatorname{Ker}(\varepsilon_3 \circ \chi_\gamma), \quad \gamma_{\rho,\sigma}'' = \operatorname{Ker}(\varepsilon_2 \circ \chi_\gamma) $$

Soit maintenant
(28)
$$ G_{\rho,\sigma} = (B_{\rho,\sigma} \cdot \widehat{\mathfrak{S}}) / L_0 $$
(1/2 direct)
$$ \simeq S_0 G_{\rho,\sigma} $$
où 
(29)
$$ L_0 \xrightarrow{i_{L_0, B}} B_{\rho,\sigma} \cdot \widehat{\mathfrak{S}} $$
$$ \ell \mapsto (i(\ell), \ell^{-1}) $$
dont l'image est invariante. On a un 
diagramme commutatif 
(30)
\begin{tikzcd}
L_0 \ar[r, hook] \ar[dr] & \widehat{\mathfrak{S}} \ar[dr] \\
& B_{\rho,\sigma} \ar[r] & G_{\rho,\sigma}
\end{tikzcd}
et 
$$ \operatorname{Ker}(\widehat{\mathfrak{S}} \to G_{\rho,\sigma}) \simeq \operatorname{Ker}(L_0 \to B_{\rho,\sigma}) $$
$$ \simeq L_0 \cap \text{Centre}(\widehat{\mathfrak{S}}) \text{ est trivial vu } \dots (20), $$
ce qui [unclear] justifie les 
suites exactes 
(31)
\begin{tikzcd}
1 \ar[r] & L_0 \ar[r] \ar[d, hook] & B_{\rho,\sigma} \ar[r, "\delta_B"] \ar[d] & \gamma_{\rho,\sigma} \ar[r] \ar[d, equal] & 1 \\
1 \ar[r] & \widehat{\mathfrak{S}} \ar[r] & G_{\rho,\sigma} \ar[r, "\delta_G"'] & \gamma_{\rho,\sigma} 
\end{tikzcd}
$$ [ \xrightarrow{\chi_G} \hat{\mathbb{Z}}^\times ] $$

\newpage

%% ===== Page 537 =====

Posons
$$ (32) \qquad \left\{ \begin{matrix} Z_\rho = Cent_{G_{\rho,\sigma}}(\rho) \\ Z_\sigma = Cent_{G_{\rho,\sigma}}(\sigma) \end{matrix} \right. $$
alors les hom. induits par $\delta_G$
$$ (33) \qquad Z_\rho \xrightarrow{\delta_\rho} \gamma_{\rho,\sigma} \, , \quad Z_\sigma \xrightarrow{\delta_\sigma} \gamma_{\rho,\sigma} $$
ont des images qui contiennent
$\gamma_{\rho,\sigma}'$ et $\gamma_{\rho,\sigma}''$ respectivement, et [unclear]

[unclear: D'ailleurs l'image de $Z_\rho$ est contenue dans $\gamma_{\rho,\sigma}'$]

$\rho$ et $\rho^{-1}$ [unclear: en] conjugués dans $\widehat{G}$, [unclear: il suffit]
L'image de $Z_\sigma$ est
n'est autre : $\gamma_{\rho,\sigma}''$ car $\sigma$ et $\sigma^{-1}=-\sigma$
[unclear: sont] conjugués dans $\widehat{G}$. Notons p.ex.
pour $\widehat{G} = \widehat{Sl(2,\mathbb{Z})}$, $\sigma^2 = 1$,
la dernière condition est par
suite : $Z_\sigma \to \gamma_{\rho,\sigma}$ [unclear: surjective],
visiblement elle [unclear: est fausse]
[unclear] $\sigma \mapsto -\sigma$ par la
[unclear: symétrie de] Steinberg p.ex.) -
le [unclear: problème] se pose-t-il pareillement dans les
$\widehat{Sl(2,\mathbb{Z})}^\sim = \widehat{\mathfrak{S}}_{1,1}$ ?

[MARGIN: La constatation p.ex. que pour $\widehat{G} = \widehat{Sl(2,\mathbb{Z})}/(\pm 1)$, on a [unclear] s'ensuit de [unclear]. Celui-ci donne $Z_\rho \simeq \widehat{Sl(2,\mathbb{Z})}$ (d'où $Z_\rho \to \gamma_{\rho,\sigma}'$ est surjective, $\gamma_{\rho,\sigma}' \simeq \operatorname{Sl}(2,\hat{\mathbb{Z}})$) et $Z_\sigma$ s'envoye de même sur $P\operatorname{Gl}(2,\hat{\mathbb{Z}})$. L'image de $Z_\sigma$ dans $\gamma_{\rho,\sigma}''$ est d'indice 2. Mais d'où vient la [unclear] ?]

\newpage

%% ===== Page 538 =====

572

Proposition 1°) L'image $\delta_{\rho}(Z_{\rho})$ contient $\gamma_{\rho, \sigma}'$,
et est réduite à $\gamma_{\rho, \sigma}'$ ssi $\rho$ et $\rho^{-1}$ non conju-
gués dans $\overline{\mathbb{G}}$ - donc elle est égale
à $\gamma_{\rho, \sigma}$ i.e. $\delta_{\rho}: Z_{\rho} \to \gamma_{\rho, \sigma}$ épi ssi $\rho$ et $\rho^{-1}$ sont conjugués
dans $\overline{\mathbb{G}}$.

2°) L'image $\delta_{\sigma}(Z_{\sigma})$ contient $\gamma_{\rho, \sigma}''$, et elle
est réduite à $\gamma_{\rho, \sigma}''$ ssi $\sigma$ et $\sigma'=-\sigma$ non
conjugués dans $\overline{\mathbb{G}}$ - donc elle est égale
à $\gamma_{\rho, \sigma}$ i.e. $\delta_{\sigma}: Z_{\sigma} \to \gamma_{\rho, \sigma}$ épi ssi
$\sigma$ et $\sigma'=-\sigma$ sont conjugués dans $\overline{\mathbb{G}}$.

Corollaire! Si $\overline{\mathbb{G}}$ contient comme quotient
$\operatorname{Sl}(2, \mathbb{Z}/3\mathbb{Z})'$ avec $\rho = \begin{pmatrix} 0 & 1 \\ -1 & 1 \end{pmatrix}$, $\sigma = \begin{pmatrix} 0 & -1 \\ 1 & 0 \end{pmatrix}$
(mod $\pm 1$), alors l'image de $\delta_{\rho}$ est $\gamma_{\rho, \sigma}'$.

En effet, on voit que $\rho$ et $\rho^{-1}$ ne sont pas
conjugués dans $\overline{\mathbb{G}}$, car ils ne le sont
pas dans $\operatorname{Sl}(2, \mathbb{Z}/3\mathbb{Z})'$ - vérifions le :
[unclear: crossed out text] $\rho$ et $-\rho^{-1}$ ont des traces distinctes
$1$ et $-1$, dans $\mathbb{F}_3$, ils ne sont pas conjugués.
On a $(\tau_{\infty})(\rho) = \rho^{-1}$, donc les éléments
de $\operatorname{Gl}(2, \mathbb{Z}/3\mathbb{Z})$ qui conjuguent
$\rho$ en $\rho^{-1}$ sont de la forme $f \tau_{\infty}$, où
$f$ [unclear] $\det(f) = 1$ i.e. $\det(f \tau_{\infty}) = -1$
(car $\det(\tau_{\infty}) = -1$) i.e. $f \tau_{\infty} \notin Sl$, impossible.

\newpage

%% ===== Page 539 =====

Corollaire 2 Si $\sigma^2=1$, alors $Z_\sigma \to \gamma_{p,\sigma}$ est surjectif. Si par contre $\underline{G}$ contient comme quotient [unclear: rayé] $\operatorname{Sl}(2, \mathbb{Z}/4\mathbb{Z})$ (sans $\pm 1!$), avec $\rho,\sigma$ comme d'habitude alors $\operatorname{Im}(\delta_\sigma) = \gamma^{''o}_{p,\sigma}$.

Il faut voir que $\sigma$ et $\sigma^{-1}=-\sigma$ ne sont conjugués dans $\operatorname{Sl}(2, \mathbb{Z}/4\mathbb{Z})$, en utilisant que $\tau_\infty(\sigma)=\sigma^{-1}$, donc [unclear: it] $f$ les élts de $\operatorname{Gl}(2, \mathbb{Z}/4\mathbb{Z})$ qui conjuguent $\sigma$ en $\sigma^{-1}$ 
[MARGIN: $g=f\tau_\infty$] 
sont les $f = \tau_\infty m$, et dire que $det(g)=1$ signifie $det f = -1$ i.e. $m^2=-1$, impossible.

Posons [unclear: rayé]
(34) $\left\{ \begin{array}{l} S Z_\rho = Z_\rho \cap \underline{G}^\circ = Cent_{\underline{G}^\circ}(\rho) \\ S Z_\sigma = Z_\sigma \cap \underline{G}^\circ = Cent_{\underline{G}^\circ}(\sigma) \\ S B_{\rho,\sigma} = \ker \delta_B = B_{\rho,\sigma} \cap \underline{G}^\circ = L_0 \end{array} \right.$

[unclear: texte rayé]
(35) $Soit M_{\rho,\sigma} \subset B_{\rho,\sigma} \times Z_\rho \times Z_\sigma \times G_{\rho,\sigma}$
$= \{ (u,a,b,U) | \delta_p(u) = \delta_G(U) = \delta_{\rho,\sigma}(b) = \epsilon_3(u) \delta_p(a) \}$

\newpage

%% ===== Page 540 =====

539

Notons que d'après (2), (4) on a
(35) $\tau_B \overset{dfn}{=} (\tau, -1) \in B_{p, \sigma}$
Et évidemment $\tau_B \notin L_0 = B_{p, \sigma} \cap \overline{\mathfrak{S}}$. Donc
on a un plongement

[MARGIN: NB $\overline{\mathfrak{S}}'$ est un quotient de $\operatorname{Gl}(2, \mathbb{Z})^\wedge$, et l'y [unclear: relevé] de $\tau$ est l'élément $\tau_\infty = \begin{pmatrix} -1 & 0 \\ 0 & 1 \end{pmatrix}$ de $\operatorname{Gl}(2, \mathbb{Z})^\wedge$]

(36) $\overline{\mathfrak{S}}' = \{1, \tau\} \cdot \overline{\mathfrak{S}} \hookrightarrow G_{p, \sigma}$
$\tau \mapsto \tau_B$
$g \mapsto g \quad (g \in \overline{\mathfrak{S}})$

donc l'automorphisme $\tau$ de $\overline{\mathfrak{S}}$ sera considéré comme
induit par $Int(\tau_B)$, i.e. $\tau_B$ se considére
lui-même comme un élément de $G_{p, \sigma}$.
De l'égalité
$\tau_B(\sigma) = \sigma^{-1} \ (= -\sigma)$
qui montre que $\tau_B$ normalise le sous-groupe $L_0$ engendré par $\sigma$
([unclear: isomorphe] à un quotient de $\hat{\mathbb{Z}} / 4 \hat{\mathbb{Z}}$), 
notant $Z_\sigma = Cent_{G_{p, \sigma}}(L_0)$, introduisons

(37) $N_\sigma = \{1, \tau_B\} \cdot Z_\sigma \quad \text{ (½ direct)}$

et le hom. canonique

(38) $N_\sigma \xrightarrow{i_\sigma} G_{p, \sigma}$
$\tau_B \mapsto \tau_B, \ x \mapsto x \ (x \in Z_\sigma)$

L'image de $i_\sigma$ est le normalisateur de $\sigma$ dans $G_{p, \sigma}$, car un [unclear: élément]
$g \in G_{p, \sigma}$ qui normalise $L_0$ doit transformer
$\sigma$ en un générateur de $L_0$, i.e. en $\sigma$ ou $\sigma^{-1}$;
dans le premier cas $g \in Z_\sigma$, dans
le deuxième $g(\sigma) = \tau_B(\sigma)$ i.e. $\tau_B^{-1} g \in Z_\sigma$
i.e. $g \in \tau_B Z_\sigma$. L'hom. (38) est injectif.

\newpage

%% ===== Page 541 =====

[crossed out: $ssi$] $\tau_B \notin Z_\sigma$ i.e. $\sigma \neq \sigma^{-1}$ i.e. $\sigma^2 \neq 1$
i.e. $"-1" \neq 1$.

Puisqu'on y est, profitons de la formule
$$ \tau_B(\rho) = \sigma(\rho^{-1}) $$
pour introduire
(39) $$ \tau' \overset{df}{=} \sigma^{-1} \tau \in \tilde{\mathfrak{S}}' \subset G_{\rho, \sigma} $$
(qui correspond à l'élément
(40) $$ \tau'_\infty = \sigma_\infty^{-1} \tau_\infty = - \sigma_\infty \tau_\infty = \begin{pmatrix} 0 & 1 \\ 1 & 0 \end{pmatrix} \in \operatorname{Gl}(2, \mathbb{Z}) $$ )
[MARGIN: $\tilde{\mathfrak{S}}_{1,1}$]

[MARGIN: NB On a [unclear] donc $\tau'$ [unclear] on peut considérer $\tau'$ comme él. d'o. 2 de $N_\rho$, et $N_\rho = \{1, \tau'\} . Z_\rho$ (produit s. direct)]
On aura donc
(40') $$ \tau'^2 = 1, \quad \tau'(\rho) = \rho^{-1} $$
Ainsi $\tau'$ normalise le sous-groupe $L_\rho$ (donc
on peut [crossed out: s'en servir] [crossed out: engendré par] $\rho$, donc il normalise le centralisateur de
$L_\rho$, et on peut construire
(41) $$ N_\rho \overset{df}{=} \{1, \tau'\} . Z_\rho $$
(semi-direct)
et on a des hom. canoniques
(42) $$ \begin{cases} N_\rho \xrightarrow{i_\rho} G_{\rho, \sigma} \\ \tau' \mapsto \tau', \ x \mapsto x \text{ pour } x \in Z_\rho \end{cases} $$
L'image en est [crossed out: $sss$] le normalisateur du
s.g. $L_\rho$ dans $G_{\rho, \sigma}$ - on le voit comme tout à l'heure pour $i_\sigma$,
vu que $\rho, \rho^{-1}$ sont les seuls générateurs de $L_\rho$. L'hom.
$i_\rho$ est injectif ssi $\tau' \notin Z_\rho$ i.e. $\rho \neq \rho^{-1}$
i.e. $\rho^2 \neq 1$ i.e. $\rho \neq -1$ [crossed out: implique] (NB $\rho = -1$ 
équivaut à ce que $\rho$ commute à $\sigma$, et que $\tilde{\mathfrak{S}}_{1,1}$
soit un quotient de $\mathbb{Z}/2\mathbb{Z}$, ce qui permet d'avoir $\rho=1$
et $\sigma=3$, mais si $2\rho=0$ dans $\hat{\mathbb{Z}}$ on a bien $\rho=-1$
dans $\mathbb{Z}/4\mathbb{Z}$ avec $\rho$ correspondant à $2$, $\sigma$ à $3=-1$ ... ).
* On a fait ce qu'il fallait pour éviter des épimorphismes.

[MARGIN: 
(43) $$ \begin{cases} N_\rho \xrightarrow{j_\rho = \delta \circ i_\rho} \bar{G}_{\rho, \sigma} \\ N_\sigma \xrightarrow{j_\sigma = \delta \circ i_\sigma} \bar{G}_{\rho, \sigma} \end{cases} $$
épimorphismes
]

\newpage

%% ===== Page 542 =====

[MARGIN: 
(46) 
$Sg_{\rho,\sigma}^{\#} \xrightarrow{\text{épi.}} \Sigma_{\rho,\sigma}^{\#}$]

(44)
$$Sg_{\rho,\sigma}^{\#} \subset \hat{C}^* \times \hat{C}^* \times \hat{Z}^*$$
$$= \left\{ (\alpha, \beta, \mu) \mid \alpha \equiv \tau^{\nu_3(\mu)} \mod \hat{\mathbb{Q}} \text{ i.e. } \chi(\alpha) = \varepsilon_3(\mu) , \text{ et } \alpha \rho(\alpha)^{-1} = \varepsilon_2(\mu) \beta \sigma(\beta)^{-1} l_1^\nu, \text{ où } \nu = (\mu-1)/2 \text{ (} l_1 = \varepsilon_1^2 \text{)} \right\}$$

d'où
(45)
$$Sg_{\rho,\sigma}^{\#} \longrightarrow \Sigma_{\rho,\sigma} \subset \hat{C} \times \hat{C} \times \hat{Z}^*$$
$$(\alpha, \beta, \mu) \longmapsto ( \alpha \rho(\alpha)^{-1} \varepsilon_{\text{unclear}}(\mu) \rho, \beta \sigma(\beta)^{-1}, \mu )$$

On désigne par $G_{\rho,\sigma}^{\#}$ l'image inverse de
$\Sigma_{\rho,\sigma}$ (cf [unclear: var.] p. 570) $= B_{\rho,\sigma}$. Par constructi-
on on a une application surjective continue
On considère $Sg_{\rho,\sigma}^{\#}$ comme contenu
dans un groupe plus gros

(47)
$$G_{\rho,\sigma}^{\#} \subset G_{\rho,\sigma} \times G_{\rho,\sigma} \times \hat{Z}^*$$
$$= \left\{ (\alpha, \beta, \mu) \mid \chi(\alpha) = \varepsilon_3(\mu) \chi(\beta) ; \alpha \rho(\alpha)^{-1} = \varepsilon_2(\mu) \beta \sigma(\beta)^{-1} l_1^\nu \right\}$$

de sorte que l'on aura
(48)
$$Sg_{\rho,\sigma}^{\#} = G_{\rho,\sigma}^{\#} \cap ( \hat{C}^* \times \hat{C}^* \times \hat{Z}^* )$$
$$= \left\{ (\alpha, \beta, \mu) \in G_{\rho,\sigma}^{\#} \mid \delta(\beta) = 1, \delta_p(\alpha \tau^{\nu_3(\mu)}) = 1 \right\}$$

\newpage

%% ===== Page 543 =====

et on a encore une application canonique

(49)
$$ g_{\rho,\sigma}^\# \longrightarrow \Sigma_{\rho,\sigma} $$
$$ (\alpha, \beta, \mu) \longmapsto ([\alpha, \rho], \varepsilon_2(\mu)[\beta, \sigma], \mu) \ . $$

On note $g_{\rho,\sigma}^{\sharp\sharp}$ l'image inverse de $\Sigma_{\rho,\sigma}^{eff}$ , et on trouve

(50)
$$ Sg_{\rho,\sigma}^{\sharp\sharp} = g_{\rho,\sigma}^{\sharp\sharp} \cap (\underline{\mathfrak{S}}' \times \underline{\mathfrak{S}} \times \hat{\mathbb{Z}}^*) $$
$$ = \left\{ (\alpha, \beta, \mu) \in g_{\rho,\sigma}^{\sharp\sharp} \mid \begin{matrix} \delta_\sigma(\beta) = 1 \\ \delta_\rho(\alpha \tau^{\nu_3(\mu)}) = 1 \end{matrix} \right\} $$

On a une section de

(51)
\begin{tikzcd}
g_{\rho,\sigma}^{\sharp\sharp} \ar[r] & \Sigma_{\rho,\sigma}^{eff} \\
& B_{\rho,\sigma} \ar[u, "\simeq"'] \ar[ul, "section"]
\end{tikzcd}

[unclear: marginal text to the right of diagram]

donnée par

(52)
$$ u \longmapsto (\varepsilon u, u \tau^{\nu_2(\mu)}, \mu = \chi_D(u)) $$

puisqu'on a (posant $\varepsilon = \varepsilon_2(\mu)$ [unclear])

$$ [u,\rho] = u \rho u^{-1} \rho^{-1} = \xi $$
$$ \varepsilon [u\tau^{\nu_2(\mu)}, \sigma] = \varepsilon u \tau^{\nu_2} \sigma (u\tau^{\nu_2})^{-1} \sigma^{-1} = \varepsilon u (\tau^{\nu_2} \sigma \tau^{-\nu_2} \sigma^{-1}) \sigma(u^{-1}) \sigma^{-1} = \varepsilon u \sigma(u^{-1}) = \eta $$

OK. Donc à $(\alpha, \beta, \mu)$ correspond dans $\Sigma_{\rho,\sigma}^{eff}$

$$ (y_1, y_2, \mu) = ([u,\rho], [u,\sigma], \mu) $$ exactement [unclear]

\newpage

%% ===== Page 544 =====

575

[crossed out: les] $G_{\rho, \sigma}(u, \mu) =$

(53) $\left\{ (\overbrace{u a_0^{-1}}^\alpha, \overbrace{u \tau'^{1/2} b_0^{-1}}^\beta, \mu) \mid a_0 \in Z_\rho, b_0 \in Z_\sigma, \nu_\tau = \nu_\tau(\mu) \right\}$

$(= (u a_0^{-1}, u b^{-1}, \mu))$ avec $b = b_0 \tau'^{-1/2}$

qu'on écrit aussi

(54) $\left\{ (\overbrace{u a_0^{-1}}^\alpha, \overbrace{u i_\sigma(b)}^\beta, \mu) \mid a_0 \in Z_\rho, b \in N_\sigma, \underbrace{\varepsilon_\sigma(b) = \varepsilon_\tau(\mu)}_{(\text{cf } (37)(38))} \right\}$

où $\varepsilon_\sigma : N_\sigma \to \{\pm 1\}$ est défini sur (37) comme
l'homom. canon. d'image d'un $\mathbb{Z}$ et rentre dans la suite ex.

(55) $1 \to Z_\sigma \to N_\sigma \xrightarrow{\varepsilon_\sigma} \underset{\simeq \{1, \tau'\}}{\{\pm 1\}} \to 1$

Si on cherche l'un des éléments de
$S G_{\rho, \sigma}^{\text{eff}}$ - donné en $(\gamma, \gamma, \mu)$, il faut de
plus imposer les conditions

(56) $\alpha = u a_0^{-1} \tau'^{\nu_\tau(\mu)} \in \tilde{C}$, $u i_\sigma(b) \in \tilde{C}$

i.e.
$\delta_{\tilde{C}}(u a_0^{-1} \tau'^{\nu_\tau(\mu)}) = 1$, $\delta_{\tilde{C}}(u i_\sigma(b)) = 1$

soit encore

(57) $\delta_{\tilde{C}}(u) = \underbrace{\delta_{\tilde{C}}(b)}_{\text{cf } (43)} = \underbrace{\delta_{\tilde{C}}(\tau'^{1/2} a_0)}_{\text{cf } (43)}$

Posons donc
$a = \tau'^{1/2} a_0$ dans $N_\rho$ (donc $a_0 = \tau'^{-1/2} a$)

on voit que

(58) $G_{\rho, \sigma}(u, \mu) = \left\{ (u a^{-1} \tau'^{1/2}, u b^{-1}, \mu) \mid \begin{array}{l} a \in N_\rho, b \in N_\sigma \text{ tels que } \\ \varepsilon_\rho(a) = \varepsilon_\tau(\mu), \varepsilon_\sigma(b) = \varepsilon_\tau(\mu) \\ \delta_\rho(a) = \delta_\sigma(b) = \delta_{\tilde{C}}(u) \end{array} \right\}$

Introduisons donc le groupe

(59) [crossed out: $M_{\rho, \sigma} = B_{\rho, \sigma} \times$] $M_{\rho, \sigma} = N_\rho \times_{\mathcal{G}_{\rho, \sigma}} N_\sigma$

et les homomorphismes composés

(60)
[crossed out: $M_{\rho, \sigma}$] $M_{\rho, \sigma} \to N_\rho \xrightarrow{\varepsilon_\rho} \{\pm 1\}$
[crossed out: $M_{\rho, \sigma}$] $M_{\rho, \sigma} \to N_\sigma \xrightarrow{\varepsilon_\sigma} \{\pm 1\}$
encore notés $\varepsilon_\rho, \varepsilon_\sigma$ (p. abus de notations)

\newpage

%% ===== Page 545 =====

et
(57) $$M_{\rho, [\sigma]} \xrightarrow{\delta_M} \Gamma_{\rho, \sigma} \xrightarrow{\chi_\gamma} \hat{\mathbb{Z}}^* \quad \text{ soit } \chi_{S\mathcal{M}} \text{ (ou } \chi_{M_{\rho, \sigma}})$$

considérons
(58) $$M_{\rho, [\sigma]} \xrightarrow{\varepsilon_{\rho, \sigma} (= \varepsilon_{\rho}, \varepsilon_{\sigma})} \{ \pm 1 \} \times \{ \pm 1 \}$$
et
$$(\varepsilon_{2, 3}: x \mapsto (\varepsilon_2(\chi(x)), \varepsilon_3(\chi(x)))$$
(59) $$M_{\rho, [\sigma]} \longrightarrow \{ \pm 1 \} \times \{ \pm 1 \}$$

et considérons le sous-groupe des coïncidences des deux hom. précédents
(60) $$M_{\rho, [\sigma]}^{eff} = \{ x \in M_{\rho, [\sigma]} \mid \varepsilon_{\rho, \sigma}(x) = \varepsilon_{2, 3}(x) \}$$
$$= \{ (u, a, b) \mid \delta_{\rho}(a) = \delta_{\sigma}(b) \text{ soit } \Delta \in \Gamma_{\rho, \sigma} \text{; } \varepsilon_{\rho}(a) = \varepsilon_2(\chi_\gamma(\Delta)), \varepsilon_{\sigma}(b) = \varepsilon_3(\chi_\gamma(\Delta)) \}$$

On a une application naturelle
(61) $$M_{\rho, [\sigma]}^{eff} \longrightarrow G' \times G \times \hat{\mathbb{Z}}^*$$
$$(u, a, b) \longmapsto (\alpha, \beta, \mu), \text{ où } \alpha = u a^{-1} \tau'^{-1} \mathfrak{p}_{\rho}(a)$$
$$\beta = u b^{-1}$$

[MARGIN: NB on a posé $\mathfrak{p}_{\rho}(a) \in \hat{\mathbb{Z}} / 2$ tel que $(-1)^{\mathfrak{p}_{\rho}(a)} = \varepsilon_{\rho}(a)$ et]
[MARGIN: NB on identifie $\Gamma_{\rho} = \Gamma_{\sigma}$ t.q. pour $a \in N_{\rho}, b \in N_{\sigma}$ on ait $\delta_{\rho}(a) = \delta_{\sigma}(b) = \Delta \in \Gamma_{\rho, \sigma}$]

Le fait ce qu'il fallait prouver pour que cell-ci soit d'accord [unclear: il y] ait un isomorphisme
$S\mathcal{G}_{\rho, \sigma}^{eff}$, donne en dessous une bijection
(62) $$M_{\rho, [\sigma]}^{eff} \xrightarrow{\sim} S\mathcal{G}_{\rho, \sigma}^{eff}$$
les formules pour [unclear]

En fait (61) définit une application
(63) $$M_{\rho, [\sigma]} \longrightarrow G' \times G \times \hat{\mathbb{Z}}^*$$
et pour un $(u, a, b) \in M_{\rho, [\sigma]}$, son image $(\alpha, \beta, \mu)$
est lié à $u \in \Gamma_{\rho, \sigma}$ par les formules

\newpage

%% ===== Page 546 =====

(63)
$$u(\rho) = int(\alpha)(\rho) \quad \text{ (à-un-signe-près : } \rho^{\varepsilon'}, \text{ où } \varepsilon' = \varepsilon_{\rho}(\alpha))$$
$$u(\sigma) = \varepsilon int(\beta)(\sigma) \quad \text{ (à-un-signe-près : } \sigma^{\varepsilon}, \text{ où } \varepsilon = \varepsilon_{\sigma}(\beta))$$

(mais sans qu'on ait nécessairement $\varepsilon' = \varepsilon_3(\mu), \varepsilon = \varepsilon_2(\mu)$).

Ceci suggère une définition d'un groupe un
peu plus gros que $B_{\rho,\sigma}$, dans lequel $B_{\rho,\sigma}$ sera
comme sous-groupe d'indice $1, 2 \dots 4$ — et qui
aura l'avantage, dans les "bons cas", d'être le
groupe des pts rationnels d'un schéma en
groupes rassemble sur $\mathbb{Z}$. D'autre part,
pour [unclear: éliminer] le signe $\varepsilon$ dans (63),
(qu'on présume $=$ [unclear] $(u,\alpha,\beta) \in \mathcal{M}_{\rho,\sigma}[0]^{eff}$ ),
j'ai envie de remplacer $\beta$ par $\beta' = \tau^\nu \beta$, où
$\nu \in \mathbb{Z}/2$ tel que
(64)
$$\beta' = \tau^\nu \beta \qquad (\text{où } \nu \in \mathbb{Z}/2 \text{ tel que } (-1)^\nu = \varepsilon)$$

De sorte qu'on aura
(65)
$$u(\sigma) = int(\beta')(\sigma) \quad \text{i.e.} \quad [u, \sigma] = \beta' \sigma (\beta')^{-1} \ .$$

Ceci est lié à la convention qui me
dit [unclear], que nous repérions par un
couple $(\alpha, \beta) \in \pi_{0,3}^{\wedge \ \prime} \times \pi_{0,3}^{\wedge \ \prime}$ , non disons mais
opéré par un couple $(\alpha, \beta') \in \pi_{0,3}^{\wedge \ \prime} \times \pi_{0,3}^{\wedge \ \prime}$ , où
(66)
$$\chi(\alpha) = \varepsilon_3(\mu) , \chi(\beta') = \varepsilon_2(\mu) \quad (\text{où } \mu = \chi(u) \text{ dans le multiplicateur})$$

ce qui donne pour l'opération de $u$ sur $\widehat{\mathcal{B}}_{0,3}^{\wedge}$
$$ \widehat{\mathcal{B}}_{0,3}^{\wedge +} = \operatorname{Sl}(2,\mathbb{Z})^\wedge = \widehat{SL}(2,\mathbb{Z}) / \pm 1 $$
un [unclear] par en-taille les formules partielles :
(67)
$$
\begin{cases}
u(\rho) = int(\alpha).\rho , \ u(\sigma) = int(\beta')\sigma \quad (= [unclear]) \\
u(\rho \sigma) = [unclear] \quad (= \beta' \sigma (\beta')^{-1} \rho^\nu) \quad \nu = \frac{\mu-1}{2}
\end{cases}
$$

\newpage

%% ===== Page 547 =====

mais pour l'opération dans [crossed out]
$$ \mathcal{C}_{1,1}^{+\wedge} \simeq \operatorname{Sl}(2,\Z)^\wedge $$ lui-même
donne des formules plus simples, à savoir , au
lieu de
(68) $$ \begin{cases} u(\rho) = \varepsilon. int(\alpha) \rho \quad u(\sigma) = \varepsilon. int(\beta) \sigma \\ \alpha \rho c \alpha \tau^{-1} = \varepsilon \beta \sigma (\beta)^{-1} h_1^\nu \end{cases} $$ )
les formules plus symétriques
(69) $$ \begin{cases} u(\rho) = int(\alpha) \rho \quad , \quad u(\sigma) = int(\beta) \sigma \\ \alpha \rho c \alpha \tau^{-1} = \beta \sigma (\beta)^{-1} \sigma^{-1} h_1^\nu \end{cases} $$ ,
où le signe $\varepsilon = \varepsilon_1(\mu)$ disparait. Il
apparaît dans les formules
(70) $$ \chi(\alpha) = \varepsilon_3(\mu), \chi(\beta) = \varepsilon_2(\mu) $$.

Il me semble en avoir vu assez pour
pouvoir reprendre "à main levée" les
constructions précédentes et les achever, avec
des notations plus adéquates.
Le groupe $G'$ [MARGIN: quotient de $\mathcal{C}_{1,1}^\wedge \simeq \operatorname{Sl}(2,\Z)^\wedge$] perd d'ailleurs une
importance croissante dans nos calculs, c'est
lui qui mérite vrai la notation sigle
$\tilde{G}$, le sous-groupe fermé engendré par $\rho, \sigma$
devenant $\tilde{G}^+$ - il [unclear: s'envoie] : $\mathcal{C}_{1,1}^{+\wedge} \simeq \operatorname{Sl}(2,\Z)^{+\wedge}$.
On s'aperçoit d'ailleurs que $\tau \notin \tilde{G}^+$ (c. à d. l'image
de $\tau \infty$), mais qu'il [unclear] $=$, ou
(71) $$ \tau' \notin \tilde{G}^+ $$

\newpage

%% ===== Page 548 =====

C'est [unclear] et [unclear], qui conduit
à la fois [unclear] et [unclear] les formules
remarquables

(*) $\tau'(\rho) = \rho^{-1}$, $\tau'(\sigma) = \sigma^{-1}$

qui prend le pas sur $\tau$, qui satisfait à

(40) $\tau(\varepsilon_0) = \varepsilon_0^{-1}$, $\tau(\varepsilon_1) = \varepsilon_1^{-1}$ et même $\tau(\sigma) = \sigma^{-1}$

et est évidemment très [unclear: astucieux] ! ),
dont l'action sur $\rho$
$\tau(\rho) = \sigma(\rho^{-1})$
est incommode.

[MARGIN: Néanmoins $\tau'$ a sur $\tau$ le désavantage sérieux de ne pas se [unclear] dans le groupe $\hat{\Gamma}_{1,1}^+ \simeq \operatorname{Sl}(2, \mathbb{Z})^\wedge$, la relation $\tau(\rho) = \sigma(\rho^{-1})$ contredit $\rho^3 = \sigma^2$ ! Alors que $\tau'$ ne contredit rien, mais $\tau'$ a ce désavantage. De plus, $\tau' \notin$ [unclear]]

10) Soit donc $\mathcal{G}$ un groupe profini, muni d'
un hom. surjectif

(1) $\hat{\Gamma}_{1,1} \simeq \operatorname{Gl}(2, \mathbb{Z})^\wedge \overset{épi}{\longrightarrow} \mathcal{G}$

ce que nous exprimerons par la donnée
de trois générateurs topologiques $\rho, \sigma, \tau' \in \mathcal{G}$,
satisfaisant à

(2) $\rho^6 = \sigma^4 = \tau'^2 = 1$, $\rho^3 = \sigma^2$ (noté $-1$), $\tau' \rho \tau'^{-1} = \rho^{-1}$, $\tau' \sigma \tau'^{-1} = \sigma^{-1}$.

On désigne par $\mathcal{G}^+$ le sous-groupe fermé engen-
dré par $\rho, \sigma$, qui est donc l'image de

$\hat{\Gamma}_{1,1}^+ \simeq \operatorname{Sl}(2, \mathbb{Z})^\wedge$
)

\newpage

%% ===== Page 549 =====

et on suppose que $\mathbb{G}^+ \neq \mathbb{G}$, donc

(3) $$\mathbb{G} \simeq \{1, \tau'\} \cdot \mathbb{G}^+$$
$$(1/2 \text{ direct})$$

et on a une suite exacte canonique

(4) $$1 \longrightarrow \mathbb{G}^+ \longrightarrow \mathbb{G} \longrightarrow \{\pm 1\} \longrightarrow 1$$

On désignera aussi par $\nu_\mathbb{G}$ la surjection

(5) $$\mathbb{G} \xrightarrow{\varepsilon_\mathbb{G}} \{\pm 1\} \simeq \mathbb{Z}/2\mathbb{Z}$$
$$(-1)^\nu \leftarrow \nu$$

De façon générale, on désignera (par un léger abus de langage) par une même lettre des éléments de $\operatorname{Gl}(2, \hat{\mathbb{Z}})$ - tels $\varepsilon_0, \varepsilon_1, \tau=\tau_\infty, \tau'=\tau'_\infty$, $\sigma=\sigma_\infty, \rho=\rho_\infty$ etc - et leurs images dans $\mathbb{G}$.

2º) Soit 
[MARGIN: Les groupes $Z_{\rho,\sigma}$, $Z_{\rho,\tau}$]

(6) $$Z_{\rho,\sigma} \subset \operatorname{Aut}(\mathbb{G}^+) \times \hat{\mathbb{Z}}^* \times \{\pm 1\} \times \{\pm 1\}$$
$$\| dfn$$
$$ \{ (u, \mu, \varepsilon, \varepsilon') \mid \begin{array}{l} u(\varepsilon_0) = {\varepsilon_0}^\mu \\ u(\rho) \text{ conjugué dans } \mathbb{G}^+ \text{ à } \rho^{\varepsilon'} \\ u(\sigma) \text{ \hspace{2cm} à } \sigma^{-\varepsilon} \end{array} \}$$

C'est un sous-groupe du groupe produit.

Posons

(7) $$\begin{cases} L_\rho = \text{sous-groupe de } \mathbb{G} \text{ engendré par } \rho \\ L_\sigma = \text{ \hspace{4cm} } \sigma \end{cases}$$

de sorte qu'on a deux épimorphismes

(8) $$\begin{cases} \hat{\mathbb{Z}}/6\mathbb{Z} \xrightarrow{\text{épi}} L_\rho \quad n \mapsto \rho^n \\ \hat{\mathbb{Z}}/4\mathbb{Z} \longrightarrow L_\sigma \quad n \mapsto \sigma^n \end{cases}$$

\newpage

%% ===== Page 550 =====

un diagramme commutatif
et des homomorphismes

(9)
\begin{tikzcd}
& L_\rho \ar[dr] & \\
\{\pm 1\} \ar[ur] \ar[dr] & & \check{\mathbb{G}}^+ \\
& L_\sigma \ar[ur] &
\end{tikzcd}

- Si $\check{\mathbb{G}}$ est "assez gros", on aura d'ailleurs
que $\{\pm 1\} \hookrightarrow \check{\mathbb{G}}^+$ i.e. "$-1$" $\neq \mathbb{1}$ dans $\check{\mathbb{G}}^+$, et
même
$$L_\rho \cap L_\sigma = \{\pm 1\}$$

Les conditions (6) impliquent donc que

(10)
$$
\begin{cases}
u(L_\rho) \text{ } \check{\mathbb{G}}^+ \text{-conjugué à } L_\rho \\
u(L_\sigma) \text{ } \check{\mathbb{G}}^+ \text{-conjugué à } L_\sigma
\end{cases}
$$

Si $\check{\mathbb{G}}^+$ est "assez gros" pour que $\rho$ et $\rho^{-1}$, $\sigma$ et $\sigma^{-1}$
n'y soient pas conjugués, alors un élément
$(u,\mu,\varepsilon,\varepsilon')$ de $Z_{\rho,\sigma}$ est connu par la seule
connaissance de $(u,\mu)$. Si d'autre part l'hom.

(11)
\begin{tikzcd}
\hat{\mathbb{Z}} \ar[r] & \check{\mathbb{G}}^+ \\
1 \ar[r, mapsto] & \varepsilon_0
\end{tikzcd}

est injectif, alors $\mu$ est connu quand on connait
$u$. Donc si $\check{\mathbb{G}}^+$ est assez gros pour satisfaire
aux trois conditions envisagées (p.ex. pour
qu'on ait un "homomorphisme de bases"

(12)
$$\check{\mathbb{G}} \longrightarrow \operatorname{Gl}(2, \hat{\mathbb{Z}})$$

(transformant $\rho, \sigma, \tau$ en les éléments de ce nom),
alors un élément de $Z_{\rho,\sigma}$ est connu par la seule

\newpage

%% ===== Page 551 =====

connaissance de $u \in \operatorname{Aut}(\widehat{G}^+)$.

On a un hom.

$$ (13) \quad L_0 \xrightarrow{i_{L_0, Z}} Z_{\rho, \sigma} $$
$$ l \mapsto (int(l), 1, 1, 1) $$

dont le noyau est [unclear: évidemment], le sous-groupe $L_0 \cap Cent(\widehat{G}^+)$. Notons que l'on a

$$ (14) \quad L_0 \cap Cent(\widehat{G}^+) \subset L_0 \cap Cent_{\widehat{G}^+}(\rho) = L_0 \cap L_1 $$

car si $l \in L_0$ et $\rho$ commute à $l$, i.e. $l = \rho(l)$, [crossed out: $l = p(l)$ $\rho(L_0) = L_1$,] on a $l \in L_1$. Donc si
$L_0 \cap L_1 = \{ \pm 1 \}$ (ce qui est le cas si (12) existe,
ou s'il existe seulement

$$ (11') \quad \widehat{G} \to \operatorname{Gl}(2, \hat{\mathbb{Z}})' \quad \text{modulaire} \quad ) $$

alors on a

$$ (15) \quad L_0 \cap Cent(\widehat{G}^+) = \{1\} \quad , $$

et (13) est injectif. Nous allons le supposer, pour
fixer les idées.

On pose

$$ (16) \quad \Gamma_{\rho, \sigma} = Z_{\rho, \sigma} / Im L_0 \quad , $$

de sorte qu'on aura une suite exacte canonique

$$ (17) \quad 1 \to L_0 \xrightarrow{i_{L_0, Z}} Z_{\rho, \sigma} \xrightarrow{\pi_Z} \Gamma_{\rho, \sigma} \to 1 \quad . $$

Les hom. canoniques de project. $(u, \mu, \varepsilon, \varepsilon') \mapsto \mu, \varepsilon, \varepsilon'$
se factorisent par $\Gamma_{\rho, \sigma}$ et fournissent
des homomorphismes canoniques

\newpage

%% ===== Page 552 =====

(18) 
$$ \gamma_{p,\sigma} \xrightarrow{\chi_\gamma} \hat{\mathbb{Z}}^* $$
$$ \gamma_{p,\sigma} \xrightarrow{\varepsilon_\gamma} \{ \pm 1 \} $$
$$ \gamma_{p,\sigma} \xrightarrow{\varepsilon'_\gamma} \{ \pm 1 \} $$
[MARGIN: 599] soit $\gamma_{p,\sigma} \xrightarrow{(\chi_\gamma, \varepsilon_\gamma, \varepsilon'_\gamma)} \hat{\mathbb{Z}}^* \times \{ \pm 1 \} \times \{ \pm 1 \}$

les composés de ceux-ci avec $Z_{p,\sigma} \to \gamma_{p,\sigma}$
sont notés
(19) $$ \chi_Z = \chi_\gamma \circ \delta_Z , \quad \varepsilon_Z = \varepsilon_\gamma \circ \delta_Z , \quad \varepsilon'_Z = \varepsilon'_\gamma \circ \delta_Z , $$
de sorte que
(20) $$ \chi_Z(u, \mu, \varepsilon, \varepsilon') = \mu , \quad \varepsilon_Z(u, \mu, \varepsilon, \varepsilon') = \varepsilon , \quad \varepsilon'_Z(u, \mu, \varepsilon, \varepsilon') = \varepsilon' $$

On aurait envie de désigner par $\gamma_{p,\sigma}^0$ le
noyau de $(\varepsilon_\gamma, \varepsilon'_\gamma)$, par $\gamma_{p,\sigma}^1$ celui de $\chi_\gamma$,
par $\gamma_{p,\sigma}^{0+}$ leur intersection, et d'introduire
par les notations $Z_{p,\sigma}^0$, $Z_{p,\sigma}^1$, $Z_{p,\sigma}^{0+}$ leurs
images inverses dans $Z_{p,\sigma}$, i.e. les noyaux respect.
de $(\varepsilon_Z, \varepsilon'_Z)$, de $\chi_Z$ et de $(\chi_Z, \varepsilon_Z, \varepsilon'_Z)$.

[MARGIN: Les groupes $SG_{p,\sigma}$, $G_{p,\sigma}$]
3°) L'homomorphisme 
(21) $$ Z_{p,\sigma} \longrightarrow \operatorname{Aut}(\hat{\mathfrak{S}}^+) $$
$$ (u,\mu,\varepsilon,\varepsilon') \longmapsto u $$

d'où $Z_{p,\sigma}$ opère sur $\hat{\mathfrak{S}}^+$, et on peut
donc considérer le produit $1/2$ direct
(22) $$ SG_{p,\sigma} = Z_{p,\sigma} \underset{(1/2\text{-direct})}{\cdot} \hat{\mathfrak{S}}^+ $$

\newpage

%% ===== Page 553 =====

On a un hom.
(23)
$$L_0 \xrightarrow{i_{L_0, SG}} S_0 G_{\rho,\sigma} = Z_{\rho,\sigma} \cdot \widehat{C}^+$$
$$l \longmapsto i_{L_0,Z}(l) \cdot l^{-1}$$
dont l'image est encore invariante, d'où
(24)
$$G_{\rho,\sigma} = S_0 G_{\rho,\sigma} / \text{Im } i_{L_0, SG} ,$$
de sorte qu'on a une suite exacte
(25)
$$1 \longrightarrow L_0 \xrightarrow{i_{L_0, SG}} S_0 G_{\rho,\sigma} \longrightarrow G_{\rho,\sigma} \longrightarrow 1$$
et un diagramme (d'homomorphismes)
(26)
\begin{tikzcd}
& Z_{\rho,\sigma} \ar[dr, "i_Z"] \\
L_0 \ar[ur, "i_{L_0,Z}"] \ar[dr, "i_{L_0,\widehat{C}}"] & & G_{\rho,\sigma} \\
& \widehat{C}^+ \ar[ur, "i_{\widehat{C}^+}"']
\end{tikzcd}
( a priori $\operatorname{Ker}(\widehat{C}^+ \to G_{\rho,\sigma}) \simeq \operatorname{Ker}(L_0 \to Z_{\rho,\sigma})$
$\simeq L_0 \cap \text{Cent}(\widehat{C}^+)$, mais par hyp. (15)
il se réduit à $\{1\}$ ), s'insérant dans un
diagramme de suites exactes
(27)
\begin{tikzcd}
1 \ar[r] & L_0 \ar[r, "i_{L_0,Z}"] \ar[d, "i_{L_0,\widehat{C}}"'] & Z_{\rho,\sigma} \ar[r, "\delta_Z"] \ar[d] & \Gamma_{\rho,\sigma} \ar[r] \ar[d, equal] & 1 \\
1 \ar[r] & \widehat{C}^+ \ar[r, "i_{\widehat{C}^+}"'] & G_{\rho,\sigma} \ar[r, "\delta_G"'] & \Gamma_{\rho,\sigma} \ar[r] & 1
\end{tikzcd}
On posera
(28)
$$\chi_G = \chi_\Gamma \delta_G , \varepsilon_G = \varepsilon_\Gamma \delta_G , \varepsilon'_G = \varepsilon'_\Gamma \delta_G$$

\newpage

%% ===== Page 554 =====

d'où un hom.
(29) $$G_{\rho,\sigma} \xrightarrow{(\chi_{G}, \varepsilon_{G}, \varepsilon'_{G})} \hat{\mathbb{Z}}^\times \times \{\pm 1\} \times \{\pm 1\}$$

nous conduisant à introduire $G_{\rho,\sigma}^+, G_{\rho,\sigma}^-, G_{\rho,\sigma}^{+0}$ comme d'habitude.

L'injection $i_{\mathfrak{G}^+} : \mathfrak{G}^+ \hookrightarrow G_{\rho,\sigma}$ peut se prolonger en un plongement [unclear: bijectif] (noté $i_{\mathfrak{G}}$)
(30) $$i_{\mathfrak{G}} : \mathfrak{G} \hookrightarrow G_{\rho,\sigma}$$
que je définis par sa valeur sur $\tau = \sigma \tau'$, par
(31) $$i_{\mathfrak{G}}(\tau) = i_Z(\tau_Z)$$

[MARGIN: Si $\tau_{\mathfrak{G}}$ désigne l'image de $\tau = \sigma \tau' \in \mathfrak{G}$ par $i_{\mathfrak{G}}$ dans $G_{\rho,\sigma}$, l'id. (31) s'écrit $\tau_{\mathfrak{G}} = i_Z(\tau_Z)$ dans $G_{\rho,\sigma}$]
(32) $$ \begin{cases} \tau_Z \in Z_{\rho,\sigma} \text{ défini comme} \\ \tau_Z = (\text{int}(\tau) \mid \mathfrak{G}^+, -1, -1, -1) \\ \qquad \in \operatorname{Aut}(\mathfrak{G}^+) \times \hat{\mathbb{Z}}^\times \times \{\pm 1\} \times \{\pm 1\} \end{cases} $$

Désormais nous identifions $\mathfrak{G}$ à un sous-groupe de $G_{\rho,\sigma}$, et nous désignons par les mêmes lettres des éléments de $\mathfrak{G}$ tels que $(... , \tau, \tau', \rho, \sigma, \xi_0, \xi_1, \text{ etc} )$ et leurs images dans $G_{\rho,\sigma}$. On aura d'ailleurs

[MARGIN: Noter que $\tau_{\mathfrak{G}}^2 \notin Z_{\rho,\sigma}$ par contre $\tau_{\mathfrak{G}} \dots$ [crossed out]]
(32bis) $$\chi_{G}|\mathfrak{G} = \varepsilon_{G}|\mathfrak{G} = \varepsilon'_{G}|\mathfrak{G} = \varepsilon_{\mathfrak{G}} : \mathfrak{G} \to \{\pm 1\} \text{ (trivial sur } \mathfrak{G}^+, \text{ -1 sur } \mathfrak{G} \smallsetminus \mathfrak{G}^+)$$

4º) On [unclear: a déjà] noté que l'élément $\tau$ de $G_{\rho,\sigma}$ normalise $L_\rho$ et $L_\sigma$, donc normalise aussi

[MARGIN: Les groupes $Z_\rho, Z_\sigma, N_\rho, N_\sigma \dots SN_\rho, SN_\sigma \dots$]
(33) $$ \begin{cases} Z_\rho \overset{\text{df}}{=} \text{Cent}_{G_{\rho,\sigma}}(\rho) = \text{Cent}_{G_{\rho,\sigma}}(L_\rho) \\ Z_\sigma \overset{\text{df}}{=} \text{Cent}_{G_{\rho,\sigma}}(\sigma) = \text{Cent}_{G_{\rho,\sigma}}(L_\sigma) \end{cases} $$

\newpage

%% ===== Page 555 =====

dans $\{1,\tau'\}$ opère sur ces groupes, et on
peut considérer

(34)
$$
\begin{cases}
N_\rho = \{1,\tau'\}. Z_\rho \\
N_\sigma = \{1,\tau'\}. Z_\sigma
\end{cases} \quad \text{(produits semi-directs).}
$$

On a des suites exactes canoniques

(35)
$$1 \to Z_\rho \to N_\rho \xrightarrow{\varepsilon_\rho} \{\pm 1\} \to 1$$
$$1 \to Z_\sigma \to N_\sigma \xrightarrow{\varepsilon_\sigma} \{\pm 1\} \to 1$$

et un diagramme commutatif d'hom.
canoniques

(36)
\begin{tikzcd}
& D_\rho \ar[r] & N_\rho \ar[dr, "i_\rho"] & \\
\{1,\tau'\} \times \{\pm 1\} \ar[ur] \ar[dr] & & & G_{\rho,\sigma} \\
& D_\sigma \ar[r] & N_\sigma \ar[ur, "i_\sigma"'] &
\end{tikzcd}

où

(37)
$$
\begin{cases}
D_\rho \simeq D_6 \subset \operatorname{Gl}(2, \mathbb{Z}) & \text{engendré par } \rho, \tau' \\
D_\sigma \simeq D_4 \subset \operatorname{Gl}(2, \mathbb{Z}) & \text{------- } \sigma, \tau' \\
\{1,\tau'\} \times \{\pm 1\} = D_\rho \cap D_\sigma \text{ dans } \operatorname{Gl}(2, \mathbb{Z})
\end{cases}
$$

dans lesquels

(38)
$$
\begin{cases}
D_\rho = \{1,\tau'\}.L_\rho & \text{(où } L_\rho \text{ est d'ordre } 6 \text{ dans } \bar{\mathbb{S}}^+ \text{)} \\
D_\sigma = \{1,\tau'\}.L_\sigma & \text{( ------- } \sigma \text{ d'ordre } 4 \text{ --------)}
\end{cases}
$$

et les hom.

$$D_\rho \to N_\rho, \quad D_\sigma \to N_\sigma$$

sont évidents. Les hom. $i_\rho, i_\sigma$ sont inj., bien sûr.

Rappelons que

(39)
$$
\begin{cases}
i_\rho \text{ i-jectif} \iff \tau' \notin Z_\rho \iff \rho^2 \neq 1 \iff \rho \neq \pm 1 \iff \bar{\mathbb{S}}^+ \text{ est [unclear]} \\
i_\sigma \text{ i-jectif} \iff \tau' \notin Z_\sigma \iff \sigma^2 \neq 1 \iff \text{[unclear]}
\end{cases}
$$

\newpage

%% ===== Page 556 =====

[MARGIN: Si $g \in G_{\rho,\sigma}$ et $\varepsilon = \varepsilon_G(g)$, $g^{-1} \rho g = \rho^\varepsilon$, $int(g)(\rho) = \rho^\varepsilon$ donc $int(g)(\sigma) = \sigma^\varepsilon$
[unclear]
]

Faisant opérer $G_{\rho,\sigma}$ sur $\widehat{\mathbb{C}}^+$, sous-groupe distingué, via autom. int., on voit que l'image de $G_{\rho,\sigma}$ dans le gr. des [unclear] de $\widehat{\mathbb{C}}^+$ est [unclear]
[crossed out text]
(posant
(40) $N_\rho^* = \{ x \in N_\rho \mid \varepsilon_\rho(x) = \varepsilon_G(i_\rho(x)) \}$, $N_\sigma^* = \{ x \in N_\sigma \mid \varepsilon_\sigma(x) = \varepsilon_G(i_\sigma(x)) \}$
les homomorphismes composés

$$ N_\rho^* \xrightarrow{i_\rho} G_{\rho,\sigma} \xrightarrow{\delta_a} \Gamma_{\rho,\sigma} , \delta_\rho^* = \delta_G \circ i_\rho^* $$
$$ N_\sigma^* \xrightarrow{i_\sigma} G_{\rho,\sigma} \xrightarrow{\delta_a} \Gamma_{\rho,\sigma} , \delta_\sigma^* = \delta_G \circ i_\sigma^* $$
sont des épimorphismes.

a) Si $\rho^2 \neq 1$, alors on a commut. [unclear] des
diagrammes de suites exactes
(41)
\begin{tikzcd}
1 \ar[r] & Z_\rho \ar[r] \ar[d, "j_\rho"'] & N_\rho^* \ar[r, "\varepsilon_\rho"] \ar[d, "\delta_\rho^* \text{ épi}"] & \{\pm 1\} \ar[r] \ar[d, equal] & 1 \\
1 \ar[r] & \Gamma_{\rho,\sigma}' \ar[r] & \Gamma_{\rho,\sigma} \ar[r, "\varepsilon_\Gamma'"'] & \{\pm 1\} \ar[r] & 1
\end{tikzcd}
où l'on a posé
(42) $$ \Gamma_{\rho,\sigma}' = \operatorname{Ker}(\Gamma_{\rho,\sigma} \xrightarrow{\varepsilon_\Gamma'} \{\pm 1\}) ; $$

b) Si $\sigma^2 \neq 1$ (cf (39)), alors on a [unclear] diag. comm. de suites exactes
(43)
\begin{tikzcd}
1 \ar[r] & Z_\sigma \ar[r] \ar[d, "j_\sigma"'] & N_\sigma^* \ar[r, "\varepsilon_\sigma"] \ar[d, "\delta_\sigma^* \text{ épi}"] & \{\pm 1\} \ar[r] \ar[d, equal] & 1 \\
1 \ar[r] & \Gamma_{\rho,\sigma}'' \ar[r] & \Gamma_{\rho,\sigma} \ar[r, "\varepsilon_\Gamma''"'] & \{\pm 1\} \ar[r] & 1
\end{tikzcd}
où
(44) $$ \Gamma_{\rho,\sigma}'' = \operatorname{Ker}(\Gamma_{\rho,\sigma} \xrightarrow{\varepsilon_\Gamma''} \{\pm 1\}) . $$

\newpage

%% ===== Page 557 =====

NB Les notations $\varepsilon_\gamma, \varepsilon_{\gamma'} : \Gamma_{\rho,\sigma} \to \{\pm 1\}$
pourraient prêter à confusion, la première
exprimant l'action d'un $u$ sur $L_\sigma$ (non sur
$L_\rho$). Une notation un peu plus lourde,
mais ne prêtant pas à confusion, serait
$\varepsilon_{\gamma,\rho}, \varepsilon_{\gamma,\sigma}$ - ou on écrira seulement $\varepsilon_\rho, \varepsilon_\sigma$,
la confusion avec les hom de ces noms
$N_\rho^* \to \{\pm 1\}$, $N_\sigma^* \to \{\pm 1\}$ étant exclue lorsque
$\rho^2 \neq 1, \sigma^2 \neq 1$, à cause des compatibilités (41) (43).

On a donc construit, à l'aide
de (1) i.e. de $\widehat{\mathcal{L}}$ munie de $\rho, \sigma, \tau'$, des
groupes $C^+, G_{\rho,\sigma}, \Gamma_{\rho,\sigma}, Z_{\rho,\sigma}^*, N_\rho^*, N_\sigma^*$, s'insérant
dans un diagramme commutatif de
groupes

(45)
\begin{tikzcd}
L_0 \simeq S Z_{\rho,\sigma} \ar[r, hook] \ar[d, hook] \ar[dd, hook, bend right=40] & Z_{\rho,\sigma}^* \ar[d, hook] \ar[dd, hook, bend left=40] \ar[dddr, dashed, "\delta_Z"] & & \\
S N_\rho^* \ar[r, hook] \ar[dd, hook, bend right=40] & N_\rho^* \ar[dd, hook, bend left=40] \ar[ddr, dashed, "\delta_\rho"'] & & \hat{\mathbb{Z}}^* \times \{\pm 1\} \times \{\pm 1\} \\
S N_\sigma^* \ar[r, hook] \ar[d, hook] & N_\sigma^* \ar[d, hook] \ar[dr, dashed, "\delta_\sigma"'] & & \\
1 \ar[r] & C^+ \ar[loop below, "G"'] \ar[r, "i^+"'] & G_{\rho,\sigma} \ar[r, "\delta_G"'] & \Gamma_{\rho,\sigma} \ar[r] \ar[uuu, hook, "(\chi{,} \varepsilon_{\gamma,\rho}{,} \varepsilon_{\gamma,\sigma})"'] & 1
\end{tikzcd}

où on a posé

(46)
$$ S Z_{\rho,\sigma} = i_{L_0, Z}(L_0) = \ker(\delta_Z : Z_{\rho,\sigma}^* \to \Gamma_{\rho,\sigma}) $$
$$ S N_\rho^* = \ker(\delta_\rho : N_\rho^* \to \Gamma_{\rho,\sigma}) = \text{Cent}_{C^+}(\rho) $$
$$ S N_\sigma^* = \ker(\delta_\sigma : N_\sigma^* \to \Gamma_{\rho,\sigma}) = \text{Cent}_{C^+}(\sigma) $$

[MARGIN: NB On peut d'ailleurs les noter $S Z_\rho, S Z_\sigma$ au lieu de $S N_\rho^*, S N_\sigma^*$.]

\newpage

%% ===== Page 558 =====

[MARGIN: 
On a 
$\operatorname{Im}(N_\rho \to \widehat{C}^+) = Nor_{\widehat{C}^+}(L_\rho)$
(46 bis) $\operatorname{Im}(S N_\rho^* \to \widehat{C}^+) = Nor_{\widehat{C}^+}(L_\rho^+)$
(ou $L_\rho^+ = L_\rho$ si $\rho^2 \neq 1$ resp. $\sigma^2 \neq 1$)
$N_\rho \cap \widehat{C} = Nor_{\widehat{C}}(L_\rho)$
]

qui a des suites exactes

(47)
$$ \left\{ \begin{array}{l} 1 \longrightarrow S N_\rho^* \longrightarrow N_\rho^* \stackrel{\delta_\rho}{\longrightarrow} \Gamma_{\rho,\sigma} \longrightarrow 1 \\ 1 \longrightarrow S N_\sigma^* \longrightarrow N_\sigma^* \stackrel{\delta_\sigma}{\longrightarrow} \Gamma_{\rho,\sigma} \longrightarrow 1 \end{array} \right. $$

s'envoient dans la suite exacte

$$1 \longrightarrow \widehat{C}^+ \longrightarrow G_{\rho,\sigma} \longrightarrow \Gamma_{\rho,\sigma} \longrightarrow 1 \qquad ,$$

de façon canonique : (47). De plus, les
hom. $i_\rho : N_\rho^* \to G_{\rho,\sigma}$ et $i_\sigma : N_\sigma^* \to G_{\rho,\sigma}$ sont
en réalité des mon. : ils sont injectifs (cf (39)),
ce qu'on a indiqué sur le diagramme (45)
par les signes $\hookrightarrow$ (signe d'inclusions pointillés).

On introduit aussi

(48)
$$ \left\{ \begin{array}{l} \chi_\rho = \chi_{\Gamma} \circ \delta_\rho : N_\rho^* \longrightarrow \hat{\mathbb{Z}}^* \\ \chi_\sigma = \chi_{\Gamma} \circ \delta_\sigma : N_\sigma^* \longrightarrow \hat{\mathbb{Z}}^* \end{array} \right. \qquad , $$

dont les noyaux (qui contiennent $S N_\rho^*$
resp. $S N_\sigma^*$, et qu'on a gardé de confondre
avec ceux-ci) sont notés comme de juste

$N_\rho^{**}, N_\sigma^{**}$.

NB Il faudrait encore envisager les composées

(49)
$$ \left\{ \begin{array}{l} N_\rho^* \longrightarrow \Gamma_{\rho,\sigma} \stackrel{\varepsilon_{\rho,\sigma}}{\longrightarrow} \{\pm 1\} \\ N_\sigma^* \longrightarrow \Gamma_{\rho,\sigma} \stackrel{\varepsilon_{\rho,\sigma}}{\longrightarrow} \{\pm 1\} \end{array} \right. $$

on verra que dans les "bons cas", la restriction
de ces hom. à $Z_\rho, Z_\sigma$ respectivement sont [unclear: surj.]

\newpage

%% ===== Page 559 =====

583

[MARGIN: Cancelé / En révisant (15) / dans le cas où / une certaine [unclear] / $Norm$ de $(L_\rho)$ / et idem pour $L_\sigma, L_0$]

On pose

(50)
$$
\begin{cases} 
S N_\rho^* = \operatorname{Ker}( N_\rho^* \xrightarrow{\delta_\rho} \gamma_{\rho,\sigma} ) \doteq \operatorname{Norm}_{\hat{\mathbb{Z}}}(L_\rho) \\
S N_\sigma^* = \operatorname{Ker}( N_\sigma^* \xrightarrow{\delta_\sigma} \gamma_{\rho,\sigma} ) = \operatorname{Norm}_{\hat{\mathbb{Z}}}(L_\sigma) \\
S B_{\rho,\sigma} = \operatorname{Ker}( B_{\rho,\sigma} \xrightarrow{\delta_B} \gamma_{\rho,\sigma} ) = [unclear] \cap \operatorname{Norm}_{\hat{\mathbb{Z}}}(L_0) 
\end{cases}
$$

[MARGIN: Les groupes $S_0 M_{\rho,\sigma}, M_{\rho,\sigma}, S_0 \Gamma_{\rho,\sigma}, \Gamma_{\rho,\sigma}$]

$$
\begin{cases} 
S_0 M_{\rho,\sigma} = Z_{\rho,\sigma} \times_{\gamma_{\rho,\sigma}} N_\rho^* \times_{\gamma_{\rho,\sigma}} N_\sigma^* \times_{\gamma_{\rho,\sigma}} C_{\rho,\sigma} \\
M_{\rho,\sigma} = N_\rho^* \times_{\gamma_{\rho,\sigma}} N_\sigma^* \times_{\gamma_{\rho,\sigma}} C_{\rho,\sigma} 
\end{cases}
$$

(51)
$$
\begin{cases} 
S_0 \Gamma_{\rho,\sigma} = Z_{\rho,\sigma} \times_{\gamma_{\rho,\sigma}} N_\rho^* \times_{\gamma_{\rho,\sigma}} N_\sigma^* \\
\Gamma_{\rho,\sigma} = N_\rho^* \times_{\gamma_{\rho,\sigma}} N_\sigma^* 
\end{cases}
$$

d'où un diagramme de suites ex. [unclear]

(52)
\begin{tikzcd}
1 \ar[r] & C^+ \ar[r] \ar[d, equal] & S_0 M_{\rho,\sigma} \ar[r] \ar[d] & S_0 \Gamma_{\rho,\sigma} \ar[r] \ar[d] & 1 \\
1 \ar[r] & C^+ \ar[r] & M_{\rho,\sigma} \ar[r] \ar[d] & \Gamma_{\rho,\sigma} \ar[r] \ar[d] & 1 \\
& & 1 & 1
\end{tikzcd}

et des suites ex.

(53)
$$ 1 \to S Z_{\rho,\sigma} \times S N_\rho^* \times S N_\sigma^* \times C^+ \xrightarrow{\simeq L_0} S_0 M_{\rho,\sigma} \to \gamma_{\rho,\sigma} \to 1 $$

(54)
$$ 1 \to S N_\rho^* \times S N_\sigma^* \times C^+ \to M_{\rho,\sigma} \to \gamma_{\rho,\sigma} \to 1 $$

(55)
$$ 1 \to S Z_{\rho,\sigma} \times S N_\rho^* \times S N_\sigma^* \xrightarrow{\simeq L_0} S_0 \Gamma_{\rho,\sigma} \to \gamma_{\rho,\sigma} \to 1 $$

(56)
$$ 1 \to S N_\rho^* \times S N_\sigma^* \to \Gamma_{\rho,\sigma} \to \gamma_{\rho,\sigma} \to 1 $$

\newpage

%% ===== Page 560 =====

On a aussi

(57) $$S_0 \mathcal{M}_{\rho,\sigma} \simeq \mathcal{M}_{\rho,\sigma} \times_{\gamma_{\rho,\sigma}} \mathbb{Z}_{\rho,\sigma}$$
(58) $$S_0 \Gamma_{\rho,\sigma} \simeq \Gamma_{\rho,\sigma} \times_{\gamma_{\rho,\sigma}} \mathbb{Z}_{\rho,\sigma}$$

i.e. les ext. $S_0 \mathcal{M}_{\rho,\sigma}$, $S_0 \Gamma_{\rho,\sigma}$ de $\mathcal{M}_{\rho,\sigma}, \Gamma_{\rho,\sigma}$
par $L_0$ se déduisent de l'extension $\mathbb{Z}_{\rho,\sigma}$ de
$\gamma_{\rho,\sigma}$ par $L_0$ par image inverse par
$\mathcal{M}_{\rho,\sigma} \to \gamma_{\rho,\sigma}$, $\Gamma_{\rho,\sigma} \to \gamma_{\rho,\sigma}$ - et --

aussi

(59) $$S_0 \mathcal{M}_{\rho,\sigma} \simeq \mathcal{M}_{\rho,\sigma} \times_{\Gamma_{\rho,\sigma}} S_0 \Gamma_{\rho,\sigma}$$

i.e. l'extension $S_0 \mathcal{M}_{\rho,\sigma}$ de $\mathcal{M}_{\rho,\sigma}$ par $L_0$
se déduit de l'extension $S_0 \Gamma_{\rho,\sigma}$ de
$\Gamma_{\rho,\sigma}$ par $L_0$, par image inverse par
$\mathcal{M}_{\rho,\sigma} \to \Gamma_{\rho,\sigma}$.

On introduit un élément

(60) $\tau_M \in S_0 \mathcal{M}_{\rho,\sigma} \quad (\tau_M = (\overset{a}{\tau'_Z}, \overset{b}{\tau'_N}, \overset{a}{\tau''_N}, \overset{b}{\tau'_G}))$

[unclear: et isomorphie entre les groupes introduits, je vais les indiquer ci-dessous]

où $\tau'_Z, \tau'_N, \tau''_N, \tau'_G$
"relevés" de $\tau'$, pris resp. dans $\mathbb{Z}_{\rho,\sigma}$, $N_\rho^*, N_\sigma^*, C_{\rho,\sigma}$. Ceux-
ci ont bien tous la même image dans $\gamma_{\rho,\sigma}$ et [unclear]
des $\gamma_{\rho,\sigma}$, donc définissent un él. de $S_0 \mathcal{M}_{\rho,\sigma}$, dont
l'opération sur $\widetilde{C}^+$ est celle de $\tau$. D'où un plon-
gement

(61) $\begin{tikzcd} \widetilde{C} \ar[r, hook] & S_0 \mathcal{M}_{\rho,\sigma} \end{tikzcd} \quad \begin{matrix} : \text{identifiant } \widetilde{C}^+ \\ \tau \mapsto \tau_M \end{matrix}$

\newpage

%% ===== Page 561 =====

584

(62)
\begin{tikzcd}
L_0 \times \widehat{SN}^*_\rho \times \widehat{SN}^*_\sigma \ar[rrr, "pr"] \ar[dr, "pr"'] \ar[dd, "pr"'] & & & L_0 \ar[d, "\theta_Z"] & \\
& L_0 \times \widehat{SN}_\rho \times \widehat{SN}^*_\sigma \ar[dl, "pr"'] & S_0 \Gamma_{\rho,\sigma} \ar[ur, "pr"'] \ar[r] \ar[dr, "p_\Gamma"'] & Z_{\rho,\sigma} \ar[d, "\delta_Z"] & \\
S\widehat{N}_\rho \times S\widehat{N}_\sigma \ltimes \widehat{C}^+ \ar[r, "pr"] & S_0 M_{\rho,\sigma} \ar[ur] \ar[d, "p_M"'] \ar[r, "\theta_\rho"] & N_\rho^* \ar[r, "\delta_\rho"] & \Gamma_{\rho,\sigma} \ar[r, "\varepsilon_\rho=\varepsilon_\rho' \mapsto \pm 1", shift left=1.5ex] \ar[r, "\varepsilon_{\rho,\sigma}=\varepsilon_\sigma \mapsto \pm 1"', shift right=1.5ex] \ar[dd, "\delta_G"'] & \hat{\mathbb{Z}}^* \\
\widehat{C}^+ \ar[u, hook, "inj"] \ar[r, dashed] \ar[d] & M_{\rho,\sigma} \ar[urr] \ar[r, "\theta_\sigma"'] \ar[drr, "\theta_G"'] & N_\sigma^* \ar[ur, "\delta_\sigma"'] & & \\
\{1,\tau\} \ltimes \widehat{C}^+ \ar[rrr, dashed, "\overline{\theta}_G"'] & & & G_{\rho,\sigma} &
\end{tikzcd}

On a résumé certaines des relations entre les groupes introduits jusqu'à présent dans le diagramme commutatif (62), que je commente brièvement.
a) Toutes les flèches sont soit injectives ($\hookrightarrow$) soit surjectives ($\twoheadrightarrow$), à l'exception peut-être des flèches issues de $\Gamma_{\rho,\sigma}$.
b) Les dix carrés (six faces d'un cube, et quatre carrés I, II, III, IV) y contenues sont commutatifs.
c) Les flèches marquées $pr$ sont des projections canoniques de groupes produits sur des produits partiels. La flèche marquée $inj$ $\widehat{C}^+ \hookrightarrow S\widehat{N}_\rho \times S\widehat{N}_\sigma \ltimes \widehat

\newpage

%% ===== Page 562 =====

e) Pour ne pas surcharger le diagramme,
je n'ai pas introduit dans ce diagramme des
notations pour certaines flèches composées utiles,
telles

$$ (63) \left\{ \begin{array}{l} \delta_\Gamma : \Gamma_{\rho,\sigma} \to G_{\rho,\sigma} \\ \delta_{S_0\Gamma} : S_0\Gamma_{\rho,\sigma} \to G_{\rho,\sigma} \\ \delta_M : M_{\rho,\sigma} \to G_{\rho,\sigma} \\ \delta_{S_0M} : S_0M_{\rho,\sigma} \to G_{\rho,\sigma} \end{array} \right. $$

et

$$ (64) \left\{ \begin{array}{l} \chi_\Gamma = \chi_{G_{\rho,\sigma}} \circ \delta_\Gamma : \Gamma_{\rho,\sigma} \to \hat{\mathbb{Z}}^\times \\ \chi_{S_0\Gamma} = \chi_{G_{\rho,\sigma}} \circ \delta_{S_0\Gamma} : S_0\Gamma_{\rho,\sigma} \to \hat{\mathbb{Z}}^\times \\ \chi_M = \chi_{G_{\rho,\sigma}} \circ \delta_M : M_{\rho,\sigma} \to \hat{\mathbb{Z}}^\times \\ \chi_{S_0M} = \chi_{G_{\rho,\sigma}} \circ \delta_{S_0M} : S_0M_{\rho,\sigma} \to \hat{\mathbb{Z}}^\times \end{array} \right. $$

et les composées $\varepsilon_{?,\rho}, \varepsilon_{?,\sigma} : ? \to \{\pm 1\}$, où $?$ peut
prendre toute valeur $Z, G$ (déjà introduits) $N_\rho^*, N_\sigma^*$
(on a aussi $\varepsilon_{N_\rho^*,\rho} = \varepsilon_\rho$ [unclear], $\varepsilon_{N_\sigma^*,\sigma} = \varepsilon_\sigma$ [unclear]),
$\Gamma, S_0\Gamma, M, S_0M$. Le noyau de $\chi_?$ est noté $?^+$
[MARGIN: il y a une petite impropriété de notations car $N_\rho^*, N_\sigma^*$ déjà utilisés dans un tout autre sens]
celui de $\varepsilon_{?,\rho}$ est noté $?'$, celui de $\varepsilon_{?,\sigma}$ est noté $?''$ -
on a $?^\circ = ?' \cap ?''$, $?^{+'} = ?^+ \cap ?'$, $?^{+''} = ?^+ \cap ?''$,
$?^{+\circ} = ?^+ \cap ?^\circ$, $S? = \text{Ker } \delta_?$, et donc

(65

\newpage

%% ===== Page 563 =====

585

Le noyau de $\varepsilon_{\hat{\mathbb{Z}}}$ n'est autre que $\hat{\mathbb{Z}}^+$. En tant
que sous-groupe de $M_{\rho,\sigma}$ (resp. de $\Gamma_{\rho,\sigma}$) $\hat{\mathbb{Z}}^+$
s'identifie à l'image inverse dans ce groupe \textsuperscript{d'un sous-groupe}
du sous-groupe $\{1, \tau'_\rho\}$ de $\bar{\Gamma}_{\rho,\sigma}$ (resp. $\{1, \bar{\tau}'_\rho\}$
de $\bar{\Gamma}_{\rho,\sigma}$), où $\tau'_\rho, \bar{\tau}'_\rho$ désignent les
images de $\tau' \in \widehat{G}$ dans $\bar{\Gamma}_{\rho,\sigma}$ et $\bar{\Gamma}_{\rho,\sigma}$ respectivement.

g) Le diagramme (62) a comme partie
principale

(67)
\begin{tikzcd}
& \bullet \ar[r] \ar[dr] & \bullet \ar[dr] \ar[rr, dashed] & & \\
\bullet \ar[ur] \ar[r] \ar[dr] & \bullet \ar[r] & \bullet \ar[r] \ar[rr, dashed] & \bullet & {[ \ \ ]} \\
& \bullet \ar[ur] \ar[r] & \bullet \ar[ur] \ar[rr, dashed] & &
\end{tikzcd}

Tous les groupes marqués dans ce \textsuperscript{dernier} diagramme, à
l'exception de $\hat{\mathbb{Z}}^*$, $\{\pm 1\}$, $\{\pm 1\}$ (mis entre crochets
$[ \ \ ]$) s'envoient dans $\bar{\Gamma}_{\rho,\sigma}$, qui est le terme
terminal de ce diagramme, relatif à $8$ groupes
(en ne comptant pas \sout{[unclear]} \textsuperscript{l'image inverse par une flèche pointillée}) dont $\bar{\Gamma}_{\rho,\sigma}$, déduits
de quatre d'entre eux $Z_{\rho,\sigma}$, $N_\rho^*$, $N_\sigma^*$, $G_{\rho,\sigma}$
(qui s'envoient dans $\bar{\Gamma}_{\rho,\sigma}$ par des épis) en prenant
des produits fibrés sur $\bar{\Gamma}_{\rho,\sigma}$ de certains
d'entre eux. Le terme initial étant $S_0 M_{\rho,\sigma}$,
qui est le produit fibré des quatre facteurs
précédents. En prenant les noyaux des hom.
de ces \sout{quatre} neuf groupes dans $\bar{\Gamma}_{\rho,\sigma}$, on trouve
un diagramme du même type combinatoire, relié

562

\newpage

%% ===== Page 564 =====

neuf groupes, dont $\{1\} = \operatorname{Ker}(\bar{\Gamma}_{p,\sigma} \to \bar{\Gamma}_{p,\sigma})$, formés de
$L_0 = S Z_{p,\sigma}$, $S N_p^*$, $S N_\sigma^*$ et $\hat{G}^+ = S G_{p,\sigma}$ et de certains
produits de ces groupes (dont le produit sur l'ens
d'indice vide, savoir $\{1\} = \operatorname{Ker}(\bar{\Gamma}_{p,\sigma} \xrightarrow{id} \bar{\Gamma}_{p,\sigma})$),
avec comme flèches [unclear] des projections
canoniques. Nous n'avons inséré dans (62) seule-
ment cinq parmi ces groupes produits, avec
les injections canoniques dans les produits
fibrés dont ils proviennent, et on s'est
dispensé d'y faire figurer égalem.
$S N_p^*, S N_\sigma^*, S G_{p,\sigma} = \hat{G}^+$, et les inclusions
canoniques $S N_p^* \hookrightarrow N_p^*, S N_\sigma^* \hookrightarrow N_\sigma^*, \hat{G}^+ = S G_{p,\sigma} \hookrightarrow G_{p,\sigma}$,
et les [unclear] des projections des
produits partiels de ces groupes et de $L_0$,
y compris les projections sur ces trois groupes, et
les hom de ces groupes et de $L_0$ dans le
groupe unité.

\newpage

%% ===== Page 565 =====

60) Lien avec $\Sigma_{p,\sigma}, S\mathcal{G}_{p,\sigma}$.

On pose

(68) $$ \Sigma_{p,\sigma} = \mathcal{G}^+ \times \mathcal{G}^+ \times \hat{\mathbb{Z}}^* \times \{\pm 1\} \times \{\pm 1\} $$
$$ \left\{ (\xi,\eta,\mu,\epsilon,\epsilon') \mid \xi \text{ } p\text{-conjugué à } p^{\epsilon'} \text{, } \eta \text{ } \sigma\text{-conjugué à } \sigma^\epsilon \text{, } \xi = \eta \varepsilon_1^{\mu-1} \right\} $$
[unclear: $p_1^\nu = v^{\frac{H-1}{2}}$]

NB Si $p$ non conjugué à $p^{-1}$ dans $\mathcal{G}^+$, $\epsilon'$ connu quand on connait $\xi, \eta$ - de même si $\sigma$ non conjugué à $\sigma^{-1}$ (donc $\sigma^2 \neq 1$), $\epsilon$ est déterminé par $\xi, \eta$. De même $\hat{\mathbb{Z}} \to L_0$ ($u \mapsto \varepsilon_0^{-u}$) est injectif, donc $\mu$ est déterminé par $\xi, \eta$, et la relation $\xi = \eta \varepsilon_1^{\mu-1}$ s'écrit aussi (sans faire intervenir $\mu$)

(69) $$ \eta^{-1} \xi \in L_1 $$.

Soit d'autre part

(70) $$ S\mathcal{G}_{p,\sigma} = \mathcal{G} \times \mathcal{G} \times \hat{\mathbb{Z}}^* $$
$$ \left\{ (\alpha, \beta, \mu) \mid [\alpha, p] = [\beta, \sigma] \varepsilon_1^{\mu-1} \right\} $$

On a une application

(71) $$ S\mathcal{G}_{p,\sigma} \longrightarrow \Sigma_{p,\sigma} $$
$$ (\alpha, \beta, \mu) \longmapsto (\xi, \eta, \mu, \epsilon, \epsilon') $$

avec
$$ \begin{cases} \xi = [\alpha, p] \\ \eta = [\beta, \sigma] \\ \mu \text{ le même} \\ \epsilon' = \varepsilon_{\mathcal{G}}(\alpha) \\ \epsilon = \varepsilon_{\mathcal{G}}(\beta) \end{cases} $$

[MARGIN: NB $\xi = \alpha p \alpha^{-1} \in \mathcal{G}^+$, car on a $\varepsilon_{\mathcal{G}}(\xi) = \varepsilon(p) = 1$ car $p$ opère trivialement sur $\mathcal{G}/\mathcal{G}^+ \simeq \{\pm 1\}$, puisque $p \in \mathcal{G}^+$. $\eta = \beta \sigma \beta^{-1} \in \mathcal{G}^+$]

qui est évidemment surjective. 

On définit d'autre part

\newpage

%% ===== Page 566 =====

(72)
$$ \Sigma_{\rho,\sigma}^{eff} \subset \Sigma_{\rho,\sigma} $$
$$ = \left\{ (\xi,\eta,\mu,\varepsilon',\varepsilon) \in \Sigma_{\rho,\sigma} \left| \exists u \in \operatorname{Aut}(\hat{G}^+) \text{ avec } \begin{cases} u(\rho) = \xi \rho \\ u(\sigma) = \eta \sigma \end{cases} \text{ et induisant l'autom. } \mu \text{ de } \hat{\mathbb{Z}} \text{ et } [unclear] \right. \right\} $$

[MARGIN: i.e. $[u,\rho]=\xi , [u,\sigma]=\eta$]

(cet $u$ sera alors unique). Il y a correspondance
1-1 entre les éléments de $\Sigma_{\rho,\sigma}^{eff}$ et les
systèmes $(u,\mu,\varepsilon',\varepsilon) \in \operatorname{Aut}(\hat{G}^+) \times \hat{\mathbb{Z}}^* \times \{\pm 1\} \times \{\pm 1\}$,
tels que l'on ait 
$\begin{cases} u(\rho) \text{ est } \hat{G}^+\text{-conjugué à } \rho^{\varepsilon'} \\ u(\sigma) \qquad \text{"} \qquad \qquad \text{ à } \sigma^{\varepsilon} \\ u(\varepsilon_0) = \varepsilon_0^\mu \quad (\text{qui exprime } \xi = \eta l_1^V) \end{cases}$

— or cet ensemble n'est autre que le
sous-groupe $\mathbb{B}_{\rho,\sigma}$ de $\operatorname{Aut}(\hat{G}^+) \times \hat{\mathbb{Z}}^* \times \{\pm 1\} \times \{\pm 1\}$ ! (cf (6) p. 577 !)
On trouve ainsi une bijection

(73)
$$ \mathbb{B}_{\rho,\sigma} \overset{\sim}{\longrightarrow} \Sigma_{\rho,\sigma}^{eff} $$
$$ u \longmapsto (\xi,\eta,\mu,\varepsilon',\varepsilon) \text{ avec } \begin{cases} \xi = [u,\rho] \\ \eta = [u,\sigma] \\ \mu = \mu(u) \\ \varepsilon' = \varepsilon'(u) \\ \varepsilon = \varepsilon(u) \end{cases} $$

On désigne par $SG_{\rho,\sigma}^{eff}$ l'image inverse de $\Sigma_{\rho,\sigma}^{eff}$ dans $SG_{\rho,\sigma}$ par (71), donc on aura

(74)
$$ SG_{\rho,\sigma}^{eff} \subset SG_{\rho,\sigma} $$
$$ = \left\{ (\alpha,\beta,\mu) \in \hat{G} \times \hat{G} \times \hat{\mathbb{Z}}^* \left| \begin{matrix} [\alpha,\rho] = [\beta,\sigma] \varepsilon_1^{\mu-1} \\ \exists u \in \operatorname{Aut}(\hat{G}^+) \text{ tel qu'on ait} \\ \begin{cases} u(\rho) = ad(\alpha)(\rho) & (= \xi \rho, \text{ avec } \xi = [\alpha,\rho]) \\ u(\sigma) = ad(\beta)(\sigma) & (= \eta \sigma, \text{ avec } \eta = [\beta,\sigma]) \end{cases} \end{matrix} \right. \right\} $$

de sorte qu'on a un hom. surjectif induit par (71)

(75)
$$ SG_{\rho,\sigma}^{eff} \longrightarrow \Sigma_{\rho,\sigma}^{eff} $$

\newpage

%% ===== Page 567 =====

587

On a une application canonique

$$ (76) \quad S_0\Gamma_{\rho,\sigma} = \underline{\Pi}_{\rho,\sigma} \times_{\Gamma_{\rho,\sigma}} N_{\rho}^* \times_{\Gamma_{\rho,\sigma}} N_{\sigma}^* \longrightarrow SG_{\rho,\sigma}^{eff} $$
$$ x=(u, a, b) \longmapsto (\alpha, \beta, \mu) \quad \text{avec} \quad \left\{ \begin{matrix} \alpha = u a^{-1} \tau'^{\nu'(x)} \in Z_\rho \\ \beta = u b^{-1} \tau^{\nu(x)} \in Z_\sigma \\ \mu = \chi(x) \end{matrix} \right. $$

[MARGIN: pour être plus concret, il faudrait écrire $\alpha = u i_\rho(a)^{-1} \tau'^{\nu'(x)}$, $\beta = u i_\sigma(b)^{-1} \tau^{\nu(x)}$]

où on a défini $\nu'(x), \nu(x) \in \mathbb{Z}/2\mathbb{Z}$ par
$$ (-1)^{\nu'(x)} = \varepsilon'(x), \quad (-1)^{\nu(x)} = \varepsilon(x) $$

Par définition du $S_0\Gamma_{\rho,\sigma}$ comme produit fibré, on a en effet $u a^{-1}, u b^{-1} \in \widehat{C}^+$, donc $\alpha, \beta \in \widehat{C}$.

[MARGIN: NB On a d'ailleurs le fait que pour $a \in N_\rho$ [unclear] $\alpha \in Z_\rho$ [unclear] Ainsi pour [unclear]]

on a d'ailleurs
$$ \alpha(\rho) = u \big( a^{-1} \big( \underset{\rho^{\varepsilon'(x)}}{\tau'^{\nu'(x)}(\rho)} \big) \big) = u(\rho) $$
et de même
$$ \beta(\sigma) = u(\sigma) $$
i.e. $\quad [\alpha,\rho] = [u,\rho]$
$\quad [\beta,\sigma] = [u,\sigma]$

de sorte que d'ailleurs la condition $\underset{[u,\rho]}{[\alpha,\rho]} = \underset{[u,\sigma]}{[\beta,\sigma]} \dots$ ne fait qu'exprimer la condition que $u(\varepsilon_0) = \varepsilon_0^\mu$ qui intervient dans la définition de $Z_{\rho,\sigma}$.

En fait l'application envisagée est injective, car [unclear] de voir que connaissant $(\alpha,\beta,\mu)$ on a

$$ (77) \quad \nu'(x) = \nu_{\widehat{C}}(\alpha), \quad \nu(x) = \nu_{\widehat{C}}(\beta), \quad \chi(x) = \chi(\mu) = \mu $$

d'où on tire $a = \tau'^{\nu'(x)} \alpha^{-1} u$, $\quad b = \tau^{\nu(x)} \beta^{-1} u$

de sorte si $u$ est connu, on connait également $a$ et $b$ donc $x$. D'autre part pour déterminer $u \in Z_{\rho,\sigma} \subset \operatorname{Aut}(\widehat{C}^+) \times \hat{\mathbb{Z}}^\times \times \{\pm 1\}$ il suffit en vertu de (78) de le déterminer entant qu'élément de $\operatorname{Aut}(\widehat{C}^+)$ qu'il définit, ce qui est [unclear].

\newpage

%% ===== Page 568 =====

immédiat par (77).

(79) $$S\Gamma_{\rho,\sigma} \hookrightarrow SG_{\rho,\sigma}^\#$$

il reste à voir que cette application est surjective.
On a fait ce qu'il fallait pour. Notons
d'abord que le composé
$$S\Gamma_{\rho,\sigma} \hookrightarrow SG_{\rho,\sigma}^\# \xrightarrow{eff \; it^c} \Sigma_{\rho,\sigma}^\# (\xrightarrow{\sim} \underline{Z}_{\rho,\sigma})$$
n'est autre, par (77) et (71), que

(80) $x=(u,a,b) \longmapsto (\xi, y, \mu, \varepsilon', \varepsilon) \quad \begin{cases} \xi = [u,\rho] \text{ avec } u(\rho)=\xi \rho \\ y = [u,\sigma] \text{ avec } u(\sigma)=y \sigma \\ \mu = \chi_{S\Gamma}(x) \\ \varepsilon' = \varepsilon'(x) \\ \varepsilon = \varepsilon(x) \end{cases}$

donc son image composé avec $\Sigma_{\rho,\sigma}^\# \xrightarrow{\sim} Z_{\rho,\sigma}$
n'est autre que la projection canonique
$$S\Gamma_{\rho,\sigma} \longrightarrow Z_{\rho,\sigma}$$
$$(u,a,b) \longmapsto u$$
qui est surjective, de noyau $SN_{\rho}^* \times SN_{\sigma}^* = SZ_{\rho}^* \times SZ_{\sigma}^*$.
Considérons d'autre part l'opération

(81) $\begin{cases} \text{opération de } SN_{\rho}^* \times SN_{\sigma}^* \\ \text{sur } SG_{\rho,\sigma} \subset \hat{C}^* \times \hat{C}^* \times \hat{\mathbb{Z}}^* \\ \{(\alpha,\beta,\mu) \mid [\alpha,\rho] = [\beta,\sigma] \varepsilon_1^{\dots}\} \end{cases} \left| \begin{cases} (a_0,b_0) \in SN_{\rho}^* \times SN_{\sigma}^* \\ \text{opère par} \\ (\alpha,\beta,\mu) \longmapsto (\alpha a_0, \beta b_0, \mu) \end{cases} \right.$

Il est clair qu'il opère fidèlement librement, et
que le quotient de $SG_{\rho,\sigma}$ par cette action
s'identifie à $Z_{\rho,\sigma}$, donc on a le résultat
identique pour $S\Gamma_{\rho,\sigma}$ et $\Sigma_{\rho,\sigma}^\#$, ...
(82) [unclear: equation involving $S\Gamma_{\rho,\sigma}$, $\Sigma_{\rho,\sigma}^\#$ and action]

\newpage

%% ===== Page 569 =====

[unclear: crossed out text at top]

(82) $$ \varepsilon = (-1)^v, \quad \varepsilon' = (-1)^{v'} $$

fixés, l'opération précédente en $(a_0, b_0)$ correspond,
par l'application (76), à l'opération

(83) $$ (u, a, b) \longmapsto (u, \tau^u(a_0^{-1})a, \tau^u(b_0^{-1})b) $$

i.e. à la translation à gauche par
$$(\tau^u, \tau^u)(a_0^{-1}, b_0^{-1}) = (a_0, b_0)^{-1}$$
D'où la surjectivité, et l'
isomorphisme cherché

(84) $$ S_0 \Gamma_{\rho, \sigma} \overset{\sim}{\longrightarrow} S \mathcal{G}_{\rho, \sigma}^{eff} \quad (cf \ 75) $$

La translation à gauche par $l \in L_0$ correspond, [MARGIN: dans $S_0 \Gamma_{\rho, \sigma}$]
par la bijection précédente, à l'opération

(85) $$ (\alpha, \beta, \mu) \longmapsto (l \alpha, l \beta, \mu) $$

dont il était clair a priori qu'elle stabilise
$S \mathcal{G}_{\rho, \sigma}^{eff}$. Passant au quotient, on trouve
donc

(86) $$ \Gamma_{\rho, \sigma} \overset{\sim}{\longrightarrow} L_0 \backslash S \mathcal{G}_{\rho, \sigma}^{eff} . $$

Enfin, considérons

(87) $$ \begin{cases} \begin{aligned} S_0 M_{\rho, \sigma} &\longrightarrow S \mathcal{G}_{\rho, \sigma}^{eff} \times \underline{\hat{C}}^+ = \left\{ (\alpha, \beta, \mu, f) \in \hat{C} \times \hat{C} \times \hat{\mathbb{Z}} \times \hat{C} \mid [\alpha, \rho] = [\beta, \sigma] \varepsilon^{\mu-1} \right\} \\ x = (u, a, b, U) &\longmapsto (\alpha, \beta, \mu, f) \quad \text{avec} \quad \begin{cases} \alpha = u a^{-1} \tau^u(a) \\ \beta = u b^{-1} \tau^u(b) \\ \mu = \chi_{SM}(x) \quad (= \chi_\rho(u) = f_\rho(a) = f_\sigma(b) = \chi_\sigma(u)) \\ f = u U^{-1} \end{cases} \end{aligned} \end{cases} $$

Cette application est bijective (trivial par (84)).

\newpage

%% ===== Page 570 =====

7) Fonctorialités

Considérons
(88) $\hat{G} \twoheadrightarrow \hat{G}'$ hom. surjectif de groupes profinis

soient $\rho', \sigma', \tau', \dots$ les images de $\rho, \sigma, \tau$,
$\tau' = \tau_{\rho'\sigma'}$ dans $\hat{G}'$.

Prop. Conditions équivalentes sur $\varphi$ :

a) L'application induite par $\varphi$
(89) $\Sigma_{\rho,\sigma} \to \Sigma_{\rho',\sigma'}$
envoie $\Sigma_{\rho,\sigma}^{eff}$ dans $\Sigma_{\rho',\sigma'}^{eff}$

b) Pour tout $u \in Aut \hat{C}_{>}^+$ qui normalise $L_0$,
et tel que $u(L_\rho) \underset{\hat{C}^+}{\sim} L_\rho$, $u(L_\sigma) \underset{\hat{C}^+}{\sim} L_\sigma$,
il existe $u' \in Aut (\hat{C}'^+_{>})$ qui rend commutatif
(90)
\begin{tikzcd}
\hat{C}^+ \ar[r, "\varphi"] \ar[d, "u"'] & \hat{C}'^+ \ar[d, "u'"] \\
\hat{C}^+ \ar[r, "\varphi"'] & \hat{C}'^+
\end{tikzcd}
en d'autres termes, $u$ respecte $Ker \varphi$.

b') idem, en supposant seulement, à la place que
$u$ normalise $L_0$, que $u(L_0) \underset{\hat{C}^+}{\sim} L_0$ (plus les autres conditions)

c) $\exists$ hom. de groupes
(91) $\varphi_Z : Z_{\rho,\sigma} \to Z_{\rho',\sigma'}$
compatible avec $j, j_\rho, j_\sigma$ et tel que $\forall u \in Z_{\rho,\sigma}$, le
$u'$ qui lui est associé soit comme dans (90).

[MARGIN: alors il induit $\Gamma_{\rho,\sigma} \to \Gamma_{\rho',\sigma'}$] (NB $\varphi_Z$ uniquement déterminé par $\varphi$, car est compatible avec $j : L_0 \to Z$)

c') $\exists$ hom de groupes
(92) $\varphi_{\hat{G}} : G_{\rho,\sigma} \to G_{\rho',\sigma'}$
compatible avec $j_a, j_b, j'_c$ et satisfaisant à la
condition de compatibilité (90). (NB $\varphi_{\hat{G}}$ a [unclear])

\newpage

%% ===== Page 571 =====

589

est alors unique, et compatible avec $i_{Z, \underline{G}}$, $i_{\underline{G}, S}$ [unclear]
[crossed out] $\underline{N}_{\underline{P}}$ dans $\underline{N}'_{\underline{P}}$, $\underline{N}_{\underline{\sigma}}$ dans $\underline{N}'_{\underline{\sigma}}$, qui se prolonge
en $\varphi_{\underline{P}} : \underline{N}_{\underline{P}}^\sim \to \underline{N}'^{\sim}_{\underline{P}}$ (remarque $\underline{N}_{\underline{P}}^\sim = \underline{N}_{\underline{P}} \rtimes \dots$ et se
prolonge en $\varphi_{\underline{\sigma}} : \underline{N}_{\underline{\sigma}}^{\sim \underline{G}} \to \underline{N}'^{\sim \underline{G}'}_{\underline{\sigma}}$ (remarque $\underline{N}_{\underline{\sigma}}^{\sim \underline{G}} = \underline{N}_{\underline{\sigma}} \rtimes \dots$))
[crossed out]
On dira alors que $\varphi$ est "prolongeable".
Corollaire On peut alors prolonger $\varphi$ en
(93) $\varphi_{\mathcal{M}} : S_0 \underline{G}_{\underline{P}, \underline{\sigma}} \longrightarrow S_0 \underline{G}'_{\underline{P}, \underline{\sigma}}$ ,
compatible avec $\theta_Z, \theta_P, \theta_\sigma, \theta_G$
compte tenu de $\varphi_Z, \varphi_{\underline{P}}, \varphi_{\underline{\sigma}}, \varphi_{\underline{G}}$, et ceci
de façon "compatible avec tout", notamment
avec les diagrammes (45) (62).

Exemple Supposons $\underline{\Sigma}_{\underline{P}, \underline{\sigma}}^{\underline{G}'} = \underline{\Sigma}_{\underline{P}, \underline{\sigma}}^{\underline{G}\flat}$ , alors
les conditions de la prop. sont automatiques
satisfaites (quel que soit l'hom. surjectif
$\underline{G} \to \underline{G}'$).

[MARGIN: Il faut [unclear] se donner $\tau_{\underline{Z}}^\natural \in Z_{\underline{P}, \underline{\sigma}}^\natural$] Rmq. Il se peut qu'on se donne un
sous-groupe $Z_{\underline{P}, \underline{\sigma}}^\natural \subset Z_{\underline{P}, \underline{\sigma}}$, et qu'on
ne demande que la condition a) de la prop.
seulement pour des $u \in Z_{\underline{P}, \underline{\sigma}}^\natural$. Mais à l'aide
de $Z_{\underline{P}, \underline{\sigma}}^\natural$, on construit
(94) $\underline{G}_{\underline{P}, \underline{\sigma}}^\natural = Z_{\underline{P}, \underline{\sigma}}^\natural . \underline{G}^+$, $Z_{\underline{P}}^\natural = Z_{\underline{P}} \cap \underline{G}_{\underline{P}, \underline{\sigma}}^\natural$, $Z_{\underline{\sigma}}^\natural = Z_{\underline{\sigma}} \cap \underline{G}_{\underline{P}, \underline{\sigma}}^\natural$
des $\underline{N}_{\underline{P}}^\natural$ et $\underline{N}_{\underline{P}}^{*\natural}$, $\underline{N}_{\underline{\sigma}}^\natural$ et $\underline{N}_{\underline{\sigma}}^{*\natural}$, $\Gamma_{\underline{P}, \underline{\sigma}}^\natural$, puis
$S_0 \underline{M}_{\underline{P}}^\natural$ et $\underline{M}_{\underline{P}}^\natural$. On a une variante évidente de la
[crossed out] ... variante évidente de la

\newpage

%% ===== Page 572 =====

proposition, [unclear: plus ou moins] des conditions de
prolong. de $\varphi$ de $\underline{S}_0 \underline{M}_{\rho,\sigma}^\natural \longrightarrow \underline{S}_0 \underline{M}_{\rho,\sigma}^{\underline{G}'}$.

Cas universel $\underline{G} = \widehat{Gl(2,\mathbb{Z})}^\wedge$, $\rho = \begin{pmatrix} 0 & 1 \\ -1 & 1 \end{pmatrix}$, $\sigma = \begin{pmatrix} 0 & -1 \\ 1 & 0 \end{pmatrix}$
$\tau' = \begin{pmatrix} 0 & 1 \\ 1 & 0 \end{pmatrix}$, $(\tau = \begin{pmatrix} -1 & 0 \\ 0 & +1 \end{pmatrix}$, $\xi_0 = \begin{pmatrix} 1 & -1 \\ 0 & 1 \end{pmatrix}$,

[unclear: crossed out math] ~~$\xi_1 = \begin{pmatrix} 1 & 0 \\ 1 & 1 \end{pmatrix}$.~~

Comme $\rho$

\newpage

%% ===== Page 573 =====

8°) Le cas universel $\hat{G} = \operatorname{Gl}(2,\mathbb{Z})^\wedge$.  (590)

Considérons la suite exacte

(1) 
\begin{tikzcd}
1 \ar[r] & \{\pm 1\} \ar[r] & \mathcal{N}_{1,1}^\sim \ar[r] \ar[d, equal] & \mathcal{M}_{0,3}^\sim \ar[r] & 1 \\
& & \mathcal{M} & &
\end{tikzcd}

[MARGIN:
$\hat{G} = \hat{G}_{1,1} \simeq \operatorname{Gl}(2,\mathbb{Z})^\wedge$
$\hat{G}^+ = \hat{\Gamma}_{1,1} \simeq \operatorname{Sl}(2,\mathbb{Z})^\wedge$
]

$\{\pm 1\} = Centr(\mathcal{M}) = Centr(\hat{G}^+)$

d'où par (1) une application

(3) $\mathcal{M}_{0,3}^\sim \to \operatorname{Aut}(\mathcal{M}) \xrightarrow{\text{restriction à } \hat{G}^+} \operatorname{Aut}(\hat{G}^+)$

D'ailleurs, dans le cas [unclear] de $\hat{G}^+$ [unclear] 
[unclear] un centre, [unclear: par suite] d'où

(4) $\operatorname{Aut}(\hat{G}^+) \to \operatorname{Aut}(\hat{G}^+ / \pm 1 \simeq \hat{G}_0^+)$

et le composé de (3) avec (4) est l'homomorphisme canonique :

(5)
\begin{tikzcd}
\mathcal{M}_{0,3}^\sim \ar[r] \ar[dr] & \operatorname{Aut}(\hat{G}^+) \ar[d] \\
& \operatorname{Aut}(\hat{G}_0^+ = \hat{G}^+ / \pm 1)
\end{tikzcd}

Cela montre que (3) est injectif, puisque $\mathcal{M}_{0,3}^\sim \to \operatorname{Aut}(\hat{G}_0^+)$ l'est. L'application

Remarque L'application $\operatorname{Aut}(\hat{G}^+) \to \operatorname{Aut}(\hat{G}_0^+)$ n'est pas injective par contre, et les éléments du noyau correspondent aux

hom. $\hat{G}^+ \to Centr(\hat{G}^+) = \{\pm 1\}$

\newpage

%% ===== Page 574 =====

ou encore aux hom. de

(6) $$ \widehat{C}^+_{ab} \longrightarrow \{\pm 1\} \quad \text{ou encore} \quad (\widehat{C}^+_{ab})_2 \longrightarrow \{\pm 1\} $$

Or on voit que

$$ \widehat{C}^+_{ab} = \{ \rho, \sigma \mid \rho^6 = \sigma^4 = [\rho, \sigma] = 1, \rho^3 = \sigma^2 \} = \mathbb{Z}/12\mathbb{Z} $$
[MARGIN: $\rho \mapsto 2 \bmod 12$, $\sigma \mapsto 3 \bmod 12$]

donc

(7) $$ \left\{ \begin{aligned} (\widehat{C}^+_{ab})_2 &\xrightarrow{\sim} \mathbb{Z}/2\mathbb{Z} \\ \rho &\longmapsto 0 \\ \sigma &\longmapsto 1 \end{aligned} \right. $$

Cet homomorphisme correspond à l'hom.
$\widehat{C}^+ \longrightarrow \{\pm 1\}$ qui s'obtient aussi comme
composé :

(8)
\begin{tikzcd}
\widehat{C}^+ \ar[r] & \mathfrak{S}_3 = \widehat{C}^+/\widehat{\pi} \ar[r, "sgn"] & \{\pm 1\}
\end{tikzcd}
[MARGIN: sous $\mathfrak{S}_3$ : groupe symétrique à trois lettres. sous $\widehat{\pi}$ : im. inv. de $\widehat{\Pi}_{0,3}$ dans $\widehat{C}^+$]

Ainsi le seul élément $\Theta$ différent de $1$
du noyau de (4) est donné par

(9) $$ \Theta : x \longmapsto x \cdot \text{sgn}(x) $$

On notera, en posant

(10) $$ \widehat{L}_0 = \varepsilon_0^{\hat{\mathbb{Z}}}, \quad \widehat{L}''_0 = \widehat{L}_0 \cap \widehat{\pi} = \widehat{L}_0^2 $$

que le s-groupe $\Theta(\widehat{L}_0)$ de $\widehat{C}^+$
n'est pas conjugué à $\widehat{L}_0$. En effet

(11) $$ \Theta(\varepsilon_0) = -\varepsilon_0 \quad (\text{car } \varepsilon_0 = \sigma \rho, \text{ sgn}(\varepsilon_0) = -1) $$

et on vérifie que $\varepsilon_0$ n'est pas conjugué à
$(-\varepsilon_0)^k = \tau_0 \varepsilon_0^k$ pour quelque $k \in \hat{\mathbb{Z}}^*$ (car on voit

\newpage

%% ===== Page 575 =====

591

Il s'ensuit, en passant à $\hat{\mathfrak{S}}/\Pi^\wedge$
[unclear crossed out: $\hat{\mathfrak{S}}_{0,3} / \Pi_{0,3}$] et notant que $l_0 = \varepsilon_0^2$ serait
conjugué : $f \varepsilon_0 f^{-1} = \varepsilon_0^{-1}$, ce qui est exclu
($l_0 \neq 1$, or $\varepsilon_0$ et $-\varepsilon_0$ ne sont pas conjugués
même dans $\operatorname{Sl}(2, \hat{\mathbb{Z}})$! ). On va
s'intéresser justement aux élts $f \in \hat{\mathfrak{S}}^{+\wedge}$
qui transforment $L_0^\pi$ en un sous-groupe
conjugué. On notera que par contre
(12) $\theta(L_0) = L_0$ donc $\theta(L_0^\pi) = L_0^\pi$

Proposition L'hom. (3)
$$ \tilde{M}_{0,3} \longrightarrow \operatorname{Aut}(\hat{\mathfrak{S}}^+) $$
est injectif, et son image est formée des autom de $\hat{\mathfrak{S}}^{+\wedge}$ qui satisfont les deux
conditions suivantes :
a) $u$ stabilise $\Pi^\wedge = \hat{\mathfrak{S}}^{+\wedge}[2] = \operatorname{Ker}(\hat{\mathfrak{S}}^{+\wedge} \to \operatorname{Sl}(2, \mathbb{Z}/2\mathbb{Z}))$
b) $u$ transforme $L_0^\pi$ en un sous-groupe de $\hat{\mathfrak{S}}^{+\wedge}$
$\hat{\mathfrak{S}}^{+\wedge}$-conjugué de $L_0^\pi$ i.e. $\exists \mu \in \hat{\mathbb{Z}}^*$ et $f \in \hat{\mathfrak{S}}^{+\wedge}$
tels qu'on ait
(13) $u(\varepsilon_0) = int(f)(\varepsilon_0^\mu)$

Dém. L'injectivité a déjà été vue [unclear crossed out]
que les conditions a) b) sont nécessaires pour
qu'un $u$ [unclear crossed out] provienne d'un $u_0 \in \tilde{M}_{0,3}$, Pour a)
c'est trivial, puisque $u_0$ stabilise $\Pi_{0,3}^\wedge$, donc son

\newpage

%% ===== Page 576 =====

image inverse dans $\widehat{\mathfrak{S}}^{+\wedge}$. Pour la condition
b), on voit qu'on se ramène au cas où $u_0 \in M_{0,3}^\sim[0]$
(quitte à conjuguer par un élément de $\widehat{\mathfrak{S}}_{0,3}^\sim \simeq \widehat{\mathfrak{S}}^{+\wedge} / \pm 1$)
dans ce cas on a par construction de $M_{1,1}^\sim$ à partir
de $M_{0,3}^\sim$ (cf. III, p. 665...)
(14) $$u_0(\varepsilon_0) = \varepsilon_0^\mu$$
où $\mu \in \hat{\mathbb{Z}}^\times$ est le multiplicateur, OK.

Inversement, soit $u$ un autom. de $\widehat{\mathfrak{S}}^{+\wedge}$ qui
satisfait a), b), considérons l'autom. $\tilde{u}$ de
$\widehat{\mathfrak{S}}^{+\wedge} / \pm 1 \simeq \widehat{\mathfrak{S}}_{0,3}^\sim$ qu'il induit. On voit donc
que $\tilde{u}(\widehat{\Pi}_{0,3}) = \widehat{\Pi}_{0,3}$, et que $\tilde{u}(\varepsilon_0)$ est $\widehat{\mathfrak{S}}_{0,3}^\sim$-con-
jugué à un $\varepsilon_0^\mu$ ($\mu \in \hat{\mathbb{Z}}^\times$). Donc par
définition, on a $\tilde{u} \in M_{0,3}^\sim$ (cf. page 390 ff).
Soit $u'$ l'autom de $\widehat{\mathfrak{S}}^{+\wedge}$ induit par $\tilde{u}$
(via B)) - alors $u'$ et $u$ induisent le même
autom. de $\widehat{\mathfrak{S}}_{0,3}^\sim = \widehat{\mathfrak{S}}^{+\wedge} / \pm 1$, donc pour
on a $u' = u \circ \Theta$, où $\Theta$ [unclear: $= u^{-1}u'$ written below] est un autom. qui
1) induit l'identité sur $\widehat{\mathfrak{S}}^{+\wedge} / \pm 1$, et
2) transforme $L_0^\pi$ en un sous-groupe conjugué -
on a vu que cela implique $\Theta = 1$, cqfd.

[crossed out text]
Considérons maintenant la suite
exacte
(15) $$1 \longrightarrow \pm 1 \longrightarrow \Pi^\wedge \longrightarrow \widehat{\Pi}_{0,3} \longrightarrow 0$$
[under $\Pi^\wedge$: $= \widehat{\mathfrak{S}}^{+\wedge}[2]$]

\newpage

%% ===== Page 577 =====

592

et considérons le scindage donné par

$$ (16) \quad \begin{cases} \Pi_{0,3} \longrightarrow \pi = \widehat{S}^+[2] = \operatorname{Ker}(\operatorname{Sl}(2,\mathbb{Z}) \to \operatorname{Sl}(2,\mathbb{Z}/2)) \\ l_i \longmapsto \lambda_i \overset{def}{=} -l_i \end{cases} $$

Il y a un tel hom., en vertu de

$$ (17) \quad \begin{cases} (-l_0)(-l_1)(-l_\infty) = 1 \quad dans \quad \operatorname{Sl}(2,\mathbb{Z}) \quad i.e. \\ l_\infty l_1 l_0 = -1 \end{cases} $$

et il est évident que c'est un scindage. On en déduit un scindage pour $\widehat{\Pi} \to \widehat{\Pi}_{0,3}$.

On désigne par $\Pi_0$ l'image de (16), i.e.

$$ (18) \quad \begin{cases} \Pi_0 = ss\text{-groupe de } \pi = \widehat{S}^+[2] \text{ engendré par} \\ \lambda_0, \lambda_1, \lambda_\infty \quad (\text{où } \lambda_i = -l_i) \end{cases} $$

et on a donc un ss-groupe $\widehat{\Pi}_0$ de $\widehat{\Pi}$.

Proposition a) Le sous-groupe $\Pi_0$ de $\pi$ est distingué d'indice 2, et il est distingué dans $\widehat{S}^+$.

b) Le ss-groupe $\widehat{\Pi}_0$ de $\widehat{\Pi}$ est distingué d'indice 2 et il est distingué dans $M$ - ou, ce qui revient au même, stable sous l'action de $M/\pm 1 \simeq \mathcal{M}_{0,3}^\sim$.

Dém. Comme

$$ (19) \quad \begin{cases} \rho(\lambda_i) = \lambda_{i+1} & \sigma(\lambda_\infty) = \sigma(\lambda_1 \lambda_0)^{-1} = \lambda_1^{-1} \lambda_0^{-1} \\ \sigma(\lambda_0) = \lambda_1^{-1}, \sigma(\lambda_1) = \lambda_0^{-1} & [\text{unclear crossed out text}] \\ \tau(\lambda_0) = \lambda_0^{-1}, \tau(\lambda_1) = \lambda_1^{-1} & \tau(\lambda_\infty) = \tau(\lambda_1 \lambda_0)^{-1} = \lambda_0 \lambda_1 \end{cases} $$

on voit que $\Pi_0$ est normalisé par $\rho, \sigma, \tau$ donc par $\widehat{S}$. Qu'il soit d'indice 2 est évident. Donc $\widehat{\Pi}_0$ est d'indice 2 dans $\widehat{\Pi}$ et normalisé par $\widehat{S}$. Pour voir qu'il est

\newpage

%% ===== Page 578 =====

normalisé par $M$ tout entier, comme

$$M = M(\omega) . \widehat{\mathbb{S}}^{+ \wedge}$$

il revient à vérifier qu'il est normalisé
par $M(0) \simeq M_{0,3}(0)$. Or on a, pour

$$u \in M_{0,3}(0)$$

(20) $$ \begin{cases} u(\rho) = int(\alpha) \rho \\ u(\sigma) = \varepsilon_2 int(\beta) \sigma \\ u(\varepsilon_0) = \varepsilon_0^k \end{cases} \qquad \begin{cases} \varepsilon_2 = \varepsilon(u) \\ \alpha, \beta \in \widehat{\Pi} . \{1, \tau\} \\ \beta \in \widehat{\pi} \end{cases} $$

d'où

(21) $$ \begin{cases} u(l_0) = l_0^k \\ u(-l_0) = (-l_0)^k \text{ i.e. } u(\lambda_0) = \lambda_0^k \quad (\text{car } (-1)^k = -1) \end{cases} $$

on a $u(\lambda_0) \in \widehat{\Pi}_0^\wedge$, il reste à prouver que

$u(\lambda_1) \in \widehat{\Pi}_0^\wedge$, or

$$u(\lambda_1) = u(\rho(\lambda_0)) = (u\rho)(\lambda_0) = int(\alpha)(\rho) (\underbrace{u(\lambda_0)}_{\lambda_0^k}) = int(\alpha)(\lambda_1^k)$$

et comme $\widehat{\Pi}_0^\wedge$ est normalisé par $\widehat{\mathbb{S}}^\wedge$ et en
particulier par $\alpha$, on trouve bien $u(\lambda_1) \in \widehat{\Pi}_0^\wedge$.

Corollaire. On a

(22) $$ \begin{cases} \widehat{\Pi} \simeq \widehat{\Pi}_0 \times \mu \\ \widehat{\overline{\Pi}} \simeq \widehat{\Pi}_0^\wedge \times \mu \end{cases} $$

où $\mu = \pm 1_3 = Center(\widehat{\mathbb{S}}^{+\wedge})$

C'est trivial, puisqu'on a un scindage de
l'extension centrale (15) dans $\widehat{\Pi}_0$.

\newpage

%% ===== Page 579 =====

Soient
$$
(23) \quad
\begin{cases}
\mu_\tau = \{1, \tau\} \subset \hat{\mathbb{S}} & D_\tau = \mu_\tau \times \mu \subset \hat{\mathbb{S}} \\
\mu_\sigma = \{1, \sigma\} \subset \hat{\mathbb{S}} & D_\sigma = \mu_\sigma \times \mu \subset \hat{\mathbb{S}}
\end{cases}
$$

de sorte que l'on a

$$
(24) \quad
\begin{cases}
D_\tau \cap \hat{\mathbb{S}} = D_\tau \cap \pi = \mu & D_\tau \cap \pi_0 = \{1\} \\
D_\sigma \cap \hat{\mathbb{S}} = D_\sigma \cap \pi = \mu & D_\sigma \cap \pi_0 = \{1\}
\end{cases}
$$

[MARGIN: idem en $\hat{S}$, $\hat{\pi}$ remplaçant $\mathbb{S}^\wedge$, $\pi$]

On pose

$$
(25) \quad
\begin{cases}
\pi_\tau = \mu_\tau . \pi = D_\tau . \pi_0 \ , \ \pi_{0, \tau} = \mu_\tau . \pi_0 \\
\pi_\sigma = \mu_\sigma . \pi = D_\sigma . \pi_0 \ , \ \pi_{0, \sigma} = \mu_\sigma . \pi_0
\end{cases}
$$

de sorte que l'on a

$$
(26) \quad
\begin{cases}
\pi_\tau = \pi_{0, \tau} \times \mu \\
\pi_\sigma = \pi_{0, \sigma} \times \mu
\end{cases}
$$

et idem en mettant des $^\wedge$.

Normalisation de $\alpha$, $\beta$
Les relations 20) ne déterminent $\alpha$, $\beta$ qu'a "modulo signe". On lève cette ambiguïté en prenant
$$
(27) \quad \alpha \in \pi_{0, \tau} \ , \ \beta \in \pi_{0}
$$

Le groupe $\hat{\mathbb{S}} / \pi_0$. Il est décrit par
gén. $\rho, \sigma, \tau$ et relations

$$
(28) \quad
\begin{cases}
\rho^6 = \sigma^4 = \tau^2 = 1 \ , \ \rho^3 = \sigma^2 \\
\tau(\sigma) = \sigma^{-1} \\
\tau(\rho) = \rho
\end{cases}
$$

$$
(29) \quad \tau(\rho) = (-l_1^{-1}) \rho
$$
donc la dernière relation (28), s'écrit $-l_1^{-1} = 1$ i.e. $l_1 = -1$
(et implique $\lambda_0 = \bar{\lambda}_0 = -1$ en [unclear]).

\newpage

%% ===== Page 580 =====

En termes de $\rho, \sigma, \tau'$ on trouve [unclear] $=$
(30)
$$
\left\{ \begin{array}{l}
\rho^6 = \sigma^4 = \tau'^2 = 1 \quad , \quad \rho^3 = \sigma^2 \\
\tau'(\rho) = \rho^{-1}, \tau'(\sigma) = \sigma^{-1} \\
\sigma(\rho) = \rho^{-1}
\end{array} \right.
$$

Le sous-groupe $\widehat{\mathfrak{S}}^+/\Pi_0$ a la présentation
en $\rho, \sigma$
(31)
$$
\left\{ \begin{array}{l}
\rho^6 = \sigma^4 = 1 \quad , \quad \rho^3 = \sigma^2 \\
\sigma(\rho) = \rho^{-1}
\end{array} \right.
$$

Quelle est l'opération de $\mathcal{M}_{0,\tilde{3}} \simeq \mathcal{M} / \pm 1$
sur $\widehat{\mathfrak{S}}^+/\Pi_0$ ? [unclear: On a la] structure
déduite de
(32)
$$
1 \to \mu \to \widehat{\mathfrak{S}}^+/\Pi_0 \to \widehat{\mathfrak{S}}^+/\mu \to 1 \\
\quad \parallel \qquad \qquad \qquad \qquad \qquad \qquad \parallel \\
\quad \pm 1 \qquad \qquad \qquad \qquad \qquad \qquad \mathfrak{S}_3
$$

et les aut. de [unclear] en question conservent cette
structure d'extension. Les éléments de
$\mathcal{M}_{0,\tilde{3}} \simeq \mathcal{M} / \pm 1$ opèrent en réduisant à l'
identité sur le quotient $\mathfrak{S}_3$, et bien sûr
aussi sur $\mu$ - or un tel automorphisme
d'extension correspond à un hom.
(33)
$$
\mathfrak{S}_3 \to \mu \quad \text{i.e.} \quad \mathfrak{S}_{3, ab} \to \mu \\
\simeq \qquad \qquad \qquad \qquad \qquad \qquad \\
\mathfrak{S}_3/\mathfrak{A}_3 \xrightarrow{\sim} \{\pm 1\} = \mu
$$

donc le groupe des autom. de $\widehat{\mathfrak{S}}^+/\Pi_0$ indui-
sant l'identité sur $\mathfrak{S}_3$ est $\simeq \mathbb{Z}/2\mathbb{Z}$, le
générateur $\Theta$ est donné par
(34)
$$
\Theta(x) = sg(x) \cdot x \quad \text{donc } \left\{ \begin{array}{l} \Theta(\rho) = \rho \\ \Theta(\sigma) = - \sigma \end{array} \right.
$$

\newpage

%% ===== Page 581 =====

Comparons avec les formules (20), on trouve
$$ \Theta^{\varepsilon_2(u)} : \begin{cases} \rho \mapsto \rho \\ \sigma \mapsto \varepsilon_2(u) \sigma \end{cases} $$
(35)
i.e.
(35 bis) $[\dot{u}, \dot{\sigma}] = \varepsilon_2(u)$ dans $\hat{\mathfrak{S}}^{+\wedge}/\Pi_0$
$\dot{u}(\dot{\sigma})\dot{\sigma}^{-1}$

[MARGIN: 
Cela s'écrit aussi
$[\dot{u}, \dot{x}] = sg(x)$ dans $\hat{\mathfrak{S}}^{+\wedge}/\Pi_0$
(car $\mathcal{M}_{0,3}^\sim$ opère triv. 
sur l'abélianisé
de $\hat{M}$ ... )
]

[unclear: Je voudrais]
Reprenons la définition $\hat{\mathfrak{S}}^{+\wedge}$ :
(36) $Z = Z_{\hat{\mathfrak{S}}, \Gamma} \overset{df}{=} \left\{ u \in \operatorname{Aut}(\hat{\mathfrak{S}}^{+\wedge}) \mid u(\rho) \text{ conjugué à } \rho, u(\sigma) \text{ conjugué à } \sigma, u(l_0) = l_0 \right\}$

de sorte qu'on aura, pour $u \in Z$
(37) $\begin{cases} u(\rho) \text{ } \hat{\mathfrak{S}}\text{-conjugué } : \rho^{\varepsilon'} \\ u(\sigma) \text{ } \hat{\mathfrak{S}}\text{-conjugué } : \sigma^\varepsilon \quad (\varepsilon = \pm 1) \\ u(l_0) = l_0^H \end{cases}$

pour $\varepsilon' \in \hat{\mathbb{Z}}^\times$ et $\varepsilon \in \{\pm 1\}$ bien déterminés,

Car $Z \to \hat{\mathfrak{S}}$ est injectif, (car $\rho$ n'est pas $\hat{\mathfrak{S}}$-conjugué à $\sigma$ ni $\sigma^{-1}$). Je dis qu'un tel $u$ provient de $\mathcal{M}_{0,3}^\sim$ - il reste à vérifier, en vertu de la prop..., que $u$ stabilise $\hat{\Pi}^\wedge$, qui est engendré par $l_0, l_1, l_\infty$.
Or on aura
$$u(l_0) = l_0^H \in \hat{\Pi}^\wedge$$
et
$$u(l_1) = u_\rho(l_0) = \text{int}(f\rho^{\varepsilon'})(u(l_0)) = \text{int}(f\rho^{\varepsilon'})(l_0^H)$$
or $\hat{\mathfrak{S}}^{+\wedge}$ [unclear] donc $(f\rho^{\varepsilon'})$ normalise $\hat{\Pi}^\wedge$,
donc $u(l_1) \in \hat{\Pi}^\wedge$
et de même
$u(-1) = -1$, d'où $u(l_\infty) = - u(l_1 l_0)^{-1} = -u(l_0^{-1})u(l_1)^{-1}$
donc $u(l_\infty) \in \hat{\Pi}^\wedge$, OK.
Si on identifie $u$ à un élément de $\mathcal{M}_{0,3}^\sim$, on voit que c'est en fait un élément de [unclear: la].

\newpage

%% ===== Page 582 =====

[MARGIN: Attention :
suite à ce qui précède, l'intro (et la notation) $\mathcal{M}_{0,3}^\sim[0] \subset \mathcal{M}_{0,3}^\sim$
(et $\mathcal{M}_{0,3}[0] \subset \mathcal{M}_{0,3}$)
n'a de sens que si $\dots$ est $\dots$
Or il est utile de s'en servir pour d'autres $\dots$
par ex. pour l'image de l'homomorphisme de $B_3$
dans $\mathcal{M}_{0,3}^\sim$ $\dots$ (cf. détails p. [unclear])
On va donc corriger le tir plus bas.
Cf. diag p. [unclear] :
$Z_{\rho,\sigma}^\sim = M[0] = \mathcal{M} \cap \mathcal{M}_{0,3}^\sim[0]$ ]

$u \in \mathcal{M}_{0,3}^\sim[0] \simeq \mathcal{M}^![\Gamma_0] \subset \mathcal{M}$

donne (30), de façon plus précise

$u(\rho) = Int(\alpha)\rho \quad \alpha \in \Pi_{0,\hat{Z}}, \beta \in \Pi_0^\wedge$
$u(\sigma) = \varepsilon_2 Int(\beta)\sigma \quad \mu \in \hat{\mathbb{Z}}^*$
$u(\varepsilon_0) = \varepsilon_0^\mu \quad (\varepsilon_1 = \varepsilon_1(\mu), \varepsilon_2(\alpha) = \varepsilon_3(\mu))$

qui en tous cas impliquent 37). En fait

(38) les formules pour $u(\rho), u(\sigma), u(\varepsilon_0)$ déterminent
de façon unique $\mu \in \hat{\mathbb{Z}}^*, \varepsilon_2 \in \{\pm 1\}, \alpha \in \Pi_{0,\hat{Z}}, \beta \in \Pi_0^\wedge$,

et on a nécessairement alors

(39) $\varepsilon_2 = \varepsilon_2(\mu), \varepsilon_2(\alpha) = \varepsilon_3(\mu)$

où $\varepsilon_2$ est l'hom que l'on déduit de la suite exacte

(40) $1 \to \Pi_0^\wedge \to \Pi_{0,\hat{Z}} \xrightarrow{\varepsilon_2} \{\pm 1\} \to 1$ .

Ainsi on a prouvé

Proposition Dans le cas universel ($\overline{G} = \underline{G}^{+\wedge}$)
on a $Z_{\rho,\sigma}^\sim \simeq M'[0] \overset{\text{déf}}{=} \mathcal{M}_{0,3}^\sim[0]$ ,
les composantes $u, \mu, \varepsilon$ d'un él. de
$Z_{\rho,\sigma}^\sim$ [unclear: correspondent resp.] En termes de
$u_0 \in M'[0]$ : $u_0$ identifié à un él. de $\mathcal{M}(\underline{G}^{+\wedge})$
$\mu = \chi(u_0)$ (le multiplicateur de $u_0$), $\varepsilon = \varepsilon_3(\mu), \varepsilon' = \varepsilon_2(\mu)$.

Corollaire On a $G_{\rho,\sigma} \simeq \mathcal{M} = N_{1,1}^\sim$ ,

$\Gamma_{\rho,\sigma} \simeq \mathcal{M}(\underline{G}^{+\wedge}) = N_{1,1}^\sim(\underline{G}^{+\wedge}) \simeq \mathcal{M}_{0,3}^\sim / \underline{G}_{0,3}^{+\wedge}$
$\simeq \Gamma_{\mathbb{Q}}^\sim$.

\newpage

%% ===== Page 583 =====

Remarque. Il est tentant d'écrire (aux [unclear: asymétries] près) les relations (38) sous la forme

(41) $$ \begin{cases} (\alpha^{-1} u)(\rho) = \rho & \text{i.e. } \alpha^{-1} u \in Z_\rho^\flat \\ (\beta^{-1} u)(\sigma) = \varepsilon_2 \sigma & \text{soit } \beta^{-1} u \in N_\sigma' \end{cases} $$

et d'introduire alors

(42) $$ \begin{cases} a = \alpha^{-1} u \in Z_\rho^\flat, & \text{d'où } u = \alpha a, \quad \alpha = u a^{-1} \\ b = \beta^{-1} u \in N_\sigma', & \text{d'où } u = \beta b, \quad \beta = u b^{-1} \end{cases} $$

et les développements du § 48 (rubrique XI) peuvent faire s'attendre que les quantités $a$, $b$, pour $u \in M[0]$ [unclear], dépendent multiplicativement de $u$. Nous allons voir cependant qu'il n'en est rien. Soit en effet
$$ \delta : G \xrightarrow{\sim} M \to \gamma_{\rho \sigma} \simeq \Gamma_Q $$
l'hom. canonique, de sorte que l'on a

(43) $$ \delta(u) = \delta(b) = \tau_\gamma^{\nu_3} \delta(a) $$

en posant
$$ \tau_\gamma = \delta(\tau) \quad , \quad \varepsilon_3 = (-1)^{\nu_3} \quad (\nu_3 \in \mathbb{Z}/2\mathbb{Z}) $$

Si $u'$ correspond à $b'$, $a'$, $\nu_3'$, ... on a

(43 bis) $$ \delta(u') = \delta(b') = \tau_\gamma^{\nu_3'} \delta(a') $$

et si $u'' = u' u$ correspond à $a''$, $b''$, $\nu_3''$ :

(43 ter) $$ \delta(u'') = \delta(b'') = \tau_\gamma^{\nu_3''} \delta(a'') $$

En fait, on a

(44) $$ b'' = b' b \quad , \quad \nu_2'' = \nu_2' + \nu_2 \quad \text{et} \quad \nu_3'' = \nu_3' + \nu_3 $$
(sc. $\varepsilon_2'' = \varepsilon_2' \varepsilon_2$ et $\varepsilon_3'' = \varepsilon_3' \varepsilon_3$)

\newpage

%% ===== Page 584 =====

On a
$$u(\sigma) = \varepsilon_2 int(\beta)\sigma$$
$$u'(\sigma) = \varepsilon'_2 int(\beta')\sigma$$

d'où
$$(u'u)(\sigma) = u'(u(\sigma)) = \varepsilon_2 u'(int(\beta)\sigma) = \varepsilon_2 \varepsilon'_2 int(u'(\beta)\beta')\sigma$$

avec $u'(\beta)\beta' \in \hat{\Pi}_0$, ce qui par l'unicité
prouve
(45) $$\beta'' = u'(\beta)\beta' \quad , \quad \varepsilon''_2 = \varepsilon'_2 \varepsilon_2$$

d'où on conclut bien
$$b'' = \underset{\beta'^{-1}u'}{b'} \underset{\beta^{-1}u}{b}$$

i.e.
$$\beta''^{-1}u'' = \beta'^{-1}u' \beta^{-1}u \quad \text{i.e.} \quad \beta'' = \beta' u'(\beta)$$
i.e.
$$\beta'' = u'(\beta)\beta' \quad \text{ce qui n'est autre que } 45).$$

La relation $\varepsilon''_3 = \varepsilon'_3 \varepsilon_3$ est aussi évidente, p. ex.
vu la relation (qui contient $\varepsilon_3$, et $\varepsilon'_3, \varepsilon''_3$)
(46) $$u(\rho) \text{ est } \underline{conj.} \text{ à } \rho^{\varepsilon_3}$$
Il reste donc à examiner la relation
hypothétique
$$a'' = a' a$$
qui implique
(*) $$\delta(a'') = \delta(a') \delta(a)$$
mais comparant (43) (43 bis) (43 ter) on trouve
la relation
[MARGIN: i. e.]
$$\delta(u')\delta(u) = \tau_\gamma^{\nu'_3} \delta(a') \tau_\gamma^{\nu_3} \delta(a)$$
$$\delta(u'') = \tau_\gamma^{\nu''_3} \delta(a'')$$

\newpage

%% ===== Page 585 =====

de sorte que la relation (*) équivaut à
$$ \tau_{\gamma}^{\nu_3' + \nu_3} \delta(a'') = \tau_{\gamma}^{\nu_3'} \delta(a') \tau_{\gamma}^{\nu_3} \delta(a) $$
ou encore
[unclear: crossed out equation]
i.e.
(47) $$ \delta(a'') = \tau_{\gamma}^{\nu_3}(\delta(a')) . \delta(a) $$

Dans la relation (*), équivaut à
(48) $$ \delta(a') = \tau_{\gamma}^{\nu_3}(\delta(a')) $$ i.e. $[\delta(a'), \tau_{\gamma}^{\nu_3}] = 1$ ou
i.e. $[\delta(u'), \tau_{\gamma}^{\nu_3}] = 1$

[MARGIN: puisque $\delta(a') = \tau_{\gamma}^{\nu_3}(\delta(u'))$]

Cela montre que si $\nu_3 = 1$ i.e. $\varepsilon_3 = -1$, alors
la relation (*) équivaut à la commutativité
de $\delta(u')$ avec $\tau_{\gamma}$, condition
du [unclear] sur un élément de
$\hat{\Gamma}_Q^\sim$ qui, pour un élément de $\hat{\Gamma}_Q$ lui-
-même, implique que $u' \in \{1, \tau_{\gamma}\}$
- et par des arguments heuristiques, on
a vu qu'il devait en être ainsi pour
tous éléments de $\hat{\Gamma}_Q^\sim$. Il apparait que
le succès de la présentation du §48 $\overline{XI}$
tient au fait que dans la situation
[MARGIN: Il se trouve que le centre $\Pi_0^\sim$ est réduit à $1$]
envisagée, le groupe $\Gamma_{1,1} \simeq \hat{\Z}^*$
était commutatif, (En dehors de la com-
mutativité $a \rightleftharpoons a'u \in Z_{\hat{\Gamma}} \hookrightarrow G_{\Gamma}$, ce qui a permis multi-
plicativement de l'élément $u \in \underline{M}[0]$, seule-
ment à condition de se borner à des $u \in \underline{M}[0]'$.

\newpage

%% ===== Page 586 =====

Nous allons essayer de secouer la noire, en trouvant un homomorphisme canonique
$$ \hat{M}[0] \overset{!}{\longrightarrow} \hat{Z}_{p,\sigma} \times_{\hat{\Gamma}_{p,\sigma}} \hat{N}_{p}^* \times_{\hat{\Gamma}_{p,\sigma}} \hat{N}_{\sigma}^* $$

[MARGIN: NB On a $\hat{S}_{0}\Gamma_{p,\sigma} / \hat{L}_0 \simeq \hat{S}_{0}\Gamma_{p,\sigma} / \hat{\Gamma}_{0, \dots}^2 \simeq \hat{\Gamma}_{p,\sigma}$]

Notons que ce deuxième membre est formé des triples
$$ (49) \quad \{ (u, a, b) \ | \ u = fb = f'a, \text{ avec } f, f' \in \dots \} $$
(sous $u \in \hat{M}[0]$, sous $a \in \hat{N}_p^*$, sous $b \in \hat{N}_\sigma^*$) ce qui implique $\varepsilon_2(u) = \varepsilon_0(b)$, $\varepsilon_3(u) = \varepsilon_0(a)$.

et pour un tel triple, on aura
$$ (49 \text{ bis}) \quad \begin{cases} u(\sigma) = int(fb)(\sigma) = int(f) \ int(b)(\sigma) = int(f) (\sigma^{\varepsilon_2(u)}) \\ u(p) = int(f'a)(p^{\varepsilon_3(u)}) = int(f') \varepsilon_3(a) (p) = int(f' \varepsilon_3(\tau))(p) \end{cases} $$

où on pose $\varepsilon_3(\tau) \overset{dfn}{=} \begin{cases} 1 & si \ \varepsilon_3=+1 \\ \tau^2 & si \ \varepsilon_3=-1 \end{cases}$ i.e. [unclear]

Si $u$ correspond à $\alpha, \beta$ (cf 38) ), ces formules (49 bis) s'écrivent - vu les relations
$$ \begin{cases} int(f)(\sigma) = int(\beta)(\sigma) & i.e. \ f = \beta \sigma^i \ (i \in \hat{\mathbb{Z}} / 4\hat{\mathbb{Z}}) \\ int(f' \varepsilon_3(\tau))(p) = int(\alpha)(p) & i.e. \ f' \varepsilon_3(\tau) = \alpha p^j \ (j \in \hat{\mathbb{Z}} / 6\hat{\mathbb{Z}}) \end{cases} $$

Dans ce [unclear]

Proposition. Pour un $u \in \hat{Z}_{p,\sigma}^! = \hat{M}[0] \simeq \hat{H}_{\tau_0, \tilde{s}_0} [unclear]$ fixé, les $(u, a, b) \in \hat{S}_0 \Gamma_{p,\sigma}$ au dessus de $u$ sont exactement les éléments de la forme

[MARGIN: où $\alpha, \beta, \varepsilon_2, \varepsilon_3$ sont définis par $u$ (38)]

$$ (50) \quad \begin{cases} a = f'^{-1} u \\ b = f^{-1} u \end{cases} \quad \text{avec } \begin{matrix} f' = \alpha p^j \varepsilon_3(\tau') & (j \in \hat{\mathbb{Z}}/6\hat{\mathbb{Z}}) \\ f = \beta \sigma^i & (i \in \hat{\mathbb{Z}}/4\hat{\mathbb{Z}}) \end{matrix} $$

Considérons alors les éléments $(u, a, b)$ au dessus de $u$ ayant l'expression la plus simple en termes de $\alpha, \beta$, en prenant $i = 0, j = 0$, correspondant donc aux choix :

\newpage

%% ===== Page 587 =====

597

$$
(51) \qquad \begin{cases} a = \varepsilon_3(\tau') \alpha^{-1} u & \text{i.e. } \alpha = u a^{-1} \varepsilon_3(\tau') \\ b = \beta^{-1} u & \text{i.e. } \beta = u b^{-1} \end{cases}
$$

Je me demande si cette section ensembliste
de $S_0 \Gamma_{p, \sigma}^{!\sim}$ au dessus de $\tilde{Z}_{p, \sigma}^! = M^![\Gamma_0]$ est
multiplicative, i.e. ce qui revient [unclear: au même]
si $\alpha$ et $\beta$ définis par (51) dépendent multi-
plicativement de $u$.

Proposition. Soit $(u, a, b) \in S_0 \Gamma_{p, \sigma}^{!\sim}$ (cf (50)).
Pour qu'on ait $i=0, j=0$, i.e. pour qu'on
ait (51) - i.e. pour que $(u, a, b)$
appartienne à la section ensembliste précédente -
il f. et i. qu'on ait les relations

[MARGIN: on aura posé $\varepsilon_3(\sigma) = \begin{cases} 1 \text{ si } \varepsilon_3 = +1 \\ -1 \text{ si } \varepsilon_3 = -1 \end{cases}$]
$$
(52) \qquad u \equiv a \equiv \varepsilon_3(\sigma) b \pmod{\hat{\Pi}_0}
$$

[MARGIN: N.B. $\tau'\tau = \sigma^{-1}$ \\ $\tau^{-1}\sigma^{-1} = \tau$]
Démonstration. Comme $\beta = u b^{-1}$ et $\alpha = u a^{-1} \varepsilon_3(\tau')$
sont tous deux dans $\hat{\Pi}_0$, la condition est
nécessaire. Pour la suffisance, on note
que (50) s'écrit, vu les relations de $\hat{\Pi}_0$,

[MARGIN: (53) $f, f'\varepsilon_3(\sigma) \in \hat{\Pi}_0$]
que $f = u a^{-1}$ i.e. $f' \varepsilon_3(\sigma) = u a^{-1} \varepsilon_3(\sigma)$ [unclear: crossed out text]
et réduit mod $\hat{\Pi}_0$ dans $\mathfrak{S}/\hat{\Pi}_0$
qui s'écrit $\sigma^i = 1$ et $\alpha p^j \varepsilon_3(\tau' \sigma^{-1}) = 1$,
dernière relation s'écrit aussi $\alpha p^j \varepsilon_3(\tau') = 1$ i.e.
cette dernière relation s'écrit $\varepsilon_3(\tau) \alpha p^j = 1$
donne la relation $\varepsilon_3(\tau) \alpha p^j = 1$ ou $\varepsilon_3(\tau) \alpha = 1$, d'où
$\varepsilon_3(\tau) \alpha = 1$ donc la condition est
$\sigma^i = p^j = 1 \quad \text{i.e.} \quad i=0, j=0$
cqfd.

Corollaire. La section envisagée est multiplicative.

\newpage

%% ===== Page 588 =====

En même temps, ceci nous suggère
une modification des constructions des
questions précédentes, en remplaçant les
produits fibrés sur $V_{g,r} = M_{g,r} \dots$ par
des produits fibrés sur $V_{g,r}$ le quotient [unclear: plus fin] $C_{g,r} / \Pi_0$, qui est un
rev. de degré 12 de celui-ci.

Soient maintenant $u, b$ correspondant à $a, b$
$\varepsilon_3$ et $u'$, correspondant à $a', b', \varepsilon'_3$ ,
soit $u'' = u'u$, correspondant à $a'', b'' = b'b$ et
$\varepsilon''_3 = \varepsilon'_3 \varepsilon_3$, de sorte qu'on aura

$$u \equiv \varepsilon_3(\sigma) a \quad (\hat{\Pi}_0)$$
$$u' \equiv \varepsilon'_3(\sigma) a' \quad (\hat{\Pi}_0)$$
$$u'' \equiv \varepsilon''_3(\sigma) a'' \quad (\hat{\Pi}_0)$$
$$||$$
$$u'u \equiv \varepsilon \varepsilon'_3(\sigma) \varepsilon_3(\sigma) a'' \quad \text{où } \varepsilon \in \{\pm 1\}, \varepsilon = -1 \text{ ssi } \varepsilon_3 = -1 \text{ et } \varepsilon'_3 = -1$$
$$|||$$
$$(\varepsilon'_3(\sigma)a') (\varepsilon_3(\sigma)a)$$

[MARGIN: 
de $\sigma(u') \equiv u' \quad (\hat{\Pi}_0)$
$\sigma(u) \equiv -u \quad (\hat{\Pi}_0)$
i.e. ]

on trouve que

$$a' \varepsilon_3(\sigma) a \equiv \varepsilon \varepsilon_3(\sigma) a'' \quad (\hat{\Pi}_0) \quad i.e.$$
$$a'' \equiv \varepsilon \varepsilon_3(\sigma)(a') . a$$

[MARGIN: N.B. $\varepsilon_3(\sigma)$ et $\varepsilon_3(\sigma^{-1})$ (53) opéreront...]

Cette relation est asymétrique
$$a'' = a' a \quad \text{équivaut à} \quad a'' \equiv a' a \quad (\hat{\Pi}_0)$$
soit 
vrai $\iff \varepsilon_3(\sigma)(a') \equiv \varepsilon a' \quad (\hat{\Pi}_0)$

(54) $$\varepsilon_3(\sigma)(a') \equiv \varepsilon a' \quad (\hat{\Pi}_0)$$

[MARGIN:
1) $\varepsilon_3=1, \varepsilon'_3=1 : a'' \equiv a'a$
2) $\varepsilon_3=1, \varepsilon'_3=-1 : a'' \equiv a'a$
3) $\varepsilon_3=-1, \varepsilon'_3=1 : a'' \equiv \sigma(a')a$
4) $\varepsilon_3=-1, \varepsilon'_3=-1 : a'' \equiv -\sigma(a')a$]

Cette condition a l'air bien curieuse [unclear: ... cette]
condition est [unclear: automatiquement] vérifiée

587

\newpage

%% ===== Page 589 =====

(54) $$ \varepsilon_3(\sigma)(a') \equiv \varepsilon a' \quad (\hat{\Pi}_0) \quad \text{où } \begin{cases} \varepsilon \in \{\pm 1\}, \varepsilon = -1 \iff \varepsilon_3 = \varepsilon'_3 = -1 \\ \varepsilon_3(\sigma) = \begin{cases} id / 1 & \text{si } \varepsilon_3 = +1 \\ \sigma & \text{si } \varepsilon_3 = -1 \end{cases} \end{cases}$$

Cette condition (54) s'explicite suivant les cas :

si $\varepsilon_3 = 1$, alors $\varepsilon = 1$ et (54) [unclear: est vérifiée]

(55) si $\varepsilon_3 = -1$, alors (54) s'écrit $\sigma(a') \equiv \varepsilon'_3 \cdot a' \quad (\hat{\Pi}_0)$ i.e.
$$[\dot{\sigma}, \dot{a}'] = \varepsilon'_3$$
où les $\dot{}$ désignent les classes dans $M/\hat{\Pi}_0$
$$ a' = \varepsilon'_3(\sigma)^{-1} u' = \begin{cases} u' & \text{si } \varepsilon'_3 = 1 \\ \sigma^{-1} u' & \text{si } \varepsilon'_3 = -1 \end{cases} $$

[unclear: Comme on] trouve que la condition (55, b) équivaut à :

(56) $[\dot{\sigma}, \dot{u}'] = \varepsilon'_3$ dans $M/\hat{\Pi}_0$, i.e.
$$[\dot{\sigma}, \dot{u}'] = \varepsilon'_3$$
(valable aussi bien pour $\varepsilon'_3 = 1$ que $\varepsilon'_3 = -1$).
Or on a vu (35) que l'on a dans $\hat{\Pi}_0$
$$\dot{u}'(\dot{\sigma}) = \varepsilon_2(u') \dot{\sigma}$$
i.e. $$[\dot{u}', \dot{\sigma}] = \varepsilon_2(u')$$ ou encore

(57) $[\dot{\sigma}, \dot{u}'] = \varepsilon_2(u')$ .

Comparant (56) et (57), on voit que (56) prend la forme équivalente

(58) $\varepsilon_2(u') = \varepsilon'_3$

On a donc prouvé :

[MARGIN: NB $\begin{cases} \varepsilon_3 = \varepsilon_3(u) = \varepsilon_2(u) \\ \varepsilon'_3 = \varepsilon_3(u') = \varepsilon_2(u') \\ \varepsilon"_3 = \varepsilon_3(u") = \varepsilon_2(u") \end{cases}$]

Proposition. Soient $u, u'$ et $u" = u' u \in M_{[\hat{\Pi}_0]} = \hat{\mathbb{Z}}^*_{>0}$,
correspondant à $\varepsilon_3, \varepsilon'_3, \varepsilon"_3 = \varepsilon'_3 \varepsilon_3 \in \{\pm 1\}$, et :
$a = \varepsilon_3(\sigma) \alpha^{-1} u$, $a' = \varepsilon'_3(\sigma) \alpha'^{-1} u'$, $a" = \varepsilon"_3(\sigma) \alpha"^{-1} u"$,
avec $\alpha, \alpha', \alpha" \in \mathcal{N}_{\bar{k}}$, $\alpha \equiv \varepsilon_3(\sigma) u, \alpha' \equiv \varepsilon'_3(\sigma) u', \dots$
Pour que les [unclear: conditions soient] équivalentes :

\newpage

%% ===== Page 590 =====

a) On a $a'' = a' a$

b) On a $a'' \equiv a' a \quad (\hat{\Pi}_0)$
i.e. $a(\rho)=\rho$,
[MARGIN: b) On a [unclear, crossed out: $\varepsilon_3=1 \quad \text{ou} \quad \sigma(a') \equiv \varepsilon'_3 a' \quad (\hat{\Pi}_0)$] ]

b') On a $\varepsilon_3 = 1$ \quad (ou \quad $\dot{a}'(\sigma) \equiv \varepsilon'_3 \sigma \quad (\hat{\Pi}_0)$

c) On a $\varepsilon_3=1 \quad$ (ou [unclear, crossed out] $\equiv \varepsilon'_3 u' \quad (\hat{\Pi}_0)$

c') On a $\varepsilon_3=1$ i.e. $a(\rho)=\rho$, \quad ou \quad $u'(\sigma) \equiv \varepsilon'_3 u' \quad (\hat{\Pi}_0)$

d) On a $\boxed{\varepsilon_3=1 \quad \text{ou} \quad \varepsilon'_2 = \varepsilon'_3 \quad \text{i.e. } \varepsilon'_2\varepsilon'_3=1}$

Corollaire \quad Au-dessus des sous-
groupes
(59) 
$$ \begin{cases} 
Z_{\rho,\sigma}^{!} = M^{!}_{[\Gamma_0]} \stackrel{df}{=} \{ u \in Z_{\hat{\rho},\hat{\sigma}}^! \mid \varepsilon_3(u) = 1 \} \\ 
Z_{\rho,\sigma}^{!\text{bis}} = M^{!\text{bis}}_{[\Gamma_0]} \stackrel{df}{=} \{ u \in Z_{\hat{\rho},\hat{\sigma}}^! \mid \varepsilon_2(u) = \varepsilon_3(u) \} 
\end{cases} $$
l'application section ensembliste
(60) $A : Z_{\rho,\sigma}^{!} \to S_0 \Gamma_{\rho,\sigma}^{!}$
$M^{!}_{[\Gamma_0]}$
$u \mapsto (u, a=\varepsilon_3(\tau)\alpha^{-1}u, b=\beta^{-1}u)$
(cf prop p. 597)
est une section multiplicative. Mais
elle n'est pas multiplicative au
dessus de $M^{!\text{bis}}_{[\Gamma_0]}$ tout entier.

\newpage

%% ===== Page 591 =====

Remarque les éléments de l'im. [unclear: pr] inverse $S_0 \Gamma_{\rho,\sigma}^!$ de $Z_{\rho,\sigma}^!$ dans $S_0 \Gamma_{\rho,\sigma}^!$ qui appartiennent à la section envisagée sont caractérisés (par vertu de la prop. p. 598) par la relation 
[MARGIN: (en plus de $\varepsilon_3(u)=1$ qui caractérise les éléments de $S_0 \Gamma_{\rho,\sigma}^!$)]
(61) $$u \equiv b \equiv a \pmod{\hat{\Pi}_0}$$

Il me [unclear: parait] évident que c'est un sous-groupe. Je n'ai pas trouvé de jolie expression analytique pour l'appartenance d'un élément de $S_0 \Gamma_{\rho,\sigma}^!$ (caractérisé par la relation $\varepsilon_2(u)\varepsilon_3(u)=1$) à la section correspondante.

Proposition Soit $S_0^\pi \Gamma_{\rho,\sigma}^! \subset S_0 \Gamma_{\rho,\sigma}^!$ le sous-ensemble de $S_0 \Gamma_{\rho,\sigma}^!$ [unclear: défini par]

(62) $S_0^\pi \Gamma_{\rho,\sigma}^! \subset S_0 \Gamma_{\rho,\sigma}^!$
$\parallel$
$\{ (u,a,b) \in \mathcal{M}_{[\sigma]} \mid u \equiv b \ (\hat{\Pi}_0), \ u \equiv \varepsilon_3(\sigma) a \pmod{\hat{\Pi} = \hat{\Pi}_0 \rtimes \{\pm 1\}} \}$

Alors
a) $S_0^\pi \Gamma_{\rho,\sigma}^!$ est un sous-groupe de $S_0 \Gamma_{\rho,\sigma}^!$
b) L'homomorphisme $S_0^\pi \Gamma_{\rho,\sigma}^! \to Z_{\rho,\sigma}^! = \mathcal{M}_{[\sigma]}^!$,
$(u,a,b) \mapsto u$, est surjectif, et son noyau est le sous-groupe, isomorphe à $\mu = \{\pm 1\}$, des éléments $(1, \varepsilon, 1)$ avec $\varepsilon \in \mu = [unclear]$

\newpage

%% ===== Page 592 =====

Dém Le fait qu'on a un sous-groupe
résulte des calculs déjà faits p. 597!, 598,
la condition pour que le produit de deux
éléments $(u',a',b')$ $(u,a,b)$ de $S\hat{\Gamma}_{p,\sigma}^!$ y soit
s'exprime encore par la formule
(56) $[\dot{\sigma}, \dot{u}'] = \dot{\varepsilon}'_3$ dans $M/\hat{\Pi}$ (au lieu
de $M/\hat{\Pi}_0$)
[unclear: on a encore]
(56 bis) $[\dot{u}', \dot{\sigma}] = \dot{\varepsilon}'_3$ avec $\dot{\varepsilon}'_3 = \dot{\varepsilon}_3$
$\dot{u}' (\dot{\sigma}) \dot{\sigma}^{-1} = \dot{\varepsilon}'_3$
Or [unclear] "W-structure" tombent dans le sous-groupe
central $\hat{\Pi}/\hat{\Pi}_0$, donc ils sont triviaux dans
$\widehat{S}/\hat{\Pi}$, donc la relation (56 bis) est trivialement
satisfaite.

La surjectivité de l'application $S\hat{\Gamma}_{p,\sigma}^! \to Z_{p,\sigma}^!$
provient du fait que $S\hat{\Gamma}_{p,\sigma}^!$ contient la
section ensembliste $A(Z_{p,\sigma}^!)$ de la prop.
p. 597. Le noyau est formé des
$(1, a, b)$ avec $a \in N^*_{\Gamma} \cap \hat{\Pi}$, $b \in N^*_{\sigma} \cap \hat{\Pi}_0$.
Or $N^*_{\Gamma} \cap \widehat{S} = L_{\Gamma}$, $N^*_{\sigma} \cap \widehat{S} = L_{\sigma}$, et
(57) $L_{\Gamma} \cap \hat{\Pi} = L_{\sigma} \cap \hat{\Pi} = \mu$, $L_{\Gamma} \cap \hat{\Pi}_0 = L_{\sigma} \cap \hat{\Pi}_0 = \{1\}$

\newpage

%% ===== Page 593 =====

600

d'où la caractérisation du noyau.

Corollaire Considérons $\hat{S}_{\rho,\sigma}^!$ comme
une extension de $\hat{Z}_{\rho,\sigma}^! = \hat{M}[\tau_0]$ par
$L_\rho \times L_\sigma \simeq \mathbb{Z}/6\mathbb{Z} \times \mathbb{Z}/4\mathbb{Z}$, donc
(en écrivant $L_\rho \simeq \mathbb{Z}/3\mathbb{Z} \times \mathbb{Z}/2\mathbb{Z}$)
une extension de 
$$\underbrace{\mathbb{Z}/2\mathbb{Z}}_{\mu} \times \underbrace{\mathbb{Z}/3\mathbb{Z}}_{L_\rho} \times \underbrace{\mathbb{Z}/4\mathbb{Z}}_{L_\sigma}.$$

Cette extension correspond donc canoniquement
à trois extensions, par $\mu, L_\rho, L_\sigma$
respectivement. L'extension par $L_\rho \times L_\sigma$
admet deux sections canoniques.
Il reste une extension (c'est juste $\hat{S}_{\rho,\sigma}^!$) de $\hat{Z}_{\rho,\sigma}^!$ par $\mu$.

Cette extension correspond justement à
l'ambiguïté de signes pour choisir
un $\alpha \in \hat{\pi}$ au dessus de $\tilde{\alpha} \in \hat{\pi}/\pm 1$
défini par $u \in \hat{Z}_{\rho,\sigma}^! = \hat{M}[\tau_0]$, qui se reflète
par l'ambiguïté à choisir le signe
de $a \in \hat{N}_\rho$, lié à $\alpha$ par (51)
[MARGIN: (58)] $a^2 = \varepsilon_3(u)^2 u$, $\alpha = u a^{-1} \varepsilon_3(u)$.
Nous avons pu lever cette ambiguïté
par l'exigence $\alpha \in \hat{\Pi}_{\hat{\mathbb{Z}}} = \hat{\Pi}_{\{2\}} \cdot \hat{\Pi}_3$, et nous

\newpage

%% ===== Page 594 =====

venons de constater que ce choix, pour
canonique qu'il soit, n'est pas multiplica-
tif. En même temps, on trouve
l'explicitation complète de $\alpha''$ (normalisé
canonique par la condition $\alpha'' \in \hat{\Pi}_{0\tilde{2}}$)
associé à $u'' = u'u$, en termes de
$\alpha'$ (associé à $u'$) et $\alpha$ (associé à $u$)
également normalisés, (En fait, la page
p 598 nous dit dans quels cas
[unclear] on a

(59) $$a'' = a'a$$

savoir ssi on a

(59 bis) $\varepsilon_3(u)=1$ ou $\varepsilon_3(u')\varepsilon_3(u'')=1$ .

Mais on sait maintenant que dans
les deux cas

(60) $$a'' = \varepsilon a'a \quad \varepsilon \in \{\pm 1\}$$

et on trouve maintenant que $\varepsilon = +1$
ssi on est dans le cas (59 bis), $-1$
dans le cas contraire i.e. où $\varepsilon_3(u)=-1$
et $\varepsilon_3(u')\varepsilon_3(u'')=+1$. En résumé -

Corollaire Soient $u, u' \in \Lambda'[0]$, d'où $u''=u'u$,
soient $\alpha, \alpha', \alpha'' \in \hat{\Pi}_{0\tilde{2}}$ associés (définis
par [unclear: a=1 resp.] etc ) et posons

\newpage

%% ===== Page 595 =====

(61) $$a = \varepsilon_3(\tau') \alpha^{-1} u, \quad a' = \varepsilon'_3(\tau') \alpha'^{-1} u', \quad a'' = \varepsilon''_3(\tau') \alpha''^{-1} u''$$

où $\varepsilon_3(\tau') = \{$ $1$ si $\varepsilon_3 = 1$, $\tau'$ si $\varepsilon_3 = -1$, et de même pour $\varepsilon'_3(\tau'), \varepsilon''_3(\tau')$,
et où $\varepsilon_3 = \varepsilon_3(u) = \varepsilon_3(\chi(u))$, et de même pour $\varepsilon'_3, \varepsilon''_3$
— d'où $\varepsilon''_3 = \varepsilon'_3 \varepsilon_3$. On a donc

(62) $$a'' = \varepsilon a' a , \qquad \text{où } \varepsilon \in \{ \pm 1 \}$$
et
(63) $\varepsilon = +1$ ssi $\varepsilon_3 = 1$ [unclear] $\varepsilon'_3 \varepsilon_3 = 1$ d'où
$\varepsilon = -1$ ssi $\varepsilon_3 = -1$ et $\varepsilon'_3 = \varepsilon_3 = -1$

[unclear crossed out text]
[unclear crossed out text] $\alpha'' = u'' a''^{-1} \varepsilon''_3(\tau')$
$= \varepsilon u' u a^{-1} a'^{-1} \varepsilon''_3(\tau') \varepsilon'_3(\tau')$
[MARGIN: Pas vrai en général]

(64) $$\alpha'' = \varepsilon u'(\alpha) \alpha' \qquad \text{où } \varepsilon \text{ donné par (63).}$$

En effet $\left\{ \begin{array}{l} \alpha = u a^{-1} \varepsilon_3(\tau') \\ \alpha' = u' a'^{-1} \varepsilon'_3(\tau') \end{array} \right.$ d'où
$u'(\alpha) = u' u a^{-1} \varepsilon_3(\tau') u'^{-1}$
$u'(\alpha) \alpha' = u' u a^{-1} \varepsilon_3(\tau') u'^{-1} u' a'^{-1} \varepsilon'_3(\tau')$

La formule (64), en vertu de (*), prend
donc la forme
$$a'^{-1} \varepsilon_3(\tau') = \varepsilon_3(\tau') a'^{-1}$$
i.e.
$$[ \varepsilon_3(\tau') , a'^{-1} ] = 1$$

elle n'est pas toujours vérifiée si $\varepsilon_3 = -1$
i.e. $\varepsilon_3(\tau') = \tau'$, [unclear] elle représente une
restriction draconienne sur $a'$. La formule
toujours correcte est

(65) $$\alpha'' = ( \varepsilon u'(\alpha) \alpha' ) . \varepsilon'_3(\tau') ( [a', \varepsilon_3(\tau')] )$$
[unclear: best guess: $[ a'_3(\tau') (a') , \varepsilon_3(\tau') ]$ written below]

594

\newpage

%% ===== Page 596 =====

[MARGIN: sic! Voilà le monstre qui m'a fait frémir [unclear: pendant des mois] m'a l'air peut-être bien doux !]

Comme $\varepsilon_3'(\tau') a' = \alpha'^{-1} u'$ (61), la

formule un peu alambiquée :

$$ (65) \quad \alpha'' = (\varepsilon \ u'(\alpha) \ \alpha') . \big[ \underset{\varepsilon_3'(\tau) a'}{\alpha'^{-1} u'}, \varepsilon_3(\tau) \big] $$

(où $\varepsilon \in \{\pm 1\}$ déf. dans (63))

Je veux : examiner la question

si l'extension de $\hat{S}_0^\pi \Gamma_{\bar{p}, \sigma}$ de $\hat{Z}_{\bar{p}, \sigma}^1 = M(\bar{c}_0)$

par $\mu = \{\pm 1\}$ est triviale ou non, tenant

compte du fait que ses restrictions

à $\hat{Z}_{\bar{p}, \sigma}^{1'}$ et $\hat{Z}_{\bar{p}, \sigma}^{1'''}$ sont triviales.

Soient $\hat{S}_0^\pi \Gamma_{\bar{p}, \sigma}^{1'}$, $\hat{S}_0^\pi \Gamma_{\bar{p}, \sigma}^{1'''}$ les images inv.

de $\hat{Z}_{\bar{p}, \sigma}^{1'}$, $\hat{Z}_{\bar{p}, \sigma}^{1'''}$ dans $\hat{S}_0^\pi \Gamma_{\bar{p}, \sigma}$, d'où

$$ (66) \quad \begin{cases} \hat{S}_0^\pi \Gamma_{\bar{p}, \sigma}^{1'} = \hat{S}_0^\pi \Gamma_{\bar{p}, \sigma} \cap \Gamma_{\bar{p}, \sigma}^{1'} \\ \hat{S}_0^\pi \Gamma_{\bar{p}, \sigma}^{1'''} = \hat{S}_0^\pi \Gamma_{\bar{p}, \sigma} \cap \Gamma_{\bar{p}, \sigma}^{1'''} \end{cases} $$

et on a des suites exactes d'ext. centrales

$$ (67) \quad \begin{cases} 1 \to \mu \to \hat{S}_0^\pi \Gamma_{\bar{p}, \sigma}^{1'} \to \hat{Z}_{\bar{p}, \sigma}^{1'} \to 1 \\ 1 \to \mu \to \hat{S}_0^\pi \Gamma_{\bar{p}, \sigma}^{1'''} \to \hat{Z}_{\bar{p}, \sigma}^{1'''} \to 1 \end{cases} $$

qui sont scindées par des sous-groupes

très opportuns $\hat{S}_0^{\pi_0} \Gamma_{\bar{p}, \sigma}^{1'}$, $\hat{S}_0^{\pi_0} \Gamma_{\bar{p}, \sigma}^{1'''}$

de sorte que

$$ (68) \quad \begin{cases} \hat{S}_0^\pi \Gamma_{\bar{p}, \sigma}^{1'} \simeq \hat{S}_0^{\pi_0} \Gamma_{\bar{p}, \sigma}^{1'} \times \mu \\ \hat{S}_0^\pi \Gamma_{\bar{p}, \sigma}^{1'''} \simeq \hat{S}_0^{\pi_0} \Gamma_{\bar{p}, \sigma}^{1'''} \times \mu \end{cases} $$

\newpage

%% ===== Page 597 =====

Je me pose la question si ces sous-groupes

[MARGIN: ss-groupes distingués dans $S_0^\pi \Gamma_{\rho,\sigma}$]

(69) $$ \begin{cases} S_0^{\pi_0} \Gamma_{\rho,\sigma}^{!} \subset S_0^\pi \Gamma_{\rho,\sigma}^{!} \subset S_0^\pi \Gamma_{\rho,\sigma} \\ S_0^{\pi_0} \Gamma_{\rho,\sigma}^{'''} \subset S_0^\pi \Gamma_{\rho,\sigma}^{'''} \subset S_0^\pi \Gamma_{\rho,\sigma} \end{cases} $$

sont invariants dans $S_0^\pi \Gamma_{\rho,\sigma}$, ce qui
signifierait (vu [unclear] que ces
sous-groupes) que l'extension
$S_0^\pi \Gamma_{\rho,\sigma}$ de $Z_{\rho,\sigma}$ par $\mu$, scindé sur
$Z_{\rho,\sigma}^!$ et sur $Z_{\rho,\sigma}^{'''}$, provient en fait
(via lesdits scindages) d'extensions de
$Z_{\rho,\sigma}^! / Z_{\rho,\sigma}$ ($\simeq \pm 1$) et de $Z_{\rho,\sigma}^{'''} / Z_{\rho,\sigma}$ ($\simeq \pm 1$)
par $\mu$, qu'il faudrait ensuite
caractériser (comme triviales ou non triviales).

Comme les sous-groupes en question
sont distingués dans $S_0^\pi \Gamma_{\rho,\sigma}^{!}$ resp $S_0^\pi \Gamma_{\rho,\sigma}^{'''}$, il
si $x$ est un élém de $S_0^\pi \Gamma_{\rho,\sigma}$ n'étant
pas dans $S_0^\pi \Gamma_{\rho,\sigma}^{!}$ resp $S_0^\pi \Gamma_{\rho,\sigma}^{'''}$, dire que
$S_0^{\pi_0} \Gamma_{\rho,\sigma}^{!}$ (resp. $S_0^{\pi_0} \Gamma_{\rho,\sigma}^{'''}$) est distingué dans
$S_0^\pi \Gamma_{\rho,\sigma}$ signifie simplement qu'il
est invariant par $x$.

1º) Cas de $S_0^{\pi_0} \Gamma_{\rho,\sigma}^{!} \subset S_0^\pi \Gamma_{\rho,\sigma}^{!}$. On prendra
$x = (u_0, a_0, b_0)$ avec $u_0 = \tau$, d'où $\varepsilon_3(u_0) = \varepsilon_1(u_0) = \chi(u_0) = -1$

(70) $$ \begin{cases} x = \bar{\tau}_2 \text{ , qui } \notin S_0^\pi \Gamma_{\rho,\sigma}^! \\ \alpha_0 = 1, \beta_0 = 1 \text{ d'où } a_0 = \varepsilon_3'(\tau) \alpha_0^{-1} u_0 = \tau', b_0 = \beta_0^{-1} u_0 = \tau \end{cases} $$
$$ x = (u_0 = \tau, a_0 = \tau', b_0 = \tau) $$

\newpage

%% ===== Page 598 =====

donc ...
$x \in \widehat{S}_0^{\Pi_0}\widehat{\Gamma}^{\prime,\prime\prime} \setminus \widehat{S}_0^{\Pi}\widehat{\Gamma}^{\prime\prime}$ [unclear: cross outs] et $x \notin \widehat{S}_0^{\Pi}\widehat{\Gamma}^{\prime\prime}$
et il faut voir si $x$ normalise $\widehat{S}_0^{\Pi_0}\widehat{\Gamma}^{\prime\prime}$.
Soit
(71) $y = (u,a,b) \in \widehat{S}_0^{\Pi_0}\widehat{\Gamma}^{\prime\prime} \quad i.e. \quad \begin{cases} u \equiv a \equiv b \quad (\Pi_0) \\ \varepsilon_3(u) = 1 \end{cases}$
et calculons
$$x y x^{-1} = x y x = (\tau(u), \tau'(a), \tau(b))$$
est-il dans $\widehat{S}_0^{\Pi_0}\widehat{\Gamma}^{\prime\prime}$ (et on a vu dans $\widehat{S}_0^{\Pi}\widehat{\Gamma}^{\prime\prime}$) i.e. a-t-on
(72) $\tau(u) \equiv \tau'(a) \equiv \tau(b) \quad (\Pi_0) \quad ?$
La relation $\tau(u) \equiv \tau(b) (\equiv \tau(a))$ est conséquence
imm. de 71), la $2^e$ équivaut à
$\tau'(a) \equiv \tau(a) \quad (\Pi_0) \quad i.e. \quad \sigma^{-1}(\tau(a)) = \tau(a)$
i.e.
[MARGIN: NB]
(73) $[\sigma, \tau(a)] = 1$ dans $\widehat{\mathfrak{S}}/\Pi_0$
Ou encore [unclear: se récrivant] en vertu de (35 bis)
$a = \varepsilon_0(\tau') \alpha^{-1} u$, d'où $\tau(a) = \tau(\alpha)^{-1} \tau(u)$
$\uparrow \varepsilon_2(\tau(a))=1 \quad i.e. \quad \varepsilon_2(a) = 1$
ou enfin (comme $a \equiv u \quad (\Pi_0)$)
(74) $\varepsilon_2(u) = 1$.
Cette condition n'est pas conséquence de
la condition $\varepsilon_3(u) = 1$, de sorte
que $\widehat{S}_0^{\Pi_0}\widehat{\Gamma}^{\prime\prime}$ n'est pas invariant dans $\widehat{S}_0^{\Pi}\widehat{\Gamma}^{\prime\prime}$
$\widehat{S}_0^{\Pi_0}\widehat{\Gamma}^{\prime,\prime\prime}$ est-il invariant dans $\widehat{S}_0^{\Pi}\widehat{\Gamma}^{\prime,\prime\prime}$ ?
Soit [unclear cross out]
(75) $x = (u_0, a_0, b_0) \in \widehat{S}_0^{\Pi_0}\widehat{\Gamma}^{\prime,\prime\prime} \setminus \widehat{S}_0^{\Pi_0}\widehat{\Gamma}^{\prime\prime} \quad d'où \quad \begin{cases} \varepsilon_3(u_0) = 1 \\ \varepsilon_1(u_0) = -1 \end{cases}$
$u_0 \equiv b_0 \equiv \tau a_0 \quad (\Pi_0)$

\newpage

%% ===== Page 599 =====

et étudions son action sur
$$ (76) \quad y = (u,a,b) \in S_0^{\Pi_0} \Gamma^{!!!} \quad \text{i.e.} \begin{cases} \varepsilon_3(u) \varepsilon_3(a)=1 \\ u \equiv b \equiv \varepsilon_3(\sigma) a \quad (\Pi_0) \end{cases} $$

On a
$$ (77) \quad x y x^{-1} = (u_0(u), a_0(a), b_0(b)) $$
et la relation $x y x^{-1} \in S_0^{\Pi_0} \Gamma^{!!!}$ s'écrit
[MARGIN: N.B. $\varepsilon_3(u_0(u)) = \varepsilon_3(u)$]
$$ (78) \quad u_0(u) \equiv b_0(b) \equiv \varepsilon_3(\sigma) . a_0(a) \quad . $$
Or les relations 75) b) et 76) b) impliquent
aussitôt
$$ u_0(u) \equiv b_0(b) \equiv a_0(\varepsilon_3(\sigma) . a) $$
$$ \equiv \varepsilon_3[a_0(\sigma)] . a_0(a) $$
donc la relation (78) équivaut à
$$ (79) \quad \varepsilon_3[\sigma] \equiv \varepsilon_3[a_0(\sigma)] $$
ce qui est évidemment toujours vérifié si $\varepsilon_3=1$,
et pour $\varepsilon_3=-1$ s'écrit :
$$ (80) \quad \sigma \equiv a_0(\sigma) \quad (\hat{\Pi}_0) $$
i.e. $\varepsilon_2(a_0)=1$ i.e. $\varepsilon_1(u_0)=1$,
relation qui est fausse par l'hyp. 75 a).
Donc on trouve encore $S_0^{\Pi_0}\Gamma^{!!!}$ n'est
pas invariant dans $S_0^\pi\Gamma^!$. Pour en finir, posons

$$ (81) \quad \begin{cases} S_0\Gamma^\circ = \operatorname{Ker}( S_0^\pi\Gamma^\circ \xrightarrow{(\varepsilon_2,\varepsilon_3)} \mu \times \mu ) & = S_0\Gamma' \cap S_0\Gamma'' \\ S_0^\pi\Gamma^\circ = S_0\Pi_0 \cap S_0^\pi\Gamma^! , & = S_0^\pi\Gamma' \cap S_0^\pi\Gamma^{'''} \\ S_0^{\Pi_0}\Gamma^\circ = S_0\Pi_0 \cap S_0^{\Pi_0}\Gamma^{!!} , \end{cases} $$
de sorte qu'on a une suite exacte
$$ (82) \quad 1 \to \mu \to S_0^{\Pi_0}\Gamma^\circ \to S_0\Gamma^\circ \to 1 \ , $$

\newpage

%% ===== Page 600 =====

[unclear: produit] semi-direct par $S_0^{\pi_0}\Gamma_0$ ...
[unclear: tr... d...]
(82) $$S_0^{\pi}\Gamma_0 = S_0^{\pi_0}\Gamma_0 \times \mu$$

Ceci posé, les calculs précédents montrent que $S_0^{\pi_0}\Gamma'_0$, qui est d'indice 8 dans $S_0^{\pi}\Gamma'$, est un sous-groupe invariant. En résumé :

Proposition Considérons les sous-groupes (d'indices 4, 4, 8) $S_0^{\pi_0}\Gamma'$, $S_0^{\pi_0}\Gamma'''$, $S_0^{\pi_0}\Gamma_0$ de $S_0^{\pi}\Gamma'$ qui sont les images [unclear] respectives à $Z'$, $Z'''$, $Z'_0$ de la section co-absolue(?) $A$ définie par (51) (52) de $S_0^{\pi}\Gamma'$ sur $Z = M[0]$ — [unclear: qui opère] par multiplication par $\mu$ sur ces sous-groupes, se déduisent donc des scindages des images inverses $S_0^{\pi}\Gamma'$, $S_0^{\pi}\Gamma'''$, $S_0^{\pi}\Gamma_0$ de $Z'$, $Z'''$, $Z'_0$ dans $S_0^{\pi}\Gamma'$ :
(83) $$S_0^{\pi}\Gamma' \simeq S_0^{\pi_0}\Gamma' \times \mu, \quad S_0^{\pi}\Gamma''' \simeq S_0^{\pi_0}\Gamma''' \times \mu, \quad S_0^{\pi}\Gamma_0 \simeq S_0^{\pi_0}\Gamma_0 \times \mu$$

[MARGIN: ce qui [unclear: assure] qu'il n'existe pas de section [unclear] au dessus des [unclear: éclat.] de $Z'$ par $\Pi_0$ qui coïncide avec $A$ sur $Z'''$]

On a de plus ceci :

a) $S_0^{\pi_0}\Gamma'$, $S_0^{\pi_0}\Gamma'''$ ne sont pas distingués dans $S_0^{\pi}\Gamma'$

b) $S_0^{\pi_0}\Gamma_0$ est distingué dans $S_0^{\pi}\Gamma'$, de sorte que l'ext. $S_0^{\pi}\Gamma'$ de $Z' = M'[0]$ par $\mu$ est [unclear: canoniquement isom. à] : l'image inverse d'une extension centrale $S_0^{\pi}\Gamma' / S_0^{\pi_0}\Gamma_0$ de $Z' / Z'_0$ par $\mu \times \mu$ par $\mu$.

\newpage

%% ===== Page 601 =====

604

Examinons cette extension centrale
$$ (84) \quad 1 \to \mu \to \varepsilon \xrightarrow{(\varepsilon_3, \varepsilon_2)} \mu \times \mu \to 1 $$
$$ \simeq S_0^\pi \Gamma_{\rho,\sigma}^! / S_0^{\pi_0} \Gamma_{\rho,\sigma}^0 $$
s'insérant dans le diagramme d'extensions,

$$ (85) $$
\begin{tikzcd}
1 \ar[r] & \mu \ar[r] \ar[d, equal] & \varepsilon \ar[r] \ar[d, hook] & \mu \times \mu \ar[r] & 1 \\
1 \ar[r] & \mu \ar[r] \ar[d] & S_0^\pi \Gamma_{\rho,\sigma}^! \ar[r] \ar[d] & Z_{\rho,\sigma}^! \simeq M[0] \ar[u, "(\varepsilon_3{,}\varepsilon_2)"'] \ar[d] & \\
1 \ar[r] & L_\rho \times L_\sigma \ar[r] & S_0 \Gamma_{\rho,\sigma} \ar[r] & Z_{\rho,\sigma} \ar[r] & 1
\end{tikzcd}

L'intérêt de (85) est qu'il permet de récupérer l'extension $S_0^\pi \Gamma_{\rho,\sigma}^!$ de $Z_{\rho,\sigma}^!$ par $L_\rho \times L_\sigma$, par la loi usuelle - contravariante dans le foncteur "extensions", via

$$ (85') \quad Z_{\rho,\sigma}^! \xrightarrow{(\varepsilon_3, \varepsilon_2)} \mu \times \mu \quad , \quad \mu \hookrightarrow L_\rho \times L_\sigma $$
$$ -1 \mapsto (\rho^3 \sigma^2) $$
("hom. diagonal")

ainsi que la sous-extension $S_0^\pi \Gamma_{\rho,\sigma}^!$ de celle ci (de $Z_{\rho,\sigma}^!$ par $\mu$).

Etudions d'abord l'ordre des éléments de $\varepsilon$ - ils sont évidemment diviseurs de 4, sont ils tous d'ordre 2 (... 1)? Soit

[MARGIN: NB Pour l'intuition dans ce qui suit, $\underline{u} = (u,a,b) \in S_0^\pi \Gamma^!$ (dans $Z, \hat{N}_\rho^\pi, \hat{N}_\sigma^\pi$) i.e. $u \equiv b \, (\hat{\Pi}_0)$ et $u \equiv \varepsilon_3[\sigma] a \, (\hat{\Pi})$ (87). On a $u^2 \in S_0^{\pi_0} \Gamma_0$. Evidemment $\underline{u} = (u^2, a', b') \in S_0^\pi \Gamma_0$]

\newpage

%% ===== Page 602 =====

il reste à voir si $\underline{u}^2 \in S_0^{\Pi_0}\Gamma^!$, i.e. si ...
[unclear: ~~nous savions $b^2 \equiv b^2 \pmod{\Pi_0}$ et $u^2 \equiv a^2 \pmod{\hat{\Pi}}$~~]

— ce qui est clair et exprime simplement
$\underline{u}^2 \in S_0^\Pi \Gamma$ — mais même
(88) $$u^2 \equiv a^2 \pmod{\hat{\Pi}_0} \quad (?)$$

La relation 
$$u \equiv \varepsilon_3[\sigma].a \pmod{\hat{\Pi}}$$
s'écrit aussi
$$u \equiv \pm \varepsilon_3[\sigma].a \pmod{\hat{\Pi}_0}$$
ce qui implique
$$u^2 \equiv \varepsilon_3[\sigma] a \varepsilon_3[\sigma] a \pmod{\hat{\Pi}_0}$$
ce qu'on peut écrire
(89) $$u^2 \equiv \varepsilon_3 \varepsilon_3[\sigma](a) . a = \begin{cases} a^2 & \text{si } \varepsilon_3 = 1 \\ - \sigma(a).a & \text{si } \varepsilon_3 = -1 \end{cases} \pmod{\hat{\Pi}_0}$$

Donc (88) est vraie si $\varepsilon_3 = 1$, tandis que
pour $\varepsilon_3 = -1$, elle équivaut à
$$\sigma(a) \equiv -a \pmod{\hat{\Pi}_0}$$
i.e.
$$[\sigma, a] = -1 \quad \text{dans} \quad \mathcal{C} / \hat{\Pi}_0$$
i.e.
$$\varepsilon_2(a) = -1 \quad \text{i.e.} \quad \varepsilon_2(x)\varepsilon_3(x) = 1 \ .$$
(où $\varepsilon_2(a) = \varepsilon_2(x)$)

Donc on trouve :

Proposition Soit $x \in \mathcal{E}$ (cf (84)). Pour que
l'— on ait $x^2 = 1$, il f. et s. qu'—
on ait $\varepsilon_3(x) = 1$ ou $\varepsilon_2(x)\varepsilon_3(x) = 1$, i.e. que
~~l'on ait $x \in \mathcal{E}' \cup \mathcal{E}''$~~
(90) $$x \in \mathcal{E}' \cup \mathcal{E}'''$$

\newpage

%% ===== Page 603 =====

où
(91) $\mathcal{E}' = \{x \in \mathcal{E} \mid \varepsilon_3(x) = 1\}$, $\mathcal{E}''' = \{x \in \mathcal{E} \mid \varepsilon_2(x)\varepsilon_3(x) = 1\}$

Corollaire 1 $\mathcal{E}$ contient exactement deux éléments
d'ordre 4, qui sont les deux éléments de $\mathcal{E}$ au
dessus de l'élément $(-1, 1)$ de $\mu \times \mu$.
Ces deux éléments sont inverses l'un de
l'autre, et sont les deux générateurs du
seul sous-groupe cyclique d'ordre 4 de $\mathcal{E}$,
qui n'est autre que 
(92) $\mathcal{E}'' = \{x \in \mathcal{E} \mid \varepsilon_2(x) = 1\}$

Corollaire 2 L'extension $\mathcal{E}$ n'est pas scindée
car étant centrale, elle serait isomorphe
à $(\mu \times \mu) \times \mu$, dont tous les éléments seraient d'ordre
qui divise 2.

Nous pouvons regarder maintenant
$\mathcal{E}$ comme extension de $\mathcal{E}/\mathcal{E}'' \simeq \mathbb{Z}/2\mathbb{Z} = \mu$ par
$\mathcal{E}'' (\simeq \mathbb{Z}/4\mathbb{Z})$. Mais il faudrait
d'abord expliciter l'opération de $\mathcal{E}/\mathcal{E}'' = \mu$
sur $\mathcal{E}''$ — l'élément $-1$ (d'ordre 2) de $\mu$
induit un autom. involutif de $\mathcal{E}''$
d'ailleurs $\operatorname{Aut}(\mathcal{E}'') \simeq (\mathbb{Z}/4\mathbb{Z})^* \simeq \{\pm 1\}$, donc
l'action est la multiplication par un
signe $\pm 1$ qu'il faut déterminer. On va
prendre l'élément de $\mathcal{E}$ provenant
de $x = (u_0=\tau, a_0=\tau', b_0=\tau) \in \tilde{S}_0^{\bullet}\Gamma_0'''$ (cf (70))
pour le faire opérer sur l'image d'un
$u = (u, a, b) \in \tilde{S}_0^{\bullet}\Gamma_0'''$ avec $\varepsilon_3(u)=-1, \varepsilon_2(u)=1$
(93)
(qui soit un générateur de $\mathcal{E}''$). 
[MARGIN: $(u \equiv b \mod \hat{\Pi}_0$
$u \equiv \sigma a \mod \hat{\Pi})$]

\newpage

%% ===== Page 604 =====

On aura donc
$\underline{x}(\underline{u}) = (\tau(u), \tau(a), \tau(b)) \equiv$ [crossed out: $(1, \varepsilon, \varepsilon)(u, a, b)$]
$(u, \varepsilon a, \varepsilon b)$

mod $\tilde{S}_0^{\pi_0} \Gamma^0$, avec un signe $\varepsilon \in \{\pm 1\}$ à déterminer.

La [unclear: condition] de [unclear: a-groupe] s'écrit
$(u \tau(u)^{-1}, \varepsilon a \tau(a)^{-1}, \varepsilon b \tau(b)^{-1}) \in \tilde{S}_0^{\pi_0} \Gamma^0$
qui, mod $\tilde{S}_0^{\pi_0}$ est évident (et clair
quel que soit le choix du signe
$\varepsilon$), c'est la condition de congruence qui
distingue $\tilde{S}_0^{\pi_0} \Gamma^1$ de $\tilde{S}_0^{\pi} \Gamma^1$ qui est
essentielle, à savoir

(94) $u \tau(u)^{-1} \equiv \varepsilon a \tau(a)^{-1} \quad (\hat{\Pi}_0)$ soit

[crossed out: $u \tau u^{-1} \equiv \varepsilon a \tau(a)^{-1}$]

La difficulté avec cette approche vient
du fait que les seuls éléments de
$M$ qu'on sache expliciter sont ceux qui
proviennent de $G \simeq \operatorname{Gl}(2, \mathbb{Z})$ lui-même, or
pour ceux-ci on a $\varepsilon_2(u)\varepsilon_3(u) = 1$, alors
que le $u$ cherché ci-dessus doit vérifier
$\varepsilon_2(u) = 1, \varepsilon_3(u) = -1$ donc $\varepsilon_2(u)\varepsilon_3(u) = -1$. On
préfère donc pour travailler dans le
[crossed out: $\tilde{S}_0^{\pi_0}$] $\operatorname{Gl}(2, \hat{\mathbb{Z}})$ de $M$, et
y prendre

(95) $u = \begin{pmatrix} \mu & 0 \\ 0 & 1 \end{pmatrix}$ avec $\mu \in \hat{\mathbb{Z}}^*, \begin{cases} \mu \equiv -1\ (3) \\ \mu \equiv 1\ (4) \end{cases}$ i.e. $\mu \equiv 5\ (12)$

tandis que $\tau$ correspond à $c$
[unclear: avec] $\tau = \begin{pmatrix} -1 & 0 \\ 0 & 1 \end{pmatrix}$

\newpage

%% ===== Page 605 =====

on aura de toutes façons
$$u \tau u \tau^{-1} = 1$$
soit d'autre part dans $\operatorname{Gl}(2, \hat{\mathbb{Z}})$
(95) $\begin{cases} a = \tau' a_0 & a_0 = c \rho + d \quad c^2+cd+d^2 \in \hat{\mathbb{Z}}^* \\ & \tau' = \begin{pmatrix} 0 & 1 \\ 1 & 0 \end{pmatrix}, \rho = \begin{pmatrix} 0 & -1 \\ 1 & 1 \end{pmatrix} \end{cases}$

d'où
(97) $a = \begin{pmatrix} 0 & 1 \\ 1 & 0 \end{pmatrix} \begin{pmatrix} d & -c \\ c & c+d \end{pmatrix} = \begin{pmatrix} c & c+d \\ d & -c \end{pmatrix}$

et
[MARGIN: un peu plus bas ...] (98) $\mu = -(c^2+cd+d^2)$

On va travailler dans le sous-groupe
$\hat{\Pi}_0$ dans $\operatorname{Sl}(2, \hat{\mathbb{Z}})$ engendré par
$-l_{\infty} = \begin{pmatrix} 1 & -2 \\ 0 & 1 \end{pmatrix}, -l_1 = -\begin{pmatrix} 1 & 0 \\ 2 & 1 \end{pmatrix}, -l_0 = -\begin{pmatrix} 3 & -2 \\ 2 & -1 \end{pmatrix}$

et dans le groupe quotient de
$\operatorname{Gl}(2, \hat{\mathbb{Z}})$ par $\hat{\Pi}_0$, qui s'identifie au
$\hat{\mathbb{Z}}^* \times_{\{\pm 1\}} \operatorname{Sl}(2, \hat{\mathbb{Z}})/\hat{\Pi}_0 \simeq \operatorname{Sl}(2, \mathbb{Z})/\Pi_0 \simeq \left\{ (\rho, \sigma) \mid \begin{smallmatrix} \rho^6 = \sigma^4 = 1 \\ \rho^3 = \sigma^2 \\ \sigma \rho \sigma^{-1} = \rho^{-1} \end{smallmatrix} \right\},$
comme un produit $1/2$ direct quotient :
$$T_0 = \left\{ \begin{pmatrix} \zeta & 0 \\ 0 & 1 \end{pmatrix} \mid \zeta \in \hat{\mathbb{Z}}^* \right\} \overset{\text{det}}{\simeq} \hat{\mathbb{Z}}^*$$
opérant par l'intermédiaire de $\varepsilon_2(\bar{\zeta})$
cf p. 593, 599 - formules (34)(35)(35 bis). D'où
à poser $\sigma a \equiv u \quad (\hat{\Pi}_0) \quad \text{où } \sigma = \begin{pmatrix} 0 & -1 \\ 1 & 0 \end{pmatrix}$
(99)
d'où
(100) $\sigma a = \begin{pmatrix} -d & -c \\ -c & c+d \end{pmatrix}$

\newpage

%% ===== Page 606 =====

La relation 199), [rayé: il suffirait d'avoir $\det \sigma a \equiv \det u \ (\operatorname{Sl}(2,\hat{\mathbb{Z}}))$ i.e. $\det \sigma a = \det u$ i.e. (98)] s'écrit

(102) $$ \sigma a \equiv u \quad (2) $$

soit et cette double congruence s'écrit

(103) $$ c \equiv 0 \quad (2) \quad d \equiv 1 \quad (2) \qquad \left( -d \equiv \mu \quad (2) \text{ en conséquence on a } \mu \equiv 1 \equiv -1 \ (2) \right) $$

[Texte rayé :
NB On a pu [unclear] la 
[unclear] on a $\sigma a \equiv u \ (4)$ , [unclear]
[unclear] $-d \equiv -1 \ (4)$, 
qui est contraire à l'hypothèse (98) $\mu \equiv 1 \ (4)$.
[unclear]
(102 bis) $$ \sigma a \equiv u \quad (2) $$
(103 bis) $$ c \equiv 0 \quad (2), \quad d \equiv 1 \quad (2) $$
qui entraîne $-d \equiv \mu \quad (2)$
par conséquent on a $\mu \equiv 1 \equiv -1 \ (2)$). Ce
Revenant à la relation (94), on trouve
]

ceci :

Lemme. Soient $c,d \in \hat{\mathbb{Z}}$ satisfaisant et
(104) $$ \mu = -(c^2+cd+d^2) \in \hat{\mathbb{Z}}^* $$
[MARGIN: NB On a au reste $\mu \equiv 1 \ (3)$]

(105) $$ \mu \equiv 1 \quad (4) \quad \text{(i.e. } c^2+cd+d^2 \equiv -1 \ (4) \text{)} $$

(106) $$ c \equiv 0 \quad (2), \quad d \equiv 1 \quad (2) $$

Soit 
(107) $$ a = \tau'(c\rho+d) = \begin{pmatrix} -c & c+d \\ d & c \end{pmatrix} \in \operatorname{Gl}(2,\hat{\mathbb{Z}}) $$
(noter de déterminant $\mu$). Alors on a une congruence

(108) $$ a \tau'(a)^{-1} \equiv \varepsilon \quad \bmod \hat{\Pi}_0 \subset \operatorname{Gl}(2,\hat{\mathbb{Z}}) $$
où le signe $\varepsilon \in \{ \pm 1 \}$ ne dépend pas des choix faits pour $c, d$.

\newpage

%% ===== Page 607 =====

607

C'est le moment de le calculer ! Reprenons

$$a_0 = c \rho + d \in Z_{\rho}$$

$$a = \tau' a_0 , \quad \tau'(a) = a_0 \tau' , \quad a \tau'(a)^{-1} = \tau' a_0 \tau'^{-1} a_0^{-1} = \tau'(a_0) a_0^{-1}$$

où

[MARGIN:
On prend
dans $M_2(\hat{\mathbb{Z}})$
$\tau u = (Tr u) - u$
d'où $\tau \rho = \rho^{-1}$
et $\tau u = u^{-1}$ si $\det u = 1$
(plus général^t ...)]

$$ \begin{cases} \tau'(a_0) = c \rho^{-1} + d = {}^t a_0 \\ a_0^{-1} = {}^t a_0 / \det a_0 = - \frac{1}{\mu} (c \rho^{-1} + d) \end{cases} $$

donc

$$ a \tau'(a)^{-1} = \tau'(a_0) a_0^{-1} = - \frac{1}{\mu} (c \rho^{-1} + d)^2 = \frac{1}{c^2 + c d + d^2} (c \rho^{-1} + d)^2 $$

Or
$$ (c \rho^{-1} + d)^2 = \underset{-\rho}{c^2 \rho^{-2}} + \underset{1-\rho}{2cd \rho^{-1}} + d^2 = -(c^2+2cd)\rho + (2cd+d^2) $$

Or comme
$$ \begin{cases} c \equiv 0 \pmod 2 \\ d \equiv 1 \pmod 2 \end{cases} \implies \begin{matrix} c^2, 2cd \equiv 0 \pmod 4 \\ d^2 \equiv 1 \pmod 4 \end{matrix} $$
donc

(*) [unclear: crossed out text]
$$ (c \rho^{-1} + d)^2 \equiv 1 \pmod 4 $$
d'autre part
[unclear: crossed out text]
et comme par hypothèse
$$ c^2 + cd + d^2 \equiv -1 \pmod 4 $$

on aura

(109) $$ a \tau'(a)^{-1} = \frac{1}{c^2+cd+d^2} (c \rho^{-1} + d)^2 \equiv -1 \pmod 4 $$

Donc on a prouvé en fait

Proposition. Soient $c,d \in \hat{\mathbb{Z}}$ tels que
$\mu = c^2+cd+d^2 \in \hat{\mathbb{Z}}^\times$, soit
$$ a = \tau' (c \rho + d) = \begin{pmatrix} -c & c+d \\ d & c \end{pmatrix} \in \operatorname{Gl}(2, \hat{\mathbb{Z}}) , $$

[MARGIN: NB $\sigma = \begin{pmatrix} 1 & 0 \\ 0 & -1 \end{pmatrix}$]
[unclear: crossed out text]
$$ a \equiv \begin{pmatrix} 0 & 1 \\ 1 & 0 \end{pmatrix} \pmod 2 \quad \text{i.e.} $$
$$ c \equiv 0 \pmod 2, \quad d \equiv 1 \pmod 2 . $$

\newpage

%% ===== Page 608 =====

Alors on a (pour $\tau' = \begin{pmatrix} -1 & 0 \\ 0 & 1 \end{pmatrix}$)
$$a \tau'(a)^{-1} \equiv \varepsilon_2(\frac{c}{2}) \quad (4)$$
(= fonction [unclear] de la congruence mod $\hat{\pi}_0$ ).

Corollaire L'extension $\mathcal{E}$ de $\mu \times \mu$ par $\mu$ (84),
où on a $\mu \hookrightarrow \mathcal{E}'' \simeq \mathbb{Z}/4\mathbb{Z}$, est non commutative.

Mais voici une façon de le voir sans
calcul. Si l'extension était commutative,
sa classe d'isomorphie serait donnée par un $Ext^1_\mathbb{Z}$ (ou un
$Ext^1_{\mathbb{Z}/4\mathbb{Z}}$ ). Or le foncteur $Ext(-, \mu)$
est additif. Comme la restriction de
l'extension aux sous-groupes $\mathcal{E}'/\mu$ et $\mathcal{E}''/\mu$ de
$\mu \times \mu$ est triviale, et que $\mu \times \mu$ est la somme
directe de ces sous-groupes (qui sont resp.
l'axe et le sous-groupe diagonal?), il s'ensui-
vrait que l'extension serait triviale, ce
qu'on a vu n'est pas le cas.

La structure du groupe $\mathcal{E}$
s'explicite aisément, fix. un tel généra-
teur de $\mathcal{E}/\mathcal{E}'' \simeq \mu \vdash \mathcal{E}'' \simeq \mathbb{Z}/4\mathbb{Z}$. Soit
choisissons un relèvement quelconque de $\mathcal{E} \setminus \mathcal{E}''$;
il est d'ordre 2 (les seuls éléments d'ordre 4
étant dans $\mathcal{E}''$) de sorte que l'on a

\newpage

%% ===== Page 609 =====

(110) $$ \varepsilon \simeq \{1, x\} \cdot \varepsilon'' $$
($1/2$ direct)

Notons que si $x$ opère trivialement sur $\varepsilon''$,
c.à.d. si $\varepsilon$ était isomorphe au produit,
alors pour un él. $u$ de $\varepsilon''$ d'ordre 4,
$x.u$ serait d'ordre 4, ce qui est absurde
(troisième démonstration du fait que
$\varepsilon/\varepsilon'' \simeq \mu$ opère non trivialement sur $\varepsilon''$). Or
du seul fait que l'op. de $\{1, x\}$ sur $\varepsilon''$ est
non triviale, elle est connue,
la structure de $\varepsilon$ est
déterminée en fait --

(111) $$ \varepsilon \simeq D_4 $$

Il y a un choix privilégié de $x$,
en prenant pour $x$ l'él.

(112) $$ \begin{cases} \tau_\xi = \tau_{S_0 P} \pmod{\widehat{S_0 P}^0} , \text{ où bien sûr} \\ \tau_{S_0 P} = (\tau, \tau', \tau) \end{cases} $$

de sorte qu'on a l' -- une représentation
canonique

(113) $$ \varepsilon \simeq \{1, \tau_\xi\} \cdot \varepsilon'' $$
produit $1/2$ direct

[MARGIN: $\varepsilon'' \simeq \mathbb{Z}/4\mathbb{Z}$
$\tau_\xi(u) = u^{-1}$
pour $u \in \varepsilon''$.]

L'hom. $\varepsilon \to \mu \times \mu$ se récupère par
ses restrictions. - [unclear: on a d'une part]

(114) $$ \mu \simeq \{1, \tau_\xi\} \hookrightarrow \mu \times \mu \quad \tau_\xi \mapsto (-1, -1) $$
[MARGIN: (hom diagonal)]

$$ \varepsilon'' \to \mu \times \mu $$
l'unique
hom. non trivial $\varepsilon'' \twoheadrightarrow \mu \hookrightarrow \mu \times \mu$
déduit de l'isom $\varepsilon''_2 = \varepsilon''/\varepsilon''^2 \stackrel{\sim}{\rightarrow} \mu$

\newpage

%% ===== Page 610 =====

Le fait que $\mathcal{E}$ ne soit pas commutative
s'exprime aussi par le fait que l'hom.
(115) $$ \mathcal{E}_{ab} \xrightarrow{\sim} \mu \times \mu $$
est un iso, i.e. que le sous-groupe noyau
de $\mathcal{E} \to \mu \times \mu$ est [unclear: juste] le sous-groupe
des commutateurs.

On a [unclear: vu] que l'extension $\mathcal{E}$ de $\mu \times \mu$
n'est pas scindée, mais qu'en est-il de
son image inverse, l'extension $\tilde{S}_0^{!\Gamma}$ de
$\hat{Z}' \simeq \hat{M}[0]$ par $\mu$ ? A priori, on est
assuré qu'il n'y a pas de scindage de
cette extension qui prolonge le scindage
canonique de $\Delta$ dont on dispose au-dessus
de $Z^\circ = \operatorname{Ker}(\hat{Z}' \xrightarrow{(\varepsilon_3, \varepsilon)} \mu \times \mu)$. Je présume
que ce $\tilde{S}_0^{!\Gamma}$, extension de $\hat{Z}'$ par $\mu$,
est non triviale, et que l'extension qu'elle
induit sur son gr.
$$ \Gamma'_Q = \Gamma_Q \underset{[unclear]}{\cap} \hat{Z}' \subset \hat{M}(0) \subset \hat{M}'[0] = \hat{Z}' $$

Je serais intrigué d'identifier, l'ext.
plus précisément, l'extension pour $\Gamma_Q$ qu'
elle induit, qui correspond à une
classe canonique
(115') $$ c \in H^2(\text{Spec } \mathbb{Q}, \mu_2) \subset \text{Br}(\mathbb{Q}) $$

\newpage

%% ===== Page 611 =====

sauf erreur on aura

(115) $$c = \{ "i" , "j" \} , \quad \xi_{"i"} \smile \xi_{"j"} \in H^2(\mathbb{Q}, \mu_2)$$
(cup-produit)

sont les éléments qui décrivent respectivement
les extensions $\mathbb{Q}(\sqrt[3]{1}) = \mathbb{Q}(\sqrt{-3})$, et
$\mathbb{Q}(\sqrt{3})$ - contenues l'une et l'autre ,
ainsi que $\mathbb{Q}(\sqrt{-1}) = \mathbb{Q}(i)$ , dans l'extension
biquadratique $\mathbb{Q}(i, j) = \mathbb{Q}(\sqrt[12]{1})$, le corps
des racines 12ièmes de l'unité dont le groupe
de Galois est $(\mathbb{Z}/12\mathbb{Z})^* \simeq (\mathbb{Z}/3\mathbb{Z})^* \times (\mathbb{Z}/4\mathbb{Z})^* = \mu \times \mu$.

Il doit être trivial - quand on connait un
peu les fondements - que $c$ ainsi
explicité est $\neq 0$, et n'est pas non plus
scindé en passant à $\mathbb{Q}(i)$, le troisième
larron.

Considérons maintenant le
$\mu$-groupe central

[MARGIN: NB $Z \simeq M_{\Gamma_Q}$
$\Gamma \simeq \Gamma_{\tilde{Q}}$]

(116) $$\mu \xrightarrow{i_{S, \Gamma}} S_0 \Gamma_{pr}^! = Z^!_{pr} \times N_{\hat{p}}^* \times N_{\hat{\sigma}}^*$$
$$\varepsilon \longmapsto (1, \varepsilon, 1)$$

à valeurs dans
$$\operatorname{Ker}(S_0 \Gamma_{pr}^! \to Z^!_{pr}) \simeq \mu_q \times \mu_\sigma ,$$

donne un hom. canonique de relèvement

(117) $$M_{/\sigma}^! Z_{pr}^! \simeq S_0^\pi \Gamma_{pr}^! / \mu \longrightarrow S_0 \Gamma_{pr}^!(\mu) ,$$

\newpage

%% ===== Page 612 =====

[MARGIN: Ceci mérite d'être explicité en un théorème récapitulatif]
qui s'explicite au travers des deux hom.
canoniques fondamentaux

(118) $$Z_{p,\sigma}^{!} \xrightarrow{\underline{a}} N_{\rho}^{*}/\mu \ , \ Z_{p,\sigma}^{!} \xrightarrow{\underline{b}} N_{\sigma}^{*}$$
$$\simeq$$
$$M^{\Gamma_Q}$$
$$\simeq$$
$$\tilde{\Gamma}_Q \cdot \underline{\pi}$$

où, je rappelle
(119) $$ \begin{cases} N_{\rho}^{*} = M = N_{1,1} \simeq \tilde{\Gamma}_{Q} \cdot \mathfrak{S}^+_{0(3, \mathbb{Z})^\wedge} \\ \parallel \\ \{u \in M \mid u(\rho) = \rho^{\varepsilon_3(u)}\} \\ \\ N_{\sigma}^* \subset M \\ \parallel \\ \{u \in M \mid u(\sigma) = \sigma^{\varepsilon_2(u)} \quad (= \varepsilon_1(u)\sigma ) \} \end{cases} $$

Les applications $\underline{a}, \underline{b}$ sont données par
les formules (cf 51)
(120) $$ \begin{cases} \underline{a}(u) = a = \varepsilon_3[2'] \alpha^{-1} u \quad , \quad \underline{b}(u) = b = \beta^{-1} u \qquad i.e. \\ \alpha = u a^{-1} \varepsilon_3[2'] \quad , \quad \beta = u b^{-1} \end{cases} $$
où $\alpha \in \Pi_{0 \tau}^{\wedge}$, $\beta \in \Pi_0^{\wedge}$ sont caractérisés par
(121) $$ u(\rho) = int(\alpha)(\rho), \quad u(\sigma) = \varepsilon_2(\sigma) int(\beta)(\sigma) $$

Les quantités [unclear: associées à un terme] $u \in Z^!$ ca-
ractérisent ce dernier aux signes près par les équations
(122) $$ u \equiv b \pmod{\Pi_0^\wedge} \quad , \quad u \equiv \varepsilon_3[2'] a \pmod{\Pi_\tau^\wedge} $$
qui impliquent à fortiori [unclear]
que [unclear]

\newpage

%% ===== Page 613 =====

(123)
\begin{tikzcd}
& N^*_\rho \ar[dr, "\delta_\rho"] \\
Z^!_{\rho,\sigma} \ar[ur, dashed] \ar[r, dashed, "\delta_Z"] \ar[dr, dashed] & \dots & \gamma_{\rho,\sigma} \simeq \tilde{\Gamma}_Q \\
& N^*_\sigma \ar[ur, "\delta_\sigma"']
\end{tikzcd}

Considérons maintenant les deux
[unclear: crossed out text]
Il serait temps
de rectifier la
confusion (p. 594') entre $Z = Z_{\rho,\sigma} \subset \tilde{M}_{0,3}^\sim$
et $Z^! = Z^!_{\rho,\sigma} = Z \cap \tilde{M}_{0,3}^!$. On a en fait
$$ Z_{\rho,\sigma} \simeq \tilde{M}(0) \cdot \tilde{L}_0^\otimes \subset \mathcal{M} $$
(124) \hfill (où $\tilde{L}_0^\otimes$ est un $\frac{1}{2}$-tour)
$$ \searrow \tilde{M}_{0,3}^\sim(0) \cdot \tilde{L}_0^{\otimes_{0,3}} \subset \tilde{M}_{0,3}^\sim $$

(125) $$ Z^!_{\rho,\sigma} \simeq Z_{\rho,\sigma} \cap M^! = M(0) \cdot L_0^\pi \subset M $$
\hfill (où $L_0^\pi$ est un $\pi$-tour)
$$ \searrow M_{0,3}(0) \cdot L_0^{\pi_{0,3}} \subset M_{0,3} $$

Le lien entre les deux est donné
par l' e.g. [unclear]

(126)
\begin{tikzcd}
1 \ar[r] & L_0^\pi \ar[r] \ar[d, hook] & Z^!_{\rho,\sigma} \ar[r] \ar[d, hook, "p"] & \gamma_{\rho,\sigma} \ar[r] \ar[d, equal] & 1 \\
1 \ar[r] & \tilde{L}_0^\otimes \ar[r] & Z_{\rho,\sigma} \ar[r] & \gamma_{\rho,\sigma} \ar[r] \ar[d, equal] & 1 \\
& & & \simeq M(0) \simeq \tilde{\Gamma}_Q
\end{tikzcd}

où $L_0^\pi \subset \tilde{L}_0^\otimes$
L'inclusion d'indice 2 de
(127) $$ L_0^\pi = l_0^{\hat{\mathbb{Z}}} \subset l_0^{\frac{1}{2}\hat{\mathbb{Z}}} = \tilde{L}_0^\otimes $$

Bien entendu, c'est bien sur $Z^!_{\rho,\sigma}$, non sur
$Z_{\rho,\sigma}$, qu'est défini (117), (118) d'où (123).

\newpage

%% ===== Page 614 =====

Je m'intéresse à la restriction des
hom. (123) au sous-groupe $L_0^\pi$ de $Z_{\rho,\sigma}^!$
la concrétisation (122) montre que
l'hom. induit
$$L_0^\pi \longrightarrow N_\rho^*/\{\pm 1\}$$
est trivial (car son terme en $m$ ne
dépend que de $m \bmod \pi^\rho$), tandis que l'hom.
induit
$$L_0^\pi \longrightarrow N_\sigma^*$$
se factorise par $L_0^\pi/L_0^{\pi_0} \simeq \mathbb{Z}/2\mathbb{Z}$ (car
le terme en $m$ ne dépend que de
$m \bmod \pi_0^\sigma$). Sa valeur $b$ pour $l_0$ est
déterminée par les conditions
$$b \in N_\sigma^*, \quad b \equiv l_0 \pmod{\pi_0}$$
qui sont satisfaites évidemment par $b=-1$. Donc
il y a lieu d'introduire aussi, quitte
[unclear: à modifier un peu le diagr.] central
canonique
$$(128) \quad \mu \hookrightarrow N_\sigma^*$$
$$-1 \longmapsto -1$$
et une commutativité dans

(129)
\begin{tikzcd}
L_0^\pi \ar[r, hook] \ar[d] & Z_{\rho,\sigma}^! \ar[d] \\
L_0^\pi/L_0^{\pi_0} \simeq \mu \ar[r, hook] & N_\sigma^*
\end{tikzcd}

[unclear: Donnons un tronçon], à partir des données (123)

\newpage

%% ===== Page 615 =====

611

de $Z_{\rho,\sigma}^!$, par passage au quotient deux hom.
canoniques
$$ (130) \qquad \begin{cases} \gamma_{\rho,\sigma} \longrightarrow N_\rho^* / \{\pm 1\} \\ \Vert \ \tilde{\Pi}_Q \\ \gamma_{\rho,\sigma} \longrightarrow N_\sigma^* / \{\pm 1\} \end{cases} $$
qui sont en fait des sections des
$$ \begin{matrix} N_\rho^* / \{\pm 1\} \longrightarrow \gamma_{\rho,\sigma} \\ \Vert \ \tilde{\Gamma}_Q \simeq M(0) \\ N_\sigma^* / \{\pm 1\} \longrightarrow \gamma_{\rho,\sigma} \end{matrix} $$
de sorte qu'on trouve des décompositions
$$ (131) \qquad \begin{cases} N_\rho^* / \pm 1 \simeq \gamma_{\rho,\sigma} \ltimes S Z_\rho / \{\pm 1\} \\ \qquad \qquad \quad = L_\rho / \{\pm 1\} \simeq \mathbb{Z}/3\mathbb{Z} \\ N_\sigma^* / \pm 1 \simeq \gamma_{\rho,\sigma} \cdot S Z_\sigma / \pm 1 \\ \qquad \qquad \quad = L_\sigma / \pm 1 \simeq \mathbb{Z}/2\mathbb{Z} \end{cases} $$
(produits semi-directs). L'opération de
$\gamma_{\rho,\sigma}$ sur $L_\rho / \{\pm 1\} \simeq \mathbb{Z}/3\mathbb{Z}$ est donnée par le
caractère quadratique $\varepsilon_\rho$, celle sur $L_\sigma / \{\pm 1\} \simeq \mathbb{Z}/2\mathbb{Z}$
est donnée par le caractère quadratique $\varepsilon_\sigma$.

La symétrie de présentation des deux
lignes (130) est un peu bidon - p. ex.
l'hom. $\tilde{\Pi}_Q = \gamma_{\rho,\sigma} \longrightarrow N_\sigma^* / \{\pm 1\}$ se remonte
de façon évidente en
(132)
\begin{tikzcd}
\tilde{\Pi}_Q = \gamma_{\rho,\sigma} \ar[r] \ar[dr, "\sim"'] & N_\sigma^* \\
& M(0) \subset M[0] = Z^! \ar[u, hook]
\end{tikzcd}

\newpage

%% ===== Page 616 =====

Une autre façon de le voir est de noter que

(132) $$ \gamma_{\rho, \sigma} \simeq \mathcal{M}_{0}^{!} [0] / L_{0}^{\pi_0} $$
où $$ \mathcal{M}_{0}^{!} [0] \overset{dfn}{=} \mathcal{M}(0) \cdot L_{0}^{\pi_0} \subset \mathcal{M}^{!} [0] = \mathcal{M}(0) \cdot L_0^{!} $$
[MARGIN: sous-groupe d'indice 2]

or sur $L_0^{\pi_0}$, l'hom. induit par
$$ \mathcal{M}^{!} [0] = Z_{\rho,\sigma}^! \to \mathcal{N}_{\sigma}^* $$
est trivial. Ainsi on a un isomorphisme
semi-direct

(133) $$ 1 \to L_{\sigma} \to \mathcal{N}_{\sigma}^* \to \underset{\underset{\Pi_{Q}^{\sim}}{\parallel}}{\gamma_{\rho,\sigma}} \to 1 $$

en un produit $1/2$-direct

$$ \mathcal{N}_{\sigma}^* \simeq \gamma_{\rho,\sigma} \cdot L_{\sigma} $$

qu'il vaut mieux écrire en concrétisant
la sous-groupe section en vu que
sous-groupe de $G_{\rho,\sigma} \simeq \mathcal{M}$, soit donc

(134) $$ \mathcal{N}_{00}^{! *} \overset{dfn}{=} \operatorname{Im}(\mathcal{M}_0^{!} [0] \to \mathcal{N}_{\rho}^*) $$
$$ = \operatorname{Im}(\mathcal{M}[0] \hookrightarrow \mathcal{N}_{\rho}^*) $$
$$ = \{ \beta^{-1} u \mid u \in \mathcal{M}(0) , \beta \in \hat{\Pi}_0 \text{ tels que } u(\sigma) = \varepsilon_3(u) sgn_{\rho}(\beta) (\sigma) \} $$
$$ = \{ v \in \mathcal{N}_{\sigma}^* \mid \exists \beta \in \hat{\Pi}_0 \text{ tel que } \beta v \in \mathcal{M}[0] = \operatorname{Norm}_{\mathcal{M}}(L_0) \} $$

[MARGIN: NB On définit $\mathcal{M}^!$ comme l'image inverse de $\mathcal{M}_3^!$ dans le noyau de l'opération canonique $\mathcal{M} \overset{\varepsilon_3}{\to} \hat{\mathbb{Z}}^* / \pm 1 \simeq \operatorname{Aut}(\hat{\mathbb{Z}} / \pm 1)$]

Posant $dfn$

(135) $$ \mathcal{M}_0^! = \mathcal{M}(0) \cdot \Pi_0 \underset{\text{indice 2}}{\subset} \mathcal{M}^! = \mathcal{M}(0) \cdot \Pi $$

\newpage

%% ===== Page 617 =====

de sorte qu'on a

(135) $$ \mathcal{M}' \simeq \mathcal{M}'_0 \times \mu $$ ,

[MARGIN: NB on a un $\mathcal{M}'_0 \simeq \mathcal{M}_{0,3}$]

on a donc

(137) $ \left\{ \begin{array}{l} \mathcal{N}^{*!}_0 = \mathcal{N}^*_0 \cap \mathcal{M}'! \\ \mathcal{N}^{*!}_{0,0} \overset{dfn}{=} \mathcal{N}^*_0 \cap \mathcal{M}'_0 \end{array} \right. $ et on trouve que $ \mathcal{N}^{*!}_0 = \mathcal{N}^{*!}_{0,0} \times \mu $ ,

(138) $$ \mathcal{N}^{*!}_{0,0} \overset{\sim}{\longrightarrow} \gamma_{\bar{p}, \sigma} = \widetilde{\Gamma}_Q $$

Posons

(139) $$ \mathcal{M}(1/2) = \mathcal{N}^{*!}_{0,0} \quad \text{cf (137)} $$

on aura donc que $\mathcal{M}(1/2)$ est une section pour $\mathcal{N}^*_0$ extension de $\gamma_{\bar{p}, \sigma}$ par $L_0$, et en fait pour $\mathcal{M}$ ext. de $\gamma_{\bar{p}, \sigma}$ par $\hat{\mathbb{Z}}^{+\wedge}$ on a des formules pour $\mathcal{M}'_0, \mathcal{M}'$

(140) $ \begin{cases} \mathcal{M}'_0 \simeq \mathcal{M}(1/2) \cdot \hat{\Pi}_0 \\ \mathcal{M}' \simeq \mathcal{M}(1/2) \cdot \hat{\Pi} \\ \mathcal{M} \simeq \mathcal{M}(1/2) \cdot \hat{\mathbb{Z}}^{+\wedge} \\ \mathcal{N}^*_0 \simeq \mathcal{M}(1/2) \cdot L_0 \end{cases} $ $\leftarrow$ attention, $\mathcal{M}'$ est inv. dans $\mathcal{M}$ et $\mathcal{M}'_0$ est inv. dans $\mathcal{M}'$, mais non dans $\mathcal{M}$ cf plus bas

[ Si on veut plus expliciter on introduirait

(141) $$ \mathcal{N}^{*!}_0 \overset{dfn}{=} \mathcal{N}^*_0 \cap \mathcal{M}'! = \{ u \in \mathcal{N}^*_0 \mid \exists \alpha \in \hat{\Pi}_0 \text{ t.q. } \underbrace{\alpha v}_{u} \in \mathcal{M}(0) \} $$

c'est à dire l'un des $v$ qu'on peut associer à un $u \in \mathcal{M}_0 \subset \mathcal{M}'_0[\sigma] = \mathbb{Z}^1_{\bar{p}, \sigma}$, à l'aide de $\alpha \in \hat{\Pi}_0$ tel que $\alpha^{-1} u \in \mathcal{N}^*_0$ i.e. $u(p) = \text{classe}(p^{\varepsilon_3})$ avec $\varepsilon_3 = \varepsilon_3(u) = \varepsilon_3(v)$. Mais comme $u(p)$ doit [unclear] l'image de $\mathbb{Z}_3^+ = \hat{\mathbb{Z}}^+ / \hat{\Pi}$ par $p$, il faut pour cela ait $\varepsilon_3 = 1$, et alors $\alpha$ est l'élément $\alpha(u)$ associé à $u$ et $v = \alpha^{-1} u$ est l'image de $u$ dans $\mathcal{N}^*_0$. Ainsi ... ]

\newpage

%% ===== Page 618 =====

trouve

Proposition L'hom. $\mathcal{N}_{\hat{p}}^{*!} \rightarrow \gamma_{\hat{p},\sigma} = \tilde{\Gamma}_{\mathbb{Q}}^\sim$ est
injectif, et a pour image $\gamma_{\hat{p},\sigma}^\circ = \{\gamma \in \gamma_{\hat{p},\sigma} \mid \varepsilon_3(\gamma)=1\}$

Corollaire $\mathcal{N}_{\hat{p}}^{*!} = Z_{\hat{p}}^{*!}$ i.e. $\mathcal{N}_{\hat{p}}^{*!} = \mathcal{Z}_{\hat{p}}^* \cap M_!$

Il est naturel de poser

(142) $\mathcal{N}_{\hat{p}\circ}^{*!} = Z_{\hat{p}\circ}^! = \mathcal{M}(-j)$.

Comme $\mathcal{N}_{\hat{p}}^{*!} \supset \mu$, on voit que l'on a

(142)
$$ \begin{cases} \mathcal{N}_{\hat{p}}^{*!} = \mathcal{N}_{\hat{p}\circ}^{*!} \dot{\times} \mu \quad (\text{donc } \mathcal{N}_{\hat{p}}^{*!} = Z_{\hat{p}}^!) \\ Z_{\hat{p}}^! = Z_{\hat{p}\circ}^! \times \mu \end{cases} $$

Le sous-groupe $Z_{\hat{p}\circ}^! = \mathcal{M}(-j)$ est une section
de $\mathcal{M}'$ sur $\gamma_{\hat{p},\sigma} \simeq \tilde{\Gamma}_{\mathbb{Q}}^\sim$, et définit
des factorisations de $\mathcal{M}'_!, \mathcal{M}', \mathcal{M}^*, Z_{\hat{p}}^!, Z_{\hat{p}}$
avec des décompositions en produit $1/2$ direct :

(143)
$$ \begin{cases} \mathcal{M}'_! \simeq \mathcal{M}(-j) \cdot \hat{\Pi}_0 \\ \mathcal{M}' \simeq \mathcal{M}(-j) \cdot \hat{\Pi} \\ \mathcal{M}^* \simeq \mathcal{M}(-j) \cdot \widehat{\mathfrak{S}} \\ Z_{\hat{p}}^! \simeq \mathcal{M}(-j) \times \mu \\ Z_{\hat{p}} \simeq \mathcal{M}(-j) \times L_p \end{cases} $$

Notons maintenant que
$$ \mu_2 = \{1, z'\} \subset \widehat{\mathfrak{S}} $$
normalise $L_{\hat{p}}$, donc $Z_{\hat{p}}$, donc $Z_{\hat{p}}^! = Z_{\hat{p}} \cap M_!$

\newpage

%% ===== Page 619 =====

613

Revenons au diagramme (130), pour
dire que l'hom
$$ \overline{\gamma}_{\rho,\sigma} \longrightarrow \mathcal{N}_{\rho}^*/\{\pm 1\} $$
$$ \| \wr $$
$$ \overline{\Gamma}_{\tilde{Q}}^\sim $$
ne se relève pas en un hom de $\overline{\gamma}_{\rho,\sigma}$
dans $\mathcal{N}_{\rho}^*$ — en d'autres termes, que
l'extension de $\overline{\gamma}_{\rho,\sigma}$ par $\{\pm 1\} = \mu$ qu'elle
définit n'est pas triviale. Pour identifier
cette extension, considérons le diagramme

[MARGIN: NB]
(144)
\begin{tikzcd}
\mathcal{M}'[0] = \mathcal{Z}_{\rho,\sigma}^! \ar[r, "can"] & \overline{\gamma}_{\rho,\sigma} = \mathcal{Z}_{\rho,\sigma}^! / L_0^\pi \ar[r, "\text{hom (130)}"] & \mathcal{N}_{\rho}^* / \{ \pm 1 \} \\
S_0 \Gamma_{\rho,\sigma}^! \ar[u, "can \ (u,v,b) \mapsto b"] \ar[ur, "can \atop {(\text{passage au} \atop \text{quotient})}"'] & & \\
S_0'' \Gamma_{\rho,\sigma}^! \ar[u, "incl."] \ar[ur, "incl."'] \ar[rr] & & \mathcal{N}_{\rho}^* \ar[uu]
\end{tikzcd}

qui montre que l'on a un diagramme
de carrés cartésiens

(145)
\begin{tikzcd}
\mathcal{M}(0) \ar[r, hook] \ar[dr, "\text{I}" description, phantom] & \mathcal{M}'[0] = \mathcal{Z}_{\rho,\sigma}^! \ar[r, "\text{épi}"] \ar[dr, "\text{II}" description, phantom] & \overline{\gamma}_{\rho,\sigma} \ar[r, hook] \ar[dr, "\text{III}" description, phantom] & \mathcal{N}_{\rho}^* / \{ \pm 1 \} \\
\mathcal{M}(0)^\sim \ar[u] \ar[r, hook] & S_0'' \Gamma_{\rho,\sigma}^! \ar[u, "2"'] \ar[r] & \overline{\gamma}_{\rho,\sigma}^\sim \ar[u] \ar[r, hook] & \mathcal{N}_{\rho}^* \ar[u, "2"']
\end{tikzcd}

où $\mathcal{M}(0)^\sim$ est défini comme l'extension de
$\mathcal{M}(0)$ par $\mu$ induite par l'ext. $S_0'' \Gamma_{\rho,\sigma}^!$ de $\mathcal{M}'[0]$,
qui s'identifie de même à l'image inv. de

\newpage

%% ===== Page 620 =====

l'ext. $\tilde{\gamma}_{p,\sigma}$ et $\gamma_{p,\sigma}$ par $\mu$, par l'isom. canon.
$\mathcal{M}(\infty) \overset{\sim}{\longrightarrow} \gamma_{p,\sigma}$. Donc l'ext. $\tilde{\gamma}_{p,\sigma}$ de $\gamma_{p,\sigma}$
par $\mu$ s'identifie nat. à l'ext. $\mathcal{M}(\infty)^\sim$ de
$\mathcal{M}(\infty)$ par $\mu$, dont on s'est convaincu plus
haut qu'elle n'était pas triviale.

Voici une autre façon d'obtenir
l'hom canonique $\gamma_{p,\sigma} \rightarrow N_p^\ast / \pm 1$,
et l'extension $\tilde{\gamma}_{p,\sigma}$ de $\gamma_{p,\sigma}$. Considérons
(146) $$N_p^{\ast !} = N_p^\ast \cap \mathcal{M}!$$

On a alors
(147) $$N_p^{\ast !} \cap \widehat{S}^+ = \operatorname{Ker}(N_p^{\ast !} \rightarrow \underbrace{\gamma_{p,\sigma}}_{\Pi_{\bar{Q}}}) = L_p \cap \Pi = \mu$$

d'autre part [unclear: que l'on a]
(148) $$N_p^{\ast !} \longrightarrow \gamma_{p,\sigma} \quad \text{est épi. (par } \widehat{H})$$

[unclear: Pour cela] il suffit de voir que $\forall$
$u \in \mathcal{M}(\infty)$, il existe un élément de
$N_p^{\ast !}$ qui lui est congru $\mod \widehat{S}^+$. Or on a
$u(\rho) = \text{int}(\alpha)(\rho)$ avec $\alpha \in \widehat{\Pi}_{\bar{Q}}$ [deleted text],

si $\varepsilon_3(u) = 1$ i.e. $\alpha \in \widehat{\Pi}$, $\alpha^{-1} u \in Z_p \subset N_p^\ast$
fait l'affaire. Si $\varepsilon_3(u) = -1$
on a $\alpha = \alpha_0 \tau$, $\alpha_0 \in \widehat{\Pi}$
et on a
$$u(\rho) = \alpha(\rho) = \alpha_0(\tau(\rho)) = \alpha_0(\sigma(\rho^{-1}))$$ donc
donc
$$(\sigma^{-1} \alpha_0^{-1} u)(\rho) = \rho^{-1} \quad \text{donc} \quad \sigma^{-1} \alpha_0^{-1} u \in N_p^{\ast \delta}$$
OK.

\newpage

%% ===== Page 621 =====

614

Donc on a une suite exacte

$$ (149) \quad 1 \to \mu \to \underset{(= \tilde{\gamma}_{\rho,\sigma})}{Z_{\rho}^{*!}} \to \underset{\simeq \tilde{\Pi_Q}}{\gamma_{\rho,\sigma}} \to 1 $$

qui fait de $Z_{\rho}^{*!}$ une $\mu$-ext. de $\gamma_{\rho,\sigma}$ ,

d'où un hom. can.

$$ \gamma_{\rho,\sigma} \simeq Z_{\rho}^{*!}/\mu \to N_{\hat{\rho}}^*/\mu $$

qui, bien entendu, n'est autre que
le premier hom. (140).

## 9) Récapitulation sur le cas universel.

Finalement, j'ai l'impression que
j'ai fait encore de longs détours et
contorsions, pour finalement un
peu me perdre dans les sables, et perdre un
peu le fil. Je vais donc récapituler
ce qui me semble essentiel.

[MARGIN: a]
(1) $Z_{\rho,\sigma} \simeq M(0) . L_0^\pi = M(0)$ [unclear: opérant sur $\hat{C}^+$ par opp. induite par l'identité]

en distinguant donc un sous-groupe
d'indice 2
(2) $Z_{\rho,\sigma}^! \simeq M(0) . L_0^\pi$.

\newpage

%% ===== Page 622 =====

On a

(3) $$ \widehat{G}_{\rho,\sigma} \simeq \hat{M} $$

et, plus précisément, le diagramme

(4)
\begin{tikzcd}
1 \ar[r] & L_0 \ar[r] \ar[d, hook] & Z_{\rho,\sigma} \ar[r] \ar[d, hook] & \Gamma_{\rho,\sigma} \ar[r] \ar[d, equal] & 1 \\
1 \ar[r] & \widehat{C}^+_{\rho,\sigma} \ar[r] & \widehat{G}_{\rho,\sigma} \ar[r] & \Gamma_{\rho,\sigma} \ar[r] & 1
\end{tikzcd}

s'identifie au diagramme

(4 bis)
\begin{tikzcd}
1 \ar[r] & L_0 \ar[r] \ar[d, hook] & \tilde{M}[0] \ar[r] \ar[d, hook] & \tilde{\Gamma}_Q \ar[r] \ar[d, "\simeq"'] & 1 \\
1 \ar[r] & \widehat{C}^+ \ar[r] & \hat{M} \ar[r] & \tilde{\Gamma}_Q \ar[r] & 1
\end{tikzcd}

D'ailleurs, les hom.

(5)
\begin{tikzcd}
& \hat{\mathbb{Z}}^* \\
\Gamma_{\rho,\sigma} \ar[ur, "\varepsilon_{\rho} = \varepsilon_f"] \ar[dr, "\varepsilon_{\sigma}"'] \\
& \mu
\end{tikzcd}

s'identifient à :

(5 bis)
\begin{tikzcd}
& \mu \\
\tilde{\Gamma}_Q \ar[ur, "\varepsilon_3"] \ar[r, "\chi"] \ar[dr, "\varepsilon_2"'] & \hat{\mathbb{Z}}^* \\
& \mu
\end{tikzcd}

b) On a

(6) $$ N_0^* = \{ u \in \hat{M} \mid u(\sigma) = \varepsilon_3(u) \sigma^{\varepsilon_2(u)} \} \quad \begin{array}{l} (= \varepsilon_3(u) \cdot \sigma) \\ = \varepsilon_3(\tau)(\sigma) \\ = \varepsilon_3(\tau')(\sigma) \end{array} $$

[unclear]

(7)
$$
\begin{cases}
N_0^* \cap \widehat{G}^{\wedge} = D_0 = \{1, \tau \} \cdot L_0 = \{1, \tau' \} \cdot L_0 \\
N_0^* \cap \widehat{G}^{+\wedge} = L_0 \quad (\simeq \hat{\mathbb{Z}}) \\
N_0^* \cap \Pi = \mu \\
N_0^* \cap \Pi_+ = \{ 1 \}
\end{cases}
$$

\newpage

%% ===== Page 623 =====

Posons

$$ \begin{cases} \mathcal{M}_0^! = \mathcal{M}(0) . \Pi_0 \subset \mathcal{M}^! = \mathcal{M}(0) . \Pi = \operatorname{Ker}(\mathcal{M} \to \mathfrak{S}_3) \\ \text{de sorte que } \mathcal{M}^! \simeq \mathcal{M}_0^! \times \mu \end{cases} $$

et

$$ \begin{cases} \mathcal{M}(1/2) \stackrel{df}{=} \mathcal{N}_\sigma^* \cap \mathcal{M}_0^! \\ \mathcal{N}_\sigma^{*!} = \mathcal{N}_\sigma^* \cap \mathcal{M}^! \end{cases} $$

alors l'hom. $\mathcal{M} \to \mathcal{T}_{\rho, \sigma}$ induit un

iso
(10) $$ \mathcal{M}(1/2) = \mathcal{N}_{\sigma 0}^* \stackrel{\sim}{\longrightarrow} \mathcal{T}_{\rho, \sigma} $$

donnant des scindages en produits semi-directs

(11) $$ \begin{cases} \mathcal{M} \simeq \mathcal{M}(1/2) . \widehat{\mathfrak{S}}^{+\wedge} \\ \mathcal{M}^! \simeq \mathcal{M}(1/2) . \Pi \\ \mathcal{M}_0^! \simeq \mathcal{M}(1/2) . \Pi_0 \\ \mathcal{N}_\sigma^{*!} \simeq \mathcal{M}(1/2) \times \mu \\ \mathcal{N}_\sigma^* \simeq \mathcal{M}(1/2) . L_\sigma \end{cases} , \quad \text{et} $$

(12) $$ \mathcal{T}_{\rho, \sigma} \stackrel{\sim}{\longleftarrow} \mathcal{M}(1/2) = \mathcal{N}_{\sigma 0}^* \hookrightarrow \mathcal{N}_\sigma^* . $$

c) On a

(13) $$ \mathcal{N}_\rho^* = \left\{ u \in \mathcal{M} \mid u(\rho) = \rho^{\varepsilon_3(u)} \quad (= \varepsilon_3[\tau_1](\rho)) \right\} $$

et

(14) $$ \begin{cases} \mathcal{N}_\rho^* \cap \widehat{\mathfrak{S}}^\wedge = D_\rho = \{1, \tau_1\} . L_\rho \\ \mathcal{N}_\rho^* \cap \widehat{\mathfrak{S}}^{+\wedge} = L_\rho \quad (\simeq \hat{\mathbb{Z}}) \\ \mathcal{N}_\rho^* \cap \Pi = \mu \\ \mathcal{N}_\rho^* \cap \Pi_0 = \{1\} \end{cases} $$

\newpage

%% ===== Page 624 =====

Soit

(15) $$ N_{\rho}^{*!} \overset{dfn}{=} N_{\rho}^{*} \cap M^! \quad , \quad N_{\rho}^{\wedge!} = N_{\rho}^{\wedge} \cap M^! $$

on trouve

(16) $$ N_{\rho}^{*!} \subset Z_{\rho} = Cent_{M}(\rho) = Cent_{M}(L_{\rho}) $$

et on montre ais.

(17) $$ N_{\rho, 0}^{\wedge !} = Z_{\rho, 0}^{!} = M^{!}[-J] . $$

On voit que l'hom. $M \to \gamma_{\rho, \sigma} = \Gamma_{\mathbb{Q}}^{\sim}$

induit un iso

(18) $$ M^{!}[-J] = Z_{\rho, 0}^{!} \overset{\sim}{\longrightarrow} \gamma_{\rho, \sigma}^{\sim} = \Gamma_{\mathbb{Q}}^{\sim} $$
$$ \| dfn $$
$$ \{ \gamma \in \gamma_{\rho, \sigma} | \varepsilon_3(\gamma) = 1 \} $$

d'où l'on déduit des scindages en

produits semi-directs (voire produits directs)

(19)
$$ M^{!} \simeq M^{!}[-J] \cdot \widehat{\mathfrak{S}}^{+\wedge} $$
$$ M^{!!} \simeq M^{!}[-J] \cdot \pi $$
$$ M_0^{!} \simeq M^{!}[-J] \cdot \pi_0 $$
$$ N_{\rho}^{*!} \simeq M^{!}[-J] \times \mu $$
$$ N_{\rho}^{\wedge!} \simeq M^{!}[-J] \times L_{\rho} $$

[Texte barré et raturé : La comparaison entre $M\{0\} \simeq \gamma_{\rho, \sigma}$ et $M[-J] \subset N_{\rho}^{*!} \overset{\sim}{\longrightarrow} \gamma_{\rho, \sigma}$ donne par l'opération (18) $M^{!} [\pi_0] \longrightarrow$ ]

\newpage

%% ===== Page 625 =====

516

Soit maintenant

(18) $$ \mu_{\tau} = \{1, \tau'\} \subset [unclear: \widehat{\mathfrak{S}}^\wedge] $$

alors $\mu_{\tau}$ normalise $L_{\rho}$, car
$$ \tau'(\rho) = \rho^{-1} \quad \text{d'où} \quad \tau'(L_{\rho}) = L_{\rho} $$
et il normalise $Z_{\rho} = Cent_M(L_{\rho})$, d'où
$$ Z_{\rho}^! = Z_{\rho} \cap M^! \quad \text{(puisque } M^! \text{ est normal dans } M) = Z_{\rho, 0}^! \rtimes \mu = M'(-j) \rtimes \mu. $$

[rayé: D'ailleurs $\mu_{\tau} \cap Z_{\rho}^! = \mu_{\tau} \cap M^! = \{1\}$. Sur]

(21) $$ N_{\rho}^? = \mu_{\tau} \cdot Z_{\rho}^! = \mu_{\tau} \cdot (Z_{\rho, 0}^! \rtimes \mu) \quad [rayé: illisible] $$
(produit 1/2 direct) $$ M'(-j) \rtimes \mu \simeq \Gamma_Q^\sim \rtimes \mu $$

On fera attention que contrairement à ce que pourrait suggérer cette écriture, $\mu_{\tau}$ ne normalise pas $Z_{\rho, 0}^!$ - sinon on aurait une décomposition en produit
$$ N_{\rho}^? = N_{\rho, 0}^? \times \mu, \quad \text{où} \quad N_{\rho, 0}^? = [rayé: \dots] Z_{\rho, 0}^! \cdot \mu_{\tau}, $$
or il n'en est rien. On a une suite exacte

(22) $$ 1 \to \mu \to N_{\rho}^? \to \gamma_{\rho, \sigma} \to 1 $$
$$ \parallel $$
$$ \Gamma_Q^\sim $$

de sorte que $N_{\rho}^?$ est une quasi-section de multiplicateur 2 de l'extension
$N_{\rho}^\wedge \to \gamma_{\rho, \sigma}$, - ou $M \to \gamma_{\rho, \sigma} = \Gamma_Q^\sim$.

Ainsi on trouve

\newpage

%% ===== Page 626 =====

(23)
$$ \begin{cases} M = N_p^? \wedge_{\mu_2} \widehat{C}^{+ \wedge} \\ M_p^? = N_p^? \wedge_{\mu_2} Z_p^! \simeq N_p^? \times L_p \end{cases} $$
[under $L_p$]: $\simeq \mathbb{Z}/2\mathbb{Z}$

La [unclear word crossed out] (22) de $\gamma_{\rho, \sigma}$ par $\mu$
l'extension induite $\gamma_{\rho, \sigma}$ dans la section précé-
dente, elle est [unclear words crossed out]
Elle donne naissance à un iso-
morphisme

(24) $$ \gamma_{\rho, \sigma} \simeq N_p^? / \mu \xrightarrow{\sim} N_p^* / \mu . $$

Remarque. Pour bien faire il faudrait
expliciter l'opération de $\tau'$ sur
$Z_p^! \simeq Z_{p,0}^! \times \mu \simeq \widehat{\Gamma}_Q^{\sim} \times \mu$, en termes de
l'opération de $\tau'$ [crossed out word] de $\tau'$ sur
$\widehat{\Gamma}_Q^{\sim}$ par autom. int. (NB $\tau'$ n'est
défini que par l'interm. de $\widehat{\Gamma}_Q$, $\widehat{\Gamma}_Q^{\sim}$, "à con-
jugaison près") et d'un hom. de
[MARGIN: de d'une "correction quadratique" de $\widehat{\Gamma}_Q^{\sim}$]
$\widehat{\Gamma}_Q^{\sim}$ dans $\mu$ - il y a fort à parier que
celui-ci est justement $\varepsilon_2$ - de sorte
que $\Gamma_Q^{\sim} \simeq \Gamma_Q^{\sim} \wedge \Gamma_Q^{\sim}$ apparaît
comme invariante par $\tau'$ ... En somme,
il y aurait lieu de reprendre l'étude
de l'extension faite dans la section précédente,
en l'épurant du point de vue des [unclear: $\alpha, \beta, \gamma \dots$]

\newpage

%% ===== Page 627 =====

dans l'optique actuelle, sans doute plus
simple. Il y aurait lieu en même temps
d'expliciter une autre section partielle,
qui mériterait serait au dessus de $\widetilde{\Gamma_\mathbb{Q}}$,
$$ \Gamma_\mathbb{Q}^{'' \sim} = \{ x \in \widetilde{\Gamma_\mathbb{Q}} \mid \varepsilon_1(x)\varepsilon_2(x) = 1 \} $$
et qui mériterait peut-être la notation

[MARGIN: ? $\mathcal{M}(i)$ th un cas ... identif. !]
$\mathcal{M}(\sqrt{j})$ ou du moins $\mathcal{M}(q)$, pour
une génératrice $q$ convenable de $\mathbb{Q}(\sqrt{j})$ ???

Je me rends compte que je suis en train
de déconner - ayant un peu perdu contact
avec le contenu géométrique de mes
calculs. Il faut garder à l'esprit que
(25) $$ \mathcal{M}_0^! \xrightarrow{\sim} \mathcal{M}_{0,3}^{! \sim} $$
correspond, à un $\sim$ près, au $\Pi_1$ de $\mathcal{U}_{0,3 \mathbb{Q}}$,
les sections de $\mathcal{M}_0^! \to \Gamma_\mathbb{Q}^\sim \simeq \mathcal{G}_{j, \sigma}$ correspondent
donc (au $\sim$ près) : elles ou $\Pi_1(\mathcal{U}_{0,3 \mathbb{Q}})$ sur $\Gamma_\mathbb{Q}$,
donc ("moralement") aux points de $\mathcal{U}_{0,3 \mathbb{Q}}$
— quand on a des sections partielles, aux
points de $\mathcal{U}_{0,3 \mathbb{Q}}$ à valeurs dans des extensions
finies de $\mathbb{Q}$. Ainsi, la section envisagée
dans b) correspond bien à un pt $1/2$ de $\mathcal{U}_{0,3 \mathbb{Q}}$,
la quasi-section envisagée $\mathcal{M}(-j)$ ici dans c) correspond à
un point $-j$ de $\mathcal{U}_{0,3 \mathbb{Q}}$, qui est dans $\mathcal{U}_{0,3 \mathbb{Q}}(\mathbb{Q}[\sqrt{j}])$,
d'où la notation $\mathcal{M}(\sqrt{j})$ raisonnable. Dans ce $\mathcal{M}_{\hat{\Gamma}}^\sim$
elle n'est pas $\subset \mathcal{M}_0^!$ - on en divise par $\mu$ ;
on trouve après division par $\mu$ une section

\newpage

%% ===== Page 628 =====

de $M_{0,3}^\sim$ sur $\Gamma_{\mathbb{Q}}^\sim$, mais on n'a pas une
section de $M_{0,3}^\sim$. Or à un $\sim$ près, $M_{0,3}^\sim$
est le $\pi_1$ de $M_{1,1 \mathbb{Q}}$, schéma modulaire sur $\mathbb{Q}$
des courbes elliptiques mod. symétrie.
Donc la donnée de cette section correspond à
la donnée d'une telle "courbe elliptique
mod. symétrie" sur $\mathbb{Q}$. Passant au sous-
groupe $\Gamma_{\mathbb{Q}}'^\sim$ - ce qui correspond à l'ex-
tension du corps de base de $\mathbb{Q}$ à $\mathbb{Q}(j) = \mathbb{Q}(\sqrt[3]{1})$ -
on retrouve l'image dans $M_{0,3}'^\sim$ de la quasi
section pointée $M(-j)$ - cela montre
donc qu'il s'agit toujours de la "même" courbe
elliptique. J'ai bien l'impression cependant
que sur $\Gamma_{\mathbb{Q}}^\sim$ la section, cette section de
$M_{0,3}^\sim$ ne provient pas d'une section
de $M_{1,1}^\sim$ sur $\Gamma_{\mathbb{Q}}^\sim$ - i.e. justement que l'ext. $N_{\mathbb{Q}}^*$
de $\Gamma_{\mathbb{Q}}^\sim$ par $\mu_2$ n'est pas triviale. Cela signifierait
il qu'il n'existe pas de courbe elliptique
déf / $\mathbb{Q}$ d'invariant $0$ ? (Chose qui semble absurde
- car doit pouvoir sur un corps tr. donner n'importe
[MARGIN: ?] quel invariant ...). De façon précise, on a le
diagramme de schémas modulaires sur $\mathbb{Q}$

(26)
\begin{tikzcd}
\text{sch. mod. } M_{1,1} \text{ des courbes ell.} \ar[d, "\text{gerbe [unclear]}"'] \ar[dr] & \\
\text{sch. mod. des } M'_{1,1} \text{ courbes ellipt. mod. symétrie} \ar[r] & E' \text{ schéma modulaire grossier}
\end{tikzcd}

Le fait qu'un point de $M'_{1,1}(\mathbb{Q})$ ne se re-

\newpage

%% ===== Page 629 =====

montre qu'un [unclear: point] de $M_{1,1}$, [unclear: vu comme]
par [unclear: son] image dans $E'$ [unclear] se remonte pas
à $M_{1,1}$. De façon précise, la [unclear: digression]
s'identifie à :

(27)
$$M_{1,1} \simeq (U_{0,3}, \tilde{\mathfrak{S}}_3)$$
$$\downarrow$$
$$M'_{1,1} \simeq (U_{0,3}, \mathfrak{S}_3) \longrightarrow U_{0,3}/\mathfrak{S}_3 = U_{0,3}/\tilde{\mathfrak{S}}_3 = E'$$

où $\tilde{\mathfrak{S}}_3$ est l'extension de $\mathfrak{S}_3$ par $\mu_2$, définie par

(28) $$\tilde{\mathfrak{S}}_3 = \hat{\mathfrak{S}}^+/\Pi_0 \simeq \operatorname{Sl}(2, \mathbb{Z}) / \text{sous-groupe invariant engendré par } -e_0$$
$$\simeq \operatorname{Sl}(2, \mathbb{Z}/4\mathbb{Z}) / \text{sous-groupe (d'ordre 4) engendré par } -e_0, \dots$$

en faisant opérer $\tilde{\mathfrak{S}}_3$ sur $U_{0,3}$ par l'intermédiaire de $\mathfrak{S}_3$. Un pt de $M'_{1,1}$ s'identifie à un $\mathfrak{S}_3$-torseur $T$ sur $\mathbb{Q}$, muni d'un morphisme

$$T \longrightarrow U_{0,3}$$

compatible avec les actions de $\mathfrak{S}_3$, les points de $M_{1,1}$ s'identifiant à un $\tilde{\mathfrak{S}}_3$-torseur $\tilde{T}$ sur $\mathbb{Q}$, muni d'un morphisme

(30) $$\tilde{T} \longrightarrow U_{0,3}$$

compatible avec l'opération de $\tilde{\mathfrak{S}}_3$. Désignant par

$$T = \tilde{T}/\mu_2$$

le $\mathfrak{S}_3$-torseur correspondant, la donnée de (30) équivaut à la donnée du [unclear: corps]. Dans la

[MARGIN: Tout le baratin sur $\tilde{\mathfrak{S}}_3 \subset \operatorname{Sl}(2, \mathbb{Z}/4\mathbb{Z})$ ... [unclear] d'étudier les points d'ordre 4 via $\operatorname{Gl}(2, \mathbb{Z}/4\mathbb{Z})$ pour la "courbe elliptique" ...]

\newpage

%% ===== Page 630 =====

relèvements à $M_{1,1}$ d'un point de $M_{1,1}'$ à / 0
s'identifient aux $\widetilde{\mathfrak{S}}_3$-torseurs $\widetilde{T}$ qui relèvent le
$\mathfrak{S}_3$-torseur $T$, i.e. (à iso près) ils correspondent
[MARGIN: en supp. choisi un pt base de $T(\bar{\mathbb{Q}})$]
aux hom. $\Gamma_{\mathbb{Q}} \to \widetilde{\mathfrak{S}}_3$ qui relèvent l'hom
$\Gamma_{\mathbb{Q}} \to \mathfrak{S}_3$ qui définit $T$ :

\begin{tikzcd}
& \widetilde{\mathfrak{S}}_3 \ar[d] \\
\Gamma_{\mathbb{Q}} \ar[r] \ar[ur, dashed] & \mathfrak{S}_3
\end{tikzcd}

l'obstruction est, comme il est d'usage, une
extension de $\Gamma_{\mathbb{Q}}$ par $\mu_2$, i.e. un élément
de $H^2(\mathbb{Q}, \mu_2) \simeq {}_2\text{Br}(\mathbb{Q})$. Dans le cas
qui nous occupe l'image de $T$ dans $U_{0,3}$
est formé de $\{ -j, -j^2 \}$ qui est connexe sur $\mathbb{Q}$,
[MARGIN: condition d'[unclear] plus bas]
ce qui entraine $\Gamma_{\mathbb{Q}} \to \mathfrak{S}_3$ est surjectif. De façon
générale, en [unclear] le th. de [unclear]
galoisienne en termes de scindages
de $M_{0,3}$ sur $\Gamma_{\mathbb{Q}}$, l'hom. $\Gamma_{\mathbb{Q}} \to \mathfrak{S}_3$
est le composé
$$ \Gamma_{\mathbb{Q}} \xrightarrow{\kappa} M_{0,3} \to \mathfrak{S}_3 ,$$
où $\kappa$ est une section de
$$ M_{0,3}^1 = \operatorname{Ker}(M_{0,3} \to \mathfrak{S}_3) $$
sur $\Gamma_{\mathbb{Q}}$. Ici une section de
$M_{1,1}'$ au dessus de $\Gamma_{\mathbb{Q}}$ donne une section de
$M_{1,1}' \xrightarrow{\text{au dessus de}} \dots \Gamma_{\mathbb{Q}}'$ d'ordre 2,
donc $\kappa(\Gamma_{\mathbb{Q}}) = 1$, on voit que

\newpage

%% ===== Page 631 =====

c'est le sous-groupe $\{1, \tilde{\sigma}_\infty\}$ de $\widetilde{\mathfrak{S}}_3$. (Car l'image de
$\tau' \in M_{0,3}$ dans $\widetilde{\mathfrak{S}}_3$ est (comme celle de $\tau = 1$)
est égale à celle de $\sigma = \sigma_\infty$, i.e. $\tilde{\sigma}_\infty$. Ce signifie
que l'hom.
$$T \simeq T_0 \wedge_{\{1, \tilde{\sigma}_\infty\}} \widetilde{\mathfrak{S}}_3$$
i.e. $T$ se déduit d'un $\mu$-torseur (i.e. d'une
ext. quadratique de $\mathbb{Q}$, savoir justement
$\mathbb{Q}(\sqrt{j}) = \mathbb{Q}(j)$) par extension du groupe
$$ \mu \hookrightarrow \widetilde{\mathfrak{S}}_3 \quad (-1 \mapsto \tilde{\sigma}_\infty)$$
d'opérateurs. Donc, prenant l'image inverse
de $\mu = \{1, \tilde{\sigma}_\infty\}$ dans $\widetilde{\mathfrak{S}}_3$, on trouve une
extension de $\mu$ par $\mu$ - ça ne peut guère être

[MARGIN: oui, c'est clair puisqu'on a $\tilde{\sigma}_\infty$ dans $\widetilde{\mathfrak{S}}_3 = [unclear]$ est bien d'ordre 4 donc $\mathcal{E}'' \simeq \mathbb{Z}/4\mathbb{Z}$ canoniquement - ]

que $\mathbb{Z}/4\mathbb{Z}$ ! - que j'ai envie de noter
$\mathcal{E}''$ :
$$1 \to \mu \to \overset{\mathbb{Z}/4\mathbb{Z}}{\mathcal{E}''} \to \mu \to 1$$
et le pb d'obstruction est de
trouver un remontage

\begin{tikzcd}
\mathcal{E}'' \ar[r] & \mu \\
\Gamma_{\mathbb{Q}} \ar[u, dashed] \ar[ur] &
\end{tikzcd}

[MARGIN: C'est évident - [unclear: il n'y a qu'à voir] $(\Gamma_{\mathbb{Q}})_{ab} \simeq \hat{\mathbb{Z}}^\times \supset \mathbb{Z}_3^\times$, or on a $(\mathbb{Z}_3^\times)_4 \simeq \mathbb{Z}/2\mathbb{Z}$ (pas de r. q. dans $\mathbb{Z}_3$ [unclear]) - ]

de l'obstruction - s'interprète comme l'extension
de $\Gamma_{\mathbb{Q}}$ par $\mu$, image inverse de l'extension
$\mathcal{E}''$. (Il faudrait regarder pourquoi, en termes

\newpage

%% ===== Page 632 =====

classiques, cette obstruction est éliminée, quitte à
passer au sous-groupe $\Gamma_{\mathbb{Q}}''$ correspondant au
corps $\mathbb{Q}(\sqrt{3})$

Ainsi, le point $\mathbb{Q}$-rationnel envisagé de
$M_{1,1}$ ne remonte pas en un point $\mathbb{Q}$-
rationnel de $M_{1,1}$ . Par contre, son [unclear]
i.e. $0$ dans $E' \simeq U_{0,3} / \widetilde{S}_3$, doit bien en
remonter - i.e. il doit bien exister un
$\widetilde{S}_3$-torseur $\widetilde{P}$, et un $\widetilde{S}_3$-hom
$$\widetilde{P} \longrightarrow U_{0,3} \quad i.e. \quad P \longrightarrow U_{0,3}$$
avec $P = \widetilde{P}/\mu$.

Considérons en effet le schéma [unclear]
$$Q = U_{0,3} \times_{E'} \{0\}$$
l'image inverse de $0$, soit
$$P_0 = Q_{red}$$
c'est un sous-schéma fini, qui sur un corps
$K \supset \mathbb{Q}(j)$, se scinde en une somme de
deux pts, et qui sur $\mathbb{Q}$ est connexe!
- c'est un schéma isomorphe à
$\operatorname{Spec}(\mathbb{Q}(\sqrt{-3}))$. Il est stable par $\widetilde{S}_3$,
qui y opère via
$$\widetilde{S}_3 \xrightarrow{sgn} \mu$$
et l'opération galoisienne de $\mu$ sur le $\mu$-torseur $P_0$.

\newpage

%% ===== Page 633 =====

Bref ... le diagramme

(31)
\begin{tikzcd}
\tilde{\mathfrak{S}}_3 \ar[r, "can"] & \mathfrak{S}_3 \ar[r, "sgn"] & \mu \\
& \Gamma_{\mathbb{Q}} \ar[u] \ar[ur] &
\end{tikzcd}

[MARGIN: NB le noyau de $\tilde{\mathfrak{S}}_3 \to \mu$ est une extension centrale de $\mathbb{Z}/3\mathbb{Z}$ par $\mu = \mathbb{Z}/2\mathbb{Z}$, elle est donc triviale et isom. : $\mathbb{Z}/6\mathbb{Z}$ (unique [unclear])]

On voit donc que la catégorie des relèvements
d'un pt $\mathbb{Q}$-rationnel $e$ de $E'$ en un
pt $\mathbb{Q}$-rationnel de $M_{1,1}^\sim$ resp. $M_{1,1}$, est
équivalente à : celle des $\tilde{\mathfrak{S}}_3$-torseurs $\tilde{P}$, resp.
des $\mathfrak{S}_3$-torseurs $P$ sur $\mathbb{Q}$, qui relèvent le
$\mu$-torseur $P_0$. On avait choisi un
$P=T$ de façon particulièrement triviale,
en utilisant le scindage
$$ \mu \xrightarrow{\sim} \{1, \sigma_3\} \subset \mathfrak{S}_3 $$
de $\mathfrak{S}_3 \to \mu$ — mais ce torseur $T$ ne se
remonte pas en un $\tilde{T}$. Les classes d'isomor-
phie de pts de $M_{1,1}^\sim$ resp. de $M_{1,1}$ au-dessus
du pt $e$ de $E'$, correspondent, à [unclear] près
(mod noyaux de $\tilde{\mathfrak{S}}_3 \to \mu$ et $\mathfrak{S}_3 \to \mu$),
aux classes de conjugaisons (des relèvements de
$\Gamma_{\mathbb{Q}} \to \mu$ en $\Gamma_{\mathbb{Q}} \to \mathfrak{S}_3$ resp. en $\Gamma_{\mathbb{Q}} \to \tilde{\mathfrak{S}}_3$).

Considérons un relèvement
$$ \Gamma_{\mathbb{Q}} \to \tilde{\mathfrak{S}}_3 $$
de ce type... un sous-groupe de $\tilde{\mathfrak{S}}_3$, qui n'est

\newpage

%% ===== Page 634 =====

$\{1,\dot{\sigma}_3\}$ - ou un de ses conjugués $\{1,\sigma'_3\} , \{1,\sigma''_3\}$
ou c'est $\mathcal{G}_3$ tout entier. Dans le premier
cas, on voit que $s_Q$ ne se remonte pas
en $\tilde{\mathcal{G}}_3$. Donc pour [unclear: y arriver]
il faut prendre un épi $\Gamma_{\mathbb{Q}} \to \mathcal{G}_3$ -
ce qui sort [unclear: du cadre des] ext. abéliennes sur
$\mathbb{Q}$, p.ex entrer dans le domaine des
extensions abéliennes de $\mathbb{Q}(j)$.

Examiner en termes des groupes de
Galois s'il existe un tel relèvement, soit
$\Gamma_{\mathbb{Q}} \to \tilde{\mathcal{G}}_3$, et quand, comment les obtenir,
apparait comme un exercice plaisant
et délectable de théorie des corps de
classes, que je ne vais pas poursuivre
pour le moment, pour ne pas éterniser
cette digression. Revenant par la suite
sur ce point, il me faudrait en même
temps expliciter les relations entre les
points de vue remontages d'hom. de $\Gamma_{\mathbb{Q}}$,
et scindages de l'extension $\mathcal{M}_{0,3}$ de $\Gamma_{\mathbb{Q}}$
par $\hat{\mathcal{B}}_{0,3}$, ou $\mathcal{M}_{1,1}$ de $\Gamma_{\mathbb{Q}}$ par $\hat{\mathcal{G}}^\sim \simeq \operatorname{Sl}(2,\mathbb{Z})^\wedge$.

\newpage

%% ===== Page 635 =====

621

Revenons au sous-groupe [crossed out: unclear]
$$ \mathcal{N}_p^? \underset{4\downarrow}{\overset{3}{\subset}} \mathcal{N}_p^* \subset \mathcal{M} $$
$$ Z_p^! = \mathcal{M}^!(-\bar{\jmath}) $$
contenant les [crossed out: sections] sous-groupes partiels $\mathcal{M}[-j] \simeq \hat{\Gamma}_Q^\sim$.

Soit
(32) $$ \mathcal{M}^\circ[-j] \overset{2}{\subset} \mathcal{M}^![-\bar{\jmath}] $$
$$ Ker ( \mathcal{M}[-j] \xrightarrow{\varepsilon_2} \mu ) = \mathcal{N}_p^* \cap \mathcal{M}^{\circ !} $$
qui a [crossed out: unclear] une section partielle au dessus
(33) $$ \hat{\Gamma}_Q^{\circ \sim} = Ker ( \hat{\Gamma}_Q^\sim \xrightarrow{(\varepsilon_2, \varepsilon_3)} \mu \times \mu ) $$
$$ = Ker ( \hat{\Gamma}_Q^\sim \longrightarrow \hat{\mathbb{Z}}^* \longrightarrow (\mathbb{Z}/12\mathbb{Z})^* ) $$

[crossed out: Considérons] Modulo une vérification que je vais faire plus bas, $\tau'$ normalise $\mathcal{M}^\circ[-j]$, considérons le
(34) $$ \mathcal{M}'''[-j] \overset{d\acute{e}f}{=} \mu_2 \cdot \mathcal{M}^\circ[-j] \quad (\mu_2 = \{1, \tau'\}) $$
(semi-direct)
c'est un sous-groupe d'indice 4 de $\mathcal{N}_p^?$,
et on a maintenant
(35) $$ \mathcal{M}'''[-j] \overset{\sim}{\longrightarrow} \hat{\Gamma}_Q^\sim \quad , $$
c'est une section partielle au dessus de $\hat{\Gamma}_Q^\sim$, d'où des isomorphismes

\newpage

%% ===== Page 636 =====

(36)
$$
\begin{cases}
\mathcal{M}''' \simeq \mathcal{M}'''(-\overline{j}) . C^{\pm} \\
\mathcal{N}_p^* \simeq \mathcal{M}'''(-\overline{j}) . L_p
\end{cases}
$$

On a le diagramme récapitulatif
d'inclusions de sous-groupes de $\mathcal{N}_p^*$

(37)
\begin{tikzcd}
& & Z_p \ar[ddr, "2"] \\
& \mathcal{M}'(-\overline{j}) = Z_{\hat{p}}! \ar[ur, "6"] \ar[dr, "4"] \\
\mathcal{M}^0(-\overline{j}) \ar[ur, "2"] \ar[dr, "2"'] & & \mathcal{N}_p^? \ar[r, "3"] & \mathcal{N}_p^* \\
& \mathcal{M}'''(-\overline{j}) \ar[ur, "4"']
\end{tikzcd}

Notons que l'inclusion

(38) $$ \underset{\simeq \mathbb{D}_6}{\mathbb{D}_p} \overset{df}{=} \underset{\simeq \mathbb{Z}/2 \cdot \mathbb{Z}/6\mathbb{Z}}{\mu_{\pm} \cdot L_p} \subset \mathcal{N}_p^* $$

induit un isomorphisme d'

(39) $$ \mathbb{D}_p \overset{\sim}{\longrightarrow} \mathcal{N}_p^* / \text{[unclear]} $$

\newpage

%% ===== Page 637 =====

(36) 
$$ \begin{cases} M''' \simeq M'''(-\bar{J}) . C_2^{+\wedge} \\ N_{\rho}^* \simeq M'''(-\bar{J}) . L_{\rho} \end{cases} $$

Diagramme récapitulatif des sous-groupes de $N_{\rho}^*$ obtenus jusqu'à présent :

(37)
\begin{tikzcd}
& M^0(-\bar{J}) \ar[rr, "2"] \ar[dl, "2"'] && \underset{=Z^0_{\rho}}{M'(-\bar{J})} \ar[rr, "2"] \ar[dl, "2"'] && M'(\bar{J})_{xt} \ar[rr, "3"] \ar[dl, "2"'] && Z_{\rho} \ar[dl, "2"'] \ar[dd, "2"] \\
M'''(-\bar{J}) \ar[rr, "2"] && M''(-\bar{J}) \times \mu_C \ar[rr, "2"] && N_{\rho}^{\pm C} \ar[rr, "3"] && N_{\rho}^* \ar[dr] \\
&& \mu_{\pm 1 C} \ar[ul, "1"] \ar[rr, "2"] && D_{\pm 1 C} \ar[ul, "\mu_C"'] \ar[rr, "3"] && D_{\rho} \ar[ul, "2"'] \ar[ur, "2"'] & L_{\rho}
\end{tikzcd}

où nous omettons les inclusions relatives
à 2 "faux sous-groupes", comme les
sous-groupes d'indice fini de $N_{\rho}^*$,
intermédiaires entre celui-ci et $M^0(-\bar{J})$
(d'indice 24 dans $N_{\rho}^*$). Je vais examiner
les questions d'isomorphismes de ces sous-groupes
les uns dans les autres, et la structure
des groupes quotients.

Je vais faire un diagramme qui
montrera bien les conglomérations entre
les groupes en question.

\newpage

%% ===== Page 638 =====

(38)
\begin{tikzcd}[row sep=1.5em, column sep=1.5em]
& Z_\rho \ar[rrr, "2"] & & & N_\rho^* \\
& & & Z_{\rho^0}^! = M'(-J) \ar[ull, "3"'] \ar[rrr, "2"] \ar[dlll, "2" near end] & & & M'(-J).\mu_{\tau'} \ast \ar[ull, "3"'] \ar[dlll, "2" near end] \\
Z_\rho^1 = M'(-J) \rtimes \mu \ar[ur, "2"] \ar[rrr, "2" near start] & & & N_\rho^{\prime !} \ar[ur, "2"] \ar[rr] & & M''(-J) \ar[ur, "2"'] \\
& Z_{\rho^0} \ar[uu, dashed, "3"] \ar[rrr, dashed, "2" near start] \ar[dl, "2" near end] & & & N_\rho' \ar[uu, dashed, "3"] \ar[dl, "2" near end] \\
M^0(-J) \rtimes \mu \ar[uu, "2"] \ar[rrr, "2"] & & & M''(-J) \rtimes \mu \ar[uu, "2"] \ar[rr, "2"] & & \mu_{\tau'} \ar[uu, "2"'] \\
& H \ar[uu, dashed] \ar[rrr, dashed, "2"] & & & C \ar[uu, dashed] \ar[rr, dashed, "2"] & & D_{\tau'} \ar[uu, dashed] \\
& & & & & & \\
& L_\rho \ar[uuu, dashed] \ar[rrr, dashed] & & & D_\rho \ar[uuu, dashed] \ar[uurr, dashed, "3"']
\end{tikzcd}

[MARGIN: $\ast$ Attention $M'(-J).\mu_{\tau'}$ n'est pas un sous-groupe, car $\mu_{\tau'}$ ne normalise pas $M'(-J)$.]

0) Sous-groupes d'indice fini de $Z_\rho$ -
dans le diagramme il y a une portion

(39)
\begin{tikzcd}
Z_\rho^1 = M'(-J) \rtimes \mu \ar[d, "2"'] & M'(-J) = Z_\rho^! = N_\rho^! \ar[l, hook', "2"'] \ar[d, "2"] \\
M^0(-J) \rtimes \mu & M^0(-J) = Z_{\rho^0}^! \ar[

\newpage

%% ===== Page 639 =====

623

même invariants dans $N_{\hat{\rho}}^*$, puisque
[crossed out: les $Z_{\hat{\rho}}, M^!, \ker \varepsilon_2$] sont stabilisés par $N_{\hat{\rho}}^*$.
D'ailleurs, comme $Z_{\hat{\rho}0}^! = M^!(-j)$ est d'indice 2
dans $Z_{\hat{\rho}}^! = M^!(-j) \times \mu$, pour montrer qu'il est
distingué dans $Z_{\hat{\rho}}$, il suffit de prouver qu'il est
normalisé par $\rho$ (puisque $Z_{\hat{\rho}} = \langle \rho \rangle \cdot Z_{\hat{\rho}}^!$ - c'est
$\pm$ trivial!), or examinons de façon générale
pour $x \in M$ si $x$ normalise $M^!_0$
$$M^!_0 \overset{d\acute{e}f}{=} M(0) \cdot \Pi_0$$
puisqu'il normalise $\Pi_0$, il suffit de
voir si pour $u \in M(0)$, on a
$$x u x^{-1} \in M^!_0$$
ce qui équivaut encore à
$$x u x^{-1} \equiv u \pmod{\Pi_0}$$
i.e.
$$[x, u] \in \Pi_0 \quad (u \in M^!_0)$$
or on sait que pour $x \in \widehat{\mathcal{B}}_3$, on a
(E) $$ \underbrace{[x, u]}_{x u x^{-1} u^{-1}} \equiv (sg(x))^{\varepsilon_2(u)} \pmod{\Pi_0} $$

$$sg : \widehat{\mathcal{B}}_3 \overset{can}{\longrightarrow} \mathfrak{S}_3 \overset{sg}{\longrightarrow} \mu$$

[crossed out: Cela montre que]
Lemme Si $x \in \widehat{\mathcal{B}}_3$, $u \in M^!_0$, alors [crossed out text] conditions équivalentes :
a) $x u x^{-1} \in M^!_0$
b) $x u x^{-1} \equiv u \pmod{\Pi_0}$ i.e. $[x, u] \in \Pi_0$
c) $(sg(x))^{\varepsilon_2(u)} = 1$ i.e. $sg(x) = 1$ ou $\varepsilon_2(u) = 1$

\newpage

%% ===== Page 640 =====

[crossed out: Pour $x=\rho$, cela montre que la condition]

Corollaire 1 Soit $x \in \widehat{\mathcal{B}}^{+\wedge}$. Pour que $x$ normalise $\mathcal{M}^!_0$, il faut et il suffit qu'on ait $sg(x)=+1$.

Corollaire 2 Le sous-groupe $\mathcal{M}^{\prime \prime !} = \mathcal{M}^{\prime \prime (0)} . \Pi_0$ est invariant donc normalisé par $\widehat{\mathcal{B}}^{+\wedge}$.

Cela montre en particulier que $\rho$ normalise $\mathcal{M}^!_0$, donc aussi $M'(-j) = Z'_p = Z_p \cap \mathcal{M}^!_0$, qui est donc invariant dans $Z_p$. Il en est de même de son intersection avec $\mathcal{M}^0(-j) \times \mu$, soit $\mathcal{M}^0(-j)$. (Ou encore ; ce groupe est égal : $Z_p \cap \mathcal{M}^{\prime \prime !}$, et $\mathcal{M}^{\prime \prime !}$ est normalisé par $\rho$... ). Ainsi les quatre groupes (31) sont distingués dans $Z_p$.

Pour voir s'ils sont distingués dans $\mathcal{N}_p^* = \langle \tau' \rangle Z_p$, il suffit de voir s'ils sont normalisés par $\tau'$. C'est clair pour $Z'_p = Z_p \cap \mathcal{M}^!_0$ - d'où le fait que $\mathcal{N}^{\prime *}_p = \langle \tau' \rangle Z'_p$ est un sous-groupe de $\mathcal{N}^*_p$. Pour $\mathcal{M}^0(-j) = Z_p \cap \mathcal{M}^{\prime \prime !}$, il [crossed out: en est de même]. Pour voir si les autres de ces sous-groupes de [unclear] sont normalisés, il faudrait donc expliciter l'action de $\tau'$ [crossed out: dessus] sur $Z'_p = \mathcal{M}'(-j) \times \mu$.

\newpage

%% ===== Page 641 =====

623

Notons que $\tau \in M(\infty) \subset M_0^!$ normalise $M_0^!$, et il opère sur $Z_\rho$ via $M^! = M_0^! \times \mu$ (avec $M_0^! \simeq \Gamma_Q \cdot \Pi_0$) par l'opération produit de la conjugaison par $\tau$ dans $M_0^!$, et l'identité dans $\mu$. On a d'autre part $\tau' = \tau \sigma$, on a vu que

(40) $$ \sigma(x) \equiv \varepsilon_\tau(x) x \quad (\Pi_0) \quad x \in M $$

donc

(41) $$ \tau'(x) \equiv \varepsilon_\tau(x) \tau(x) \quad (\Pi_0) \quad [MARGIN: \iff \tau'(x) \in M_0^!] \quad x \in M $$

d'où il s'ensuit

(42) $$ \text{si } x \in M_0^! \text{, alors } \tau'(x) \equiv x \quad (\Pi_0) \iff \varepsilon_\tau(x) = 1 $$

en particulier, $M_0^{!!}$ est normalisé par $\tau'$.
Donc si on applique ceci aux éléments de $Z_\rho^1 = M'(-J)$ dans $Z_\rho^! = M'(-J) \times \mu$, on trouve

(43) $$ \text{Soit } x \in Z_\rho^1 = M'(-J) \text{, alors } \tau'(x) \in Z_\rho^1 \iff \varepsilon_\tau(x) = 1 $$

Corollaire. $Z_\rho^1 = M'(-J)$ n'est pas normalisé par $\tau'$, et donc n'est pas invariant dans $N_\rho^*$, mais $M^0(-J) = Z_\rho \cap M_0^{!!}$ est normalisé par $\tau'$, donc invariant [unclear: dans le produit par $\mu$].

Donc le seul des quatre sous-groupes (39) de $Z_\rho$ qui ne soit distingué dans $N_\rho$ est $M'(-J) = Z_\rho^1$. Les autres le sont.

\newpage

%% ===== Page 642 =====

maintenant, par produit $\frac{1}{2}$ direct avec $L_{\hat{\rho}}$,
à des sous-groupes de $N_{\hat{\rho}}$, invariants dans $N_{\hat{\rho}}$.
Le plus petit de tous ces sous-groupes
est $M^\circ(-J)$, qui est d'indice $24$
dans $N_{\hat{\rho}}^*$. Le groupe quotient 
$$N_{\hat{\rho}}^* / M^\circ(-J)$$
est donc un groupe d'ordre 24 [MARGIN: $= H_{24} \cdot L_{\hat{\rho}}$] qui
reçoit d'ailleurs le groupe $D_{\hat{\rho}} \subset N_{\hat{\rho}}^*$,
normalisateur de $L_{\hat{\rho}}$ dans $\operatorname{Gl}(2, \hat{\mathbb{Z}})$,
groupe isomorphe à $D_6$, et d'ordre
12. Le noyau de
(44) $$D_{\hat{\rho}} \hookrightarrow N_{\hat{\rho}}^* / M^\circ(-J)$$
est évidemment contenu dans $L_{\hat{\rho}} = D_{\hat{\rho}} \cap Z_{\hat{\rho}}$
et il est réduit à 1, car
précisément 
$$L_{\hat{\rho}} \cap Z_0! \subset L_{\hat{\rho}} \cap M_0^* = L_{\hat{\rho}} \cap \Pi_0 = \{1\} \ .$$
(d'où 12)
Ainsi $D_{\hat{\rho}}$ apparait comme un sous-
groupe d'indice 2 de $N_{\hat{\rho}}^* / M^\circ(-J)$,
tout comme il est sous-groupe d'indice
2 de $\hat{G} / \Pi_0$. Il serait tentant

\newpage

%% ===== Page 643 =====

alors d'établir un isomorphisme
$$ N_{\hat{\Gamma}}^* / M^\circ(-J) \overset{?}{\simeq} \hat{\mathbb{S}}^* / \Pi_0 $$
compatible avec un plongement de $D_f$.
Essayons-le aussi :

Proposition Considérons l'hom. composé
(45)
\begin{tikzcd}
M \ar[r, "\theta_M"] & \operatorname{Gl}(2, \hat{\mathbb{Z}}) \ar[r] & \operatorname{Gl}(2, \mathbb{Z}/4\mathbb{Z}) \ar[r, "\text{épi. de}", "\text{degré 4}"'] & \hat{\mathbb{S}}^* / \Pi_0 \\
& & \operatorname{Gl}(2, \mathbb{Z}) / \mathfrak{S}[4] \ar[u, "\simeq"'] \ar[ur, "\simeq \hat{\mathbb{S}} / \hat{\Pi}_0"'] & 
\end{tikzcd}
qui induit sur $\hat{\mathbb{S}}^* \subset M$ l'homomorphisme canonique $\hat{\mathbb{S}}^* \hookrightarrow \hat{\mathbb{S}}^* / \Pi_0$. Alors l'hom. induit
$$ N_{\hat{\Gamma}}^* \longrightarrow \hat{\mathbb{S}} / \Pi_0 $$
a comme noyau $M^\circ(-J)$, et induit par passage au quotient un isomorphisme
(46)
$$ N_{\hat{\Gamma}}^* / M^\circ(-J) \overset{\sim}{\longrightarrow} \hat{\mathbb{S}} / \Pi_0 $$

Démonstration. Les éléments de $M^\circ(-J)$ sont les éléments de la forme
$$ \left\{ v = \alpha^{-1} u \ \left| \begin{array}{l} u \in M(0) \ ,\ \alpha \in \Pi_0 \\ u(p) = \alpha(p) \quad (\text{donc } \varepsilon_3(u) = 1) \\ \varepsilon_2(u) = 1 \end{array} \right. \right\} $$
ils sont en corr. 1-1 avec les éléments de
$$ M \hat{(0)} = \{ u \in M(0) \mid \varepsilon_2(u) = \varepsilon_3(u) = 1 \} \ . \quad \text{L'image de} $$

\newpage

%% ===== Page 644 =====

$v$ dans $\operatorname{Gl}(2,\hat{\mathbb{Z}})/\{\pm 1\} \Pi_0$, / [unclear] dans $\hat{\mathfrak{S}}/\Pi_0$,
est égale à celle de $u$. Donc l'assertion
que $M^\circ(-j)$ est dans le noyau de
$M \to \hat{\mathfrak{S}}/\Pi_0$, équivaut au fait que
$M^\circ(0)$ est dans ce noyau, donc que
l'on a un diagramme commut.

[MARGIN: NB. Ceci achève la démonstration, $\varepsilon_3$ n'intervenant pas dans l'exp. explicite de $M \to \hat{\mathfrak{S}}/\Pi_0$, cf plus bas...]

\begin{tikzcd}
M(0) \ar[r] \ar[d, "(\varepsilon_1,\varepsilon_3)"'] & M \ar[d, "(45)"] \\
\mu \times \mu \ar[r, "?"] & \hat{\mathfrak{S}}/\Pi_0
\end{tikzcd}

Notons que le composé
$$ M \to \hat{\mathfrak{S}}/\Pi_0 \to (\widehat{\mathbb{Z}/4\mathbb{Z}})^* \simeq \{\pm 1\} $$
n'est autre que $\varepsilon_2$. Notons aussi que
par la compatibilité

\begin{tikzcd}
\hat{\mathfrak{S}} \ar[r, hook] \ar[dr, "can"'] & M \ar[d, "(45)"] \\
& \hat{\mathfrak{S}}/\Pi_0
\end{tikzcd}

l'opération de $M$ sur $\hat{\mathfrak{S}}/\Pi_0$ déduite de
l'hom. (45) n'est autre que l'opération
déduite, par des autom. int. par passage au
quotient par $\Pi_0$, opération qui, pour $u \in M$,
est en particulier sur $M_0 \simeq M^\circ$, est
donnée par $\varepsilon_2(u)$ (cf p. 594). Ainsi

\newpage

%% ===== Page 645 =====

525

On voit que $\mathbb{M}_!^{'''}$ s'envoie dans le centre de $\hat{\mathcal{S}}\{\Pi_0\}$, qui est contenu dans $\hat{\mathcal{S}}^+/\Pi_0$, et en fait égal à l'image $\mu$ de $\hat{\pi}/\Pi_0$ dans $\hat{\mathcal{S}}^+/\Pi_0$ (puisque le centre de $\hat{\mathcal{S}}^+/\hat{\pi} \simeq \hat{\mathcal{S}}_g$ est trivial). Donc on a une factorisation
$$ \mathbb{M}_!^{'''} \longrightarrow \mu \hookrightarrow_{centre} \hat{\mathcal{S}}/\Pi_0 $$
et il faudrait vérifier que cet hom $\mathbb{M}_!^{'''} \to \mu$ n'est autre que $\varepsilon_3$, qui est donc trivial sur $\mathbb{M}_!^0$ et a fortiori sur $\mathbb{M}^0(0)$. [unclear: En fait. ce sera vrai que en remplaçant par $\operatorname{Ker}(\varepsilon_3|\mathbb{M}_!^{'''})$ (qui entier]

[MARGIN: - plutôt $\mu_3$ !) mais $\mathbb{M}_!^{'''}$ - on y reviendra plus bas.]

Admettons provisoirement ceci, pour trouver donc une factorisation de
$$ \mathcal{N}_{\hat{p}}^* \longrightarrow \hat{\mathcal{S}}/\Pi_0 $$
$$ (*) \quad \mathcal{N}_{\hat{p}}^*/\mathbb{M}^0(-j) \longrightarrow \hat{\mathcal{S}}/\Pi_0 $$
et comme l'hom. composé
$$ D_{\hat{p}} \hookrightarrow \mathcal{N}_{\hat{p}}^*/\mathbb{M}^0(-j) \longrightarrow \hat{\mathcal{S}}/\Pi_0 $$
est injectif, on voit que le noyau de $(*)$ est un sous-groupe de $\mathcal{N}_{\hat{p}}^*/\mathbb{M}^0(-j)$ dont l'intersection avec $D_{\hat{p}}$ est $\{1\}$ - il est donc d'ordre 1 ou d'ordre 2 (puisque

\newpage

%% ===== Page 646 =====

$D_{\underline{P}}$ est d'indice 2 dans $N_{\underline{P}}^* / M^{\cdot}(-J)$).

Je vais mieux énoncer la
Proposition Considérons la filtration
de $\widehat{S}/\Pi_0$ par $\widehat{S}/\widehat{S}^+ \simeq \mu$,
$\widehat{S}^+/\Pi \simeq \widehat{S}_3$, et $\Pi/\Pi_0 \simeq \mu$. Considé-
rons la filtration de $M$ par $M/M'' \simeq \mu^{\varepsilon_2}$,
par $M''/M^{0!} \simeq \widehat{S}_3$, et $M^{0!}/M^0 \simeq \mu^{\varepsilon_3}$.
Alors l'hom. can.
$$M \longrightarrow \widehat{S}/\Pi_0$$
est compatible avec les filtrations,
et induit des isomorphismes évidents
sur les facteurs. En particulier,
le noyau de $M \longrightarrow \widehat{S}/\Pi_0$ est $M^0$ !

Admettons ce point, on trouve
donc
$$\operatorname{Ker}(N_{\underline{P}}^* \longrightarrow \widehat{S}/\Pi_0) = N_{\underline{P}}^* \cap M^{0!}$$
et comme $N_{\underline{P}}^* \cap M^{0!} = Z_{\underline{P}}^{0!} = M^{\cdot}(-J)$,
on trouve que le noyau cherché est
égal : $M^{\cdot}(-J) \cap M'' = M^{\cdot}(-J)$, cqfd.
Donc l'hom en (*) est bien un hom.
injectif (*), et comme les groupes

\newpage

%% ===== Page 647 =====

en présence ont même ordre 24, cet hom.
est bien un iso, cqfd.

d) Application à la structure de $\mathcal{M}_{\rho, \sigma}$ et à
la définition d'un hom. [crossed out text] $\mathcal{M} \to \mathcal{M}_{\rho, \sigma}/\mu_2$.

Rappelons que
(48) $$\mathcal{M}_{\rho, \sigma} \overset{def}{=} N_{\rho}^* \times_{\gamma_{\rho, \sigma}} N_{\sigma}^* \times_{\gamma_{\rho, \sigma}} G_{\rho, \sigma}$$
[crossed out text]
(49) $$ \begin{cases} G_{\rho, \sigma} \simeq \mathcal{M} = \underbrace{\mathcal{M}_{0,3}}_{\simeq \mathcal{M}_{0,3}(0)} \cdot \widehat{\mathfrak{S}}^{+\sim} \\ \gamma_{\rho, \sigma} \overset{df}{=} G_{\rho, \sigma}/\widehat{\mathfrak{S}}^{+\sim} \simeq \Gamma_Q^{\sim} (\simeq \mathcal{M}(0)) , \\ N_{\rho}^* = \operatorname{Norm}_{\mathcal{M}}(L_{\rho}), \quad N_{\sigma}^* = \operatorname{Norm}_{\mathcal{M}}(L_{\sigma}) \end{cases}$$

D'autre part, nous avons introduit les
sous-groupes remarquables
(50) $$ \begin{cases} Z_{\sigma}^! = \mathcal{M}(1/2) \subset N_{\sigma}^* \quad (p. 615 \ (9)) \\ N_{\rho}^? = Z_{\rho}^! \cdot \mu_2 \subset N_{\rho}^* \quad (p. 616 \ (21)) \end{cases} $$

Le premier de ces groupes ne contient
pas $\mu_2$, le deuxième le contient :
(51) $$ \begin{cases} Z_{\sigma}^! \cap \mu_2 = \{ 1 \} \\ N_{\rho}^? \supset \mu_2 \end{cases} $$
et
(52) $$ \begin{cases} Z_{\sigma}^! \overset{\sim}{\to} \gamma_{\rho, \sigma} = \Gamma_Q^{\sim} \\ \text{[crossed out text]} 1 \to \mu_2 \to N_{\rho}^? \to \gamma_{\rho, \sigma} \to 1 \end{cases} \quad \text{exactes} $$

\newpage

%% ===== Page 648 =====

Soit alors

(53) $$ \mathcal{M}_{\rho,\sigma}^\natural \subset \mathcal{M}_{\rho,\sigma} $$
$$ ||def $$
$$ N_\rho^? \times_{\gamma_{\rho,\sigma}} Z_0^! \times_{\gamma_{\rho,\sigma}} G_{\rho,\sigma} $$
(avec $\simeq \mathcal{M}$ sous $G_{\rho,\sigma}$)

On a un [unclear: sous-groupe central] canonique

[unclear: texte raturé]

(54) $$ \mu \times \mu \times \mu \hookrightarrow \mathcal{M}_{\rho,\sigma} $$
$$ (\varepsilon, \varepsilon', \varepsilon'') \longmapsto (\varepsilon, \varepsilon', \varepsilon'') $$

provenant des plongements canoniques

$$ \mu \hookrightarrow N_\rho^*, \quad \mu \hookrightarrow N_\sigma^*, \quad \mu \hookrightarrow G_{\rho,\sigma} $$
(avec une flèche sous le dernier terme pointant vers $\hookrightarrow \tilde{G}_{\rho,\sigma}^+ \simeq \tilde{\mathcal{G}}$)

et on a

$$ \mathcal{M}_{\rho,\sigma}^\natural \cap (\mu \times \mu \times \mu) = (1 \times \mu \times \mu) $$

d'où

(55) $$ \mu \times \mu \hookrightarrow \mathcal{M}_{\rho,\sigma}^\natural $$

On désigne par $\mu_\rho$ le sous-groupe central $\mu$ de $N_\rho$ - et $N_\rho^?$, par [unclear] dans $\mathcal{M}_{\rho,\sigma}^\natural$. Il est clair alors par (52) que l'on a une suite exacte

(55') $$ 1 \longrightarrow \mu_\rho \longrightarrow \mathcal{M}_{\rho,\sigma}^\natural \longrightarrow G_{\rho,\sigma} \longrightarrow 1 $$
(avec $\simeq \mathcal{M}$ sous $G_{\rho,\sigma}$)

d'un des types canoniques

647

\newpage

%% ===== Page 649 =====

(56)
$$M \overset{\sim}{\longrightarrow} M_{\rho, \sigma}^\natural/\mu_{\rho} \hookrightarrow M_{\rho, \sigma}/\mu_{\rho}$$

[MARGIN: On posera donc
$M_{\rho, \sigma}^+ = M_{\rho, \sigma}^\natural/\mu_{\rho}$]

Evidemment le sous-groupe $M_{\rho, \sigma}^\natural$ de $M_{\rho, \sigma}$
contient $G_{\rho, \sigma}^+ \simeq S G_{\rho, \sigma}$ (qui en fait...
en fait on a
un diagramme de suites exactes

(57)
\begin{tikzcd}
1 \ar[r] & \mu_{\rho} \times G_{\rho, \sigma}^+ \ar[r] \ar[d] & M_{\rho, \sigma}^\natural \ar[r] \ar[d] & \Gamma_{\rho, \sigma}^\natural \ar[r] \ar[d] & 1 \\
1 \ar[r] & \underset{L_\rho}{S\hat{\mathbb{Z}}} \times \underset{L_\sigma}{S\hat{\mathbb{Z}}} \times G_{\rho, \sigma}^+ \ar[r] & M_{\rho, \sigma} \ar[r] & \Gamma_{\rho, \sigma} \ar[r] & 1
\end{tikzcd}

Dans l'isomorphisme (56), le sous-groupe
$G^{+ \wedge}$ de $M$ correspond au sous-groupe
$\mu_\rho G_{\rho, \sigma}^+/\mu_\rho \simeq G_{\rho, \sigma}^+$, il s'envoie isomorphiquement
sur le sous-groupe $G_{\rho, \sigma}^+$ de $M_{\rho, \sigma}/\mu_{\rho}$.
On trouve ainsi

(58)
\begin{tikzcd}
& & & \Gamma_{\rho, \sigma} \ar[d, "\simeq"'] & \\
1 \ar[r] & G^{+ \wedge} \ar[r] \ar[d, "\simeq"'] & M \ar[r] \ar[d, "\simeq"'] & \hat{\Gamma}_Q \ar[r] \ar[d, "\simeq"'] & 1 \\
1 \ar[r] & G_{\rho, \sigma}^+ \ar[r] \ar[d, equal] & M_{\rho, \sigma}^\natural/\mu_\rho \ar[r] \ar[d, hook] & \Gamma_{\rho, \sigma}^\natural/\mu_\rho \ar[r] \ar[d, hook] & 1 \\
1 \ar[r] & G_{\rho, \sigma}^+ \ar[r] & M_{\rho, \sigma}/\mu_\rho \ar[r] & \Gamma_{\rho, \sigma}/\mu_\rho \ar[r] & 1
\end{tikzcd}

Ainsi, se trouve prouvé que l'hom.
canonique épi.

(59)
$$1 \to \frac{S\hat{\mathbb{Z}} \times S\hat{\mathbb{Z}}}{L_\rho \times L_\sigma} \to \Gamma_{\rho, \sigma} \to \bar{\Gamma}_{\rho, \sigma} \to 1$$

admet un "bisecteur" $\Gamma_{\rho, \sigma}^\natural$ (dont
l'intersection avec $\frac{S\hat{\mathbb{Z}} \times S\hat{\mathbb{Z}}}{L_\rho \times L_\sigma}$ est $\mu_\rho \times 1$)

\newpage

%% ===== Page 650 =====

$$
(60) \left\{ \begin{array}{l} \Gamma_{\rho,\sigma} \simeq \tilde{\Gamma}_{\rho,\sigma} \cdot \underset{\hat{\mathbb{Z}}\rho \times \hat{\mathbb{Z}}\sigma = L_\rho \times L_\sigma}{S\Gamma_{\rho,\sigma}} \quad , \quad \tilde{\Gamma}_{\rho,\sigma} \cap S\Gamma_{\rho,\sigma} = \mu_\rho \\ \Gamma_{\rho,\sigma} / \mu_\rho \simeq (\tilde{\Gamma}_{\rho,\sigma} / \mu_\rho) \cdot (L_\rho / \mu_\rho \times L_\sigma) \end{array} \right.
$$
$$
\gamma_{\rho,\sigma} \simeq \tilde{\Gamma}_Q
$$

Pour dire ce que deviennent les sous-
groupes remarquables
$$M(0), M[0], M'[0], M''[0]$$
$$M^G[0] \quad M^{\prime G}[0] \quad M^{\prime\prime G}[0]$$
$$M \rtimes L_0 \quad M \rtimes L'_0 \quad M \rtimes L''_0$$
$\mathcal{M}(1/2)$ [unclear: $Z^\circ_0 !$], $\mathcal{M}'(-j)$, $\mathcal{M}''(-j)$, $N_{\hat{\rho}}^?$ de $M$ par
les homomorphismes $M \overset{\sim}{\rightarrow} M_{\rho,\sigma}/\mu_\rho \hookrightarrow M_{\rho,\sigma}$,
il faudrait d'abord introduire les
groupes adéquats dans les $M_{\rho,\sigma}$ génériques.
Nous allons donc revoir un peu ces groupes
maintenant.

\newpage

%% ===== Page 651 =====

VI Les groupes $M_{\rho,\sigma}$ génériques: récapitulation, raffinement.

1º) Données. Ce sont celles de V 1º) (p. 577).
En plus de $\widehat{G}$ (pro-fini) et du $\operatorname{Gl}(2,\mathbb{Z})^\wedge \to \widehat{G}$
hom. continu surjectif, i.e. des générateurs
$\rho, \sigma, \tau$ (ou $\rho, \sigma, \tau^\prime$) de $\widehat{G}$, on se donne
aussi un sous-groupe ouvert
(1) $$ \widehat{G}^\natural \subset \widehat{G}^+ $$ , invariant dans $\widehat{G}$.
Dans le cas universel $\widehat{G} = \operatorname{Gl}(2,\mathbb{Z})^\wedge$,
la donnée de $\widehat{G}^\natural$ peut correspondre au
point de vue où on regarde le sous-
groupe $M^\natural$ de $M$ [unclear: crossed out text]
et qui opèrent trivialement dans $\widehat{G}^+ / \widehat{G}^\natural$ - c'est
[unclear: crossed out text] ouvert, et définit
dans $\Gamma_{\mathbb{Q}}^\sim = \mathcal{M} / \widehat{G}^+$ un sous-groupe [unclear]
$\Gamma_{\mathbb{Q}}^{\natural \sim}$ ouvert, dont on peut se proposer d'étu-
dier les représentations dans des groupes
profinis convenables. Un cas intéressant
sera celui où
$$ \widehat{G}^\natural = \Pi_0 \text{, donc } M^\natural = M^{\prime\prime}, $$
ou $$ \widehat{G}^\natural = \Pi \text{, donc } M^\natural = M^\prime . $$
Si $$ \widehat{G}^\natural = \widehat{G} \text{, donc } M^\natural = M , $$
on retrouve les constructions précédentes.

\newpage

%% ===== Page 652 =====

Un cas intéressant [unclear] est celui où $\mathfrak{G}$ est le groupe des pts $\cdot$ adéliques continus d'un schéma en groupes sur $\mathbb{Z}$ (dans ce cas, il est vrai que $\rho, \sigma, \tau$ n'engendrent pas tjs $\mathfrak{G}$ que modulo groupe fini, i.e. ils engendrent p.ê. à eux seuls un sous-groupe ouvert. Dans ce cas, il faudrait plutôt paraphraser les constructions qui suivent dans le cas schématique. J'y viendrai par la suite. Un choix intéressant d'un $\mathfrak{G}^{\natural}$ est alors la composante neutre de $\mathfrak{G}$.

J'hésite à introduire des notations $Z_{\mathfrak{G}}^{\natural}, M_{\mathfrak{G}}^{\natural}$ etc - pour ne pas surcharger des notations déjà lourdes. Donc je préfère utiliser les notations plus simples $Z_{\mathfrak{G}}, M_{\mathfrak{G}}$ etc - et sous-entendre que dans ces constructions, le groupe $\mathfrak{G}^{\natural}$ intervient - quitte à réintroduire la notation $\natural$ en cas de besoin.

On pose
(2) $\mathfrak{G}_{\pm}^{\natural} = \{ \pm 1 \} \cdot \mathfrak{G}^{\natural} \subset \mathfrak{G}$.

\newpage

%% ===== Page 653 =====

2°) $Z_{\rho,\sigma}$, $\gamma_{\rho,\sigma}$.

Soit

(4) $$Z_{\rho,\sigma} \subset \operatorname{Aut}(\hat{G}^+) \times \hat{\mathbb{Z}}^* \times \mu \times \mu$$
$$|| dfn$$
$$(u,\mu,\varepsilon_2,\varepsilon_3) \begin{cases} a) \exists \alpha \in \widehat{\mathfrak{S}}_g, \text{ avec } \varepsilon_2(\alpha)=\varepsilon_2, u(\rho)=\alpha(\rho) \\ b) \exists \beta \in \widehat{\mathfrak{S}}_g, \text{ avec } \varepsilon_3(\beta)=\varepsilon_3, u(\sigma)=\beta(\sigma) \\ c) u(\varepsilon_0) = \varepsilon_0^\mu \\ d) u \text{ normalise } \widehat{\mathfrak{S}}_g / Im \pi_1 \text{ et induit l'identité sur } (\widehat{\mathfrak{S}}_g / Im \pi_1) \end{cases}$$

[MARGIN left: cf dem suiv/ condition partic dans d) ([unclear]) que $\varepsilon_0=1$ [unclear] se traduit par $\mu = \varepsilon_3 \varepsilon_0 / \sigma$]
[MARGIN right: (65) i.e. modulo le th. d'isom. une cert. condition]

Proposition. $Z_{\rho,\sigma}$ est un sous-groupe du groupe produit.

Dém. Les conditions c), d) définissent un sous-groupe.

la condition b) signifie aussi que

(5) b') $u(\sigma)$ est $\widehat{\mathfrak{S}}_g$-conjugué à $\sigma^{\varepsilon_3}$

et sous cette forme il est clair que la condition c') définit un sous-groupe. Je ne pense pas que la condition a) définisse à elle seule un sous-groupe. Mais si $\mathscr{G}_{(b)}$ désigne le sous-groupe du produit défini par les conditions b'),..., on a dans $\mathscr{G}_{(b)}$ un sous-groupe de $\mathscr{G}_{(b)}$ qui provient... Soient $(u,\mu,\varepsilon_2,\varepsilon_3)$ et $(u',\mu',\varepsilon_2',\varepsilon_3')$, et $\alpha, \alpha' \in \widehat{\mathfrak{S}}_g$, de sorte que l'on a ...

\newpage

%% ===== Page 654 =====

(6) $u(\rho) = \alpha(\rho) \quad , \quad u'(\rho) = \alpha'(\rho) \quad , \quad \text{avec } \begin{matrix} \varepsilon_{\underline{C}}(\alpha) = \varepsilon_3 \\ \varepsilon_{\underline{C}}(\alpha') = \varepsilon'_3 \end{matrix}$ .

Prouvons que $u'u \in G_{(a,b)}$ - i.e. qu'il
satisfait à son tour à (b). Si $\varepsilon_3 = 1$ i.e. $\alpha \in \widehat{\underline{C}}_{\underline{G}}^+$
$u'(\alpha)$ est défini et on a

$(u'u)(\rho) = u'(u(\rho)) = u'(\alpha(\rho)) = u'(\alpha) (u'(\rho)) = u'(\alpha) (\alpha'(\rho))$
$= (u'(\alpha) \alpha') (\rho) = \alpha'' (\rho) \quad , \quad \text{avec}$

(7) $\alpha'' = u'(\alpha) \alpha' \in \widehat{\underline{C}}_{\underline{G}} \qquad \text{(car } \varepsilon_3 = 1$
$\varepsilon_{\underline{C}}(\alpha'') = \varepsilon_{\underline{C}}(\alpha') = \varepsilon'_3 = \varepsilon_3 \varepsilon'_3$
(puisque $u'(\alpha) \in \widehat{\underline{C}}_{\underline{G}}^+ \implies \varepsilon_{\underline{C}}(u'(\alpha)) = 1$, et $\varepsilon_3=1$.)

Dans ce cas on a gagné. Si $\varepsilon_3 = -1$, on
écrit

$\alpha = \alpha_0 \tau \qquad \alpha_0 \in \widehat{\underline{C}}_{\underline{G}}^+$

et on aura

$u(\rho) = (\alpha_0 \tau)(\rho) = \alpha_0 (\tau(\rho)) = \alpha_0 (\sigma^{-1}(\rho^{-1})) = (\alpha_0 \sigma)(\rho^{-1})$

où maintenant $\alpha_0 \sigma \in \widehat{\underline{C}}_{\underline{G}}^+$, donc $u'(\alpha_0 \sigma)$ est
défini, et on aura

$(u'u)(\rho) = u'(u(\rho)) = u'(\alpha_0 \sigma) (\underset{\alpha'(\rho^{-1})}{u'(\rho^{-1})}) = (u'(\alpha_0 \sigma) \alpha') (\rho^{-1})$

et comme $\rho^{-1} = \sigma \tau (\rho)$
on trouve

(8) $(u'u)(\rho) = \alpha'' (\rho)$
où
$$ \begin{cases} \alpha'' = \varepsilon u'(\alpha_0) u'(\sigma) \alpha' \sigma \tau \\ \varepsilon \in \mu = \{\pm 1\} \subset \text{centre } \widehat{\underline{C}}_{\underline{G}}^+ \end{cases} $$
qu'on aura choisi de façon que $\alpha''$ (si possible) soit
dans $\widehat{\underline{C}}_{\underline{G}}^+$. Notons que tous les facteurs de
ce produit sont dans $\widehat{\underline{C}}_{\underline{G}}^+$, sauf $\tau$, et donc on a $\alpha'' \in \widehat{\underline{C}}_{\underline{G}}$.

\newpage

%% ===== Page 655 =====

Réduisons mod $\widetilde{\mathfrak{S}}_\tau^+$, la relation $\alpha'' \in \widetilde{\mathfrak{S}}_\tau^+$ s'écrit
$\dot{\alpha}'' \in \{1, \dot{\tau}\}$ i.e. $\dot{\varepsilon}\ u'(\sigma)^\bullet \dot{\alpha}' \dot{\sigma} \dot{\tau} \cdot \dot{\tau} \in \{1, \dot{\tau}\}$ soit $\dot{\varepsilon} u'(\sigma)^\bullet \dot{\alpha}' \dot{\sigma} \in \{1, \dot{\tau}\}$
[unclear: on distingue deux cas]

1°) $\varepsilon'_3 = 1$ i.e. $\dot{\alpha}' = 1$, la relation s'écrit
$(*) \qquad \dot{\varepsilon}\ u'(\sigma)^\bullet \dot{\sigma} = 1$

2°) $\varepsilon'_3 = -1$ i.e. $\dot{\alpha}' = \dot{\tau}$, comme $\tau \sigma = - \sigma \tau$, la relation
s'écrit
$(**) \qquad -\dot{\varepsilon}\ u'(\sigma)^\bullet \dot{\sigma} = 1$ ,

dans les deux cas, la relation s'écrit
(9) $\qquad u'(\sigma)^\bullet = (\varepsilon \varepsilon'_3)^\bullet \dot{\sigma}^{-1}$

[crossed out: supposons maintenant qu'on a b) pour $u'$]
On utilise l'hyp. b) pour $u'$, qui s'écrit
$u'(\sigma)^\bullet = \dot{\varepsilon}'_2 \dot{\sigma}^{-1}$

on voit donc qu'il suffit de choisir $\varepsilon$ de façon que $\varepsilon \varepsilon'_3 = \varepsilon'_2$ i.e.
$\varepsilon = \varepsilon'_2 \varepsilon'_3$

dans ce cas,
[crossed out: (4)] $\qquad \alpha'' = \varepsilon'_2 \varepsilon'_3\ u'(\alpha \tau) \alpha' \sigma \tau$

i.e.
(10) $\qquad \alpha'' = \varepsilon'_2 \varepsilon'_3\ u'(\alpha \tau)\ (\alpha' \tau) \qquad | \text{ cas } \varepsilon_3 = -1$

formule qui a un sens, car $\alpha \tau \in \widetilde{\mathfrak{S}}^+$ donc
$u'(\alpha \tau)$ est défini, et l'argument précédent
prouve que l'élément en question satisfait
(11) $\qquad \alpha'' \in \widetilde{\mathfrak{S}}_\tau^+, \varepsilon_{\widetilde{\mathfrak{S}}}(\alpha'') = \varepsilon_{\widetilde{\mathfrak{S}}}(\alpha)\ \varepsilon_{\widetilde{\mathfrak{S}}}(\alpha')$
en distinguant les cas (*), (**) de (9).
Mème conclusion dans le cas premier où $\varepsilon_{\widetilde{\mathfrak{S}}}(\alpha) = 1$ (7).
Cela achève la démonstration.

\newpage

%% ===== Page 656 =====

Remarques 1) Pour que la condition $b)$ soit
satisfaite, on n'a pas besoin de la condition
$c)$ dans l'hypothèse fixe, il suffit que $u$ dans $\mathcal{G} \xrightarrow{\sim} \mathcal{G}$
normalise l'image de $L_0$ i.e. qu'on ait
$$u(\sigma) \equiv \pm \sigma \quad (2 \mathcal{G})$$

2) Si $\mathcal{G}$ est "sans gras", l'hom. de
projection est injectif
(12) $$Z_{\mathcal{G},\sigma} \hookrightarrow \operatorname{Aut}(\mathcal{G}^+)$$ (cf p. 578)

Exemples dans les cas universels
a) $\mathcal{G} = \mathcal{G}^{\pm} = \widehat{Sl}(2, \mathbb{Z})^{\wedge}$. On a vu la
condition (4) d) est vide, les conditions
a), b) [unclear] [unclear] en fait que
$u(\rho)$ conjugué à $\rho^{\pm 1}$, $u(\sigma)$ conjugué à $\sigma^{\pm 1}$, i.e.
[unclear: on retombe] sur la définition (b) p. 578, qui
donne
$$Z_{\mathcal{G},\sigma} = \mathcal{M}(0) . L_0^{\pm} = \mathcal{M}^{\pm}[0].$$

[MARGIN: NB Application - plutôt
$Z_{\mathcal{G},\sigma} \subset M_{0,3}$
$\widetilde{M}_{0,3} \simeq M_{0,3}[0]$
$\widetilde{M}(0) . L_0 \simeq \widetilde{M}(0) . L_0$]

(13) $$Z_{\mathcal{G},\sigma} = \mathcal{M}(0) . L_0^{\pm} = \mathcal{M}^{\pm}[0] = M[0]$$

[MARGIN: dans ce cas il est immédiat de trouver au moins trinôme b) $\dots Z_{\mathcal{G},\sigma} \to G \xrightarrow{\sim} \operatorname{Aut}_{ext} \dots$]

b) $\mathcal{G} = \Pi = \Pi_0 \times \mu$. On trouve
(14) $$Z_{\mathcal{G},\sigma} = M(0) . L_0^{\Pi} = M^{\Pi}[0] = M'[0]$$
(c'est le $Z'_{1,0}$ de V $\S$ 8).
(On a $M'[0] \simeq M''[0]$)

[MARGIN: expliciter l'inverse [unclear: équivalences] qui peut être [unclear] (1,1)]

c) $\mathcal{G} = \Pi_0$. Les conditions a) et b)
résultent des cas précédents, puisque
$\Pi = \Pi_0 \times \mu$ et $\mu$ est central, pour ce qui est de (4).

\newpage

%% ===== Page 657 =====

condition d) devient nettement plus stricte, on
trouve par ex. que $Z_{\widehat{\mathbb{S}}, 0} \cap \widehat{\mathbb{S}}^+ = L_0^{\Pi_0}$ (d'indice 2 dans
le groupe correspond. $L_0^\pi$ du cas précédent),
et comme l'opération de $M(0)$ sur $\widehat{\mathbb{S}}^+ / \Pi_0$
se fait par le quotient $\varepsilon_2$, on voit que l'on a

[MARGIN: maintenant la trivialisation est plus petite !]

$$Z_{\widehat{\mathbb{S}}, 0} = M''(0) \cdot L_0^{\Pi_0} = M''[0]_0$$

[MARGIN: il sera aussi utile de définir $Z_{\widehat{\mathbb{S}}, 0}$ en n'imposant d) que sur $\widehat{\mathbb{S}}^+$, c'est-à-dire une condition de commutativité dans le groupe [unclear] (qui s'envoie sur $\widehat{\mathbb{S}}^+/\Pi_0 = \mu$). Dans ce cas on devrait obtenir $Z_{\widehat{\mathbb{S}}, 0} \cdot \pi = M[0]_0 \cdot \pi$. (Mais je vois sur le moment que la minoration de d) est inutile, compte tenu de d))]

[MARGIN: cette remarque sur [unclear] devient inutile]

[MARGIN: Mais on a d'ailleurs de façon inconditionnelle à (4) (en utilisant la cond. b) précédente pour $M''[0]_0$) on aurait plutôt $Z_{\widehat{\mathbb{S}}, 0} = M''[0]_0$.]

$$(15) \begin{cases} M''(0) = \operatorname{Ker} (M(0) \xrightarrow{\varepsilon_2} \mu) \\ M''[0]_0 = \operatorname{Ker} (M[0]_0 \xrightarrow{\varepsilon_2} \mu) \\ M[0]_0 = M(0) \cdot L_0^{\Pi_0} \end{cases}$$

Ces deux derniers exemples suggèrent l'opportunité
dans le cas où on disposerait d'une [unclear: décomposition] 
en produit direct :

(16) $$\widehat{\mathbb{S}} = \widehat{\mathbb{S}}^+ \times \mu \quad \text{(où par ex. le sous-groupe central est } \mu = \{1, \rho^2\} = \{1, \sigma^2\})$$

de définir $Z_{\widehat{\mathbb{S}}, 0}$ en termes de $\widehat{\mathbb{S}}$ et des
choix de $\alpha, \beta$ de la définition (4)
dans $\widehat{\mathbb{S}}$ :

(17) $$\widehat{\mathbb{S}}_\tau \stackrel{dfn}{=} \mu_\tau \cdot \widehat{\mathbb{S}}^+ .$$

Remarque On est tenté de remplacer la
déf (4) a) b) par la définition d'apparence plus
simple de la p. 578, qui rend évident qu'on a
[unclear]
Mais elle n'est pas compatible
avec $\dots$

\newpage

%% ===== Page 658 =====

qui en se restreignant : dans $(u, \mu, \varepsilon_1, \varepsilon_3)$ ou $\varepsilon_3 = 1$,
ce $j^{-1}$ n'est pas égal à $j$ ! Il est vrai que
dans le cas universel p.ex.) [unclear: que] $u(j)$ est
conjugué dans $\underline{G}^+$ à $j^{\varepsilon_3}$, mais en général
il ne le sera pas par un élément (disons)
de $\pi$.

[MARGIN: NB On a
$L_0^{\underline{G}} \supset L_0 \cap \underline{G}^+$ et l'
inclusion peut être stricte,
(exemple p.ex. f)
ci-dessus... $\underline{G}^+ = L_0^{\pi_0}$
et $L_0^+ = L_0^\pi, L_0 \cap \underline{G}^+ = L_0^{\pi_0}$)
ce qui
s'y [unclear]]

On pose
(18) $L_0^{\underline{G}} = \{ l \in \underline{L}_0^{\underline{G}} | (int(l), 1, 1, 1) \in Z_{\rho,\sigma} \}$
(i.e. $L_0^{\underline{G}} \subset \underline{G}^+$ est l'ens
des $\varepsilon_0^n, n \in \hat{\mathbb{Z}}$)
et on a un hom. canonique

(19) $L_0^{\underline{G}} \longrightarrow Z_{\rho,\sigma}$
$$l \longmapsto (int(l), 1, 1, 1)$$

in-jectif ss.

(20) $L_0^{\underline{G}} \cap Centr(\underline{G}^+) = \{ 1 \}$

l'image invariante. On pose

(21) $\Gamma_{\rho,\sigma} = Z_{\rho,\sigma} / L_0^{\underline{G}}$ ,

d'où

(22) $$1 \longrightarrow L_0^{\underline{G}} \longrightarrow Z_{\rho,\sigma} \overset{\delta_Z}{\longrightarrow} \Gamma_{\rho,\sigma} \longrightarrow 1$$ ,

(23) 
\begin{tikzcd}
Z_{\rho,\sigma} \ar[r] & \Gamma_{\rho,\sigma} \ar[r, "\chi_\Gamma"] \ar[dr, "\varepsilon_\Gamma"'] & \hat{\mathbb{Z}}^\times \\
& & \mu
\end{tikzcd}

déduit de $(u, \mu, \varepsilon_1, \varepsilon_3) \longmapsto (\varepsilon_3 \chi, \mu)$ .

Remarque. On désignera par
(24) $\Gamma_{\rho,\sigma}' = \operatorname{Ker}(\varepsilon_\Gamma), \Gamma_{\rho,\sigma}'' = \operatorname{Ker}(\chi_\Gamma), \Gamma_{\rho,\sigma}''' = \operatorname{Ker}(\varepsilon_\Gamma, \chi_\Gamma)$
$Z_{\rho,\sigma}' = \dots, Z_{\rho,\sigma}'' = \dots, Z_{\rho,\sigma}''' = \dots$ (l'image inverse dans $Z_{\rho,\sigma}$)

\newpage

%% ===== Page 659 =====

On pose
$$ \gamma'_{\rho,\sigma} = \text{Ker } \varepsilon_3 \gamma , \gamma''_{\rho,\sigma} = \text{Ker } \varepsilon_1 \gamma , \gamma'''_{\rho,\sigma} = \text{Ker } \varepsilon_2 \rho \varepsilon_3 \gamma , \gamma^\circ_{\rho,\sigma} = \gamma' \cap \gamma'' \cap \gamma''' = \gamma''' \cap \gamma' $$
$$ Z'_{\rho,\sigma}, Z''_{\rho,\sigma}, Z'''_{\rho,\sigma}, Z^\circ_{\rho,\sigma} \text{ sont les images inv. ds } Z_{\rho,\sigma} \text{ de } \gamma' \text{ etc.} $$

Ce sont donc des sous groupes d'indices [unclear: resp 2] ou quatre de $\gamma_{\rho,\sigma}$ resp $Z_{\rho,\sigma}$ - l'indice 4 ne pouvant se présenter que pour $\gamma^\circ_{\rho,\sigma}, Z^\circ_{\rho,\sigma}$.

Remarques 1) On a donc
$$ Z''_{\rho,\sigma} = \left\{ (u,\mu) \in \operatorname{Aut}(\widehat{G}^+) \times \hat{\mathbb{Z}}^* \left| \begin{array}{l} u(\rho) \text{ est } \widehat{G}\text{-conj. à } \rho \\ u(\sigma) \text{ \quad --- \quad --- \quad } \sigma \\ u(\varepsilon_0) = \varepsilon_0 \mu \end{array} \right. \right\} $$

[MARGIN: Plus généralement $Z'''_{\rho,\sigma} = \left\{ (u,\mu,\xi) \in \operatorname{Aut}(\dots) \times \dots \mid a)\ c)\ e)\ u(\rho) \equiv \rho \pmod{\widehat{G}^\natural} ; f)\ u(\varepsilon_0) = \varepsilon_0 \mu \right\}$]

[unclear: indépendamment des sous-groupes]. La condition d) de (4) - stabilité de $\widehat{G}^\natural$ sous $u$ - est conséquence ici de a) b) c), car on aura
$$ u(\rho) \equiv \rho \pmod{\widehat{G}^\natural} $$
$$ u(\sigma) \equiv \sigma \pmod{\widehat{G}^\natural} $$
ce qui implique
(26 bis) $$ u(x) \equiv x \pmod{\widehat{G}^\natural} \quad \forall x \in \widehat{G}^+ $$
et la stabilité de $\widehat{G}^\natural$ et même l'identité ds $\widehat{G}^+ / \widehat{G}^\natural$.

[MARGIN: c'est un vrai non-sens de ma part ! (cf 4) d) en bas de la page ...) il n'y a pas d'élément de $\widehat{G}^+$ de parité -1 !]

2) On a encore l'élément remarquable
(26) $$ \tau_Z \in Z_{\rho,\sigma} , \quad \tau_Z = (\tau, -1, -1, -1) $$

3) Les groupoïdes $G_{\rho,\sigma}, G_{\rho,\sigma}', S_0 G_{\rho,\sigma}, G^\natural_{\rho,\sigma}$.

Procedant comme page 577, on trouve
$$ Z_{\rho,\sigma} \to \operatorname{Aut}(\widehat{G}^+) $$
d'où deux produits $\frac{1}{2}$ directs
(27) $$ S_0 G_{\rho,\sigma} = Z_{\rho,\sigma} \cdot \widehat{G}^+, \quad S_0 G^\natural_{\rho,\sigma} = Z_{\rho,\sigma} \cdot \widehat{G}^\natural $$
et un hom.
(28) $$ L^\natural_0 \hookrightarrow S_0 G_{\rho,\sigma} \subset S_0 G_{\rho,\sigma}' $$
$$ l \mapsto (\text{int}(l), l^{-1}) $$

\newpage

%% ===== Page 660 =====

d'image invariante, ce qui permet de poser

$$ (29) \quad G_{\rho, \sigma} \overset{\text{déf}}{=} S_0 G_{\rho, \sigma} / \operatorname{Im} L_0^\sigma \supset G_{\rho, \sigma}^\sigma = S_0 G_{\rho, \sigma}^\sigma / \operatorname{Im} L_0^\sigma $$
$$ Z_{\rho, \sigma} \cdot C^+ $$

d'où suites exactes

(30)
\begin{tikzcd}
1 \ar[r] & L_0^\sigma \ar[r] \ar[d, equal] & S_0 G_{\rho, \sigma}^\sigma \ar[r] \ar[d, hook] & G_{\rho, \sigma}^\sigma \ar[r] \ar[d, hook] & 1 \\
1 \ar[r] & L_0^\sigma \ar[r] & S_0 G_{\rho, \sigma} \ar[r] & G_{\rho, \sigma} \ar[r] & 1
\end{tikzcd}

et un diag. commutatif

(31)
\begin{tikzcd}
& Z_{\rho, \sigma} \ar[dr, "i_Z"] \\
L_0^\sigma \ar[ur, "i_{L_0, Z}"] \ar[dr, "i_{L_0, C^\sigma}"'] & & G_{\rho, \sigma}^\sigma \ar[r, hook] & G_{\rho, \sigma} \\
& C^\sigma \ar[ur, "i_{C^\sigma}"] \ar[rr] & & C^+ \ar[u, "i_{C^+}"']
\end{tikzcd}

s'insérant dans des diag. commutatifs et
suites exactes

(32)
\begin{tikzcd}
1 \ar[r] & L_0^\sigma \ar[r] \ar[d] & Z_{\rho, \sigma} \ar[r, "\delta_Z"] \ar[d] & \gamma_{\rho, \sigma} \ar[r] \ar[d, equal] & 1 \\
1 \ar[r] & C^\sigma \ar[r] \ar[d] & G_{\rho, \sigma}^\sigma \ar[r, "\delta_G"] \ar[d] & \gamma_{\rho, \sigma} \ar[r] \ar[d, equal] & 1 \\
1 \ar[r] & C^+ \ar[r] & G_{\rho, \sigma} \ar[r] & \gamma_{\rho, \sigma} \ar[r] & 1
\end{tikzcd}

On considère les opérations

(33)
\begin{tikzcd}
G_{\rho, \sigma} \ar[r, "\delta_G"] \ar[dr, "\chi_G"'] & \gamma_{\rho, \sigma} \ar[r, "\varepsilon_\gamma"] \ar[d, "\chi_\gamma"] & \{ \pm 1 \} \\
& \hat{\mathbb{Z}}^* \ar[ur, "\varepsilon"'] &
\end{tikzcd}

avec
$$ \chi_G = \chi_\gamma \delta_G, \quad \varepsilon_G = \varepsilon_\gamma \delta_G, \quad \varepsilon_\gamma = \varepsilon \chi_\gamma $$

On a maintenant

[MARGIN: idem pour $G_{\rho, \sigma}^+$...]
(34)
$$ G_{\rho, \sigma}', G_{\rho, \sigma}'', G_{\rho, \sigma}''', G_{\rho, \sigma}^\circ \quad \text{comme dans } (24) $$
$$ G_{\rho, \sigma}^+ = \operatorname{Ker} \chi_G, \ G_{\rho, \sigma}'^+, G_{\rho, \sigma}''^+, G_{\rho, \sigma}'''^+, G_{\rho, \sigma}^{\circ +} $$
$$ (\gamma_{\rho, \sigma}^+ = \operatorname{Ker} \chi_\gamma \text{ etc.} \dots) $$

Si
$$ (35) \quad \tau_G \in G_{\rho, \sigma}^\sigma, \tau_\gamma \in \gamma_{\rho, \sigma} $$

\newpage

%% ===== Page 661 =====

633

désignons les images de $\tau_{\Sigma}$ (26) dans $G_{\rho, \sigma}$, $\gamma_{\rho, \sigma}$,
-- pour faire l'identification

(36)
$$ \begin{cases} \underline{G}^{\natural} \simeq \text{image inverse de } \{1, \tau_{\Sigma}\} \subset \gamma_{\rho, \sigma} \text{ dans } G_{\rho, \sigma} \text{, par } G_{\rho, \sigma} \to \gamma_{\rho, \sigma} \\ \underline{G}^{\natural} \simeq \text{image inv. de } \{1, \tau_{\Sigma}\} \subset \underline{G}_{\rho, \sigma} \text{ dans } G^{\natural}_{\rho, \sigma} , \end{cases} $$

et

(37) $$ \xi_{\underline{G}} = \chi_G | \underline{G}^{\natural} = \varepsilon_Z | \underline{G}^{\natural} = \varepsilon_G | \underline{G}^{\natural} : \underline{G}^{\natural} \to \begin{cases} \text{triv sur } \underline{G}^+ \\ -1 \text{ sur } \tau \dots \end{cases} $$

Considérons l'application canonique

(38)
$$ G_{\rho, \sigma}^{\natural} \to \operatorname{Aut}(\underline{G}^+) \times \hat{\mathbb{Z}}^* \times \mu \times \mu $$
$$ \underline{u} \longmapsto (int(u) | \underline{G}^+, \chi_G(u), \varepsilon_Z(u), \varepsilon_G(u)) $$

Proposition C'est un hom. de groupes. L'image de $G_{\rho, \sigma}^{\natural}$ est formée des systèmes 
$$ (u, \mu, \varepsilon_Z, \varepsilon_G) \in \operatorname{Aut}(\underline{G}^+) \times \hat{\mathbb{Z}}^* \times \mu \times \mu $$
satisfaisant les conditions

(39)
a) $u(\rho_0)$ est $\underline{G}^+$-conj. à $\rho_0^{\mu}[\tau]\rho$
b) $u(\sigma_0)$ est $\underline{G}^+$-conj. à $\sigma_0^{\mu}[\tau]\sigma$
c) $u(\varepsilon_0)$ est $\underline{G}^+$-conjugué à $\varepsilon_0$
d) $u$ stabilise $\underline{C}_S$ (centralisateur de $r_0, s_0$) et $\varepsilon_Z = \varepsilon_G$

Son noyau est isomorphe à $\underline{Z}^{0+}$

(40) $$ \underline{Z}^{0+} = \left\{ f \in \underline{G}_{\rho, \sigma}^+ \mid \underline{u} = (u, 1, 1, 1) \in \underline{Z}^{0+}, \quad f u = u f \right\} $$

[unclear: crossed out block]

Corollaire Soit 
(41) $$ Z_0 = Cent_{\underline{G}_{\rho, \sigma}^+} (L_0), \quad Z^S = Z_0 \cap \underline{G}^S \quad (= \{ f \in \underline{G}^S \mid int(f) \dots \}) $$
de sorte que $L_0^S$ s'identifie : un sous-groupe

\newpage

%% ===== Page 662 =====

central de $Z_0^\mathcal{G}$. Considérons l'application

(42) $$Z_0^\mathcal{G} \longrightarrow S.G_{\rho,\sigma} = Z_{\rho,\sigma} \cdot \mathfrak{S}^\mathcal{G}$$
$$f \longmapsto f^{-1} \cdot u(f) \in \mathfrak{S}^\mathcal{G} Z_{\rho,\sigma} \quad \text{où } u(f) = (int(f), 1, 1) \in Z_{\rho,\sigma}$$

Cette application est un hom. de groupes,
induit sur $L_0^\mathcal{G} \subset Z_0^\mathcal{G}$ l'homomorphisme canonique $i_{L_0,\mathcal{G}}$
(28), et $L_0^\mathcal{G} \subset Z_0^\mathcal{G}$ est l'image inverse de
$i_{L_0,\mathcal{G}}(L_0^\mathcal{G}) = \operatorname{Ker}(S.G_{\rho,\sigma} \longrightarrow G_{\rho,\sigma})$. On
trouve ainsi un hom. injectif

(43) $$Z_0^\mathcal{G} / L_0^\mathcal{G} \hookrightarrow G_{\rho,\sigma}$$

dont l'image est le noyau de (38), i.e. on
a une suite exacte canonique

(44) $$1 \longrightarrow L_0^\mathcal{G} \longrightarrow Z_0^\mathcal{G} \longrightarrow G_{\rho,\sigma} \longrightarrow \tilde{G}_{\rho,\sigma} \longrightarrow 1$$

où $\tilde{G}_{\rho,\sigma} \subset \operatorname{Aut}(\underline{\mathcal{E}}^+) \times \hat{\mathbb{Z}}^* \times \mu \times \mu$
est formé des $(u, \mu, \varepsilon_1, \varepsilon_2)$ satisfaisant à
les conditions (39).

Remarque. Cette proposition fait apparaître
de quelle façon la $\mathcal{G}$-construction fait
de $G_{\rho,\sigma}$ un [unclear] plus fin que celle [unclear] :
en le définissant simplement comme un
groupe d'automorph. de $\underline{\mathcal{E}}^+$ (i.e. on a besoin
des $f, \varepsilon, \varepsilon'$ par dessus le marché ... ).

\newpage

%% ===== Page 663 =====

On considérera l'hom (en fait effectivement injectif par (38), même dans des cas ou on n'a pas d'hypothèses du type $\mathfrak{S}^+ = \operatorname{Sl}(2, \hat{\mathbb{Z}})$)

(45) $$ G_{\rho,\sigma}^{\mathcal{G}} \longrightarrow \operatorname{Aut}(\mathfrak{S}^+) \times \mu \times \mu $$
$$ u \longmapsto (u | \mathfrak{S}^+, \varepsilon_1(u), \varepsilon_2(u)) $$

induit par (38), dont l'image est formée des systèmes $(u, \varepsilon_1, \varepsilon_2)$ tels que
a) $u(\rho)$ est $\mathfrak{S}^+$-conjugué à $\varepsilon_1[\Gamma](\rho)$
(46) b) $u(\sigma)$ est $\mathfrak{S}^+$-conjugué à $\varepsilon_2[\Gamma](\sigma)$
c) $u$ normalise un $\mathfrak{S}^+$-conjugué de $L_0$
d) $u$ stabilise $\mathfrak{S}^{\Delta}$ (automatique si $\varepsilon_1=\varepsilon_2$, ou [unclear: $\mathfrak{S}$ ab])

Son noyau est
(47) $$ \left\{ f \in G_{\rho,\sigma}^{\mathcal{G}} \mid y = (u, \mu, 1, 1) \in Z_{\rho,\sigma} \text{ t.q. } u = \text{int}(f) \right\} $$
$$ f \in \mathcal{G}^{\Delta} $$

Pour apprécier la structure du noyau, soit
(48) $$ SN_0^{\mathcal{G}} = \operatorname{Norm}_{\mathcal{G}}(L_0), \qquad SN_0^{\Delta} = SN_0 \cap \mathcal{G}^{\Delta} $$
$$ = \{ f \in \mathcal{G}^{\Delta} \mid \exists \mu \in \hat{\mathbb{Z}}^* \text{ t.q. } (\text{int}(f), \mu, 1, 1) \in Z_{\rho,\sigma} \} $$

et soit
(49) $$ \tilde{SN}_0 \subset SN_0 \times \hat{\mathbb{Z}}^* , \qquad \tilde{SN}_0^{\Delta} = \tilde{SN}_0 \cap (SN_0^{\Delta} \times \hat{\mathbb{Z}}^*) $$
$$ = \{ (u, \mu) \mid u(s_0) = s_0^\mu \} $$

[MARGIN: idem pour $\tilde{SN}_0$]

(50) $$ \tilde{SN}_0^{\Delta} \longrightarrow SN_0^{\Delta} \qquad \text{si } \hat{\mathbb{Z}} \xrightarrow{\sim} L_0 \text{ est injectif.} $$
$$ (u, \mu) \longmapsto u $$

On a un homomorphisme canonique (injectif)
(51) $$ \tilde{SN}_0^{\Delta} \longrightarrow S_0 G_{\rho,\sigma}^{\mathcal{G}} = Z_{\rho,\sigma} \cdot \mathcal{G}^{\Delta} $$
$$ (f, \mu) \longmapsto f^{-1} \cdot \underbrace{(\text{int}(f), \mu, 1, 1)}_{u} = u \cdot f^{-1} $$

\newpage

%% ===== Page 664 =====

qui induit sur $L_0^\mathcal{G}$ l'hom canonique $\nu^{(28)}$ (et
qui prolonge (42) ) - il est encore vrai
que l'image inverse de $L_0^\mathcal{G}$ (i.e. de $i_{L_0,\mathcal{G}}(L_0^\mathcal{G})$)
est $L_0^\mathcal{G} \subset \widetilde{SN}_0^\mathcal{G}$, et on trouve ainsi une
suite exacte canonique

(52)
\begin{tikzcd}
1 \ar[r] & L_0^\mathcal{G} \ar[r] & \widetilde{SN}_0^\mathcal{G} \ar[r] & G_{\rho,\sigma}^\mathcal{G} \ar[r] & \tilde{G}_{\rho,\sigma}^\mathcal{G} \ar[r] \ar[d, hook] & 1 \\
& & & & \operatorname{Aut}(\underline{C}^+) \times \mu \times \mu
\end{tikzcd}

[unclear: où] $\tilde{G}_{\rho,\sigma}^\mathcal{G}$ est formé des $(u,\varepsilon_2,\varepsilon_3)$ satis-
faisant (46) - i.e. ...

(52 bis) $$\widetilde{SN}_0^\mathcal{G} / L_0^\mathcal{G} \simeq \operatorname{Ker}(G_{\rho,\sigma}^\mathcal{G} \to \operatorname{Aut}(\underline{C}^+) \times \mu \times \mu)$$ .

\newpage

%% ===== Page 665 =====

$4^\circ)$ $N_\rho$, $N_\sigma$, $Z_\rho$, $Z_\sigma$, $N_\rho^S$, $N_\sigma^S$, $Z_\rho^S$, $Z_\sigma^S$

On pose
$$(53) \qquad \begin{cases} N_\rho = \{ a \in G_{\rho,\sigma} \mid a(\rho) = \rho^{\varepsilon_3}, \text{où } \varepsilon_3 = \varepsilon_3(a) \} \\ N_\sigma = \{ b \in G_{\rho,\sigma} \mid b(\sigma) = \sigma^{\varepsilon_2}, \text{où } \varepsilon_2 = \varepsilon_2(b) \} \end{cases}$$

[MARGIN: *] Il n'y a pas lieu finalement de traîner des [unclear: tildes] en donnant une définition de $N_\rho, N_\sigma$ qui donnerait des groupes plus gros. (Revoir p 580, 581 pour des interversions [unclear: idiotes] à ce sujet...)

On pose
$$(54) \qquad \begin{cases} Z_\rho = Cent_{G_{\rho,\sigma}}(\rho) = Cent_{G_{\rho,\sigma}}(U_\rho) = Ker (N_\rho \xrightarrow{\varepsilon_3} \mu) \\ Z_\sigma = Cent_{G_{\rho,\sigma}}(\sigma) = Cent_{G_{\rho,\sigma}}(U_\sigma) = Ker (N_\sigma \xrightarrow{\varepsilon_2} \mu) \end{cases}$$

[MARGIN: - on définit de façon similaire $S N_\rho, S N_\sigma$, mais on a idiotement $S N_\rho = S Z_\rho$, $S N_\sigma = S Z_\sigma$ -]
$$(55) \qquad \begin{cases} S Z_\rho = Z_\rho \cap G^+ = Ker (Z_\rho \xrightarrow{\delta_Z} \Gamma_{\rho,\sigma}) \\ \quad = Cent_{G^+}(\rho) = Cent_{G^+}(U_\rho) \\ S Z_\sigma = Z_\sigma \cap G^+ = Ker (Z_\sigma \xrightarrow{\delta_Z} \Gamma_{\rho,\sigma}) \\ \quad = Cent_{G^+}(\sigma) = Cent_{G^+}(U_\sigma) \end{cases}$$

[MARGIN: suites exactes d'[unclear] scindées par [unclear] - la deuxième scindée par $\mu_2$ ]
$$(56) \qquad \begin{matrix} 1 \longrightarrow Z_\rho \longrightarrow N_\rho \xrightarrow{\varepsilon_3} \mu \longrightarrow 1 \\ 1 \longrightarrow Z_\sigma \longrightarrow N_\sigma \xrightarrow{\varepsilon_2} \mu \longrightarrow 1 \end{matrix}$$

(le fait qu'on peut mettre des $1$ au bout provient du fait que $\tau' \in N_\rho \cap N_\sigma$ satisfait
$$(57) \qquad \varepsilon_2(\tau') = \varepsilon_3(\tau') = -1$$
[unclear: On a un] morphisme de groupes $(p 580 !)$ )

\newpage

%% ===== Page 666 =====

\begin{tikzcd}
& D_\rho \ar[rr] & & \underset{= \mu_{\tau'} \cdot S Z_\rho}{N_\rho \cap \widehat{G}} \ar[rr, hook] \ar[dr, hook] & & N_\rho \ar[dd, hook] \\
(58) \quad D_{\tau'} = \mu_{\tau'} \times \mu & & & & \widehat{G} \ar[rr, hook] & & G_{\rho,\sigma} \\
& D_\sigma \ar[rr] & & \underset{= \mu_{\tau'} \cdot S Z_\sigma}{N_\sigma \cap \widehat{G}} \ar[rr, hook] \ar[ur, hook] & & N_\sigma \ar[uu, hook]
\end{tikzcd}

ou mieux

(59)
\begin{tikzcd}[row sep=1.5em, column sep=1.5em]
& & & D_\rho \ar[rr] & & S N_\rho \ar[drr, hook] \ar[rrr, hook] & & & N_\rho \ar[drr, hook] \\
\mu_{\tau'} \ar[rr] & & D_{\tau'} \ar[ur, "2"] \ar[dr, "3"'] & & & & & \widehat{G} \ar[rrr, hook] & & & G_{\rho,\sigma} \\
& & & D_\sigma \ar[rr] & & S N_\sigma \ar[urr, hook] \ar[rrr, hook] & & & N_\sigma \ar[urr, hook] \\
& & & L_\rho \ar[uuu, hook] \ar[rr] & & S Z_\rho \ar[uuu, hook] \ar[drr, hook] \ar[rrr, hook] & & & Z_\rho \ar[uuu, hook] \ar[uurr, hook, crossing over] \ar[drr, dashed] \\
1 \ar[uuuu] \ar[rr] & & \mu \ar[uuuu, hook] \ar[ur, "2"] \ar[dr, "3"'] & & & & & G^+ \ar[uuuu, hook] \ar[uuurrr, hook, crossing over] \ar[rrr, dashed] & & & G_{\rho,\sigma} \ar[uuu, hook] \\
& & & L_\sigma \ar[uuu, hook] \ar[rr] & & S Z_\sigma \ar[uuu, hook] \ar[urr, hook] \ar[rrr, hook] & & & Z_\sigma \ar[uuu, hook] \ar[uuuurr, hook, crossing over] \ar[urr, dashed]
\end{tikzcd}

On définit

\newpage

%% ===== Page 667 =====

(62) 
$$
\left\{
\begin{tikzcd}
1 \ar[r] & SZ_\rho^{\mathcal{G}} \ar[r] \ar[d] & N_\rho^{\mathcal{G}} \ar[r, "\delta_\rho^{\mathcal{G}}"] \ar[d] & \gamma_{\rho,\sigma} \ar[r] \ar[d, equal] & 1 \\
1 \ar[r] & SZ_\rho \ar[r] & N_\rho \ar[r, "\delta_\rho"'] & \gamma_{\rho,\sigma} \ar[r] & 1 \\
1 \ar[r] & SZ_\sigma^{\mathcal{G}} \ar[r] \ar[d] & N_\sigma^{\mathcal{G}} \ar[r, "\delta_\sigma^{\mathcal{G}}"] \ar[d] & \gamma_{\rho,\sigma} \ar[r] \ar[d, equal] & 1 \\
1 \ar[r] & SZ_\sigma \ar[r] & N_\sigma \ar[r, "\delta_\sigma"'] & \gamma_{\rho,\sigma} \ar[r] & 1
\end{tikzcd}
\right.
$$

(63) 
$$
\left\{
\begin{tikzcd}
1 \ar[r] & Z_\rho^{\mathcal{G}} \ar[r] \ar[d, "\text{épi}"'] & N_\rho^{\mathcal{G}} \ar[r, "\varepsilon_{3\rho}^{\mathcal{G}}"] \ar[d, "\text{épi}"'] & \mu \ar[r] \ar[d, equal] & 1 \\
1 \ar[r] & \gamma_{\rho,\sigma}^{\prime} \ar[r] & \gamma_{\rho,\sigma} \ar[r, "\varepsilon_{3\rho}"'] & \mu \ar[r] & 1 \\
1 \ar[r] & Z_\sigma^{\mathcal{G}} \ar[r] \ar[d, "\text{épi}"'] & N_\sigma^{\mathcal{G}} \ar[r, "\varepsilon_{3\sigma}^{\mathcal{G}}"] \ar[d, "\text{épi}"'] & \mu \ar[r] \ar[d, equal] & 1 \\
1 \ar[r] & \gamma_{\rho,\sigma}^{\prime\prime} \ar[r] & \gamma_{\rho,\sigma} \ar[r, "\varepsilon_{3\sigma}"'] & \mu \ar[r] & 1
\end{tikzcd}
\right.
$$

Ces diagrammes [unclear: s'intègrent] [unclear: exact] [unclear: dans] ( cf. p. [unclear] ) .
On a de même un diagramme

(64) 
\begin{tikzcd}
& L_0^\natural = S Z_{\rho,\sigma}^C \ar[rr] \ar[dl] \ar[d] & & Z_{\rho,\sigma}^C \ar[dl] \ar[d] \ar[dr, "\varepsilon_{3\gamma}"] \\
S Z_{\rho}^C \ar[d] \ar[rr] & & N_{\rho}^C \ar[d] \ar[rr] & & \mu \ar[r] & \hat{\mathbb{Z}}^* \ar[dl, "\mu"'] \\
S Z_{\sigma}^C \ar[d] \ar[rr] & & N_{\sigma}^C \ar[d] \ar[rru] \\
1 \ar[r] & C^\natural \ar[rr] \ar[d] & & G_{\rho,\sigma} \ar[rr] \ar[d] & & \gamma_{\rho,\sigma} \ar[r] \ar[d] & 1 \\
1 \ar[r] & C^+ \ar[rr] & & G \ar[rr] & & \mu_\tau \ar[r] & 1
\end{tikzcd}

( cf p. 581 (45) ) , et diagrammes analogues pour $G$.

\newpage

%% ===== Page 668 =====

[unclear: Exemple...]

Je voudrais maintenant définir l'analogue
des sous-groupes $M(1/2)$, $M(-j)$ (resp $M^{(1)}(-j)$), $N_p^\natural$ de p. 615 ss,
ce qui impliquera qu'on ait un sous-groupe $C_{p,\sigma}^\natural$
jouant le rôle du $M_0$ - l'idéal serait que
ce soit $G_{p,\sigma}^\natural$ lui-même. Mais dans le cas
universel $\underline{C}^\natural = \underline{C}(\dots)^\wedge$, $\underline{C}_+^\natural = \Pi_0$, d'où
$\Z_{p,\sigma} = M(0)' = M(0) . L_0^\pi$, (car $L_0^\natural \simeq L_0^\pi$) on trouve $G_{p,\sigma} \simeq M$,
$G_{p,\sigma}^\natural = $ sous-groupe de $G_{p,\sigma} \simeq M$ engendré par
$\underline{C}_+^\natural = \Pi_0$ et par $\Z_{p,\sigma} = M(0) . L_0^\pi$
$= M(0) . \Pi = M^1$, et non $M(0) . \Pi_0 = M^1_+$.

L'ennui provient finalement du fait
que c'est $\Z_{p,\sigma}$ qui est trop gros, on
aimerait définir dans $\Z_{p,\sigma}$ un sous-groupe
qui joue le rôle de $M(0) . \Pi_0$. Il n'y a
pas de difficulté à [unclear] définir
dans $\Z_{p,\sigma}$ un sous-groupe qui
corresponde à $M(0)$

(65) $$Z_{p,\sigma}(0) \subset Z_{p,\sigma}$$

qui donne un scindage

(66) $$Z_{p,\sigma} \simeq Z_{p,\sigma}(0) . L_0^\pi \qquad (\text{1/2 direct})$$

et on poserait alors

(67) $$G_{p,\sigma}^\natural \simeq Z_{p,\sigma}(0) . \underline{C}_+^\natural \qquad (\text{1/2 direct})$$

\newpage

%% ===== Page 669 =====

et on définirait
$$
(68) \qquad \begin{cases}
G_{\rho,\sigma}^{\natural\natural} \overset{dfn}{=} Z_{\rho,\sigma}(0) . \underline{S}^{\natural} \subset \underset{\underset{\textstyle Z_{\rho,\sigma}(0) . \underline{S}^+}{\Vert}}{G_{\rho,\sigma}^{\natural}} \subset G_{\rho,\sigma} \\
Z_{\rho,\sigma}^{\natural} = Z_{\rho,\sigma} \\
G_{\rho,\sigma}^{\natural\natural} \cap Z_{\rho,\sigma} = Z_{\rho,\sigma}(0) . (L_0 \cap \underline{S}^{\natural})
\end{cases}
$$

Rappelons une façon standard d'obtenir un scindage de con. de $Z_{\rho,\sigma}$ sur $\Gamma_{\rho,\sigma}$, quand on dispose d'un hom.
$$
(69) \qquad \underline{S}^+ \to \operatorname{Sl}(2,\hat{\mathbb{Z}})' = \operatorname{Sl}(2,\hat{\mathbb{Z}})/\{\pm 1\}
$$
commutant à l'action de $\pi$, ce qui induit un hom.
$$
(70) \qquad Z_{\rho,\sigma} \to B_0' = \left\{ \left(\begin{matrix} \mu & \lambda \\ 0 & 1 \end{matrix}\right) \mid \mu \in \hat{\mathbb{Z}}^\times, \lambda \in \hat{\mathbb{Z}} \right\}
$$
(qui [unclear: n' l'était au] § 48)
induisant le
$$
(71) \qquad \begin{cases}
L_0^\natural \longrightarrow L_0' = \left\{ \left(\begin{matrix} 1 & \lambda \\ 0 & 1 \end{matrix}\right) \mid \lambda \in \hat{\mathbb{Z}} \right\} \\
\rho^n \longmapsto \left(\begin{matrix} 1 & n \\ 0 & 1 \end{matrix}\right)
\end{cases}
$$
plus précisément un diagramme de suites exactes

(72)
\begin{tikzcd}
1 \ar[r] & L_0^\natural \ar[r] \ar[d, "\rho^n \mapsto n"'] & Z_{\rho,\sigma} \ar[r] \ar[d] & \Gamma_{\rho,\sigma} \ar[r] \ar[d, "\chi_\sigma"'] & 1 \\
1 \ar[r] & \hat{\mathbb{Z}} \ar[r, "\lambda \mapsto \left(\begin{smallmatrix} 1 & \lambda \\ 0 & 1 \end{smallmatrix}\right)"'] & B_0' \ar[r, "\text{det}", "\left(\begin{smallmatrix} \mu & \lambda \\ 0 & 1 \end{smallmatrix}\right) \mapsto \mu"'] & \hat{\mathbb{Z}}^\times \ar[r] & 1
\end{tikzcd}

Si l'hom. (71) a un scindage, alors tout

\newpage

%% ===== Page 670 =====

scindage de la $2^{ieme}$ suite exacte s'en déduit par
image inverse de la première, [unclear: tel]
Mais dans les cas typiques qui nous [unclear: intéressent] , $L_0^\natural$ (qui n'
est pas un petit groupe comme "le [unclear]"
de $G_{\rho,\sigma}^\natural$ !) - n'a pas [unclear]
pour donner un iso de (71) ! Mais
on peut aussi faire une hypothèse
supplémentaire - savoir que l'image
inverse par (70) du "bon" [unclear: sous-groupe]

(73)
$$T_0' = \left\{ \begin{pmatrix} \mu & 0 \\ 0 & 1 \end{pmatrix} \mid \mu \in \hat{\mathbb{Z}}^\times \right\} \subset B_0'$$

qui est un sous-groupe

(74)
$$Z_{\rho,\sigma}(\infty) \subset Z_{\rho,\sigma}$$
$$= \left\{ u \in Z_{\rho,\sigma} \mid \text{son image } \begin{pmatrix} \mu & \lambda \\ 0 & 1 \end{pmatrix} \text{ est dans } T_0', \text{ i.e. } \lambda = 0 \right\}$$

satisfasse à

(75)
$$Z_{\rho,\sigma}(\infty) \cap L_0^\natural = \{ 1 \} \quad .)$$

cela en fait un [unclear: sous-groupe scindé],
i.e. que $Z_{\rho,\sigma}(\infty) \to \Upsilon_{\rho,\sigma}$ soit épi i.e. iso,
i.e. que l'on a bien

(76)
$$Z_{\rho,\sigma} = Z_{\rho,\sigma}(\infty) \cdot L_0^\natural \quad .$$

J'ai de bonnes raisons de penser que
(dès que (69) existe) cette condition est
satisfaite dans tous les cas vraiment utiles,
quitte, s'il le faut, à [unclear: au choix d'un]

\newpage

%% ===== Page 671 =====

638

$Z_{f, \sigma}(\omega)$ a la bonne propriété d'ailleurs d'être
"fonctoriel" en un sens évident (cf. plus bas
la question de fonctorialité), ce qui sera
essentiel pour avoir un bon formalisme.
Ceci nous permet donc de définir un
sous-groupe $(G_{f, \sigma}^{\natural})_{T^{\sim}} \subset G_{f, \sigma}^{\natural}$ (68). Si sa
définition n'était ainsi un peu trop lourdin-
guée, une def. "plus fine" qui serait meilleure,
et mériterait en fait la notation plus
simple $G_{f, \sigma}^{\natural}$.

Autre façon de procéder ? On cherche des
restrictions adéquates dans
la définition de $Z_{f, \sigma}$, qui [unclear: donne $L_{\sigma}^{\natural} = L_{\sigma} \cap \mathfrak{S}^{\natural}$, i.e. que]
impliquent que (pour $l \in L_{\sigma}$, on a
$$((1, l), 1, 1, 1) \in Z_{f, \sigma} \iff l \in \mathfrak{S}^{\natural}$$
Mais dans le cas universel, où on a $\mathfrak{S}^{\natural} = \Pi_0$,
prenant $l = l_0$, [unclear: l'image de] $\mathfrak{S}^{\natural} / [\mathfrak{S}^{\natural}, \mathfrak{S}^{\natural}] = \mathfrak{S} / \Pi_0$
est l'extension centrale non triviale de $\mathfrak{S} / \Pi_0$,
l'homomorphisme qu'il définit dans $\mathfrak{S}^{+} / \Pi_0$
(où on a $\mathfrak{S} / \Pi_0$) est trivial - donc on n'a
pas par cette [unclear: voie de] propriétés de cette situation
qu'on arrivera à l'exorciser !

Mais peut-être est-il prématuré d'essayer
déjà de "disséquer" nos groupes $G_{f, \sigma}^{\natural}, M_{f, \sigma}^{\natural}$

\newpage

%% ===== Page 672 =====

"aussi fin que possible" - peut-être y a-t-il
[crossed out: encore] des erreurs d'orientation assez grossières, qui
se révèleront par l'examen de cas
particuliers. Donc je vais renoncer
pour l'instant à continuer l'inventaire de
de l'émergence des polygones, et m'en
tiendrai donc là [crossed out: dessus] en suspens, quitte
à y revenir par la suite et à raffiner
les constructions présentes, qu'il faudra
cependant reprendre à l'occasion de celles
du § 48 $\underline{\overline{XI}}$ .

\newpage

%% ===== Page 673 =====

$5^{\circ}$) $S_0 M_{\rho,\sigma}^\natural, M_{\rho,\sigma}^\natural, S_0 \Gamma_{\rho,\sigma}^\natural, \Gamma_{\rho,\sigma}^\natural \dots$

On posera comme dans p. 583

(77)
$$ \begin{cases} S_0 M_{\rho,\sigma} = Z_{\rho,\sigma} \times_\gamma N_\rho^\natural \times_\gamma N_\sigma^\natural \times_\gamma C_{\rho,\sigma} \\ S_0 M_{\rho,\sigma}^\natural = Z_{\rho,\sigma} \times_\gamma N_\rho^\natural \times_\gamma N_\sigma^\natural \times_\gamma C_{\rho,\sigma}^\natural \end{cases} $$

[unclear: (78')]
$$ \begin{cases} S_0 \Gamma_{\rho,\sigma} = Z_{\rho,\sigma} \times_\gamma N_\rho^\natural \times N_\sigma^\natural \\ \Gamma_{\rho,\sigma} = N_\rho^\natural \times_\gamma N_\sigma^\natural \end{cases} $$

(78)
$$ \begin{cases} M_{\rho,\sigma} \simeq N_\rho^\natural \times N_\sigma^\natural \times C_{\rho,\sigma} \\ M_{\rho,\sigma}^\natural \simeq N_\rho^\natural \times N_\sigma^\natural \times C_{\rho,\sigma}^\natural \end{cases} $$

d'où un diagramme de suites exactes

(79)
\begin{tikzcd}[row sep=2em, column sep=1.5em]
& 1 \ar[rr] \ar[dl] & & C^{\natural,+} \ar[rr] \ar[dd, equal] \ar[dl] & & S_0 M_{\rho,\sigma}^\natural \ar[rr] \ar[dd] \ar[dl] & & S_0 \Gamma_{\rho,\sigma}^\natural \ar[rr] \ar[dd] \ar[dl] & & 1 \ar[dl] \\
1 \ar[rr] & & C^+ \ar[rr] \ar[dd, equal] & & S_0 M_{\rho,\sigma} \ar[rr] \ar[dd] & & S_0 \Gamma_{\rho,\sigma} \ar[rr] \ar[dd] & & 1 \\
& 1 \ar[rr] \ar[dl] & & C^{\natural,+} \ar[rr] \ar[dl] & & M_{\rho,\sigma}^\natural \ar[rr] \ar[dl] & & \Gamma_{\rho,\sigma}^\natural \ar[rr] \ar[dl] & & 1 \ar[dl] \\
1 \ar[rr] & & C^+ \ar[rr] & & M_{\rho,\sigma} \ar[rr] & & \Gamma_{\rho,\sigma} \ar[rr] & & 1
\end{tikzcd}

et des suites exactes

(80)
\begin{tikzcd}[row sep=1.5em]
1 \ar[r] & Z_{\rho,\sigma}^\natural \times S N_\rho^\natural \times S N_\sigma^\natural \times C_{\rho,\sigma}^\natural \ar[r] \ar[d] & S_0 M_{\rho,\sigma}^\natural \ar[r] \ar[d] & \overline{\Gamma}_{\rho,\sigma} \ar[r] & 1 \\
1 \ar[r] & Z_{\rho,\sigma} \times S N_\rho^\natural \times S N_\sigma^\natural \times C_{\rho,\sigma} \ar[r] & S_0 M_{\rho,\sigma} \ar[r] & \overline{\Gamma}_{\rho,\sigma} \ar[r] & 1
\end{tikzcd}

(82)
\begin{tikzcd}[row sep=1.5em]
1 \ar[r] & S N_\rho^\natural \times S N_\sigma^\natural \times C_{\rho,\sigma}^\natural \ar[r] \ar[d] & M_{\rho,\sigma}^\natural \ar[r] \ar[d] & \overline{\Gamma}_{\rho,\sigma} \ar[r] & 1 \\
1 \ar[r] & S N_\rho^\natural \times S N_\sigma^\natural \times C_{\rho,\sigma} \ar[r] & M_{\rho,\sigma} \ar[r] & \overline{\Gamma}_{\rho,\sigma} \ar[r] & 1
\end{tikzcd}

(83)
$$ 1 \to S N_\rho^\natural \times S N_\sigma^\natural \to S_0 \Gamma_{\rho,\sigma}^\natural \to \overline{\Gamma}_{\rho,\sigma} \to 1 $$

(84)
$$ 1 \to S N_\rho^\natural \times S N_\sigma^\natural \to \Gamma_{\rho,\sigma} \to \overline{\Gamma}_{\rho,\sigma} \to 1 $$

\newpage

%% ===== Page 674 =====

(85)
$$S_0 M_{\rho,\sigma} \simeq M_{\rho,\sigma} \times_{\gamma_{\rho,\sigma}} Z_{\rho,\sigma}$$
$$S_0 M_{\rho,\sigma}^{\mathcal{G}} \simeq M_{\rho,\sigma}^{\mathcal{G}} \times_{\gamma_{\rho,\sigma}^{\mathcal{G}}} Z_{\rho,\sigma}$$

(86)
[unclear: scribbled] $$S_0 \Gamma_{\rho,\sigma} \simeq \Gamma_{\rho,\sigma} \times_{\overline{\gamma}_{\rho,\sigma}} Z_{\rho,\sigma}$$

(87)
$$S_0 M_{\rho,\sigma} \simeq M_{\rho,\sigma} \times_{\Gamma_{\rho,\sigma}} S_0 \Gamma_{\rho,\sigma}$$
$$S_0 M_{\rho,\sigma}^{\mathcal{G}} \simeq M_{\rho,\sigma}^{\mathcal{G}} \times_{\Gamma_{\rho,\sigma}^{\mathcal{G}}} S_0 \Gamma_{\rho,\sigma}$$

Remq [unclear: Cette défilée] de groupes donne le mal de mer. Je ne suis pas sur que les variantes $S_0$ ($S_0 M$, $S_0 M^{\mathcal{G}}$, $S_0 \Gamma$) soient très utiles.
Reste la compétition entre des invariants avec ou sans $\mathcal{G}$, surtout entre les $M_{\rho,\sigma}$ et $M_{\rho,\sigma}^{\mathcal{G}}$ — dont la "différence" n'est "liée" qu'entre $\mathcal{C}^+$ et $\mathcal{C}^{\mathcal{G}}$, et plutôt $G_{\rho,\sigma}$ et $G_{\rho,\sigma}^{\mathcal{G}}$, de facon précise $M_{\rho,\sigma}$ est l'image dans $M_{\rho,\sigma}^{\mathcal{G}}$ de $G_{\rho,\sigma}$ dans $G_{\rho,\sigma}^{\mathcal{G}}$ par des isomorphismes

(88)
$$ \mathcal{C}^+/\mathcal{C}^{\mathcal{G}} \xrightarrow{\sim} G_{\rho,\sigma}/G_{\rho,\sigma}^{\mathcal{G}} \xrightarrow{\sim} M_{\rho,\sigma}/M_{\rho,\sigma}^{\mathcal{G}} . $$

Les groupes $G_{\rho,\sigma}$, $M_{\rho,\sigma}$ jouent vis à $G_{\rho,\sigma}^{\mathcal{G}}$, $M_{\rho,\sigma}^{\mathcal{G}}$ le rôle de fourre-touts commodes pour nous permettre de considérer les actions de $\rho,\sigma$ sur $G_{\rho,\sigma}^{\mathcal{G}}$, $M_{\rho,\sigma}^{\mathcal{G}}$ comme induites par des autom. intérieurs i.e. des [unclear: éléments de nos groupes]

\newpage

%% ===== Page 675 =====

640

environnants $G_{\rho,\sigma}$, $M_{\rho,\sigma}$.

On a un clivage

(89) $$ \tau'_M \in M^{\mathcal{G}}_{\rho,\sigma} $$ [MARGIN: au dessus de $\tau'_\gamma \in \Gamma_{\rho,\sigma}$]
$$ (\tau'_\rho, \tau'_\sigma, \tau'_G) $$
$$ \cap \quad \cap \quad \cap $$
$$ N^\mathcal{G}_\rho \quad N^\mathcal{G}_\sigma \quad G^\mathcal{G}_{\rho,\sigma} $$

d'où des plongements canoniques

(90) $$ \mathcal{C}^\mathcal{G} \hookrightarrow M^\mathcal{G}_{\rho,\sigma} \qquad \mathcal{C}^\mathcal{G} = \text{im. inv. dans } M^\mathcal{G}_{\rho,\sigma} \text{ de } \{1,\tau_\gamma\} \subset \Gamma_{\rho,\sigma} $$
$$ \mathcal{C} \hookrightarrow M_{\rho,\sigma} \qquad \mathcal{C} = \text{im. inv. dans } M_{\rho,\sigma} \text{ de } \{1,\tau_\gamma\} \subset \Gamma_{\rho,\sigma} $$

On peut reprendre la discussion
[unclear: boutique] que au p. 584 , avec des $\mathcal{G}$
: la def ...

Je vais voir si je peux me dispenser
justement de redévelopper les liens avec
$\tilde{Z}_{\rho,\sigma}$, $S G_{\rho,\sigma}$ etc - quitte à y revenir
au besoin par la suite.

\newpage

%% ===== Page 676 =====

6°) Fonctorialités.

[MARGIN: NB jusqu'ici je n'avais pas utilisé l'hyp. que $\phi$ surjectif vu que dans (91)]

(cf 1 ; 5°))

cf. p. 588. En plus de l'hom.
Je n'examinerai pas [unclear] cf p. 643' ss...

(91)
\begin{tikzcd}
\mathcal{C} \ar[r, "\phi"] & \mathcal{C}'
\end{tikzcd}

supposons que $\phi$ applique $\mathcal{C}^\natural$ dans $\mathcal{C}'^\natural$.

On désigne ici par $Z_{\rho,\sigma}$, $M_{\rho,\sigma}$ etc les invariants associés à $\mathcal{C}$, $Z'_{\rho',\sigma'}$ etc ceux de $\mathcal{C}'$. On a des conditions équivalentes comme dans loc. cit - où on va laisser tomber la considération de $\bar{Z}_{\rho,\sigma}$ ... qui signifiait que pour tout
$$\underline{u} = (u, u_1, \varepsilon_2, \varepsilon_3) \in Z_{\rho,\sigma}$$
il existe un autom $u'$ de $\mathcal{C}'^+$ rendant commutatif

(92)
\begin{tikzcd}
\mathcal{C}^+ \ar[r, "\phi"] \ar[d, "u"'] & \mathcal{C}'^+ \ar[d, "u'"] \\
\mathcal{C}^+ \ar[r, "\phi"'] & \mathcal{C}'^+
\end{tikzcd}

lequel $u'$ sera unique. Je veux donc = à dire que les $u \in \operatorname{Aut}(\mathcal{C}^+)$ qui proviennent de $Z_{\rho,\sigma}$, i.e. qui normalisent $L_\sigma$, et qui transforment $k_{\rho,\sigma}$ en des $k_{\rho',\sigma'}$ [unclear] d'invariants $\mathcal{C}^\natural$ conjoints à $\rho,\sigma$, stabilisent le noyau de $\phi$ (91).

On trouve alors

\newpage

%% ===== Page 677 =====

641

$$ (93) \quad Z_{\rho,\sigma} \xrightarrow{\varphi_Z} Z'_{\rho',\sigma'} $$
compatible avec les actions de $Z_{\rho,\sigma}$, $Z'_{\rho',\sigma'}$ sur $\hat{G}^+$ resp. $\hat{G}^{\prime +}$, et avec les hom.
$$ L_0^\natural \subset Z_{\rho,\sigma} \rightarrow \hat{G}^+ , \quad L_0^{\prime \natural} \subset Z'_{\rho',\sigma'} \rightarrow \hat{G}^{\prime +} , $$
d'où par la construction standard un hom.
$$ (94) \quad G_{\rho,\sigma} \xrightarrow{\varphi_G} G'_{\rho',\sigma'} $$
induisant
$$ (95) \quad G_{\rho,\sigma}^{\natural} \xrightarrow{\varphi_{G^\natural}} G_{\rho',\sigma'}^{\prime \natural} $$
d'où un hom des déguisements (64) relatif à $\underline{(\rho,\sigma)}$, dans le déguisement analogue relatif à $\underline{(\rho',\sigma')}$, avec [unclear].
Je n'en dirai pas plus pour [unclear] que toutes les constructions des nos précédents, relatives à $\rho,\sigma$, s' [unclear] dans celles relatives à $\rho',\sigma'$, $(\hat{G}')$.

Remarque. Supposons que $\hat{G}'$ satisfasse à la condition suivante : Pour tout couple d'éléments $\alpha',\beta' \in \hat{G}'_{\underline{E}}$, telles que l'on ait

[MARGIN: (96) [crossed out: (98)]]
$$ [\alpha'_p] = [\beta'_\sigma] \varepsilon_1^{\prime -1} \quad \text{pour } p \in \hat{\mathbb{Z}} \text{ cons.} $$
il existe un $u' \in \operatorname{Aut}(\hat{G}^{\prime +})$ tel que

[MARGIN: (97) [crossed out: (99)]]
$$ u'(\rho) = \alpha'(\rho), \; u'(\sigma) = \beta'(\sigma) $$
($\alpha' \in Cent(\rho)$, $\beta' \in Cent(\sigma)$)

\newpage

%% ===== Page 678 =====

la relation de base (96) s'exprimera en
termes de $u'$ par la condition

(99) $$u'(L_0^{\mathbb{S}'}) = L_0^{\mathbb{S}'}$$ i.e. $u'$ normalise $\mathbb{S}'$.

Alors le $(\rho,\sigma)$-hom $(\mathbb{S}, \mathbb{S}^+) \to (\mathbb{S}', \mathbb{S}'^+)$ est
nécessairement admissible.

Nous [unclear: verrons] que ceci est vrai au
cas universel - mais pour pouvoir en dire
tant en général, je vais d'abord
parler des sous-groupes
$$M(\infty), M(\rho), M(\sigma) \subset M_{\rho,\sigma}.$$

\newpage

%% ===== Page 679 =====

[MARGIN: Les groupes $M_{\rho,\sigma}[0]$, $M_{\rho,\sigma}(\rho)$, $M_{\rho,\sigma}(\sigma)$.]

Dans le groupe
(100)
$$S_0 M_{\rho,\sigma}^\natural = Z_{\rho,\sigma} \times_{\Gamma_{\rho,\sigma}} N_\rho^\natural \times_{\Gamma_{\rho,\sigma}} N_\sigma^\natural \times_{\Gamma_{\rho,\sigma}} G_{\rho,\sigma}^\natural$$
$$= \{ (u, a, b, U) \in Z_{\rho,\sigma} \times N_\rho^\natural \times N_\sigma^\natural \times G_{\rho,\sigma}^\natural \mid \delta_Z(u) = \delta_\rho(a) = \delta_\sigma(b) = \delta(U) \}$$

on introduit trois sous-groupes

(101)
$$S_0 M_{\rho,\sigma}[0] \subset S_0 M_{\rho,\sigma}^\natural \quad \text{par condition } U = i_Z(u)$$
$$S_0 M_{\rho,\sigma}(\rho) \subset S_0 M_{\rho,\sigma}^\natural \quad \text{par condition } U = i_\rho(a)$$
$$S_0 M_{\rho,\sigma}(\sigma) \subset S_0 M_{\rho,\sigma}^\natural \quad \text{par condition } U = i_\sigma(b)$$

en utilisant maintenant les trois hom. can.

(102)
$$i_Z : Z_{\rho,\sigma} \hookrightarrow G_{\rho,\sigma}$$
$$i_\rho : N_\rho \hookrightarrow G_{\rho,\sigma}$$
$$i_\sigma : N_\sigma \hookrightarrow G_{\rho,\sigma}$$

Ces trois groupes sont manifestement tous
trois isomorphes à : 
$$Z_{\rho,\sigma} \times_{\Gamma_{\rho,\sigma}} N_\rho^\natural \times_{\Gamma_{\rho,\sigma}} N_\sigma^\natural = S_0 \Gamma_{\rho,\sigma}$$
de façon précise, dans la suite exacte (58)
$$1 \longrightarrow C^\natural \longrightarrow S_0 M_{\rho,\sigma} \longrightarrow S_0 \Gamma_{\rho,\sigma} \longrightarrow 1$$
l'hom. canonique $S_0 M_{\rho,\sigma} \longrightarrow S_0 \Gamma_{\rho,\sigma}$
induit des isom. entre les groupes (101)
et $S_0 \Gamma_{\rho,\sigma}$, ces trois groupes sont donc des
sous-groupes sections. (de $S_0 M_{\rho,\sigma}$)

On définit $M_{\rho,\sigma}[0]$, $M_{\rho,\sigma}(\rho)$ et $M_{\rho,\sigma}(\sigma)$
(ou $M_{\rho,\sigma}(-)$) comme les images de
$S_0 M_{\rho,\sigma}[0]$, $S_0 M_{\rho,\sigma}(\rho)$ et $S_0 M_{\rho,\sigma}(\sigma)$ par l'hom.
canonique $S_0 M_{\rho,\sigma} \longrightarrow M_{\rho,\sigma}$, ce qui

\newpage

%% ===== Page 680 =====

revient à poser

(103) 
$$ \begin{cases} 
M_{\rho,\sigma}[0] = \{ (a,b,U) \in M_{\rho,\sigma}^\natural \mid U \in Z_{\rho,\sigma} \} \\
M_{\rho,\sigma}(\rho) = \{ (a,b,U) \in M_{\rho,\sigma}^\natural \mid U = a \} \\
M_{\rho,\sigma}(\sigma) = \{ (a,b,U) \in M_{\rho,\sigma}^\natural \mid U = b \}
\end{cases} $$

Donc les sous-groupes $M_{\rho,\sigma}(\rho), M_{\rho,\sigma}(\sigma)$ de $M_{\rho,\sigma}^\natural$ sont des groupes entiers au dessus de $\Gamma_{\rho,\sigma}^\natural \simeq N_\rho^\natural \times_{\Gamma_{\rho,\sigma}} N_\sigma^\natural$, donc

(104) 
$$ \begin{cases} 
M_{\rho,\sigma}^\natural \simeq M_{\rho,\sigma}(\rho) \cdot C^\natural \quad (1/2 \text{ directe}) \\
M_{\rho,\sigma}^\natural \simeq M_{\rho,\sigma}(\sigma) \cdot C^\natural \quad \quad \quad )
\end{cases} $$

et de même

(105) 
$$ \begin{cases} 
M_{\rho,\sigma} \simeq M_{\rho,\sigma}(\rho) \cdot C^+ \quad (1/2 \text{ directe}) \\
M_{\rho,\sigma} \simeq M_{\rho,\sigma}(\sigma) \cdot C^+
\end{cases} $$

D'autre part, on a

(106)
\begin{tikzcd}
1 \ar[r] & L_0^\natural \ar[r] \ar[dd, equal] & M_{\rho,\sigma}[0] \ar[r] & \Gamma_{\rho,\sigma} \ar[r] \ar[dd, equal] & 1 \\
& & S_0 M_{\rho,\sigma}[0] \ar[u, "\simeq"'] \ar[d, "\simeq"] & & \\
1 \ar[r] & L_0^\natural \ar[r] & S_0 \Gamma_{\rho,\sigma} \ar[r] & \Gamma_{\rho,\sigma} \ar[r] & 1
\end{tikzcd}

donc l'extension $M_{\rho,\sigma}[0]$ de $\Gamma_{\rho,\sigma}$ par $L_0^\natural$ s'identifie à l'ext. $S_0 \Gamma_{\rho,\sigma}$ de $\Gamma_{\rho,\sigma}$ par $L_0^\natural$, i.e. à l'image inverse de l'extension $Z_{\rho,\sigma}$ de $\Gamma_{\rho,\sigma}$ par $L_0^\natural$, via l'hom. canonique

\newpage

%% ===== Page 681 =====

643

(107) $$ \Gamma_{\rho,\sigma} = N_{\rho}^\mathcal{G} \times_{\gamma_{\rho,\sigma}} N_{\sigma}^\mathcal{G} \longrightarrow \gamma_{\rho,\sigma} \quad .$$

Donc la donnée d'une section de $Z_{\rho,\sigma}$ sur $\gamma_{\rho,\sigma}$ (cf à ce sujet p. 636! ff),
en définit une de $Sol_{\rho,\sigma}$ sur $\Gamma_{\rho,\sigma}$,
donc de $M_{\rho,\sigma}[0]$ sur $\Gamma_{\rho,\sigma}$. Quand une
telle section est fixée, on déduit un
[unclear: groupe section pur]
$$M_{\rho,\sigma}(0) \subset M_{\rho,\sigma}^\mathcal{G} \quad ,$$

donc on a

(108)
$$
\begin{cases}
M_{\rho,\sigma}[0] \simeq M_{\rho,\sigma}(0) . L_0^\mathcal{G} \\
M_{\rho,\sigma}^\mathcal{G} \simeq M_{\rho,\sigma}(0) . \widehat{G}^\mathcal{G} \\
M_{\rho,\sigma} \simeq M_{\rho,\sigma}(0) . \widehat{G}
\end{cases}
$$

\newpage

%% ===== Page 682 =====

8°) Retour sur le cas universel, et l'hom.
$$ \mathcal{M}^\natural \longrightarrow \mathcal{M}_{p,\sigma}^\natural / \mu $$

Nous supposons maintenant que $\underline{\mathfrak{S}}$ correspond au cas universel $\underline{E\ell\ell}(?,?)^{\wedge}$, avec un $\underline{\mathfrak{S}}$ quelconque. On va noter les cas universels par des indices $\circ$ : $\underline{\mathfrak{S}}_\circ, \underline{\mathbb{G}}_\circ^+, \underline{\mathbb{G}}_\circ^\natural$ etc, et $\underline{\mathfrak{S}}', \underline{\mathfrak{S}}'^\natural$ etc le déduit universel $\underline{\mathfrak{S}}, \underline{\mathfrak{S}}^\natural$ etc.

[MARGIN: - $Z_{p,\sigma}$ etc... On écrira $\underline{\mathbb{G}}^\natural_\circ$ au lieu de $\underline{\mathbb{G}}_{\circ}^{\natural}$ etc.]

Je pose

[MARGIN: NB $\mathcal{M} \supset \mathcal{M}^\natural$]

$$ (107) \quad \mathcal{M}^\natural = \left\{ u \in \mathcal{M} \left| \begin{array}{l} a) \; u(p) \text{ est } \underline{\mathfrak{S}}^\natural \text{-conjugué à } \varepsilon_3[\tau](p) \\ b) \; u(\sigma) \text{ \qquad \qquad \qquad \quad \;} \varepsilon_2[\tau](\sigma) \\ c) \; u(\varepsilon_0) \text{ \qquad \qquad \qquad \quad \;} \varepsilon_0^\mu \\ d) \; u \text{ stabilise } \mathcal{M}^\natural \end{array} \right. \right\} $$

bien sûr 
$$ \mu = \chi(u), \quad \varepsilon_2 = \varepsilon_2(u) = \varepsilon_2(\mu), \quad \varepsilon_3 = \varepsilon_3(u) = \varepsilon_3(\mu) $$

Je dis que $\mathcal{M}^\natural$ est un sous-groupe de $\mathcal{M}$. On peut le voir en reprenant les calculs de la p. 629 ff. Je préfère procéder ainsi.

\newpage

%% ===== Page 683 =====

[unclear: On a un isomorphisme] $= \mathcal{M}_{0,3} . \widehat{L}_0^\natural$

(108) $Z_{\widehat{\mathcal{G}}, \sigma} \hookrightarrow \mathcal{M}[\pi_0] \left( \simeq \mathcal{M}_{0,3}^\sim [\pi_0] = \mathcal{M}_{0,3}(\sim) . L_{0,3}^\pi \right)$

l'inclusion qui est une égalité si
$\widehat{\mathcal{G}}^\natural = \widehat{\mathcal{G}}^+$, et qui identifie $Z_{\widehat{\mathcal{G}}, \sigma}$ à
$\mathcal{M}[\pi_0] = \mathcal{M}_{0,3} . L_0^\pi$ si $\widehat{\mathcal{G}}^\natural = \pi \ (= \pi_0)$.

[MARGIN: [unclear: Cette explicitation en sens inverse des vues.]]

[ - en fait, si $\widehat{\mathcal{G}}^\natural \subset \pi_0$, alors

(109) $Z_{\widehat{\mathcal{G}}, \sigma} = \left\{ \begin{array}{l} u = u_0 . l \\ \pitchfork \\ \mathcal{M}_{0,3} \boxtimes L_0^\pi \end{array} \Bigg| \begin{array}{l} l(p) = \alpha(p) \text{ avec } \alpha \in \widehat{\mathcal{G}}^\natural \\ l(\sigma) = \beta(\sigma) \text{ avec } \beta \in \widehat{\mathcal{G}}^\natural \\ \text{et stabilisent } \widehat{\mathcal{G}}^\natural_0 \end{array} \right\}$

et on voit que $\alpha, \beta$ de $a), b)$ sont
déterminés par $u = u_0 l$ de façon unique

(dès que l'on suppose seulement
$u \in \mathcal{M}[\pi_0]$ sans plus, et $l \in L_0^\pi$)

$\alpha, \beta \in \pi_0$, bref on sait que c'est un
sous-groupe de $(\hat{\mathbb{Z}}^\times)_+ \times \hat{\mathbb{Z}}^\times \times \mu_2$ (- en fait, c'est $\operatorname{Aut}(\widehat{\mathcal{G}}^\natural_+)$)
(prop. p. 629). Considérons

[unclear: l'homomorphisme de] $Z_{\widehat{\mathcal{G}}, \sigma}$ [unclear: dans] $\operatorname{Aut}(\widehat{\mathcal{G}}^\natural_+)$

via $\mathcal{M}_{0,3}^\sim[\pi_0]$, s'envoie dans $\mathcal{M}_{0,3}^\sim[\pi_0]$

sur le sous-groupe, donc dans

$\mathcal{M}[\pi_0]^\sim \simeq \mathcal{M}_{0,3}^\sim [\pi_0]$ par un hom.

de groupes. ]

\newpage

%% ===== Page 684 =====

Considérons le système

$$ (110) \quad (Z_{\rho_0,\sigma_0}, \underline{L}_0^\zeta \hookrightarrow Z_{\rho_0,\sigma_0}, Z_{\rho_0,\sigma_0} \to \operatorname{Aut}(\underline{G}_0^+), \underline{L}_0^\zeta \hookrightarrow \underline{G}_0^+) $$

(cf "digression" plus bas 9°) qui donne

naissance à

$$ (110') \quad \underline{G}_{\rho_0,\sigma_0}, \underline{G}_0^+ \hookrightarrow \underline{G}_{\rho_0,\sigma_0}, Z_{\rho_0,\sigma_0} \hookrightarrow \underline{G}_{\rho_0,\sigma_0} $$

et considérons le système analogue à (110*)

$$ (111) \quad (M[0]^\natural, \underline{L}_0^\pi \hookrightarrow M[0]^\natural, M[0]^\natural \to \operatorname{Aut}(\underline{G}_0^+), \underline{L}_0^\pi \hookrightarrow \underline{G}_0^+) $$
[en dessous de $M[0]^\natural$ : $\underset{\sim}{M}_{0,3} . \underline{L}_0^\pi$]
[en dessous de $\underline{L}_0^\pi \hookrightarrow$ : (engendré par $\varepsilon_0$)]

qui donne naissance (loc. cit 9°)

à la situation

$$ (111') \quad (M, \underline{G}_0^+ \hookrightarrow M, M[0]^\natural \hookrightarrow M) . $$

De plus : (108) donne un hom.
de (110) dans (111), correspondant
à l'identité sur $\underline{G}_0^+$, (108) sur $Z_{\rho_0,\sigma_0}$
l'inclusion sur $\underline{L}_0^\zeta \hookrightarrow \underline{L}_0^\pi$. Il
en résulte un hom. de (110') dans
(111') , en particulier un hom.
canonique de suites exactes

\newpage

%% ===== Page 685 =====

(112)
\begin{tikzcd}
1 \ar[r] & \widehat{\mathfrak{S}}_0^+ \ar[r] \ar[d] & G_{\rho_0,\sigma_0} \ar[r] \ar[d] & \gamma_{\rho_0,\sigma_0} \ar[r] \ar[d] & 1 \\
1 \ar[r] & \widehat{\mathfrak{S}}_0^+ \ar[r] & \mathcal{M} \ar[r] & \tilde{\Gamma}_Q \ar[r] & 1
\end{tikzcd}

où 
\begin{tikzcd}
\gamma_{\rho_0,\sigma_0} \ar[r] \ar[d, "\simeq"'] & \tilde{\Gamma}_Q \ar[d, "\simeq"] \\
Z_{\rho_0,\sigma_0}/L_0^\natural & M[0]/L_0^\natural
\end{tikzcd}
est l'homo. induit par
(108), et s'identifie à :

(113) $Z_{\rho_0,\sigma_0}/L_0^\natural \hookrightarrow M[0]/L_0^\natural$

Par définition de $L_0^\natural$ comme inter-
sect. $Z_{\rho_0,\sigma_0} \cap \widehat{\mathfrak{S}}_0^+$ ds [unclear: $\mathcal{M}$], il résulte que
(113) est inj. donc $G_{\rho_0,\sigma_0} \to \mathcal{M}$
est toujours injectif (c'est un iso. si
$\widehat{\mathfrak{S}}_0^+ \supset \Pi_0$ ...). On déduit par
composition de (112) un hom.
de suites exactes

(114)
\begin{tikzcd}
1 \ar[r] & \widehat{\mathfrak{S}}_0^\natural \ar[r] \ar[d, hook] & G_{\rho_0,\sigma_0}^\natural \ar[r] \ar[d, hook] & \gamma_{\rho_0,\sigma_0} \ar[r] \ar[d] & 1 \\
1 \ar[r] & \widehat{\mathfrak{S}}_0^+ \ar[r] & \mathcal{M} \ar[r] & \tilde{\Gamma}_Q \ar[r] & 1
\end{tikzcd}

On voit alors immédiatement que
l'im. de $G_{\rho_0,\sigma_0}^\natural$ ds $\mathcal{M}$ n'est
autre que $\mathcal{M}^\natural$ (107). On a donc

\newpage

%% ===== Page 686 =====

[MARGIN: à vérifier (115) le relèvement $u_0$ de $u$ ($u_0 \in M(\sigma)$) satisfait les conditions a) b) c) d) de (107)]
[MARGIN: NB (où on a $\gamma_{\rho_0,\sigma_0} = \tilde{\rho}_0 \tilde{\sigma}_0$) (entièrement déterminé par $\tilde{\rho}_Q$) $\gamma_0 = \tilde{\rho}_Q \times \text{l'élément } \tilde{\sigma}_Q$]

$$M^\natural \simeq G_{\rho_0,\sigma_0}^\natural$$

définies par (107)

$$\gamma_{\rho_0,\sigma_0} \hookrightarrow \tilde{\Gamma}_Q \simeq \Gamma_{\tilde{Q}}$$

qui montre en même temps dans la foulée que $M^\natural$ est bien un sous-groupe de $M$.

On a maintenant

(117)
$$M_{\rho_0,\sigma_0}^\natural \simeq M_{\rho_0}^\natural \times_{\gamma_{\rho_0,\sigma_0}} N_{\sigma_0}^\natural \times_{\gamma_{\rho_0,\sigma_0}} G_{\rho_0,\sigma_0}^\natural$$

où
$$G_{\rho_0,\sigma_0}^\natural \simeq M^\natural$$
$$N_{\rho_0}^\natural = \{ u \in M^\natural \mid u(\rho) = \rho^{\varepsilon_1(u)} \}$$
$$N_{\sigma_0}^\natural = \{ u \in M^\natural \mid u(\sigma) = \sigma^{\varepsilon_2(u)} \}$$

et finalement

(118)
$$M_{\rho_0,\sigma_0}^\natural \simeq N_{M^\natural}(\rho) \times_{M^\natural} N_{M^\natural}(\sigma) \times_{M^\natural} M^\natural$$
$$... \simeq \gamma_{\rho_0,\sigma_0} \subset \tilde{\Gamma}_Q \text{ explicité dans (116)}$$

d'où [unclear: un homomorphisme] de suites exactes

(119)
$$1 \to SZ_{\rho_0} \times SZ_{\sigma_0} \to M_{\rho_0,\sigma_0}^\natural \to M^\natural \to 1$$

[MARGIN: au-dessus de $SZ_{\rho_0} \times SZ_{\sigma_0}$ il y a $L_\rho \times L_\sigma$]

où le noyau $SZ_{\rho_0} \times SZ_{\sigma_0}$ est un groupe très petit, [unclear: s.v.n.]

\newpage

%% ===== Page 687 =====

686

$$ (120) \quad \operatorname{Ker}(M_{\rho_0, \sigma_0}^\natural \xrightarrow{\psi} M^\natural) \simeq S\Z_{\rho_0} \times S\Z_{\sigma_0} \simeq L_{\rho} \times L_{\sigma} \simeq \Z/l\Z \times \Z/m\Z $$

Nous allons améliorer encore le [unclear: dévissage]
en découpant dans $M_{\rho_0, \sigma_0}^\natural$ un groupe
encore plus petit, qui soit une ext.
de $M^\natural$ par $\mu$. [crossed out: Pour ceci on définit]

[crossed out: $N_{\rho_0}^\natural = N_{\rho}^\natural \cap M^\natural = M_{...}$]

Notons que

$$ (121) \quad M_{\rho_0, \sigma_0}^\natural = N_{\rho_0}^\natural \times_{\Gamma_0^\natural} N_{\sigma_0}^\natural \times_{\Gamma_0^\natural} M^\natural $$
$$ \cap $$
$$ N_{\rho_0} \times_{\Gamma_0} N_{\sigma_0} \times_{\Gamma_0} M \qquad \hookrightarrow \Gamma_0 = \tilde{\Gamma}_{\Q} $$

où ce dernier produit fibré n'est
autre que la quantité $M_{\rho_0, \sigma_0}^{(\natural ')}$ qui
correspondrait à un [unclear: $M_{\rho_0, \sigma_0}^\natural$] de
$\mathbb{G}' = \mathbb{G}_0^+$ tout entier (noté $M_{\rho, \sigma}$
à la p. 625 et suivantes). Dans
ces pages j'ai défini un sous-groupe,
qui était noté $M_{\rho, \sigma}^\natural$ à la p. 626,
et que je vais plutôt noter maintenant
$M_{\rho_0, \sigma_0}^{\natural '}$, défini comme

\newpage

%% ===== Page 688 =====

[unclear: crossed out text]

(122) $\mathcal{M}^? = \mathcal{N}_{\rho_0}^? \times_{\Gamma_0} Z_0^! \times_{\Gamma_0} \mathcal{M}$

où, on rappelle

(123) 
$\mathcal{N}_{\rho_0} = \operatorname{Norm}_{\mathcal{M}}(L_{\rho_0})$, $\mathcal{N}_{\sigma_0} = \operatorname{Norm}_{\mathcal{M}}(L_{\sigma_0})$
[unclear: crossed out] $Z_0 = Cent_{\mathcal{M}}(L_{\sigma_0}) (= Cent_{\mathcal{N}_{\sigma_0}}(L_{\sigma_0}))$
$Z_0^! = Z_0 \cap \mathcal{M}^! \ (= \mathcal{N}_{\sigma_0} \cap \mathcal{M}^!) \subset \mathcal{N}_{\sigma_0}$
$\mathcal{N}_{\rho_0}^? = Z_{\rho_0}^! \cdot \mu_{\tau} \subset \mathcal{N}_{\rho_0}$, où $Z_{\rho_0}^! = Z_{\rho_0} \cap \mathcal{M}^!$

de sorte qu'on a bien

[MARGIN: indépendant des choix d'un $\underline{\tau}$ ]

$\mathcal{M}^? \subset \mathcal{N}_{\rho_0} \times_{\Gamma_0} \mathcal{N}_{\sigma_0} \times_{\Gamma_0} \mathcal{M}$.

Ceci rappelé, on va poser

(124) $\mathcal{M}_{\rho_0, \sigma_0}^{\natural ?} \subset \mathcal{M}_{\rho_0, \sigma_0}^{\natural}$
$||$
$\mathcal{M}_{\rho_0, \sigma_0}^{\natural} \cap \mathcal{M}^?$ (intersection dans $\mathcal{N}_{\rho_0} \times_{\Gamma_0} \mathcal{N}_{\sigma_0} \times_{\Gamma_0} \mathcal{M}$)

[unclear: crossed out text]
et considérer l'hom. induit par

(119)

(125)
[DIAGRAM: Commutative square]
$\mathcal{M}_{\rho_0, \sigma_0}^{\natural ?} \longrightarrow \mathcal{M}^{\natural}$
$\downarrow \ \ \ \ \ \ \ \ \ \ \ \ \downarrow$
$\mathcal{M}^? \longrightarrow \mathcal{M}$

\newpage

%% ===== Page 689 =====

647

On a vu ($p. 626'_x (55)$) que $\mathcal{M}^? \to \mathcal{M}$
fait de $\mathcal{M}^?$ une extension de $\mathcal{M}$ de noyau
$\mathcal{H}_{\rho}$. Donc le noyau de $\mathcal{M}_{\rho_0,\sigma_0}^{\natural ?} \to \mathcal{M}^{\natural}$
est égal à $\mathcal{H}_{\rho} \cap \mathcal{M}_{\rho_0,\sigma_0}^{\natural ?} =$ or il est clair que
d'après (119), l'on a $\mathcal{H}_{\rho} \subset \mathcal{M}_{\rho_0,\sigma_0}^{\natural}$, donc $\mathcal{H}_{\rho} \subset \mathcal{M}_{\rho_0,\sigma_0}^{\natural ?}$.
On trouve donc un [unclear] des
suites exactes

[DIAGRAM: commutative diagram of exact sequences]
$$ 1 \to \mathcal{H}_{\rho} \to \mathcal{M}_{\rho_0,\sigma_0}^{\natural ?} \to \mathcal{M}^{\natural} $$
(126)   $$ || \hspace{3cm} \downarrow \hspace{2cm} \downarrow $$
$$ 1 \to \mathcal{H}_{\rho} \to \mathcal{M}^? \to \mathcal{M} \to 1 $$

et il reste à voir si la [unclear: flèche]

(ou $\mathcal{M}_{\rho_0,\sigma_0}^{\natural ?} \to \mathcal{M}^{\natural}$)
est un épimorphisme, et sinon, caractériser
son image. Notons que
la définition (124) donne [unclear]
compte tenu des expressions (118) pour
$\mathcal{M}_{\rho_0,\sigma_0}^{\natural}$ et (122) pour $\mathcal{M}^?$ :

[MARGIN: $\gamma_0 \subset \tilde{\Gamma}_Q$ \\ $||$ \\ $I_{\omega}(\mathcal{M}^\natural)$]

(127) $\mathcal{M}_{\rho_0,\sigma_0}^{\natural ?} = \mathcal{N}_{\rho_0}^{\natural ?} \times_{\mathcal{Z}_{\sigma_0}^{\natural ?}} \mathcal{Z}_{\sigma_0}^{\natural !} \times_{\text{[unclear: scribbled]}} \mathcal{M}^{\natural}$

où l'on a posé

(128) $\mathcal{N}_{\rho_0}^{\natural ?} = \mathcal{N}_{\rho_0}^? \cap \mathcal{N}_{\rho_0}^{\natural} = \mathcal{N}_{\rho_0}^? \cap \mathcal{M}^{\natural}$
$\mathcal{Z}_{\sigma_0}^{\natural !} = \mathcal{Z}_{\sigma_0}^! \cap \mathcal{N}_{\sigma_0}^{\natural} = \mathcal{Z}_{\sigma_0}^! \cap \mathcal{M}^{\natural}$

\newpage

%% ===== Page 690 =====

et l'hom (*) n'est autre que la
troisième projection du produit fibré.
Donc l'image de (*) n'est autre que l'
image inverse, dans $M^\natural$ par $M^\natural \to \Gamma_0 = \Gamma_{f_0,\sigma_0}$
de l'image de
$$(**) \quad N_{f_0}^{\natural ?} \times_{\Gamma_0^\natural} N_{\sigma_0}^{\natural !} \to \Gamma_0^\natural (\subset \Gamma_0)$$

Il reste donc à montrer que cet hom.
est surjectif, ou encore que les
hom.
$$N_{f_0}^{\natural ?} \to \Gamma_0^\natural$$
$$N_{\sigma_0}^{\natural !} \to \Gamma_0^\natural$$
sont surjectifs (le premier sera
alors une ext. de $\Gamma_0^\natural$ par $\mu_f$, le
deuxième un is.). Soit donc
$u \in \Gamma_0^\natural$, défini par un $u_{\mathcal{M}} \in \mathcal{M}_{f_0,\sigma_0}^\natural \simeq G_{f_0,\sigma_0}^\natural$
satisfaisant aux conditions a) b) c) d)
de (107). On voit donc qu'il existe
$\beta \in \underline{S}_{M^\natural}$ tel que $\beta^{-1} u_0 \in N(\sigma_0)$.

Ici j'ai arrêté les vérifications, pensant y revenir
dans les jours suivants, il s'avère qu'en fait c'est
pas urgent - il importait plutôt de dégager la ou
les idées "forces" ou "principes" pour l'étude
[unclear: structurale] du [unclear: syst.] $M_{\rho_0,\sigma_0}$ ....

\newpage

%% ===== Page 691 =====

698

9°) Digression de formalisme des groupes.

Soient
(1) $$G, \underline{G} \subset G, N \xrightarrow{\varphi_N} G, d'où SN = \varphi_N^{-1}(\underline{G}) \subset N$$
la situation d'un groupe $G$, sous-groupe
invariant $\underline{G}$, et un hom. de gr. $N$ tel que
(1') $$N \xrightarrow{epi} G/\underline{G} \text{ i.e. } N/SN \simeq G/\underline{G}$$
On veut reconstituer cette situation
à partir de la suivante
(2) $$\underline{G}, N, SN \subset N, N \xrightarrow{\theta} \operatorname{Aut}(\underline{G}),$$
$$SN \xrightarrow{\varphi = \varphi_{SN}} \underline{G}$$
où $\underline{G}, N$ sont des groupes, $SN$ un
sous-groupe distingué de $N$, $N \to \operatorname{Aut}(\underline{G})$
une action de groupe i.e. un hom. de
$N$ sur $\underline{G}$, enfin $SN \xrightarrow{\varphi} \underline{G}$ un hom.
de gr. [unclear] :
a) l'action de $SN$ sur $\underline{G}$ induite par celle de $N$,
(3) est celle déduite de $SN \xrightarrow{\varphi} \underline{G}$,
via l'op. adjoint de $\underline{G}$
b) $SN \xrightarrow{\varphi} \underline{G}$ est compatible avec actions de $N$

[unclear: On voit] comment (1) définit (2),
[unclear: et montrons] comment un syst.
(2) permet de reconstituer un syst. (1).

\newpage

%% ===== Page 692 =====

On part donc de $\mathfrak{G}$, ... on a
(3) $S_0 G = N . \mathfrak{G}$ (prod. s. direct)

on définit un hom.
(4) $$S N \overset{i}{\longrightarrow} S_0 G$$
$$u \longmapsto u . \varphi(u)^{-1}$$

Vérifions que c'est un hom.

[MARGIN: 
NB on a, pour $u \in SN$, $x \in \mathfrak{G}$
(par hyp (a))
$(*)\ u(x) = u x u^{-1} = \varphi(u) x \varphi(u)^{-1}$
dans $S_0 G = N . \mathfrak{G}$
donc si $x = \varphi(u)^{-1}$, ...
$u(x) = u x u^{-1} = x$
i.e. $x \varphi(u)^{-1}$ et $u$ commutent]

[unclear crossed out line] $i(u u') = (u' u) \varphi(u' u)^{-1}$
$i(u) i(u') = u \varphi(u)^{-1} u' \varphi(u')^{-1} = u u' (u'^{-1} \varphi(u)^{-1} u') \varphi(u')^{-1}$
$= u u' \varphi(u)^{-1} \varphi(u')^{-1}$
$= u u' ( \dots ) = (u u') \varphi(u u')^{-1}$ ok.

On n'a pas encore utilisé le fait que $SN$
est distingué dans $N$, et l'hyp b) que
l'image de (4) est distinguée dans $S_0 G$,
car si $a \in N$, $u \in SN$
$a (u . \varphi(u)^{-1}) a^{-1} = a u a^{-1} . a \varphi(u)^{-1} a^{-1}$
$\in SN \quad = \varphi(a u^{-1} a^{-1}) = \varphi(a u a^{-1})^{-1}$ ok.
i.e. l'image de (4) commute à $N$
dans $SN$ et dans $S_0 G \supset N$ -- donc l'image
de $SN$ est centralisée par $N \subset S_0 G$. Il
faut encore vérifier qu'elle est normalisée
également par $\mathfrak{G}$, soit $x \in \mathfrak{G}$, on a :
$x (u \varphi(u)^{-1}) x^{-1} = (x u x^{-1}) (x \varphi(u)^{-1}) x^{-1}$
L'image de $SN \xrightarrow{i} S_0 G$ commute à $\mathfrak{G}$, i.e.

\newpage

%% ===== Page 693 =====

$\forall u \in SN$, $u \varphi(u)^{-1}$ induit l'opération triviale sur
$\tilde{G}$ - ce qui est en effet trivial par (3).

On peut donc poser
(6) $$G \stackrel{df}{=} S_0 G / i(SN)$$
on trouve alors un diagramme
commutatif
(7)
[DIAGRAM: commutative diamond with vertices $N$ (top), $G$ (right), $\tilde{G}$ (bottom), $SN$ (left). Arrows are $SN \hookrightarrow N$, $N \xrightarrow{\varphi_N} G$, $SN \xrightarrow{\varphi_{SN}} \tilde{G}$, $\tilde{G} \to G$]

on voit que l'image de $\tilde{G}$ dans $G$ est
un ss-groupe invariant, et que l'on
a un isom. sur les quotients :
(8) $$N/SN \stackrel{\sim}{\longrightarrow} G/\tilde{G}$$ ,
ce qui implique que l'on a un iso
(9) $$\operatorname{Ker}(\varphi_{SN} : SN \to \tilde{G}) \stackrel{\sim}{\longrightarrow} \operatorname{Ker}(N \stackrel{\varphi_N}{\longrightarrow} G)$$ ;

donc que
(10) $\varphi_N : N \to G$ est un pl-g-s
$\iff$
$\varphi_{SN} : SN \to \tilde{G}$ est un pl-g-s.

On trouve ainsi des équiv. de catégories
[MARGIN: et cat. $\mathcal{F}$]
[unclear: prouvant l'une du texte, entre...]

\newpage

%% ===== Page 694 =====

[unclear: situations] (4), et celle des situations (1).

(b) Considérons maintenant une situation

$$ (11) \left\{ \begin{array}{l} G, \mathfrak{G} \subset G, (\varphi_i : N_i \to G)_{i \in I} \\ d'où \text{ des } SN_i = \varphi_i^{-1}(\mathfrak{G}) \end{array} \right. $$

avec la condition

$$ (12) \quad \forall i \in I, \varphi_i : N_i \to G/\mathfrak{G} \text{ est épi} $$

pour - la reconstituer à partir de

$$ (13) \quad \mathfrak{G}, (N_i, SN_i, \varphi_i, \Phi_i)_{i \in I} $$

où $(N_i, SN_i, \varphi_i, \Phi_i)$ satisfait les conditions

(3) a) et b) - - bien sûr $\forall i \in I$

$$ (14) \quad \Phi_i : N_i \to \operatorname{Aut}(\mathfrak{G}), \varphi_i : SN_i \to \mathfrak{G}. $$

Il est tentant de reconstituer
$G$ à l'aide des $\mathfrak{G}$ et des $N_i$,
comme une sorte de "somme amalgamée"
relative à :

[DIAGRAM: $SN_i$ pointant vers $\mathfrak{G}$ via $\varphi_{i}$ et vers $N_i$ via un plongement. $\mathfrak{G}$ et $N_i$ pointant vers $G$ via des flèches pointillées.]

693

\newpage

%% ===== Page 695 =====

[crossed out: Proposition Partons de (11) d'où (13)]

Mais il est plus vraisemblable de dire,

que pour tout $i$, on reconstitue via

$N_i, SN_i, \varphi_i, \Theta_i$ la situation (pour $i$ fixé)

(15) [DIAGRAM: carré commutatif avec flèches horizontales $N_i \to G_i (\subset G)$ et $SN_i \to \mathfrak{G}$, et inclusions verticales $SN_i \hookrightarrow N_i$ et $\mathfrak{G} \hookrightarrow G_i$]

et qu'on reconstitue (11) lui-même

en se donnant un syst. transitif d'isom.

entre les $G_i$ ($i \in I$), induisant

l'identité sur $\mathfrak{G}$. Une telle

donnée [unclear: implique qu'on se donne] la

donnée d'un syst. transitif d'iso.

entre les

(16) $$ \Gamma_i \overset{def}{=} N_i / SN_i $$

- qui s'identifie donc à un

= groupe

(17) $\Gamma$, groupe quotient commun des $N_i$ (par les $SN_i$).

Identifions maintenant les $G_i / \mathfrak{G} = \Gamma_i$

à $\Gamma$

(18) $$ G_i / \mathfrak{G} \simeq \Gamma , $$

et formons le produit fibré

\newpage

%% ===== Page 696 =====

[MARGIN: $g$ est extension de $\gamma$ par $\mathfrak{G}^I$]

(19)
$$g = \prod_{i \in I} G_i \supset \mathfrak{G}^I$$
(produit fibré sur $\gamma$)

La donnée d'un syst. transitif d'iso.
entre les $G_i$, induisant l'identité sur les
sous-groupes $\mathfrak{G}$, et compatible avec le
système transitif d'iso. déjà envisagé
entre les $\gamma_i = N_i / SN_i \simeq G_i / \mathfrak{G}$, équivaut
alors à la donnée d'un sous-groupe
$G$ de $g$, satisfaisant les conditions

(20)
a) $G \cap \mathfrak{G}^I = \text{diagonale } \delta(\mathfrak{G}) \text{ dans } \mathfrak{G}^I$
b) La projection $G \to \gamma$ est épi

En d'autres termes, on a construit,
à l'aide de (13) et d'un syst. transi-
tif d'iso. entre les $\gamma_i$, une extension
$g$ de leur "valeur commune" $\gamma$,
par $\mathfrak{G}^I$ - et la structure supplé-
mentaire, correspondant à la donnée
(11), est la "restriction" de cette
extension en une extension de $\gamma$
par le sous-groupe diagonal de $\mathfrak{G}^I$.
